\documentclass[UTF8,12pt,a4paper,fontset=fandol]{ctexart}

% 支持从项目根目录、高等代数目录或当前习题集目录读取共享样式。
\makeatletter
\def\input@path{{../}{../../}{../../../}{../../../../}}
\makeatother
\usepackage{researchnotes}
\usepackage{arydshln,mathdots} % 矩阵分块虚线与反向对角省略号。
\newtheorem{exercise}{习题}[section]

\title{高等代数习题集}
\author{}
\date{}

\begin{document}
\maketitle
\section*{编选说明}
\addcontentsline{toc}{section}{编选说明}
本书面向修读本科高等代数或线性代数课程的读者，围绕向量空间、线性映射、矩阵、多项式和内积组织练习。正文共 493 道题，其中 397 道选自 yhwu-is 及项目贡献者的《线性代数：未竟之美》，96 道为编者补充的计算题；附编另收录方而静（szdytom）的 154 道概念与证明题解，全书合计 647 道题，每题紧接参考解答。正文适合随课程按主题训练，附编适合集中练习定义、反例与证明。

来源标记 L 后依次为原讲号、分组和组内题号，“例”表示原讲义正文中的例题顺序号；例如 L10-B-3 表示原讲义第 10 讲 B 组第 3 题。S 表示本选编另行构造的计算题，其数据并非网上原题。P 后依次为附编原书小节及题号，例如 P-3C-2 表示原书 3C 第 2 题。来源编号用于查找出处，不是本书习题编号；每题的来源链接保留原定位。文中“原讲义”指《线性代数：未竟之美》，“原书”指 Axler 第四版教材。

本次整理保留所有题目，将原来插入的 PDF 附编改为可编辑正文，并修正已发现的题设、证明、记号与计算问题。全书的内容和来源说明均在本文件内，除项目根目录统一使用的 \texttt{researchnotes.sty} 及常规宏包外，不依赖本目录中的 PDF、图片或其他数据文件。原材料中的辅助示意图不再单独附入，相应论证以公式和文字完整表达。

本书是学习用选编，未经原作者逐题审定。此次完成题答配对、环境与括号、标签与引用、文件依赖的源码检查，并对修订条目及补充计算题进行针对性核验；未对所有来源题解作逐题独立审定。修订信息和来源版本记录保留在文末及源码注释中。

\subsection*{学习路线与使用方法}
初学时可先完成第~\ref{sec:topic1} 至第~\ref{sec:topic5} 节，掌握方程求解、子空间、基与维数、核与像，再进入第~\ref{sec:topic7} 至第~\ref{sec:topic11} 节，把抽象线性映射与矩阵计算联系起来。第~\ref{sec:topic6} 节的对偶空间可在熟悉线性泛函后单独学习；第~\ref{sec:topic13} 节的多项式是后续特征理论、对角化与若当标准形的代数基础。最后学习内积、正交化、二次型及应用。

做计算题时先明确矩阵尺寸、基的顺序和参数例外值，再用代回、秩、行列式或分解恒等式检验结果；做证明题时先列出假设，分别处理零向量、零空间和退化阶数等边界情况。遇到“证明或证伪”，可先在低维空间中寻找反例，再决定是否证明一般命题。参考解答展示一种可核查的路径，不要求所有题目采用同一种方法。

\subsection*{来源与许可}
正文原讲义项目为 \url{https://github.com/yhwu-is/Linear-Algebra-Left-Undone}，作者为 yhwu-is 及项目贡献者，采用知识共享署名—非商业性使用—相同方式共享 4.0 国际许可协议（CC BY-NC-SA 4.0）。本选编对该材料的编排、修订及补充沿用此许可；保留署名不表示原作者认可或审定本选编。材料按现状提供。许可全文见 \url{https://creativecommons.org/licenses/by-nc-sa/4.0/}，本地保存的许可文本同时收入本文件的注释。

附编原题作者为 Sheldon Axler，题干中文表述参考吴俊达、何阳的第四版中文译本，参考解答作者为方而静。原题与译本发布页为 \url{https://linear.axler.net/}，采用 CC BY-NC 4.0；题解项目为 \url{https://github.com/szdytom/LADRSolutions}，采用 CC BY-NC-SA 4.0。原题、译文和题解保留各自署名及许可，不因合并文件而改变。原材料检索与下载日期为 2026 年 10 月 9 日，本次整理使用此前保存的源码。

\subsection*{统一记号与必要约定}
正文未另行指定时，标量域 $\mathbf F$ 为数域，即 $\mathbf C$ 的子域；涉及大小比较、正定性和欧氏空间时取实数域。$r(A)$ 表示秩（rank），$E$、$E_n$ 或 $I$ 表示相应阶数的单位矩阵，$O$ 表示零矩阵。$\operatorname{span}(v_1,\ldots,v_k)$ 表示线性张成空间（span），$\ker T$、$\operatorname{im} T$ 分别表示线性映射 $T$ 的核（kernel）与像（image），$\mathcal L(V,W)$ 表示从 $V$ 到 $W$ 的线性映射空间。

正文沿用 $\mathbf F[x]_n$ 表示次数严格小于 $n$ 的多项式连同零多项式，维数为 $n$；附编沿用原书的 $\mathcal P_m(\mathbf F)$ 表示次数不超过 $m$ 的多项式连同零多项式，维数为 $m+1$。因此 $\mathcal P_m(\mathbf F)=\mathbf F[x]_{m+1}$，两种下标约定不能混用。附编中的 $\mathbf F$ 限定为 $\mathbf R$ 或 $\mathbf C$，$V^S$ 表示从集合 $S$ 到 $V$ 的全部函数构成的空间，$\mathbf F^\infty$ 表示标量序列空间。

矩阵的伴随矩阵记为 $\operatorname{adj}(A)$（adjugate），对偶空间记为 $V^*=\mathcal L(V,\mathbf F)$（dual space），内积空间中的伴随算子记为 $T^*$（adjoint）；这三种对象必须区分。若 $A=(a_{ij})$，则大写 $A_{ij}$ 在行列式题中表示 $a_{ij}$ 的代数余子式，故 $\operatorname{adj}(A)=(A_{ji})$。复内积对第一个变量线性、对第二个变量共轭线性。有限维维数公式、可对角化判据及内积表示结论应在题设条件下使用；涉及无限维空间的基与递归选取时，沿用通常的选择公理约定。
\clearpage
\tableofcontents
\clearpage
\section{预备知识与方程求解}\label{sec:topic1}
本节收录 3 道题。从关系的判别和高斯消元入手，区分等价类、齐次解空间与非齐次解集；求通解时核对自由变量的个数。
\begin{exercise}\label{prob:L1-A-2}
判断下列关系是否为等价关系；若是，描述等价类并写出商集。
\begin{enumerate}
\item 在 $\mathbf R$ 上，$xRy$ 当且仅当 $x-y$ 是无理数或为零。
\item 在 $S=[-1,1]^2$ 上，$(x_1,y_1)R(x_2,y_2)$ 当且仅当 $x_1=x_2$。
\item 在 $S=\mathbf R^2$ 上，$(x_1,y_1)R(x_2,y_2)$ 当且仅当 $x_1^2+y_1^2=x_2^2+y_2^2$。
\item 在平面上所有直线的集合 $L$ 上，分别定义 $l_1R_1l_2\iff l_1\perp l_2$，以及 $l_1R_2l_2\iff l_1,l_2$ 恰有一个公共点。
\item 在 $S=\{(x,y)\in\mathbf R^2:xy\ne0\}$ 上，$(x_1,y_1)R(x_2,y_2)$ 当且仅当 $x_1x_2>0$ 且 $y_1y_2>0$。
\item 在 $S=\{a+b\mathrm i:a,b\in\mathbf R,\ a\ne0\}$ 上，$(a_1+b_1\mathrm i)R(a_2+b_2\mathrm i)$ 当且仅当 $a_1a_2>0$。
\item 设 $A=\{a_1,a_2,a_3,a_4\}$ 的四个元素互异。在幂集 $\mathcal P(A)$ 上，$A_1RA_2\iff |A_1|=|A_2|$。
\item 在 $S=\mathbf N_{>0}\times\mathbf N_{>0}$ 上，$(m_1,n_1)R(m_2,n_2)\iff m_1n_2=m_2n_1$。
\end{enumerate}
\end{exercise}
\noindent 来源：\href{https://github.com/yhwu-is/Linear-Algebra-Left-Undone/blob/new/%E8%AE%B2%E4%B9%89/%E4%B8%93%E9%A2%98/1%20%E9%A2%84%E5%A4%87%E7%9F%A5%E8%AF%86.tex}{L1-A-2}。
\begin{solve}
等价关系须同时满足自反性、对称性和传递性。
\begin{enumerate}
\item 不是。取 $x=0$、$y=\sqrt2$、$z=1$，则 $xRy$、$yRz$ 成立，但 $xRz$ 不成立，故缺少传递性。
\item 是。它按第一坐标相等分类，商集为 $S/R=\{\{(c,y):-1\le y\le1\}:c\in[-1,1]\}$。
\item 是。它按到原点的距离分类，商集为 $S/R=\{C_r:r\ge0\}$，其中 $C_r=\{(x,y):x^2+y^2=r^2\}$；$C_0=\{(0,0)\}$ 是退化情形。
\item 均不是。直线不垂直于自身，也不与自身恰有一个公共点，所以两种关系都不满足自反性。
\item 是。等价类恰为四个开象限，即 $Q_{\varepsilon,\delta}=\{(x,y):\varepsilon x>0,\delta y>0\}$（$\varepsilon,\delta\in\{1,-1\}$）；商集由这四个集合组成。
\item 是。等价类为 $H_+=\{z:\operatorname{Re}z>0\}$ 与 $H_-=\{z:\operatorname{Re}z<0\}$，故 $S/R=\{H_+,H_-\}$。
\item 是。商集为 $\{S_0,S_1,S_2,S_3,S_4\}$，其中 $S_k=\{B\subseteq A:|B|=k\}$。注意 $S_0=\{\varnothing\}$。
\item 是。因分母为正，题中条件等价于 $m_1/n_1=m_2/n_2$，故商集为 $\{C_q:q\in\mathbf Q_{>0}\}$，其中 $C_q=\{(m,n)\in S:m/n=q\}$。
\end{enumerate}
第 (2)、(3)、(5)、(6)、(7)、(8) 项都可归结为某个函数取值相等，因相等关系具有上述三条性质，故确为等价关系。
\end{solve}

\begin{exercise}\label{prob:L1-A-3}
求齐次线性方程组$\begin{cases}
                x_1+x_2+x_3+4x_4-3x_5=0   \\
                2x_1+x_2+3x_3+5x_4-5x_5=0 \\
                x_1-x_2+3x_3-2x_4-x_5=0   \\
                3x_1+x_2+5x_3+6x_4-7x_5=0
            \end{cases}$的通解.
\end{exercise}
\noindent 来源：\href{https://github.com/yhwu-is/Linear-Algebra-Left-Undone/blob/new/%E8%AE%B2%E4%B9%89/%E4%B8%93%E9%A2%98/1%20%E9%A2%84%E5%A4%87%E7%9F%A5%E8%AF%86.tex}{L1-A-3}。
\begin{solve}
\begin{align*}
                A ={} &
                \begin{pmatrix}
                    1 & 1  & 1 & 4  & -3 \\
                    2 & 1  & 3 & 5  & -5 \\
                    1 & -1 & 3 & -2 & -1 \\
                    3 & 1  & 5 & 6  & -7
                \end{pmatrix}
                \rightarrow
                \begin{pmatrix}
                    1 & 1  & 1 & 4  & -3 \\
                    0 & -1 & 1 & -3 & 1  \\
                    0 & -2 & 2 & -6 & 2  \\
                    0 & -2 & 2 & -6 & 2
                \end{pmatrix} \rightarrow \\
                      &
                \begin{pmatrix}
                    1 & 1 & 1  & 4 & -3 \\
                    0 & 1 & -1 & 3 & -1 \\
                    0 & 0 & 0  & 0 & 0  \\
                    0 & 0 & 0  & 0 & 0
                \end{pmatrix}
                \rightarrow
                \begin{pmatrix}
                    1 & 0 & 2  & 1 & -2 \\
                    0 & 1 & -1 & 3 & -1 \\
                    0 & 0 & 0  & 0 & 0  \\
                    0 & 0 & 0  & 0 & 0
                \end{pmatrix}.
            \end{align*}
            令 $ x_3 = k_1, x_4 = k_2, x_5 = k_3$，有 $x_1 = -2k_1 - k_2 + 2k_3,\enspace\allowbreak x_2 = k_1 - 3k_2 + k_3 $，则
            \[ X = (x_1, x_2, x_3, x_4, x_5)^\mathrm{T} = k_1 \begin{pmatrix} -2 \\ 1 \\ 1 \\ 0 \\ 0 \end{pmatrix} + k_2 \begin{pmatrix} -1 \\ -3 \\ 0 \\ 1 \\ 0 \end{pmatrix} + k_3 \begin{pmatrix} 2 \\ 1 \\ 0 \\ 0 \\ 1 \end{pmatrix} \qquad k_1, k_2, k_3 \in \mathbf{R} \]
\end{solve}

\begin{exercise}\label{prob:L1-A-4}
求非齐次线性方程组$\begin{cases}
                x_1-x_2+2x_3-2x_4+3x_5=1     \\
                2x_1-x_2+5x_3-9x_4+8x_5=-1   \\
                3x_1-2x_2+7x_3-11x_4+11x_5=0 \\
                x_1-x_2+x_3-x_4+3x_5=3
            \end{cases}$的通解.
\end{exercise}
\noindent 来源：\href{https://github.com/yhwu-is/Linear-Algebra-Left-Undone/blob/new/%E8%AE%B2%E4%B9%89/%E4%B8%93%E9%A2%98/1%20%E9%A2%84%E5%A4%87%E7%9F%A5%E8%AF%86.tex}{L1-A-4}。
\begin{solve}
\begin{align*}
                \bar{A} ={} &
                \begin{pmatrix}
                    1 & -1 & 2 & -2  & 3  & 1  \\
                    2 & -1 & 5 & -9  & 8  & -1 \\
                    3 & -2 & 7 & -11 & 11 & 0  \\
                    1 & -1 & 1 & -1  & 3  & 3
                \end{pmatrix}
                \rightarrow
                \begin{pmatrix}
                    1 & -1 & 2  & -2 & 3 & 1  \\
                    0 & 1  & 1  & -5 & 2 & -3 \\
                    0 & 1  & 1  & -5 & 2 & -3 \\
                    0 & 0  & -1 & 1  & 0 & 2
                \end{pmatrix} \rightarrow \\
                            &
                \begin{pmatrix}
                    1 & -1 & 2  & -2 & 3 & 1  \\
                    0 & 1  & 1  & -5 & 2 & -3 \\
                    0 & 0  & -1 & 1  & 0 & 2  \\
                    0 & 0  & 0  & 0  & 0 & 0
                \end{pmatrix}
                \rightarrow
                \begin{pmatrix}
                    1 & 0 & 0 & -4 & 5 & 4  \\
                    0 & 1 & 0 & -4 & 2 & -1 \\
                    0 & 0 & 1 & -1 & 0 & -2 \\
                    0 & 0 & 0 & 0  & 0 & 0
                \end{pmatrix}.
            \end{align*}
            令 $ x_4 = k_1, x_5 = k_2 $，有 $x_1 = 4k_1 - 5k_2 + 4,\enspace\allowbreak x_2 = 4k_1 - 2k_2 - 1,\enspace\allowbreak x_3 = k_1 -2 $，则
            \[ X = (x_1, x_2, x_3, x_4, x_5)^\mathrm{T} = k_1 \begin{pmatrix} 4 \\ 4 \\ 1 \\ 1 \\ 0 \end{pmatrix} + k_2 \begin{pmatrix} -5 \\ -2 \\ 0 \\ 0 \\ 1 \end{pmatrix} + \begin{pmatrix} 4 \\ -1 \\ -2 \\ 0 \\ 0 \end{pmatrix} \qquad k_1, k_2 \in \mathbf{R} \]
\end{solve}

\section{线性空间与公理检验}\label{sec:topic2}
本节收录 5 道题。先检验零元、运算封闭性和分配律，再使用子空间判别法；同一集合能否成为向量空间取决于数域及给定运算。
\begin{exercise}\label{prob:L2-A-1}
检验下列集合对指定的加法和数乘运算是否构成实数域上的线性空间.
        \begin{enumerate}
            \item 有理数集$\mathbf{Q}$对普通的数的加法和乘法；

            \item 集合$\mathbf{R}^2$对通常的向量加法和如下定义的数量乘法：$\lambda\cdot(x,y)=(\lambda x,y)$；

            \item $\mathbf{R}_+^n$（即$n$元正实数向量）对如下定义的加法和数乘运算：
                  \begin{gather*}
                      (a_1,\ldots,a_n)+(b_1,\ldots,b_n)=(a_1b_1,\ldots,a_nb_n) \\
                      \lambda\cdot(a_1,\ldots,a_n)=(a_1^\lambda,\ldots,a_n^\lambda)
                  \end{gather*}

            \item 集合 $V$ 为区间 $[a, b]$ 上所有函数值 $\geqslant 0$ 的实变量函数，即：
                  \[V=\{f \mid f(x) \geqslant 0, \forall x \in [a, b] \}\]
                  对通常的函数加法和数与函数的乘法，即：
                  \begin{gather*}
                      (f \oplus g)(x) = f(x) + g(x) \\
                      (\lambda \circ f)(x) = \lambda f(x)
                  \end{gather*}

            \item \[V_1= \{f \mid x \in \mathbf{R}, f(x) \in \mathbf{R}, f(-x)=-f(x)\}\]
                  \[V_2= \{f \mid x \in \mathbf{R}, f(x) \in \mathbf{R}, f(0)=1, f(-x)=f(x)\}\]
                  对题(4)所定义的加法和数量乘法.

            \item $V = \{f \mid x \in \mathbf{R}, f(x) \in \mathbf{C}, f(-x)= \overline{f(x)} \}$.
                  对题(4)所定义的加法和数量乘法.
        \end{enumerate}
\end{exercise}
\noindent 来源：\href{https://github.com/yhwu-is/Linear-Algebra-Left-Undone/blob/new/%E8%AE%B2%E4%B9%89/%E4%B8%93%E9%A2%98/2%20%E7%BA%BF%E6%80%A7%E7%A9%BA%E9%97%B4.tex}{L2-A-1}。
\begin{solve}
\begin{enumerate}
\item 不是。实数标量 $\sqrt2$ 与有理数 $1$ 的乘积不在 $\mathbf Q$ 中，数乘不封闭。
\item 不是。例如对 $v=(0,1)$，题定义给出 $0\cdot v=(0,1)\ne(0,0)$，违反向量空间中零标量乘向量必为零向量的性质。
\item 是。映射 $\Phi\colon\mathbf R_+^n\to\mathbf R^n$，$\Phi(a_1,\ldots,a_n)=(\log a_1,\ldots,\log a_n)$ 为双射，且将题定义的加法、数乘分别变为通常的向量加法、数乘。因此全部公理由 $\mathbf R^n$ 继承；零向量为 $(1,\ldots,1)$，$(a_1,\ldots,a_n)$ 的加法逆元为 $(a_1^{-1},\ldots,a_n^{-1})$。
\item 不是。常值函数 $f=1$ 属于 $V$，而 $(-1)f=-1$ 不属于 $V$，故数乘不封闭。仅凭 $\lambda f\le0$ 不能对所有 $f$ 得出不属于 $V$，因为零函数是例外。
\item $V_1$ 是：零函数属于它，奇函数的线性组合仍为奇函数，所以它是全体实值函数空间的子空间。$V_2$ 不是：其中函数必须满足 $f(0)=1$，故不含零函数。
\item 是。它包含零函数；对 $f,g\in V$ 及 $a,b\in\mathbf R$，有
\[(af+bg)(-x)=a\overline{f(x)}+b\overline{g(x)}=\overline{(af+bg)(x)},\]
故对实线性组合封闭，是全体复值函数构成的实向量空间的子空间。这里实数域不可随意换成复数域，例如 $f=1$ 符合条件，而 $\mathrm if$ 不符合条件。
\end{enumerate}
\end{solve}

\begin{exercise}\label{prob:L2-A-2}
判断下列子集是否为给定线性空间的子空间：
        \begin{enumerate}
            \item $W = \{(x_1,\ldots,x_n) \in F^n \mid a_1 x_1+\cdots +a_n x_n =0\}$, 其中 $a_1, \ldots, a_n$ 为域 $F$ 中的固定数量.

            \item $W_1 = \{(x,1,0) \in \mathbf{R}^3 \}$, $W_2 = \{(x,y,0) \in \mathbf{R}^3\}$.

            \item $W_1 = \{(x,y,z) \in \mathbf{R}^3 \mid x-3y+z = 0\}$, $W_2 = \{(x,y,z) \in \mathbf{R}^3 \mid x-3y+z = 1\}$.

            \item $W_1 = \left\{(x,y,z) \in \mathbf{R}^3 \middle|\ \dfrac{x}{2} = \dfrac{y-4}{1} = \dfrac{z-1}{-3}\right\}$,

                  $W_2 = \{(x,y,z) \in \mathbf{R}^3 \mid x-y = 0, x+y+z = 0\}$.

            \item $W_1 = \{p(x) \in R[x] \mid p(1) = 0\}$, $W_2 = \{p(x) \in \mathbf R[x]_n \mid p(1) = p(0)\}$（此题主要就是要判断满足一定条件的多项式是否构成子空间）.

            \item $W = \{f \in F(-\infty, +\infty) \mid f(-x)=f(x), \forall x \in R\}$, 其中 $F(-\infty, +\infty)$ 是所有定义在 $(-\infty, +\infty)$ 上的实值函数对通常的函数加法及数与函数的乘法在实数域上构成的线性空间.
        \end{enumerate}
\end{exercise}
\noindent 来源：\href{https://github.com/yhwu-is/Linear-Algebra-Left-Undone/blob/new/%E8%AE%B2%E4%B9%89/%E4%B8%93%E9%A2%98/2%20%E7%BA%BF%E6%80%A7%E7%A9%BA%E9%97%B4.tex}{L2-A-2}。
\begin{solve}
\begin{enumerate}
                \item 对于集合 \( W = \{(x_1, \ldots, x_n) \in F^n \mid a_1x_1 + \cdots + a_nx_n = 0\} \)：

                首先，我们来判断该集合是否为线性空间的子空间. 判断一个集合是否为线性空间的子空间需要满足以下三个条件：

                \begin{enumerate}
                    \item \text{零向量在其中}：

                    当 \( x_1 = x_2 = \cdots = x_n = 0 \) 时，方程 \( a_1x_1 + \cdots + a_nx_n = 0 \) 显然成立，因此零向量 \( (0, 0, \ldots, 0) \) 属于 \( W \).

                    \item \text{封闭性（加法）}：

                    假设 \( \mathbf{v} = (x_1, \ldots, x_n) \) 和 \( \mathbf{w} = (y_1, \ldots, y_n) \) 属于 \( W \)，即：
                    \[
                    a_1x_1 + \cdots + a_nx_n = 0 \quad \text{且} \quad a_1y_1 + \cdots + a_ny_n = 0.
                    \]
                    那么，对于 \( \mathbf{v} + \mathbf{w} = (x_1 + y_1, \ldots, x_n + y_n) \)，有：
                    \[
                    a_1(x_1 + y_1) + \cdots + a_n(x_n + y_n) = (a_1x_1 + \cdots + a_nx_n) + (a_1y_1 + \cdots + a_ny_n) = 0 + 0 = 0.
                    \]
                    因此，\( \mathbf{v} + \mathbf{w} \in W \).

                    \item \text{封闭性（数乘）}：

                    对于任意 \( \mathbf{v} = (x_1, \ldots, x_n) \in W \) 和任意 \( \lambda \in F \)，有：
                    \[
                    a_1(\lambda x_1) + \cdots + a_n(\lambda x_n) = \lambda (a_1x_1 + \cdots + a_nx_n) = \lambda \cdot 0 = 0.
                    \]
                    因此，\( \lambda \mathbf{v} \in W \).
                \end{enumerate}

                由于上述三个条件均满足，集合 \( W \) 是 \( F^n \) 的一个子空间.

                \item 对于集合 \( W_1 = \{(x, 1, 0) \in \mathbb{R}^3\} \) 和 \( W_2 = \{(x, y, 0) \in \mathbb{R}^3\} \)：

                \begin{enumerate}
                    \item \text{对于 \( W_1 \)}：

                    \begin{itemize}
                        \item \text{零向量在其中}：零向量 \( (0, 0, 0) \) 不在 \( W_1 \) 中，因为 \( W_1 \) 中的第二个分量始终为 1，因此零向量不在 \( W_1 \) 中.
                        \item 因为零向量不在 \( W_1 \) 中，所以 \( W_1 \) 不是子空间.
                    \end{itemize}

                    \item \text{对于 \( W_2 \)}：

                    \begin{itemize}
                        \item \text{零向量在其中}：当 \( x = 0, y = 0 \) 时，向量 \( (0, 0, 0) \in W_2 \).
                        \item \text{封闭性（加法）}：假设 \( \mathbf{v} = (x_1, y_1, 0) \in W_2 \) 和 \( \mathbf{w} = (x_2, y_2, 0) \in W_2 \)，则：
                        \[
                        \mathbf{v} + \mathbf{w} = (x_1 + x_2, y_1 + y_2, 0) \in W_2.
                        \]
                        \item \text{封闭性（数乘）}：对于 \( \lambda \in \mathbb{R} \) 和 \( \mathbf{v} = (x, y, 0) \in W_2 \)，有：
                        \[
                        \lambda \mathbf{v} = (\lambda x, \lambda y, 0) \in W_2.
                        \]
                    \end{itemize}
                \end{enumerate}

                因此，\( W_2 \) 是 \( \mathbb{R}^3 \) 的一个子空间.

                \item 对于集合 $W_1 = \{(x, y, z) \in \mathbb{R}^3 \mid x - 3y + z = 0\}$ 和 $W_2 = \{(x, y, z) \in \mathbb{R}^3 \mid x - 3y + z = 1\}$：
                \begin{enumerate}
                    \item \text{零向量在其中}：当 $x = 0, y = 0, z = 0$ 时，方程 $x - 3y + z = 0$ 显然成立，因此零向量 $(0, 0, 0) \in W_1$.
                    \item \text{封闭性（加法）}：假设 $\mathbf{v} = (x_1, y_1, z_1) \in W_1$ 和 $\mathbf{w} = (x_2, y_2, z_2) \in W_1$，则：
                    \[
                    (x_1 - 3y_1 + z_1) = 0 \quad \text{且} \quad (x_2 - 3y_2 + z_2) = 0.
                    \]
                    对于 $\mathbf{v} + \mathbf{w} = (x_1 + x_2, y_1 + y_2, z_1 + z_2)$，有：
                    \[
                    (x_1 + x_2) - 3(y_1 + y_2) + (z_1 + z_2) = (x_1 - 3y_1 + z_1) + (x_2 - 3y_2 + z_2) = 0 + 0 = 0.
                    \]
                    因此，$\mathbf{v} + \mathbf{w} \in W_1$.
                    \item \text{封闭性（数乘）}：对于任意 $\lambda \in \mathbb{R}$ 和 $\mathbf{v} = (x, y, z) \in W_1$，有：
                    \[
                    \lambda x - 3(\lambda y) + \lambda z = \lambda(x - 3y + z) = \lambda \cdot 0 = 0.
                    \]
                    因此，$\lambda \mathbf{v} \in W_1$.
                \end{enumerate}
                综上所述，$W_1$ 是 $\mathbb{R}^3$ 的一个子空间.
                \begin{enumerate}
                    \item \text{零向量不在其中}：若 $(0, 0, 0) \in W_2$，则必须满足 $0 - 3 \cdot 0 + 0 = 1$，这显然不成立，因此零向量不在 $W_2$ 中.
                \end{enumerate}
                因为零向量不在 $W_2$ 中，所以 $W_2$ 不是子空间.

                \item 对于集合 $W_1 = \left\{(x, y, z) \in \mathbb{R}^3 \mid \frac{x}{2} = \frac{y - 4}{1} = \frac{z - 1}{-3} \right\}$ 和 $W_2 = \{(x, y, z) \in \mathbb{R}^3 \mid x - y = 0, x + y + z = 0\}$：

                这个集合的方程表示一个直线的参数方程. 通常直线不是子空间，因为它不通过原点. 因此，$W_1$ 不是子空间.

                \begin{enumerate}
                    \item \text{零向量在其中}：当 $x = 0, y = 0, z = 0$ 时，方程 $x - y = 0$ 和 $x + y + z = 0$ 都成立，因此零向量 $(0, 0, 0) \in W_2$.
                    \item \text{封闭性（加法）}：假设 $\mathbf{v} = (x_1, y_1, z_1) \in W_2$ 和 $\mathbf{w} = (x_2, y_2, z_2) \in W_2$，即：
                    \[
                    x_1 - y_1 = 0, \quad x_1 + y_1 + z_1 = 0
                    \]
                    \[
                    x_2 - y_2 = 0, \quad x_2 + y_2 + z_2 = 0
                    \]
                    对于 $\mathbf{v} + \mathbf{w} = (x_1 + x_2, y_1 + y_2, z_1 + z_2)$，有：
                    \[
                    (x_1 + x_2) - (y_1 + y_2) = (x_1 - y_1) + (x_2 - y_2) = 0 + 0 = 0,
                    \]
                    \[
                    (x_1 + x_2) + (y_1 + y_2) + (z_1 + z_2) = (x_1 + y_1 + z_1) + (x_2 + y_2 + z_2) = 0 + 0 = 0.
                    \]
                    因此，$\mathbf{v} + \mathbf{w} \in W_2$.
                    \item \text{封闭性（数乘）}：对于任意 $\lambda \in \mathbb{R}$ 和 $\mathbf{v} = (x, y, z) \in W_2$，有：
                    \[
                    \lambda x - \lambda y = \lambda (x - y) = \lambda \cdot 0 = 0,
                    \]
                    \[
                    \lambda x + \lambda y + \lambda z = \lambda (x + y + z) = \lambda \cdot 0 = 0.
                    \]
                    因此，$\lambda \mathbf{v} \in W_2$.
                \end{enumerate}
                综上所述，$W_2$ 是 $\mathbb{R}^3$ 的一个子空间.

                \item 对于集合 $W_1 = \{p(x) \in \mathbb{R}[x] \mid p(1) = 0\}$ 和 $W_2 = \{P(x) \in \mathbb{R}[x]_n \mid p(1) = p(0)\}$：

                \begin{enumerate}
                    \item \text{零多项式在其中}：零多项式 $p(x) = 0$ 满足 $p(1) = 0$，因此零多项式 $p(x) = 0 \in W_1$.
                    \item \text{封闭性（加法）}：假设 $p(x), q(x) \in W_1$，即 $p(1) = 0$ 且 $q(1) = 0$，那么对于 $p(x) + q(x)$，有：
                    \[
                    (p(x) + q(x))(1) = p(1) + q(1) = 0 + 0 = 0,
                    \]
                    因此 $p(x) + q(x) \in W_1$.
                    \item \text{封闭性（数乘）}：对于任意实数 $\lambda$ 和 $p(x) \in W_1$，有：
                    \[
                    (\lambda p(x))(1) = \lambda p(1) = \lambda \cdot 0 = 0,
                    \]
                    因此 $\lambda p(x) \in W_1$.
                \end{enumerate}
                综上所述，$W_1$ 是一个子空间.

                \begin{enumerate}
                    \item \text{零多项式在其中}：零多项式 $p(x) = 0$ 显然满足 $p(1) = p(0)$，因此零多项式 $p(x) = 0 \in W_2$.
                    \item \text{封闭性（加法）}：假设 $p(x), q(x) \in W_2$，即 $p(1) = p(0)$ 且 $q(1) = q(0)$，那么对于 $p(x) + q(x)$，有：
                    \[
                    (p(x) + q(x))(1) = p(1) + q(1), \quad (p(x) + q(x))(0) = p(0) + q(0).
                    \]
                    因为 $p(1) = p(0)$ 和 $q(1) = q(0)$，所以 $p(1) + q(1) = p(0) + q(0)$，即 $p(x) + q(x) \in W_2$.
                    \item \text{封闭性（数乘）}：对于任意实数 $\lambda$ 和 $p(x) \in W_2$，有：
                    \[
                    (\lambda p(x))(1) = \lambda p(1), \quad (\lambda p(x))(0) = \lambda p(0).
                    \]
                    因为 $p(1) = p(0)$，所以 $\lambda p(1) = \lambda p(0)$，即 $\lambda p(x) \in W_2$.
                \end{enumerate}
                综上所述，$W_2$ 是一个子空间.

                \item 对于集合 $W = \{f \in F(-\infty, +\infty) \mid f(-x) = f(x), \forall x \in \mathbb{R}\}$：

                \begin{enumerate}
                    \item \text{零函数在其中}：零函数 $f(x) = 0$ 显然满足 $f(-x) = f(x)$，因此零函数 $f(x) = 0 \in W$.
                    \item \text{封闭性（加法）}：假设 $f(x), g(x) \in W$，即 $f(-x) = f(x)$ 且 $g(-x) = g(x)$，那么对于 $f(x) + g(x)$，有：
                    \[
                    (f + g)(-x) = f(-x) + g(-x) = f(x) + g(x) = (f + g)(x),
                    \]
                    因此 $f(x) + g(x) \in W$.
                    \item \text{封闭性（数乘）}：对于任意实数 $\lambda$ 和 $f(x) \in W$，有：
                    \[
                    (\lambda f)(-x) = \lambda f(-x) = \lambda f(x) = (\lambda f)(x),
                    \]
                    因此 $\lambda f(x) \in W$.
                \end{enumerate}
                综上所述，$W$ 是一个子空间.

            \end{enumerate}
\end{solve}

\begin{exercise}\label{prob:L2-B-3}
设$V$是一个线性空间，$W$是$V$的子集，证明：$W$是$V$的子空间$\iff \operatorname{span}(W)=W$.
\end{exercise}
\noindent 来源：\href{https://github.com/yhwu-is/Linear-Algebra-Left-Undone/blob/new/%E8%AE%B2%E4%B9%89/%E4%B8%93%E9%A2%98/2%20%E7%BA%BF%E6%80%A7%E7%A9%BA%E9%97%B4.tex}{L2-B-3}。
\begin{solve}
$W$ 是 $V$ 的子空间等价于 $\forall \alpha_1,\ldots,\alpha_k\in W$，$\forall \lambda_1,\ldots,\lambda_k\in \mathbf{F}$（$\mathbf{F}$ 是 $V$ 对应数域）有 $\lambda_1\alpha_1+\cdots+\lambda_k\alpha_k\in W$. 根据线性扩张的定义，以上描述等价于 $\operatorname{span}(W)\subseteq W$，又 $\operatorname{span}(W)$ 是 $W$ 的线性扩张，即 $W\subseteq\operatorname{span}(W)$，故$\operatorname{span}(W)\subseteq W \iff \operatorname{span}(W)=W$. 综上 $W$ 是 $V$ 的子空间，得证.
\end{solve}

\begin{exercise}\label{prob:L2-C-1}
设$\mathbf{E}$是域$\mathbf{F}$的一个子域.
        \begin{enumerate}
            \item 证明：$\mathbf{F}$关于自身的加法和乘法构成一个$\mathbf{E}$上的向量空间，并举一例；

            \item 举例说明：$\mathbf{E}\enspace(\mathbf{E}\neq \mathbf{F})$不是$\mathbf{F}$上的线性空间；

            \item 证明：若$V$是$\mathbf{F}$上的一个线性空间，则$V$也是$\mathbf{E}$上的一个线性空间.
        \end{enumerate}
\end{exercise}
\noindent 来源：\href{https://github.com/yhwu-is/Linear-Algebra-Left-Undone/blob/new/%E8%AE%B2%E4%B9%89/%E4%B8%93%E9%A2%98/2%20%E7%BA%BF%E6%80%A7%E7%A9%BA%E9%97%B4.tex}{L2-C-1}。
\begin{solve}
\begin{enumerate}
\item $\mathbf F$ 的加法构成交换群，且对 $a\in\mathbf E$、$x\in\mathbf F$ 有 $ax\in\mathbf F$。数乘结合律、两条分配律和 $1x=x$ 均由域的运算律继承，所以 $\mathbf F$ 是 $\mathbf E$ 上的向量空间。例如 $\mathbf C$ 是 $\mathbf R$ 上的二维向量空间，基为 $1,\mathrm i$。
\item 按域中原有乘法定义数乘时，$\mathbf R$ 不构成 $\mathbf C$ 上的向量空间，因为 $\mathrm i\cdot1\notin\mathbf R$。一般地，取 $a\in\mathbf F\setminus\mathbf E$，则 $a\cdot1\notin\mathbf E$，所以真子域对这种数乘不封闭。
\item 保留 $V$ 的加法，将数乘中的标量限制为 $\mathbf E$ 的元素。数乘仍封闭，所有向量空间公理都是原来在 $\mathbf F$ 上成立的恒等式，限制标量后仍成立。这一操作称为限制标量。\end{enumerate}
\end{solve}

\begin{exercise}\label{prob:L2-例-6}
回答下列关于子空间的判定问题：
    \begin{enumerate}
        \item \label{item:2:常见子空间:1}
              说明$\mathbf{R}[x]_2$是$\mathbf{R}[x]_3$的子空间；

        \item \label{item:2:常见子空间:2}
              判断$W_1=\left\{(x,y,z) \;\middle|\; \dfrac{x}{3}=\dfrac{y}{2}=z\right\},\enspace W_2=\{(x,y,z) \mid x+y+z=1,\enspace x-y+z=1\}$是否为$\mathbf{R}^3$的子空间；

        \item \label{item:2:常见子空间:3}
              试说明齐次线性方程组$AX=0$的解集是线性空间$\mathbf{F}^n$的一个子空间（我们称之为齐次线性方程组的\emph{解空间}），但非齐次线性方程组的解不再构成线性空间.
    \end{enumerate}
\end{exercise}
\noindent 来源：\href{https://github.com/yhwu-is/Linear-Algebra-Left-Undone/blob/new/%E8%AE%B2%E4%B9%89/%E4%B8%93%E9%A2%98/2%20%E7%BA%BF%E6%80%A7%E7%A9%BA%E9%97%B4.tex}{L2-例-6}。
\begin{solve}
\begin{enumerate}
        \item 只需证明$\mathbf{R}[x]_2 \subseteq \mathbf{R}[x]_3$，以及$\mathbf{R}[x]_2$对$\mathbf{R}[x]_3$中的加法和数乘封闭即可.

              $\forall v \in \mathbf{R}[x]_2$，可被写作$v=a+bx,a,b \in \mathbf{R}$. 又有$\mathbf{R}[x]_3=\{a+bx+cx^2,a,b,c \in \mathbf{R}\}$，取$c=0$，有$v=a+bx \in \mathbf{R}[x]_3$，因此$\mathbf{R}[x]_2 \subseteq \mathbf{R}[x]_3$.

              对于$\mathbf{R}[x]_3$中的加法和数乘：
              \[mv_1+nv_2=m(a_1+b_1x)+n(a_2+b_2x)=(ma_1+na_2)+(mb_1+nb_2)x \in \mathbf{R}[x]_3\]
              所以$\mathbf{R}[x]_2$是$\mathbf{R}[x]_3$的子空间.

        \item 对于此类给出条件求解子空间的问题，实际上很容易理解，因为$W_1$实际上就是所有满足方程组
              \[\begin{cases}
                      x-3z=0 \\
                      y-2z=0
                  \end{cases}\]
              的向量，而$W_2$实际上就是所有满足方程组
              \[\begin{cases}
                      x+y+z=1 \\
                      x-y+z=1
                  \end{cases}\]
              的向量. 因此我们只需要验证这两个方程组的解集是否对加法和数乘封闭即可.

              对 $W_1$: 引入参数$t$，
              \[W_1=\left\{(3t,2t,t) \;\middle|\; \frac{x}{3} = \frac{y}{2} = z = t\right\}\]
              对于$\forall v_1, v_2 \in W_1, v_1 = (3t_1, 2t_1, t_1), v_2 = (3t_2, 2t_2, t_2)$，有
              \begin{align*}
                  av_1 + bv_2 & = (3at_1 + 3bt_2, 2at_1 + 2bt_2, at_1 + bt_2)           \\
                              & = (3(at_1 + bt_2), 2(at_1 + bt_2), at_1 + bt_2) \in W_1
              \end{align*}
              故$W_1$封闭，是 $\mathbf{R}^3$ 的子空间.

              对 $W_2$: 有反例. 取 $u_1 = (1, 0, 0), u_2 = (0, 0, 1) \in W_2$，但 $W_1 + W_2 = (1, 0, 1)$ 不满足 $x + y + z = 1$，故 $W_2$ 不封闭，不是 $\mathbf{R}^3$ 的子空间.

        \item 设齐次线性方程组 $AX=0$ 的解构成的集合是 $W_1$，$\forall X_1, X_2 \in W_1$，有 $AX_1 = AX_2 = 0$，所以 $\forall a, b \in \mathbf{F}$，
              \[A(a X_1 + b X_2) = A(a X_1) + A(b X_2) = a AX_1 + b AX_2 = 0\]
              故 $W_1$ 封闭，是 $\mathbf{F}^n$ 的子空间.

              设非齐次线性方程组 $AX = \beta,\enspace \beta \in \mathbf{F}^m,\enspace \beta \neq 0$ 的解构成的集合是 $W_2$，$\forall X_1, X_2 \in W_2$，有 $AX_1 = AX_2 = \beta$，所以 $A(X_1 + X_2) = AX_1 + AX_2 = 2\beta \neq \beta$. 故 $W_2$ 不封闭，不是 $\mathbf{F}^n$ 的子空间.
    \end{enumerate}
\end{solve}

\section{线性无关、基与维数}\label{sec:topic3}
本节收录 26 道题。把线性关系、张成空间与坐标表示联系起来。用矩阵判别向量组时，明确向量按行还是按列排列，并由主元数判断维数。
\begin{exercise}\label{prob:L3-A-1}
下列命题是否正确？若正确请证明，否则举出反例.
        \begin{enumerate}
            \item 若 $\alpha_1, \ldots, \alpha_m \ (m > 2)$ 线性相关，则其中每一向量都是其余向量的线性组合；

            \item 若 $\alpha_1, \ldots, \alpha_m$ 线性无关，则其中每一向量不是其余向量的线性组合，这个命题的等价命题应如何叙述？

            \item $\alpha_1, \ldots, \alpha_m \ (m > 2)$ 线性无关的充要条件是任意两个向量都线性无关；

            \item  若 $\alpha_1, \alpha_2$ 线性相关，$ \beta_1, \beta_2$ 线性相关，则 $\alpha_1 + \beta_1, \alpha_2 + \beta_2$ 也线性相关；

            \item 若 $\alpha_1, \ldots, \alpha_n$ 线性无关，则 $\alpha_1 + \alpha_2, \alpha_2 + \alpha_3, \ldots, \alpha_{n-1} + \alpha_n, \alpha_n + \alpha_1$ 也线性无关；

            \item 若 $\alpha_1, \alpha_2, \alpha_3$ 线性相关，则 $\alpha_1 + \alpha_2, \alpha_2 + \alpha_3, \alpha_3 + \alpha_1$ 也线性相关；

            \item 设 $B = \{\alpha_1, \alpha_2, \alpha_3\}$ 是 $ \mathbf{R}^3 $ 的一组基，非零向量 $\alpha_0 \in \mathbf{R}^3$，则 $\{\alpha_0 + \alpha_1, \alpha_0 + \alpha_2, \alpha_0 + \alpha_3\}$ (其中三个向量均是非零向量) 也是 $\mathbf{R}^3$ 的一组基；

            \item  设 $B = \{\alpha_1, \alpha_2\}$ 是 $\mathbf{R}^2$ 的一组基，则 $\{\alpha_1 + \alpha_2, \alpha_1 - \alpha_2\}$ 也是 $\mathbf{R}^2$ 的一组基；

            \item 一个有限维线性空间内只含有有限个子空间；

            \item  若 $W_1, W_2$ 是 $ \mathbf{R}^n $ 的两个子空间，$B_1, B_2$ 分别是 $W_1, W_2$ 的基，则存在 $ \mathbf{R}^n $ 的一组基 $B$，使得 $B \supseteq B_1 \cup B_2$。

        \end{enumerate}
\end{exercise}
\noindent 来源：\href{https://github.com/yhwu-is/Linear-Algebra-Left-Undone/blob/new/%E8%AE%B2%E4%B9%89/%E4%B8%93%E9%A2%98/3%20%E6%9C%89%E9%99%90%E7%BB%B4%E7%BA%BF%E6%80%A7%E7%A9%BA%E9%97%B4.tex}{L3-A-1}。
\begin{solve}
\begin{enumerate}
                \item 错. 反例：$\alpha_1=(1,0),\alpha_2=(2,0),\alpha_3=(0,1)$，则 $\alpha_1,\alpha_2,\alpha_3$ 线性相关而 $\alpha_3$ 不是 $\alpha_1.\alpha_2$ 的线性组合.

                \item 对. 该命题的等价命题（逆否命题）是：若存在一个向量是其余向量的线性组合，则 $\alpha_1,\ldots,\alpha_m$ 线性相关. 这正是原讲义定理 3.1 的内容，因而成立.

                \item 错. 反例：$\alpha_1=(1,0),\alpha_2=(0,1),\alpha_3=(1,1)$，则 $\alpha_1,\alpha_2,\alpha_3$ 两两无关，而三者线性相关. 可证两两无关是向量组无关的必要条件.

                \item 错. 反例：$\alpha_1=(1,0),\alpha_2=(0,0),\beta_1=(0,0),\beta_2=(0,1)$，有 $\alpha_1,\alpha_2$ 相关，$\beta_1,\beta_2$ 相关，而 $\alpha_1+\beta_1$ 与 $\alpha_2+\beta_2$ 线性无关.

                \item 错. 若 $\alpha_1,\ldots,\alpha_n$ 线性无关，有
                    \[\lambda_1\alpha_1+\cdots+\lambda_n\alpha_n=0\implies\lambda_1=\lambda_2=\cdots=\lambda_n=0,\]
                    判断 $\alpha_1+\alpha_2,\alpha_2+\alpha_3,\ldots,\alpha_n+\alpha_1$ 是否无关. 设
                    \[\lambda_1'(\alpha_1+\alpha_2)+\cdots+\lambda_n'(\alpha_n+\alpha_1)=0,\]
                    则
                    \[(\lambda_n'+\lambda_1')\alpha_1+(\lambda_1'+\lambda_2')\alpha_2+\cdots+(\lambda_{n-1}'+\lambda_{n}')\alpha_n=0,\]
                    则
                    \[\begin{cases} \begin{aligned}
                                \lambda_n'+\lambda_1'       & = 0               \\
                                                            & \vdotswithin{ = } \\
                                \lambda_{n-1}'+\lambda_{n}' & = 0               \\
                            \end{aligned} \end{cases}.\]
                    解该方程可得 $\lambda_n'=(-1)^{n-1}\lambda_1'$，因此当 $n$ 为偶数时，上述方程组有非零解，则向量组相关，而当 $n$ 为奇数时，向量组无关. 综上，该命题不成立.

                \item 对. 由原讲义定理 3.1，不妨设 $\alpha_3$ 可由 $\alpha_1,\alpha_2$ 线性表示，则 $\alpha_1+\alpha_2,\alpha_2+\alpha_3,\alpha_3+\alpha_1$ 均可由 $\alpha_1,\alpha_2$ 线性表示，再由原讲义定理 3.3 可知，$\alpha_1+\alpha_2,\alpha_2+\alpha_3,\alpha_3+\alpha_1$ 线性相关.

                \item 错. 反例：取 $\alpha_0=\alpha_1-\alpha_2-\alpha_3$，则 $\alpha_0+\alpha_1=(\alpha_0+\alpha_2)+(\alpha_0+\alpha_3)$，三者线性相关，不是 $\mathbf{R}^3$ 的基.

                \item 对. 判断 $\alpha_1+\alpha_2$ 与 $\alpha_1-\alpha_2$ 是否无关.
                    \[\lambda_1(\alpha_1+\alpha_2)+\lambda_2(\alpha_1-\alpha_2)=0,\]
                    则有$(\lambda_1+\lambda_2)\alpha_1+(\lambda_1-\lambda_2)\alpha_2=0$，则 $\lambda_1+\lambda_2=0,\lambda_1-\lambda_2=0\implies \lambda_1=\lambda_2=0$，因此线性无关且个数等于维数，是一组基.

                \item 错. 反例：$\mathbf{R}^2$ 中过原点的直线 $L_0$ 是 $\mathbf{R}^2$ 的一个子空间. 显然这样的直线有无数条.

                \item 错. 反例：$\mathbf{R}^3$ 中，子空间 $W_1=\operatorname{span}(e_1,e_2)$，$W_2=\operatorname{span}(e_1+e_2,e_3)$，则 $B_1\cup B_2=\{e_1,e_2,e_3,e_1+e_2\}$，显然 $\mathbf{R}^3$ 中的任一组基都不可能包含四个元素.
            \end{enumerate}
\end{solve}

\begin{exercise}\label{prob:L3-A-2}
证明：如果向量组线性相关，把每个向量去掉$m$个位置一致的分量，得到的缩短组仍线性相关；如果向量组线性无关，把每个向量添加$m$个位置一致的分量，得到的加长组仍线性无关.
\end{exercise}
\noindent 来源：\href{https://github.com/yhwu-is/Linear-Algebra-Left-Undone/blob/new/%E8%AE%B2%E4%B9%89/%E4%B8%93%E9%A2%98/3%20%E6%9C%89%E9%99%90%E7%BB%B4%E7%BA%BF%E6%80%A7%E7%A9%BA%E9%97%B4.tex}{L3-A-2}。
\begin{solve}
\begin{enumerate}
                \item 若向量组线性相关，则对应该方程组有无穷多解. 去掉 $m$ 个分量，相当于删去该方程组中的任意 $m$ 行方程，依然有无穷多解. 这是因为对于原方程组的任意一个解，将其带入被削减后的方程组也依然成立. 故线性相关得证.

                \item 若向量组线性无关，对应原方程组仅有唯一解，也就是全零解. 增加 $m$ 个分量相当于增加 $m$ 个方程，依然只有唯一解，因为若出现非零解，代入原方程组对应的方程中不会成立，矛盾. 故线性无关得证.
            \end{enumerate}
\end{solve}

\begin{exercise}\label{prob:L3-A-3}
$a$取何值时，$\beta_1=(1,3,6,2)^\mathrm{T},\beta_2=(2,1,2,-1)^\mathrm{T},\beta_3=(1,-1,a,-2)^\mathrm{T}$线性无关？
\end{exercise}
\noindent 来源：\href{https://github.com/yhwu-is/Linear-Algebra-Left-Undone/blob/new/%E8%AE%B2%E4%B9%89/%E4%B8%93%E9%A2%98/3%20%E6%9C%89%E9%99%90%E7%BB%B4%E7%BA%BF%E6%80%A7%E7%A9%BA%E9%97%B4.tex}{L3-A-3}。
\begin{solve}
方程组：$x_1\beta_1+x_2\beta_2+x_3\beta_3=0$ 系数矩阵
            \[\begin{pmatrix}
                    1 & 2  & 1  \\
                    3 & 1  & -1 \\
                    6 & 2  & a  \\
                    2 & -1 &
                    -2\end{pmatrix}\rightarrow\begin{pmatrix}
                    1 & 2   & 1   \\
                    0 & -5  & -4  \\
                    0 & -10 & a-6 \\
                    0 & -5  & -4
                \end{pmatrix}\rightarrow\begin{pmatrix}
                    1 & 2  & 1   \\
                    0 & -5 & -4  \\
                    0 & 0  & a+2 \\
                    0 & 0  & 0
                \end{pmatrix},\]
            仅全零解的条件是 $a\neq-2$，此时向量组线性无关.
\end{solve}

\begin{exercise}\label{prob:L3-A-4}
设$\alpha_1,\alpha_2,\ldots,\alpha_n\in\mathbf{F}^n$. 证明：$\alpha_1,\alpha_2,\ldots,\alpha_n$线性无关的充要条件是$\mathbf{F}^n$中任一向量都可以由它们线性表示.
\end{exercise}
\noindent 来源：\href{https://github.com/yhwu-is/Linear-Algebra-Left-Undone/blob/new/%E8%AE%B2%E4%B9%89/%E4%B8%93%E9%A2%98/3%20%E6%9C%89%E9%99%90%E7%BB%B4%E7%BA%BF%E6%80%A7%E7%A9%BA%E9%97%B4.tex}{L3-A-4}。
\begin{solve}
\begin{enumerate}
                \item 必要性：$\alpha_1,\ldots,\alpha_n$ 线性无关，对于 $F^n$ 中的任一向量 $\beta$, $\alpha_1,\ldots,\alpha_n,\beta$ 的向量个数大于维数 $n$，则线性相关. 由原讲义定理 3.2，$\beta$ 可被 $\alpha_1,\ldots,\alpha_n$ 唯一表示.

                \item 充分性：由于 $F^n$ 中任意向量均可被 $\alpha_1,\ldots,\alpha_n$ 线性表示，并且向量个数等于维数. 则 $\alpha_1,\ldots,\alpha_n$ 是 $F^n$ 的一组基. 则 $\alpha_1,\ldots,\alpha_n$ 线性无关.

                      $^*$ 更详细的证明：对于 $F^n$ 的一组基 $e_1,\ldots,e_n$，其可被 $\alpha_1,\ldots,\alpha_n$ 表示. 若 $\alpha_1,\ldots,\alpha_n$ 线性相关，不妨设 $\alpha_n$ 可被 $\alpha_1,\ldots,\alpha_{n-1}$ 表示，则有 $e_1,\ldots,e_n$ 可被 $\alpha_1,\ldots,\alpha_{n-1}$ 表示. 由于 $e_1,\ldots,e_n$ 线性无关. 根据原讲义定理 3.3，$n\leqslant n-1$，矛盾. 因此得证.
            \end{enumerate}
\end{solve}

\begin{exercise}\label{prob:L3-A-5}
设$S_1=\{\alpha_1,\ldots,\alpha_s\},S_2=\{\beta_1,\ldots,\beta_t\}$是向量空间$V$的两个线性无关的子集，证明：$\alpha_1,\ldots,\alpha_s,\beta_1,\ldots,\beta_t$线性无关$\iff \operatorname{span}(S_1)\cap \operatorname{span}(S_2)=\{0\}$.
\end{exercise}
\noindent 来源：\href{https://github.com/yhwu-is/Linear-Algebra-Left-Undone/blob/new/%E8%AE%B2%E4%B9%89/%E4%B8%93%E9%A2%98/3%20%E6%9C%89%E9%99%90%E7%BB%B4%E7%BA%BF%E6%80%A7%E7%A9%BA%E9%97%B4.tex}{L3-A-5}。
\begin{solve}
\begin{enumerate}
                \item 必要性：对于 $\forall v\in\operatorname{span}(S_1) \cap \operatorname{span}(S_2)$ 有 $v=a_1\alpha_1+\cdots+a_s\alpha_s=b_1\beta_1+\cdots+b_t\beta_t$. 由于 $\alpha_1,\ldots,\alpha_s,\beta_1,\ldots,\beta_t$ 线性无关 $\implies a_1=\cdots=a_s=b_1=\cdots=b_t=0$，则 $v=0$，即 $\operatorname{span}(S_1)\cap\operatorname{span}(S_2)=\{0\}$.

                \item 充分性：考虑反证法. 如果 $\alpha_1,\ldots,\alpha_s,\beta_1,\ldots,\beta_t$ 线性相关，则存在不全为零的系数使得 $a_1\alpha_1+\cdots+a_s\alpha_s+b_1\beta_1\cdots+b_t\beta_t=0$. 因此存在一个向量 $v=a_1\alpha_1+\cdots+a_s\alpha_s=-(b_1\beta_1\cdots+b_t\beta_t)\neq 0$ 且 $v\in\operatorname{span}(S_1),v\in\operatorname{span}(S_2)$. 即存在非零向量 $v$ 属于 $\operatorname{span}(S_1),v\in\operatorname{span}(S_2)$，矛盾！则充分性得证.
            \end{enumerate}
\end{solve}

\begin{exercise}\label{prob:L3-A-7}
已知$\alpha_1=(1,2,4,3),\alpha_2=(1,-1,-6,6),\alpha_3=(-2,-1,2,-9),\alpha_4=(1,1,-2,7),\beta=(4,2,4,a)$.
        \begin{enumerate}
            \item 求子空间$\operatorname{span}(\alpha_1,\alpha_2,\alpha_3,\alpha_4)$的维数和一组基；

            \item 求$a$的值使得$\beta\in W$，并求$\beta$在 (1) 所选基下的坐标.
        \end{enumerate}
\end{exercise}
\noindent 来源：\href{https://github.com/yhwu-is/Linear-Algebra-Left-Undone/blob/new/%E8%AE%B2%E4%B9%89/%E4%B8%93%E9%A2%98/3%20%E6%9C%89%E9%99%90%E7%BB%B4%E7%BA%BF%E6%80%A7%E7%A9%BA%E9%97%B4.tex}{L3-A-7}。
\begin{solve}
\begin{enumerate}
                \item 也就是求 $\alpha_1,\alpha_2,\alpha_3,\alpha_4$ 的极大线性无关组. 利用讲义中所述求法：方程组
                      \[ a_1\alpha_1+\cdots+a_4\alpha_4=0 \]
                      对应系数矩阵 
\[\begin{pmatrix}
                              1 & 1  & -2 & 1  \\
                              2 & -1 & -1 & 1  \\
                              4 & -6 & 2  & -2 \\
                              3 & 6  & -9 & 7\end{pmatrix}\]
. 化简为行阶梯型：
\[\begin{pmatrix}
                              1 & 1  & -2 & 1  \\
                              0 & -3 & 3  & -1 \\
                              0 & 0  & 0  & 3  \\
                              0 & 0  & 0  & 0\end{pmatrix}\]
，因此 $\alpha_1,\alpha_2,\alpha_3,\alpha_4$ 有非零解，这四个向量线性相关. （其实此处已知矩阵秩为 3，即维数是 3）.

                      再选取 $\alpha_1,\alpha_2,\alpha_4$ 来求解方程 $a_1\alpha_1+a_2\alpha_2+a_4\alpha_4=0$：
                      \[\begin{pmatrix}
                              1 & 1  & 1  \\
                              2 & -1 & 1  \\
                              4 & -6 & -2 \\
                              3 & 6  & 7\end{pmatrix}\rightarrow\begin{pmatrix}
                              1 & 1  & 1  \\
                              0 & -3 & -2 \\
                              0 & 0  & 2  \\
                              0 & 0  & 0\end{pmatrix},\]
                      因此该方程组只有全零解，即 $\alpha_1,\alpha_2,\alpha_4$ 是$\alpha_1,\alpha_2,\alpha_3,\alpha_4$的极大线性无关组. 则 $\operatorname{span}(\alpha_1,\alpha_2,\alpha_4)$ 的维数是 3. 一组基是 $\alpha_1,\alpha_2,\alpha_4$ .

                \item 也就是 $a_1\alpha_1+\cdots+a_4\alpha_4=\beta$有解：增广矩阵 
\[\begin{pmatrix}
                              1 & 1  & -2 & 1  & 4 \\
                              2 & -1 & -1 & 1  & 2 \\
                              4 & -6 & 2  & -2 & 4 \\
                              3 & 6  & -9 & 7  & a\end{pmatrix}\]
化为 
\[\begin{pmatrix}
                              1 & 1  & -2 & 1         & 4   \\
                              0 & -3 & 3  & -1        & -6  \\
                              0 & 0  & 0  & -\frac 83 & 8   \\
                              0 & 0  & 0  & 0         & a-9\end{pmatrix}\]
，若方程有解，则 $a=9$. 求坐标，取 $x_3=0$ 代入得 $\beta=4\alpha_1+3\alpha_2-3\alpha_4$.
            \end{enumerate}
\end{solve}

\begin{exercise}\label{prob:L3-A-9}
证明：$B=\{1,x-a,(x-a)^2\}\enspace(a\neq 0)$是$\mathbf{R}[x]_3$的一组基，并求$\mathbf{R}[x]_3$的自然基$\{1,x,x^2\}$中每个向量关于基$B$的坐标.
\end{exercise}
\noindent 来源：\href{https://github.com/yhwu-is/Linear-Algebra-Left-Undone/blob/new/%E8%AE%B2%E4%B9%89/%E4%B8%93%E9%A2%98/3%20%E6%9C%89%E9%99%90%E7%BB%B4%E7%BA%BF%E6%80%A7%E7%A9%BA%E9%97%B4.tex}{L3-A-9}。
\begin{solve}
把这三个多项式关于自然基的坐标作为列，得到
\[P=\begin{pmatrix}1&-a&a^2\\0&1&-2a\\0&0&1\end{pmatrix}.\]
由于 $\det P=1$，三个多项式线性无关，因而构成三维空间的一组基。实际上 $a=0$ 时结论也成立。由
\[1=1,\qquad x=a+(x-a),\qquad x^2=a^2+2a(x-a)+(x-a)^2\]
得到所求坐标依次为 $(1,0,0)^{\mathrm T}$、$(a,1,0)^{\mathrm T}$、$(a^2,2a,1)^{\mathrm T}$。
\end{solve}

\begin{exercise}\label{prob:L3-A-10}
已知向量组$A=\{\alpha_1,\alpha_2,\alpha_3\},\enspace B=\{\alpha_1,\alpha_2,\alpha_3,\alpha_4\},\enspace C=\{\alpha_1,\alpha_2,\alpha_3,\alpha_5\}$的秩分别为$r(A)=r(B)=3,\enspace r(C)=4$. 证明：$\{\alpha_1,\alpha_2,\alpha_3,\alpha_5-\alpha_4\}$的秩为4.
\end{exercise}
\noindent 来源：\href{https://github.com/yhwu-is/Linear-Algebra-Left-Undone/blob/new/%E8%AE%B2%E4%B9%89/%E4%B8%93%E9%A2%98/3%20%E6%9C%89%E9%99%90%E7%BB%B4%E7%BA%BF%E6%80%A7%E7%A9%BA%E9%97%B4.tex}{L3-A-10}。
\begin{solve}
等价于证明 $\alpha_1,\alpha_2,\alpha_3,\alpha_5-\alpha_4$ 线性无关. 即求解
            \begin{equation}
                \lambda_1\alpha_1+\lambda_2\alpha_2+\lambda_3\alpha_3+\lambda_4(\alpha_5-\alpha_4)=0. \tag{*} \label{eq:3:A.10}
            \end{equation}

            由于 $r(A)=r(B)=3$ 可得 $A$ 线性无关. $B$ 线性相关. 由原讲义定理 3.2 得 $\alpha_4$ 可由 $\alpha_1,\alpha_2,\alpha_3$ 唯一表示：$\alpha_4=\mu_1\alpha_1+\mu_2\alpha_2+\mu_3\alpha_3$. 则代入 式~\eqref{eq:3:A.10}. 有
            \[(\lambda_1-\mu_1\lambda_4)\alpha_1+(\lambda_2-\mu_2\lambda_4)\alpha_2+(\lambda_3-\mu_3\lambda_4)\alpha_3+\lambda_4\alpha_5=0,\]
            因为 $r(C)=4$，$\alpha_1,\alpha_2,\alpha_3,\alpha_5$ 线性无关. 有 $\lambda_4=0,\lambda_1=\mu_1\lambda_4=0,\lambda_2=\mu_2\lambda_4=0,\lambda_3=\mu_3\lambda_4=0$. 故 $\alpha_1,\alpha_2,\alpha_3,\alpha_5-\alpha_4$ 线性无关. 原题得证.
\end{solve}

\begin{exercise}\label{prob:L3-A-11}
设向量组$\alpha_1,\alpha_2,\ldots,\alpha_s$的秩为$r$. 在其中任取$m$个向量$\alpha_{i1},\alpha_{i2},\ldots,\alpha_{im}$，证明：向量组$\alpha_{i1},\alpha_{i2},\ldots,\alpha_{im}$的秩$\geqslant r+m-s$.
\end{exercise}
\noindent 来源：\href{https://github.com/yhwu-is/Linear-Algebra-Left-Undone/blob/new/%E8%AE%B2%E4%B9%89/%E4%B8%93%E9%A2%98/3%20%E6%9C%89%E9%99%90%E7%BB%B4%E7%BA%BF%E6%80%A7%E7%A9%BA%E9%97%B4.tex}{L3-A-11}。
\begin{solve}
相当于从 $\alpha_1,\ldots,\alpha_s$ 向量中选取 $s-m$ 个向量丢弃，剩余向量的秩：
            \[r(\alpha_{i1},\ldots,\alpha_{im})\geqslant r-(s-m) =r+m-s.\]
\end{solve}

\begin{exercise}\label{prob:L3-B-1}
已知$\alpha_1\neq 0$，则$\alpha_1,\alpha_2,\ldots,\alpha_n$线性相关的充要条件是存在$i\enspace(2 \leqslant i \leqslant n)$使得$\alpha_i$可由$\alpha_1,\alpha_2,\ldots,\alpha_{i-1}$线性表示，且表示法唯一.
\end{exercise}
\noindent 来源：\href{https://github.com/yhwu-is/Linear-Algebra-Left-Undone/blob/new/%E8%AE%B2%E4%B9%89/%E4%B8%93%E9%A2%98/3%20%E6%9C%89%E9%99%90%E7%BB%B4%E7%BA%BF%E6%80%A7%E7%A9%BA%E9%97%B4.tex}{L3-B-1}。
\begin{solve}
充分性显然成立，下证必要性：由于 $\alpha_1,\ldots,\alpha_n$ 线性相关，则存在 $m$，其能使得 $\alpha_1,\ldots,\alpha_m$ 线性无关的最大下标，有 $1\leqslant m<n$. 因此 $i=m+1$，$\alpha_1,\ldots,\alpha_{i-1}$ 线性无关，$\alpha_1,\ldots,\alpha_i$ 线性相关. 可得 $\alpha_i$ 可被 $\alpha_1,\ldots,\alpha_{i-1}$ 唯一表示.
\end{solve}

\begin{exercise}\label{prob:L3-B-2}
证明以下两个结论：
        \begin{enumerate}
            \item 设$U$和$W$都是$V$的非零子空间，如果$U\subseteq W$，那么$\dim U \leqslant \dim W$；

            \item 设$U$和$W$都是$V$的非零子空间，$U\subseteq W$，且$\dim U = \dim W$，则$U = W$.
        \end{enumerate}
\end{exercise}
\noindent 来源：\href{https://github.com/yhwu-is/Linear-Algebra-Left-Undone/blob/new/%E8%AE%B2%E4%B9%89/%E4%B8%93%E9%A2%98/3%20%E6%9C%89%E9%99%90%E7%BB%B4%E7%BA%BF%E6%80%A7%E7%A9%BA%E9%97%B4.tex}{L3-B-2}。
\begin{solve}
\begin{enumerate}
                \item 设 $U$ 的一组基为 $u_1,\ldots,u_m$，$W$ 的一组基为 $w_1,\ldots,w_n$. 由于 $U\subseteq W$，则 $u_1,\ldots,u_m$ 可由 $w_1,\ldots,w_n$ 线性表示，且 $u_1,\ldots,u_m$ 线性无关. 由原讲义定理 3.3 的等价命题可得 $m\leqslant n$，则 $m=\dim U\leqslant\dim W=n$ 的得证.

                \item 因为 $\dim U=\dim W$，则 $u_1,\ldots,u_m$ 也是 $W$ 的一组基. 则 $W$ 的任意向量均可由 $u_1,\ldots,u_m$ 表示，可得 $W\subseteq U$，而 $U\subseteq W$，故有 $U=W$ 得证.
            \end{enumerate}
\end{solve}

\begin{exercise}\label{prob:L3-B-3}
设向量组$\alpha_1,\alpha_2,\ldots,\alpha_n$线性无关. 证明：在向量组$\beta,\alpha_1,\alpha_2,\ldots,\alpha_n$中至多有一个向量$\alpha_i\enspace(1 \leqslant i \leqslant n)$可被其前面的$i$个向量$\beta,\alpha_1,\alpha_2,\ldots,\alpha_{i-1}$线性表示.
\end{exercise}
\noindent 来源：\href{https://github.com/yhwu-is/Linear-Algebra-Left-Undone/blob/new/%E8%AE%B2%E4%B9%89/%E4%B8%93%E9%A2%98/3%20%E6%9C%89%E9%99%90%E7%BB%B4%E7%BA%BF%E6%80%A7%E7%A9%BA%E9%97%B4.tex}{L3-B-3}。
\begin{solve}
若 $\alpha_i$ 能由其前面的向量表示，则表达式中 $\beta$ 的系数不能为零，否则 $\alpha_1,\ldots,\alpha_n$ 线性相关。因此可解出
\[\beta\in\operatorname{span}(\alpha_1,\ldots,\alpha_i).\]
若还有 $j>i$ 满足相同条件，把上式代入 $\alpha_j$ 关于 $\beta,\alpha_1,\ldots,\alpha_{j-1}$ 的表示，便得到 \[\alpha_j\in\operatorname{span}(\alpha_1,\ldots,\alpha_{j-1}),\]与已知线性无关矛盾。因此这样的下标至多有一个。
\end{solve}

\begin{exercise}\label{prob:L3-B-4}
证明：$1,e^{\lambda_1\cdot x},e^{\lambda_2\cdot x}$（$\lambda_1\neq\lambda_2$且均不为0）线性无关.
\end{exercise}
\noindent 来源：\href{https://github.com/yhwu-is/Linear-Algebra-Left-Undone/blob/new/%E8%AE%B2%E4%B9%89/%E4%B8%93%E9%A2%98/3%20%E6%9C%89%E9%99%90%E7%BB%B4%E7%BA%BF%E6%80%A7%E7%A9%BA%E9%97%B4.tex}{L3-B-4}。
\begin{solve}
设 $k_0+k_1e^{\lambda_1x}+k_2e^{\lambda_2x}=0$ 对所有实数 $x$ 成立。求导一次、两次，并在 $x=0$ 取值，得到
\[\begin{pmatrix}\lambda_1&\lambda_2\\\lambda_1^2&\lambda_2^2\end{pmatrix}\begin{pmatrix}k_1\\k_2\end{pmatrix}=0.\]
系数行列式为 $\lambda_1\lambda_2(\lambda_2-\lambda_1)\ne0$，故 $k_1=k_2=0$；代回恒等式得 $k_0=0$。于是三个函数线性无关。
\end{solve}

\begin{exercise}\label{prob:L3-B-5}
设线性空间$V(\mathbf{F})$中，向量$\beta$是$\alpha_1,\ldots,\alpha_r$的线性组合，但不是$\alpha_1,\ldots,\alpha_{r-1}$的线性组合. 证明：$\operatorname{span}(\alpha_1,\ldots,\alpha_{r-1},\alpha_r)=\operatorname{span}(\alpha_1,\ldots,\alpha_{r-1},\beta)$.
\end{exercise}
\noindent 来源：\href{https://github.com/yhwu-is/Linear-Algebra-Left-Undone/blob/new/%E8%AE%B2%E4%B9%89/%E4%B8%93%E9%A2%98/3%20%E6%9C%89%E9%99%90%E7%BB%B4%E7%BA%BF%E6%80%A7%E7%A9%BA%E9%97%B4.tex}{L3-B-5}。
\begin{solve}
只需证明 $\alpha_r$ 可以被 $\alpha_1,\ldots,\alpha_{r-1},\beta$ 表示即可. 由于 $\beta$ 是 $\alpha_1,\ldots,\alpha_r$ 的线性组合，若 $\lambda_r=0$，则 $\beta$ 是 $\alpha_1,\ldots,\alpha_{r-1}$ 的线性组合. 这与条件矛盾. 因此 \[\alpha_r=-\dfrac 1{\lambda_r}(\lambda_1\alpha_1+\cdots+\lambda_{r-1}\alpha_{r-1}-\beta),\] 则这两组向量等价，即 $\operatorname{span}(\alpha_1,\ldots,\alpha_{r-1},\alpha_r)=\operatorname{span}(\alpha_1,\ldots,\alpha_{r-1},\beta)$，得证.
\end{solve}

\begin{exercise}\label{prob:L3-B-6}
\label{item:3:正实数线性空间}
        设$\mathbf{R}^+$是所有正实数组成的集合，加法和数乘定义如下：
        \[ \forall a,b \in \mathbf{R}_+,\enspace k\in \mathbf{R}\colon\enspace a\oplus b = ab,\enspace k\odot a = a^k \]
        则 $\mathbf{R}^+$关于这一加法和数乘构成一个实线性空间. 求$\mathbf{R}^+$的一组基.
\end{exercise}
\noindent 来源：\href{https://github.com/yhwu-is/Linear-Algebra-Left-Undone/blob/new/%E8%AE%B2%E4%B9%89/%E4%B8%93%E9%A2%98/3%20%E6%9C%89%E9%99%90%E7%BB%B4%E7%BA%BF%E6%80%A7%E7%A9%BA%E9%97%B4.tex}{L3-B-6}。
\begin{solve}
分析该实线性空间，可以看出加法单位元为 1，数乘单位元为 1. 我们给出一组基：$e$，其中 $e$ 为自然对数的底数. 当然, 2，3 或者 10 都可以作为一组基. 接下来我们验证 $e$ 是 $\mathbf{R}^+$ 的基：$\forall a\in \mathbf{R}^+,\exists k=\ln a\in\mathbf{R}$，满足 $k\odot e=e^k=a$，则 $\operatorname{span}(e)=\mathbf{R}^+$ 成立. 由于该向量组只有一个元素，且并非设该空间的零元 1，则 $e$ 是线性无关的. 得证.
\end{solve}

\begin{exercise}\label{prob:L3-C-1}
已知$m$个向量$\alpha_1,\alpha_2,\ldots,\alpha_m$线性相关，但其中任意$m-1$个都线性无关，证明：
        \begin{enumerate}
            \item 若$k_1\alpha_1+\cdots+k_m\alpha_m=0$，则$k_1,\ldots,k_m$全为0或全不为0；

            \item 若以下等式成立
                  \begin{align*}
                      k_1\alpha_1+\cdots+k_m\alpha_m & =0 \\
                      l_1\alpha_1+\cdots+l_m\alpha_m & =0
                  \end{align*}
                  其中$l_1\neq 0$，证明：$\dfrac{k_1}{l_1}=\cdots=\dfrac{k_m}{l_m}$.
        \end{enumerate}
\end{exercise}
\noindent 来源：\href{https://github.com/yhwu-is/Linear-Algebra-Left-Undone/blob/new/%E8%AE%B2%E4%B9%89/%E4%B8%93%E9%A2%98/3%20%E6%9C%89%E9%99%90%E7%BB%B4%E7%BA%BF%E6%80%A7%E7%A9%BA%E9%97%B4.tex}{L3-C-1}。
\begin{solve}
\begin{enumerate}
                \item 若 不全为 0. 不妨设设至少有 $k_i=0$，则有 $k_1\alpha_1+\cdots+k_{i-1}\alpha_{i-1}+k_{i+1}\alpha_{i+1}+\cdots+k_m\alpha_m=0$，并且系数不全为 0. 因此 $\alpha_1,\ldots,\alpha_{i-1},\alpha_{i+1},\ldots,\alpha_m$ 这 $m-1$ 个向量相关，与题设矛盾. 则原题得证.

                \item $l_1\neq 0$，则 $l_2,\ldots,l_m$ 均不为 0.
                      \begin{enumerate}
                          \item 若 $k_1=\cdots=k_m=0$. 原式显然成立.

                          \item 若 $k_1,\ldots,k_m$ 不全为 0. 则
                                \begin{gather*}
                                    l_1(k_1\alpha_1+\cdots+k_m\alpha_m)=0, \\
                                    k_1(l_1\alpha_1+\cdots+l_m\alpha_m)=0.
                                \end{gather*}
                                两式相减，得
                                \[(k_2l_1-k_1l_2)\alpha_2+\cdots+(k_ml_1-k_1l_m)\alpha_m=0,\]
                                因为 $\alpha_2,\ldots,\alpha_m$ 线性无关. 则以上系数均为 0. 故$\dfrac {k_2}{l_2}=\dfrac {k_1}{l_1},\ldots,\dfrac {k_m}{l_m}=\dfrac {k_1}{l_1}$. 得证.
                      \end{enumerate}
            \end{enumerate}
\end{solve}

\begin{exercise}\label{prob:L3-C-2}
（替换定理）设$\alpha_1,\alpha_2,\ldots,\alpha_r$线性无关，且可以被$\beta_1,\beta_2,\ldots,\beta_n$线性表示，则可以将$\beta_1,\beta_2,\ldots,\beta_n$中的$r$个向量替换成$\alpha_1,\alpha_2,\ldots,\alpha_r$后得到与$\beta_1,\beta_2,\ldots,\beta_n$等价的新向量组（注：可以使用数学归纳法证明）.
\end{exercise}
\noindent 来源：\href{https://github.com/yhwu-is/Linear-Algebra-Left-Undone/blob/new/%E8%AE%B2%E4%B9%89/%E4%B8%93%E9%A2%98/3%20%E6%9C%89%E9%99%90%E7%BB%B4%E7%BA%BF%E6%80%A7%E7%A9%BA%E9%97%B4.tex}{L3-C-2}。
\begin{solve}
采用逐个替换的证明。开始时 $\beta_1,\ldots,\beta_n$ 张成一个包含全部 $\alpha_j$ 的空间。假设已经把其中 $k$ 个向量替换为 $\alpha_1,\ldots,\alpha_k$，且张成空间不变，其中 $0\le k<r$。将 $\alpha_{k+1}$ 写成当前生成组的线性组合。这个表示中至少有一个尚未替换的 $\beta_j$ 的系数非零，否则 $\alpha_{k+1}$ 在 $\alpha_1,\ldots,\alpha_k$ 的张成空间中，违背 $\alpha_1,\ldots,\alpha_r$ 线性无关。解出该 $\beta_j$，用 $\alpha_{k+1}$ 替换它，张成空间仍不变。重复 $r$ 次，就得到由全部 $\alpha_1,\ldots,\alpha_r$ 和原组中 $n-r$ 个 $\beta_j$ 组成的等价向量组。这里同时说明 $r\le n$，否则替换过程中会在放入全部 $\alpha_j$ 前耗尽可替换向量，与线性无关矛盾。
\end{solve}

\begin{exercise}\label{prob:L3-C-3}
设线性空间$V=\mathbf{F}^n$. 证明：
        \begin{enumerate}
            \item 存在$V$的子空间$W$，使得$W$的任一非零向量的分量均不为0；

            \item 若$V$的子空间$W$的任一非零向量的分量均不为0，则$\dim W=1$；

            \item 若$V$的子空间$W$的任一非零向量的零分量个数均不超过$r$，则$\dim W \leqslant r+1$.
        \end{enumerate}
\end{exercise}
\noindent 来源：\href{https://github.com/yhwu-is/Linear-Algebra-Left-Undone/blob/new/%E8%AE%B2%E4%B9%89/%E4%B8%93%E9%A2%98/3%20%E6%9C%89%E9%99%90%E7%BB%B4%E7%BA%BF%E6%80%A7%E7%A9%BA%E9%97%B4.tex}{L3-C-3}。
\begin{solve}
\begin{enumerate}
                \item $\alpha=(1,1,\ldots,1)^{\mathrm{T}}, W=\operatorname{span}(\alpha)$,显然 $W$ 是满足条件的一维子空间.

                \item 考虑反证法：若 $\dim W>1$，则 $W$ 中存在线性无关的两向量. 由条件，
                      \[a_1,\ldots,a_n,b_1,\ldots,b_n\neq 0,\]
                      可设 $a_1=kb_1,k\in F$，因此 $\alpha-k\beta=(0,a_2-kb-2,\ldots,a_n-kb_n)^{\mathrm{T}}\in W$. 且由于 $\alpha,\beta$ 无关，$\alpha-k\beta\neq 0$ 但存在分量为 0，这与条件矛盾. 故 $\dim W=1$.

                \item 考虑反证法：若 $\dim W>r+1$，则存在 $r+2$ 个线性无关的向量，设为
                      \begin{align*}
                          \alpha_1     & = (a_{11},a_{12},\ldots,a_{1(r+1)},\ldots,a_{1n})^{\mathrm{T}},                 \\
                                       & \vdotswithin{ = }                                                              \\
                          \alpha_{r+2} & = (a_{(r+2)1},a_{(r+2)2},\ldots,a_{(r+2)(r+1)},\ldots,a_{(r+2)n})^{\mathrm{T}}.
                      \end{align*}
                      取这些向量的前 $r+1$ 个分量组成新的向量组：
                      \begin{align*}
                          \beta_1     & = (a_{11},a_{12},\ldots,a_{1(r+1)})^{\mathrm{T}},             \\
                                      & \vdotswithin{ = }                                            \\
                          \beta_{r+2} & = (a_{(r+2)1},a_{(r+2)2},\ldots,a_{(r+2)(r+1)})^{\mathrm{T}}.
                      \end{align*}
                      由于 $\beta_1,\ldots,\beta_{r+2}$ 是 $r+2$ 个 $r+1$ 维向量，其必然线性相关，则存在不全为 0 的系数 $\lambda_1,\ldots,\lambda_{r+2}$，$\lambda_1\beta_1+\cdots+\lambda_{r+2}\beta_{r+2}=0$.  由于 $\alpha_1,\ldots,\alpha_{r+2}$ 线性无关知：$\lambda_1\alpha_1+\cdots+\lambda_{r+2}\alpha_{r+2}\neq 0$，且其属于 $W$. 但其前 $r+1$ 个分量均为 0，这与条件矛盾. 故 $\dim W\leqslant r+1$ 得证.
            \end{enumerate}
\end{solve}

\begin{exercise}\label{prob:L3-C-4}
设 $\mathbf{Q}(\sqrt[3]{2}) = \{a+b\sqrt[3]{2}+c\sqrt[3]{4}\mid a,b,c\in\mathbf{Q}\}$，证明 $\mathbf{Q}(\sqrt[3]{2})$ 是 $\mathbf{Q}$ 上的线性空间并求其维数.
\end{exercise}
\noindent 来源：\href{https://github.com/yhwu-is/Linear-Algebra-Left-Undone/blob/new/%E8%AE%B2%E4%B9%89/%E4%B8%93%E9%A2%98/3%20%E6%9C%89%E9%99%90%E7%BB%B4%E7%BA%BF%E6%80%A7%E7%A9%BA%E9%97%B4.tex}{L3-C-4}。
\begin{solve}
先证明 $\mathbf{Q}(\sqrt[3]{2})$ 是 $\mathbf{Q}$ 上的线性空间：

            $\mathbf{Q}(\sqrt[3]{2})$ 对通常的加法构成交换群以及其封闭性是显然的，在此不再赘述，下阐述数乘性质：
            \begin{enumerate}
                \item $\forall v = a+b\sqrt[3]{2}+c\sqrt[3]{4} \in \mathbf{Q}(\sqrt[3]{2}), \enspace 1 (a+b\sqrt[3]{2}+c\sqrt[3]{4}) = a+b\sqrt[3]{2}+c\sqrt[3]{4}$；

                \item $\forall v = a+b\sqrt[3]{2}+c\sqrt[3]{4} \in \mathbf{Q}(\sqrt[3]{2}), \forall \lambda, \mu \in \mathbf{Q}, \enspace \lambda(\mu v) = (\lambda \mu) v$；

                \item $\forall v = a+b\sqrt[3]{2}+c\sqrt[3]{4} \in \mathbf{Q}(\sqrt[3]{2}), \forall \lambda, \mu \in \mathbf{Q}, \enspace (\lambda + \mu) v = (\lambda + \mu)a + (\lambda + \mu) b\sqrt[3]{2} + (\lambda + \mu) c\sqrt[3]{4} = \lambda(a+b\sqrt[3]{2}+c\sqrt[3]{4}) + \mu(a+b\sqrt[3]{2}+c\sqrt[3]{4}) = \lambda v + \mu v$；

                \item $\forall v = a_1 + b_1\sqrt[3]{2} + c_1\sqrt[3]{4}, u = a_2 + b_2\sqrt[3]{2} + c_2\sqrt[3]{4} \in \mathbf{Q}(\sqrt[3]{2}), \forall \lambda \in \mathbf{Q}, \enspace \lambda(v+u) = \lambda (a_1 + a_2) + \lambda (b_1 + b_2)\sqrt[3]{2} + \lambda (c_1 + c_2)\sqrt[3]{4} = \lambda (a_1 + b_1\sqrt[3]{2} + c_1\sqrt[3]{4}) + \lambda (a_2 + b_2\sqrt[3]{2} + c_2\sqrt[3]{4}) = \lambda v + \lambda u$；

                \item （封闭性）$\forall v = a+b\sqrt[3]{2}+c\sqrt[3]{4} \in \mathbf{Q}(\sqrt[3]{2}), \forall \lambda \in \mathbf{Q}, \enspace \lambda v = \lambda a + \lambda b\sqrt[3]{2} + \lambda c\sqrt[3]{4} \in \mathbf{Q}(\sqrt[3]{2})$.
            \end{enumerate}
            故 $\mathbf{Q}(\sqrt[3]{2})$ 是 $\mathbf{Q}$ 上的线性空间.

            下求 $\mathbf{Q}(\sqrt[3]{2})$ 的维数. 考虑找 $\mathbf{Q}(\sqrt[3]{2})$ 的一组基，我们给出 $\{1, \sqrt[3]{2}, \sqrt[3]{4}\}$，下证其为 $\mathbf{Q}(\sqrt[3]{2})$ 的一组基：

            \begin{enumerate}
                \item （线性无关）令 $k_1 + k_2 \sqrt[3]{2} + k_3 \sqrt[3]{4} = 0$，其中 $k_1, k_2, k_3 \in \mathbf{Q}$，由于左边只有第一项为有理数，故有 $k_1 = 0$，进而有 $\sqrt[3]{2} (k_2 + k_3 \sqrt[3]{2}) = 0$，又可得到 $k_2 = 0$，并且 $k_3 = 0$. 故 $\{1, \sqrt[3]{2}, \sqrt[3]{4}\}$ 线性无关.

                \item （张成空间）由 $\mathbf{Q}(\sqrt[3]{2})$ 的定义易得.
            \end{enumerate}
            综上，$\mathbf{Q}(\sqrt[3]{2})$ 是 $\mathbf{Q}$ 上的线性空间，且 $\dim \mathbf{Q}(\sqrt[3]{2}) = 3$.
\end{solve}

\begin{exercise}\label{prob:L3-C-5}
设$\mathbf{K} \subseteq \mathbf{F} \subseteq \mathbf{E}$是三个数域，已知$\mathbf{F}$作为$\mathbf{K}$上的线性空间是$n$维的，$\mathbf{E}$作为$\mathbf{F}$上的线性空间是$m$维的，证明：$\mathbf{E}$作为$\mathbf{K}$上的线性空间是$mn$维的.
\end{exercise}
\noindent 来源：\href{https://github.com/yhwu-is/Linear-Algebra-Left-Undone/blob/new/%E8%AE%B2%E4%B9%89/%E4%B8%93%E9%A2%98/3%20%E6%9C%89%E9%99%90%E7%BB%B4%E7%BA%BF%E6%80%A7%E7%A9%BA%E9%97%B4.tex}{L3-C-5}。
\begin{solve}
设 $f_1,\ldots,f_n$ 是 $\mathbf F$ 作为 $\mathbf K$ 空间的一组基，$e_1,\ldots,e_m$ 是 $\mathbf E$ 作为 $\mathbf F$ 空间的一组基。考虑 $mn$ 个向量 $f_je_i$（$1\le i\le m,1\le j\le n$）。若
\[\sum_{i=1}^m\sum_{j=1}^n k_{ij}f_je_i=0,\qquad k_{ij}\in\mathbf K,\]
则由 $e_i$ 在 $\mathbf F$ 上线性无关，得每个 $\sum_jk_{ij}f_j=0$；再由 $f_j$ 在 $\mathbf K$ 上线性无关，得到所有 $k_{ij}=0$。故这 $mn$ 个向量线性无关。任取 $e\in\mathbf E$，先写成 $e=\sum_i\mu_ie_i$，再把每个 $\mu_i\in\mathbf F$ 写成 $\sum_jk_{ij}f_j$，即得 $e$ 在这些乘积向量下的线性表示。它们既张成又线性无关，所以构成基，所求维数为 $mn$。
\end{solve}

\begin{exercise}\label{prob:L3-例-1}
\begin{enumerate}[label=(\arabic*)]
        \item 判断$\mathbf{R}^3$中向量$(1,1,0),(0,1,1),(1,0,-1)$的线性相关性；

        \item 判断$\mathbf{R}^3$中向量$(1,-3,1),(-1,2,-2),(1,1,3)$的线性相关性；

        \item \label{item:3:线性相关性:3}
              判断$\mathbf{R}[x]_3$中$p_1(x)=1+x,\ p_2(x)=1-x,\ p_3(x)=x+x^2$的线性相关性；

        \item 判断连续函数全体构成的线性空间中$1,\ \sin^2x,\ \cos^2x$的线性相关性；

        \item \label{item:3:线性相关性:5}
              判断连续函数全体构成的线性空间中$1,\ 2^x,\ 2^{-x}$的线性相关性.
    \end{enumerate}
\end{exercise}
\noindent 来源：\href{https://github.com/yhwu-is/Linear-Algebra-Left-Undone/blob/new/%E8%AE%B2%E4%B9%89/%E4%B8%93%E9%A2%98/3%20%E6%9C%89%E9%99%90%E7%BB%B4%E7%BA%BF%E6%80%A7%E7%A9%BA%E9%97%B4.tex}{L3-例-1}。
\begin{solve}
\begin{enumerate}
        \item 根据定义，应求解方程
              \[\lambda_1(1,1,0) + \lambda_2(0,1,1) + \lambda_3(1,0,-1) = 0,\]
              即
              \[ \begin{cases}
                      \lambda_1 + \lambda_3 = 0 \\
                      \lambda_1 + \lambda_2 = 0 \\
                      \lambda_2 - \lambda_3 = 0
                  \end{cases} \]
              解得基础解系$k(1,-1,-1)^\mathrm{T}$，所以存在非零解，向量组线性相关.

        \item 求解方程
              \[\lambda_1(1,-3,1) + \lambda_2(-1,2,-2) + \lambda_3(1,1,3) = 0,\]
              即
              \[ \begin{cases}
                      \lambda_1 - \lambda_2 + \lambda_3 = 0    \\
                      -3\lambda_1 + 2\lambda_2 + \lambda_3 = 0 \\
                      \lambda_1 - 2\lambda_2 + 3\lambda_3 = 0
                  \end{cases} \]
              解得 $\lambda_1 = \lambda_2 = \lambda_3 = 0$，此向量组线性无关.

        \item 求解方程
              \begin{align*}
                  \lambda_1p_1(x) + \lambda_2p_2(x) + \lambda_3p_3(x)                           & = 0 \\
                  (\lambda_1 + \lambda_2) + (\lambda_1 - \lambda_2 + \lambda_3)x + \lambda_3x^2 & = 0
              \end{align*}
              所以需要求解方程组
              \[ \begin{cases}
                      \lambda_1 + \lambda_2 = 0             \\
                      \lambda_1 - \lambda_2 + \lambda_3 = 0 \\
                      \lambda_3 = 0
                  \end{cases} \]
              解得 $\lambda_1 = \lambda_2 = \lambda_3 = 0$，此向量组线性无关.

        \item 易知 $-1 + \sin^2x + \cos^2x = 0$，对应系数为 $-1, 1, 1$，不全为零，所以此向量组线性相关.

        \item 求解方程
              \[\lambda_1 + \lambda_2·2^{-x} + \lambda_3·2^x = 0\]
              很明显会发现仅凭此方程是难以求解的，方程数目不足. 注意到此方程应该对于任意的 $x$ 均成立，所以取 $x = 0, x = 1, x = -1$，得到方程组
              \[ \begin{cases}
                      \lambda_1 + \lambda_2 + \lambda_3 = 0              \\[1ex]
                      \lambda_1 + \dfrac{1}{2}\lambda_2 + 2\lambda_3 = 0 \\[1ex]
                      \lambda_1 + 2\lambda_2 + \dfrac{1}{2}\lambda_3 = 0
                  \end{cases} \]
              解得 $\lambda_1 = \lambda_2 = \lambda_3 = 0$，此向量组线性无关.
    \end{enumerate}
\end{solve}

\begin{exercise}\label{prob:L3-例-2}
判断 $\mathbf{R}^3$ 中向量组 $\{\alpha_1, \alpha_2, \alpha_3\}$ 和 $\{\beta_1, \beta_2, \beta_3\}$ 的线性相关性, 其中

              \[
                  \alpha_1 = (1, 1, 0), \quad \alpha_2 = (0, 1, 1), \quad \alpha_3 = (1, 0, -1),
              \]
              \[
                  \beta_1 = (1, -3, 1), \quad \beta_2 = (-1, 2, -2), \quad \beta_3 = (1, 1, 3).
              \]
\end{exercise}
\noindent 来源：\href{https://github.com/yhwu-is/Linear-Algebra-Left-Undone/blob/new/%E8%AE%B2%E4%B9%89/%E4%B8%93%E9%A2%98/3%20%E6%9C%89%E9%99%90%E7%BB%B4%E7%BA%BF%E6%80%A7%E7%A9%BA%E9%97%B4.tex}{L3-例-2}。
\begin{solve}
容易看出, $\alpha_3 = \alpha_1 - \alpha_2$, 所以 $\alpha_1, \alpha_2, \alpha_3$ 线性相关. 但对 $\beta_1, \beta_2, \beta_3$ 不易看出是否有线性关系, 要按定义来判别. 设

              \begin{equation} \label{eq:线性相关性:1}
                  x_1 \beta_1 + x_2 \beta_2 + x_3 \beta_3 = 0
              \end{equation}

              即

              \[
                  x_1(1, -3, 1) + x_2(-1, 2, -2) + x_3(1, 1, 3) = (0, 0, 0).
              \]

              这个向量方程等价于以下的三个元线性齐次方程组:

              \[
                  \begin{cases}
                      x_1 - x_2 + x_3 = 0    \\
                      -3x_1 + 2x_2 + x_3 = 0 \\
                      x_1 - 2x_2 + 3x_3 = 0
                  \end{cases},
              \]

              容易解得这个方程组只有零解 $x_1 = x_2 = x_3 = 0$. 即只有全为零的 $x_1, x_2, x_3$ 才使得式~\eqref{eq:线性相关性:1} 成立, 故 $\beta_1, \beta_2, \beta_3$ 线性无关.

              一般若 $\beta_1 = (a_1, b_1, c_1), \beta_2 = (a_2, b_2, c_2), \beta_3 = (a_3, b_3, c_3)$, 则 $\beta_1, \beta_2, \beta_3$ 线性相关（线性无关）的充要条件是三元线性齐次方程组

              \[
                  \begin{cases}
                      a_1 x_1 + a_2 x_2 + a_3 x_3 = 0 \\
                      b_1 x_1 + b_2 x_2 + b_3 x_3 = 0 \\
                      c_1 x_1 + c_2 x_2 + c_3 x_3 = 0
                  \end{cases}
              \]

              有非零解（只有零解）. 用此法可得 $\mathbf{R}^3$ 中任何 4 个向量, $\mathbf{R}^n$ 中任何 $n + 1$ 个向量都线性相关.

              总的来说，我们有以下结论：

              列向量组$\alpha_1,\alpha_2,\ldots,\alpha_m$线性相关$\iff$齐次线性方程组$x_1\alpha_1+x_2\alpha_2+\cdots+x_m\alpha_m=0$有非零解;

              列向量组$\alpha_1,\alpha_2,\ldots,\alpha_m$线性无关$\iff$齐次线性方程组$x_1\alpha_1+x_2\alpha_2+\cdots+x_m\alpha_m=0$只有零解.
\end{solve}

\begin{exercise}\label{prob:L3-例-10}
已知 $\mathbf{R}^4$ 的一个子集 $S = \{a_1, a_2, a_3, a_4\}$, 其中
    \[
        a_1 = (1,1,0,1), \quad a_2 = (0,1,2,4), \quad
        a_3 = (2,1,-2,-2), \quad a_4 = (0,1,1,1).
    \]
    试求 $\operatorname{span}(S)$ 的维数及其一组基$B$。
\end{exercise}
\noindent 来源：\href{https://github.com/yhwu-is/Linear-Algebra-Left-Undone/blob/new/%E8%AE%B2%E4%B9%89/%E4%B8%93%E9%A2%98/3%20%E6%9C%89%E9%99%90%E7%BB%B4%E7%BA%BF%E6%80%A7%E7%A9%BA%E9%97%B4.tex}{L3-例-10}。
\begin{solve}
显然这一问题的关键就是求出 $S$ 的一组极大线性无关组，因为 $S$ 的任一极大线性无关组 $B$ 都是 $\operatorname{span}(S)$ 的基. 为了达到这一目标，我们首先将四个向量按列排成矩阵，即

    \begin{equation} \label{eq:极大线性无关组:1}
        \begin{pmatrix}
            1 & 0 & 2  & 0 \\
            1 & 1 & 1  & 1 \\
            0 & 2 & -2 & 1 \\
            1 & 4 & -2 & 1
        \end{pmatrix}
    \end{equation}

    对上述矩阵作初等行变换，所得的阶梯矩阵为

    \begin{equation} \label{eq:极大线性无关组:2}
        \begin{pmatrix}
            1 & 0 & 2  & 0 \\
            0 & 1 & -1 & 0 \\
            0 & 0 & 0  & 1 \\
            0 & 0 & 0  & 0
        \end{pmatrix}
    \end{equation}

    根据上面的引理，我们知道矩阵~\eqref{eq:极大线性无关组:1} 中的列的线性相关性与矩阵~\eqref{eq:极大线性无关组:2} 中的列的线性相关性是一致的，而 \eqref{eq:极大线性无关组:2} 中的经过初等变换的线性相关关系非常容易看出来：我们直接找每一行第一个非零元素所在的列对应的向量即可得到极大线性无关组，即 $S$ 的一个极大线性无关组为 $\{a_1, a_2, a_4\}$. 故 $a_1, a_2, a_4$ 是 $\operatorname{span}(S)$ 的一组基，$\operatorname{span}(S)$ 的维数为 3.
\end{solve}

\begin{exercise}\label{prob:L3-例-11}
设$V=\mathbf{R}[x]_4$，我们已有向量组$B=\{1+x,x^3+x^2+3x\}$，请将其扩充为$V$的一组基.
\end{exercise}
\noindent 来源：\href{https://github.com/yhwu-is/Linear-Algebra-Left-Undone/blob/new/%E8%AE%B2%E4%B9%89/%E4%B8%93%E9%A2%98/3%20%E6%9C%89%E9%99%90%E7%BB%B4%E7%BA%BF%E6%80%A7%E7%A9%BA%E9%97%B4.tex}{L3-例-11}。
\begin{solve}
按讲义中的方法，取定 $\mathbf{R}[x]_4$ 的自然基，给出$B$中向量的坐标
    \[\alpha_1 = (1, 1, 0, 0), \alpha_2 = (0, 3, 1, 1).\]
    再取一组基
    \[\beta_1 = (1, 0, 0, 0), \beta_2 = (0, 1, 0, 0), \beta_3 = (0, 0, 1, 0), \beta_4 = (0, 0, 0, 1).\]
    将这 6 个向量排列成矩阵，求解极大线性无关组.
    \[ \begin{pmatrix}
            1 & 0 & 1 & 0 & 0 & 0 \\
            1 & 3 & 0 & 1 & 0 & 0 \\
            0 & 1 & 0 & 0 & 1 & 0 \\
            0 & 1 & 0 & 0 & 0 & 1
        \end{pmatrix}
        \xRightarrow{\text{初等行变换}}
        \begin{pmatrix}
            1 & 0 & 1  & 0 & 0 & 0  \\
            0 & 1 & 0  & 0 & 0 & 1  \\
            0 & 0 & -1 & 1 & 0 & -3 \\
            0 & 0 & 0  & 0 & 1 & -1
        \end{pmatrix} \]
    则可取 $\alpha_1, \alpha_2, \beta_1, \beta_3$作为$V$的一组基，即$\{1 + x, 3x + x^2 + x^3, 1, x^2\}$.
\end{solve}

\begin{exercise}\label{prob:L3-例-12}
分别求$p(x)=a_0+a_1x+a_2x^2$在基$B_1=\{1,x,x^2\}$和$B_2=\{1,x-1,(x-1)^2\}$下的坐标.
\end{exercise}
\noindent 来源：\href{https://github.com/yhwu-is/Linear-Algebra-Left-Undone/blob/new/%E8%AE%B2%E4%B9%89/%E4%B8%93%E9%A2%98/3%20%E6%9C%89%E9%99%90%E7%BB%B4%E7%BA%BF%E6%80%A7%E7%A9%BA%E9%97%B4.tex}{L3-例-12}。
\begin{solve}
通过待定系数法解方程即可.

    在 $B_1$ 下：$(a_0, a_1, a_2)$. 在 $B_2$ 下：$(a_0 + a_1 + a_2, a_1 + 2a_2, a_2)$.
\end{solve}

\begin{exercise}\label{prob:L3-例-13}
求$\mathbf{R}[x]_4$中向量组$\{p_1=x^3+x+1,p_2=4x^3+2x^2+x,p_3=2x^3-2x^2+x+2,p_4=x^3+x^2+x\}$的极大线性无关组.
\end{exercise}
\noindent 来源：\href{https://github.com/yhwu-is/Linear-Algebra-Left-Undone/blob/new/%E8%AE%B2%E4%B9%89/%E4%B8%93%E9%A2%98/3%20%E6%9C%89%E9%99%90%E7%BB%B4%E7%BA%BF%E6%80%A7%E7%A9%BA%E9%97%B4.tex}{L3-例-13}。
\begin{solve}
相对于自然基 $(1,x,x^2,x^3)$，把四个多项式的坐标作为列，得
\[M=\begin{pmatrix}1&0&2&0\\1&1&1&1\\0&2&-2&1\\1&4&2&1\end{pmatrix}.\]
行变换 $R_2\leftarrow R_2-R_1$、$R_4\leftarrow R_4-R_1$ 后，再作 $R_3\leftarrow R_3-2R_2$、$R_4\leftarrow R_4-4R_2$，得到
\[\begin{pmatrix}1&0&2&0\\0&1&-1&1\\0&0&0&-1\\0&0&4&-3\end{pmatrix}.\]
后两行也线性无关，故秩为 $4$。因此原四个多项式全部线性无关，它们自身就是极大线性无关组，也是 $\mathbf R[x]_4$ 的一组基。
\end{solve}

\section{子空间、和与直和}\label{sec:topic4}
本节收录 15 道题。区分交、并、和与直和；直和要求分解唯一，仿射子集则允许平移后不经过原点。维数相加前须检查交空间。
\begin{exercise}\label{prob:L4-A-1}
设$V=\{(a_1,a_2,a_3,a_4) \mid a_1+a_2+a_3+a_4=0\}$，$W=\{(a_1,a_2,a_3,a_4) \mid a_1-a_2-a_3+a_4=0,a_1+a_2+a_3-a_4=0\}$.
        \begin{enumerate}
            \item 证明：$V$和$W$为$\mathbf{R}^4$的子空间；
            \item 分别求$V \cap W$，$V+W$以及$W$的补空间的维数与一组基.
        \end{enumerate}
\end{exercise}
\noindent 来源：\href{https://github.com/yhwu-is/Linear-Algebra-Left-Undone/blob/new/%E8%AE%B2%E4%B9%89/%E4%B8%93%E9%A2%98/4%20%E7%BA%BF%E6%80%A7%E7%A9%BA%E9%97%B4%E7%9A%84%E8%BF%90%E7%AE%97.tex}{L4-A-1}。
\begin{solve}
\begin{enumerate}
                \item 显然任意$\alpha\in v,\beta \in w$, 有 $\alpha,\beta\in \mathbf{R}^4$ 成立，因此只需证明 $v,w$ 封闭即可. 由于 $v,w$ 内向量的约束条件都是齐次线性方程，封闭性自然满足，则 $v,w$ 都是 $\mathbf{R}^4$ 的子空间.

                \item \begin{enumerate}
                          \item $V\cap W=\{(a_1,a_2,a_3,a_4) \mid a_1+a_2+a_3+a_4=0,a_1-a_2-a_3+a_4=0,a_1+a_2+a_3-a_4=0\}$ 或者是 $\{(0,a,-a,0) \mid \alpha\in \mathbf{R}\}$.

                          \item 由维数公式：$\dim V+\dim W=\dim (V\cap W)+\dim (V+W)$，易得$\dim V=3,\enspace \dim W=2,\enspace \dim (V\cap W)=1$，因此$\dim (V+W)=4$，则有$V+W=\mathbf{R}^4$.

                          \item 先得到$W$的基：列出齐次线性方程组：
                                \[\begin{cases}
                                        a_1-a_2-a_3+a_4=0 \\
                                        a_1+a_2+a_3-a_4=0\end{cases}
                                \]
                                高斯消元得到行阶梯矩阵
                                \[\begin{pmatrix}
                                        1 & 0 & 0 & 0  \\
                                        0 & 1 & 1 & -1
                                    \end{pmatrix}\]
                                秩为2，基分别为$\beta_1= (0,-1,1,0)$，$\beta_2= (0,1,0,1)$，则$W$的补空间维数是$\dim \mathbf{R}^4 - \dim W=2$.

                                利用基的扩张求其补空间的基. 也即求$\beta_1,\beta_2,e_1,e_2,e_3,e_4$的极大无关组，其中$e_1,e_2,e_3,e_4$是自然基.
                                \[\begin{pmatrix}
                                        0  & 0 & 1 & 0 & 0 & 0 \\
                                        -1 & 1 & 0 & 1 & 0 & 0 \\
                                        1  & 0 & 0 & 0 & 1 & 0 \\
                                        0  & 1 & 0 & 0 & 0 & 1
                                    \end{pmatrix}\rightarrow
                                    \begin{pmatrix}1 & 0 & 0 & 0 & 1 & 0  \\
                 0 & 1 & 0 & 1 & 1 & 0  \\
                 0 & 0 & 1 & 0 & 0 & 0  \\
                 0 & 0 & 0 & 1 & 1 & -1
                                    \end{pmatrix}\]
                                则$\beta_1,\beta_2,e_1,e_2$即是扩张后的基，因此$W$的补空间的一组基为$e_1,e_2$.
                      \end{enumerate}
            \end{enumerate}
\end{solve}

\begin{exercise}\label{prob:L4-A-2}
设 $f_1=-1+x,\ f_2=1-x^2,\ f_3=1-x^3,\ g_1=x-x^2,\ g_2=x+x^3,\ V_1=\operatorname{span}(f_1,f_2,f_3),\ V_2=\operatorname{span}(g_1,g_2)$，求：
        \begin{enumerate}
            \item $V_1+V_2$ 的基和维数；
            \item $V_1 \cap V_2$ 的基和维数；
            \item $V_2$ 在 $\mathbf{R}[x]_4$ 空间的补.
        \end{enumerate}
\end{exercise}
\noindent 来源：\href{https://github.com/yhwu-is/Linear-Algebra-Left-Undone/blob/new/%E8%AE%B2%E4%B9%89/%E4%B8%93%E9%A2%98/4%20%E7%BA%BF%E6%80%A7%E7%A9%BA%E9%97%B4%E7%9A%84%E8%BF%90%E7%AE%97.tex}{L4-A-2}。
\begin{solve}
\begin{enumerate}
                \item 先将各向量用坐标表示：
                      \begin{gather*}
                          f_1=(-1,1,0,0),f_2=(1,0,-1,0),f_3=(1,0,0,-1) \\
                          g_1=(0,1,-1,0),g_2=(0,1,0,1) \\
                          V_1+V_2=\operatorname{span}(f_1,f_2,f_3,g_1,g_2)
                      \end{gather*}
                      只需求 $f_1,f_2,f_3,g_1,g_2$ 的极大线性无关组即可.
                      \[\begin{pmatrix}
                              -1 & 1  & 1  & 0  & 0 \\
                              1  & 0  & 0  & 1  & 1 \\
                              0  & -1 & 0  & -1 & 0 \\
                              0  & 0  & -1 & 0  & 1
                          \end{pmatrix}\rightarrow
                          \begin{pmatrix}1 & -1 & -1 & 0 & 0 \\
                 0 & 1  & 1  & 1 & 1 \\
                 0 & 0  & 1  & 0 & 1 \\
                 0 & 0  & 0  & 0 & 1
                          \end{pmatrix}\]
                      则极大线性无关组为 $f_1,f_2,f_3,g_2$，因此 $V_1+V_2$ 的基为 $f_1,f_2,f_3,g_2$，维数为4.

                \item 易得$\dim V_1=3,\enspace \dim V_2=2$. 则$\dim (V_1\cap V_2)= \dim V_1+\dim V_2-\dim (V_1+V_2)=1$. 只需找到属于$V_1\cap V_2$的一个向量，其就是$V_1\cap V_2$的基.
                      \[U=\lambda_1f_1+\lambda_2f_2+\lambda_3f_3= \mu_1g_1+\mu_2g_2,\]
                      已知$g_1$可被$f_1,f_2,f_3,g_2$表示，则只需取$\mu_1=1$，再求解其它系数即解得：
                      \[\lambda_1=1,\lambda_2=1,\lambda_3=0,\mu_2=0,\]
                      因此$u=g_1=(0,1,-1,0)$或者$x-x^2$是$V_1\cap V_2$的一组基.

                \item 只需求$g_1,g_2,e_1,e_2,e_3,e_4$的极大无关组即可.
                      \[\begin{pmatrix}
                              0  & 0 & 1 & 0 & 0 & 0 \\
                              1  & 1 & 0 & 1 & 0 & 0 \\
                              -1 & 0 & 0 & 0 & 1 & 0 \\
                              0  & 1 & 0 & 0 & 0 & 1\end{pmatrix}\rightarrow
                          \begin{pmatrix}
                              1 & 1 & 0 & 1 & 0 & 0  \\
                              0 & 1 & 0 & 1 & 1 & 0  \\
                              0 & 0 & 1 & 0 & 0 & 0  \\
                              0 & 0 & 0 & 1 & 1 & -1\end{pmatrix}\]
                      则$g_1,g_2,e_1,e_2$是一组极大无关组. $V_2$在$\mathbf{R}[x]_4$的补是$\operatorname{span}(e_1,e_2)$其中$e_1=1,e_2=x$.
            \end{enumerate}
\end{solve}

\begin{exercise}\label{prob:L4-A-5}
设 $V_1,V_2$ 是线性空间 $V$ 的两个非平凡子空间，证明: $\exists \alpha \in V$,使得$\alpha \notin V_1$ 且 $\alpha \notin V_2$. 并在$\mathbf{R}^3$ 中举一例.
\end{exercise}
\noindent 来源：\href{https://github.com/yhwu-is/Linear-Algebra-Left-Undone/blob/new/%E8%AE%B2%E4%B9%89/%E4%B8%93%E9%A2%98/4%20%E7%BA%BF%E6%80%A7%E7%A9%BA%E9%97%B4%E7%9A%84%E8%BF%90%E7%AE%97.tex}{L4-A-5}。
\begin{solve}
\textbf{证:} 因为 $V_1, V_2$ 是非平凡子空间, 所以存在 $\alpha \notin V_1$. 若 $\alpha \notin V_2$, 则命题得证；若 $\alpha \in V_2$, 另有 $\beta \notin V_2$, 此时若 $\beta \notin V_1$, 命题也得证. 设 $\beta \in V_1$, 则 $\alpha \notin V_1, \alpha \in V_2, \beta \in V_1, \beta \notin V_2$.
            考虑 $\alpha + \beta \in V_1$, 由 $\beta \in V_1 \implies \alpha \in V_1$, 从而得出矛盾, 所以 $\alpha + \beta \notin V_1$. 类似可证 $\alpha + \beta \notin V_2$. 于是命题成立.
\end{solve}

\begin{exercise}\label{prob:L4-A-6}
设 $V_1,\ldots,V_m (m > 2)$ 是线性空间$V$的$m$个非平凡子空间，证明：$V$ 中存在一个同时不属于任何一个$V_i(1 \leqslant i \leqslant m)$的向量. 并在$\mathbf{R}^3$ 中举一例.
\end{exercise}
\noindent 来源：\href{https://github.com/yhwu-is/Linear-Algebra-Left-Undone/blob/new/%E8%AE%B2%E4%B9%89/%E4%B8%93%E9%A2%98/4%20%E7%BA%BF%E6%80%A7%E7%A9%BA%E9%97%B4%E7%9A%84%E8%BF%90%E7%AE%97.tex}{L4-A-6}。
\begin{solve}
\label{eg:4:A:5}
          \textbf{证:} （归纳法）当 $r=2$ 时由上题知成立.
          假设 $r = k-1$ 时也成立, 即存在 $\alpha \in V$
          使得 $\alpha$ 同时不属于 $V_i (i=1,2,\ldots,k-1)$
          下面证明 $r=k$ 时也成立. 显然当此 $\alpha \notin V_k$ 则 $\alpha$ 即为所求.
           若 $\alpha \in V_k$, 由于 $V_k$ 是 $V$ 的真子空间, 故有 $\beta \notin V_k$, 易知对任意 $V_i (i=1,2,\ldots,k-1)$ 都至多有一个 $k_i$, 使得 $\beta + k_i \alpha \in V_i; i=1,2,\ldots,k-1$, 我们取异于 $k_i$ 的数 $k$, 则必有 $\beta + k \alpha \notin V_i; i=1,2,\ldots,k-1$. 又由于 $\alpha \in V_k, \beta \notin V_k$ 故有 $\beta + k \alpha \notin V_k$, 故 $\beta + k \alpha \notin V_k \enspace (i = 1,2,\ldots,k)$.
\end{solve}

\begin{exercise}\label{prob:L4-B-1}
已知$V_1$是线性方程组\[\begin{cases}
                3x_1+4x_2-5x_3+7x_4=0 \\
                4x_1+11x_2-13x_3+16x_4=0
            \end{cases}\]
        的解空间，$V_2$是线性方程组\[\begin{cases}
                2x_1-3x_2+3x_3-2x_4=0 \\
                7x_1-2x_2+x_3+3x_4=0
            \end{cases}\]
        的解空间，分别求$V_1 \cap V_2$与$V_1+V_2$的基和维数.
\end{exercise}
\noindent 来源：\href{https://github.com/yhwu-is/Linear-Algebra-Left-Undone/blob/new/%E8%AE%B2%E4%B9%89/%E4%B8%93%E9%A2%98/4%20%E7%BA%BF%E6%80%A7%E7%A9%BA%E9%97%B4%E7%9A%84%E8%BF%90%E7%AE%97.tex}{L4-B-1}。
\begin{solve}
\begin{enumerate}
                \item $V_1\cap V_2$ 即是两组方程组合并后的解空间.
                      \[\begin{pmatrix}
                              3 & 4  & -5  & 7  \\
                              4 & 11 & -13 & 16 \\
                              2 & -3 & 3   & -2 \\
                              7 & -2 & 1   & 3
                          \end{pmatrix}\rightarrow
                          \begin{pmatrix}1 & 7   & -8 & 9   \\
                 0 & -17 & 19 & -20 \\
                 0 & 0   & 0  & 0   \\
                 0 & 0   & 0  & 0
                          \end{pmatrix}\]
                      则解向量$u_1=\left(\frac {3}{17},\frac{19}{17},1,0\right)^\mathrm{T},u_2=\left(\frac{13}{17},-\frac{20}{17},0,1\right)^\mathrm{T}$. $u_1,u_2$是$V_1\cap V_2$的基，则其维数为2.

                \item 易得$\dim V_1=2,\dim V_2=2$，则由维数公式$\dim (V_1+V_2)=\dim V_1+\dim V_2-\dim (V_1\cap V_2)= 2$又$V_1\cap V_2\subseteq V_1+V_2$，且二者维数相等. 则$V_1\cap V_2=V_1+V_2$，亦即$V_1=V_2$. 所以$V_1+V_2$的基也是$u_1,u_2$，同$V_1\cap V_2$.
            \end{enumerate}
\end{solve}

\begin{exercise}\label{prob:L4-B-2}
设$W_1,W_2$是线性空间$V(\mathbf{F})$的两个子空间. 证明以下命题等价：
        \begin{enumerate}
            \item $W_1 \cup W_2$为$V$ 的子空间；

            \item $W_1 \subseteq W_2$或$W_2 \subseteq W_1$；

            \item $W_1 \cup W_2=W_1+W_2$.
        \end{enumerate}
\end{exercise}
\noindent 来源：\href{https://github.com/yhwu-is/Linear-Algebra-Left-Undone/blob/new/%E8%AE%B2%E4%B9%89/%E4%B8%93%E9%A2%98/4%20%E7%BA%BF%E6%80%A7%E7%A9%BA%E9%97%B4%E7%9A%84%E8%BF%90%E7%AE%97.tex}{L4-B-2}。
\begin{solve}
先证 (1) 蕴含 (2)。若两子空间互不包含，可取 $u\in W_1\setminus W_2$ 和 $v\in W_2\setminus W_1$。若 $u+v\in W_1$，则 $v=(u+v)-u\in W_1$，矛盾；同理 $u+v\notin W_2$。因此并集不对加法封闭，与 (1) 矛盾。再证 (2) 蕴含 (3)：若 $W_1\subseteq W_2$，则并与和都等于 $W_2$；另一种包含关系同理。最后，和空间 $W_1+W_2$ 总是子空间，故 (3) 蕴含 (1)。这构成了完整的循环证明。
\end{solve}

\begin{exercise}\label{prob:L4-B-3}
证明：每个 $n$ 维空间均可以表示为 $n$ 个一维子空间的直和.
\end{exercise}
\noindent 来源：\href{https://github.com/yhwu-is/Linear-Algebra-Left-Undone/blob/new/%E8%AE%B2%E4%B9%89/%E4%B8%93%E9%A2%98/4%20%E7%BA%BF%E6%80%A7%E7%A9%BA%E9%97%B4%E7%9A%84%E8%BF%90%E7%AE%97.tex}{L4-B-3}。
\begin{solve}
我们很自然地想到，可以任取 $n$ 维空间 $V$ 的一组基 $\alpha_1,\alpha_2,\ldots,\alpha_n$，然后取 $V_i = \operatorname{span}(\alpha_i)$. 如果 $V = V_1 \oplus V_2 \oplus \cdots \oplus V_n$，我们就成功地为每个 $n$ 维空间找到了写成 $n$ 个一维子空间直和的分解形式. 下面我们来证明这一构造的正确性：
            \begin{enumerate}
                \item
                \item
            \end{enumerate}
\end{solve}

\begin{exercise}\label{prob:L4-B-4}
设$V$是域$\mathbf{F}$上的$n$维线性空间，$\alpha_1,\alpha_2,\ldots,\alpha_n$是$V$的一组基，且
        \begin{gather*}
            V_1=\operatorname{span}(\alpha_1+2\alpha_2+\cdots+n\alpha_n) \\
            V_2=\left\{k_1\alpha_1+k_2\alpha_2+\cdots+k_n\alpha_n \;\middle|\; k_1+\dfrac{k_2}{2}+\cdots+\dfrac{k_n}{n}=0\right\}
        \end{gather*}
        证明：
        \begin{enumerate}
            \item $V_2$是$V$的子空间；

            \item $V=V_1\oplus V_2$.
        \end{enumerate}
\end{exercise}
\noindent 来源：\href{https://github.com/yhwu-is/Linear-Algebra-Left-Undone/blob/new/%E8%AE%B2%E4%B9%89/%E4%B8%93%E9%A2%98/4%20%E7%BA%BF%E6%80%A7%E7%A9%BA%E9%97%B4%E7%9A%84%E8%BF%90%E7%AE%97.tex}{L4-B-4}。
\begin{solve}
\begin{enumerate}
                \item 显然$V_2\subseteq V$，只需证明$V_2$封闭即可. 对于$v=k_1\alpha_1+\cdots+k_n\alpha_n,v'=k_1'\alpha_1+\cdots+k_n'\alpha_n$. $\lambda v+\mu v'=(\lambda k_1+\mu k_1')\alpha_1+\cdots+(\lambda k_n+\mu k'_n)\alpha_n$，因此  $\lambda v+\mu v'\in V_2\in V_2$，$V_2$ 封闭，则$V_2$是$V$的子空间.

                \item 设$W=V_1+V_2$. 对任意 $v\in V_1,v=\lambda \alpha_1+2\lambda\alpha_2+\cdots+n\lambda\alpha_n$，则 $\lambda+2\dfrac{\lambda}2+\allowbreak\cdots+n\dfrac{\lambda}n=n\lambda$，因此当$v\ne 0$时，$v\not\in v_2$，则 $V_1\cap V_2=\{0\}$，有$W=V_1\oplus V_2$.

                      下证$W=V$，易得$\dim V_1=1$. 对于$V_2$，写成坐标形式，则其是一个方程组的解空间，系数矩阵 
\[A=\begin{pmatrix}
                              1 & \frac 12 & \cdots & \frac 1n \\
                              0 &          &        &          \\
                                & \ddots   &        &          \\
                                &          & \ddots &          \\
                                &          &        & 0
                          \end{pmatrix}\]
. 方程组为 $AK=0$，则 $\dim V_2=n-r(A)=n-1$，因此$\dim W=\dim V_1+\dim V_2=n=\dim V$，又$W\subseteq V$,则$W=V$得证，综上$V=V_1\oplus V_2$.
            \end{enumerate}
\end{solve}

\begin{exercise}\label{prob:L4-B-6}
已知$V_1,V_2$是有限维线性空间$V$的子空间，且$\dim(V_1+V_2)=\dim(V_1 \cap V_2)+1$. 证明：要么$V_1 \subseteq V_2$，要么$V_2 \subseteq V_1$.
\end{exercise}
\noindent 来源：\href{https://github.com/yhwu-is/Linear-Algebra-Left-Undone/blob/new/%E8%AE%B2%E4%B9%89/%E4%B8%93%E9%A2%98/4%20%E7%BA%BF%E6%80%A7%E7%A9%BA%E9%97%B4%E7%9A%84%E8%BF%90%E7%AE%97.tex}{L4-B-6}。
\begin{solve}
\begin{enumerate}
                \item 法一：反证法. 如果$\exists v_1\in V_1,v_1\not\in V_2$,且 $\exists v_2\in V_2,v_2\not\in V_1$，则有 $v_1,v_2\not \in V_1\cap V_2$ 并且 $v_1,v_2$ 线性无关. 设 $A=\{\alpha_1,\ldots ,\alpha_n\}$是 $V_1\cap V_2$的基. 由基扩张，设$B=\{\alpha_1,\ldots ,\alpha_n,\alpha_{n+1},\ldots ,\alpha_{n+k}\}$是$V_1+V_2$的基. 由于$v_1,v_2$不在$V_1\cap V_2$中，因此$A$无法表出$v_1,v_2$. 则$\alpha_1,\ldots ,\alpha_n,v_1,v_2$ 线性无关. 又$v_1,v_2\in V_1+V_2$，则$\alpha_1,\ldots ,\alpha_n,v_1,v_2$可由$\alpha_1,\ldots ,\alpha_{n+k}$ 线性表出. 由定理 $3.3$ 得 $n+2\leqslant n+k$，即 $\dim (V_1+V_2)\ge\dim (V_1\cap V_2)+2$，这与条件矛盾. 因此，$\forall v_1\in V_1,v_1\in V_2$ 或 $\forall v_2\in V_2,v_2\in V_1$. 即 $V_1\subseteq V_2$ 或 $V_2\subseteq V_1$，得证.

                \item 法二：维数公式. $\dim V_1+\dim V_2=\dim  (V_1+V_2)+\dim  (V_1\cap V_2)=2\dim  (V_1\cap V_2)+1$，即 $(\dim V_1-\dim  (V_1\cap V_2))+ (\dim V_2-\dim  (V_1\cap V_2))=1$. 于是，要么 $\dim V_1-\dim  (V_1\cap V_2)=0$，要么 $\dim V_2-\dim  (V_1\cap V_2)=0$，即 $V_1\subseteq V_2$ 或 $V_2\subseteq V_1$，得证.
            \end{enumerate}
\end{solve}

\begin{exercise}\label{prob:L4-B-7}
判断下列说法是否正确：
        \begin{enumerate}
            \item 若$V \subseteq V_1 \cup V_2 \cup \cdots \cup V_s$，则$V=(V_1 \cap V)\cup(V_2 \cap V)\cup\cdots\cup(V_s \cap V)$；

            \item 若$V \subseteq V_1+V_2+\cdots +V_s$，则$V=(V_1 \cap V)+(V_2 \cap V)+\cdots+(V_s \cap V)$.
        \end{enumerate}
\end{exercise}
\noindent 来源：\href{https://github.com/yhwu-is/Linear-Algebra-Left-Undone/blob/new/%E8%AE%B2%E4%B9%89/%E4%B8%93%E9%A2%98/4%20%E7%BA%BF%E6%80%A7%E7%A9%BA%E9%97%B4%E7%9A%84%E8%BF%90%E7%AE%97.tex}{L4-B-7}。
\begin{solve}
\begin{enumerate}
              \item 正确，若$\alpha\in V$，则$\alpha \in V_1\cup V_2\cdots \cup V_s$，即$\alpha \in V\cap (V_1\cup V_2\cdots \cup V_s)$，则$\alpha \in (V\cap V_1)\cup \cdots \cup (V\cap V_s)$，则$V\subseteq(V_1\cap V)\cup \cdots \cup (V_s\cap V)$. 另一边类似的同样成立，则有 $V=(V_1\cap V)\cup \cdots \cup (V_s\cap V)$ 得证.

              \item 错误，反例：设 $V_1,V_2,V$ 是平面$K$上三条过原点$O$的不重合直线，则$V\subseteq K=V_1+V_2$，但$V\cap V_1=\{0\},V\cap V_2=\{0\}$，$V\ne (V\cap V_1)+(V\cap V_2)$.
          \end{enumerate}
\end{solve}

\begin{exercise}\label{prob:L4-B-9}
设$A_1$和$A_2$均为$V$的仿射子集，证明：$A_1\cap A_2$是$V$的仿射子集或空集（可推广至任意交）.
\end{exercise}
\noindent 来源：\href{https://github.com/yhwu-is/Linear-Algebra-Left-Undone/blob/new/%E8%AE%B2%E4%B9%89/%E4%B8%93%E9%A2%98/4%20%E7%BA%BF%E6%80%A7%E7%A9%BA%E9%97%B4%E7%9A%84%E8%BF%90%E7%AE%97.tex}{L4-B-9}。
\begin{solve}
若交集为空，结论成立。否则取 $v\in A_1\cap A_2$。由于两个仿射子集都经过 $v$，可写成 $A_1=v+U_1$、$A_2=v+U_2$，其中 $U_1,U_2$ 是线性子空间。对任意 $x\in V$，
\[x\in A_1\cap A_2\iff x-v\in U_1\cap U_2.\]
故 $A_1\cap A_2=v+(U_1\cap U_2)$，是仿射子集。任意一族非空交集也可选取同一交点并照此证明；单点集通常只是仿射子集，仅当该点为零时才是线性子空间。
\end{solve}

\begin{exercise}\label{prob:L4-C-1}
设$W_0,W_1,W_2,\ldots,W_s$是线性空间$V$的$s+1$个非平凡子空间，且$W_0 \subseteq W_1 \cup W_2 \cup \cdots \cup W_s$. 证明：必存在$i$使得$W_0\subseteq W_i$.
\end{exercise}
\noindent 来源：\href{https://github.com/yhwu-is/Linear-Algebra-Left-Undone/blob/new/%E8%AE%B2%E4%B9%89/%E4%B8%93%E9%A2%98/4%20%E7%BA%BF%E6%80%A7%E7%A9%BA%E9%97%B4%E7%9A%84%E8%BF%90%E7%AE%97.tex}{L4-C-1}。
\begin{solve}
数域是无限域。先证明无限域上的向量空间 $U$ 不可能是有限个真子空间 $H_1,\ldots,H_s$ 的并。对 $s$ 归纳：$s=1$ 时结论就是“真子空间不等于全空间”。设 $s\ge2$，由归纳假设可取 $u\notin H_1\cup\cdots\cup H_{s-1}$，另取 $v\notin H_s$。对 $i<s$，直线 $\{v+tu:t\in\mathbf F\}$ 与 $H_i$ 至多有一个交点，否则两个交点相减将推出 $u\in H_i$。对 $H_s$，若 $u\in H_s$ 则无交点，否则同样至多一个交点。因为域无限，可避开这有限多个参数，得到不在任何 $H_i$ 中的向量，完成归纳。
现在在 $U=W_0$ 中应用结论。由题设，$W_0=\bigcup_{i=1}^s(W_0\cap W_i)$，所以至少一个交空间不是 $W_0$ 的真子空间，即 $W_0=W_0\cap W_i$，从而 $W_0\subseteq W_i$。
\end{solve}

\begin{exercise}\label{prob:L4-C-4}
设$v_1,\ldots,v_m\in V$. 令
        \[A=\{\lambda_1v_1+\cdots+\lambda_mv_m \mid \lambda_1,\ldots,\lambda_m\in\mathbf{F}\text{~且~}\lambda_1+\cdots+\lambda_m=1\}.\]
        证明：
        \begin{enumerate}
            \item $A$是$V$的仿射子集；

            \item $V$的每个包含$v_1,\ldots,v_m$的仿射子集均包含$A$；

            \item 存在某个$v\in V$和$V$的子空间$U$使得$A=v+U$且$\dim U\leqslant m-1$.
        \end{enumerate}
\end{exercise}
\noindent 来源：\href{https://github.com/yhwu-is/Linear-Algebra-Left-Undone/blob/new/%E8%AE%B2%E4%B9%89/%E4%B8%93%E9%A2%98/4%20%E7%BA%BF%E6%80%A7%E7%A9%BA%E9%97%B4%E7%9A%84%E8%BF%90%E7%AE%97.tex}{L4-C-4}。
\begin{solve}
\begin{enumerate}
            \item 对于 $v=\lambda_1 v_1+\cdots+\lambda_m v_m\in A$ 和 $w=\eta_1 v_1+\cdots+\eta_m v_m\in A$, 且 $\lambda_1,\cdots,\lambda_m\in  \mathbf{F}$, $\lambda_1+\cdots+\lambda_m=1$ , $\eta_1,\cdots,\eta_m\in \mathbf{F}$, $\eta_1+\cdots+\eta_m=1$. 对于任意的$\lambda\in \mathbf{F}$, 我们有
            \begin{align*}
                \lambda v+(1-\lambda) w=\sum_{i=1}^m(\lambda\lambda_i+(1-\lambda)\eta_i)v_i.
            \end{align*}
            注意到
            \[ \sum_{i=1}^m(\lambda\lambda_i+(1-\lambda)\eta_i)=\lambda\sum_{i=1}^m\lambda_i+(1-\lambda)\sum_{i=1}^n\eta_i=\lambda+(1-\lambda)=1, \]
            我们有$\lambda v+(1-\lambda) w\in A$. 故 $A$ 是 $V$ 的一个仿射子集
        \item 我们使用数学归纳法来证明，对于任意 $V$ 的仿射子集 $W$ ,其包含$v_1$, $\cdots$, $v_m$$k\le m$,
            如果 $\lambda_1+\cdots+\lambda_k=1$,
            我们有 \[ \sum_{j=1}^k\lambda_jv_j\in W. \]

            当$k=1$ 和 $k=2$ 时，结论是显然的.
            假设 $k$ 时成立,
            接下来对于 $k+1$ ($k+1\leqslant m$).
            我们假设$\lambda_1+\cdots+\lambda_{k+1}=1$.
            若 $\lambda_{k+1}=1$, 那么\[\sum_{j=1}^{k+1}\lambda_jv_j=v_{k+1}\in W.\]
            若$\lambda_{k+1}\ne 1$,那么\[ \frac{1}{1-\lambda_{k+1}}(\lambda_1+\cdots+\lambda_k)=1. \]
            由归纳假设，我们有\[ \frac{1}{1-\lambda_{k+1}}(\lambda_1v_1+\cdots+\lambda_kv_k)\in W. \]
            也有 \[ (1-a_{k+1})\left(\frac{1}{1-\lambda_{k+1}}(\lambda_1v_1+\cdots+\lambda_kv_k)\right)+a_{k+1}v_{k+1}\in W, \]
            即 \[ \lambda_1v_1+\cdots+\lambda_{k+1}v_{k+1}\in W. \]
            故而 $\lambda_1,\cdots,\lambda_m\in  \mathbf{F}$, $\lambda_1+\cdots+\lambda_m=1$,
            \[\lambda_1 v_1+\cdots+\lambda_m v_m\in W,\]有$A\subset W$.
        \item  注意到$\lambda_1+\cdots+\lambda_m=1$, 我们有\[ \lambda_1 v_1+\cdots+\lambda_m v_m=v_1+\lambda_2(v_2-v_1)+\cdots+\lambda_m(v_m-v_1). \]所以 $A\subset v_1+\operatorname{span}(v_2-v_1,\cdots,v_m-v_1)$.
            类似的，对于任意\[v\in v_1+\operatorname{span}(v_2-v_1,\cdots,v_m-v_1),\]$v$ 可以被写为 \[ v_1+\sum_{i=2}^m \lambda_i(v_i-v_1)=(1-\lambda_2-\cdots-\lambda_m)v_1+\sum_{i=2}^m \lambda_i v_i \]
            对于 $\lambda_2$, $\cdots$, $\lambda_m\in \mathbf{F}$. 注意到\[ (1-\lambda_2-\cdots-\lambda_m)+\sum_{i=2}^m \lambda_i=1, \]
            可以推出 $v_1+\operatorname{span}(v_2-v_1,\cdots,v_m-v_1)\subset A$.
            因此\[A=v_1+\operatorname{span}(v_2-v_1,\cdots,v_m-v_1).\]
            $v=v_1, U=\operatorname{span}(v_2-v_1,\cdots,v_m-v_1)$, 有 $\dim U\le m-1$
        \end{enumerate}
\end{solve}

\begin{exercise}\label{prob:L4-例-2}
设 $\alpha_1 = (1, 0, -1, 0)$，$\alpha_2 = (0, 1, 2, 1)$，$\alpha_3 = (2, 1, 0, 1)$，是四维实行向量空间 $V$ 中的向量，他们张成的子空间为 $V_1$；又设向量 $\beta_1 = (-1, 1, 1, 1)$，$\beta_2 = (1, -1, -3, -1)$，$\beta_3 = (-1, 1, -1, 1)$ 张成的子空间为 $V_2$，求 $V_1$ 和 $V_2$ 的交与和的基.
\end{exercise}
\noindent 来源：\href{https://github.com/yhwu-is/Linear-Algebra-Left-Undone/blob/new/%E8%AE%B2%E4%B9%89/%E4%B8%93%E9%A2%98/4%20%E7%BA%BF%E6%80%A7%E7%A9%BA%E9%97%B4%E7%9A%84%E8%BF%90%E7%AE%97.tex}{L4-例-2}。
\begin{solve}
\begin{enumerate}
        \item 方法一.  $V_1 +V_2$ 是由 $\alpha_i$ 和 $\beta_i$ 生成的，因此只需要求出这 $6$ 个向量的极大线性无关组即可. 将这 $6$ 个向量按列分块方式拼成矩阵，并用初等行变换将其化为阶梯形矩阵：
              \begin{align*}
                  \begin{pmatrix}
                      1  & 0 & 2 & -1 & 1  & -1 \\
                      0  & 1 & 1 & 1  & -1 & 1  \\
                      -1 & 2 & 0 & 1  & -3 & -1 \\
                      0  & 1 & 1 & 1  & -1 & 1
                  \end{pmatrix}
                  \xrightarrow{}
                  \begin{pmatrix}
                      1 & 0 & 2 & -1 & 1  & -1 \\
                      0 & 1 & 1 & 1  & -1 & 1  \\
                      0 & 2 & 2 & 0  & -2 & -2 \\
                      0 & 0 & 0 & 0  & 0  & 0
                  \end{pmatrix} \\
                  \xrightarrow{}
                  \begin{pmatrix}
                      1 & 0 & 2 & -1 & 1  & -1 \\
                      0 & 1 & 1 & 1  & -1 & 1  \\
                      0 & 0 & 0 & -2 & 0  & -4 \\
                      0 & 0 & 0 & 0  & 0  & 0
                  \end{pmatrix}
              \end{align*}

              所以就可以取 $\alpha_1$，$\alpha_2$，$\beta_1$ 为 $V_1 + V_2$ 的基（不唯一）.

              下面再来取 $V_1\cap V_2$ 的基，首先注意到 $\alpha_1$，$\alpha_2$ 是 $V_1$ 的基（从上面的矩阵即可看出），又不难验证 $\beta_1$，$\beta_2$ 是 $V_2$ 的基，$V_2$ 中的向量可以表示为 $\beta_1$，$\beta_2$ 的线性组合. 假设 $t_1\beta_1 + t_2\beta_2$ 属于 $V_1$，则向量组 $\alpha_1, \alpha_2, t_1\beta_1 + t_2\beta_2$ 和向量组 $\alpha_1, \alpha_2$ 的秩相等（因为 $\alpha_1, \alpha_2$ 是 $V_1$ 的基）. 因此，我们可以用矩阵方法来求出参数 $t_1, t_2$. 注意到
              \[ \begin{pmatrix}
                      1  & 0 & -t + t_2   \\
                      0  & 1 & t_1 - t_2  \\
                      -1 & 2 & t_1 - 3t_2 \\
                      0  & 1 & -t_1 - t_2
                  \end{pmatrix} \xrightarrow{} \begin{pmatrix}
                      1 & 0 & -t + t_2  \\
                      0 & 1 & t_1 - t_2 \\
                      0 & 2 & -2t_2     \\
                      0 & 0 & 0
                  \end{pmatrix} \xrightarrow{} \begin{pmatrix}
                      1 & 0 & -t + t_2  \\
                      0 & 1 & t_1 - t_2 \\
                      0 & 0 & -2t_1     \\
                      0 & 0 & 0
                  \end{pmatrix} \]

              所以可以得出当且仅当 $t_1 = 0$ 时 $t_1\beta_1 + t_2\beta_2$ 属于 $V_1$，所以 $V_1 \cap V_2$ 的基可取为 $\beta_2$.

        \item 方法二. 求 $V_1 + V_2$ 的基同方法一，现用解线性方程组的方法来求 $V_1 \cap V_2$ 的基. 因为 $\alpha_1$，$\alpha_2$ 是 $V_1$ 的基，$\beta_1$，$\beta_2$ 是 $V_2$ 的基，故对任一向量 $\gamma \in V_1 \cap V_2$，$\gamma = x_1\alpha_1 + x_2\alpha_2 = -x_3\beta_1 - x_4\beta_2$. 因此，求向量 $\gamma$ 等价于求解线性方程组

              \[ x_1\alpha_1 + x_2\alpha_2 + x_3\beta_1 + x_4\beta_2 = 0. \]

              上述线性方程的通解是 $(x_1, x_2, x_3, x_4) = k(-1, 1, 0, 1)$，从而 $\gamma = -k(\alpha_1 - \alpha_2) = -k\beta_2 (k \in \mathbf{R})$，于是 $\beta_2$ 是 $V_1 \cap V_2$ 的基.
    \end{enumerate}
\end{solve}

\begin{exercise}\label{prob:L4-例-8}
设$A$是$\mathbf{R}$上的$2\times 3$矩阵：
    \[A=\begin{pmatrix}
            1 & -1 & 2 \\ 1 & 0 & -1
        \end{pmatrix}.\]
    \begin{enumerate}
        \item 求齐次线性方程组$AX=0$的解空间$W$的一组基；

        \item 求商空间$\mathbf{R}^3/W$的维数和一组基.
    \end{enumerate}
\end{exercise}
\noindent 来源：\href{https://github.com/yhwu-is/Linear-Algebra-Left-Undone/blob/new/%E8%AE%B2%E4%B9%89/%E4%B8%93%E9%A2%98/4%20%E7%BA%BF%E6%80%A7%E7%A9%BA%E9%97%B4%E7%9A%84%E8%BF%90%E7%AE%97.tex}{L4-例-8}。
\begin{solve}
\begin{enumerate}
        \item 首先求 $AX=0$ 的解，得到其基础解系为$k(1,3,1)^{\mathrm{T}},k \in \mathbf{R}$,故其解空间的一组基为$\{(1,3,1)^{\mathrm{T}}\}$.
        \item 根据原讲义“商空间的维数公式”处 的证明过程，我们将 $W$ 的基扩张为 $\mathbf{R^3}$ 的一组基，如 $\{(1,3,1)^{\mathrm{T}},(1,0,0)^{\mathrm{T}},(0,1,0)^{\mathrm{T}}\}$，然后将扩张的基中``加上'' $W$ 变成仿射子集，即可得到商空间的一组基：
        \[
            \{(1,0,0)^{\mathrm{T}}+W,(0,1,0)^{\mathrm{T}}+W\}
        \]
        维数为 $2$. 这可以从基的数量得到，也可以由 $\mathbf{R}^3$ 的维数减去 $W$ 的维数得到.

    \end{enumerate}
\end{solve}

\section{线性映射、核与像}\label{sec:topic5}
本节收录 40 道题，其中 12 道为编者补充计算题。从线性性的定义出发构造映射，求核与像时使用同一组基，并用秩与零度定理核对维数。
\begin{exercise}\label{prob:L5-A-1}
设$\sigma: V_1\to V_2$是线性映射. 证明：$\sigma(W_1)$和$\sigma^{-1}(W_2)$分别是$V_2$和$V_1$的子空间.
\end{exercise}
\noindent 来源：\href{https://github.com/yhwu-is/Linear-Algebra-Left-Undone/blob/new/%E8%AE%B2%E4%B9%89/%E4%B8%93%E9%A2%98/5%20%E7%BA%BF%E6%80%A7%E6%98%A0%E5%B0%84.tex}{L5-A-1}。
\begin{solve}
$\sigma(W_1)$是$V_2$的子空间. $\forall v_1,v_2\in\sigma(W_1)$，$\exists w_1,w_2\in W_1$使得$\sigma(w_1)=v_1$，$\sigma(w_2)=v_2$，因此
            \[\sigma(w_1+w_2)=\sigma(w_1)+\sigma(w_2)=v_1+v_2\in\sigma(W_1),\]
            \[\sigma(\lambda w_1)=\lambda\sigma(w_1)=\lambda v_1\in\sigma(W_1),\]
            因此$\sigma(W_1)$是$V_2$的子空间.

            $\sigma^{-1}(W_2)$是$V_1$的子空间. $\forall v_1,v_2\in\sigma^{-1}(W_2)$，$\sigma(v_1),\sigma(v_2)\in W_2$，因此
            \[\sigma(v_1+v_2)=\sigma(v_1)+\sigma(v_2)\in W_2,\]
            \[\sigma(\lambda v_1)=\lambda\sigma(v_1)\in W_2,\]
            因此$v_1+v_2\in\sigma^{-1}(W_2)$，$\lambda v_1\in\sigma^{-1}(W_2)$，因此$\sigma^{-1}(W_2)$是$V_1$的子空间.
\end{solve}

\begin{exercise}\label{prob:L5-A-2}
设$\sigma,\tau \in \mathcal{L}(V,V)$且$\sigma^2=\sigma$，$\tau^2=\tau$. 证明：
        \begin{enumerate}
            \item $\sigma^k=\sigma$（幂等变换）；

            \item 若$(\sigma+\tau)^2=\sigma+\tau$，则$\sigma\tau=\theta$（零变换）；

            \item 设$\sigma\tau=\tau\sigma$，则$(\sigma+\tau-\sigma\tau)^2=\sigma+\tau-\sigma\tau$.
        \end{enumerate}
\end{exercise}
\noindent 来源：\href{https://github.com/yhwu-is/Linear-Algebra-Left-Undone/blob/new/%E8%AE%B2%E4%B9%89/%E4%B8%93%E9%A2%98/5%20%E7%BA%BF%E6%80%A7%E6%98%A0%E5%B0%84.tex}{L5-A-2}。
\begin{solve}
\begin{enumerate}
                \item 使用数学归纳法证明即可.

                \item 由 $ (\sigma + \tau)^2 = \sigma + \tau $，
                      \[ (\sigma + \tau)^2 = \sigma^2 + + \sigma \tau + \tau \sigma + \tau^2 = \sigma + \tau \]
                      得
                      \begin{equation} \label{eq:5:A:2:1}
                          \sigma \tau + \tau \sigma = \theta
                      \end{equation}
                      式~\eqref{eq:5:A:2:1} 两边左乘 $ \sigma $ 得
                      \begin{align}
                          \sigma(\sigma \tau + \tau \sigma) & = \sigma^2 \tau + \sigma \tau \sigma = \sigma \theta \notag    \\
                                                            & = \sigma \tau + \sigma \tau \sigma = \theta \label{eq:5:A:2:2}
                      \end{align}
                      式~\eqref{eq:5:A:2:1} 两边右乘 $ \sigma $ 得
                      \begin{align}
                          (\sigma \tau + \tau \sigma)\sigma & = \sigma \tau \sigma + \tau \sigma^2 = \theta \sigma \notag    \\
                                                            & = \sigma \tau \sigma + \tau \sigma = \theta \label{eq:5:A:2:3}
                      \end{align}
                      式~\eqref{eq:5:A:2:3} 减去式~\eqref{eq:5:A:2:2} 得
                      \begin{equation}
                          \sigma \tau - \tau \sigma = \theta \label{eq:5:A:2:4}
                      \end{equation}
                      由式~\eqref{eq:5:A:2:1} 与式~\eqref{eq:5:A:2:4} 得
                      \[ \sigma \tau = \theta \]

                \item 若 $ \sigma \tau = \tau \sigma $，则
                      \begin{align*}
                              & (\sigma + \tau - \sigma \tau)^2                                                                                             \\
                          ={} & \sigma^2 + \tau^2 + 2 \sigma \tau - \sigma \tau \sigma - \sigma \tau^2 - \sigma^2 \tau - \tau \sigma \tau + \sigma^2 \tau^2 \\
                          ={} & \sigma + \tau + 2 \sigma \tau - \sigma \tau \sigma - \sigma \tau - \sigma \tau - \tau \sigma - \sigma \tau + \sigma \tau    \\
                          ={} & \sigma + \tau - \sigma \tau
                      \end{align*}
            \end{enumerate}
\end{solve}

\begin{exercise}\label{prob:L5-A-3}
是否存在$\mathbf{R}^3$到$\mathbf{R}^2$的线性映射$\sigma$使得$\sigma(1,-1,1)=(1,0)$，$\sigma(1,1,1)=(0,1)$？
\end{exercise}
\noindent 来源：\href{https://github.com/yhwu-is/Linear-Algebra-Left-Undone/blob/new/%E8%AE%B2%E4%B9%89/%E4%B8%93%E9%A2%98/5%20%E7%BA%BF%E6%80%A7%E6%98%A0%E5%B0%84.tex}{L5-A-3}。
\begin{solve}
存在. 设
          \begin{align}
              \sigma(1, -1, 1) & = \sigma(\vec{e}_1 - \vec{e}_2 + \vec{e}_3) \notag                 \\
                               & = \sigma(\vec{e}_1) - \sigma(\vec{e}_2) + \sigma(\vec{e}_3) \notag \\
                               & = (1, 0) = \vec{\varepsilon}_1 \label{eq:5:A:3:1}
          \end{align} \\
          \begin{align}
              \sigma(1, 1, 1) & = \sigma(\vec{e}_1 + \vec{e}_2 + \vec{e}_3) \notag                 \\
                              & = \sigma(\vec{e}_1) + \sigma(\vec{e}_2) + \sigma(\vec{e}_3) \notag \\
                              & = (0, 1) = \vec{\varepsilon}_2 \label{eq:5:A:3:2}
          \end{align}
          可取 $ \sigma(\vec{e}_3) = (0, 0) $. 联立式~\eqref{eq:5:A:3:1} 与式~\eqref{eq:5:A:3:2} 解得
          \begin{align*}
              \sigma(\vec{e}_1) & = \frac{1}{2} (\vec{\varepsilon}_1 + \vec{\varepsilon}_2) \\
              \sigma(\vec{e}_2) & = \frac{1}{2} (\vec{\varepsilon}_1 - \vec{\varepsilon}_2)
          \end{align*}
          因此存在线性映射 $ \sigma : \mathbf{R}^3 \to \mathbf{R}^2 $ 满足题设条件，其关于 $ \mathbf{R}^3 $ 基的像为
          \begin{align*}
              \sigma(\vec{e}_1) & = \frac{1}{2} (\vec{\varepsilon}_1 + \vec{\varepsilon}_2) \\
              \sigma(\vec{e}_2) & = \frac{1}{2} (\vec{\varepsilon}_1 - \vec{\varepsilon}_2) \\
              \sigma(\vec{e}_3) & = (0, 0)
          \end{align*}
\end{solve}

\begin{exercise}\label{prob:L5-A-4}
是否存在$\mathbf{R}^2$到$\mathbf{R}^3$的线性映射$\sigma$使得$\sigma(3,2)=(1,0,0)$，$\sigma(1,5)=(1,1,0)$，$\sigma(-1,4)=(1,1,1)$？
\end{exercise}
\noindent 来源：\href{https://github.com/yhwu-is/Linear-Algebra-Left-Undone/blob/new/%E8%AE%B2%E4%B9%89/%E4%B8%93%E9%A2%98/5%20%E7%BA%BF%E6%80%A7%E6%98%A0%E5%B0%84.tex}{L5-A-4}。
\begin{solve}
不存在. 不存在从 $ \mathbf{R}^2 $ 到 $ \mathbf{R}^3 $ 的满射.
\end{solve}

\begin{exercise}\label{prob:L5-A-5}
求$\sigma(x_1,x_2,\ldots,x_n)=(x_1,0,\ldots,0)$的像、核与秩.
\end{exercise}
\noindent 来源：\href{https://github.com/yhwu-is/Linear-Algebra-Left-Undone/blob/new/%E8%AE%B2%E4%B9%89/%E4%B8%93%E9%A2%98/5%20%E7%BA%BF%E6%80%A7%E6%98%A0%E5%B0%84.tex}{L5-A-5}。
\begin{solve}
\[ \sigma(x_1, x_2, \ldots, x_n) = x_1(1, 0, \ldots, 0) \]
          所以
          \begin{gather*}
              \sigma(V) = \operatorname{span}(\vec{e}_1) \\
              r(\sigma) = 1
          \end{gather*}
          由于 $ \sigma(\vec{e}_1) = \vec{e}_1,\enspace \sigma(\vec{e}_2) = \cdots = \sigma(\vec{e}_n) = \vec{0} $，所以
          \begin{gather*}
              \ker \sigma = \operatorname{span}(\vec{e}_2, \vec{e}_3, \ldots, \vec{e}_n) \\
              \dim \ker \sigma = n - 1
          \end{gather*}
\end{solve}

\begin{exercise}\label{prob:L5-A-6}
证明：同构映射的逆、复合仍然是同构映射.
\end{exercise}
\noindent 来源：\href{https://github.com/yhwu-is/Linear-Algebra-Left-Undone/blob/new/%E8%AE%B2%E4%B9%89/%E4%B8%93%E9%A2%98/5%20%E7%BA%BF%E6%80%A7%E6%98%A0%E5%B0%84.tex}{L5-A-6}。
\begin{solve}
其逆映射是同构映射是显然的.我们只证明复合映射

            对于两个同构映射 $\sigma: V_1 \to V_2, \tau: V_2 \to V_3$，我们有其核空间
            $\ker \sigma \tau$,$\tau$是单的，将$0$映射到$0$，$\sigma$也是单的，故$\sigma\tau$是单的，将$0$映射到$0$；
            $\sigma$是满的，$\tau$也是满的，故$\sigma\tau$是满的，故$\sigma\tau$是双射，是同构映射.


            如果从矩阵的角度来理解：
            同构映射对应的矩阵是可逆矩阵；
            可逆矩阵的逆矩阵，可逆矩阵的乘积矩阵仍然是可逆矩阵.
\end{solve}

\begin{exercise}\label{prob:L5-B-1}
已知
        \begin{gather*}
            \alpha_1=(1,-1),\enspace\alpha_2=(2,-1),\enspace\alpha_3=(-3,2) \\
            \beta_1=(1,0),\enspace\beta_2=(0,1),\enspace\beta_3=(1,1)
        \end{gather*}
        是否存在$\sigma\in \mathcal{L}(\mathbf{R}^2,\mathbf{R}^2)$，使得$\sigma(\alpha_i)=\beta_i,\enspace i=1,2,3$？
\end{exercise}
\noindent 来源：\href{https://github.com/yhwu-is/Linear-Algebra-Left-Undone/blob/new/%E8%AE%B2%E4%B9%89/%E4%B8%93%E9%A2%98/5%20%E7%BA%BF%E6%80%A7%E6%98%A0%E5%B0%84.tex}{L5-B-1}。
\begin{solve}
不存在. 假设存在，则由 $ \alpha_1 + \alpha_2 + \alpha_3 = \vec{0} $ 有
          \begin{align*}
              \sigma(\vec{0}) & = \sigma(\alpha_1 + \alpha_2 + \alpha_3)                 \\
                              & = \sigma(\alpha_1) + \sigma(\alpha_2) + \sigma(\alpha_3) \\
                              & = \beta_1 + \beta_2 + \beta_3 = \vec{0}
          \end{align*}
          这与 $ \beta_1 + \beta_2 + \beta_3 = (2, 2) $ 矛盾.
\end{solve}

\begin{exercise}\label{prob:L5-B-2}
设$\alpha_1,\alpha_2$是线性空间$V(\mathbf{F})$的一组基，$x_1\alpha_1+x_2\alpha_2 \in V$. 定义$T(x_1\alpha_1+x_2\alpha_2)=r_1x_1\alpha_1+r_2x_2\alpha_2$，其中$r_1,r_2$是域$\mathbf{F}$中的两个常数. 证明：$T$是$V$上的一个线性变换. 当$V=\mathbf{R}^2$时，说明$T$的几何意义.
\end{exercise}
\noindent 来源：\href{https://github.com/yhwu-is/Linear-Algebra-Left-Undone/blob/new/%E8%AE%B2%E4%B9%89/%E4%B8%93%E9%A2%98/5%20%E7%BA%BF%E6%80%A7%E6%98%A0%E5%B0%84.tex}{L5-B-2}。
\begin{solve}
$ \forall x_1, x_2, y_1, y_2, k_1, k_2 \in \mathbf{F} $，有
          \begin{align*}
                  & T(k_1(x_1 \alpha_1 + x_2 \alpha_2) + k_2(y_1 \alpha_1 + y_2 \alpha_2))                \\
              ={} & r_1(k_1 x_1 + k_2 y_1) \alpha_1 + r_2(k_1 x_2 + k_2 y_2) \alpha_2                     \\
              ={} & k_1 (r_1 x_1 \alpha_1 + r_2 x_2 \alpha_2) + k_2 (r_1 y_1 \alpha_1 + r_2 y_2 \alpha_2) \\
              ={} & k_1 T(x_1 \alpha_1 + x_2 \alpha_2) + k_2 T(y_1 \alpha_1 + y_2 \alpha_2)
          \end{align*}
          因此 $ T $ 是线性映射. 其几何意义：将 $ \mathbf{R}^2 $ 中向量沿 $ x, y $ 轴分别变为原来的 $ r_1, r_2 $ 倍.
\end{solve}

\begin{exercise}\label{prob:L5-B-3}
已知$\mathbf{R}$上的线性变换$\sigma(x_1,x_2)=(x_1-x_2,x_1+x_2)$，$\tau(x_1,x_2)=(x_1-x_2,x_2-x_1)$.
        \begin{enumerate}
            \item 求$\sigma^2(x_1,x_2)$；

            \item $\sigma$是否可逆？如可逆，求$\sigma^{-1}(x_1,x_2)$；

            \item 求$\xi\in \mathcal{L}(\mathbf{R}^2,\mathbf{R}^2)$，使得$\xi\tau=\theta$（零变换）.
        \end{enumerate}
\end{exercise}
\noindent 来源：\href{https://github.com/yhwu-is/Linear-Algebra-Left-Undone/blob/new/%E8%AE%B2%E4%B9%89/%E4%B8%93%E9%A2%98/5%20%E7%BA%BF%E6%80%A7%E6%98%A0%E5%B0%84.tex}{L5-B-3}。
\begin{solve}
\begin{enumerate}
                \item \begin{align*}
                          \sigma^2(x_1, x_2) & = \sigma(\sigma(x_1, x_2))                               \\
                                             & = \sigma(x_1 - x_2, x_1 + x_2)                           \\
                                             & = ((x_1 - x_2) - (x_1 + x_2), (x_1 - x_2) + (x_1 + x_2)) \\
                                             & = (-2 x_2, 2 x_1)
                      \end{align*}

                \item $ \sigma $ 可逆的充分必要条件是存在线性变换 $ \tau $ 使得 $ \sigma \tau = I $. 于是，$ \forall \alpha \in \mathbf{R}^2 $，令 $ (\tau \sigma)(\alpha) = \alpha $，即
                      \[ (\tau \sigma)(x_1, x_2) = \tau(x_1 - x_2, x_1 + x_2) = \tau(y_1, y_2) = (x_1, x_2) \]
                      解得
                      \begin{gather*}
                          x_1 = \frac{y_1 + y_2}{2} \\
                          x_2 = \frac{y_1 - y_2}{2}
                      \end{gather*}
                      所以
                      \[ \tau(x_1, x_2) = \sigma^{-1}(x_1, x_2) = \left(\frac{x_1 + x_2}{2}, \frac{x_1 - x_2}{2}\right) \]

                \item 当 $ \xi = \theta $ 时，显然满足 $ \xi \tau = \theta $. 当 $ \xi \neq \theta $ 时，
                      \[ \xi \tau(x_1, x_2) = \xi(x_1 - x_2, x_2 - x_1) = (0, 0) \]
                      记 $ y_1 = x_1 - x_2,\enspace y_2 = x_2 - x_1 $. 由于 $ y_1 + y_2 = 0 $，
                      \[ \xi(y_1, y_2) = (y_1 + y_2, y_1 + y_2) \]
                      即 $ \xi(x_1, x_2) = (x_1 + x_2, x_1 + x_2) $ 满足 $ \xi \tau = \theta $.
            \end{enumerate}
\end{solve}

\begin{exercise}\label{prob:L5-B-4}
已知$\mathbf{R}^3$上的两个线性变换$\sigma,\tau$为：
        \begin{gather*}
            \sigma(x_1,x_2,x_3)=(x_3,0,0) \\
            \tau(x_1,x_2,x_3)=(x_1+x_2+x_3,x_1-x_2,0)
        \end{gather*}
        \begin{enumerate}
            \item 求$r(\sigma),\enspace r(\tau),\enspace \operatorname{im}\sigma,\enspace \ker\sigma$；

            \item 求$r(\tau\sigma),\enspace r(\sigma\tau),\enspace r(\sigma+\tau)$；

            \item 求$\operatorname{im}\tau+\ker\tau$.
        \end{enumerate}
\end{exercise}
\noindent 来源：\href{https://github.com/yhwu-is/Linear-Algebra-Left-Undone/blob/new/%E8%AE%B2%E4%B9%89/%E4%B8%93%E9%A2%98/5%20%E7%BA%BF%E6%80%A7%E6%98%A0%E5%B0%84.tex}{L5-B-4}。
\begin{solve}
\begin{enumerate}
                \item \begin{align*}
                          r(\sigma)   & = 1                          \\
                          r(\tau)     & = 2                          \\
                          \operatorname{im} \sigma  & = \operatorname{span}(\vec{e}_1)            \\
                          \ker \sigma & = \operatorname{span}(\vec{e}_1, \vec{e}_2)
                      \end{align*}

                \item \begin{gather*}
                          (\sigma \tau)(x_1, x_2, x_3) = \sigma(x_1 + x_2 + x_3, x_1 - x_2, 0) = (0, 0, 0) \\
                          r(\sigma \tau) = 0
                      \end{gather*}
                      \begin{gather*}
                          (\tau \sigma)(x_1, x_2, x_3) = \tau(x_3, 0, 0) = (x_3, x_3, 0) \\
                          \tau \sigma \neq \theta = \sigma \tau
                      \end{gather*}
                      求 $ r(\tau \sigma) $，方法 1：
                      \begin{gather*}
                          (\tau \sigma)(V) = \operatorname{span}((1, 1, 0)) \\
                          r(\tau \sigma) = 1
                      \end{gather*}
                      方法 2：由 $ \tau \sigma \neq \theta $ 知
                      \[ 1 \leqslant r(\tau \sigma) \leqslant \min\{r(\tau), r(\sigma)\} = 1 \]
                      因此 $ r(\tau \sigma) = 1 $.

                      \begin{gather*}
                          \begin{aligned}
                              (\sigma + \tau)(x_1, x_2, x_3) & = \sigma(x_1, x_2, x_3) + \tau(x_1, x_2, x_3) \\
                                                             & = (x_1 + x_2 + 2 x_3, x_1 - x_2, 0)           \\
                                                             & = x_1(1, 1, 0) + x_2(1, -1, 0) + x_3(2, 0, 0)
                          \end{aligned} \\
                          r(\sigma + \tau) = r\{(1, 1, 0), (1, -1, 0), (2, 0, 0)\} = 2
                      \end{gather*}

                \item \begin{gather*}
                          \begin{aligned}
                              \operatorname{im} \tau & = \operatorname{span}((1, 1, 0), (1, -1, 0), (2, 0, 0)) \\
                                       & = \operatorname{span}((1, 1, 0), (1, -1, 0))
                          \end{aligned} \\
                          \ker \tau = \operatorname{span}((1, 1, -2))
                      \end{gather*}
                      可得
                      \[ \operatorname{im} \tau + \ker \tau = \operatorname{span}((1, 1, 0), (1, -1, 0), (1, 1, -2)) = \mathbf{R}^3 \]
            \end{enumerate}
\end{solve}

\begin{exercise}\label{prob:L5-B-5}
设 $\sigma$ 是线性空间 $V$ 上的线性变换，如果 $\sigma^{k-1}(\alpha) \neq 0$，但 $\sigma^{k}(\alpha) = 0$，证明：\\
        $\alpha,\sigma(\alpha),\ldots,\sigma^{k-1}(\alpha)\enspace(k>0)$ 线性无关（本题还有对应的矩阵版本，解法基本一致）.
\end{exercise}
\noindent 来源：\href{https://github.com/yhwu-is/Linear-Algebra-Left-Undone/blob/new/%E8%AE%B2%E4%B9%89/%E4%B8%93%E9%A2%98/5%20%E7%BA%BF%E6%80%A7%E6%98%A0%E5%B0%84.tex}{L5-B-5}。
\begin{solve}
使用反证法. 假设 $ \alpha, \sigma(\alpha), \ldots, \sigma^{k - 1}(\alpha) $ 线性相关，则存在不全为 0 的 $ l_0, l_1, \ldots, l_{k - 1} $ 使得
          \[ l_0 \alpha + l_1 \sigma(\alpha) + \cdots + l_{k - 1} \sigma^{k - 1}(\alpha) = \vec{0} \]
          记 $ i = \min\{k \mid l_k \neq 0\} $，则
          \[ l_i \sigma^{i}(\alpha) + l_{i + 1} \sigma^{i + 1}(\alpha) + \cdots + l_{k - 1} \sigma^{k - 1}(\alpha) = \vec{0} \]
          两边同时作用 $ \sigma^{k - i - 1} $. 由于 $ \sigma^{k}(\alpha) = \vec{0} $，故 $ \sigma^{k + 1}(\alpha) = \sigma^{k + 2}(\alpha) = \cdots = \vec{0} $，上式成为
          \[ l_i \sigma^{k - 1}(\alpha) = \vec{0} \]
          而 $ l_i \neq 0,\enspace \sigma^{k - 1}(\alpha) \neq \vec{0} $，矛盾！
\end{solve}

\begin{exercise}\label{prob:L5-B-6}
设$\mathbf{R}[x]_3$是次数小于3的实系数多项式和全体零多项式一起组成的集合关于多项式加法和数乘多项式运算构成的实数域上的线性空间.
        \begin{enumerate}
            \item 证明：$W=\{f(x)\in \mathbf{R}[x]_3 \mid f(1)=0\}$是$\mathbf{R}[x]_3$的一个子空间，并求$W$的维数和一组基；

            \item 定义从$\mathbf{R}[x]_3$到$\mathbf{R}$的线性映射$\sigma(f(x))=f(1)$，证明：$\sigma$为线性映射，并求$\operatorname{im}\sigma$和$\dim\ker\sigma$；

            \item 设$f,g,h \in \mathbf{R}[x]_3$且$f(1)=g(1)=h(1)=0$，证明：$f,g,h$线性相关.
        \end{enumerate}
\end{exercise}
\noindent 来源：\href{https://github.com/yhwu-is/Linear-Algebra-Left-Undone/blob/new/%E8%AE%B2%E4%B9%89/%E4%B8%93%E9%A2%98/5%20%E7%BA%BF%E6%80%A7%E6%98%A0%E5%B0%84.tex}{L5-B-6}。
\begin{solve}
\begin{enumerate}
\item 零多项式在 $W$ 中；若 $f,g\in W$ 且 $a,b\in\mathbf R$，则 $(af+bg)(1)=af(1)+bg(1)=0$，所以 $W$ 是子空间。任意 $f=ax^2+bx+c\in W$ 满足 $c=-a-b$，从而 $f=a(x^2-1)+b(x-1)$。两个生成多项式比较最高次项即可证明线性无关，故它们是一组基且 $\dim W=2$。
\item 对任意标量 $a,b$，$\sigma(af+bg)=a\sigma(f)+b\sigma(g)$，故 $\sigma$ 线性。常值多项式可取任意实数，故 $\operatorname{im}\sigma=\mathbf R$；$\ker\sigma=W$，故 $\dim\ker\sigma=2$。
\item 三个多项式均属于二维空间 $W$，而其中任意三个向量必线性相关。\end{enumerate}
\end{solve}

\begin{exercise}\label{prob:L5-B-7}
设$\sigma(p(x))=p'(x)$（求导），$\forall p(x) \in \mathbf{R}[x]_n$.
        \begin{enumerate}
            \item 证明：$\sigma$是$\mathbf{R}[x]_n$上的线性变换；

            \item 求$\sigma$的值域和$r(\sigma)$，说明$\sigma$是否可逆；

            \item 求$\sigma$的核及其维数；

            \item 求$r(\sigma)+\dim\ker\sigma$，问：$\mathbf{R}[x]_n=\ker\sigma+\operatorname{im} \sigma$是否成立.
        \end{enumerate}
\end{exercise}
\noindent 来源：\href{https://github.com/yhwu-is/Linear-Algebra-Left-Undone/blob/new/%E8%AE%B2%E4%B9%89/%E4%B8%93%E9%A2%98/5%20%E7%BA%BF%E6%80%A7%E6%98%A0%E5%B0%84.tex}{L5-B-7}。
\begin{solve}
\begin{enumerate}
                \item 证明：$ \forall p_1(x), p_2(x) \in \mathbf{R}[x]_n,\enspace \forall k_1, k_2 \in \mathbf{R} $，有
                      \begin{align*}
                          \sigma(k_1 p_1 + k_2 p_2) & = (k_1 p_1(x) + k_2 p_2(x))'                                    \\
                                                    & = k_1 p_1'(x) + k_2 p_2'(x) = k_1 \sigma(p_1) + k_2 \sigma(p_2)
                      \end{align*}
                      因此 $ \sigma $ 是 $ \mathbf{R}[x]_n $ 上的线性变换.

                \item \begin{gather*}
                          \begin{aligned}
                              \operatorname{im} \sigma & = \operatorname{span}(\sigma(1), \sigma(x), \sigma(x^2), \ldots, \sigma(x^{n - 1})) \\
                                         & = \operatorname{span}(1, 2x, \ldots, (n - 1) x^{n - 2})                             \\
                                         & = \operatorname{span}(1, x, \ldots, x^{n - 2})
                          \end{aligned} \\
                          r(\sigma) = n - 1
                      \end{gather*}
                      可知 $ \sigma $ 不是单射，因此不可逆.

                \item 由 $ \sigma(p(x)) = 0 $ 可知 $ p(x) = c $（常数）. 因此 $ \ker \sigma = \operatorname{span}(1),\enspace \dim \ker \sigma = 1 $.

                \item \begin{gather*}
                          r(\sigma) + \dim \ker \sigma = (n - 1) + 1 = n \\
                          \operatorname{im} \sigma + \ker \sigma = \operatorname{span}(1, x, \ldots, x^{n - 2}) = \operatorname{im} \sigma \neq \mathbf{R}[x]_n
                      \end{gather*}
            \end{enumerate}
\end{solve}

\begin{exercise}\label{prob:L5-B-8}
设$V$为有限维线性空间，$T\in \mathcal{L}(V,V)$且$T$不是恒等变换也不是零变换，问：下列情况是否可能发生？如果可能请举例，不可能请说明理由.
        \begin{enumerate}
            \item $\operatorname{im} T \cap \ker T = \{0\}$；

            \item $\operatorname{im} T \subseteq \ker T$；

            \item $\ker T = \operatorname{im} T$；

            \item $\ker T \subseteq \operatorname{im} T$.
        \end{enumerate}
\end{exercise}
\noindent 来源：\href{https://github.com/yhwu-is/Linear-Algebra-Left-Undone/blob/new/%E8%AE%B2%E4%B9%89/%E4%B8%93%E9%A2%98/5%20%E7%BA%BF%E6%80%A7%E6%98%A0%E5%B0%84.tex}{L5-B-8}。
\begin{solve}
\begin{enumerate}
                \item 可能. 例如 $ \sigma(x, y) = (x + y, x + y) $.
                      \begin{gather*}
                          \operatorname{im} \sigma = \operatorname{span}(\vec{e}_1 + \vec{e}_2) \\
                          \ker \sigma = \operatorname{span}(\vec{e}_1 - \vec{e}_2) \\
                          \operatorname{im} \sigma \cap \ker \sigma = \{\vec{0}\}
                      \end{gather*}

                \item 可能. 例如 $ \sigma(x, y, z) = (x - y, x - y, x - y) $.
                      \begin{gather*}
                          \operatorname{im} \sigma = \operatorname{span}((1, 1, 1)) \\
                          \ker \sigma = \operatorname{span}((1, 1, 1), (1, 1, 0)) \\
                          \operatorname{im} \sigma \subseteq \ker \sigma
                      \end{gather*}

                \item 可能. 例如 $ \sigma(x, y) = (x - y, x - y) $.
                      \begin{gather*}
                          \operatorname{im} \sigma = \ker \sigma = \operatorname{span}((1, 1))
                      \end{gather*}

                \item 可能. 例如 $ \sigma(x, y) = (x, x - y) $.
                      \begin{gather*}
                          \operatorname{im} \sigma = \mathbf{R}^2 \\
                          \ker \sigma = \{\vec{0}\} \\
                          \ker \sigma \subseteq \operatorname{im} \sigma
                      \end{gather*}
            \end{enumerate}
\end{solve}

\begin{exercise}\label{prob:L5-B-9}
若$\sigma_1,\sigma_2\in \mathcal{L}(V,V)$，判断下列说法是否正确，正确请给出证明，反之给出反例：
        \begin{enumerate}
            \item 由$r(\sigma)+\dim\ker\sigma=n$可知$V=\ker\sigma+\operatorname{im} \sigma$；

            \item 若有$\operatorname{im} \sigma \cap \ker \sigma = \{0\}$，则$V=\ker\sigma+\operatorname{im} \sigma$成立；

            \item 因为$\forall \alpha \in V$有$(\sigma_1+\sigma_2)(\alpha)=\sigma_1(\alpha)+\sigma_2(\alpha)$，所以$(\sigma_1+\sigma_2)(V)=\sigma_1(V)+\sigma_2(V)$；

            \item $(I-\sigma)(V)+\sigma(V)=V$，其中$I$为恒等映射.
        \end{enumerate}
\end{exercise}
\noindent 来源：\href{https://github.com/yhwu-is/Linear-Algebra-Left-Undone/blob/new/%E8%AE%B2%E4%B9%89/%E4%B8%93%E9%A2%98/5%20%E7%BA%BF%E6%80%A7%E6%98%A0%E5%B0%84.tex}{L5-B-9}。
\begin{solve}
\begin{enumerate}
                \item 错误. 一个反例为 $ \sigma(x_1, x_2) = (x_1 - x_2, x_1 - x_2) $，则 $ \operatorname{im} \sigma + \ker \sigma = \operatorname{span}((1, 1)) $.

                \item 正确. $ \operatorname{im} \sigma + \ker \sigma \subseteq V $，此时 $ \dim(\operatorname{im} \sigma + \ker \sigma) = \dim \operatorname{im} \sigma + \dim \ker \sigma - 0 = \dim V $，所以 $ \operatorname{im} \sigma + \ker \sigma = V $.

                \item \label{item:6:B:2:3}
                      错误. $ \forall \alpha \in V $ 有
                      \[ (\sigma_1 + \sigma_2)(\alpha) = \sigma_1(\alpha) + \sigma_2(\alpha) \in \sigma_1(V) + \sigma_2(V) \]
                      所以
                      \[ (\sigma_1 + \sigma_2)(V) \subseteq \sigma_1(V) + \sigma_2(V) \]
                      但上式中等号不一定成立. 反例：$ V = \mathbf{R}^3 $ 上的线性变换 $ \sigma_1, \sigma_2 $ 关于 $ \mathbf{R}^3 $ 的基 $ \vec{e}_1, \vec{e}_2, \vec{e}_3 $ 的像分别为
                      \begin{gather*}
                          \sigma_1(\vec{e}_1) = \vec{e}_1,\enspace \sigma_1(\vec{e}_2) = \sigma_1(\vec{e}_3) = \vec{e}_2 \\
                          \sigma_2(\vec{e}_1) = \vec{e}_1,\enspace \sigma_2(\vec{e}_2) = \sigma_2(\vec{e}_3) = \vec{e}_3 \\
                          \sigma_1(V) = \operatorname{span}(\vec{e}_1, \vec{e}_2),\enspace \sigma_2(V) = \operatorname{span}(\vec{e}_1, \vec{e}_3) \\
                          \sigma_1(V) + \sigma_2(V) = \operatorname{span}(\vec{e}_1, \vec{e}_2, \vec{e}_3) = \mathbf{R}^3
                      \end{gather*}
                      而
                      \begin{gather*}
                          (\sigma_1 + \sigma_2)(\vec{e}_1) = \vec{e}_1 + \vec{e}_1 = 2 \vec{e}_1 \\
                          (\sigma_1 + \sigma_2)(\vec{e}_2) = (\sigma_1 + \sigma_2)(\vec{e}_3) = \vec{e}_2 + \vec{e}_3 \\
                          (\sigma_1 + \sigma_2)(V) = \operatorname{span}(\vec{e}_1, \vec{e}_2 + \vec{e}_3) \neq \mathbf{R}^3
                      \end{gather*}

                \item 正确. $(I - \sigma)(V) + \sigma(V) \subseteq V$ 是显然的，另一个方向上，$ \forall \alpha \in V $ 有
                      \[ \alpha = (I - \sigma)(\alpha) + \sigma(\alpha) \in (I - \sigma)(V) + \sigma(V) \]
                      所以 $ V \subseteq (I - \sigma)(V) + \sigma(V) $. 综上有 $ V = (I - \sigma)(V) + \sigma(V) $.
            \end{enumerate}
\end{solve}

\begin{exercise}\label{prob:L5-B-10}
设$V(\mathbf{R})$是线性空间，$\sigma$是$V(\mathbf{R})$到$\mathbf{R}^3$的同构映射，且$\sigma(\alpha_1)=(1,0,1),\enspace\allowbreak\sigma(\alpha_2)=(-2,1,0),\enspace\allowbreak\sigma(\alpha_3)=(-3,2,1),\enspace\allowbreak\sigma(\alpha_4)=(1,1,2)$.
        \begin{enumerate}
            \item $\alpha_1$在$\operatorname{span}(\alpha_2,\alpha_3)$中吗？

            \item 设$W_1=\operatorname{span}(\alpha_1,\alpha_2),\enspace W_2=\operatorname{span}(\alpha_3,\alpha_4)$，求$W_1\cap W_2$.
        \end{enumerate}
\end{exercise}
\noindent 来源：\href{https://github.com/yhwu-is/Linear-Algebra-Left-Undone/blob/new/%E8%AE%B2%E4%B9%89/%E4%B8%93%E9%A2%98/5%20%E7%BA%BF%E6%80%A7%E6%98%A0%E5%B0%84.tex}{L5-B-10}。
\begin{solve}
\begin{enumerate}
                \item 在. 因为
                      \[ \sigma(\alpha_1) = -2 \sigma(\sigma_2) + \sigma(\alpha_3) = \sigma(-2 \alpha_2 + \alpha_3) \]
                      同构映射 $ \sigma $ 可逆. 所以
                      \[ \alpha_1 = -2 \alpha_2 + \alpha_3 \in \operatorname{span}(\alpha_2, \alpha_3) \]
                \item 对于注意力惊人的同学，直接观察当然可得答案；如果没什么思路，我们可以先从定义入手；
                由于同构是将基映射到基的，我们可以将对于$\alpha$的操作看做对于其像的操作，为了方便，在这里直接用 $a$ 来代替$\alpha$的像，若 $\beta \in W_1\cap W_2$，则
                \[ \beta = x_1 a_1+x_2 a_2=x'_3 a_3+x'_4 a_4 \implies  x_1 a_1+x_2 a_2-x'_3 a_3-x'_4 a_4 = 0 \]
                改写$x'_3,x'_4$为$x_3,x_4$，则转换为
                \[ x_1 a_1+x_2 a_2+x_3 a_3+x_4 a_4 = 0 \]
                对应的系数矩阵为
                \[ \begin{pmatrix}
                    1 & -2 & -3 & 1 \\
                    0 & 1 & 2 & 1\\
                    1 & 0 & 1 & 2
                \end{pmatrix} \]
                四个未知数，三个方程，解得其基础解系为
                \[ \begin{pmatrix}
                    -1 \\
                    -2 \\
                    1 \\
                    0
                \end{pmatrix} \]

                故$\beta = k(-a_1+-2a_2=-a_3)$其中$k\in \mathbf{R}$.故所求交空间的基为$\alpha_3$(注意要对应回$V(R)$中)
            \end{enumerate}
\end{solve}

\begin{exercise}\label{prob:L5-B-11}
设$c_1,c_2,\ldots,c_n$是$n$个互异的实常数. 证明：$\mathbf{R}[x]_n$到$\mathbf{R}^n$的一个映射$\sigma$：
        \[\sigma(p(x))=(p(c_1),p(c_2),\ldots,p(c_n))\]
        是$\mathbf{R}[x]_n$到$\mathbf{R}^n$的一个同构映射.
\end{exercise}
\noindent 来源：\href{https://github.com/yhwu-is/Linear-Algebra-Left-Undone/blob/new/%E8%AE%B2%E4%B9%89/%E4%B8%93%E9%A2%98/5%20%E7%BA%BF%E6%80%A7%E6%98%A0%E5%B0%84.tex}{L5-B-11}。
\begin{solve}
我们仅对 $ n = 3 $ 的情况给出证明. 

          先证 $ \sigma $ 是线性映射. $ \forall p(x), q(x) \in \mathbf{R}[x]_3,\enspace \forall k_1, k_2 \in \mathbf{R} $ 有
          \begin{align*}
                  & \sigma(k_1 p(x) + k_2 q(x))                                                 \\
              ={} & (k_1 p(c_1) + k_2 q(c_1), k_1 p(c_2) + k_2 q(c_2), k_1 p(c_3) + k_2 q(c_3)) \\
              ={} & k_1(p(c_1), p(c_2), p(c_3)) + k_1(q(c_1), q(c_2), q(c_3))                   \\
              ={} & k_1 \sigma(p(x)) + k_2 \sigma(q(x))
          \end{align*}

          再证 $ \sigma $ 是双射，即 $ \forall (d_1, d_2, d_3) \in \mathbf{R}^3 $，存在唯一的
          \[ p(x) = a + bx + cx^2 \in \mathbf{R}[x]_3 \]
          使 $ \sigma(p(x)) = (d_1, d_2, d_3) $. 根据
          \[ \sigma(p(x)) = (p(c_1), p(c_2), p(c_3)) \]
          以及 $ \sigma(p(x)) = (d_1, d_2, d_3) $，有
          \[ \begin{cases}
                  a + bc_1 + cc_1^2 = d_1 \\
                  a + bc_2 + cc_2^2 = d_2 \\
                  a + bc_3 + cc_3^3 = d_3
              \end{cases} \]
          方程组是关于未知元 $ a, b, c $ 的三元线性非齐次方程组，其中 $ c_1, c_2, c_3 $ 是互异的实常数. 用高斯-若当消元法，易将其增广矩阵变换为下列阶梯形矩阵，即
          \begin{align}
                             & \begin{pmatrix}[ccc|c]
                                   1 & c_1 & c_1^2 & d_1 \\
                                   1 & c_2 & c_2^2 & d_2 \\
                                   1 & c_3 & c_3^3 & d_3
                               \end{pmatrix} \notag                                                                    \\
              \implies \quad & \begin{pmatrix}[ccc|c]
                                   1 & c_1 & c_1^2     & d_1                                                         \\[1ex]
                                   0 & 1   & c_2 + c_1 & \dfrac{d_1 - d_2}{c_1 - c_2}                                \\[2ex]
                                   0 & 0   & c_3 - c_2 & \dfrac{d_3 - d_1}{c_3 - c_1} - \dfrac{d_2 - d_1}{c_2 - c_1}
                               \end{pmatrix} \label{item:6:B:6:1}
          \end{align}
          阶梯性矩阵~\eqref{item:6:B:6:1} （其中 $ c_3 - c_2, c_3 - c_1, c_2 - c_1 $ 均为非零常数）对应的方程组有唯一解 $ a, b, c $，即存在唯一的
          \[ p(x) = a + bx + cx^2 \in \mathbf{R}[x]_3 \]
          使得 $ \sigma(p(x)) = (d_1, d_2, d_3) $ 成立. 所以 $ \sigma $ 是线性双射，即 $ \mathbf{R}[x]_3 $ 到 $ \mathbf{R}^3 $ 的同构映射.
\end{solve}

\begin{exercise}\label{prob:L5-B-12}
设$\sigma$和$\tau$分别为有限维线性空间$U\to V$和$V\to W$的线性映射，证明
        \[\dim\ker\sigma+\dim(\operatorname{im}\sigma\cap\ker\tau)=\dim\ker(\tau\sigma).\]
\end{exercise}
\noindent 来源：\href{https://github.com/yhwu-is/Linear-Algebra-Left-Undone/blob/new/%E8%AE%B2%E4%B9%89/%E4%B8%93%E9%A2%98/5%20%E7%BA%BF%E6%80%A7%E6%98%A0%E5%B0%84.tex}{L5-B-12}。
\begin{solve}
我们先观察等式的右边，在 $\ker \tau \sigma$中的元素 $u$,会满足 $\tau \sigma (u) = 0 $,仔细想想什么元素会满足这个式子
            \begin{itemize}
                \item $\sigma(u)=0$的元素，亦即$u\in \ker \sigma$；
                \item $v=\sigma(u)\neq 0$，但是$\tau(v)=0$的元素，亦即$v\in \operatorname{im} \sigma \cap \ker \tau$.
            \end{itemize}
            这恰好对应左边的两项，所以等式显然是成立的.
\end{solve}

\begin{exercise}\label{prob:L5-C-1}
设 $V(\mathbf{F})$ 是一个 $n$ 维线性空间，$\sigma \in \mathcal{L}(V,V)$. 证明：
        \begin{enumerate}
            \item 在 $\mathbf{F}[x]$ 中有一个次数不高于 $n^2$ 的非零多项式 $p(x)$ 使 $p(\sigma) = \theta$；

            \item $\sigma$ 可逆$\iff$有一常数项不为 0 的多项式 $p(x)$ 使 $p(\sigma) = \theta$.
        \end{enumerate}
\end{exercise}
\noindent 来源：\href{https://github.com/yhwu-is/Linear-Algebra-Left-Undone/blob/new/%E8%AE%B2%E4%B9%89/%E4%B8%93%E9%A2%98/5%20%E7%BA%BF%E6%80%A7%E6%98%A0%E5%B0%84.tex}{L5-C-1}。
\begin{solve}
\begin{enumerate}
                \item \label{item:5:C:1:1}
                      $ \mathcal{L}(V, V) $ 是 $ n^2 $ 维线性空间，其中 $ n^2 + 1 $ 个元素 $ I, \sigma, \sigma^2, \ldots, \sigma^{n^2} $ 线性相关，即 $ \exists a_i\enspace (i = 1, 2, \ldots, n^2) $ 使得
                      \[ a_0 I + a_1 \sigma + \ldots + a_{n^2} \sigma^{n^2} = \theta \]
                      于是 $ \exists p(x) = a_0 + a_1 x + \ldots + a_{n^2} x^{n^2} \in \mathbf{F}[x] $ 使得 $ p(\sigma) = \theta $.

                \item 必要性：设有一常数项不为 0 的多项式
                      \[ p(x) = a_0 + a_1 x + \cdots + a_k x^k \qquad a_0 \neq 0 \]
                      满足
                      \[ p(\sigma) = a_0 I + a_1 \sigma + \cdots + a_k \sigma^k = \theta \]
                      所以
                      \[ \sigma(a_1 I + a_2 \sigma + \cdots + a_k \sigma^{k - 1}) = -a_0 I \]
                      因此
                      \[ \sigma^{-1} = -a_0^{-1} (a_1 I + a_2 \sigma + \cdots + a_k \sigma^{k - 1}) \]

                      充分性：由 \eqref{item:5:C:1:1} 可知存在次数 $ k \leqslant n^2 $ 的多项式
                      \[ p(x) = a_0 + a_1 x + a_2 x^2 + \cdots + a_k x^k \]
                      满足
                      \[ p(\sigma) = a_0 I + a_1 \sigma + \cdots + a_k \sigma^k = \theta \]
                      若 $ a_0 \neq 0 $，则 $ p(x) $ 即为所求多项式.

                      若 $ a_0 = a_1 = \cdots = a_{i - 1} = 0,\enspace a_i \neq 0 $，即
                      \[ a_i I + a_{i + 1} \sigma + \cdots + a_k \sigma^{k - i} = \theta \]
                      于是
                      \[ P(x) = a_i + a_{i + 1} x + \cdots + a_k x^{k - i} \qquad a_i \neq 0 \]
                      为所求的多项式.
            \end{enumerate}
\end{solve}

\begin{exercise}\label{prob:L5-C-2}
已知$\sigma_1,\sigma_2,\ldots,\sigma_s$是线性空间$V$上的$s$个两两不同的线性变换，证明：在$V$中必存在向量$\alpha$使得$\sigma_1(\alpha),\sigma_2(\alpha),\ldots,\sigma_s(\alpha)$也两两不同.
\end{exercise}
\noindent 来源：\href{https://github.com/yhwu-is/Linear-Algebra-Left-Undone/blob/new/%E8%AE%B2%E4%B9%89/%E4%B8%93%E9%A2%98/5%20%E7%BA%BF%E6%80%A7%E6%98%A0%E5%B0%84.tex}{L5-C-2}。
\begin{solve}
$\alpha \in V$使得$\sigma_i(\alpha) \neq \sigma_j(\alpha)$当且仅当$\alpha \notin \ker(\sigma_i - \sigma_j)$,我们可以只考虑该核空间不是零空间的情况，因为我们所取出的$\alpha$显然不可能是零，
            又因为我们的线性映射是两两不相同的，所以该核空间一定是一个真子空间. 根据覆盖定理，对于$m$个非平凡子空间，我们总可以取到一个不属于任意一个这些子空间的向量，这样我们就找到了所需要的$\alpha$.
\end{solve}

\begin{exercise}\label{prob:L5-C-3}
设$V$是一个$n$维线性空间，$V=W_1\oplus W_2,\enspace\sigma\in \mathcal{L}(V,V)$. 证明：$\sigma$可逆$\iff V=\sigma(W_1)+\sigma(W_2)$.
\end{exercise}
\noindent 来源：\href{https://github.com/yhwu-is/Linear-Algebra-Left-Undone/blob/new/%E8%AE%B2%E4%B9%89/%E4%B8%93%E9%A2%98/5%20%E7%BA%BF%E6%80%A7%E6%98%A0%E5%B0%84.tex}{L5-C-3}。
\begin{solve}
证明：必要性：$ \forall \alpha \in V $，由 $ \sigma $ 可逆，存在唯一的 $ \beta \in V $ 使得 $ \sigma(\beta) = \alpha $ 且 $ \beta = \beta_1 + \beta_2 $，其中 $ \beta_1 \in W_1 $，$ \beta_2 \in W_2 $. 于是
          \[ \alpha = \sigma(\beta) = \sigma(\beta_1) + \sigma(\beta_2) \in \sigma(W_1) + \sigma(W_2) \]
          所以 $ V \subseteq \sigma(W_1) + \sigma(W_2) $.

          $ \sigma(W_1), \sigma(W_2) $ 都是 $ V $ 的子空间，它们的和也是 $ V $ 的子空间. 所以 $ \sigma(W_1) + \sigma(W_2) \subseteq V $，故 $ V = \sigma(W_1) + \sigma(W_2) $.

          充分性：$ \forall \alpha \in V = \sigma(W_1) + \sigma(W_2) $，$ \exists \alpha_i \in \sigma(W_i) $ 且 $ \exists \beta_i \in W_i $ 使 $ \alpha_i = \sigma(\beta_i) $（$ i = 1, 2 $）使得
          \[ \alpha = \alpha_1 + \alpha_2 = \sigma(\beta_1) + \sigma(\beta_2) = \sigma(\beta_1 + \beta_2) = \sigma(\beta) \]
          其中 $ \beta = \beta_1 + \beta_2 \in W_1 + W_2 = V $，所以 $ \sigma $ 为满射.

          由于 $ n $ 维线性空间上的线性映射为满射时也必为单射，从而必是双射，所以 $ \sigma $ 可逆.
\end{solve}

\begin{exercise}\label{prob:L5-C-5}
设$V_1$是有限维线性空间，$\sigma,\tau\in \mathcal{L}(V_1,V_2)$，则
        \[r(\sigma+\tau) \leqslant r(\sigma)+r(\tau).\]

        事实上前两题的结论在后续章节矩阵的秩中都会涉及，此处有兴趣的同学可以尝试从线性映射的角度理解这两个秩不等式. 如果无法找到合适方式，可以考虑化为矩阵进行证明.
\end{exercise}
\noindent 来源：\href{https://github.com/yhwu-is/Linear-Algebra-Left-Undone/blob/new/%E8%AE%B2%E4%B9%89/%E4%B8%93%E9%A2%98/5%20%E7%BA%BF%E6%80%A7%E6%98%A0%E5%B0%84.tex}{L5-C-5}。
\begin{solve}
证明：由于 $ \forall \beta \in (\sigma + \tau)(V_1),\enspace \exists \alpha \in V_1 $ 使 $ \beta = (\sigma + \tau)(\alpha) = \sigma(\alpha) + \tau(\alpha) \in \sigma(V_1) + \tau(V_1) $，所以
          \[ (\sigma + \tau)(V_1) \subseteq \sigma(V_1) + \tau(V_1) \]
          因此
          \begin{align*}
              \dim(\sigma + \tau)(V_1) & \leqslant \dim(\sigma(V_1) + \tau(V_1))     \\
                                       & \leqslant \dim \sigma(V_1) + \dim \tau(V_1)
          \end{align*}
\end{solve}

\begin{exercise}\label{prob:L5-C-6}
设$\sigma\in \mathcal{L}(V,V)$，$\dim V=n$，且$\sigma^2=\sigma$，$I$是$V$上的恒等变换. 证明：
        \begin{enumerate}
            \item $(I-\sigma)(V) \subseteq \ker\sigma$；

            \item $r(I-\sigma)+r(\sigma)=n$.
        \end{enumerate}
\end{exercise}
\noindent 来源：\href{https://github.com/yhwu-is/Linear-Algebra-Left-Undone/blob/new/%E8%AE%B2%E4%B9%89/%E4%B8%93%E9%A2%98/5%20%E7%BA%BF%E6%80%A7%E6%98%A0%E5%B0%84.tex}{L5-C-6}。
\begin{solve}
\begin{enumerate}
                \item \label{item:5:C:6:1}
                      $ \forall \sigma \in \mathcal{L}(V, V) $，则 $ I - \sigma \in \mathcal{L}(V, V) $. $ \forall \alpha \in (I - \sigma)(V),\enspace \exists \beta \in V $，有
                      \begin{gather*}
                          \alpha = (I - \sigma)(\beta) = \beta - \sigma(\beta) \\
                          \sigma(\alpha) = \sigma(\beta - \sigma(\beta)) = \sigma(\beta) - \sigma^2(\beta)
                      \end{gather*}
                      而由于 $ \sigma^2 = \sigma $，所以 $ \sigma(\alpha) = \vec{0} $，于是 $ \alpha \in \ker \sigma $，因此 $ (I - \sigma)(V) \subseteq \ker \sigma $.

                \item 利用 $ r(\sigma + \tau) \leqslant r(\sigma) + r(\tau) $ 和 $ r(\sigma) + \dim \ker \sigma = n $，由 \eqref{item:5:C:6:1} 可得
                      \begin{equation} \label{eq:5:C:6:2:1}
                          r(I - \sigma) + r(\sigma) \leqslant n
                      \end{equation}
                      又因为
                      \begin{equation} \label{eq:5:C:6:2:2}
                          r(I - \sigma) + r(\sigma) \geqslant r(I - \sigma + \sigma) = r(I) = n
                      \end{equation}
                      于是由式~\eqref{eq:5:C:6:2:1} 和式~\eqref{eq:5:C:6:2:2} 即可得到 $ r(I - \sigma) + r(\sigma) = n $.
            \end{enumerate}
\end{solve}

\begin{exercise}\label{prob:L5-C-7}
设$V$是一个$n$维线性空间，$\sigma\in \mathcal{L}(V,V)$，且$\sigma^2=\theta$（零映射）. 证明：
        \begin{enumerate}
            \item $ \dim \operatorname{im} \sigma \leqslant \dfrac{n}{2} $；

            \item 设$A$是$\sigma$在某组基下的矩阵，则方程组$AX=0$的基础解系至少有$\dfrac{n}{2}$个解.
        \end{enumerate}
\end{exercise}
\noindent 来源：\href{https://github.com/yhwu-is/Linear-Algebra-Left-Undone/blob/new/%E8%AE%B2%E4%B9%89/%E4%B8%93%E9%A2%98/5%20%E7%BA%BF%E6%80%A7%E6%98%A0%E5%B0%84.tex}{L5-C-7}。
\begin{solve}
证明：\begin{enumerate}
                \item \label{item:5:C:7:1}
                      $ \forall \alpha \in \operatorname{im} \sigma,\enspace \exists \beta \in V $ 使得 $ \sigma(\beta) = \alpha $. 由 $ \sigma^2 = \theta $ 可得 $ \sigma(\alpha) = \sigma^2(\beta) = \vec{0} $，因此 $ \alpha \in \ker \sigma $，从而 $ \operatorname{im} \sigma \subseteq \ker \sigma $. 于是我们得到
                      \[ n = \dim \operatorname{im} \sigma + \dim \ker \sigma \geqslant 2 \dim \operatorname{im} \sigma \]
                      即 $ \dim \operatorname{im} \sigma \leqslant \dfrac{n}{2} $.

                \item 由 \eqref{item:5:C:7:1} 可知，方程组 $ AX = \vec{0} $ 的基础解系含有 $ n - r(A) = \dim \ker \sigma \geqslant \dfrac{n}{2} $ 个解向量，所以结论成立.
            \end{enumerate}
\end{solve}

\begin{exercise}\label{prob:L5-例-4}
写出下列映射的出发空间和到达空间，并判断其是否为线性映射：
    \begin{enumerate}
        \item $\sigma(x_1,x_2)=(x_1-x_2,x_1,x_1+x_2)$；

        \item $\sigma(x_1,x_2)=(x_1x_2,x_1+x_2)$；

        \item $\sigma(p(x))=p(x+1)-p(x),\enspace\forall p(x) \in \mathbf{R}[x]_n$；

        \item $\sigma(p(x))=p(a),\enspace\forall p(x)$，其中$a$为常数；

        \item $\sigma(\xi)=2\xi+\xi_0$，其中$\xi_0$是线性空间$V$中的一个固定向量.
    \end{enumerate}
\end{exercise}
\noindent 来源：\href{https://github.com/yhwu-is/Linear-Algebra-Left-Undone/blob/new/%E8%AE%B2%E4%B9%89/%E4%B8%93%E9%A2%98/5%20%E7%BA%BF%E6%80%A7%E6%98%A0%E5%B0%84.tex}{L5-例-4}。
\begin{solve}
\begin{enumerate}
        \item 出发空间为 $ \mathbf{R}^2 $，到达空间为 $ \mathbf{R}^3 $. $ \sigma $ 是线性映射.

              $ \forall (x_1, x_2), (y_1, y_2) \in \mathbf{R}^2,\enspace k_1, k_2 \in \mathbf{R} $, 有
              \begin{align*}
                      & \sigma(k_1(x_1, x_2) + k_2(y_1, y_2))                                                                     \\
                  ={} & ((k_1 x_1 + k_2 y_1) - (k_1 x_2 + k_2 y_2), k_1 x_1 + k_2 y_1, (k_1 x_1 + k_2 y_1) + (k_1 x_2 + k_2 y_2)) \\
                  ={} & k_1(x_1 - x_2, x_1, x_1 + x_2) + k_2(y_1 - y_2, y_1, y_1 + y_2)                                           \\
                  ={} & k_1 \sigma(x_1, x_2) + k_2 \sigma(y_1, y_2)
              \end{align*}

        \item 出发空间为 $ \mathbf{R}^2 $，到达空间为 $ \mathbf{R}^2 $. $ \sigma $ 不是线性映射.

              $ \forall (x_1, x_2), (y_1, y_2) \in \mathbf{R}^2,\enspace k_1, k_2 \in \mathbf{R} $, 有
              \begin{align*}
                  \sigma((x_1, x_2) + (y_1, y_2))
                   & = \sigma(x_1 + y_1, x_2 + y_2)                                        \\
                   & = ((x_1 + y_1)(x_2 + y_2), ((x_1 + y_1) + (x_2 + y_2))                \\
                   & = (x_1 x_2 + x_1 y_2 + y_1 x_2 + y_1 y_2, x_1 + y_1 + x_2 + y_2)      \\
                   & = (x_1 x_2 + y_1 y_2, x_1 + y_1 + x_2 + y_2) + (x_1 y_2 + y_1 x_2, 0) \\
                   & = \sigma(x_1, x_2) + \sigma(y_1, y_2) + (x_1 y_2 + y_1 x_2, 0)        \\
                   & \neq \sigma(x_1, x_2) + \sigma(y_1, y_2)
              \end{align*}

        \item 出发空间为 $ \mathbf{R}[x]_n $，到达空间为 $ \mathbf{R}[x]_{n - 1} $. $ \sigma $ 是线性映射.

              $ \forall p_1(x), p_2(x) \in \mathbf{R}[x]_n,\enspace k_1, k_2 \in \mathbf{R} $, 有
              \begin{align*}
                      & \sigma(k_1 p_1(x) + k_2 p_2(x))                     \\
                  ={} & (k_1 p_1 + k_2 p_2)(x + 1) - (k_1 p_1 + k_2 p_2)(x) \\
                  ={} & k_1(p_1(x + 1) - p_1(x)) + k_2(p_2(x + 1) - p_2(x)) \\
                  ={} & k_1 \sigma(p_1(x)) + k_2 \sigma(p_2(x))
              \end{align*}

        \item 出发空间为 $ \mathbf{R}[x] $，到达空间为 $ \mathbf{R}[x] $. $ \sigma $ 是线性映射.

              $ \forall p_1(x), p_2(x) \in \mathbf{R}[x],\enspace k_1, k_2 \in \mathbf{R} $, 有
              \begin{align*}
                      & \sigma(k_1 p_1(x) + k_2 p_2(x))         \\
                  ={} & (k_1 p_1 + k_2 p_2)(a)                  \\
                  ={} & k_1 p_1(a) + k_2 p_2(a)                 \\
                  ={} & k_1 \sigma(p_1(x)) + k_2 \sigma(p_2(x))
              \end{align*}

        \item 出发空间为 $ V $，到达空间为 $ V $. 当 $ \xi_0 = \vec{0} $ 时， $ \sigma(\xi) = 2 \xi $ 是线性映射.

              $ \forall \xi_1, \xi_2 \in V $, 有
              \begin{align*}
                  \sigma(\xi_1 + \xi_2) & = 2(\xi_1 + \xi_2)              \\
                                        & = 2 \xi_1 + 2 \xi_2             \\
                                        & = \sigma(\xi_1) + \sigma(\xi_2)
              \end{align*}

              当 $ \xi_0 \neq \vec{0} $ 时， $ \sigma $ 不是线性映射.

              $ \forall \xi_1, \xi_2 \in V $, 有
              \begin{align*}
                  \sigma(\xi_1 + \xi_2) & = 2(\xi_1 + \xi_2) + \xi_0              \\
                                        & = 2 \xi_1 + 2 \xi_2 + \xi_0             \\
                                        & = \sigma(\xi_1) + \sigma(\xi_2) - \xi_0 \\
                                        & \neq \sigma(\xi_1) + \sigma(\xi_2)
              \end{align*}
    \end{enumerate}
\end{solve}

\begin{exercise}\label{prob:L5-例-7}
是否存在$\mathbf{R}^3$到$\mathbf{R}^2$的线性映射$\sigma$使得$\sigma(1,-1,1)^{\mathrm{T}}=(1,0)^{\mathrm{T}},\enspace\sigma(1,1,1)^{\mathrm{T}}=(0,1)^{\mathrm{T}}$？
\end{exercise}
\noindent 来源：\href{https://github.com/yhwu-is/Linear-Algebra-Left-Undone/blob/new/%E8%AE%B2%E4%B9%89/%E4%B8%93%E9%A2%98/5%20%E7%BA%BF%E6%80%A7%E6%98%A0%E5%B0%84.tex}{L5-例-7}。
\begin{solve}
事实上这里的定义完全不违反上述的必要条件，因此我们考虑证明这一线性映射存在. 根据原讲义“线性映射构造”处，我们考虑构造出$\sigma$在一组基（我们取自然基）下的像，然后就可以根据原讲义“线性映射构造”处 知道这一映射一定存在.

    实际上，根据题给条件有
    \begin{gather*}
        \sigma(1,-1,1)^{\mathrm{T}}
        =\sigma(e_1-e_2+e_3)
        =\sigma(e_1)-\sigma(e_2)+\sigma(e_3)
        =(1,0)^{\mathrm{T}} \\
        \sigma(1,1,1)^{\mathrm{T}}
        =\sigma(e_1+e_2+e_3)
        =\sigma(e_1)+\sigma(e_2)+\sigma(e_3)
        =(0,1)^{\mathrm{T}}
    \end{gather*}
    我们希望解出$\sigma(e_1),\sigma(e_2),\sigma(e_3)$，这样就可以直接根据原讲义“线性映射构造”处 构造出这一线性映射. 但这一方程组只有2个方程却有三个未知量. 事实上我们可以任意定义$\sigma(e_3)=(0,0)$，然后解方程组得到
    \[\sigma(e_1)=\dfrac{1}{2} (1,1)^{\mathrm{T}},\enspace \sigma(e_2)=\dfrac{1}{2} (-1,1)^{\mathrm{T}}\]
    又由$\sigma(e_3)=(0,0)^{\mathrm{T}}$，根据原讲义“线性映射构造”处，满足题目条件的线性映射存在.
\end{solve}

\begin{exercise}\label{prob:L5-例-10}
指出下面各组内的两个线性空间是否同构，若同构可以进一步思考同构映射的构造：
    \begin{enumerate}
        \item 最高次不超过$n-1$的多项式构成的线性空间$\mathbf{R}[x]_n$与$\mathbf{R}^n$；

        \item 全体复数在实数域上的线性空间$\mathbf{C}(\mathbf{R})$与$\mathbf{R}^2$；

        \item 全体二元复向量$\mathbf{C}^2$在实数域上构成的线性空间$\mathbf{C}^2(\mathbf{R})$与$\mathbf{R}[x]_4$；

        \item 全体二元复向量$\mathbf{C}^2$在复数域上构成的线性空间$\mathbf{C}^2(\mathbf{C})$与$\mathcal{L}(\mathbf{R}^4,\mathbf{R})$.
    \end{enumerate}
\end{exercise}
\noindent 来源：\href{https://github.com/yhwu-is/Linear-Algebra-Left-Undone/blob/new/%E8%AE%B2%E4%B9%89/%E4%B8%93%E9%A2%98/5%20%E7%BA%BF%E6%80%A7%E6%98%A0%E5%B0%84.tex}{L5-例-10}。
\begin{solve}
\begin{enumerate}
        \item 同构，因为二者维数均为$n$. 同构映射非常简单，因为$\mathbf{R}[x]_n$在基$\{1,x,\ldots,x^{n-1}\}$下的坐标就在$\mathbf{R}^n$中，因此同构映射就是这一坐标映射：
              \[\sigma:\mathbf{R}[x]_n\to\mathbf{R}^n,\quad a_0+a_1x+\cdots+a_{n-1}x^{n-1}\mapsto(a_0,a_1,\ldots,a_{n-1}).\]
              我们很容易验证这一映射是线性双射，因此是同构映射.

        \item 同构，因为二者维数均为2（回忆原讲义“不同数域的维数”处）. 这一同构映射同上一小问理，$\mathbf{C}(\mathbf{R})$在基$\{1,i\}$下的坐标就在$\mathbf{R}^2$中，因此同构映射就是这一坐标映射：
              \[\sigma:\mathbf{C}(\mathbf{R})\to\mathbf{R}^2,\quad a+b\mathrm{i}\mapsto(a,b).\]
              事实上这就是将复数在二维平面中的表示，我们很容易验证这一映射是线性双射，因此是同构映射..

        \item 同构，因为二者维数均为4. 这一同构映射也非常简单，因为$\mathbf{C}^2(\mathbf{R})$在基
              \[\{(1,0),(\mathrm{i},0),(0,1),(0,\mathrm{i})\}\]
              下的坐标和$\mathbf{R}[x]_4$在基$\{1,x,x^2,x^3\}$下的坐标都在$\mathrm{R}^n$中，可以将它们一一对应，因此同构映射就是这一坐标映射：
              \[\sigma:\mathbf{C}^2(\mathbf{R})\to\mathbf{R}[x]_4,\quad (a+bi,c+di)\mapsto a+bx+cx^2+dx^3.\]
              我们很容易验证这一映射是线性双射，因此是同构映射.

        \item 不同构，因为$\mathbf{C}^2(\mathbf{C})$的维数为2，而$\mathcal{L}(\mathbf{R}^4,\mathbf{R})$的维数为4.
    \end{enumerate}
\end{solve}

\begin{exercise}\label{prob:L5-例-11}
求积空间$\mathbf{R}[x]_3\times\mathbf{R}^2$的一组基.
\end{exercise}
\noindent 来源：\href{https://github.com/yhwu-is/Linear-Algebra-Left-Undone/blob/new/%E8%AE%B2%E4%B9%89/%E4%B8%93%E9%A2%98/5%20%E7%BA%BF%E6%80%A7%E6%98%A0%E5%B0%84.tex}{L5-例-11}。
\begin{solve}
我们知道$\mathbf{R}[x]_3$的一组基为$1,x,x^2$，而$\mathbf{R}^2$的一组基为$(1,0),(0,1)$. 很自然的想法是：我们可以先取$\mathbf{R}[x]_3$的一组基，$\mathbf{R}^2$的位置置零，然后反之取$\mathbf{R}^2$的一组基，$\mathbf{R}[x]_3$的位置置零，即$(1,(0,0)),\ (x,(0,0)),\ (x^2,(0,0)),\ (0,(1,0)),\ (0,(0,1))$. 我们很容易可以证明上述向量组满足基的两个条件：线性无关和张成空间.
\end{solve}

\begin{exercise}\label{prob:S-POLYMAP-1}
在 $V=\mathbf R[x]_3$ 上定义 $T(f)(x)=f(x+1)+(-2)f'(x)$。证明 $T$ 线性，求它在基 $(1,x,x^2)$ 下的矩阵、核和像，并计算 $T(2x^{2}+1)$。
\end{exercise}
\noindent 编者补充：S-POLYMAP-1；相关主题见原讲义对应讲，题目数据与解答由本选编另行构造。
\begin{solve}
平移代入与求导均保持线性组合，故 $T(af+bg)=aT(f)+bT(g)$。分别计算基向量的像：
\[T(1)=1,\quad T(x)=x+(-1),\quad T(x^2)=x^2+(-2)x+1.\]
把这些像的坐标按列排列，得到 $[T]=\begin{pmatrix}1&-1&1\\0&1&-2\\0&0&1\end{pmatrix}$。这个上三角矩阵的对角元全为 $1$，所以可逆，$\ker T=\{0\}$、$\operatorname{im} T=V$，秩为 $3$。所求多项式的坐标为
\[[T]\begin{pmatrix}1\\0\\2\end{pmatrix}=\begin{pmatrix}3\\-4\\2\end{pmatrix},\]
因此 $T(2x^{2}+1)=2x^{2}-4x+3$。也可直接将 $x+1$ 代入原多项式后加上导数项核对。
\end{solve}

\begin{exercise}\label{prob:S-PROJ-1}
设 $T(x)=Ax$，其中 $A=\begin{pmatrix}1&0&0\\0&1&0\\1&3&0\end{pmatrix}$。验证 $T^2=T$，求 $\operatorname{im} T$、$\ker T$ 的基，并把 $x=\begin{pmatrix}1\\1\\1\end{pmatrix}$ 唯一分解为像空间向量与核向量之和。判断该投影是否一定是正交投影。
\end{exercise}
\noindent 编者补充：S-PROJ-1；相关主题见原讲义对应讲，题目数据与解答由本选编另行构造。
\begin{solve}
取 $P=\begin{pmatrix}1&-3&0\\0&1&0\\1&0&1\end{pmatrix}$，有 $A=P\operatorname{diag}(1,1,0)P^{-1}$。因为对角矩阵的平方仍等于自身，$A^2=A$。像的一组基为 $P$ 的前两列，核的一组基为其第三列，故维数分别为 $2$ 与 $1$。对于题给向量，
\[u=Ax=\begin{pmatrix}1\\1\\4\end{pmatrix},\qquad v=(I-A)x=\begin{pmatrix}0\\0\\-3\end{pmatrix},\qquad x=u+v.\]
由 $A(I-A)=O$ 得 $v\in\ker T$。若一个向量同时在像与核中，则它既满足 $Az=z$ 又满足 $Az=0$，只能为零，故分解唯一。幂等性只保证沿核向像的投影；在标准实内积下，正交投影还须 $A^{\mathrm T}=A$，应另外检查这一条件。
\end{solve}

\begin{exercise}\label{prob:S-POLYMAP-2}
在 $V=\mathbf R[x]_3$ 上定义 $T(f)(x)=f(x+2)+(-1)f'(x)$。证明 $T$ 线性，求它在基 $(1,x,x^2)$ 下的矩阵、核和像，并计算 $T(2x^{2}-x+2)$。
\end{exercise}
\noindent 编者补充：S-POLYMAP-2；相关主题见原讲义对应讲，题目数据与解答由本选编另行构造。
\begin{solve}
平移代入与求导均保持线性组合，故 $T(af+bg)=aT(f)+bT(g)$。分别计算基向量的像：
\[T(1)=1,\quad T(x)=x+(1),\quad T(x^2)=x^2+(2)x+4.\]
把这些像的坐标按列排列，得到 $[T]=\begin{pmatrix}1&1&4\\0&1&2\\0&0&1\end{pmatrix}$。这个上三角矩阵的对角元全为 $1$，所以可逆，$\ker T=\{0\}$、$\operatorname{im} T=V$，秩为 $3$。所求多项式的坐标为
\[[T]\begin{pmatrix}2\\-1\\2\end{pmatrix}=\begin{pmatrix}9\\3\\2\end{pmatrix},\]
因此 $T(2x^{2}-x+2)=2x^{2}+3x+9$。也可直接将 $x+2$ 代入原多项式后加上导数项核对。
\end{solve}

\begin{exercise}\label{prob:S-PROJ-2}
设 $T(x)=Ax$，其中 $A=\begin{pmatrix}1&0&2\\0&1&-2\\0&0&0\end{pmatrix}$。验证 $T^2=T$，求 $\operatorname{im} T$、$\ker T$ 的基，并把 $x=\begin{pmatrix}2\\1\\0\end{pmatrix}$ 唯一分解为像空间向量与核向量之和。判断该投影是否一定是正交投影。
\end{exercise}
\noindent 编者补充：S-PROJ-2；相关主题见原讲义对应讲，题目数据与解答由本选编另行构造。
\begin{solve}
取 $P=\begin{pmatrix}1&-1&-2\\0&1&2\\0&0&1\end{pmatrix}$，有 $A=P\operatorname{diag}(1,1,0)P^{-1}$。因为对角矩阵的平方仍等于自身，$A^2=A$。像的一组基为 $P$ 的前两列，核的一组基为其第三列，故维数分别为 $2$ 与 $1$。对于题给向量，
\[u=Ax=\begin{pmatrix}2\\1\\0\end{pmatrix},\qquad v=(I-A)x=\begin{pmatrix}0\\0\\0\end{pmatrix},\qquad x=u+v.\]
由 $A(I-A)=O$ 得 $v\in\ker T$。若一个向量同时在像与核中，则它既满足 $Az=z$ 又满足 $Az=0$，只能为零，故分解唯一。幂等性只保证沿核向像的投影；在标准实内积下，正交投影还须 $A^{\mathrm T}=A$，应另外检查这一条件。
\end{solve}

\begin{exercise}\label{prob:S-POLYMAP-3}
在 $V=\mathbf R[x]_3$ 上定义 $T(f)(x)=f(x+3)+(0)f'(x)$。证明 $T$ 线性，求它在基 $(1,x,x^2)$ 下的矩阵、核和像，并计算 $T(2x^{2}-2x+3)$。
\end{exercise}
\noindent 编者补充：S-POLYMAP-3；相关主题见原讲义对应讲，题目数据与解答由本选编另行构造。
\begin{solve}
平移代入与求导均保持线性组合，故 $T(af+bg)=aT(f)+bT(g)$。分别计算基向量的像：
\[T(1)=1,\quad T(x)=x+(3),\quad T(x^2)=x^2+(6)x+9.\]
把这些像的坐标按列排列，得到 $[T]=\begin{pmatrix}1&3&9\\0&1&6\\0&0&1\end{pmatrix}$。这个上三角矩阵的对角元全为 $1$，所以可逆，$\ker T=\{0\}$、$\operatorname{im} T=V$，秩为 $3$。所求多项式的坐标为
\[[T]\begin{pmatrix}3\\-2\\2\end{pmatrix}=\begin{pmatrix}15\\10\\2\end{pmatrix},\]
因此 $T(2x^{2}-2x+3)=2x^{2}+10x+15$。也可直接将 $x+3$ 代入原多项式后加上导数项核对。
\end{solve}

\begin{exercise}\label{prob:S-PROJ-3}
设 $T(x)=Ax$，其中 $A=\begin{pmatrix}1&0&4\\0&1&2\\0&0&0\end{pmatrix}$。验证 $T^2=T$，求 $\operatorname{im} T$、$\ker T$ 的基，并把 $x=\begin{pmatrix}3\\1\\-1\end{pmatrix}$ 唯一分解为像空间向量与核向量之和。判断该投影是否一定是正交投影。
\end{exercise}
\noindent 编者补充：S-PROJ-3；相关主题见原讲义对应讲，题目数据与解答由本选编另行构造。
\begin{solve}
取 $P=\begin{pmatrix}1&-2&-4\\1&-1&-2\\0&0&1\end{pmatrix}$，有 $A=P\operatorname{diag}(1,1,0)P^{-1}$。因为对角矩阵的平方仍等于自身，$A^2=A$。像的一组基为 $P$ 的前两列，核的一组基为其第三列，故维数分别为 $2$ 与 $1$。对于题给向量，
\[u=Ax=\begin{pmatrix}-1\\-1\\0\end{pmatrix},\qquad v=(I-A)x=\begin{pmatrix}4\\2\\-1\end{pmatrix},\qquad x=u+v.\]
由 $A(I-A)=O$ 得 $v\in\ker T$。若一个向量同时在像与核中，则它既满足 $Az=z$ 又满足 $Az=0$，只能为零，故分解唯一。幂等性只保证沿核向像的投影；在标准实内积下，正交投影还须 $A^{\mathrm T}=A$，应另外检查这一条件。
\end{solve}

\begin{exercise}\label{prob:S-POLYMAP-4}
在 $V=\mathbf R[x]_3$ 上定义 $T(f)(x)=f(x+4)+(1)f'(x)$。证明 $T$ 线性，求它在基 $(1,x,x^2)$ 下的矩阵、核和像，并计算 $T(2x^{2}-3x+4)$。
\end{exercise}
\noindent 编者补充：S-POLYMAP-4；相关主题见原讲义对应讲，题目数据与解答由本选编另行构造。
\begin{solve}
平移代入与求导均保持线性组合，故 $T(af+bg)=aT(f)+bT(g)$。分别计算基向量的像：
\[T(1)=1,\quad T(x)=x+(5),\quad T(x^2)=x^2+(10)x+16.\]
把这些像的坐标按列排列，得到 $[T]=\begin{pmatrix}1&5&16\\0&1&10\\0&0&1\end{pmatrix}$。这个上三角矩阵的对角元全为 $1$，所以可逆，$\ker T=\{0\}$、$\operatorname{im} T=V$，秩为 $3$。所求多项式的坐标为
\[[T]\begin{pmatrix}4\\-3\\2\end{pmatrix}=\begin{pmatrix}21\\17\\2\end{pmatrix},\]
因此 $T(2x^{2}-3x+4)=2x^{2}+17x+21$。也可直接将 $x+4$ 代入原多项式后加上导数项核对。
\end{solve}

\begin{exercise}\label{prob:S-PROJ-4}
设 $T(x)=Ax$，其中 $A=\begin{pmatrix}1&0&2\\0&1&2\\0&0&0\end{pmatrix}$。验证 $T^2=T$，求 $\operatorname{im} T$、$\ker T$ 的基，并把 $x=\begin{pmatrix}4\\1\\-2\end{pmatrix}$ 唯一分解为像空间向量与核向量之和。判断该投影是否一定是正交投影。
\end{exercise}
\noindent 编者补充：S-PROJ-4；相关主题见原讲义对应讲，题目数据与解答由本选编另行构造。
\begin{solve}
取 $P=\begin{pmatrix}1&0&-2\\0&1&-2\\0&0&1\end{pmatrix}$，有 $A=P\operatorname{diag}(1,1,0)P^{-1}$。因为对角矩阵的平方仍等于自身，$A^2=A$。像的一组基为 $P$ 的前两列，核的一组基为其第三列，故维数分别为 $2$ 与 $1$。对于题给向量，
\[u=Ax=\begin{pmatrix}0\\-3\\0\end{pmatrix},\qquad v=(I-A)x=\begin{pmatrix}4\\4\\-2\end{pmatrix},\qquad x=u+v.\]
由 $A(I-A)=O$ 得 $v\in\ker T$。若一个向量同时在像与核中，则它既满足 $Az=z$ 又满足 $Az=0$，只能为零，故分解唯一。幂等性只保证沿核向像的投影；在标准实内积下，正交投影还须 $A^{\mathrm T}=A$，应另外检查这一条件。
\end{solve}

\begin{exercise}\label{prob:S-POLYMAP-5}
在 $V=\mathbf R[x]_3$ 上定义 $T(f)(x)=f(x+5)+(2)f'(x)$。证明 $T$ 线性，求它在基 $(1,x,x^2)$ 下的矩阵、核和像，并计算 $T(2x^{2}-4x+5)$。
\end{exercise}
\noindent 编者补充：S-POLYMAP-5；相关主题见原讲义对应讲，题目数据与解答由本选编另行构造。
\begin{solve}
平移代入与求导均保持线性组合，故 $T(af+bg)=aT(f)+bT(g)$。分别计算基向量的像：
\[T(1)=1,\quad T(x)=x+(7),\quad T(x^2)=x^2+(14)x+25.\]
把这些像的坐标按列排列，得到 $[T]=\begin{pmatrix}1&7&25\\0&1&14\\0&0&1\end{pmatrix}$。这个上三角矩阵的对角元全为 $1$，所以可逆，$\ker T=\{0\}$、$\operatorname{im} T=V$，秩为 $3$。所求多项式的坐标为
\[[T]\begin{pmatrix}5\\-4\\2\end{pmatrix}=\begin{pmatrix}27\\24\\2\end{pmatrix},\]
因此 $T(2x^{2}-4x+5)=2x^{2}+24x+27$。也可直接将 $x+5$ 代入原多项式后加上导数项核对。
\end{solve}

\begin{exercise}\label{prob:S-PROJ-5}
设 $T(x)=Ax$，其中 $A=\begin{pmatrix}1&6&-3\\0&1&0\\0&2&0\end{pmatrix}$。验证 $T^2=T$，求 $\operatorname{im} T$、$\ker T$ 的基，并把 $x=\begin{pmatrix}5\\1\\-3\end{pmatrix}$ 唯一分解为像空间向量与核向量之和。判断该投影是否一定是正交投影。
\end{exercise}
\noindent 编者补充：S-PROJ-5；相关主题见原讲义对应讲，题目数据与解答由本选编另行构造。
\begin{solve}
取 $P=\begin{pmatrix}1&4&3\\0&1&0\\0&2&1\end{pmatrix}$，有 $A=P\operatorname{diag}(1,1,0)P^{-1}$。因为对角矩阵的平方仍等于自身，$A^2=A$。像的一组基为 $P$ 的前两列，核的一组基为其第三列，故维数分别为 $2$ 与 $1$。对于题给向量，
\[u=Ax=\begin{pmatrix}20\\1\\2\end{pmatrix},\qquad v=(I-A)x=\begin{pmatrix}-15\\0\\-5\end{pmatrix},\qquad x=u+v.\]
由 $A(I-A)=O$ 得 $v\in\ker T$。若一个向量同时在像与核中，则它既满足 $Az=z$ 又满足 $Az=0$，只能为零，故分解唯一。幂等性只保证沿核向像的投影；在标准实内积下，正交投影还须 $A^{\mathrm T}=A$，应另外检查这一条件。
\end{solve}

\begin{exercise}\label{prob:S-POLYMAP-6}
在 $V=\mathbf R[x]_3$ 上定义 $T(f)(x)=f(x+6)+(3)f'(x)$。证明 $T$ 线性，求它在基 $(1,x,x^2)$ 下的矩阵、核和像，并计算 $T(2x^{2}-5x+6)$。
\end{exercise}
\noindent 编者补充：S-POLYMAP-6；相关主题见原讲义对应讲，题目数据与解答由本选编另行构造。
\begin{solve}
平移代入与求导均保持线性组合，故 $T(af+bg)=aT(f)+bT(g)$。分别计算基向量的像：
\[T(1)=1,\quad T(x)=x+(9),\quad T(x^2)=x^2+(18)x+36.\]
把这些像的坐标按列排列，得到 $[T]=\begin{pmatrix}1&9&36\\0&1&18\\0&0&1\end{pmatrix}$。这个上三角矩阵的对角元全为 $1$，所以可逆，$\ker T=\{0\}$、$\operatorname{im} T=V$，秩为 $3$。所求多项式的坐标为
\[[T]\begin{pmatrix}6\\-5\\2\end{pmatrix}=\begin{pmatrix}33\\31\\2\end{pmatrix},\]
因此 $T(2x^{2}-5x+6)=2x^{2}+31x+33$。也可直接将 $x+6$ 代入原多项式后加上导数项核对。
\end{solve}

\begin{exercise}\label{prob:S-PROJ-6}
设 $T(x)=Ax$，其中 $A=\begin{pmatrix}1&0&0\\0&1&0\\-1&-1&0\end{pmatrix}$。验证 $T^2=T$，求 $\operatorname{im} T$、$\ker T$ 的基，并把 $x=\begin{pmatrix}6\\1\\-4\end{pmatrix}$ 唯一分解为像空间向量与核向量之和。判断该投影是否一定是正交投影。
\end{exercise}
\noindent 编者补充：S-PROJ-6；相关主题见原讲义对应讲，题目数据与解答由本选编另行构造。
\begin{solve}
取 $P=\begin{pmatrix}1&0&0\\0&1&0\\-1&-1&1\end{pmatrix}$，有 $A=P\operatorname{diag}(1,1,0)P^{-1}$。因为对角矩阵的平方仍等于自身，$A^2=A$。像的一组基为 $P$ 的前两列，核的一组基为其第三列，故维数分别为 $2$ 与 $1$。对于题给向量，
\[u=Ax=\begin{pmatrix}6\\1\\-7\end{pmatrix},\qquad v=(I-A)x=\begin{pmatrix}0\\0\\3\end{pmatrix},\qquad x=u+v.\]
由 $A(I-A)=O$ 得 $v\in\ker T$。若一个向量同时在像与核中，则它既满足 $Az=z$ 又满足 $Az=0$，只能为零，故分解唯一。幂等性只保证沿核向像的投影；在标准实内积下，正交投影还须 $A^{\mathrm T}=A$，应另外检查这一条件。
\end{solve}

\section{对偶空间}\label{sec:topic6}
本节收录 1 道题。线性泛函是取标量值的线性映射；对偶基的条件为 $f_i(g_j)=\delta_{ij}$，可转化为求逆矩阵。
\begin{exercise}\label{prob:L6-例-1}
设 $V=\mathbf R[x]_3$，定义线性泛函
\[f_1(g)=\int_0^1g(x)\,\mathrm dx,\quad f_2(g)=\int_0^2g(x)\,\mathrm dx,\quad f_3(g)=\int_0^{-1}g(x)\,\mathrm dx.\]
证明 $f_1,f_2,f_3$ 构成对偶空间 $V^*=\mathcal L(V,\mathbf R)$ 的一组基，并求 $V$ 的一组基 $g_1,g_2,g_3$，使 $f_i(g_j)=\delta_{ij}$。
\end{exercise}
\noindent 来源：\href{https://github.com/yhwu-is/Linear-Algebra-Left-Undone/blob/new/%E8%AE%B2%E4%B9%89/%E4%B8%93%E9%A2%98/6%20%E5%AF%B9%E5%81%B6%E7%A9%BA%E9%97%B4.tex}{L6-例-1}。
\begin{solve}
积分的线性性保证 $f_i\in V^*$。对 $g(x)=a+bx+cx^2$，有
\[\begin{pmatrix}f_1(g)\\f_2(g)\\f_3(g)\end{pmatrix}=M\begin{pmatrix}a\\b\\c\end{pmatrix},\qquad
M=\begin{pmatrix}1&1/2&1/3\\2&2&8/3\\-1&1/2&-1/3\end{pmatrix}.\]
直接计算 $\det M=-2\ne0$，故三个泛函线性无关；又 $\dim V^*=3$，它们构成一组基。所求 $g_j$ 的坐标是 $M^{-1}$ 的第 $j$ 列，由
\[M^{-1}=\begin{pmatrix}1&-1/6&-1/3\\1&0&1\\-3/2&1/2&-1/2\end{pmatrix}\]
得到
\[g_1=1+x-\frac32x^2,\quad g_2=-\frac16+\frac12x^2,\quad g_3=-\frac13+x-\frac12x^2.\]
等式 $MM^{-1}=I_3$ 同时验证 $f_i(g_j)=\delta_{ij}$；$M^{-1}$ 可逆也保证这三个多项式为基。
\end{solve}

\section{矩阵运算与线性映射的矩阵表示}\label{sec:topic7}
本节收录 32 道题，其中 6 道为编者补充计算题。矩阵的第 $j$ 列记录第 $j$ 个基向量的像坐标。换基时先写出向量与坐标的关系，避免颠倒过渡矩阵方向。
\begin{exercise}\label{prob:L7-A-1}
证明：若$AB=BA$，$AC=CA$，则$A,B,C$为同阶方阵，且
        \[A(BC)=(BC)A,\enspace A(B+C)=(B+C)A.\]
\end{exercise}
\noindent 来源：\href{https://github.com/yhwu-is/Linear-Algebra-Left-Undone/blob/new/%E8%AE%B2%E4%B9%89/%E4%B8%93%E9%A2%98/7%20%E7%BA%BF%E6%80%A7%E6%98%A0%E5%B0%84%E7%9F%A9%E9%98%B5%E8%A1%A8%E7%A4%BA%E4%B8%8E%E7%9F%A9%E9%98%B5%E8%BF%90%E7%AE%97%E5%9F%BA%E7%A1%80.tex}{L7-A-1}。
\begin{solve}
由 $AB=BA$ 得 $A(BC)=(AB)C=(BA)C=B(AC)=BCA=(BC)A$，由$AC=CA$ 得 $A(B+C)=AB+AC=BA+CA=B+C$.
\end{solve}

\begin{exercise}\label{prob:L7-A-2}
$A,B$都是$n$阶矩阵，求下列等式成立的充分条件：
        \begin{enumerate}
            \item $(A+B)^3=A^3+3A^2B+3AB^2+B^3$；

            \item $(A+B)(A-B)=A^2-B^2$.
        \end{enumerate}
\end{exercise}
\noindent 来源：\href{https://github.com/yhwu-is/Linear-Algebra-Left-Undone/blob/new/%E8%AE%B2%E4%B9%89/%E4%B8%93%E9%A2%98/7%20%E7%BA%BF%E6%80%A7%E6%98%A0%E5%B0%84%E7%9F%A9%E9%98%B5%E8%A1%A8%E7%A4%BA%E4%B8%8E%E7%9F%A9%E9%98%B5%E8%BF%90%E7%AE%97%E5%9F%BA%E7%A1%80.tex}{L7-A-2}。
\begin{solve}
\begin{enumerate}
                \item $(A+B)^3=A^3+3A^2B+3AB^2+B^3$的充分条件是$AB=BA$；

                \item $(A+B)(A-B)=A^2-B^2$的充分条件是$AB=BA$.
            \end{enumerate}
\end{solve}

\begin{exercise}\label{prob:L7-A-3}
设$A$是$n$阶方阵且$A^n=O$，证明：
        \[(E_n-A)(E_n+A+A^2+\cdots+A^{n-1})=E_n.\]
\end{exercise}
\noindent 来源：\href{https://github.com/yhwu-is/Linear-Algebra-Left-Undone/blob/new/%E8%AE%B2%E4%B9%89/%E4%B8%93%E9%A2%98/7%20%E7%BA%BF%E6%80%A7%E6%98%A0%E5%B0%84%E7%9F%A9%E9%98%B5%E8%A1%A8%E7%A4%BA%E4%B8%8E%E7%9F%A9%E9%98%B5%E8%BF%90%E7%AE%97%E5%9F%BA%E7%A1%80.tex}{L7-A-3}。
\begin{solve}
由$A^n=O$得$A^n-E_n^n=-E_n$，即$(A-E_n)(A^{n-1}+A^{n-2}+\cdots+A+E_n)=-E_n$，即$(E_n-A)(E_n+A+A^2+\cdots+A^{n-1})=E_n$.
\end{solve}

\begin{exercise}\label{prob:L7-A-4}
证明：若线性映射$\sigma \in \mathcal{L}(V_1,V_2)$可逆，则其逆映射唯一.
\end{exercise}
\noindent 来源：\href{https://github.com/yhwu-is/Linear-Algebra-Left-Undone/blob/new/%E8%AE%B2%E4%B9%89/%E4%B8%93%E9%A2%98/7%20%E7%BA%BF%E6%80%A7%E6%98%A0%E5%B0%84%E7%9F%A9%E9%98%B5%E8%A1%A8%E7%A4%BA%E4%B8%8E%E7%9F%A9%E9%98%B5%E8%BF%90%E7%AE%97%E5%9F%BA%E7%A1%80.tex}{L7-A-4}。
\begin{solve}
设$\sigma$的逆映射为$\tau_1,\tau_2$，则$\sigma\tau_1=\sigma\tau_2=E$，则$\sigma(\tau_1-\tau_2)=0$.

            因为 $\sigma$ 可逆，所以对于任意 $x \neq 0$，有 $\sigma(x) \neq 0$.

            而 $\forall x, \sigma((\tau_1-\tau_2)(x))=0$，故 $\tau_1-\tau_2=0, \tau_1=\tau_2$.
\end{solve}

\begin{exercise}\label{prob:L7-A-5}
证明：有一行元素或一列元素全为0的$n$阶方阵必定不可逆.
\end{exercise}
\noindent 来源：\href{https://github.com/yhwu-is/Linear-Algebra-Left-Undone/blob/new/%E8%AE%B2%E4%B9%89/%E4%B8%93%E9%A2%98/7%20%E7%BA%BF%E6%80%A7%E6%98%A0%E5%B0%84%E7%9F%A9%E9%98%B5%E8%A1%A8%E7%A4%BA%E4%B8%8E%E7%9F%A9%E9%98%B5%E8%BF%90%E7%AE%97%E5%9F%BA%E7%A1%80.tex}{L7-A-5}。
\begin{solve}
设$A$的第$i$行全为0，则$A$的第$i$行与$A$的任意行线性相关，记 $A$ 为 $n$ 阶矩阵，则$r(A) < n$.

            故对任意 $n$ 阶矩阵 $B$，有 $r(AB) \leqslant r(A) < n = r(E_n)$，故 $A$ 不可逆.
\end{solve}

\begin{exercise}\label{prob:L7-A-6}
设$\alpha,\beta$为三维列向量，且$\alpha\beta^\mathrm{T}=\begin{pmatrix}
                -1 & 2  & 1  \\
                1  & -2 & -1 \\
                2  & -4 & -2
            \end{pmatrix}$，求$\alpha^\mathrm{T}\beta$.
\end{exercise}
\noindent 来源：\href{https://github.com/yhwu-is/Linear-Algebra-Left-Undone/blob/new/%E8%AE%B2%E4%B9%89/%E4%B8%93%E9%A2%98/7%20%E7%BA%BF%E6%80%A7%E6%98%A0%E5%B0%84%E7%9F%A9%E9%98%B5%E8%A1%A8%E7%A4%BA%E4%B8%8E%E7%9F%A9%E9%98%B5%E8%BF%90%E7%AE%97%E5%9F%BA%E7%A1%80.tex}{L7-A-6}。
\begin{solve}
记 $\alpha = (a_1,a_2,a_3)$, $\beta = (b_1,b_2,b_3)$，则 
\[\alpha\beta^\mathrm{T} = \begin{pmatrix}
                    a_1b_1 & a_1b_2 & a_1b_3 \\
                    a_2b_1 & a_2b_2 & a_2b_3 \\
                    a_3b_1 & a_3b_2 & a_3b_3
                \end{pmatrix}\]
,
            则有：$\text{tr}(\alpha\beta^\mathrm{T}) = a_1b_1+a_2b_2+a_3b_3 = \alpha^\mathrm{T}\beta$ ，因此 $\alpha^\mathrm{T}\beta = -1-2-2 = -5$.
\end{solve}

\begin{exercise}\label{prob:L7-B-1}
已知$f_1=1-x,f_2=1+x^2,f_3=x+2x^2$是$\mathbf{R}[x]_3$中三个元素，$\sigma$是$\mathbf{R}[x]_3$上的线性变换且满足$\sigma(f_1)=2+x^2,\sigma(f_2)=x,\sigma(f_3)=1+x+x^2$.
        \begin{enumerate}
            \item 证明：$f_1,f_2,f_3$构成$\mathbf{R}[x]_3$的一组基；

            \item 求$\sigma$在基$f_1,f_2,f_3$下的矩阵；

            \item 设$f=1+2x+3x^2$，求$\sigma(f)$.
        \end{enumerate}
\end{exercise}
\noindent 来源：\href{https://github.com/yhwu-is/Linear-Algebra-Left-Undone/blob/new/%E8%AE%B2%E4%B9%89/%E4%B8%93%E9%A2%98/7%20%E7%BA%BF%E6%80%A7%E6%98%A0%E5%B0%84%E7%9F%A9%E9%98%B5%E8%A1%A8%E7%A4%BA%E4%B8%8E%E7%9F%A9%E9%98%B5%E8%BF%90%E7%AE%97%E5%9F%BA%E7%A1%80.tex}{L7-B-1}。
\begin{solve}
\begin{enumerate}
                \item 只需要证明$f_1,f_2,f_3$是线性无关的即可，事实上，这也十分显然(考虑其在$1,x,x^2$下的坐标表示线性无关),想必读者可以自行证明

                \item 方法一：设 $ (\sigma(f_1), \sigma(f_2), \sigma(f_3)) = (f_1, f_2, f_3) A $，则由
                      \[ (\sigma(f_1), \sigma(f_2), \sigma(f_3)) = (1, x, x^2) \begin{pmatrix}
                              2 & 0 & 1 \\
                              0 & 1 & 1 \\
                              1 & 0 & 1
                          \end{pmatrix} = (1, x, x^2) \begin{pmatrix}
                              1  & 1 & 0 \\
                              -1 & 0 & 1 \\
                              0  & 1 & 2
                          \end{pmatrix} A \]
                      可得
                      \begin{align*}
                          A & = \begin{pmatrix}
                                    1  & 1 & 0 \\
                                    -1 & 0 & 1 \\
                                    0  & 1 & 2
                                \end{pmatrix}^{-1}
                          \begin{pmatrix}
                              2 & 0 & 1 \\
                              0 & 1 & 1 \\
                              1 & 0 & 1
                          \end{pmatrix}           \\
                            & = \begin{pmatrix}
                                    -1 & -2 & 1  \\
                                    2  & 2  & -1 \\
                                    -1 & -1 & 1
                                \end{pmatrix}
                          \begin{pmatrix}
                              2 & 0 & 1 \\
                              0 & 1 & 1 \\
                              1 & 0 & 1
                          \end{pmatrix}
                          = \begin{pmatrix}
                                -1 & -2 & -2 \\
                                3  & 2  & 3  \\
                                -1 & -1 & -1
                            \end{pmatrix}
                      \end{align*}

                      方法二：通过待定系数法解方程组
                      \[ \begin{cases}
                              \sigma(f_1) = -f_1 + 3 f_2 - f_3   \\
                              \sigma(f_2) = -2 f_1 + 2 f_2 - f_3 \\
                              \sigma(f_3) = -2 f_1 + 3 f_2 - f_3
                          \end{cases} \]
                      解得
                      \[ (\sigma(f_1), \sigma(f_2), \sigma(f_3)) = (f_1, f_2, f_3) \begin{pmatrix}
                              -1 & -2 & -2 \\
                              3  & 2  & 3  \\
                              -1 & -1 & -1
                          \end{pmatrix} \]

                \item \begin{align*}
                          \sigma(f) & = \sigma\left((1, x, x^2)
                          \begin{pmatrix} 1 \\ 2 \\ 3 \end{pmatrix}\right)= \sigma\left((f_1, f_2, f_3)
                          \begin{pmatrix}
                              1  & 1 & 0 \\
                              -1 & 0 & 1 \\
                              0  & 1 & 2
                          \end{pmatrix}^{-1} \begin{pmatrix} 1 \\ 2 \\ 3 \end{pmatrix}\right) \\
                                    & = (\sigma(f_1), \sigma(f_2), \sigma(f_3))
                          \begin{pmatrix}
                              1  & 1 & 0 \\
                              -1 & 0 & 1 \\
                              0  & 1 & 2
                          \end{pmatrix}^{-1} \begin{pmatrix} 1 \\ 2 \\ 3 \end{pmatrix}        \\
                                    & = (1, x, x^2)
                          \begin{pmatrix}
                              2 & 0 & 1 \\
                              0 & 1 & 1 \\
                              1 & 0 & 1
                          \end{pmatrix}
                          \begin{pmatrix}
                              1  & 1 & 0 \\
                              -1 & 0 & 1 \\
                              0  & 1 & 2
                          \end{pmatrix}^{-1}
                          \begin{pmatrix} 1 \\ 2 \\ 3 \end{pmatrix}                           \\
                                    & = (1, x, x^2)
                          \begin{pmatrix} -4 \\ 3 \\ 2 \end{pmatrix}                          \\
                                    & = 2x^2 + 3x - 4
                      \end{align*}

                      或也可采用待定系数法求出 $ f = -2 f_1 + 3 f_2 $，所以 $ \sigma(f) = -2 \sigma(f_1) + 3 \sigma(f_2) = 2x^2 + 3x - 4 $.
            \end{enumerate}
\end{solve}

\begin{exercise}\label{prob:L7-B-3}
设 $\mathbf{R}[x]_4$ 是数域 $\mathbf{R}$ 上次数小于 4 的多项式所构成的线性空间（约定零多项式次数为 $-\infty$）. $\mathbf{M}_2(\mathbf{R})$ 是 $\mathbf{R}$ 上 2 阶方阵所构成的线性空间. 定义 $T \colon \mathbf{R}[x]_4 \to \mathbf{M}_2(\mathbf{R})$ 如下：对 $f(x) \in \mathbf{R}[x]_4$，
        \[T(f(x))=\begin{pmatrix}f(0) & f(1) \\ f(-1) & f(0)\end{pmatrix}\]
        \begin{enumerate}
            \item 求出 $T$ 的核空间 $N(T)$ 和像空间 $R(T)$；

            \item 求$T$在$\mathbf{R}[x]_4$和$\mathbf{M}_2(\mathbf{R})$的基下的矩阵表示.
        \end{enumerate}
\end{exercise}
\noindent 来源：\href{https://github.com/yhwu-is/Linear-Algebra-Left-Undone/blob/new/%E8%AE%B2%E4%B9%89/%E4%B8%93%E9%A2%98/7%20%E7%BA%BF%E6%80%A7%E6%98%A0%E5%B0%84%E7%9F%A9%E9%98%B5%E8%A1%A8%E7%A4%BA%E4%B8%8E%E7%9F%A9%E9%98%B5%E8%BF%90%E7%AE%97%E5%9F%BA%E7%A1%80.tex}{L7-B-3}。
\begin{solve}
\begin{enumerate}
                \item \label{item:7:B:5:1}
                      求核空间，即求使得 $ T(f(x)) $ 为零矩阵的 $ f(x) $ 构成的空间. 设 $ f(x) = ax^3 + bx^2 + cx + d $，则有
                      \[ \begin{cases}
                              f(0) = d = 0             \\
                              f(1) = a + b + c + d = 0 \\
                              f(-1) = -a + b -c + d = 0
                          \end{cases} \]
                      解方程，令 $ a = t $ 有
                      \[ \begin{cases}
                              a = t  \\
                              b = 0  \\
                              c = -t \\
                              d = 0
                          \end{cases} \implies f(x) = t(x^3 - x) \]
                      故 $ N(T) = \operatorname{span}(x^3 - x) $.

                      求像空间，取 $ \mathbf{R}[x]_4 $ 的自然基 $ 1, x, x^2, x^3 $.
                      \[ \begin{matrix}
                              T(1) = \begin{pmatrix} 1 & 1 \\ 1 & 1 \end{pmatrix}   & T(x) = \begin{pmatrix} 0 & 1 \\ -1 & 0 \end{pmatrix}   \\
                              T(x^2) = \begin{pmatrix} 0 & 1 \\ 1 & 0 \end{pmatrix} & T(x^3) = \begin{pmatrix} 0 & 1 \\ -1 & 0 \end{pmatrix}
                          \end{matrix} \]
                      求它们的极大线性无关组. 我们发现 $ T(x) = T(x^3) $，故先舍弃 $ T(x^3) $，然后令
                      \begin{align*}
                                     & k_1 T(1) + k_2 T(x) + k_3 T(x^2) = 0                                             \\
                          \implies{} & \begin{pmatrix} k_1 & k_1 + k_2 + k_3 \\ k_1 - k_2 + k_3 & k_1 \end{pmatrix} = 0 \\
                          \implies{} & \begin{cases}
                                           k_1 = 0             \\
                                           k_1 + k_2 + k_3 = 0 \\
                                           k_1 - k_2 + k_3 = 0 \\
                                       \end{cases}                                                              \\
                          \implies{} & k_1 = k_2 = k_3 = 0
                      \end{align*}
                      故 $ T(1), T(x), T(x^2) $ 线性无关. 故 $ R(T) = \operatorname{span}\left(\begin{pmatrix} 1 & 1 \\ 1 & 1 \end{pmatrix}, \begin{pmatrix} 0 & 1 \\ -1 & 0 \end{pmatrix}, \begin{pmatrix} 0 & 1 \\ 1 & 0 \end{pmatrix}\right) $.

                \item 我们只需要将$2 \times 2$矩阵的第二列拼到第一列下方就可以变为我们熟悉的列向量形式，所以其矩阵表示也是简单的
                      \[
                        T(1,x,x^2,x^3)=(e_1,e_2,e_3,e_4)\begin{pmatrix}
                            1 & 0 & 0 & 0 \\
                            1 & -1 & 1 & -1 \\
                            1 & 1 & 1 & 1 \\
                            1 & 0 & 0 & 0
                        \end{pmatrix}
                      \]

                      其中$e_1=\begin{pmatrix}1&0\\0&0\end{pmatrix}$,$e_2=\begin{pmatrix}0&0\\1&0\end{pmatrix}$,
                      $e_3=\begin{pmatrix}0&1\\0&0\end{pmatrix}$,$e_4=\begin{pmatrix}0&0\\0&1\end{pmatrix}$.
            \end{enumerate}
\end{solve}

\begin{exercise}\label{prob:L7-B-4}
设$A=\begin{pmatrix}
                1 & -1 & -1 \\ -1 & 1 & 1 \\ 0 & -4 & 2
            \end{pmatrix},\enspace\xi_1=(-1,1,-2)^\mathrm{T}$.
        \begin{enumerate}
            \item 求满足$A\xi_2=\xi_1$及$A^2\xi_3=\xi_1$的所有$\xi_2,\xi_3$；

            \item 证明：$\xi_1,\xi_2,\xi_3$线性无关.
        \end{enumerate}
\end{exercise}
\noindent 来源：\href{https://github.com/yhwu-is/Linear-Algebra-Left-Undone/blob/new/%E8%AE%B2%E4%B9%89/%E4%B8%93%E9%A2%98/7%20%E7%BA%BF%E6%80%A7%E6%98%A0%E5%B0%84%E7%9F%A9%E9%98%B5%E8%A1%A8%E7%A4%BA%E4%B8%8E%E7%9F%A9%E9%98%B5%E8%BF%90%E7%AE%97%E5%9F%BA%E7%A1%80.tex}{L7-B-4}。
\begin{solve}
\begin{enumerate}
                \item 对非齐次线性方程组$AX=\xi_1$，\\
                      
\[\bar{A}=\begin{pmatrix}
                              1  & -1 & -1 & -1 \\
                              -1 & 1  & 1  & 1  \\
                              0  & -4 & -2 & -2
                          \end{pmatrix}
                          \rightarrow\begin{pmatrix}
                              1 & -1 & -1       & -1       \\
                              0 & 1  & \frac 12 & \frac 12 \\
                              0 & 0  & 0        & 0
                          \end{pmatrix}
                          \rightarrow\begin{pmatrix}
                              1 & 0 & -\frac 12 & -\frac 12 \\
                              0 & 1 & \frac 12  & \frac 12  \\
                              0 & 0 & 0         & 0
                          \end{pmatrix}\]
，则\\
                      
\[\xi_2=C_1\begin{pmatrix}
                              \frac 12  \\
                              -\frac 12 \\
                              1
                          \end{pmatrix} + \begin{pmatrix}
                              -\frac 12 \\
                              \frac 12  \\
                              0
                          \end{pmatrix}=\dfrac 12\begin{pmatrix}
                              C_1 - 1  \\
                              -C_1 + 1 \\
                              2C_1
                          \end{pmatrix}\]
（其中$C_1$为任意常数）.\\
                      $A^2=\begin{pmatrix}
                              2  & 2  & 0 \\
                              -2 & -2 & 0 \\
                              4  & 4  & 0
                          \end{pmatrix}$，对齐次线性方程组$A^2X=\xi_1$，\\
                      
\[\bar{B}=\begin{pmatrix}
                              A^2 & \xi_1
                          \end{pmatrix}=\begin{pmatrix}
                              2  & 2  & 0 & 1  \\
                              -2 & -2 & 0 & 1  \\
                              4  & 4  & 0 & -2
                          \end{pmatrix}
                          \rightarrow\begin{pmatrix}
                              1 & 1 & 0 & -\frac 12 \\
                              0 & 0 & 0 & 0         \\
                              0 & 0 & 0 & 0
                          \end{pmatrix}\]
，\\
                      则$A^2X=\xi_1$的通解\\
                      
\[\xi_3=C_2\begin{pmatrix}
                              -1 \\
                              1  \\
                              0
                          \end{pmatrix}+C_3\begin{pmatrix}
                              0 \\
                              0 \\
                              1
                          \end{pmatrix}+\begin{pmatrix}
                              -\frac 12 \\
                              0         \\
                              0
                          \end{pmatrix}=\begin{pmatrix}
                              -C_2 - \frac 12 \\
                              C_2             \\
                              C_3
                          \end{pmatrix}\]
（其中$C_2, C_3$为任意常数）.
                \item
                      因为$|\xi_1,\xi_2,\xi_3|=\dfrac 12\begin{vmatrix}
                              -1 & C_1 - 1  & -C_2 - \frac 12 \\
                              1  & -C_1 + 1 & C_2             \\
                              -2 & 2C_1     & C_3
                          \end{vmatrix}=-\dfrac 12\neq 0$，\\
                      所以$\xi_1,\xi_2,\xi_3$线性无关.
            \end{enumerate}
\end{solve}

\begin{exercise}\label{prob:L7-B-5}
已知$V$为$n$维线性空间，$\sigma\in \mathcal{L}(V,V)$，且$\ker\sigma=\operatorname{im} \sigma$，证明：
        \begin{enumerate}
            \item $n$为偶数；

            \item 存在$V$的一组基$\alpha_1,\ldots,\alpha_n$使得
                  \[\sigma(\alpha_1,\ldots,\alpha_n)=(\alpha_1,\ldots,\alpha_n)\begin{pmatrix}
                          O & E_{\frac{n}{2}} \\ O & O
                      \end{pmatrix}.\]
        \end{enumerate}
\end{exercise}
\noindent 来源：\href{https://github.com/yhwu-is/Linear-Algebra-Left-Undone/blob/new/%E8%AE%B2%E4%B9%89/%E4%B8%93%E9%A2%98/7%20%E7%BA%BF%E6%80%A7%E6%98%A0%E5%B0%84%E7%9F%A9%E9%98%B5%E8%A1%A8%E7%A4%BA%E4%B8%8E%E7%9F%A9%E9%98%B5%E8%BF%90%E7%AE%97%E5%9F%BA%E7%A1%80.tex}{L7-B-5}。
\begin{solve}
证明：\begin{enumerate}
                \item 由于 $ \ker \sigma = \operatorname{im} \sigma $，由 $ \dim \operatorname{im} \sigma + \dim \ker \sigma = \dim V $ 可得.

                \item 设 $ \beta_1, \ldots, \beta_n $ 为 $ V $ 的一组基，则
                      \[ \operatorname{im} \sigma = \operatorname{span}(\sigma(\beta_1), \ldots, \sigma(\beta_n)) = \ker \sigma \]
                      设 $ \sigma(\beta_1), \ldots, \sigma(\beta_{\frac{n}{2}}) $ 为 $ \operatorname{im} \sigma $ 的基，则可以证明
                      \[ \sigma(\beta_1), \ldots, \sigma(\beta_{\frac{n}{2}}), \beta_1, \ldots, \beta_{\frac{n}{2}} \]
                      线性无关，且 $ \sigma $ 在此基下的矩阵即为所求的形式.
            \end{enumerate}
\end{solve}

\begin{exercise}\label{prob:L7-B-6}
若$f(x)$是$x$的实系数$m$次多项式：
        \[f(x)=a_mx^m+a_{m-1}x^{m-1}+\cdots+a_1x+a_0\]
        则有矩阵多项式：
        \[f(A)=a_mA^m+a_{m-1}A^{m-1}+\cdots+a_1A+a_0E\]
        其中 $A^0=E$.
        \begin{enumerate}
            \item 若$A$为对角矩阵$B=\begin{pmatrix}
                          \lambda_1 & 0 \\ 0 & \lambda_2
                      \end{pmatrix}$，证明：$f(A)=\begin{pmatrix}
                          f(\lambda_1) & 0 \\ 0 & f(\lambda_2)
                      \end{pmatrix}$；

            \item 若$A=PBP^{-1}$，证明：$f(A)=Pf(B)P^{-1}$.
        \end{enumerate}
\end{exercise}
\noindent 来源：\href{https://github.com/yhwu-is/Linear-Algebra-Left-Undone/blob/new/%E8%AE%B2%E4%B9%89/%E4%B8%93%E9%A2%98/7%20%E7%BA%BF%E6%80%A7%E6%98%A0%E5%B0%84%E7%9F%A9%E9%98%B5%E8%A1%A8%E7%A4%BA%E4%B8%8E%E7%9F%A9%E9%98%B5%E8%BF%90%E7%AE%97%E5%9F%BA%E7%A1%80.tex}{L7-B-6}。
\begin{solve}
\begin{enumerate}
                \item 因为 $A^k=\begin{pmatrix}\lambda_1^k & 0 \\ 0 & \lambda_2^k\end{pmatrix}$，所以 $f(A)=a_mA^m+a_{m-1}A^{m-1}+\cdots+a_1A+a_0E = \begin{pmatrix}f(\lambda_1) & 0 \\ 0 & f(\lambda_2)\end{pmatrix}$.
                \item $A=PBP^{-1}$，则 $A^2=(PBP^{-1})(PBP^{-1})=PB^2P^{-1}$. 由归纳法得 $A^k=PB^kP^{-1}$，于是
                      \begin{align*}
                          f(A) & = a_mPB^mP^{-1}+a_{m-1}PB^{m-1}P^{-1}+\cdots+a_0                         \\
                               & =P\begin{pmatrix}f(\lambda_1) & 0 \\ 0 & f(\lambda_2)\end{pmatrix}P^{-1} \\
                               & =Pf(B)P^{-1}
                      \end{align*}
            \end{enumerate}
\end{solve}

\begin{exercise}\label{prob:L7-B-14}
设$\mathbf{F}$为数域，$V_1=\{A\in\mathbf{F}^{n\times n} \mid A^\mathrm{T}=A\},\enspace
            V_2=\{A\in\mathbf{F}^{n\times n} \mid A^\mathrm{T}=-A\},\enspace V_3=\{A\in\mathbf{F}^{n\times n} \mid A\text{~为上三角矩阵}\}$.
        \begin{enumerate}
            \item 证明：$V_1,V_2,V_3$都是$\mathbf{F}^{n\times n}$的子空间；

            \item 证明：$\mathbf{F}^{n\times n}=V_1+V_3$但不为直和，$\mathbf{F}^{n\times n}=V_2\oplus V_3$.
        \end{enumerate}
\end{exercise}
\noindent 来源：\href{https://github.com/yhwu-is/Linear-Algebra-Left-Undone/blob/new/%E8%AE%B2%E4%B9%89/%E4%B8%93%E9%A2%98/7%20%E7%BA%BF%E6%80%A7%E6%98%A0%E5%B0%84%E7%9F%A9%E9%98%B5%E8%A1%A8%E7%A4%BA%E4%B8%8E%E7%9F%A9%E9%98%B5%E8%BF%90%E7%AE%97%E5%9F%BA%E7%A1%80.tex}{L7-B-14}。
\begin{solve}
三个集合均含零矩阵，而且它们的定义条件都对线性组合封闭：转置保持线性组合，上三角矩阵相加或数乘仍是上三角矩阵，所以它们是子空间。
对任意 $A=(a_{ij})$，定义对称矩阵 $S$，使其对角线及下三角部分与 $A$ 相同，即 $s_{ij}=a_{ij}$（$i\ge j$），$s_{ij}=a_{ji}$（$i<j$）。于是 $A-S$ 为严格上三角矩阵，说明 $V_1+V_3$ 是全空间。但 $I\in V_1\cap V_3$，故该和不是直和。
另定义反对称矩阵 $K$：$k_{ij}=a_{ij}$（$i>j$）、$k_{ij}=-a_{ji}$（$i<j$）、$k_{ii}=0$。那么 $A-K$ 是上三角矩阵，说明 $V_2+V_3$ 为全空间。一个既上三角又反对称的矩阵，其严格下三角元、严格上三角元均为零；数域上 $2\ne0$，对角元也为零。因此交为 $\{O\}$，从而 $\mathbf F^{n\times n}=V_2\oplus V_3$。
\end{solve}

\begin{exercise}\label{prob:L7-B-15}
设\[W_1=\left\{\begin{pmatrix}
            x & -x \\ y & z
        \end{pmatrix} \;\middle|\; x,y,z\in \mathbf{F} \right\},W_2=\left\{\begin{pmatrix}
            a & b \\ -a & c
        \end{pmatrix} \;\middle|\; a,b,c\in \mathbf{F} \right\}.\]

        \begin{enumerate}
            \item 证明：$W_1,W_2$是$\mathbf{M}_2(\mathbf{F})$的子空间，并求$\dim W_1,\dim W_2,\dim(W_1+W_2),\dim(W_1\cap W_2)$;

            \item 求$W_1\cap W_2$的一组基，并求$A=\begin{pmatrix}
                        3 & -3 \\ -3 & 1
                    \end{pmatrix}$关于这组基的坐标.
        \end{enumerate}
\end{exercise}
\noindent 来源：\href{https://github.com/yhwu-is/Linear-Algebra-Left-Undone/blob/new/%E8%AE%B2%E4%B9%89/%E4%B8%93%E9%A2%98/7%20%E7%BA%BF%E6%80%A7%E6%98%A0%E5%B0%84%E7%9F%A9%E9%98%B5%E8%A1%A8%E7%A4%BA%E4%B8%8E%E7%9F%A9%E9%98%B5%E8%BF%90%E7%AE%97%E5%9F%BA%E7%A1%80.tex}{L7-B-15}。
\begin{solve}
\begin{enumerate}
                \item 只需证明 $W_1$ 封闭即可. 对于 $A=\begin{pmatrix}x & -x \\ y & z\end{pmatrix},A'=\begin{pmatrix} x' & -x' \\ y' & z' \end{pmatrix}$ 有 
\[\lambda A+\mu A'=\begin{pmatrix}
                            \lambda x+\mu x' & -(\lambda x+\mu x') \\
                            \lambda y+\mu y  & \lambda z+\mu z
                        \end{pmatrix}\]
，则$W_1$封闭，$W_1$是$\mathbf{M}_2(\mathbf{F})$的子空间.	$W_2$同理可证.

                \item $W_1$的基：$B_1=\begin{pmatrix}1&-1\\0&0\end{pmatrix},B_2=\begin{pmatrix}0&0\\1&0\end{pmatrix},B_3=\begin{pmatrix}0&0\\0&1\end{pmatrix}$.

                    $W_2$的基：$B_1=\begin{pmatrix}1&0\\-1&0\end{pmatrix},B_2=\begin{pmatrix}0&1\\0&0\end{pmatrix},B_3=\begin{pmatrix}0&0\\0&1\end{pmatrix}$.

                    则$\dim W_1=3,\dim W_2=3$. 要求$\dim  (W_1+W_2)$，只需求$B_1,B_2,\ldots ,B_6$的极大无关组即可. 可知 $B_1,B_2,B_3,B_4$是极大线性无关组.	$\dim(W_1+W_2)=4$. 根据维数公式$\dim (W_1\cap W_2)=\dim W_1+\dim W_2-\dim (W_1+W_2)=2$.

                \item $W_1\cap W_2=\left\{\begin{pmatrix}x&-x\\-x&y\end{pmatrix} \mid x,y\in \mathbf{F}\right\}$. 则一组基为$E_1=\begin{pmatrix}1&-1\\-1&0\end{pmatrix},E_2=\begin{pmatrix}0&0\\0&1\end{pmatrix}$. $A$ 的坐标即为 $(3,1)$.
            \end{enumerate}
\end{solve}

\begin{exercise}\label{prob:L7-C-2}
设 $A,B,C$ 为二阶复方阵，且 $A,B,C$ 在 $\mathbf{M}_2(\mathbf{C})$ 中线性无关. 证明：存在$z_1,z_2,z_3 \in \mathbf{C}$使得 $z_1A+z_2B+z_3C$ 为可逆矩阵.
\end{exercise}
\noindent 来源：\href{https://github.com/yhwu-is/Linear-Algebra-Left-Undone/blob/new/%E8%AE%B2%E4%B9%89/%E4%B8%93%E9%A2%98/7%20%E7%BA%BF%E6%80%A7%E6%98%A0%E5%B0%84%E7%9F%A9%E9%98%B5%E8%A1%A8%E7%A4%BA%E4%B8%8E%E7%9F%A9%E9%98%B5%E8%BF%90%E7%AE%97%E5%9F%BA%E7%A1%80.tex}{L7-C-2}。
\begin{solve}
首先设
            \[A=\begin{pmatrix}a_1 & a_2 \\ a_3 & a_4\end{pmatrix},B=\begin{pmatrix}b_1 & b_2 \\ b_3 & b_4\end{pmatrix},C=\begin{pmatrix}c_1 & c_2 \\ c_3 & c_4\end{pmatrix}.\]
            由于 $A,B,C$ 在 $\mathbf{M}_2(\mathbf{C})$ 中线性无关，所以将 $A,B,C$ 的元素排为一列，可知矩阵
            \[\begin{pmatrix}a_1 & b_1 & c_1 \\ a_2 & b_2 & c_2 \\ a_3 & b_3 & c_3 \\ a_4 & b_4 & c_4\end{pmatrix}\]
            的秩为 3，这里不妨设前三个行向量线性无关，即有
            \[D=\begin{pmatrix}a_1 & b_1 & c_1 \\ a_2 & b_2 & c_2 \\ a_3 & b_3 & c_3\end{pmatrix}\]
            为可逆矩阵.

            另外，注意到对任意的 $x_1,x_2,x_3$，有
            \[x_1A+x_2B+x_3C=\begin{pmatrix}a_1x_1+b_1x_2+c_1x_3 & a_2x_1+b_2x_2+c_2x_3 \\ a_3x_1+b_3x_2+c_3x_3 & a_4x_1+b_4x_2+c_4x_3\end{pmatrix}.\]
            现在考虑方程组
            \[\begin{cases}a_1x_1+b_1x_2+c_1x_3 = 0 \\ a_2x_1+b_2x_2+c_2x_3 = 1 \\ a_3x_1+b_3x_2+c_3x_3 = 1\end{cases}\]
            其系数矩阵为 $D$，这是一个可逆矩阵，所以上述方程存在唯一解，不妨记为 $(x_1',x_2',x_3')$，此时就有 $\lvert x_1'A+x_2'B+x_3'C \rvert = \begin{vmatrix}0 & 1 \\ 1 & \ast\end{vmatrix} = -1$.

            所以 $x_1'A+x_2'B+x_3'C$ 为可逆矩阵.
\end{solve}

\begin{exercise}\label{prob:L7-例-1}
已知$\sigma \in \mathcal{L}(\mathbf{R}^3,\mathbf{R}^3)$且$\sigma(x_1,x_2,x_3)^{\mathrm{T}}=(x_1+x_2,x_1-x_3, x_2)^{\mathrm{T}}$
    \begin{enumerate}
        \item 求$\sigma$的像空间和核空间；

        \item 求$\sigma$关于$\mathbf{R}^3$自然基的矩阵.
    \end{enumerate}
\end{exercise}
\noindent 来源：\href{https://github.com/yhwu-is/Linear-Algebra-Left-Undone/blob/new/%E8%AE%B2%E4%B9%89/%E4%B8%93%E9%A2%98/7%20%E7%BA%BF%E6%80%A7%E6%98%A0%E5%B0%84%E7%9F%A9%E9%98%B5%E8%A1%A8%E7%A4%BA%E4%B8%8E%E7%9F%A9%E9%98%B5%E8%BF%90%E7%AE%97%E5%9F%BA%E7%A1%80.tex}{L7-例-1}。
\begin{solve}
\begin{enumerate}
        \item 求像空间和核空间的方法我们在之前已经介绍过，为了计算方便取$\mathbf{R}^3$的自然基$e_1,e_2,e_3$计算有：
              \[
                \operatorname{im}\sigma
                =\operatorname{span}(\sigma(e_1),\sigma(e_2),\sigma(e_3))
                =\operatorname{span}(
                        (1,1,0)^{\mathrm{T}},
                        (1,0,1)^{\mathrm{T}},
                        (0,-1,0)^{\mathrm{T}}
                    )
                =\mathbf{R}^3
              \]
              对于核空间，解方程$\sigma(\alpha)=0$即可. 我们也可以用更简洁的方式书写：
              \[
                \ker\sigma
                =\{
                    (x_1,x_2,x_3)^{\mathrm{T}} \mid
                    \sigma(x_1,x_2,x_3)^{\mathrm{T}}
                    =(0,0,0)^{\mathrm{T}}
                \}
                =\{(0,0,0)^{\mathrm{T}}\}
              \]
              即方程只有零解，核空间可以记为$\ker\sigma=\{0\}$（只含零元的空间的一般记法）.

        \item 根据原讲义“线性映射矩阵表示”处，应先写出$\sigma$在出发空间一组基（按题目要求是$\mathbf{R}^3$自然基）下的像，并将像表示为到达空间基（按题目要求是$\mathbf{R}^3$自然基）的线性组合，即
            \begin{gather*}
                \sigma(e_1) = (1,1,0)^{\mathrm{T}}
                =e_1+e_2=(e_1,e_2,e_3)\begin{pmatrix}
                    1 \\ 1 \\ 0
                \end{pmatrix} \\
                \sigma(e_2) = (1,0,1)^{\mathrm{T}}
                =e_1+e_3=(e_1,e_2,e_3)\begin{pmatrix}
                    1 \\ 0 \\ 1
                \end{pmatrix} \\
                \sigma(e_3) = (0,-1,0)^{\mathrm{T}}
                =-e_2=(e_1,e_2,e_3)\begin{pmatrix}
                    0 \\ -1 \\ 0
                \end{pmatrix}
            \end{gather*}
            接下来把坐标依次按列称矩阵就得到了本题需要求解的矩阵：
            \[
                \mathbf{M}(\sigma)=\begin{pmatrix}
                    1 & 1 & 0  \\
                    1 & 0 & -1 \\
                    0 & 1 & 0
                \end{pmatrix}
            \]
    \end{enumerate}
\end{solve}

\begin{exercise}\label{prob:L7-例-2}
设$\sigma\colon\mathbf{R}^2\to\mathbf{R}^2$是绕原点逆时针旋转$\theta$角的变换，求$\sigma$在$\mathbf{R}^2$的自然基下的矩阵表示.
\end{exercise}
\noindent 来源：\href{https://github.com/yhwu-is/Linear-Algebra-Left-Undone/blob/new/%E8%AE%B2%E4%B9%89/%E4%B8%93%E9%A2%98/7%20%E7%BA%BF%E6%80%A7%E6%98%A0%E5%B0%84%E7%9F%A9%E9%98%B5%E8%A1%A8%E7%A4%BA%E4%B8%8E%E7%9F%A9%E9%98%B5%E8%BF%90%E7%AE%97%E5%9F%BA%E7%A1%80.tex}{L7-例-2}。
\begin{solve}
求解的过程是很自然简单的，我们只需要考虑$\sigma$在常用基$e_1,e_2$下的像，即
    \[
    \begin{cases}
        \sigma(e_1)=\cos\theta\cdot e_1+\sin\theta\cdot e_2=(e_1,e_2)\begin{pmatrix}
            \cos\theta \\ \sin\theta
        \end{pmatrix} \\
        \sigma(e_2)=-\sin\theta\cdot e_1+\cos\theta\cdot e_2=(e_1,e_2)\begin{pmatrix}
            -\sin\theta \\ \cos\theta
        \end{pmatrix}
    \end{cases}
    \]
    故
    \[\mathbf{M}(\sigma)=\begin{pmatrix}
        \cos\theta & -\sin\theta \\
        \sin\theta & \cos\theta
    \end{pmatrix}.\]
\end{solve}

\begin{exercise}\label{prob:L7-例-3}
设$A=\begin{pmatrix}
            1 & 0 & -1 \\
            1 & 1 & -3
        \end{pmatrix}, B=\begin{pmatrix}
            0 & 3 \\
            1 & 2 \\
            3 & 1
        \end{pmatrix}$，求$AB$和$BA$.
\end{exercise}
\noindent 来源：\href{https://github.com/yhwu-is/Linear-Algebra-Left-Undone/blob/new/%E8%AE%B2%E4%B9%89/%E4%B8%93%E9%A2%98/7%20%E7%BA%BF%E6%80%A7%E6%98%A0%E5%B0%84%E7%9F%A9%E9%98%B5%E8%A1%A8%E7%A4%BA%E4%B8%8E%E7%9F%A9%E9%98%B5%E8%BF%90%E7%AE%97%E5%9F%BA%E7%A1%80.tex}{L7-例-3}。
\begin{solve}
直接根据矩阵乘法定义计算即可：
    \begin{gather*}
        AB = \begin{pmatrix}
            1 \times 0 + 0 \times 1 + (-1) \times 3 & 1 \times 3 + 0 \times 2 + (-1) \times 1 \\
            1 \times 0 + 1 \times 1 + (-3) \times 3 & 1 \times 3 + 1 \times 2 + (-3) \times 1
        \end{pmatrix} = \begin{pmatrix}
                -3 & 2 \\
                -8 & 2
            \end{pmatrix} \\
        BA = \begin{pmatrix}
            0 \times 1 + 3 \times 1 & 0 \times 0 + 3 \times 1 & 0 \times (-1) + 3 \times (-3) \\
            1 \times 1 + 2 \times 1 & 1 \times 0 + 2 \times 1 & 1 \times (-1) + 2 \times (-3) \\
            3 \times 1 + 1 \times 1 & 3 \times 0 + 1 \times 1 & 3 \times (-1) + 1 \times (-3)
        \end{pmatrix} = \begin{pmatrix}
                3 & 3 & -9 \\
                3 & 2 & -7 \\
                4 & 1 & -6
            \end{pmatrix}
    \end{gather*}
\end{solve}

\begin{exercise}\label{prob:L7-例-6}
展开矩阵多项式$(A+\lambda E)^n$.
\end{exercise}
\noindent 来源：\href{https://github.com/yhwu-is/Linear-Algebra-Left-Undone/blob/new/%E8%AE%B2%E4%B9%89/%E4%B8%93%E9%A2%98/7%20%E7%BA%BF%E6%80%A7%E6%98%A0%E5%B0%84%E7%9F%A9%E9%98%B5%E8%A1%A8%E7%A4%BA%E4%B8%8E%E7%9F%A9%E9%98%B5%E8%BF%90%E7%AE%97%E5%9F%BA%E7%A1%80.tex}{L7-例-6}。
\begin{solve}
由于$A$与$E$是可交换的，并且$A^nE^m=A^n$显然成立，因此我们结合中学学过的二项式展开，得到结果：
    \begin{align*}
        (A+\lambda E)^n & =\sum_{i=0}^nC_n^iA^i(\lambda E)^{n-i}    \\
                        & =\sum_{i=0}^nC_n^i\lambda^{n-i}A^iE^{n-i} \\
                        & =\sum_{i=0}^nC_n^i\lambda^{n-i}A^i.
    \end{align*}
\end{solve}

\begin{exercise}\label{prob:L7-例-7}
设$f(x),g(x) \in \mathbf{F}[x],\enspace A,B \in \mathbf{M}_n(\mathbf{F})$. 证明：
    \begin{enumerate}
        \item $f(A)g(A)=g(A)f(A)$；

        \item 如果$AB=BA$，则$f(A)g(B)=g(B)f(A)$；
    \end{enumerate}
\end{exercise}
\noindent 来源：\href{https://github.com/yhwu-is/Linear-Algebra-Left-Undone/blob/new/%E8%AE%B2%E4%B9%89/%E4%B8%93%E9%A2%98/7%20%E7%BA%BF%E6%80%A7%E6%98%A0%E5%B0%84%E7%9F%A9%E9%98%B5%E8%A1%A8%E7%A4%BA%E4%B8%8E%E7%9F%A9%E9%98%B5%E8%BF%90%E7%AE%97%E5%9F%BA%E7%A1%80.tex}{L7-例-7}。
\begin{solve}
我们可以直接证明第二点，因为第一点是第二点的特例. 设$f(x)=\displaystyle\sum_{i=0}^ma_ix^i$，$g(x)=\displaystyle\sum_{j=0}^sb_jx^j$，$A^0=B^0=E$，则
    \begin{align*}
        f(A)g(B) & =\sum_{i=0}^ma_iA^i\cdot \sum_{j=0}^sb_jB^j\quad(AB=BA)                            \\
                 & =\sum_{k=0}^{m+s}\sum_{i+j=k}a_ib_jA^iB^j=\sum_{k=0}^{m+s}\sum_{i+j=k}b_ja_iB^jA^i \\
                 & =g(B)f(A).
    \end{align*}
    事实上由于$A\cdot A=A\cdot A$，因此$f(A)g(A)=g(A)f(A)$只是上面证明的结论的特例.
\end{solve}

\begin{exercise}\label{prob:L7-例-9}
已知3阶矩阵$A=\begin{pmatrix}
            1 & 0 & 1 \\ 0 & -1 & 0 \\ -1 & 1 & -1
        \end{pmatrix}$. 定义$\mathbf{R}^{3 \times 3}$上的线性变换$\sigma(X)=AX,\enspace X \in \mathbf{R}^{3 \times 3}$. 求$\sigma$的像和核.
\end{exercise}
\noindent 来源：\href{https://github.com/yhwu-is/Linear-Algebra-Left-Undone/blob/new/%E8%AE%B2%E4%B9%89/%E4%B8%93%E9%A2%98/7%20%E7%BA%BF%E6%80%A7%E6%98%A0%E5%B0%84%E7%9F%A9%E9%98%B5%E8%A1%A8%E7%A4%BA%E4%B8%8E%E7%9F%A9%E9%98%B5%E8%BF%90%E7%AE%97%E5%9F%BA%E7%A1%80.tex}{L7-例-9}。
\begin{solve}
核空间求解较为简单，我们先求解核空间. 首先求解线性方程组$AY=0$，其中$Y$为列向量，解得其基础解系为$\eta=(1,0,-1)^\mathrm{T}$.

    记$X=(X_1,X_2,X_3)$，则$X\in\ker\sigma$即$AX=(AX_1,AX_2,AX_3)=O$，即$AX_1=AX_2=AX_3=0$，因此$X_1,X_2,X_3$都能由$\eta$线性表出，故
    \[X=(k_1\eta,k_2\eta,k_3\eta)=\begin{pmatrix}
            k_1 & k_2 & k_3 \\ 0 & 0 & 0 \\ -k_1 & -k_2 & -k_3
        \end{pmatrix},\enspace k_1,k_2,k_3\in\mathbf{R},\]
    即
\[X=k_1\begin{pmatrix}
            1 & 0 & 0 \\ 0 & 0 & 0 \\ -1 & 0 & 0
        \end{pmatrix}+k_2\begin{pmatrix}
            0 & 1 & 0 \\ 0 & 0 & 0 \\ 0 & -1 & 0
        \end{pmatrix}+k_3\begin{pmatrix}
            0 & 0 & 1 \\ 0 & 0 & 0 \\ 0 & 0 & -1
        \end{pmatrix}\]
，即$\ker\sigma$中所有元素均可由这三个矩阵线性表示，并且这三个矩阵显然线性无关，因此核空间就是这三个矩阵的线性组合，且核空间维数为3.

    注意到$\mathbf{R}^{3 \times 3}$的一组基为$E_{11},E_{12},E_{13},E_{21},E_{22},E_{23},E_{31},E_{32},E_{33}$，其中$E_{ij}$表示第$i$行第$j$列元素为1，其余元素为0的矩阵，例如$E_{23}=\begin{pmatrix}
            0 & 0 & 0 \\ 0 & 0 & 1 \\ 0 & 0 & 0
        \end{pmatrix}$.

    根据$\sigma$的定义我们可以求得
    \begin{align*}
        \sigma(E_{11})&=\sigma(E_{31})=\begin{pmatrix}
            1 & 0 & 0 \\ 0 & 0 & 0 \\ -1 & 0 & 0
        \end{pmatrix}, &
        \sigma(E_{12})&=\sigma(E_{32})=\begin{pmatrix}
            0 & 1 & 0 \\ 0 & 0 & 0 \\ 0 & -1 & 0
        \end{pmatrix}, \\
        \sigma(E_{13})&=\sigma(E_{33})=\begin{pmatrix}
            0 & 0 & 1 \\ 0 & 0 & 0 \\ 0 & 0 & -1
        \end{pmatrix}, &
        \sigma(E_{21})&=\begin{pmatrix}
            0 & 0 & 0 \\ -1 & 0 & 0 \\ 1 & 0 & 0
        \end{pmatrix}, \\
        \sigma(E_{22})&=\begin{pmatrix}
            0 & 0 & 0 \\ 0 & -1 & 0 \\ 0 & 1 & 0
        \end{pmatrix}, &
        \sigma(E_{23})&=\begin{pmatrix}
            0 & 0 & 0 \\ 0 & 0 & -1 \\ 0 & 0 & 1
        \end{pmatrix}.
    \end{align*}
    所以$\sigma$的像空间为上述六个矩阵线性扩张而成的空间，即\[\operatorname{im}\sigma=\operatorname{span}\{\sigma(E_{ij}):1\le i\le2,\ 1\le j\le3\}.\]. 又由原讲义“线性映射基本定理”处可知，$\dim\operatorname{im}\sigma=9-\dim\ker\sigma=6$，因此像空间就是这六个矩阵线性扩张而成的空间.
\end{solve}

\begin{exercise}\label{prob:L7-例-10}
用解方程法求矩阵 $A = \begin{pmatrix}1 & -1 & 1 \\ 0 & 1 & 2 \\ 1 & 0 & 4\end{pmatrix}$ 的逆矩阵.
\end{exercise}
\noindent 来源：\href{https://github.com/yhwu-is/Linear-Algebra-Left-Undone/blob/new/%E8%AE%B2%E4%B9%89/%E4%B8%93%E9%A2%98/7%20%E7%BA%BF%E6%80%A7%E6%98%A0%E5%B0%84%E7%9F%A9%E9%98%B5%E8%A1%A8%E7%A4%BA%E4%B8%8E%E7%9F%A9%E9%98%B5%E8%BF%90%E7%AE%97%E5%9F%BA%E7%A1%80.tex}{L7-例-10}。
\begin{solve}
以 $A$ 为系数矩阵的非齐次线性方程组 $AX = b$，对于任意的 $b = (b_1, b_2, b_3)$，可以用高斯-若当消元法将其增广矩阵

    \[(A,b)=\left(\begin{array}{ccc:c}
                1 & -1 & 1 & b_1 \\
                0 & 1  & 2 & b_2 \\
                1 & 0  & 4 & b_3
            \end{array}\right)\] 化为：

    \[\left(\begin{array}{ccc:c}
                1 & 0 & 0 & 4b_1+4b_2-3b_3 \\
                0 & 1 & 0 & 2b_1+3b_2-2b_3 \\
                0 & 0 & 1 & -b_1-b_2+b_3
            \end{array}\right)\].

    因此，对任意的 $b$，方程组 $AX=b$ 有唯一解：
    \[X=\begin{pmatrix}x_1\\x_2\\x_3\end{pmatrix}=\begin{pmatrix}
            4b_1+4b_2-3b_3 \\2b_1+3b_2-2b_3\\-b_1-b_2+b_3
        \end{pmatrix}=\begin{pmatrix}
            4  & 4  & -3 \\
            2  & 3  & -2 \\
            -1 & -1 & 1
        \end{pmatrix}\begin{pmatrix}b_1\\b_2\\b_3\end{pmatrix}\] 即：
    \[A^{-1}=\begin{pmatrix}
            4  & 4  & -3 \\
            2  & 3  & -2 \\
            -1 & -1 & 1
        \end{pmatrix}\]
\end{solve}

\begin{exercise}\label{prob:L7-例-11}
设$\alpha=(1,-1,2)^\mathrm{T},\enspace\beta=(3,1,-2)^\mathrm{T},\enspace A=\alpha\beta^\mathrm{T}$，求$A^n$.
\end{exercise}
\noindent 来源：\href{https://github.com/yhwu-is/Linear-Algebra-Left-Undone/blob/new/%E8%AE%B2%E4%B9%89/%E4%B8%93%E9%A2%98/7%20%E7%BA%BF%E6%80%A7%E6%98%A0%E5%B0%84%E7%9F%A9%E9%98%B5%E8%A1%A8%E7%A4%BA%E4%B8%8E%E7%9F%A9%E9%98%B5%E8%BF%90%E7%AE%97%E5%9F%BA%E7%A1%80.tex}{L7-例-11}。
\begin{solve}
由于$A^2=\alpha\beta^\mathrm{T}\alpha\beta^\mathrm{T}=\alpha(\beta^\mathrm{T}\alpha)\beta^\mathrm{T}=kA$，其中$k=\beta^\mathrm{T}\alpha=-2$，则$A^2=-2A$. 则可递推得到$A^n=(-2)^{n-1}A=(-2)^{n-1}\begin{pmatrix}
            3 & 1 & -2 \\ -3 & -1 & 2 \\ 6 & 2 & -4
        \end{pmatrix}$，原因在于$A^n$展开后中间会出现$n-1$个$\beta^\mathrm{T}\alpha$.
\end{solve}

\begin{exercise}\label{prob:L7-例-12}
证明以下几点性质：
    \begin{enumerate}
        \item 反对称矩阵主对角元均为 $0$；

        \item $AA^\mathrm{T}$ 和 $A^\mathrm{T}A$ 均为对称矩阵；

        \item 设 $A, B$ 分别为 $n$ 阶对称和反对称矩阵，则 $AB + BA$ 是反对称矩阵；

        \item 对称矩阵的乘积不一定对称；

        \item 可逆的对称（反对称）矩阵的逆矩阵也是对称（反对称）矩阵.
    \end{enumerate}
\end{exercise}
\noindent 来源：\href{https://github.com/yhwu-is/Linear-Algebra-Left-Undone/blob/new/%E8%AE%B2%E4%B9%89/%E4%B8%93%E9%A2%98/7%20%E7%BA%BF%E6%80%A7%E6%98%A0%E5%B0%84%E7%9F%A9%E9%98%B5%E8%A1%A8%E7%A4%BA%E4%B8%8E%E7%9F%A9%E9%98%B5%E8%BF%90%E7%AE%97%E5%9F%BA%E7%A1%80.tex}{L7-例-12}。
\begin{solve}
\begin{enumerate}
        \item 由于$A$为反对称矩阵，因此根据定义有 $a_{ii} = -a_{ii}$，即 $a_{ii} = 0$；

        \item 由于 $(AA^\mathrm{T})^\mathrm{T} = (A^\mathrm{T})^\mathrm{T}A^\mathrm{T} = AA^\mathrm{T}$，因此 $AA^\mathrm{T}$ 为对称矩阵；同理可证 $A^\mathrm{T}A$ 为对称矩阵；

        \item 由于$A,B$分别为对称和反对称矩阵，因此$(AB+BA)^\mathrm{T}=B^\mathrm{T}A^\mathrm{T}+A^\mathrm{T}B^\mathrm{T}=-AB-BA=-(AB+BA)$，因此$AB+BA$为反对称矩阵；

        \item 注意$(AB)^\mathrm{T}=B^\mathrm{T}A^\mathrm{T}=BA$，因为矩阵乘法不一定可交换，因此$AB$不一定对称；

        \item 因为$A$可逆有$(A^{-1})^\mathrm{T}=(A^\mathrm{T})^{-1}=A^{-1}$，因此$A^{-1}$为对称矩阵；同理可证$A$反对称的情况.
    \end{enumerate}
\end{solve}

\begin{exercise}\label{prob:L7-例-14}
$a,b,c,d$是四个实数. 证明$\begin{cases}
            a^2+b^2=1 \\
            c^2+d^2=1 \\
            ac+bd=0
        \end{cases}$成立的充分必要条件是$\begin{cases}
            a^2+c^2=1 \\
            b^2+d^2=1 \\
            ab+cd=0
        \end{cases}$.
\end{exercise}
\noindent 来源：\href{https://github.com/yhwu-is/Linear-Algebra-Left-Undone/blob/new/%E8%AE%B2%E4%B9%89/%E4%B8%93%E9%A2%98/7%20%E7%BA%BF%E6%80%A7%E6%98%A0%E5%B0%84%E7%9F%A9%E9%98%B5%E8%A1%A8%E7%A4%BA%E4%B8%8E%E7%9F%A9%E9%98%B5%E8%BF%90%E7%AE%97%E5%9F%BA%E7%A1%80.tex}{L7-例-14}。
\begin{solve}
设$A=\begin{pmatrix}
            a & b \\ c & d
        \end{pmatrix}$，则有
    \[AA^\mathrm{T}=\begin{pmatrix}
            a^2+b^2 & ac+bd \\ ac+bd & c^2+d^2
        \end{pmatrix},\enspace A^\mathrm{T}A=\begin{pmatrix}
            a^2+c^2 & ab+cd \\ ab+cd & b^2+d^2
        \end{pmatrix}.\]
    因此题中的充要条件可以转化为$AA^\mathrm{T}=E$是$A^\mathrm{T}A=E$的充要条件. 这是显然的，因为$AA^\mathrm{T}=E\iff A^{-1}=A^\mathrm{T}\iff A^\mathrm{T}A=E$成立.
\end{solve}

\begin{exercise}\label{prob:L7-例-15}
设$n$阶矩阵$A$分块为$A=\begin{pmatrix}
                      B & O \\ C & D
                  \end{pmatrix}$，其中$B,D$分别为$k$阶、$m$阶矩阵，求当$B,D$可逆时的$A^{-1}$.
\end{exercise}
\noindent 来源：\href{https://github.com/yhwu-is/Linear-Algebra-Left-Undone/blob/new/%E8%AE%B2%E4%B9%89/%E4%B8%93%E9%A2%98/7%20%E7%BA%BF%E6%80%A7%E6%98%A0%E5%B0%84%E7%9F%A9%E9%98%B5%E8%A1%A8%E7%A4%BA%E4%B8%8E%E7%9F%A9%E9%98%B5%E8%BF%90%E7%AE%97%E5%9F%BA%E7%A1%80.tex}{L7-例-15}。
\begin{solve}
本题我们使用的方法非常直接，就是直接设出$A^{-1}$的形式，然后验证即可. 之后我们还会学习一种基于分块矩阵初等变换的进阶方法（事实上考试如果考察的话基本是本题的解法，分块矩阵初等变换是在教材中是小字部分）. 设$A^{-1}=\begin{pmatrix}
                      X & Y \\ Z & T
                  \end{pmatrix}$，其中$X,T$分别为$k,m$阶矩阵，那么我们有
              \[\begin{pmatrix}
                      B & O \\ C & D
                  \end{pmatrix}\begin{pmatrix}
                      X & Y \\ Z & T
                  \end{pmatrix}=\begin{pmatrix}
                      BX & BY \\ CX+DZ & CY+DT
                  \end{pmatrix}=\begin{pmatrix}
                      E_k & O \\ O & E_m
                  \end{pmatrix},\]
              又由题意$A$可逆，$B,D$可逆，因此$BX=E_k$可得$X=B^{-1}$，$BY=O$可得$Y=O$，$CY+DT=DT=E_m$可得$T=D^{-1}$，$CX+DZ=CB^{-1}+DZ=O$可得$Z=-D^{-1}CB^{-1}$，因此
              \[A^{-1}=\begin{pmatrix}
                      B^{-1} & O \\ -D^{-1}CB^{-1} & D^{-1}
                  \end{pmatrix}.\]
\end{solve}

\begin{exercise}\label{prob:L7-例-16}
设\[A=\begin{pmatrix}
            1 & 2 & 0  & 0  & 0  \\
            2 & 5 & 0  & 0  & 0  \\
            0 & 0 & -2 & 1  & 0  \\
            0 & 0 & 0  & -2 & 1  \\
            0 & 0 & 0  & 0  & -2
        \end{pmatrix},\enspace B=\begin{pmatrix}
            1  & 0 & 1 & 0 \\
            -1 & 2 & 3 & 0 \\
            1  & 2 & 0 & 4 \\
            0  & 1 & 2 & 4 \\
            0  & 0 & 1 & 4
        \end{pmatrix}\]
    利用分块矩阵的方法，求$A^2,\enspace AB,\enspace A^\mathrm{T},\enspace A^{-1}$.
\end{exercise}
\noindent 来源：\href{https://github.com/yhwu-is/Linear-Algebra-Left-Undone/blob/new/%E8%AE%B2%E4%B9%89/%E4%B8%93%E9%A2%98/7%20%E7%BA%BF%E6%80%A7%E6%98%A0%E5%B0%84%E7%9F%A9%E9%98%B5%E8%A1%A8%E7%A4%BA%E4%B8%8E%E7%9F%A9%E9%98%B5%E8%BF%90%E7%AE%97%E5%9F%BA%E7%A1%80.tex}{L7-例-16}。
\begin{solve}
按照前两行和后三行分块，$A=\operatorname{diag}(A_1,A_2)$。乘积、转置与逆分别对对角块计算；矩阵 $B$ 的行按相同位置分块。直接计算得到
\[A^2=\begin{pmatrix}5&12&0&0&0\\12&29&0&0&0\\0&0&4&-4&1\\0&0&0&4&-4\\0&0&0&0&4\end{pmatrix},\qquad AB=\begin{pmatrix}-1&4&7&0\\-3&10&17&0\\-2&-3&2&-4\\0&-2&-3&-4\\0&0&-2&-8\end{pmatrix}.\]
\[A^{\mathrm T}=\begin{pmatrix}1&2&0&0&0\\2&5&0&0&0\\0&0&-2&0&0\\0&0&1&-2&0\\0&0&0&1&-2\end{pmatrix},\qquad A^{-1}=\begin{pmatrix}5&-2&0&0&0\\-2&1&0&0&0\\0&0&-\frac{1}{2}&-\frac{1}{4}&-\frac{1}{8}\\0&0&0&-\frac{1}{2}&-\frac{1}{4}\\0&0&0&0&-\frac{1}{2}\end{pmatrix}.\]
其中二阶块的行列式为 $1$，三阶上三角块的行列式为 $-8$，故两块均可逆；验证 $AA^{-1}=I_5$ 即可检验结果。
\end{solve}

\begin{exercise}\label{prob:S-BASIS-1}
在 $\mathbf R^3$ 中，基 $\mathcal B$、$\mathcal C$ 的向量分别为 $P=\begin{pmatrix}1&-3&0\\0&1&0\\1&0&1\end{pmatrix}$、$Q=\begin{pmatrix}1&0&0\\0&1&0\\0&-2&1\end{pmatrix}$ 的列。若 $[v]_{\mathcal B}=\begin{pmatrix}-1\\1\\3\end{pmatrix}$，求 $v$、$[v]_{\mathcal C}$ 及从 $\mathcal B$ 坐标到 $\mathcal C$ 坐标的变换矩阵。
\end{exercise}
\noindent 编者补充：S-BASIS-1；相关主题见原讲义对应讲，题目数据与解答由本选编另行构造。
\begin{solve}
两矩阵行列式都为 $1$，因此各自的列确实构成基。由坐标定义，$v=P[v]_{\mathcal B}$，所以 $v=\begin{pmatrix}-4\\1\\2\end{pmatrix}$。又有 $v=Q[v]_{\mathcal C}$，故
\[[v]_{\mathcal C}=Q^{-1}P[v]_{\mathcal B},\qquad Q^{-1}P=\begin{pmatrix}1&-3&0\\0&1&0\\1&2&1\end{pmatrix}.\]
代入题给坐标得 $[v]_{\mathcal C}=\begin{pmatrix}-4\\1\\4\end{pmatrix}$。验证 $Q[v]_{\mathcal C}=P[v]_{\mathcal B}$ 可确认变换方向。反向变换为 $P^{-1}Q$，它是 $Q^{-1}P$ 的逆矩阵；一般不能把这两个乘积交换。这里同一个几何向量的坐标随基变化，但向量自身未改变。
\end{solve}

\begin{exercise}\label{prob:S-BASIS-2}
在 $\mathbf R^3$ 中，基 $\mathcal B$、$\mathcal C$ 的向量分别为 $P=\begin{pmatrix}1&-1&-2\\0&1&2\\0&0&1\end{pmatrix}$、$Q=\begin{pmatrix}1&0&0\\0&1&1\\0&0&1\end{pmatrix}$ 的列。若 $[v]_{\mathcal B}=\begin{pmatrix}0\\0\\3\end{pmatrix}$，求 $v$、$[v]_{\mathcal C}$ 及从 $\mathcal B$ 坐标到 $\mathcal C$ 坐标的变换矩阵。
\end{exercise}
\noindent 编者补充：S-BASIS-2；相关主题见原讲义对应讲，题目数据与解答由本选编另行构造。
\begin{solve}
两矩阵行列式都为 $1$，因此各自的列确实构成基。由坐标定义，$v=P[v]_{\mathcal B}$，所以 $v=\begin{pmatrix}-6\\6\\3\end{pmatrix}$。又有 $v=Q[v]_{\mathcal C}$，故
\[[v]_{\mathcal C}=Q^{-1}P[v]_{\mathcal B},\qquad Q^{-1}P=\begin{pmatrix}1&-1&-2\\0&1&1\\0&0&1\end{pmatrix}.\]
代入题给坐标得 $[v]_{\mathcal C}=\begin{pmatrix}-6\\3\\3\end{pmatrix}$。验证 $Q[v]_{\mathcal C}=P[v]_{\mathcal B}$ 可确认变换方向。反向变换为 $P^{-1}Q$，它是 $Q^{-1}P$ 的逆矩阵；一般不能把这两个乘积交换。这里同一个几何向量的坐标随基变化，但向量自身未改变。
\end{solve}

\begin{exercise}\label{prob:S-BASIS-3}
在 $\mathbf R^3$ 中，基 $\mathcal B$、$\mathcal C$ 的向量分别为 $P=\begin{pmatrix}1&-2&-4\\1&-1&-2\\0&0&1\end{pmatrix}$、$Q=\begin{pmatrix}1&-2&1\\0&2&-1\\0&-1&1\end{pmatrix}$ 的列。若 $[v]_{\mathcal B}=\begin{pmatrix}1\\-1\\3\end{pmatrix}$，求 $v$、$[v]_{\mathcal C}$ 及从 $\mathcal B$ 坐标到 $\mathcal C$ 坐标的变换矩阵。
\end{exercise}
\noindent 编者补充：S-BASIS-3；相关主题见原讲义对应讲，题目数据与解答由本选编另行构造。
\begin{solve}
两矩阵行列式都为 $1$，因此各自的列确实构成基。由坐标定义，$v=P[v]_{\mathcal B}$，所以 $v=\begin{pmatrix}-9\\-4\\3\end{pmatrix}$。又有 $v=Q[v]_{\mathcal C}$，故
\[[v]_{\mathcal C}=Q^{-1}P[v]_{\mathcal B},\qquad Q^{-1}P=\begin{pmatrix}2&-3&-6\\1&-1&-1\\1&-1&0\end{pmatrix}.\]
代入题给坐标得 $[v]_{\mathcal C}=\begin{pmatrix}-13\\-1\\2\end{pmatrix}$。验证 $Q[v]_{\mathcal C}=P[v]_{\mathcal B}$ 可确认变换方向。反向变换为 $P^{-1}Q$，它是 $Q^{-1}P$ 的逆矩阵；一般不能把这两个乘积交换。这里同一个几何向量的坐标随基变化，但向量自身未改变。
\end{solve}

\begin{exercise}\label{prob:S-BASIS-4}
在 $\mathbf R^3$ 中，基 $\mathcal B$、$\mathcal C$ 的向量分别为 $P=\begin{pmatrix}1&0&-2\\0&1&-2\\0&0&1\end{pmatrix}$、$Q=\begin{pmatrix}1&0&0\\-1&1&1\\-3&2&3\end{pmatrix}$ 的列。若 $[v]_{\mathcal B}=\begin{pmatrix}2\\-2\\3\end{pmatrix}$，求 $v$、$[v]_{\mathcal C}$ 及从 $\mathcal B$ 坐标到 $\mathcal C$ 坐标的变换矩阵。
\end{exercise}
\noindent 编者补充：S-BASIS-4；相关主题见原讲义对应讲，题目数据与解答由本选编另行构造。
\begin{solve}
两矩阵行列式都为 $1$，因此各自的列确实构成基。由坐标定义，$v=P[v]_{\mathcal B}$，所以 $v=\begin{pmatrix}-4\\-8\\3\end{pmatrix}$。又有 $v=Q[v]_{\mathcal C}$，故
\[[v]_{\mathcal C}=Q^{-1}P[v]_{\mathcal B},\qquad Q^{-1}P=\begin{pmatrix}1&0&-2\\0&3&-7\\1&-2&3\end{pmatrix}.\]
代入题给坐标得 $[v]_{\mathcal C}=\begin{pmatrix}-4\\-27\\15\end{pmatrix}$。验证 $Q[v]_{\mathcal C}=P[v]_{\mathcal B}$ 可确认变换方向。反向变换为 $P^{-1}Q$，它是 $Q^{-1}P$ 的逆矩阵；一般不能把这两个乘积交换。这里同一个几何向量的坐标随基变化，但向量自身未改变。
\end{solve}

\begin{exercise}\label{prob:S-BASIS-5}
在 $\mathbf R^3$ 中，基 $\mathcal B$、$\mathcal C$ 的向量分别为 $P=\begin{pmatrix}1&4&3\\0&1&0\\0&2&1\end{pmatrix}$、$Q=\begin{pmatrix}1&0&0\\1&1&0\\-1&-2&1\end{pmatrix}$ 的列。若 $[v]_{\mathcal B}=\begin{pmatrix}3\\-3\\3\end{pmatrix}$，求 $v$、$[v]_{\mathcal C}$ 及从 $\mathcal B$ 坐标到 $\mathcal C$ 坐标的变换矩阵。
\end{exercise}
\noindent 编者补充：S-BASIS-5；相关主题见原讲义对应讲，题目数据与解答由本选编另行构造。
\begin{solve}
两矩阵行列式都为 $1$，因此各自的列确实构成基。由坐标定义，$v=P[v]_{\mathcal B}$，所以 $v=\begin{pmatrix}0\\-3\\-3\end{pmatrix}$。又有 $v=Q[v]_{\mathcal C}$，故
\[[v]_{\mathcal C}=Q^{-1}P[v]_{\mathcal B},\qquad Q^{-1}P=\begin{pmatrix}1&4&3\\-1&-3&-3\\-1&0&-2\end{pmatrix}.\]
代入题给坐标得 $[v]_{\mathcal C}=\begin{pmatrix}0\\-3\\-9\end{pmatrix}$。验证 $Q[v]_{\mathcal C}=P[v]_{\mathcal B}$ 可确认变换方向。反向变换为 $P^{-1}Q$，它是 $Q^{-1}P$ 的逆矩阵；一般不能把这两个乘积交换。这里同一个几何向量的坐标随基变化，但向量自身未改变。
\end{solve}

\begin{exercise}\label{prob:S-BASIS-6}
在 $\mathbf R^3$ 中，基 $\mathcal B$、$\mathcal C$ 的向量分别为 $P=\begin{pmatrix}1&0&0\\0&1&0\\-1&-1&1\end{pmatrix}$、$Q=\begin{pmatrix}1&0&0\\3&1&-2\\-1&0&1\end{pmatrix}$ 的列。若 $[v]_{\mathcal B}=\begin{pmatrix}4\\-4\\3\end{pmatrix}$，求 $v$、$[v]_{\mathcal C}$ 及从 $\mathcal B$ 坐标到 $\mathcal C$ 坐标的变换矩阵。
\end{exercise}
\noindent 编者补充：S-BASIS-6；相关主题见原讲义对应讲，题目数据与解答由本选编另行构造。
\begin{solve}
两矩阵行列式都为 $1$，因此各自的列确实构成基。由坐标定义，$v=P[v]_{\mathcal B}$，所以 $v=\begin{pmatrix}4\\-4\\3\end{pmatrix}$。又有 $v=Q[v]_{\mathcal C}$，故
\[[v]_{\mathcal C}=Q^{-1}P[v]_{\mathcal B},\qquad Q^{-1}P=\begin{pmatrix}1&0&0\\-3&-1&2\\0&-1&1\end{pmatrix}.\]
代入题给坐标得 $[v]_{\mathcal C}=\begin{pmatrix}4\\-2\\7\end{pmatrix}$。验证 $Q[v]_{\mathcal C}=P[v]_{\mathcal B}$ 可确认变换方向。反向变换为 $P^{-1}Q$，它是 $Q^{-1}P$ 的逆矩阵；一般不能把这两个乘积交换。这里同一个几何向量的坐标随基变化，但向量自身未改变。
\end{solve}

\section{矩阵相抵与秩}\label{sec:topic8}
本节收录 27 道题，其中 6 道为编者补充计算题。相抵变换保持秩，但通常不保持特征值。含参消元须区分主元为零的特殊参数，不能先除以待讨论的表达式。
\begin{exercise}\label{prob:L8-A-4}
给定$\mathbf{R}^4$的两组基
        \begin{gather*}
            \alpha_1=(1,1,1,1),\ \alpha_2=(1,1,-1,-1),\ \alpha_3=(1,-1,1,-1),\ \alpha_4=(1,-1,-1,1) \\
            \beta_1=(1,1,0,1),\ \beta_2=(2,1,3,1),\ \beta_3=(1,1,0,0),\ \beta_4=(0,1,-1,-1)
        \end{gather*}
        求由基$\alpha_1,\alpha_2,\alpha_3,\alpha_4$到基$\beta_1,\beta_2,\beta_3,\beta_4$的过渡矩阵，并求向量$\xi=(1,0,0,-1)$在基$\alpha_1,\alpha_2,\alpha_3,\alpha_4$下的坐标.
\end{exercise}
\noindent 来源：\href{https://github.com/yhwu-is/Linear-Algebra-Left-Undone/blob/new/%E8%AE%B2%E4%B9%89/%E4%B8%93%E9%A2%98/8%20%E7%9B%B8%E6%8A%B5%E6%A0%87%E5%87%86%E5%BD%A2.tex}{L8-A-4}。
\begin{solve}
首先取 $ \mathbf{R}^4 $ 的自然基 $\varepsilon_1, \varepsilon_2, \varepsilon_3, \varepsilon_4$，那么
            \[ (\alpha_1, \alpha_2, \alpha_3, \alpha_4) = (\varepsilon_1, \varepsilon_2, \varepsilon_3, \varepsilon_4) A,\enspace (\beta_1, \beta_2, \beta_3, \beta_4) = (\varepsilon_1, \varepsilon_2, \varepsilon_3, \varepsilon_4) B \]
            其中
            \[ A = \begin{pmatrix}
                    1 & 1  & 1  & 1  \\
                    1 & 1  & -1 & -1 \\
                    1 & -1 & 1  & -1 \\
                    1 & -1 & -1 & 1
                \end{pmatrix},\enspace B = \begin{pmatrix}
                    1 & 2 & 1 & 0  \\
                    1 & 1 & 1 & 1  \\
                    0 & 3 & 0 & -1 \\
                    1 & 1 & 0 & -1
                \end{pmatrix} \]
            由此可知
            \[ (\beta_1, \beta_2, \beta_3, \beta_4) = (\varepsilon_1, \varepsilon_2, \varepsilon_3, \varepsilon_4) B = (\alpha_1, \alpha_2, \alpha_3, \alpha_4) A^{-1} B \]
            所以由基 $ \alpha_1, \alpha_2, \alpha_3, \alpha_4 $ 到基 $ \beta_1, \beta_2, \beta_3, \beta_4 $ 的过渡矩阵为
            \[ A^{-1} B = \dfrac{1}{4} \begin{pmatrix}
                    1 & 1  & 1  & 1  \\
                    1 & 1  & -1 & -1 \\
                    1 & -1 & 1  & -1 \\
                    1 & -1 & -1 & 1
                \end{pmatrix} \begin{pmatrix}
                    1 & 2 & 1 & 0  \\
                    1 & 1 & 1 & 1  \\
                    0 & 3 & 0 & -1 \\
                    1 & 1 & 0 & -1
                \end{pmatrix} = \dfrac{1}{4} \begin{pmatrix}
                    3  & 7  & 2 & -1 \\
                    1  & -1 & 2 & 3  \\
                    -1 & 3  & 0 & -1 \\
                    1  & -1 & 0 & -1
                \end{pmatrix} \]
            另外，由于 $ \xi $ 在基 $ \varepsilon_1, \varepsilon_2, \varepsilon_3, \varepsilon_4 $ 下的坐标为 $ X_0 = (1, 0, 0, -1)^{\mathrm{T}} $，所以在 $ \alpha_1, \alpha_2, \alpha_3, \alpha_4 $ 下的坐标为
            \[ A^{-1} X_0 = \dfrac{1}{4} \begin{pmatrix}
                    1 & 1  & 1  & 1  \\
                    1 & 1  & -1 & -1 \\
                    1 & -1 & 1  & -1 \\
                    1 & -1 & -1 & 1
                \end{pmatrix} \begin{pmatrix}
                    1 \\ 0 \\ 0 \\ -1
                \end{pmatrix} = \begin{pmatrix}
                    0 \\ \dfrac{1}{2} \\ \dfrac{1}{2} \\ 0
                \end{pmatrix} \]
\end{solve}

\begin{exercise}\label{prob:L8-A-5}
证明：矩阵添加一列（或一行），其秩或不变，或增加1.
\end{exercise}
\noindent 来源：\href{https://github.com/yhwu-is/Linear-Algebra-Left-Undone/blob/new/%E8%AE%B2%E4%B9%89/%E4%B8%93%E9%A2%98/8%20%E7%9B%B8%E6%8A%B5%E6%A0%87%E5%87%86%E5%BD%A2.tex}{L8-A-5}。
\begin{solve}
证明：考虑矩阵 $A$ 的行向量组的极大线性无关组，若添加的一行可由其极大线性无关组线性表示，则秩不变. 否则秩增加 $1$.
\end{solve}

\begin{exercise}\label{prob:L8-A-6}
设$A$是$s \times n$矩阵，$B$是$A$前$m$行构成的$m \times n$矩阵，证明：$r(B) \geqslant r(A) + m - s$.
\end{exercise}
\noindent 来源：\href{https://github.com/yhwu-is/Linear-Algebra-Left-Undone/blob/new/%E8%AE%B2%E4%B9%89/%E4%B8%93%E9%A2%98/8%20%E7%9B%B8%E6%8A%B5%E6%A0%87%E5%87%86%E5%BD%A2.tex}{L8-A-6}。
\begin{solve}
证明：设 $A$ 的行向量组为 $\alpha_1,\alpha_2,\ldots,\alpha_s$，$r(A)=r$; $B$ 的行向量组为 $\alpha_1,\alpha_2,\ldots,\alpha_m ,r(B)=k$.
            不妨设：$B$ 的行向量组的极大线性无关组为 $\alpha_1,\alpha_2,\ldots,\alpha_k,\alpha_{i_1},\ldots,\alpha_{i_{r-k}}$，其中 $\alpha_{i_1},\ldots,\alpha_{i_{r-k}}$（共 $r-k$ 个向量）是包含在 $\alpha_{m+1},\ldots,\alpha_s$（共 $s-m$ 个向量）之中的. 显然有 $r-k \leqslant s-m $，即
            \[r(B)=k\geqslant r+m-s=r(A)+m-s.\]
\end{solve}

\begin{exercise}\label{prob:L8-B-2}
设$W$是$n$维线性空间$V$的一个非平凡子空间，$W$中取一组基$\delta_1,\ldots,\delta_m$，按如下两种方式将其扩充为$V$的一组基：
        \begin{align*}
            B_1 & =\{\delta_1,\ldots,\delta_m,\alpha_{m+1},\ldots,\alpha_n\} \\
            B_2 & =\{\delta_1,\ldots,\delta_m,\beta_{m+1},\ldots,\beta_n\}
        \end{align*}
        设基$B_1$到$B_2$的过渡矩阵为$P$，求商空间$V/W$的基$\alpha_{m+1}+W,\ldots,\alpha_n+W$到$\beta_{m+1}+W,\ldots,\beta_n+W$的过渡矩阵.
\end{exercise}
\noindent 来源：\href{https://github.com/yhwu-is/Linear-Algebra-Left-Undone/blob/new/%E8%AE%B2%E4%B9%89/%E4%B8%93%E9%A2%98/8%20%E7%9B%B8%E6%8A%B5%E6%A0%87%E5%87%86%E5%BD%A2.tex}{L8-B-2}。
\begin{solve}
$V$ 的基 $B_1$ 到 $B_2$ 的过渡矩阵 $P$ 具有下述形式：
          \[P=\begin{pmatrix}E_m & B_1 \\ 0 & B_2\end{pmatrix}\]
          其中 $B_1,B_2$ 分别是域 $\mathbf{F}$ 上 $m\times (n-m),(n-m)\times (n-m)$ 矩阵，
          \[\beta_j=b_{j1}\delta_1+\cdots+b_{jm}\delta_m+b_{j,m+1}\alpha_{m+1}+\cdots+b_{jn}\alpha_n\]
          其中 $j=m+1,\ldots,n$. 于是
          \[\beta_j+W=b_{j,m+1}(\alpha_{m+1}+W)+\cdots+b_{jn}(\alpha_n+W)\]
          因此商空间 $V/W$ 的基 $\alpha_{m+1}+W,\ldots,\alpha_n+W$ 到 $\beta_{m+1}+W,\ldots,\beta_n+W$ 的过渡矩阵是 $B_2$.
\end{solve}

\begin{exercise}\label{prob:L8-B-6}
证明：当$n$为奇数时，$\alpha_1,\alpha_2,\ldots,\alpha_n$线性无关的充要条件是$\alpha_1+\alpha_2,\alpha_2+\alpha_3,\ldots,\alpha_n+\alpha_1$线性无关.
\end{exercise}
\noindent 来源：\href{https://github.com/yhwu-is/Linear-Algebra-Left-Undone/blob/new/%E8%AE%B2%E4%B9%89/%E4%B8%93%E9%A2%98/8%20%E7%9B%B8%E6%8A%B5%E6%A0%87%E5%87%86%E5%BD%A2.tex}{L8-B-6}。
\begin{solve}
证明：我们设 $ \beta_1 = \alpha_1 + \alpha_2,\enspace \ldots, \beta_{n - 1} = \alpha_{n - 1} + \alpha_n,\enspace \beta_n = \alpha_n + \alpha_1 $. 由于
            \[ (\beta_1, \beta_2, \ldots, \beta_n) = (\alpha_1, \alpha_2, \ldots, \alpha_n) A \]
            其中
            \[ A = \begin{pmatrix}
                    1 &   &        &   & 1 \\
                    1 & 1 &        &   &   \\
                        & 1 & \ddots &   &   \\
                        &   & \ddots & 1 &   \\
                        &   &        & 1 & 1
                \end{pmatrix} \]
            满足 $ |A| = 1 + (-1)^{n + 1} = 2 $，所以 $ A $ 可逆. 根据教材原讲义命题 3.10.3 可得 $ \alpha_1, \alpha_2, \ldots, \alpha_n $ 与 $ \beta_1, \beta_2, \ldots, \beta_n $ 等价. 当然，这说明 $ \alpha_1, \alpha_2, \ldots, \alpha_n $ 线性无关的充要条件是 $ \beta_1, \beta_2, \ldots, \beta_n $ 线性无关.
\end{solve}

\begin{exercise}\label{prob:L8-B-7}
设$B=\{\beta_1,\beta_2,\ldots,\beta_n\}$是实数域$\mathbf{R}$上的线性空间$V$的一组基，$T \in \mathcal{L}(V),\enspace T(\beta_1)=\beta_2,T(\beta_2)=\beta_3,\ldots,T(\beta_{n-1}) = \beta_n,T(\beta_n)=\displaystyle\sum_{i=1}^{n}a_i\beta_i(a_i \in \mathbf{R})$，求$T$关于基$B$的表示矩阵，并求在什么条件下$T$是同构映射.
\end{exercise}
\noindent 来源：\href{https://github.com/yhwu-is/Linear-Algebra-Left-Undone/blob/new/%E8%AE%B2%E4%B9%89/%E4%B8%93%E9%A2%98/8%20%E7%9B%B8%E6%8A%B5%E6%A0%87%E5%87%86%E5%BD%A2.tex}{L8-B-7}。
\begin{solve}
$ T(\beta_1, \beta_2, \ldots, \beta_n) = (\beta_1, \beta_2, \ldots, \beta_n)A $，其中
            \[ A = \begin{pmatrix}
                      &   &        &   & a_1    \\
                    1 &   &        &   & a_2    \\
                      & 1 &        &   & a_3    \\
                      &   & \ddots &   & \vdots \\
                      &   &        & 1 & a_n
                \end{pmatrix} \]
            是 $ T $ 关于基 $ B $ 的表示矩阵.

            $ T $ 是同构 $ \iff T $ 是双射 $ \iff r(T) = n $，满秩. 所以当 $ a_1 \neq 0 $ 时，$ r(T) = n $ 满秩，此时 $ T $ 是同构映射.
\end{solve}

\begin{exercise}\label{prob:L8-B-8}
设矩阵空间$\mathbf{R}^{2\times 2}$的子空间为
        \[V=\{X=(x_{ij})_{2\times 2} \mid x_{11}+x_{12}+x_{21}=0,\enspace x_{ij}\in \mathbf{R}\}\]
        V中的线性变换为$\sigma(X)=X+X^\mathrm{T}$，求$V$的一组基，使得$\sigma$在该基下的矩阵表示为对角矩阵.
\end{exercise}
\noindent 来源：\href{https://github.com/yhwu-is/Linear-Algebra-Left-Undone/blob/new/%E8%AE%B2%E4%B9%89/%E4%B8%93%E9%A2%98/8%20%E7%9B%B8%E6%8A%B5%E6%A0%87%E5%87%86%E5%BD%A2.tex}{L8-B-8}。
\begin{solve}
我们可以类比线性空间
          \[ W = \{ x \in \mathbf{R}^4 \mid x_1 + x_2 + x_3 = 0 \} \]
          $ W $ 的基只需求解线性方程组 $ x_1 + x_2 + x_3 = 0 $ 即可，得到基础解系为 $ (-1, 0, 1, 0)^{\mathrm{T}},\allowbreak (-1, 1, 0, 0)^{\mathrm{T}},\allowbreak (0, 0, 0, 1)^{\mathrm{T}} $.

          换回 $ V $，即有基为 $ A_1 = \begin{pmatrix} -1 & 0 \\ 1 & 0 \end{pmatrix}, A_2 = \begin{pmatrix} -1 & 1 \\ 0 & 0 \end{pmatrix}, A_3 = \begin{pmatrix} 0 & 0 \\ 0 & 1 \end{pmatrix} $. 而
          \begin{align*}
                         & \sigma(A_1) = \begin{pmatrix} -2 & 1 \\ 1 & 0 \end{pmatrix} = A_1 + A_2                      \\
                         & \sigma(A_2) = \begin{pmatrix} -2 & 1 \\ 1 & 0 \end{pmatrix} = A_1 + A_2                      \\
                         & \sigma(A_3) = \begin{pmatrix} -2 & 1 \\ 1 & 0 \end{pmatrix} = 2 A_3                          \\
              \implies{} & \sigma(A_1 + A_2) = 2(A_1 + A_2),\enspace \sigma(A_1 - A_2) = 0,\enspace \sigma(A_3) = 2 A_3
          \end{align*}
          取基 $ A_1 - A_2, A_1 + A_2, A_3 $，有
          \[ (\sigma(A_1 - A_2), \sigma(A_1 + A_2), \sigma(A_3)) = (A_1 - A_2, A_1 + A_2, A_3) = \begin{pmatrix}
                  0 & 0 & 0 \\
                  0 & 2 & 0 \\
                  0 & 0 & 2
              \end{pmatrix} \]
          为对角矩阵.
\end{solve}

\begin{exercise}\label{prob:L8-B-9}
设$A$为$n$阶实反对称矩阵，证明：$E-A$ 可逆.
\end{exercise}
\noindent 来源：\href{https://github.com/yhwu-is/Linear-Algebra-Left-Undone/blob/new/%E8%AE%B2%E4%B9%89/%E4%B8%93%E9%A2%98/8%20%E7%9B%B8%E6%8A%B5%E6%A0%87%E5%87%86%E5%BD%A2.tex}{L8-B-9}。
\begin{solve}
反证法，假设不可逆，则存在非零实向量$x=(x_1,\cdots,x_n)^\mathrm{T}$使得$(E-A)x=0$，即$Ax=x$，又反对称矩阵有一个性质，即对任意的$x$，$x^\mathrm{T}Ax=0$，而$x^\mathrm{T}Ax=x^\mathrm{T}x=x_1^2+\cdots+x_n^2$，又由实数条件可知$x$为零向量，矛盾.
\end{solve}

\begin{exercise}\label{prob:L8-B-10}
设
        \[B_1=\left\{\begin{pmatrix}
                1 & 0 \\ 0 & 0
            \end{pmatrix},\begin{pmatrix}
                0 & 1 \\ 0 & 0
            \end{pmatrix},\begin{pmatrix}
                0 & 0 \\ 1 & 0
            \end{pmatrix}\begin{pmatrix}
                0 & 0 \\ 0 & 1
            \end{pmatrix}\right\},\]
        \[B_2=\left\{\begin{pmatrix}
                1 & 0 \\ 0 & 0
            \end{pmatrix},\begin{pmatrix}
                1 & 1 \\ 0 & 0
            \end{pmatrix},\begin{pmatrix}
                1 & 1 \\ 1 & 0
            \end{pmatrix}\begin{pmatrix}
                1 & 1 \\ 1 & 1
            \end{pmatrix}\right\}.\]
        \begin{enumerate}
            \item 证明：$B_2$也是线性空间$\mathbf{M}_2(\mathbf{R})$的基；

            \item 求基$B_2$变为基$B_1$的变换矩阵；

            \item 求$\mathbf{M}_2(\mathbf{R})$的一组基$B_3=\{A_1,A_2,A_3,A_4\}$，使得$A_i^2=A_i,\enspace i=1,2,3,4$；

            \item 已知矩阵$A$关于基$B_2$的坐标为$(1,1,1,1)^\mathrm{T}$，求$A$关于基$B_3$的坐标.
        \end{enumerate}
\end{exercise}
\noindent 来源：\href{https://github.com/yhwu-is/Linear-Algebra-Left-Undone/blob/new/%E8%AE%B2%E4%B9%89/%E4%B8%93%E9%A2%98/8%20%E7%9B%B8%E6%8A%B5%E6%A0%87%E5%87%86%E5%BD%A2.tex}{L8-B-10}。
\begin{solve}
记 $ B_1 = \{\vec{e}_{11}, \vec{e}_{12}, \vec{e}_{21}, \vec{e}_{22}\} $，$ B_2 = \{\vec{g}_1, \vec{g}_2, \vec{g}_3, \vec{g}_4\} $.
            \begin{enumerate}
                \item 设 $ k_1 \vec{g}_1 + k_2 \vec{g}_2 + k_3 \vec{g}_3 + k_4 \vec{g}_4 = O $，即
                      \[ \begin{pmatrix} k_1 + k_2 + k_3 + k_4 & k_2 + k_3 + k_4 \\ k_3 + k_4 & k_4 \end{pmatrix} = \begin{pmatrix} 0 & 0 \\ 0 & 0 \end{pmatrix} \]
                      这个方程对应的四元齐次线性方程组的解为
                      \[ k_1 = k_2 = k_3 = k_4 = 0 \]
                      所以 $ \vec{g}_1, \vec{g}_2, \vec{g}_3, \vec{g}_4 $ 线性无关，从而是 $ \mathbf{M}_2(\mathbf{R}) $ 的一组基.

                \item \label{item:11:B:4:2}
                      $ \mathbf{M}_2(\mathbf{R}) \cong \mathbf{R}^4 $，所以 $ \vec{e}_{11}, \vec{e}_{12}, \vec{e}_{21}, \vec{e}_{22} $ 可表示为 $ \mathbf{R}^4 $ 的自然基 $ \vec{e}_1, \vec{e}_2, \vec{e}_3, \vec{e}_4 $，而 $ \vec{g}_1, \vec{g}_2, \vec{g}_3, \vec{g}_4 $ 可表示为 $ (1, 0, 0, 0)^{\mathrm{T}},\allowbreak (1, 1, 0, 0)^{\mathrm{T}},\allowbreak (1, 1, 1, 0)^{\mathrm{T}},\allowbreak (1, 1, 1, 1)^{\mathrm{T}} $. 于是由
                      \[ (\vec{g}_1, \vec{g}_2, \vec{g}_3, \vec{g}_4) = (\vec{e}_{11}, \vec{e}_{12}, \vec{e}_{21}, \vec{e}_{22}) C \]
                      可得
                      \[ (\vec{e}_{11}, \vec{e}_{12}, \vec{e}_{21}, \vec{e}_{22}) = (\vec{g}_1, \vec{g}_2, \vec{g}_3, \vec{g}_4) C^{-1} \]
                      其中
                      \[ C = \begin{pmatrix}
                              1 & 1 & 1 & 1 \\
                              0 & 1 & 1 & 1 \\
                              0 & 0 & 1 & 1 \\
                              0 & 0 & 0 & 1
                          \end{pmatrix},\enspace C^{-1} = \begin{pmatrix}
                              1 & -1 & 0  & 0  \\
                              0 & 1  & -1 & 0  \\
                              0 & 0  & 1  & -1 \\
                              0 & 0  & 0  & 1
                          \end{pmatrix} \]
                      所以基 $ B_2 $ 变换为基 $ B_1 $ 的变换矩阵为 $ C^{-1} $.

                \item 在 $ A^2 = A $ 中选取较为简单的. 例如，由
                      \[ \begin{pmatrix} a & b \\ 0 & 0 \end{pmatrix}^2 = \begin{pmatrix} a^2 & ab \\ 0 & 0 \end{pmatrix} = \begin{pmatrix} a & b \\ 0 & 0 \end{pmatrix} \]
                      取 $ a = 1 $，$ b = 0 $ 或 $ b = 1 $ 得
                      \[ A_1 = \begin{pmatrix} 1 & 0 \\ 0 & 0 \end{pmatrix},\enspace A_2 = \begin{pmatrix} 1 & 1 \\ 0 & 0 \end{pmatrix} \]
                      由
                      \[ \begin{pmatrix} 0 & 0 \\ c & d \end{pmatrix}^2 = \begin{pmatrix} 0 & 0 \\ cd & d^2 \end{pmatrix} = \begin{pmatrix} 0 & 0 \\ c & d \end{pmatrix} \]
                      取 $ d = 1 $，$ c = 0 $ 或 $ c = 1 $ 得
                      \[ A_3 = \begin{pmatrix} 0 & 0 \\ 0 & 1 \end{pmatrix},\enspace A_4 = \begin{pmatrix} 0 & 0 \\ 1 & 1 \end{pmatrix} \]
                      上面的 $ A_1, A_2, A_3, A_4 $ 均满足 $ A_i^2 = A_i \enspace(i = 1, \ldots, 4) $，而且线性无关. 所以它们是 $ \mathbf{M}_2(\mathbf{R}) $ 的一组基 $ B_3 $.

                \item 先求基 $ B_2 $ 变换为基 $ B_3 $ 的变换矩阵 $ D $，即
                      \[ (A_1, A_2, A_3, A_4) = (\vec{g}_1, \vec{g}_2, \vec{g}_3, \vec{g}_4) D \]
                      由 \eqref{item:11:B:4:2} 所述，此时有
                      \[ \begin{pmatrix}
                              1 & 1 & 0 & 0 \\
                              0 & 1 & 0 & 0 \\
                              0 & 0 & 0 & 1 \\
                              0 & 0 & 1 & 1
                          \end{pmatrix} = CD \]
                      所以
                      \[ D = C^{-1} \begin{pmatrix}
                              1 & 1 & 0 & 0 \\
                              0 & 1 & 0 & 0 \\
                              0 & 0 & 0 & 1 \\
                              0 & 0 & 1 & 1
                          \end{pmatrix} = \begin{pmatrix}
                              1 & 0 & 0  & 0  \\
                              0 & 1 & 0  & -1 \\
                              0 & 0 & -1 & 0  \\
                              0 & 0 & 1  & 1
                          \end{pmatrix} \]
                      由于矩阵 $ A $ 关于基 $ B_2 $ 的坐标为 $ X = (1, 1, 1, 1)^{\mathrm{T}} $，所以 $ A $ 关于基 $ B_3 $ 的坐标为
                      \[ Y = D^{-1} X = \begin{pmatrix}
                              1 & 0 & 0  & 0  \\
                              0 & 1 & 0  & -1 \\
                              0 & 0 & -1 & 0  \\
                              0 & 0 & 1  & 1
                          \end{pmatrix}^{-1} \begin{pmatrix} 1 \\ 1 \\ 1 \\ 1 \end{pmatrix} = \begin{pmatrix}
                              1 & 0 & 0  & 0 \\
                              0 & 1 & 1  & 1 \\
                              0 & 0 & -1 & 0 \\
                              0 & 0 & 1  & 1
                          \end{pmatrix}^{-1} \begin{pmatrix} 1 \\ 1 \\ 1 \\ 1 \end{pmatrix} = \begin{pmatrix} 1 \\ 3 \\ -1 \\ 2 \end{pmatrix} \]
            \end{enumerate}
\end{solve}

\begin{exercise}\label{prob:L8-B-11}
利用相抵标准形证明以下结论：
        \begin{enumerate}
            \item 设$B_1,B_2$为$s \times n$列满秩矩阵，证明：存在$s$阶可逆矩阵$C$使得$B_2=CB_1$；

            \item 设$B_1,B_2$为$s \times n$行满秩矩阵，证明：存在$n$阶可逆矩阵$C$使得$B_2=B_1C$；

            \item 任意秩为$r$的矩阵都可以被分解为$r$个秩为1的矩阵之和；

            \item 已知$A$是$n$阶方阵，证明：存在$n$阶方阵$B$使得$A=ABA,\enspace B=BAB$.
        \end{enumerate}
\end{exercise}
\noindent 来源：\href{https://github.com/yhwu-is/Linear-Algebra-Left-Undone/blob/new/%E8%AE%B2%E4%B9%89/%E4%B8%93%E9%A2%98/8%20%E7%9B%B8%E6%8A%B5%E6%A0%87%E5%87%86%E5%BD%A2.tex}{L8-B-11}。
\begin{solve}
\begin{enumerate}
                \item 初等变换即可.

                \item 同上.

                \item 矩阵 $A$ 秩为 $r$ 可写作 $A=P\begin{pmatrix}E_r & 0 \\ 0 & 0\end{pmatrix}Q = P(E_{11}+E_{22}+\cdots+E_{rr})Q$（$E_r$ 是 $r\times r$ 的单位矩阵，$E_{ii}$ 是 $n\times n$ 的只有第 $i$ 行 $i$ 列的这个元素为 1，其他元素均为 0 的矩阵）. 每个 $PE_{ii}Q$ 都是秩为 1 的矩阵，故得证.

                \item 记 $r(A)=r$，把 $A$ 写成 $P\begin{pmatrix}E_r & 0 \\ 0 & 0\end{pmatrix}Q$ 的形式. 构造 $B=Q^{-1}\begin{pmatrix}E_r & 0 \\ 0 & 0\end{pmatrix}P^{-1}$ 可以发现其满足条件，故得证.
            \end{enumerate}
\end{solve}

\begin{exercise}\label{prob:L8-B-12}
设 $A \in \mathbf{M}_{m \times n}(\mathbf{F})$，$r(A)=r$，$k$ 是满足条件 $r \leqslant k \leqslant n$ 的任意整数，证明存在 $n$ 阶方阵 $B$，使得 $AB=O$，且 $r(A)+r(B)=k$.
\end{exercise}
\noindent 来源：\href{https://github.com/yhwu-is/Linear-Algebra-Left-Undone/blob/new/%E8%AE%B2%E4%B9%89/%E4%B8%93%E9%A2%98/8%20%E7%9B%B8%E6%8A%B5%E6%A0%87%E5%87%86%E5%BD%A2.tex}{L8-B-12}。
\begin{solve}
$r(A)=r$ 则 $AX=0$ 的解空间维数 $\dim N(A) = n-r$. 由 $r(A)+r(B)=k$ 得 $r(B)=k-r \leqslant n-r=\dim N(A)$. 要求 $AB=O$，说明 $B$ 的列向量均为 $AX=0$ 的解，那么只需要选择合适的列向量组拼接成 $B$ 即可（这一定能做到，因为 $B$ 维数不会超过解空间维数）.
\end{solve}

\begin{exercise}\label{prob:L8-B-13}
设$A$是$m \times n$矩阵（$m \leqslant n$），$r(A)=m$，证明：存在$n \times m$矩阵$B$使得$AB=E$.
\end{exercise}
\noindent 来源：\href{https://github.com/yhwu-is/Linear-Algebra-Left-Undone/blob/new/%E8%AE%B2%E4%B9%89/%E4%B8%93%E9%A2%98/8%20%E7%9B%B8%E6%8A%B5%E6%A0%87%E5%87%86%E5%BD%A2.tex}{L8-B-13}。
\begin{solve}
由于 $A$ 是 $m\times n$ 矩阵，$r(A)=m$，可知对于矩阵 $A$ 做初等列变换，可使其前 $m$ 列变为单位矩阵，后 $n-m$ 列变为全 0 列. 因此，存在 $n$ 阶可逆矩阵 $P$ 使得
          \[AP=\begin{pmatrix}E_m & O_{m\times (n-m)}\end{pmatrix}\]
          于是\[AP(AP)^{\mathrm{T}} = \begin{pmatrix}E & O\end{pmatrix} \begin{pmatrix}E \\ O\end{pmatrix}=E_m\]
          所以存在 $B=(PP^{\mathrm{T}}A^{\mathrm{T}})$ 为 $n\times m$ 矩阵，使 $AB=E$.
\end{solve}

\begin{exercise}\label{prob:L8-B-14}
设$A,B \in \mathbf{M}_n(\mathbf{F})$，$r(A)+r(B) \leqslant n$，证明：存在可逆矩阵$M$，使得$AMB=O$.
\end{exercise}
\noindent 来源：\href{https://github.com/yhwu-is/Linear-Algebra-Left-Undone/blob/new/%E8%AE%B2%E4%B9%89/%E4%B8%93%E9%A2%98/8%20%E7%9B%B8%E6%8A%B5%E6%A0%87%E5%87%86%E5%BD%A2.tex}{L8-B-14}。
\begin{solve}
利用 $A,B$ 的相抵标准形. 存在 $n$ 阶可逆矩阵 $P_1,Q_1,P_2,Q_2$ 使得
          \[P_1AQ_1=\begin{pmatrix}E_{r_A} & O \\ O & O\end{pmatrix},P_2BQ_2=\begin{pmatrix}O & O \\ O & E_{r_B}\end{pmatrix}\]
          于是 \[AQ_1=P_1^{-1}\begin{pmatrix}E_{r_A} & O \\ O & O\end{pmatrix},P_2B=\begin{pmatrix}O & O \\ O & E_{r_B}\end{pmatrix}Q_2^{-1}\]
          所以 \[AQ_1P_2B=P_1^{-1}\begin{pmatrix}E_{r_A} & O \\ O & O\end{pmatrix}\begin{pmatrix}O & O \\ O & E_{r_B}\end{pmatrix}Q_2^{-1}=O\]
          取 $M=Q_1P_2$ 即可.
\end{solve}

\begin{exercise}\label{prob:L8-B-16}
\begin{enumerate}
            \item 设 $A$ 和 $B$ 是两个 $m \times n$矩阵，证明：存在可逆矩阵 $M$ 使得$B = AM$ 的充要条件是 $A$ 和 $B$ 的列向量可以互相表示；
            \item 设 $A$ 和 $B$ 是两个 $m \times n$ 矩阵，且存在两个方阵 $C, D$ 使得 $A = BC, B = AD$，证明：存在可逆矩阵 $M$ 使得 $B = AM$.
        \end{enumerate}
\end{exercise}
\noindent 来源：\href{https://github.com/yhwu-is/Linear-Algebra-Left-Undone/blob/new/%E8%AE%B2%E4%B9%89/%E4%B8%93%E9%A2%98/8%20%E7%9B%B8%E6%8A%B5%E6%A0%87%E5%87%86%E5%BD%A2.tex}{L8-B-16}。
\begin{solve}
\begin{enumerate}
                \item 必要性：$B = AM$，则 $B$ 的每一列都为 $A$ 的列向量的线性组合，所以 $B$ 的列向量可以由 $A$ 的列向量表示. 反之，由于 $M$ 可逆，所以 $A = BM^{-1}$，则 $A$ 的每一列都为 $B$ 的列向量的线性组合，所以 $A$ 的列向量可以由 $B$ 的列向量表示.

                充分性：由于 $A$ 与 $B$ 的列向量组等价，则 $A$ 与 $B$ 的列空间相同，令其维数为 $r$. 对 $A, B$ 分别作初等列变换，存在可逆矩阵 $Q_1, Q_2 \in \mathbf{R}^{n\times n}$，使得
                \[ AQ_1 = (A_1, 0), \qquad BQ_2 = (B_1, 0), \]
                其中 $A_1, B_1 \in \mathbf{R}^{m\times r}$ 均列满秩。又因 $A$ 与 $B$ 的列空间相同，故 $A_1$ 与 $B_1$ 是同一列空间的两组基，因此存在可逆矩阵 $M_1 \in \mathbf{R}^{r\times r}$，使得 $B_1 = A_1M_1$。

                于是
                \[ BQ_2=(B_1,0) = (A_1M_1,0) = (A_1,0) \begin{pmatrix}
                    M_1 & 0 \\ 0 & I_{n-r}
                \end{pmatrix} = AQ_1 \begin{pmatrix}
                    M_1 & 0 \\ 0 & I_{n-r}
                \end{pmatrix}. \]
                因此
                \[ B = AQ_1 \begin{pmatrix}
                    M_1 & 0\\ 0 & I_{n-r}
                \end{pmatrix} Q_2^{-1}. \]
                令
                \[ M = Q_1 \begin{pmatrix}
                    M_1 & 0 \\ 0 & I_{n-r}
                    \end{pmatrix} Q_2^{-1}, \]
                则 $M$ 可逆，且 $B = AM$。

                \item 由于 $A = BC$，故 $A$ 的每一列都为 $B$ 的列向量的线性组合，同样由于$B = AD$，故 $B$ 的每一列都为 $A$ 的列向量的线性组合，所以 $A$ 与 $B$ 的列向量可以互相表示，因此它们的过渡矩阵可逆.
            \end{enumerate}
\end{solve}

\begin{exercise}\label{prob:L8-C-1}
设矩阵$A \in \mathbf{F}^{m \times n}$，$A$的秩$r(A)=r$，定义$\mathbf{F}^{n \times p}$到$\mathbf{F}^{m \times p}$的线性映射$\sigma$，使得$\forall X \in \mathbf{F}^{n \times p}$，$\sigma(X)=AX$. 求$\sigma$核空间的维数.
\end{exercise}
\noindent 来源：\href{https://github.com/yhwu-is/Linear-Algebra-Left-Undone/blob/new/%E8%AE%B2%E4%B9%89/%E4%B8%93%E9%A2%98/8%20%E7%9B%B8%E6%8A%B5%E6%A0%87%E5%87%86%E5%BD%A2.tex}{L8-C-1}。
\begin{solve}
本题直接求核空间困难，但由于我们只需求维数，所以我们转而求像空间维数并通过维数公式求核空间维数.

          取 $ \mathbf{F}^{n \times p} $ 的自然基 $ \vec{e}_{ij} = (e_{kl})_{n \times p} $，其中 $ e_{kl} = \delta_{ik} \delta_{jl} $，即第 $ i $ 行第 $ j $ 列元素为 $ 1 $，其余元素为 $ 0 $. 则 $ \operatorname{im} \sigma = \operatorname{span}(\sigma(\vec{e}_{11}), \ldots, \sigma(\vec{e}_{np})) $.

          取 $ A $ 的列向量组 $ A = (\alpha_1, \alpha_2, \ldots, \alpha_n) $，则可将 $ \sigma(\vec{e}_{ij}) $ 排列如下：
          \[ \begin{matrix}
                  (a_1, 0, \ldots, 0) & (0, a_1, \ldots, 0) & \cdots & (0, 0, \ldots, a_1) \\
                  (a_2, 0, \ldots, 0) & (0, a_2, \ldots, 0) & \cdots & (0, 0, \ldots, a_2) \\
                  \vdots              & \vdots              & \ddots & \vdots              \\
                  (a_n, 0, \ldots, 0) & (0, a_n, \ldots, 0) & \cdots & (0, 0, \ldots, a_n)
              \end{matrix} \]
          记 $ r(A) = r $，则 $ \alpha_1, \alpha_2, \ldots, \alpha_n $ 的极大线性无关组有 $ r $ 个向量，不妨设为 $ \alpha_1, \alpha_2, \ldots, \alpha_r $. 那么下列向量：
          \[ \begin{matrix}
                  (a_{r+1}, 0, \ldots, 0) & (0, a_{r+1}, \ldots, 0) & \cdots & (0, 0, \ldots, a_{r+1}) \\
                  \vdots                  & \vdots                  & \ddots & \vdots                  \\
                  (a_n, 0, \ldots, 0)     & (0, a_n, \ldots, 0)     & \cdots & (0, 0, \ldots, a_n)
              \end{matrix} \]

          均可以被其他向量线性表出. 观察除了上述向量的剩下的向量，可以发现这 $ rp $ 个向量线性无关，它们是 $ \operatorname{im} \sigma $ 的一组基，因此 $ \dim \operatorname{im} \sigma = rp $. 由维数公式，$ \dim \ker \sigma = \dim \mathbf{F}^{n \times p} - \dim \operatorname{im} \sigma = (n - r)p $.
\end{solve}

\begin{exercise}\label{prob:L8-例-2}
已知$\beta_1=\alpha_2+\alpha_3,\enspace\beta_2=\alpha_1+\alpha_3,\enspace\beta_3=\alpha_1+\alpha_2$，证明$\alpha_1,\alpha_2,\alpha_3$与$\beta_1,\beta_2,\beta_3$等价.
\end{exercise}
\noindent 来源：\href{https://github.com/yhwu-is/Linear-Algebra-Left-Undone/blob/new/%E8%AE%B2%E4%B9%89/%E4%B8%93%E9%A2%98/8%20%E7%9B%B8%E6%8A%B5%E6%A0%87%E5%87%86%E5%BD%A2.tex}{L8-例-2}。
\begin{solve}
事实上，$(\beta_1,\beta_2,\beta_3)=(\alpha_1,\alpha_2,\alpha_3)\begin{pmatrix}
            0 & 1 & 1 \\
            1 & 0 & 1 \\
            1 & 1 & 0
        \end{pmatrix}$，而$\begin{pmatrix}
            0 & 1 & 1 \\
            1 & 0 & 1 \\
            1 & 1 & 0
        \end{pmatrix}$可逆（可以由原讲义“可逆等价条件”处（3）（4）很容易验证），故$\alpha_1,\alpha_2,\alpha_3$与$\beta_1,\beta_2,\beta_3$等价.
\end{solve}

\begin{exercise}\label{prob:L8-例-3}
已知三维线性空间 $V$ 的线性变换 $\sigma$ 关于基 $\alpha_1,\alpha_2,\alpha_3$ 所对应的矩阵为
    \[\begin{pmatrix}1 & 2 & -1 \\ 2 & 1 & 0 \\ 3 & 0 & 1\end{pmatrix}\]
    \begin{enumerate}
        \item 求 $\sigma$ 在基 $\beta_1,\beta_2,\beta_3$ 下对应的矩阵 $B$，其中：
              \[\beta_1=2\alpha_1+\alpha_2+3\alpha_3,\enspace \beta_2=\alpha_1+\alpha_2+2\alpha_3,\enspace \beta_3=-\alpha_1+\alpha_2+\alpha_3;\]

        \item 求 $\sigma$ 的值域 $\sigma(V)$ 和核 $\ker\sigma$；

        \item 把 $\sigma(V)$ 的基扩充为 $V$ 的基，并求 $\sigma$ 在这组基下对应的矩阵；

        \item 把 $\ker\sigma$ 的基扩充为 $V$ 的基，并求 $\sigma$ 在这组基下对应的矩阵.
    \end{enumerate}
\end{exercise}
\noindent 来源：\href{https://github.com/yhwu-is/Linear-Algebra-Left-Undone/blob/new/%E8%AE%B2%E4%B9%89/%E4%B8%93%E9%A2%98/8%20%E7%9B%B8%E6%8A%B5%E6%A0%87%E5%87%86%E5%BD%A2.tex}{L8-例-3}。
\begin{solve}
\begin{enumerate}
        \item 对新的一组基，使用过渡矩阵进行表达如下：
              \[(\beta_{1}, \beta_{2}, \beta_{3})=(\alpha_{1}, \alpha_{2}, \alpha_{3})\begin{pmatrix}
                      2 & 1 & -1 \\
                      1 & 1 & 1  \\
                      3 & 2 & 1
                  \end{pmatrix}=(\alpha_{1}, \alpha_{2}, \alpha_{3}) C,\]
              其中 $C$ 是可逆矩阵，且
              \[(\alpha_{1},\alpha_{2},\alpha_{3})=(\beta_{1},\beta_{2},\beta_{3}) C^{-1},\]
              使用原讲义“基的选择对变换矩阵的影响”处可以直接得到 $B=C^{-1} A C$，或者根据原讲义“基与矩阵的"结合律"”处，将上式代入已知条件得
              \[\sigma((\beta_{1}, \beta_{2}, \beta_{3}) C^{-1})=((\beta_{1}, \beta_{2}, \beta_{3}) C^{-1}) A,\]
              \[(\sigma(\beta_{1}, \beta_{2}, \beta_{3})) C^{-1}=(\beta_{1}, \beta_{2}, \beta_{3})(C^{-1} A),\]
              从而得 $\sigma(\beta_{1}, \beta_{2}, \beta_{3})=(\beta_{1}, \beta_{2}, \beta_{3})(C^{-1} A C)$，故 $\sigma$ 关于基 $\beta_{1}, \beta_{2}, \beta_{3}$ 下对应的矩阵为
              \[B=C^{-1} A C = \begin{pmatrix}
                      2 & 1 & -1 \\
                      1 & 1 & 1  \\
                      3 & 2 & 1
                  \end{pmatrix}^{-1}\begin{pmatrix}
                      1 & 2 & -1 \\
                      2 & 1 & 0  \\
                      3 & 0 & 1
                  \end{pmatrix}\begin{pmatrix}
                      2 & 1 & -1 \\
                      1 & 1 & 1  \\
                      3 & 2 & 1
                  \end{pmatrix}=\begin{pmatrix}
                      2 & 0 & 1  \\
                      0 & 2 & 1  \\
                      3 & 1 & -1
                  \end{pmatrix}.\]

        \item  $\sigma$的值域是 $A$ 列向量组的极大线性无关组，由于 $ A $ 的第 1 列可以由第 2 列和第 3 列线性表示，从而 $\sigma(V)=\operatorname{span}(2 \alpha_{1}+\alpha_{2}, -\alpha_{1}+\alpha_{3})$.

        \item 由于 $\alpha_{1}$ 不能由 $2 \alpha_{1}+\alpha_{2}$ 和 $-\alpha_{1}+\alpha_{3}$ 线性表示，可以把 $\sigma(V)$ 的基扩充为 $V$ 的基 $\{\alpha_{1}, 2 \alpha_{1}+\alpha_{2}, -\alpha_{1}+\alpha_{3}\}$，$\sigma$ 在这个基下对应的矩阵是
              \[\begin{pmatrix}
                      1 & 2 & -1 \\
                      0 & 1 & 0  \\
                      0 & 0 & 1
                  \end{pmatrix}^{-1}\begin{pmatrix}
                      1 & 2 & -1 \\
                      2 & 1 & 0  \\
                      3 & 0 & 1
                  \end{pmatrix}\begin{pmatrix}
                      1 & 2 & -1 \\
                      0 & 1 & 0  \\
                      0 & 0 & 1
                  \end{pmatrix}=\begin{pmatrix}
                      0 & 0 & 0  \\
                      2 & 5 & -2 \\
                      3 & 6 & -2
                  \end{pmatrix}.\]

        \item 由于 $\alpha_{1}, \alpha_{2}$ 不能由 $\alpha_{1}-2 \alpha_{2}-3 \alpha_{3}$ 线性表示，可以把 $\ker \sigma$ 的基扩 充为 $V$ 的基 $\{\alpha_{1}, \alpha_{2}, \alpha_{1}-2 \alpha_{2}-3 \alpha_{3}\}$，$\sigma$ 在这个基下对应的矩阵是
              \[\begin{pmatrix}
                      1 & 0 & 1  \\
                      0 & 1 & -2 \\
                      0 & 0 & -3
                  \end{pmatrix}^{-1}\begin{pmatrix}
                      1 & 2 & -1 \\
                      2 & 1 & 0  \\
                      3 & 0 & 1
                  \end{pmatrix}\begin{pmatrix}
                      1 & 0 & 1  \\
                      0 & 1 & -2 \\
                      0 & 0 & -3
                  \end{pmatrix}=\begin{pmatrix}
                      2 & 2 & 0 \\
                      0 & 1 & 0 \\
                      1 & 0 & 0
                  \end{pmatrix}.\]
    \end{enumerate}
\end{solve}

\begin{exercise}\label{prob:L8-例-4}
已知矩阵 $A=\begin{pmatrix}a & b & c \\ d & e & f \\ h & x & y\end{pmatrix}$ 的逆是 $A^{-1}=\begin{pmatrix}-1 & -2 & -1 \\ 2 & 1 & 0 \\ 0 & -3 & -1\end{pmatrix}$，\\
    $B=\begin{pmatrix}a-2b & b-3c & -c \\ d-2e & e-3f & -f \\ h-2x & x-3y & -y\end{pmatrix}$. 求矩阵 $X$ 满足：

    \[X+\left(B(A^\mathrm{T}B^2)^{-1}A^\mathrm{T}\right)^{-1}=X\left(A^2(B^\mathrm{T}A)^{-1}B^\mathrm{T}\right)^{-1}(A+B)\]
\end{exercise}
\noindent 来源：\href{https://github.com/yhwu-is/Linear-Algebra-Left-Undone/blob/new/%E8%AE%B2%E4%B9%89/%E4%B8%93%E9%A2%98/8%20%E7%9B%B8%E6%8A%B5%E6%A0%87%E5%87%86%E5%BD%A2.tex}{L8-例-4}。
\begin{solve}
注意到
    \[B=A\begin{pmatrix}
            1 &  & \\ -2 & 1 & \\ & & 1
        \end{pmatrix} \begin{pmatrix}
            1 &  & \\ & 1 & \\ & -3 & 1
        \end{pmatrix} \begin{pmatrix}
            1 &  & \\ & 1 & \\ & & -1
        \end{pmatrix},\]
    从而 $B$ 由 $A$ 经过有限次初等变换得到，因 $A$ 可逆，故 $B$ 可逆. 一方面，等式左侧可以化简为
    \[X + \left(B\left(A^\mathrm{T} B^{2}\right)^{-1} A^\mathrm{T}\right)^{-1} = X + B,\]
    另一方面，等式右侧可以化简为
    \[X\left(A^{2}\left(B^\mathrm{T} A\right)^{-1} B^\mathrm{T}\right)^{-1}(A + B) = X A^{-1}(A + B) = X + X A^{-1} B,\]
    于是 
\[X = A = (A^{-1})^{-1} = \begin{pmatrix}
            -\dfrac{1}{3} & \dfrac{1}{3} & \dfrac{1}{3}  \\[0.5em]
            \dfrac{2}{3}  & \dfrac{1}{3} & -\dfrac{2}{3} \\
            -2           & -1          & 1
        \end{pmatrix}\]
.
\end{solve}

\begin{exercise}\label{prob:L8-例-6}
设$A=\begin{pmatrix}
            1 & 0 & 2 & -4 \\ 2 & 1 & 3 & -6 \\ -1 & -1 & -1 & 2
        \end{pmatrix}$. 求
    \begin{enumerate}[label=(\arabic*)]
        \item $A$的秩$r$和相抵标准形；

        \item
              3 阶可逆矩阵$P$和 4 阶可逆矩阵$Q$使得$PAQ=\begin{pmatrix}
                      E_r & 0 \\ 0 & 0
                  \end{pmatrix}$.
    \end{enumerate}
\end{exercise}
\noindent 来源：\href{https://github.com/yhwu-is/Linear-Algebra-Left-Undone/blob/new/%E8%AE%B2%E4%B9%89/%E4%B8%93%E9%A2%98/8%20%E7%9B%B8%E6%8A%B5%E6%A0%87%E5%87%86%E5%BD%A2.tex}{L8-例-6}。
\begin{solve}
\begin{enumerate}
        \item 利用求解极大线性无关组的方法可以解得$r=2$，因此对应的相抵标准形为
              \[\begin{pmatrix}
                      1 & 0 & 0 & 0 \\ 0 & 1 & 0 & 0 \\ 0 & 0 & 0 & 0
                  \end{pmatrix}.\]
        \item 这里介绍两种方法（严格来说，两种方法是一样的）.

              法一：直接将所有行变换相乘得到$P$，所有列变换相乘得到$Q$，此处略去步骤，给出答案为
              \[P=\begin{pmatrix}
                      1 & 0 & 0 \\ -2 & 1 & 0 \\ -1 & 1 & 1
                  \end{pmatrix},\enspace Q=\begin{pmatrix}
                      1 & 0 & -2 & 4 \\ 0 & 1 & 1 & -2 \\ 0 & 0 & 1  & 0 \\ 0 & 0 & 0 & 1
                  \end{pmatrix}.\]

              法二：我们用原讲义“初等变换求逆矩阵”处中的思想，将原讲义“相抵标准形另证”中的过程具象化. 具体来说，对 $m\times n$ 阶矩阵 $A$：
              \begin{gather}
                  \left(\begin{array}{c:c}
                      A & E_m
                  \end{array}\right)\xrightarrow{\text{初等行变换}}\left(\begin{array}{c:c}
                      A_1 & P
                  \end{array}\right)\\
                  \left(\begin{array}{c}
                      A_1 \\ \hdashline E_n
                  \end{array}\right)\xrightarrow{\text{初等列变换}}\left(\begin{array}{c}
                      U_r \\ \hdashline Q
                  \end{array}\right)
              \end{gather}
              其中 $A_1$ 为（行）简化阶梯矩阵，$P = P_s \cdots P_1$，$Q = Q_1 \cdots Q_t$为所求矩阵.

              以本题为例：
              \begin{gather*}
                  (A, E) =
                  \left(\begin{array}{cccc:cccc}
                          1  & 0  & 2  & -4 & 1 & 0 & 0 \\
                          2  & 1  & 3  & -6 & 0 & 1 & 0 \\
                          -1 & -1 & -1 & 2  & 0 & 0 & 1
                      \end{array}\right) \\
                  \xrightarrow{[2] + [1] \times (-2)}
                  \left(\begin{array}{cccc:cccc}
                          1  & 0  & 2  & -4 & 1  & 0 & 0 \\
                          0  & 1  & -1 & 2  & -2 & 1 & 0 \\
                          -1 & -1 & -1 & 2  & 0  & 0 & 1
                      \end{array}\right) \\
                  \xrightarrow{\substack{[3] + [1] \\ [3] + [2]}}
                  \left(\begin{array}{cccc:cccc}
                          1 & 0 & 2  & -4 & 1  & 0 & 0 \\
                          0 & 1 & -1 & 2  & -2 & 1 & 0 \\
                          0 & 0 & 0  & 0  & -1 & 1 & 1
                      \end{array}\right) \\
                  \left(\begin{array}{c}
                      A_1 \\ \hdashline E_n
                  \end{array}\right) =
                  \left(\begin{array}{cccc}
                      1 & 0 & 2  & -4 \\
                      0 & 1 & -1 & 2  \\
                      0 & 0 & 0  & 0  \\ \hdashline
                      1 & 0 & 0  & 0  \\
                      0 & 1 & 0  & 0  \\
                      0 & 0 & 1  & 0  \\
                      0 & 0 & 0  & 1
                  \end{array}\right)
                  \xrightarrow{\substack{[3] + [1] \times (-2) \\ [3] + [2] \\ [4] + [1] \times 4 \\ [4] + [2] \times (-2)}}
                  \left(\begin{array}{cccc}
                      1 & 0 & 0  & 0  \\
                      0 & 1 & 0  & 0  \\
                      0 & 0 & 0  & 0  \\ \hdashline
                      1 & 0 & -2 & 4  \\
                      0 & 1 & 1  & -2 \\
                      0 & 0 & 1  & 0  \\
                      0 & 0 & 0  & 1
                  \end{array}\right)
              \end{gather*}
              从而可以得到 $P$，$Q$，结果同上.
    \end{enumerate}
\end{solve}

\begin{exercise}\label{prob:L8-例-7}
设$\sigma:V\to W$为线性映射，取$V$和$W$的基$\alpha_1,\alpha_2,\ldots,\alpha_n$和$\beta_1,\beta_2,\ldots,\beta_m$，且$\sigma$在这两组基下的矩阵表示为$A=(a_{ij})_{m\times n}$.
    \begin{enumerate}
        \item 若将$\alpha_i$换为$c\alpha_i$（$c$为常数），求$\sigma$在新基下的矩阵表示；

        \item 若将$\beta_i$换为$c\beta_i$（$c$为常数），求$\sigma$在新基下的矩阵表示；

        \item 若将$\alpha_i$换为$\alpha_i+k\alpha_j$（$k$为常数），求$\sigma$在新基下的矩阵表示；

        \item 若将$\beta_i$换为$\beta_i+k\beta_j$（$k$为常数），求$\sigma$在新基下的矩阵表示；

        \item 若将$\alpha_i$和$\alpha_j$对换，求$\sigma$在新基下的矩阵表示；

        \item 若将$\beta_i$和$\beta_j$对换，求$\sigma$在新基下的矩阵表示.
    \end{enumerate}
\end{exercise}
\noindent 来源：\href{https://github.com/yhwu-is/Linear-Algebra-Left-Undone/blob/new/%E8%AE%B2%E4%B9%89/%E4%B8%93%E9%A2%98/8%20%E7%9B%B8%E6%8A%B5%E6%A0%87%E5%87%86%E5%BD%A2.tex}{L8-例-7}。
\begin{solve}
直接根据原讲义“换基公式”处，写出新旧基之间的过渡矩阵. 我们很容易发现这些过渡矩阵实际上就对应于六种初等矩阵，结果如下：
    \begin{enumerate}
        \item 矩阵表示$A_1$就是$A$的第$i$列数乘$c$；

        \item 矩阵表示$A_2$就是$A$的第$i$行数乘$1/c$；

        \item 矩阵表示$A_3$就是$A$的第$i$列加上第$j$列数乘$k$；

        \item 矩阵表示$A_4$就是$A$的第$j$行减去第$i$行数乘$k$；

        \item 矩阵表示$A_5$就是$A$的第$i$列和第$j$列对换；

        \item 矩阵表示$A_6$就是$A$的第$i$行和第$j$行对换.
    \end{enumerate}
\end{solve}

\begin{exercise}\label{prob:L8-例-8}
设$A$为$n$阶方阵，证明：
    \begin{enumerate}
        \item 存在可逆矩阵$B$和幂等矩阵$C$（即满足$C^2=C$）使得$A=BC$；
        \item 存在对称矩阵$B$和可逆矩阵$C$使得$A=BC$.
    \end{enumerate}
\end{exercise}
\noindent 来源：\href{https://github.com/yhwu-is/Linear-Algebra-Left-Undone/blob/new/%E8%AE%B2%E4%B9%89/%E4%B8%93%E9%A2%98/8%20%E7%9B%B8%E6%8A%B5%E6%A0%87%E5%87%86%E5%BD%A2.tex}{L8-例-8}。
\begin{solve}
根据相抵标准形分解，设$r(A)=r$，则我们有$A=P\begin{pmatrix}
            E_r & O \\ O & O
        \end{pmatrix}Q$，其中$P$和$Q$都是可逆矩阵.
    \begin{enumerate}
        \item 第一问难点在于取出这个幂等矩阵 $C$，实际上，我们可以取$C=Q^{-1}\begin{pmatrix}
                      E_r & O \\ O & O
                  \end{pmatrix}Q$，则$C^2=C$（因为中间的$QQ^{-1}$会抵消，而相抵标准形是幂等的）. 为了使得 $A = BC$，令 $B = PQ$ 即可（因为$P$ 和 $Q$ 都可逆，其乘积也必定可逆）.

        \item 第二问难点在于取出这个对称矩阵，但相信读者在经过上一小问的历练后，应该会产生一种取到对称矩阵的直觉. 因为标准形是对称的，我们只需要再多写一个$P^\mathrm{T}$，令$B=P\begin{pmatrix}
                      E_r & O \\ O & O
                  \end{pmatrix}P^\mathrm{T}$，则$B$是对称矩阵，且$A=BC$，令$C=(P^\mathrm{T})^{-1}Q$即可（注意矩阵可逆则转置也可逆）.
    \end{enumerate}
\end{solve}

\begin{exercise}\label{prob:S-PARAM-1}
设 $P=\begin{pmatrix}1&-3&0\\0&1&0\\1&0&1\end{pmatrix}$，$A(t)=P\operatorname{diag}(1,t-1,t-3)P^{-1}$，其中 $t\in\mathbf R$。求 $A(t)$ 的行列式和秩；分别求奇异参数下的核；判断方程 $A(t)x=\begin{pmatrix}-2\\1\\2\end{pmatrix}$ 何时有解。
\end{exercise}
\noindent 编者补充：S-PARAM-1；相关主题见原讲义对应讲，题目数据与解答由本选编另行构造。
\begin{solve}
相似变换保持行列式和秩，故
\[\det A(t)=(t-1)(t-3),\qquad r(A(t))=\begin{cases}3,&t\ne1,3,\\2,&t=1\text{ 或 }t=3.\end{cases}\]
当 $t=1$ 时，核由 $P$ 的第二列 $p_2=\begin{pmatrix}-3\\1\\0\end{pmatrix}$ 张成；当 $t=3$ 时，核由第三列 $p_3=\begin{pmatrix}0\\0\\1\end{pmatrix}$ 张成。令 $y=P^{-1}x$，题给右端在这个基下的坐标是 $(1,1,1)^{\mathrm T}$，方程成为
\[y_1=1,\qquad(t-1)y_2=1,\qquad(t-3)y_3=1.\]
所以奇异参数下方程无解，其余参数下唯一解为
\[x=P\begin{pmatrix}1\\1/(t-1)\\1/(t-3)\end{pmatrix}.\]
在讨论含参消元时必须把可能为零的主元单独分类；上面的基变换直接显示了所有特殊值。
\end{solve}

\begin{exercise}\label{prob:S-PARAM-2}
设 $P=\begin{pmatrix}1&-1&-2\\0&1&2\\0&0&1\end{pmatrix}$，$A(t)=P\operatorname{diag}(1,t-2,t-4)P^{-1}$，其中 $t\in\mathbf R$。求 $A(t)$ 的行列式和秩；分别求奇异参数下的核；判断方程 $A(t)x=\begin{pmatrix}-2\\3\\1\end{pmatrix}$ 何时有解。
\end{exercise}
\noindent 编者补充：S-PARAM-2；相关主题见原讲义对应讲，题目数据与解答由本选编另行构造。
\begin{solve}
相似变换保持行列式和秩，故
\[\det A(t)=(t-2)(t-4),\qquad r(A(t))=\begin{cases}3,&t\ne2,4,\\2,&t=2\text{ 或 }t=4.\end{cases}\]
当 $t=2$ 时，核由 $P$ 的第二列 $p_2=\begin{pmatrix}-1\\1\\0\end{pmatrix}$ 张成；当 $t=4$ 时，核由第三列 $p_3=\begin{pmatrix}-2\\2\\1\end{pmatrix}$ 张成。令 $y=P^{-1}x$，题给右端在这个基下的坐标是 $(1,1,1)^{\mathrm T}$，方程成为
\[y_1=1,\qquad(t-2)y_2=1,\qquad(t-4)y_3=1.\]
所以奇异参数下方程无解，其余参数下唯一解为
\[x=P\begin{pmatrix}1\\1/(t-2)\\1/(t-4)\end{pmatrix}.\]
在讨论含参消元时必须把可能为零的主元单独分类；上面的基变换直接显示了所有特殊值。
\end{solve}

\begin{exercise}\label{prob:S-PARAM-3}
设 $P=\begin{pmatrix}1&-2&-4\\1&-1&-2\\0&0&1\end{pmatrix}$，$A(t)=P\operatorname{diag}(1,t-3,t-5)P^{-1}$，其中 $t\in\mathbf R$。求 $A(t)$ 的行列式和秩；分别求奇异参数下的核；判断方程 $A(t)x=\begin{pmatrix}-5\\-2\\1\end{pmatrix}$ 何时有解。
\end{exercise}
\noindent 编者补充：S-PARAM-3；相关主题见原讲义对应讲，题目数据与解答由本选编另行构造。
\begin{solve}
相似变换保持行列式和秩，故
\[\det A(t)=(t-3)(t-5),\qquad r(A(t))=\begin{cases}3,&t\ne3,5,\\2,&t=3\text{ 或 }t=5.\end{cases}\]
当 $t=3$ 时，核由 $P$ 的第二列 $p_2=\begin{pmatrix}-2\\-1\\0\end{pmatrix}$ 张成；当 $t=5$ 时，核由第三列 $p_3=\begin{pmatrix}-4\\-2\\1\end{pmatrix}$ 张成。令 $y=P^{-1}x$，题给右端在这个基下的坐标是 $(1,1,1)^{\mathrm T}$，方程成为
\[y_1=1,\qquad(t-3)y_2=1,\qquad(t-5)y_3=1.\]
所以奇异参数下方程无解，其余参数下唯一解为
\[x=P\begin{pmatrix}1\\1/(t-3)\\1/(t-5)\end{pmatrix}.\]
在讨论含参消元时必须把可能为零的主元单独分类；上面的基变换直接显示了所有特殊值。
\end{solve}

\begin{exercise}\label{prob:S-PARAM-4}
设 $P=\begin{pmatrix}1&0&-2\\0&1&-2\\0&0&1\end{pmatrix}$，$A(t)=P\operatorname{diag}(1,t-4,t-6)P^{-1}$，其中 $t\in\mathbf R$。求 $A(t)$ 的行列式和秩；分别求奇异参数下的核；判断方程 $A(t)x=\begin{pmatrix}-1\\-1\\1\end{pmatrix}$ 何时有解。
\end{exercise}
\noindent 编者补充：S-PARAM-4；相关主题见原讲义对应讲，题目数据与解答由本选编另行构造。
\begin{solve}
相似变换保持行列式和秩，故
\[\det A(t)=(t-4)(t-6),\qquad r(A(t))=\begin{cases}3,&t\ne4,6,\\2,&t=4\text{ 或 }t=6.\end{cases}\]
当 $t=4$ 时，核由 $P$ 的第二列 $p_2=\begin{pmatrix}0\\1\\0\end{pmatrix}$ 张成；当 $t=6$ 时，核由第三列 $p_3=\begin{pmatrix}-2\\-2\\1\end{pmatrix}$ 张成。令 $y=P^{-1}x$，题给右端在这个基下的坐标是 $(1,1,1)^{\mathrm T}$，方程成为
\[y_1=1,\qquad(t-4)y_2=1,\qquad(t-6)y_3=1.\]
所以奇异参数下方程无解，其余参数下唯一解为
\[x=P\begin{pmatrix}1\\1/(t-4)\\1/(t-6)\end{pmatrix}.\]
在讨论含参消元时必须把可能为零的主元单独分类；上面的基变换直接显示了所有特殊值。
\end{solve}

\begin{exercise}\label{prob:S-PARAM-5}
设 $P=\begin{pmatrix}1&4&3\\0&1&0\\0&2&1\end{pmatrix}$，$A(t)=P\operatorname{diag}(1,t-5,t-7)P^{-1}$，其中 $t\in\mathbf R$。求 $A(t)$ 的行列式和秩；分别求奇异参数下的核；判断方程 $A(t)x=\begin{pmatrix}8\\1\\3\end{pmatrix}$ 何时有解。
\end{exercise}
\noindent 编者补充：S-PARAM-5；相关主题见原讲义对应讲，题目数据与解答由本选编另行构造。
\begin{solve}
相似变换保持行列式和秩，故
\[\det A(t)=(t-5)(t-7),\qquad r(A(t))=\begin{cases}3,&t\ne5,7,\\2,&t=5\text{ 或 }t=7.\end{cases}\]
当 $t=5$ 时，核由 $P$ 的第二列 $p_2=\begin{pmatrix}4\\1\\2\end{pmatrix}$ 张成；当 $t=7$ 时，核由第三列 $p_3=\begin{pmatrix}3\\0\\1\end{pmatrix}$ 张成。令 $y=P^{-1}x$，题给右端在这个基下的坐标是 $(1,1,1)^{\mathrm T}$，方程成为
\[y_1=1,\qquad(t-5)y_2=1,\qquad(t-7)y_3=1.\]
所以奇异参数下方程无解，其余参数下唯一解为
\[x=P\begin{pmatrix}1\\1/(t-5)\\1/(t-7)\end{pmatrix}.\]
在讨论含参消元时必须把可能为零的主元单独分类；上面的基变换直接显示了所有特殊值。
\end{solve}

\begin{exercise}\label{prob:S-PARAM-6}
设 $P=\begin{pmatrix}1&0&0\\0&1&0\\-1&-1&1\end{pmatrix}$，$A(t)=P\operatorname{diag}(1,t-6,t-8)P^{-1}$，其中 $t\in\mathbf R$。求 $A(t)$ 的行列式和秩；分别求奇异参数下的核；判断方程 $A(t)x=\begin{pmatrix}1\\1\\-1\end{pmatrix}$ 何时有解。
\end{exercise}
\noindent 编者补充：S-PARAM-6；相关主题见原讲义对应讲，题目数据与解答由本选编另行构造。
\begin{solve}
相似变换保持行列式和秩，故
\[\det A(t)=(t-6)(t-8),\qquad r(A(t))=\begin{cases}3,&t\ne6,8,\\2,&t=6\text{ 或 }t=8.\end{cases}\]
当 $t=6$ 时，核由 $P$ 的第二列 $p_2=\begin{pmatrix}0\\1\\-1\end{pmatrix}$ 张成；当 $t=8$ 时，核由第三列 $p_3=\begin{pmatrix}0\\0\\1\end{pmatrix}$ 张成。令 $y=P^{-1}x$，题给右端在这个基下的坐标是 $(1,1,1)^{\mathrm T}$，方程成为
\[y_1=1,\qquad(t-6)y_2=1,\qquad(t-8)y_3=1.\]
所以奇异参数下方程无解，其余参数下唯一解为
\[x=P\begin{pmatrix}1\\1/(t-6)\\1/(t-8)\end{pmatrix}.\]
在讨论含参消元时必须把可能为零的主元单独分类；上面的基变换直接显示了所有特殊值。
\end{solve}

\section{分块矩阵与矩阵运算进阶}\label{sec:topic9}
本节收录 40 道题。分块运算首先检查尺寸和乘法次序；求逆公式所需的可逆条件必须明确，证明一般恒等式时还要覆盖奇异情形。
\begin{exercise}\label{prob:L9-A-1}
设方阵$A$满足$A^2-A-2E=O$，证明：
        \begin{enumerate}
            \item $A$和$E-A$都是可逆矩阵，并求它们的逆矩阵；

            \item $A+E$和$A-2E$不可能同时可逆.
        \end{enumerate}
\end{exercise}
\noindent 来源：\href{https://github.com/yhwu-is/Linear-Algebra-Left-Undone/blob/new/%E8%AE%B2%E4%B9%89/%E4%B8%93%E9%A2%98/9%20%E7%9F%A9%E9%98%B5%E8%BF%90%E7%AE%97%E8%BF%9B%E9%98%B6.tex}{L9-A-1}。
\begin{solve}
\begin{enumerate}
                \item 若 $ AB = kE \enspace(k \neq 0) $，则 $ A^{-1} = \dfrac{1}{k} B $. 由
                      \[ A^2 - A - 2E = A(A - E) - 2E = O \]
                      可得
                      \[ A(A - E) = 2E \text{~或~} A(E - A) = -2E. \]
                      所以
                      \[ A^{-1} = \frac{1}{2}(A-E),\enspace (E - A)^{-1} = -\frac{1}{2}A. \]

                \item 由
                      \[  A^2 - A - 2E = (A - 2E)(A + E) = O \]
                      可知，$ A + E $ 和 $ A - 2E $ 不能同时可逆，否则 $ A^2 - A - 2E $ 为零矩阵可逆，矛盾.
            \end{enumerate}
\end{solve}

\begin{exercise}\label{prob:L9-A-3}
若$A,B$为两个$n$阶矩阵且满足$A+B=AB$，证明：
        \begin{enumerate}
            \item $A-E$和$B-E$均可逆；

            \item $AB=BA$；

            \item $r(A)=r(B)$.
        \end{enumerate}
\end{exercise}
\noindent 来源：\href{https://github.com/yhwu-is/Linear-Algebra-Left-Undone/blob/new/%E8%AE%B2%E4%B9%89/%E4%B8%93%E9%A2%98/9%20%E7%9F%A9%E9%98%B5%E8%BF%90%E7%AE%97%E8%BF%9B%E9%98%B6.tex}{L9-A-3}。
\begin{solve}
\begin{enumerate}
                \item \begin{align*}
                          (A - E)(B - E) & = AB - AE - EB + E^2 \\
                                         & = AB - A - B + E     \\
                                         & = AB - (A + B) + E   \\
                                         & = AB - AB + E        \\
                                         & = E.
                      \end{align*}
                      所以 $ A - E $ 与 $ B - E $ 互为逆矩阵.

                \item 由于 $ A - E $ 与 $ B - E $ 互为逆矩阵，所以
                      \begin{align*}
                          (B - E)(A - E) & = E              \\
                                         & = BA - B - A + E \\
                                         & = BA - AB + E.
                      \end{align*}
                      所以 $ BA - AB = O $，即 $ AB = BA $.

                \item 由 $ A = AB - B = B(A - E) $ 可得 $ r(A) \leqslant r(B) $，同理可得 $ r(B) \leqslant r(A) $，所以 $ r(A) = r(B) $.
            \end{enumerate}
\end{solve}

\begin{exercise}\label{prob:L9-A-4}
求下列矩阵的可逆的条件与逆矩阵：
            \begin{enumerate}
                \item $\begin{pmatrix}
                        A & B \\ O & D
                    \end{pmatrix}$，其中 $A$ 是 $m$ 阶方阵，$D$ 是 $n$ 阶方阵，$B$ 是 $m \times n$ 矩阵；

                \item $\begin{pmatrix}
                        O & B \\ C & D
                    \end{pmatrix}$，其中 $B$ 是 $n$ 阶方阵，$C$ 是 $m$ 阶方阵，$D$ 是 $n \times m$ 矩阵.
            \end{enumerate}
\end{exercise}
\noindent 来源：\href{https://github.com/yhwu-is/Linear-Algebra-Left-Undone/blob/new/%E8%AE%B2%E4%B9%89/%E4%B8%93%E9%A2%98/9%20%E7%9F%A9%E9%98%B5%E8%BF%90%E7%AE%97%E8%BF%9B%E9%98%B6.tex}{L9-A-4}。
\begin{solve}
\begin{enumerate}
                \item 对于矩阵 $\begin{pmatrix}
                        A & B \\ O & D
                    \end{pmatrix}$，

                    若 $A$ 不可逆，则 $r\begin{pmatrix} A \\ O \end{pmatrix} = r(A) < m$，
                    故 $r\begin{pmatrix} A & B \\ O & D \end{pmatrix} < m+n$，原矩阵不可逆.

                    若 $D$ 不可逆，则 $r\begin{pmatrix} O & D \end{pmatrix} = r(D) < n$，
                    故 $r\begin{pmatrix} A & B \\ O & D \end{pmatrix} < m+n$，原矩阵不可逆.

                    若 $A$ 与 $D$ 均可逆，利用分块矩阵初等变换，我们有
                    \[
                        \begin{pmatrix}
                            A^{-1} & O \\ O & D^{-1}
                        \end{pmatrix} \begin{pmatrix}
                            E_m & -BD^{-1} \\ O & E_n
                        \end{pmatrix} \begin{pmatrix}
                            A & B \\ O & D
                        \end{pmatrix} = E_{m+n},
                    \]
                    故有
                    \[
                        \begin{pmatrix}
                            A & B \\ O & D
                        \end{pmatrix}^{-1} = \begin{pmatrix}
                            A^{-1} & O \\ O & D^{-1}
                        \end{pmatrix} \begin{pmatrix}
                            E_m & -BD^{-1} \\ O & E_n
                        \end{pmatrix} = \begin{pmatrix}
                            A^{-1} & -A^{-1}BD^{-1} \\ O & D^{-1}
                        \end{pmatrix}.
                    \]

                \item 对于矩阵 $\begin{pmatrix}
                        O & B \\ C & D
                    \end{pmatrix}$，

                    由于
                    \[
                        \begin{pmatrix}
                            O & E_m \\ E_n & O
                        \end{pmatrix} \begin{pmatrix}
                            O & B \\ C & D
                        \end{pmatrix} = \begin{pmatrix}
                            C & D \\ O & B
                        \end{pmatrix},
                    \]

                    而 $\begin{pmatrix} O & E_m \\ E_n & O \end{pmatrix}$ 是可逆矩阵，故
                    \[
                        \begin{pmatrix}
                            O & B \\ C & D
                        \end{pmatrix} \, \text{可逆}
                        \iff \begin{pmatrix}
                            C & D \\ O & B
                        \end{pmatrix} \, \text{可逆}
                        \iff B, C \, \text{可逆}.
                    \]

                    若 $B$ 与 $C$ 均可逆，由第一问的结论，有
                    \[
                        \left( \begin{pmatrix}
                            O & E_m \\ E_n & O
                        \end{pmatrix} \begin{pmatrix}
                            O & B \\ C & D
                        \end{pmatrix} \right)^{-1} = \begin{pmatrix}
                            C^{-1} & -C^{-1}DB^{-1} \\ O & B^{-1}
                        \end{pmatrix},
                    \]

                    故
                    \[
                        \begin{pmatrix}
                            O & B \\ C & D
                        \end{pmatrix}^{-1} = \begin{pmatrix}
                            C^{-1} & -C^{-1}DB^{-1} \\ O & B^{-1}
                        \end{pmatrix} \begin{pmatrix}
                            O & E_m \\ E_n & O
                        \end{pmatrix} = \begin{pmatrix}
                            -C^{-1}DB^{-1} & C^{-1} \\ B^{-1} & O
                        \end{pmatrix}.
                    \]

            \end{enumerate}
\end{solve}

\begin{exercise}\label{prob:L9-B-1}
设$f(x)=1+x+\cdots+x^{m-1}$，$g(x)=1-x$，$A=\begin{pmatrix}
                a & b \\ 0 & a
            \end{pmatrix}$，计算$f(A)g(A)$.
\end{exercise}
\noindent 来源：\href{https://github.com/yhwu-is/Linear-Algebra-Left-Undone/blob/new/%E8%AE%B2%E4%B9%89/%E4%B8%93%E9%A2%98/9%20%E7%9F%A9%E9%98%B5%E8%BF%90%E7%AE%97%E8%BF%9B%E9%98%B6.tex}{L9-B-1}。
\begin{solve}
\begin{align*}
                f(A) g(A) & = (E + A + A^2 + \cdots + A^{m - 1})(E - A) \\
                          & = E - A^m = \begin{pmatrix}
                                            1 - a^m & -mba^{m - 1} \\
                                            0       & 1 - a^m
                                        \end{pmatrix}.
            \end{align*}
\end{solve}

\begin{exercise}\label{prob:L9-B-2}
已知矩阵$A=\begin{pmatrix}
                1 & 0 & 4 \\ 0 & 1 & 2 \\ 0 & 1 & 2
            \end{pmatrix}$，求证：所有与$A$可交换的矩阵构成$\mathbf{M}_3(\mathbf{R})$的一个子空间，并求子空间的一组基.
\end{exercise}
\noindent 来源：\href{https://github.com/yhwu-is/Linear-Algebra-Left-Undone/blob/new/%E8%AE%B2%E4%B9%89/%E4%B8%93%E9%A2%98/9%20%E7%9F%A9%E9%98%B5%E8%BF%90%E7%AE%97%E8%BF%9B%E9%98%B6.tex}{L9-B-2}。
\begin{solve}
求与 $ A $ 可交换的矩阵等价于求与 $ A - E $ 可交换的矩阵 $ X $. 可解得
            \begin{align*}
                X & = \begin{pmatrix}
                            x_{11} & -2x_{11} - 2x_{32} + 2x_{33} & 4x_{32} \\
                            0      & -x_{32} + x_{33}             & 2x_{32} \\
                            0      & x_{32}                       & x_{33}
                        \end{pmatrix} \\
                    & = x_{11} \begin{pmatrix}
                                1 & -2 & 0 \\
                                0 & 0  & 0 \\
                                0 & 0  & 0
                            \end{pmatrix}
                + x_{32} \begin{pmatrix}
                            0 & -2 & 4 \\
                            0 & -1 & 2 \\
                            0 & 1  & 0
                        \end{pmatrix}
                + x_{33} \begin{pmatrix}
                            0 & 2 & 0 \\
                            0 & 1 & 0 \\
                            0 & 0 & 1
                        \end{pmatrix},
            \end{align*}
            其中 $ x_{11}, x_{32}, x_{33} $ 为任意实数. 由此我们也得到了 $ A $ 可交换的矩阵构成的子空间的一组基.
\end{solve}

\begin{exercise}\label{prob:L9-B-3}
已知矩阵$A=\begin{pmatrix}
                1 & 0 & 1 \\ 0 & 2 & 0 \\ 1 & 0 & 1
            \end{pmatrix}$，
        \begin{enumerate}
            \item 求所有与$A$可交换的矩阵；

            \item 若$AB+E=A^2+B$，求$B$.
        \end{enumerate}
\end{exercise}
\noindent 来源：\href{https://github.com/yhwu-is/Linear-Algebra-Left-Undone/blob/new/%E8%AE%B2%E4%B9%89/%E4%B8%93%E9%A2%98/9%20%E7%9F%A9%E9%98%B5%E8%BF%90%E7%AE%97%E8%BF%9B%E9%98%B6.tex}{L9-B-3}。
\begin{solve}
\begin{enumerate}
                \item 不妨设 $ B = (b_{ij})_{3 \times 3} $ 与 $ A $ 可交换，即 $ AB = BA $，这等价于 $ (A - E)B = B(A - E) $，即
                      \[ \begin{pmatrix}
                                &   & 1 \\
                                & 1 &   \\
                              1 &   &
                          \end{pmatrix}
                          \begin{pmatrix}
                              b_{11} & b_{12} & b_{13} \\
                              b_{21} & b_{22} & b_{23} \\
                              b_{31} & b_{32} & b_{33}
                          \end{pmatrix} =
                          \begin{pmatrix}
                              b_{11} & b_{12} & b_{13} \\
                              b_{21} & b_{22} & b_{23} \\
                              b_{31} & b_{32} & b_{33}
                          \end{pmatrix}
                          \begin{pmatrix}
                                &   & 1 \\
                                & 1 &   \\
                              1 &   &
                          \end{pmatrix}. \]
                      于是
                      \[ \begin{pmatrix}
                              b_{31} & b_{32} & b_{33} \\
                              b_{21} & b_{22} & b_{23} \\
                              b_{11} & b_{12} & b_{13}
                          \end{pmatrix} = \begin{pmatrix}
                              b_{11} & b_{12} & b_{13} \\
                              b_{21} & b_{22} & b_{23} \\
                              b_{31} & b_{32} & b_{33}
                          \end{pmatrix} \]
                      对应元素相等，可得 $ b_{31} = b_{13}, b_{32} = b_{12}, b_{33} = b_{11}, b_{21} = b_{23} $，即所有与 $ A $ 可交换的矩阵为
                      \[ \begin{pmatrix}
                              b_{11} & b_{12} & b_{13} \\
                              b_{21} & b_{22} & b_{21} \\
                              b_{13} & b_{12} & b_{11}
                          \end{pmatrix}, \]
                      其中 $ b_{11}, b_{12}, b_{13}, b_{21}, b_{22} $ 为任意实数.

                \item 由于 $ AB + E = A^2 + B $，所以
                      \[ (A - E)B = A^2 - E = (A - E)(A + E). \]
                      又由于 $ |A - E| = -1 \neq 0 $，所以 $ A - E $ 可逆，进而
                      \[ B = A + E = \begin{pmatrix}
                              2 & 0 & 1 \\
                              0 & 3 & 0 \\
                              1 & 0 & 2
                          \end{pmatrix}. \]
            \end{enumerate}
\end{solve}

\begin{exercise}\label{prob:L9-B-4}
设$A \in \mathbf{F}^{n \times n}$，令$C(A)=\{B \in \mathbf{F}^{n \times n} \mid AB=BA\}$.
        \begin{enumerate}
            \item 证明：$C(A)$为$\mathbf{F}^{n \times n}$的一个子空间；

            \item 求$C(E)$；

            \item 当$A$为对角线上元素互不相等的对角阵时，求$C(A)$的维数和一组基.
        \end{enumerate}
\end{exercise}
\noindent 来源：\href{https://github.com/yhwu-is/Linear-Algebra-Left-Undone/blob/new/%E8%AE%B2%E4%B9%89/%E4%B8%93%E9%A2%98/9%20%E7%9F%A9%E9%98%B5%E8%BF%90%E7%AE%97%E8%BF%9B%E9%98%B6.tex}{L9-B-4}。
\begin{solve}
\begin{enumerate}
                \item 首先有 $ E \in C(A) $，所以 $ C(A) $ 非空. $ \forall B_1, B_2 \in C(A), \lambda \in \mathbf{F} $，有
                      \begin{gather*}
                          A(B_1 + B_2) = AB_1 + AB_2 = B_1A + B_2A = (B_1 + B_2)A \\
                          A(\lambda B_1) = \lambda AB_1 = \lambda B_1A = (\lambda B_1)A.
                      \end{gather*}
                      所以 $ C(A) $ 是 $ \mathbf{F}^{n \times n} $ 的子空间.

                \item $ C(E) = \mathbf{F}^{n \times n} $.

                \item 我们有结论：$ C(A) $ 为全体对角矩阵构成的集合. 故 $ \dim C(A) = n $. 基矩阵 $ B_k = (b_{ij})_{n \times n} $，其中 $ b_{ij} = \delta_{ij} \delta_{jk} $. $ B_1, B_2, \ldots, B_n $ 为一组基.
            \end{enumerate}
\end{solve}

\begin{exercise}\label{prob:L9-B-5}
设$V$是域$\mathbf{F}$上的$n$阶对称矩阵关于矩阵加法和数乘运算构成的线性空间，令
        \[U=\{A\in V \mid \operatorname{tr}(A)=0\},\enspace W=\{\lambda E \mid \lambda\in\mathbf{F}\}.\]
        \begin{enumerate}
            \item 证明：$U,W$为$V$的子空间；

            \item 分别求$U,W$的一组基和维数；

            \item 证明：$V=U\oplus W$.
        \end{enumerate}
\end{exercise}
\noindent 来源：\href{https://github.com/yhwu-is/Linear-Algebra-Left-Undone/blob/new/%E8%AE%B2%E4%B9%89/%E4%B8%93%E9%A2%98/9%20%E7%9F%A9%E9%98%B5%E8%BF%90%E7%AE%97%E8%BF%9B%E9%98%B6.tex}{L9-B-5}。
\begin{solve}
\begin{enumerate}
                \item 根据定义可知$\operatorname{tr}$是线性的，有 $\operatorname{tr}(\lambda A+\mu B)=\lambda\operatorname{tr}(A)+\mu\operatorname{tr}(B)=0$，则$U$是封闭的，又$W$封闭是显然的，则 $U,W$ 是 $V$ 的子空间.

                \item 对于$W$,基为$\{E\}$，$\dim W=1$. 对于$u$,设$E_{ij}$是 $a_{ij}=1$，其余元素为0的$n$阶方阵. 则由于$u$是对称矩阵，在非对角线元素上，$u$的基包含$E_{12}+E_{21},E_{13}+E_{31},\ldots ,E_{(n-1)n}+E_{n(n-1)}$. 对于对角线元素 $a_{11}+\cdots +a_{nn}=0$. 方程解为
                      \[\begin{pmatrix}a_{11}\\ \vdots\\ a_{nn}	\end{pmatrix}=k_1\begin{pmatrix}1 \\ -1 \\ 0\\ \vdots \\0	\end{pmatrix}=k_2\begin{pmatrix}1 \\ 0 \\ -1\\ \vdots \\0	\end{pmatrix}+\cdots+k_{n-1}\begin{pmatrix}1 \\ 0 \\ 0\\ \vdots \\-1	\end{pmatrix}.\]
                      则还有$n-1$个基$E_{11}-E_{22},E_{11}-E_{33},\ldots ,E_{11}-E_{nn}$. 故 $\dim U=\dfrac{n^2-n}2+n-2=\dfrac{n^2+n-2}2$.

                \item 设$V'=U+W$. 先证明直和，即$U\cap W=\{0\}$. 这是显然的，因为$\operatorname{tr}(\lambda E)=n\lambda$，仅当 $\lambda=0$ 时 $\lambda E\in U$ 即 $U\cap W=\{0\}$得证. 又$\dim U+\dim W=\dim V'=n=\dim V$，则$V=V'=U\oplus W$得证.
            \end{enumerate}
\end{solve}

\begin{exercise}\label{prob:L9-B-9}
设$A$是$n$阶矩阵，$A^k=O$对某个正整数$k$成立，求证下列方阵可逆，并求它们的逆：
        \begin{enumerate}
            \item $E+A$；

            \item $E-A$；

            \item $E+A+\dfrac{1}{2!}A^2+\cdots+\dfrac{1}{(k-1)!}A^{k-1}$.
        \end{enumerate}
\end{exercise}
\noindent 来源：\href{https://github.com/yhwu-is/Linear-Algebra-Left-Undone/blob/new/%E8%AE%B2%E4%B9%89/%E4%B8%93%E9%A2%98/9%20%E7%9F%A9%E9%98%B5%E8%BF%90%E7%AE%97%E8%BF%9B%E9%98%B6.tex}{L9-B-9}。
\begin{solve}
\begin{enumerate}
                \item 由 $ A^k = O $ 有
                      \[ E = E - A^k = (E - A)(E + A + \cdots + A^{k - 1}). \]
                      故 $ E - A $ 可逆，且
                      \[ (E - A)^{-1} = E + A + \cdots + A^{k - 1}. \]

                \item 若 $ k $ 为偶数，设 $ k = 2m $，由 $ A^k = O $ 有
                      \[ E = E - A^{2m} = (E + A)(E - A)(E + A^2 + A^4 + \cdots + A^{2m - 2}). \]
                      故 $ E + A $ 可逆，且
                      \[ (E + A)^{-1} = (E - A)(E + A^2 + A^4 + \cdots + A^{2m - 2}). \]

                      若 $ k $ 为奇数，设 $ k = 2m + 1 $，由 $ A^k = O $ 有
                      \[ E = E - A^{2m + 1} = (E - A)(E + A + A^2 + \cdots + A^{2m - 1} + A^{2m}) \]
                      故 $ E - A $ 可逆，且
                      \[ (E - A)^{-1} = E + A + A^2 + \cdots + A^{2m - 1} + A^{2m} \]

                \item 由于 $ A^k = O $，于是
                      \[ e^A = E + A + \frac{1}{2!} A^2 + \cdots + \frac{1}{(k - 1)!} A^{k - 1}. \]
                      由 $ e^A e^{-A} = E $ 知 $ E + A + \dfrac{1}{2!} A^2 + \cdots + \dfrac{1}{(k - 1)!} A^{k - 1} $ 可逆，且
                      \[ (E + A + \frac{1}{2!} A^2 + \cdots + \frac{1}{(k - 1)!} A^{k - 1})^{-1} = e^{-A}. \]
            \end{enumerate}
\end{solve}

\begin{exercise}\label{prob:L9-B-10}
设$A=\begin{pmatrix}
                1      & \cdots & 1      \\
                \vdots & \ddots & \vdots \\
                1      & \cdots & 1
            \end{pmatrix}_{n \times n}$，求：
        \begin{enumerate}
            \item 一个二次实系数多项式$f(x)=ax^2+bx$，使得$f(A)=O$；
            \item $A^{100}$；
            \item $(A+E)^3$；
            \item $(A+E)^{-1}$.
        \end{enumerate}
\end{exercise}
\noindent 来源：\href{https://github.com/yhwu-is/Linear-Algebra-Left-Undone/blob/new/%E8%AE%B2%E4%B9%89/%E4%B8%93%E9%A2%98/9%20%E7%9F%A9%E9%98%B5%E8%BF%90%E7%AE%97%E8%BF%9B%E9%98%B6.tex}{L9-B-10}。
\begin{solve}
\begin{enumerate}
                \item 
\[A^2 = \begin{pmatrix}
                        n      & \cdots & n      \\
                        \vdots & \ddots & \vdots \\
                        n      & \cdots & n
                    \end{pmatrix}\]
，故 $A^2 - nA = O$，令 $a=1, b=-n$ 即可.

                \item 利用 $A^2 = nA$，有
                    \[
                        A^{100} = n A^{99} = n^2 A^{98} = \cdots = n^{99} A = \begin{pmatrix}
                            n^{99} & \cdots & n^{99} \\
                            \vdots & \ddots & \vdots \\
                            n^{99} & \cdots & n^{99}
                        \end{pmatrix}.
                    \]

                \item 利用 $A^2 = nA$，有
                    \begin{align*}
                        (A+E)^3 &= A^3 + 3A^2 + 3A + E \\
                                &= (n^2 + 3n + 3)A + E \\
                                &= \begin{pmatrix}
                                    n^2 + 3n + 4 & n^2 + 3n + 3 & \cdots & n^2 + 3n + 3 \\
                                    n^2 + 3n + 3 & n^2 + 3n + 4 & \cdots & n^2 + 3n + 3 \\
                                    \vdots & \vdots & \ddots & \vdots \\
                                    n^2 + 3n + 3 & n^2 + 3n + 3 & \cdots & n^2 + 3n + 4
                                \end{pmatrix}.
                    \end{align*}

                \item 为得到 $(A+E)^{-1}$，我们需要构造另一个含 $A$ 的一次多项式，该式乘以 $A+E$ 后出现 $A^2 - nA$，以利用条件 $A^2 = nA$. 观察发现所需的多项式为 $A - (n+1)E$：
                    \[
                        (A+E)(A-(n+1)E) = A^2 - nA -(n+1)E = -(n+1)E,
                    \]

                    于是
                    \[
                        (A+E)^{-1} = -\frac{1}{n+1} A + E = \begin{pmatrix}
                            \dfrac{n}{n+1}  & -\dfrac{1}{n+1} & \cdots & -\dfrac{1}{n+1} \\[2em]
                            -\dfrac{1}{n+1} & \dfrac{n}{n+1}  & \cdots & -\dfrac{1}{n+1} \\[2em]
                            \vdots          & \vdots          & \ddots & \vdots          \\[2em]
                            -\dfrac{1}{n+1} & -\dfrac{1}{n+1} & \cdots & \dfrac{n}{n+1}
                        \end{pmatrix}.
                    \]
            \end{enumerate}
\end{solve}

\begin{exercise}\label{prob:L9-B-11}
设$A,B,A+B$都是可逆矩阵，求证：$A^{-1}+B^{-1}$也是可逆矩阵.
\end{exercise}
\noindent 来源：\href{https://github.com/yhwu-is/Linear-Algebra-Left-Undone/blob/new/%E8%AE%B2%E4%B9%89/%E4%B8%93%E9%A2%98/9%20%E7%9F%A9%E9%98%B5%E8%BF%90%E7%AE%97%E8%BF%9B%E9%98%B6.tex}{L9-B-11}。
\begin{solve}
这里我们遇到的加法和逆在一起的情况，因此我们使用提逆法，不难得到 $A^{-1}+B^{-1}=A^{-1}(A+B)B^{-1}$，由此可知其逆为 $B(A+B)^{-1}A$.
\end{solve}

\begin{exercise}\label{prob:L9-B-12}
已知矩阵$A=\begin{pmatrix}1 & 2 & a  \\
            1 & 3 & 0  \\
            2 & 7 & -a\end{pmatrix}$可以通过初等列变换转化为矩阵$B=\begin{pmatrix}1  & a & 2 \\
            0  & 1 & 1 \\
            -1 & 1 & 1\end{pmatrix}$.
        \begin{enumerate}
            \item 求常数$a$；
            \item 求满足$AP=B$的可逆矩阵$P$.
        \end{enumerate}
\end{exercise}
\noindent 来源：\href{https://github.com/yhwu-is/Linear-Algebra-Left-Undone/blob/new/%E8%AE%B2%E4%B9%89/%E4%B8%93%E9%A2%98/9%20%E7%9F%A9%E9%98%B5%E8%BF%90%E7%AE%97%E8%BF%9B%E9%98%B6.tex}{L9-B-12}。
\begin{solve}
\begin{enumerate}
                \item 由题意 $r(A)=r(B)$. $A$ 可由初等变换为 $\begin{pmatrix}1 & 2 & a \\ 1 & 3 & 0 \\ 0 & 0 & 0\end{pmatrix}$，$B$ 可由初等变换为 $\begin{pmatrix}1 & a & 2 \\ 0 & 1 & 1 \\ 0 & 0 & 2-a\end{pmatrix}$. 由于秩相同故 $2-a=0,a=2$.

                \item 
\[(A,B)=\begin{pmatrix}1 & 2 & a & 1 & a & 2 \\ 1 & 3 & 0 & 0 & 1 & 1 \\ 2 & 7 & -a & -1 & 1 & 1\end{pmatrix}\rightarrow\begin{pmatrix}1 & 0 & 6 & 3 & 4 & 4 \\ 0 & 1 & -2 & -1 & -1 & -1 \\ 0 & 0 & 0 & 0 & 0 & 0\end{pmatrix}\]


                        $X = \begin{pmatrix}-6k_1+3 \\ 2k_1-1 \\ k_1\end{pmatrix},Y = \begin{pmatrix}-6k_2+4 \\ 2k_2-1 \\ k_2\end{pmatrix},Z = \begin{pmatrix}-6k_3+4 \\ 2k_3-1 \\ k_3\end{pmatrix}$

                        $\begin{pmatrix}-6k_1+3 & -6k_2+4 & -6k_3+4 \\ 2k_1-1 & 2k_2-1 & 2k_3-1 \\ k_1 & k_2 & k_3\end{pmatrix}\rightarrow \begin{pmatrix}1 & 1 & 1 \\ 0 & 1 & 1 \\ 0 & 0 & k_3-k_2\end{pmatrix}$

                        因为 $P$ 可逆，所以 $k_2\neq k_3$，故 $P=\begin{pmatrix}-6k_1+3 & -6k_2+4 & -6k_3+4 \\ 2k_1-1 & 2k_2-1 & 2k_3-1 \\ k_1 & k_2 & k_3\end{pmatrix},k_2\neq k_3,\enspace k_1,k_2,k_3 \in \mathbf{R}$.
            \end{enumerate}
\end{solve}

\begin{exercise}\label{prob:L9-B-13}
利用列向量线性相关性，证明矩阵秩不等式：\[|r(A)-r(B)|\leqslant r(A\pm B) \leqslant r(A)+r(B).\]
\end{exercise}
\noindent 来源：\href{https://github.com/yhwu-is/Linear-Algebra-Left-Undone/blob/new/%E8%AE%B2%E4%B9%89/%E4%B8%93%E9%A2%98/9%20%E7%9F%A9%E9%98%B5%E8%BF%90%E7%AE%97%E8%BF%9B%E9%98%B6.tex}{L9-B-13}。
\begin{solve}
我们先证明 $r(A \pm B) \leqslant r(A)+r(B)$：

            设 $A = (\alpha_1, \alpha_2, \ldots, \alpha_n), B = (\beta_1, \beta_2, \ldots, \beta_n)$，并设 $r(A) = r, r(B) = s$. 不妨设 $A$ 的列向量组的一个极大线性无关组为 $\alpha_1, \alpha_2, \ldots, \alpha_r$，$B$ 的列向量组的一个极大线性无关组为 $\beta_1, \beta_2, \ldots, \beta_s$.

            取 $\alpha_1, \alpha_2, \ldots, \alpha_r, \beta_1, \beta_2, \ldots, \beta_s$ 的一个极大线性无关组 $\gamma_1, \gamma_2, \ldots, \gamma_p$，则 $p \leqslant r+s$. 该极大线性无关组可以线性表示 $A$ 与 $B$ 的列向量组，因此可以线性表示 $\alpha_1 \pm \beta_1, \alpha_2 \pm \beta_2, \ldots, \alpha_n \pm \beta_n$，即 $A \pm B$ 的列向量组，故
            \[
                r(A \pm B) \leqslant p \leqslant r+s = r(A) + r(B).
            \]

            再证明 $|r(A) - r(B)| \leqslant r(A \pm B)$：

            不妨设 $r(A) \geqslant r(B)$，即证 $r(A) - r(B) \leqslant r(A \pm B)$. 由于
            \begin{align*}
                r(A) &\leqslant r(A + B) + r(-B) = r(A + B) + r(B), \\
                r(A) &\leqslant r(A - B) + r(B),
            \end{align*}

            即有
            \[
                r(A) - r(B) \leqslant r(A \pm B).
            \]

            同理可得若 $r(A) \leqslant r(B)$，则 $r(B) - r(A) \leqslant r(A \pm B)$. 综上，结论成立.
\end{solve}

\begin{exercise}\label{prob:L9-B-14}
设$B$是$3 \times 1$矩阵，$C$是$1 \times 3$矩阵，证明：$r(BC) \leqslant 1$. 反之，若$A$是秩为1的$3 \times 3$矩阵，证明：存在$3 \times 1$矩阵$B$和$1 \times 3$矩阵$C$，使得$A = BC$.
\end{exercise}
\noindent 来源：\href{https://github.com/yhwu-is/Linear-Algebra-Left-Undone/blob/new/%E8%AE%B2%E4%B9%89/%E4%B8%93%E9%A2%98/9%20%E7%9F%A9%E9%98%B5%E8%BF%90%E7%AE%97%E8%BF%9B%E9%98%B6.tex}{L9-B-14}。
\begin{solve}
若 $B$ 是 $3 \times 1$ 矩阵，$C$ 是 $1 \times 3$ 矩阵，利用秩不等式 $r(BC) \leqslant \min\{r(B), r(C)\}$，以及 $r(B) \leqslant 1, r(C) \leqslant 1$，有
            \[
                r(BC) \leqslant 1.
            \]

            反之，若 $A$ 是秩为 $1$ 的 $3 \times 3$ 矩阵，则 $A$ 任意两个列向量线性相关. 任取其中的一个非零列向量 $\beta$（$r(A)=1$ 保证其存在），则 $A$ 可表示为
            \[
                A = (a_1 \beta, a_2 \beta, a_3 \beta) = \beta (a_1 , a_2 , a_3).
            \]

            令 $B = \beta, C = (a_1 , a_2 , a_3)$，即有 $A = BC$.
\end{solve}

\begin{exercise}\label{prob:L9-B-15}
设$\alpha,\beta$为$n$维列向量，且$A=\alpha\alpha^\mathrm{T}+\beta\beta^\mathrm{T}$.
        \begin{enumerate}
            \item 证明：$r(A) \leqslant 2$；

            \item 若$\alpha,\beta$线性相关，证明：$r(A) \leqslant 1$.
        \end{enumerate}
\end{exercise}
\noindent 来源：\href{https://github.com/yhwu-is/Linear-Algebra-Left-Undone/blob/new/%E8%AE%B2%E4%B9%89/%E4%B8%93%E9%A2%98/9%20%E7%9F%A9%E9%98%B5%E8%BF%90%E7%AE%97%E8%BF%9B%E9%98%B6.tex}{L9-B-15}。
\begin{solve}
\begin{enumerate}
                \item 设 $\alpha = (a_1, a_2, \ldots, a_n)^\mathrm{T}$，$\beta = (b_1, b_2, \ldots, b_n)^\mathrm{T}$，则
                    \begin{align*}
                        \alpha \alpha^\mathrm{T} &= (a_1 \alpha, a_2 \alpha, \ldots, a_n \alpha), \\
                        \beta \beta^\mathrm{T} &= (b_1 \beta, b_2 \beta, \ldots, b_n \beta).
                    \end{align*}
                    故 $\alpha \alpha^\mathrm{T}$ 和 $\beta \beta^\mathrm{T}$ 的任意两个列向量线性相关，$r\left(\alpha \alpha^\mathrm{T}\right) \leqslant 1$，$r\left(\beta \beta^\mathrm{T}\right) \leqslant 1$，故
                    \[
                        r\left(\alpha \alpha^\mathrm{T} + \beta \beta^\mathrm{T}\right) \leqslant r\left(\alpha \alpha^\mathrm{T}\right) + r\left(\beta \beta^\mathrm{T}\right) \leqslant 2.
                    \]

                \item 若 $\alpha = 0$，则 $A = \beta \beta^\mathrm{T}$，故 $r(A) \leqslant 1$.

                    若 $\alpha \neq 0$，由于 $\alpha, \beta$ 线性相关，存在常数 $k$，使得 $\beta = k \alpha$，则
                    \[
                        A = \alpha \alpha^\mathrm{T} + (k \alpha)(k \alpha)^\mathrm{T} = (k^2 + 1) \alpha \alpha^\mathrm{T},
                    \]
                    故 $r(A) \leqslant 1$.
            \end{enumerate}
\end{solve}

\begin{exercise}\label{prob:L9-B-16}
设$A$和$B$是两个$n$阶实方阵，证明：$r(A)=r(AB)$当且仅当存在$n$阶实方阵$C$使得$A=ABC$.
\end{exercise}
\noindent 来源：\href{https://github.com/yhwu-is/Linear-Algebra-Left-Undone/blob/new/%E8%AE%B2%E4%B9%89/%E4%B8%93%E9%A2%98/9%20%E7%9F%A9%E9%98%B5%E8%BF%90%E7%AE%97%E8%BF%9B%E9%98%B6.tex}{L9-B-16}。
\begin{solve}
\begin{enumerate}
                \item 存在 $n$ 阶实方阵 $C$ 使得 $A = ABC \implies r(A) = r(AB)$：
                    由秩不等式 $r(AB) \leqslant \min\{r(A), r(B)\}$，我们有
                    \[
                        r(ABC) \leqslant r(AB) \leqslant r(A).
                    \]

                    而由 $A=ABC$ 可知 $r(A) = r(ABC)$，故 $r(A) = r(AB)$.

                \item $r(A) = r(AB) \implies$ 存在 $n$ 阶实方阵 $C$ 使得 $A = ABC$：
                    设 $A = (\alpha_1, \alpha_2, \ldots, \alpha_n)$，$B = (b_{ij})_{n \times n}$，则
                    \begin{align*}
                        AB &= (\alpha_1, \alpha_2, \ldots, \alpha_n) \begin{pmatrix}
                                  b_{11} & b_{12} & \cdots & b_{1n} \\
                                  b_{21} & b_{22} & \cdots & b_{2n} \\
                                  \vdots & \vdots & \ddots & \vdots \\
                                  b_{n1} & b_{n2} & \cdots & b_{nn}
                              \end{pmatrix} \\
                           &= \left( \sum\limits_{i=1}^n b_{i1} \alpha_i, \sum\limits_{i=1}^n b_{i2} \alpha_i, \ldots, \sum\limits_{i=1}^n b_{in} \alpha_i \right).
                    \end{align*}
                    设 $\gamma_j = \sum\limits_{i=1}^n b_{ij} \alpha_i$. 由于 $r(A) = r(AB)$，$\alpha_1, \alpha_2, \ldots, \alpha_n$ 与 $\gamma_1, \gamma_2, \ldots, \gamma_n$ 等价，故存在 $C = (k_{ij})_{n \times n}$，使得
                    \begin{align*}
                        (\alpha_1, \alpha_2, \ldots, \alpha_n)
                        &= (\gamma_1, \gamma_2, \ldots, \gamma_n) \begin{pmatrix}
                            k_{11} & k_{12} & \cdots & k_{1n} \\
                            k_{21} & k_{22} & \cdots & k_{2n} \\
                            \vdots & \vdots & \ddots & \vdots \\
                            k_{n1} & k_{n2} & \cdots & k_{nn}
                        \end{pmatrix} \\
                        &= (\gamma_1, \gamma_2, \ldots, \gamma_n) C,
                    \end{align*}
                    即 $A = ABC$.
            \end{enumerate}
\end{solve}

\begin{exercise}\label{prob:L9-B-17}
证明如下结论：
        \begin{enumerate}
            \item 若存在正整数$m$使得$r(A^m)=r(A^{m+1})$，则必有$r(A^m)=r(A^{m+1})=r(A^{m+2})=\cdots$；
            \item 设$A$为$n$阶方阵，证明：$r(A^n)=r(A^{n+1})=r(A^{n+2})=\cdots$.
        \end{enumerate}
\end{exercise}
\noindent 来源：\href{https://github.com/yhwu-is/Linear-Algebra-Left-Undone/blob/new/%E8%AE%B2%E4%B9%89/%E4%B8%93%E9%A2%98/9%20%E7%9F%A9%E9%98%B5%E8%BF%90%E7%AE%97%E8%BF%9B%E9%98%B6.tex}{L9-B-17}。
\begin{solve}
\begin{enumerate}
                \item 该问题可以转换为线性映射的像空间停止收缩的问题：
                    设 $A = M(\sigma)$，则 $r(A^m) = r(A^{m+1})$ 等价于 $\dim \operatorname{im} \sigma^m = \dim \operatorname{im} \sigma^{m+1}$.

                    任取 $k \geqslant 2$，由于 $\forall \beta \in \operatorname{im} \sigma^k$，存在 $\alpha \in \mathbf{F}^n$，使得 $\beta = \sigma^k(\alpha) = \sigma^{k-1}(\sigma(\alpha)) \in \operatorname{im} \sigma^{k-1}$，我们有 $\operatorname{im} \sigma^k \subseteq \operatorname{im} \sigma^{k-1}$，即
                    \[
                        \mathbf{F}^n \supseteq \operatorname{im} \sigma \supseteq \operatorname{im} \sigma^2 \supseteq \cdots \supseteq \operatorname{im} \sigma^m \supseteq \operatorname{im} \sigma^{m+1} \supseteq \cdots,
                    \]

                    又 $\dim \operatorname{im} \sigma^m = \dim \operatorname{im} \sigma^{m+1}$，故
                    \[
                        \operatorname{im} \sigma^{m+1} = \operatorname{im} \sigma^m.
                    \]

                    假设 $n=k-1 \enspace (k > m)$ 时，有 $\operatorname{im} \sigma^n = \operatorname{im} \sigma^{n+1}$ 成立，下证 $n=k$ 时该结论也成立：

                    $\forall \gamma \in \operatorname{im} \sigma^k$，存在 $\alpha \in \mathbf{F}^n$，使得 $\gamma = \sigma^k(\alpha) = \sigma(\sigma^{k-1}(\alpha))$. 设 $\beta = \sigma^{k-1}(\alpha) \in \operatorname{im} \sigma^{k-1} = \operatorname{im} \sigma^k$，则 $\gamma = \sigma(\beta) \in \operatorname{im} \sigma^{k+1}$，故 $\operatorname{im} \sigma^k \subseteq \operatorname{im} \sigma^{k+1}$，而另一个方向的包含关系已经证明，所以有
                    \[
                        \operatorname{im} \sigma^{k+1} = \operatorname{im} \sigma^k.
                    \]

                    因此，
                    \[
                        \operatorname{im} \sigma^m = \operatorname{im} \sigma^{m+1} = \operatorname{im} \sigma^{m+2} = \cdots.
                    \]

                    显然有
                    \[
                        r(A^m) = r(A^{m+1}) = r(A^{m+2}) = \cdots.
                    \]

                \item 反证法. 假设 $r(A^n) \neq r(A^{n+1})$，则由 $\operatorname{im} \sigma^{n+1} \subseteq \operatorname{im} \sigma^n$，有 $r(A^n) > r(A^{n+1}) \geqslant 0$.
                    若存在 $k < n$ 使得 $r(A^k) = r(A^{k+1})$，则由第一问可知 $r(A^n) = r(A^{n+1})$，矛盾，故不存在这样的 $k$，即
                    \[
                        n > r(A) > r(A^2) > \cdots > r(A^{n-1}) > r(A^n) > r(A^{n+1}) \geqslant 0,
                    \]

                    由于秩只能是非负整数，上式不可能成立. 故有 $r(A^n) = r(A^{n+1})$. 再由第一问可得
                    \[
                        r(A^n) = r(A^{n+1}) = r(A^{n+2}) = \cdots.
                    \]

            \end{enumerate}
\end{solve}

\begin{exercise}\label{prob:L9-C-3}
已知数列$\{a_n\},\{b_n\}$满足$a_0=-1,\enspace b_0=3$，且
        \[\begin{cases}
                a_n=3a_{n-1}+b_{n-1}+2^{n-1} \\
                b_n=2a_{n-1}+4b_{n-1}+2^n
            \end{cases}\]
        求$\{a_n\},\{b_n\}$的通项公式.
\end{exercise}
\noindent 来源：\href{https://github.com/yhwu-is/Linear-Algebra-Left-Undone/blob/new/%E8%AE%B2%E4%B9%89/%E4%B8%93%E9%A2%98/9%20%E7%9F%A9%E9%98%B5%E8%BF%90%E7%AE%97%E8%BF%9B%E9%98%B6.tex}{L9-C-3}。
\begin{solve}
设 $v_n = (a_n, b_n, 2^n)^\mathrm{T}$，则递推式可化为
            \[
                v_n = \begin{pmatrix}
                    3 & 1 & 1 \\ 2 & 4 & 2 \\ 0 & 0 & 2
                \end{pmatrix} v_{n-1},
            \]
            其中 $v_0 = (-1, 3, 1)^\mathrm{T}$.

            设 $A = \begin{pmatrix}
                3 & 1 & 1 \\ 2 & 4 & 2 \\ 0 & 0 & 2
            \end{pmatrix}$，则
            \[
                v_n = A^n v_0.
            \]

            下面我们求 $A^n$. 注意到
            \[
                A = 2E + \begin{pmatrix}
                    1 & 1 & 1 \\ 2 & 2 & 2 \\ 0 & 0 & 0
                  \end{pmatrix}
                  = 2E + \begin{pmatrix}
                      1 \\ 2 \\ 0
                  \end{pmatrix} \begin{pmatrix}
                      1 & 1 & 1
                  \end{pmatrix}.
            \]

            令 $\alpha = (1, 2, 0)^\mathrm{T}, \beta = (1, 1, 1), B=\alpha \beta$，则
            \[
                B^n = \left( \alpha \beta \right)^n = \alpha \left( \beta \alpha \right)^{n-1} \beta = \operatorname{tr}(B)^{n-1} \alpha \beta = 3^{n-1} B.
            \]

            此时我们有
            \begin{align*}
                A^n &= (2E + B)^n \\
                    &= \tbinom{n}{0} 2^n E + \tbinom{n}{1} 2^{n-1} B + \tbinom{n}{2} 2^{n-2} B^2 + \cdots + \tbinom{n}{n} B^n \\
                    &= 2^n E + \left(\tbinom{n}{1} 2^{n-1} 3^0 + \tbinom{n}{2} 2^{n-2} 3^1 + \cdots + \tbinom{n}{n} 2^0 3^{n-1} \right) B \\
                    &= 2^n E - \frac{2^n}{3} B + \frac{1}{3} \left(\tbinom{n}{0} 2^n 3^0 + \tbinom{n}{1} 2^{n-1} 3^1 + \cdots + \tbinom{n}{n} 2^0 3^{n-1} \right) B \\
                    &= 2^n E + \frac{5^n - 2^n}{3} B \\
                    &= \frac{1}{3} \begin{pmatrix}
                        5^n + 2 \cdot 2^n & 5^n - 2^n         & 5^n - 2^n    \\
                        2(5^n - 2^n)      & 2 \cdot 5^n + 2^n & 2(5^n - 2^n) \\
                        0                 & 0                 & 3 \cdot 2^n
                    \end{pmatrix}.
            \end{align*}

            代入 $v_n = A^n v_0$，有
            \[
                v_n = \begin{pmatrix}
                    5^n - 2 \cdot 2^n \\
                    2 \cdot 5^n + 2^n \\
                    2^n
                \end{pmatrix}.
            \]

            所以 $a_n = 5^n - 2 \cdot 2^n, b_n = 2 \cdot 5^n + 2^n$.
\end{solve}

\begin{exercise}\label{prob:L9-C-4}
设$n\geqslant 3$，证明下列矩阵不可逆：
        \[A=\begin{pmatrix}
                \cos(\alpha_1-\beta_1) & \cos(\alpha_1-\beta_2) & \cdots & \cos(\alpha_1-\beta_n) \\
                \cos(\alpha_2-\beta_1) & \cos(\alpha_2-\beta_2) & \cdots & \cos(\alpha_2-\beta_n) \\
                \vdots                 & \vdots                 & \ddots & \vdots                 \\
                \cos(\alpha_n-\beta_1) & \cos(\alpha_n-\beta_2) & \cdots & \cos(\alpha_n-\beta_n)
            \end{pmatrix},\]
        \[B=\begin{pmatrix}
                1+x_1y_1 & 1+x_1y_2 & \cdots & 1+x_1y_n \\
                1+x_2y_1 & 1+x_2y_2 & \cdots & 1+x_2y_n \\
                \vdots   & \vdots   & \ddots & \vdots   \\
                1+x_ny_1 & 1+x_ny_2 & \cdots & 1+x_ny_n
            \end{pmatrix}.\]
\end{exercise}
\noindent 来源：\href{https://github.com/yhwu-is/Linear-Algebra-Left-Undone/blob/new/%E8%AE%B2%E4%B9%89/%E4%B8%93%E9%A2%98/9%20%E7%9F%A9%E9%98%B5%E8%BF%90%E7%AE%97%E8%BF%9B%E9%98%B6.tex}{L9-C-4}。
\begin{solve}
\begin{enumerate}
                \item 可以将 $A$ 拆成一个 $n \times 2$ 矩阵与一个 $2 \times n$ 矩阵的乘积：
                    \[
                        A = \begin{pmatrix}
                            \cos \alpha_1 & \sin \alpha_1 \\
                            \cos \alpha_2 & \sin \alpha_2 \\
                            \vdots        & \vdots        \\
                            \cos \alpha_n & \sin \alpha_n
                        \end{pmatrix} \begin{pmatrix}
                            \cos \beta_1 & \cos \beta_2 & \cdots & \cos \beta_n \\
                            \sin \beta_1 & \sin \beta_2 & \cdots & \sin \beta_n
                        \end{pmatrix}.
                    \]
                    故由秩不等式 $r(AB) \leqslant \min\{r(A), r(B)\}$，有 $r(A) \leqslant 2 < n$，故 $A$ 不可逆.

                \item 可以将 $A$ 拆成两个秩为 $1$ 的矩阵的和：
                    \[
                        A = \begin{pmatrix}
                            1      & 1      & \cdots & 1      \\
                            1      & 1      & \cdots & 1      \\
                            \vdots & \vdots & \ddots & \vdots \\
                            1      & 1      & \cdots & 1
                        \end{pmatrix} + \begin{pmatrix}
                            x_1 y_1 & x_1 y_2 & \cdots & x_1 y_n \\
                            x_2 y_1 & x_2 y_2 & \cdots & x_2 y_n \\
                            \vdots  & \vdots  & \ddots & \vdots  \\
                            x_n y_1 & x_n y_2 & \cdots & x_n y_n
                        \end{pmatrix},
                    \]
                    故由秩不等式 $r(A+B) \leqslant r(A)+r(B)$，有 $r(A) \leqslant 1 + 1 = 2 < n$，故 $A$ 不可逆.
            \end{enumerate}
\end{solve}

\begin{exercise}\label{prob:L9-C-5}
证明以下两个命题：
        \begin{enumerate}
            \item 与矩阵
\[I=\begin{pmatrix}
                          0 & 1 &   &        &   \\
                            &   & 1 &        &   \\
                            &   &   & \ddots &   \\
                            &   &   &        & 1 \\
                          1 &   &   &        & 0
                      \end{pmatrix}\]
可交换的矩阵$A$都可以写成$I$的一个多项式，即$A=a_{11}E+a_{12}I+a_{13}I^2+\cdots+a_{1n}I^{n-1}$；

            \item 与矩阵
\[J=\begin{pmatrix}
                          0 & 1 &   &        &   \\
                            &   & 1 &        &   \\
                            &   &   & \ddots &   \\
                            &   &   &        & 1 \\
                            &   &   &        & 0
                      \end{pmatrix}\]
可交换的矩阵$A$都可以写成$J$的一个多项式，即$A=a_{11}E+a_{12}J+a_{13}J^2+\cdots+a_{1n}J^{n-1}$.
        \end{enumerate}
\end{exercise}
\noindent 来源：\href{https://github.com/yhwu-is/Linear-Algebra-Left-Undone/blob/new/%E8%AE%B2%E4%B9%89/%E4%B8%93%E9%A2%98/9%20%E7%9F%A9%E9%98%B5%E8%BF%90%E7%AE%97%E8%BF%9B%E9%98%B6.tex}{L9-C-5}。
\begin{solve}
\begin{enumerate}
                \item 设 
\[B = \begin{pmatrix}
                        b_{11} & b_{12} & \cdots & b_{1n} \\
                        b_{21} & b_{22} & \cdots & b_{2n} \\
                        \vdots & \vdots & \ddots & \vdots \\
                        b_{n1} & b_{n2} & \cdots & b_{nn}
                    \end{pmatrix}\]
 与 $I$ 可交换，即 $BI=IB$，则
                    \[
                        \begin{pmatrix}
                            b_{1n} & b_{11} & b_{12} & \cdots & b_{1,n-1} \\
                            b_{2n} & b_{21} & b_{22} & \cdots & b_{2,n-1} \\
                            b_{3n} & b_{31} & b_{32} & \cdots & b_{3,n-1} \\
                            \vdots & \vdots & \vdots &        & \vdots    \\
                            b_{n-1,n} & b_{n-1,1} & b_{n-1,2} & \cdots & b_{n-1,n-1} \\
                            b_{nn} & b_{n1} & b_{n2} & \cdots & b_{n,n-1}
                        \end{pmatrix} = \begin{pmatrix}
                            b_{21} & b_{22} & b_{23} & \cdots & b_{2n} \\
                            b_{31} & b_{32} & b_{33} & \cdots & b_{3n} \\
                            b_{41} & b_{42} & b_{43} & \cdots & b_{4n} \\
                            \vdots & \vdots & \vdots &        & \vdots \\
                            b_{n1} & b_{n2} & b_{n3} & \cdots & b_{nn} \\
                            b_{11} & b_{12} & b_{13} & \cdots & b_{1n}
                        \end{pmatrix},
                    \]
                    逐一对照两矩阵中的对应元素，可得
                    \begin{align*}
                        b_{1n} &= b_{21} = b_{32} = \cdots = b_{n,n-1}, \\
                        b_{1,n-1} &= b_{2n} = b_{31} = \cdots = b_{n,n-2}, \\
                        &\vdotswithin{=} \\
                        b_{11} &= b_{22} = b_{33} = \cdots = b_{nn}.
                    \end{align*}
                    第 $k$ 行等式对应了 $I^{n-k}$，故与 $I$ 可交换的矩阵都可以写成
                    \[
                        a_{11}E + a_{12}I + a_{13}I^2 + \cdots + a_{1n}I^{n-1}.
                    \]

                \item 按照与第一问相同的方法进行操作，即可得到结果.
            \end{enumerate}
\end{solve}

\begin{exercise}\label{prob:L9-C-6}
证明：若$n$阶矩阵$A$的各阶左上角子块矩阵都可逆，则存在$n$阶下三角矩阵$B$，使得$BA$为上三角矩阵.
\end{exercise}
\noindent 来源：\href{https://github.com/yhwu-is/Linear-Algebra-Left-Undone/blob/new/%E8%AE%B2%E4%B9%89/%E4%B8%93%E9%A2%98/9%20%E7%9F%A9%E9%98%B5%E8%BF%90%E7%AE%97%E8%BF%9B%E9%98%B6.tex}{L9-C-6}。
\begin{solve}
下使用数学归纳法证明该命题. 对于 $n=1$ 的情况，取 $B=E_1$ 即可保证 $BA$ 为上三角矩阵. 现在假设命题对 $n-1$ 阶满足条件的矩阵均成立，下证命题对 $n$ 阶满足条件的矩阵也成立：

            设 $A_n = \begin{pmatrix}
                A_{n-1} & \beta \\
                \alpha  & a_{nn}
            \end{pmatrix}$，其中 $A_{n-1}$ 是 $n-1$ 阶方阵，满足各阶左上角子块矩阵都可逆. 由归纳假设可知存在 $n-1$ 阶下三角矩阵 $B_{n-1}$，使得 $C_{n-1} = B_{n-1}A_{n-1}$ 为上三角矩阵. 下面我们构造 $n$ 阶下三角矩阵 $B_n$：令
            \[
                B_n = \begin{pmatrix}
                    B_{n-1} & \mathbf{0} \\
                    \gamma  & 1
                \end{pmatrix},
            \]
            我们要使
            \[
                C_n = B_n A_n = \begin{pmatrix}
                    B_{n-1} & \mathbf{0} \\
                    \gamma  & 1
                \end{pmatrix} \begin{pmatrix}
                    A_{n-1} & \beta \\
                    \alpha  & a_{nn}
                \end{pmatrix} = \begin{pmatrix}
                    C_{n-1}                 & B_{n-1} \beta         \\
                    \gamma A_{n-1} + \alpha & \gamma \beta + a_{nn}
                \end{pmatrix}
            \]
            为上三角矩阵，只需令 $\gamma A_{n-1} + \alpha = 0$ 即可. 由于 $A_{n-1}$ 可逆，故对任意 $\alpha$ 都有对应的 $\gamma$ 存在，于是 $B_n$ 存在. 因此，命题对 $n$ 阶矩阵也成立.
\end{solve}

\begin{exercise}\label{prob:L9-C-7}
设$A$是数域$\mathbf{F}$上的$n$阶可逆矩阵，把$A$和$A^{-1}$做如下分块：
        \[A=\begin{pmatrix}
                A_{11} & A_{12} \\ A_{21} & A_{22}
            \end{pmatrix},\enspace A^{-1}=\begin{pmatrix}
                B_{11} & B_{12} \\ B_{21} & B_{22}
            \end{pmatrix}\]
        其中$A_{11}$是$l \times k$矩阵，$B_{11}$是$k \times l$矩阵，$l$，$k$是小于$n$的正整数. 用$W$表示$A_{12}X=0$的解空间，$U$表示$B_{12}Y=0$的解空间，其中$X$和$Y$为列向量，证明$W$与$U$同构.
\end{exercise}
\noindent 来源：\href{https://github.com/yhwu-is/Linear-Algebra-Left-Undone/blob/new/%E8%AE%B2%E4%B9%89/%E4%B8%93%E9%A2%98/9%20%E7%9F%A9%E9%98%B5%E8%BF%90%E7%AE%97%E8%BF%9B%E9%98%B6.tex}{L9-C-7}。
\begin{solve}
证明：$ \forall \alpha \in W,\enspace A_{12} \alpha = \vec{0} $.
            \begin{align*}
                \begin{pmatrix} O_{k \times 1} \\ \alpha \end{pmatrix}
                & = A^{-1}A \begin{pmatrix} O_{k \times 1} \\ \alpha \end{pmatrix} = A^{-1} \begin{pmatrix} A_{11} & A_{12} \\ A_{21} & A_{22} \end{pmatrix} \begin{pmatrix} O_{k \times 1} \\ \alpha \end{pmatrix} = A^{-1} \begin{pmatrix} O_{l \times 1} \\ A_{22} \alpha \end{pmatrix} \\
                & = \begin{pmatrix} B_{11} & B_{12} \\ B_{21} & B_{22} \end{pmatrix} \begin{pmatrix} O_{l \times 1} \\ A_{22} \alpha \end{pmatrix} = \begin{pmatrix} B_{12} A_{22} \alpha \\ B_{22} A_{22} \alpha \end{pmatrix}.
            \end{align*}
            故 $ B_{12} A_{22} \alpha = \vec{0} $，故我们可以推测如下定义：$ \sigma \in \mathcal{L}(W,U),\enspace \sigma(\alpha) = A_{22} \alpha $.
            只需证明 $ \sigma $ 是单射且满射即可.

            单射：$ \sigma(\alpha) = \sigma(\beta) \implies A_{22}(\alpha - \beta) = \vec{0} $. 又有 $ A_{12} \alpha = A_{12} \beta = \vec{0} \implies A_{12}( \alpha - \beta) = \vec{0} $. 故 $ \begin{pmatrix} A_{12} \\ A_{22} \end{pmatrix} (\alpha - \beta) = \vec{0} $. 由于 $ \begin{pmatrix} A_{12} \\ A_{22} \end{pmatrix} $ 列满秩（$ A $ 可逆），故 $ \alpha = \beta $.

            满射：$ \forall \gamma \in U,\enspace B_{12} \gamma = \vec{0} $.
            \begin{align*}
                \begin{pmatrix} O_{l \times 1} \\ \gamma \end{pmatrix}
                & = A A^{-1} \begin{pmatrix} O_{l \times 1} \\ \gamma \end{pmatrix} = A \begin{pmatrix} B_{11} & B_{12} \\ B_{21} & B_{22} \end{pmatrix} \begin{pmatrix} O_{l \times 1} \\ \gamma \end{pmatrix}      \\
                & = A \begin{pmatrix} O_{k \times 1} \\ B_{22} \gamma \end{pmatrix} = \begin{pmatrix} A_{11} & A_{12} \\ A_{21} & A_{22} \end{pmatrix} \begin{pmatrix} O_{k \times 1} \\ B_{22} \gamma \end{pmatrix} \\
                & = \begin{pmatrix} A_{12} B_{22} \gamma \\ A_{22} B_{22} \gamma \end{pmatrix}.
            \end{align*}
            故 $ \exists B_{22} \gamma \in W,\enspace A_{22} B_{22} \gamma = \gamma \in U $.
\end{solve}

\begin{exercise}\label{prob:L9-C-8}
证明 Schur 补的商等式 $A/B=(A/C)/(B/C)$，其中 $A,B,C$ 均为可逆矩阵，且 $B$ 为 $A$ 的主子矩阵，$C$ 为 $B$ 的主子矩阵.
\end{exercise}
\noindent 来源：\href{https://github.com/yhwu-is/Linear-Algebra-Left-Undone/blob/new/%E8%AE%B2%E4%B9%89/%E4%B8%93%E9%A2%98/9%20%E7%9F%A9%E9%98%B5%E8%BF%90%E7%AE%97%E8%BF%9B%E9%98%B6.tex}{L9-C-8}。
\begin{solve}
设 $A = \begin{pmatrix}
                    A_{11} & A_{12} & A_{13} \\
                    A_{21} & A_{22} & A_{23} \\
                    A_{31} & A_{32} & A_{33}
                \end{pmatrix}$，$B = \begin{pmatrix}
                    A_{11} & A_{12} \\
                    A_{21} & A_{22}
                \end{pmatrix}$，$C = A_{11}$，则根据 Schur 补的定义，有
                \begin{align*}
                    A/C &= \begin{pmatrix}
                        A_{22} & A_{23} \\
                        A_{32} & A_{33}
                    \end{pmatrix} - \begin{pmatrix}
                        A_{21} \\ A_{31}
                    \end{pmatrix} A_{11}^{-1} \begin{pmatrix}
                        A_{12} & A_{13}
                    \end{pmatrix} \\
                    &= \begin{pmatrix}
                        A_{22} - A_{21} A_{11}^{-1} A_{12} & A_{23} - A_{21} A_{11}^{-1} A_{13} \\
                        A_{32} - A_{31} A_{11}^{-1} A_{12} & A_{33} - A_{31} A_{11}^{-1} A_{13}
                    \end{pmatrix},
                \end{align*}
                \[
                    A/B = A_{33} - \begin{pmatrix}
                        A_{31} & A_{32}
                    \end{pmatrix} \begin{pmatrix}
                        A_{11} & A_{12} \\
                        A_{21} & A_{22}
                    \end{pmatrix}^{-1} \begin{pmatrix}
                        A_{13} \\ A_{23}
                    \end{pmatrix}, \enspace
                    B/C = A_{22} - A_{21} A_{11}^{-1} A_{12}.
                \]

                设 $P_{ij} = A_{ij} - A_{i1} A_{11}^{-1} A_{1j} \enspace (i=2,3; j=2,3)$，则
                \[
                    A/C = \begin{pmatrix}
                        P_{22} & P_{23} \\
                        P_{32} & P_{33}
                    \end{pmatrix}, \enspace
                    B/C = P_{22}.
                \]
                故
                \[
                    (A/C)/(B/C) = P_{33} - P_{32} P_{22}^{-1} P_{23}.
                \]

                由分块矩阵初等变换可以求得
                \[
                    \begin{pmatrix}
                        A_{11} & A_{12} \\
                        A_{21} & A_{22}
                    \end{pmatrix}^{-1} = \begin{pmatrix}
                        A_{11}^{-1} + A_{11}^{-1} A_{12} P_{22}^{-1} A_{21} A_{11}^{-1} & -A_{11}^{-1} A_{12} P_{22}^{-1} \\
                        -P_{22}^{-1} A_{21} A_{11}^{-1} & P_{22}^{-1}
                    \end{pmatrix},
                \]
                由此可得
                \begin{align*}
                    A/B = {} & A_{33} - \begin{pmatrix}
                        A_{31} & A_{32}
                    \end{pmatrix} \begin{pmatrix}
                        A_{11}^{-1} + A_{11}^{-1} A_{12} P_{22}^{-1} A_{21} A_{11}^{-1} & -A_{11}^{-1} A_{12} P_{22}^{-1} \\
                        -P_{22}^{-1} A_{21} A_{11}^{-1} & P_{22}^{-1}
                    \end{pmatrix} \begin{pmatrix}
                        A_{13} \\ A_{23}
                    \end{pmatrix} \\
                    = {} & A_{33} - A_{31} \left(A_{11}^{-1} + A_{11}^{-1} A_{12} P_{22}^{-1} A_{21} A_{11}^{-1}\right) - A_{31} \left(-A_{11}^{-1} A_{12} P_{22}^{-1}\right) - \\
                      {} & A_{32} \left(-P_{22}^{-1} A_{21} A_{11}^{-1}\right) - A_{32} P_{22}^{-1} A_{23} \\
                    = {} & \left(A_{33} - A_{31} A_{11}^{-1} A_{13}\right) - \left(A_{32} - A_{31} A_{11}^{-1} A_{12}\right) P_{22}^{-1} \left(A_{23} - A_{21} A_{11}^{-1} A_{13}\right) \\
                    = {} & P_{33} - P_{32} P_{22}^{-1} P_{23} \\
                    = {} & (A/C)/(B/C).
                \end{align*}
\end{solve}

\begin{exercise}\label{prob:L9-C-10}
已知$A$是一个$s \times n$矩阵，证明：$r(E_n-A^\mathrm{T}A)-r(E_s-AA^\mathrm{T})=n-s$.
\end{exercise}
\noindent 来源：\href{https://github.com/yhwu-is/Linear-Algebra-Left-Undone/blob/new/%E8%AE%B2%E4%B9%89/%E4%B8%93%E9%A2%98/9%20%E7%9F%A9%E9%98%B5%E8%BF%90%E7%AE%97%E8%BF%9B%E9%98%B6.tex}{L9-C-10}。
\begin{solve}
令
            \[
                B = \begin{pmatrix}
                    E_n & A^\mathrm{T} \\
                    A   & E_s
                \end{pmatrix},
            \]
            对 $B$ 做分块倍加变化，有
            \begin{gather*}
                \begin{pmatrix}
                    E_n & A^\mathrm{T} \\
                    A   & E_s
                \end{pmatrix} \to \begin{pmatrix}
                    E_n - A^\mathrm{T}A & A^\mathrm{T} \\
                    O                   & E_s
                \end{pmatrix} \to \begin{pmatrix}
                    E_n - A^\mathrm{T}A & O   \\
                    O                   & E_s
                \end{pmatrix}, \\
                \begin{pmatrix}
                    E_n & A^\mathrm{T} \\
                    A   & E_s
                \end{pmatrix} \to \begin{pmatrix}
                    E_n & A^\mathrm{T}        \\
                    O   & E_s - AA^\mathrm{T}
                \end{pmatrix} \to \begin{pmatrix}
                    E_n & O                   \\
                    O   & E_s - AA^\mathrm{T}
                \end{pmatrix}.
            \end{gather*}
            分块倍加变换不改变矩阵的秩，故有
            \begin{align*}
                r(B) &= r\begin{pmatrix}
                    E_n - A^\mathrm{T}A & O   \\
                    O                   & E_s
                \end{pmatrix} = r(E_n - A^\mathrm{T}A) + s, \\
                r(B) &= r\begin{pmatrix}
                    E_n & O                   \\
                    O   & E_s - AA^\mathrm{T}
                \end{pmatrix} = r(E_s - AA^\mathrm{T}) + n.
            \end{align*}
            于是有 $r(E_n-A^\mathrm{T}A)-r(E_s-AA^\mathrm{T})=n-s$.
\end{solve}

\begin{exercise}\label{prob:L9-C-12}
$f(x)=f_1(x)f_2(x)$是多项式，且$f_1(x)$与$f_2(x)$互素，则$f(A)=O$的充要条件是$r(f_1(A))+r(f_2(A))=n$. （注：此题的推论非常多，如$A^2=A$，$A^n=E$等形式的结论都可以利用这个例子推导出）
\end{exercise}
\noindent 来源：\href{https://github.com/yhwu-is/Linear-Algebra-Left-Undone/blob/new/%E8%AE%B2%E4%B9%89/%E4%B8%93%E9%A2%98/9%20%E7%9F%A9%E9%98%B5%E8%BF%90%E7%AE%97%E8%BF%9B%E9%98%B6.tex}{L9-C-12}。
\begin{solve}
略有超纲，使用贝祖定理. $\exists u(x),v(x),u(x)f_1(x)+v(x)f_2(x)=1 $，于是
            \begin{align*}
                r\begin{pmatrix}
                    f_1(A) & O      \\
                    O      & f_2(A)
                \end{pmatrix}
                &= r\begin{pmatrix}
                        f_1(A) & f_1(A)u(A)+f_2(A)v(A) \\
                        O      & f_2(A)
                    \end{pmatrix} = r\begin{pmatrix}
                        f_1(A) & E_n    \\
                        O      & f_2(A)
                    \end{pmatrix} \\
                &= r\begin{pmatrix}
                        f_1(A)        & E_n \\
                        -f_2(A)f_1(A) & O
                    \end{pmatrix} = r\begin{pmatrix}
                        O    & E_n \\
                        f(A) & O
                    \end{pmatrix}.
            \end{align*}
            因此
            \[
                f(A) = O \iff
                r\begin{pmatrix}
                    f_1(A) & O      \\
                    O      & f_2(A)
                \end{pmatrix} = r\begin{pmatrix}
                    O    & E_n \\
                    f(A) & O
                \end{pmatrix} = n \iff
                r(f_1(A)) + r(f_2(A)) = n.
            \]
\end{solve}

\begin{exercise}\label{prob:L9-C-13}
设$A,B$分别为$3 \times 2$和$2 \times 3$实矩阵. 若$AB=\begin{pmatrix}
                8             & 0 & -4 \\[1ex]
                -\dfrac{3}{2} & 9 & -6 \\[1ex]
                -2            & 0 & 1
            \end{pmatrix}$，求$BA$.
\end{exercise}
\noindent 来源：\href{https://github.com/yhwu-is/Linear-Algebra-Left-Undone/blob/new/%E8%AE%B2%E4%B9%89/%E4%B8%93%E9%A2%98/9%20%E7%9F%A9%E9%98%B5%E8%BF%90%E7%AE%97%E8%BF%9B%E9%98%B6.tex}{L9-C-13}。
\begin{solve}
由于 $A$ 是列满秩矩阵，$B$ 是行满秩矩阵，知存在可逆矩阵 $P_{3\times 3},Q_{2\times 2}$ 使得
            \[A=P\begin{pmatrix}E_2 \\ O\end{pmatrix},B=\begin{pmatrix}E_2 & O\end{pmatrix}Q,\]
            于是
            \[BA=\begin{pmatrix}E_2 & O\end{pmatrix}QP\begin{pmatrix}E_2 \\ O\end{pmatrix}.\]
            由 $(AB)^2=9AB$ 有 \[P\begin{pmatrix}E_2 \\ O\end{pmatrix}\begin{pmatrix}E_2 & O\end{pmatrix}QP\begin{pmatrix}E_2 \\ O\end{pmatrix}\begin{pmatrix}E_2 & O\end{pmatrix}Q=9P\begin{pmatrix}E_2 \\ O\end{pmatrix}\begin{pmatrix}E_2 & O\end{pmatrix}Q,\]
            即
            \[\begin{pmatrix}E_2 \\ O\end{pmatrix}BA\begin{pmatrix}E_2 & O\end{pmatrix}=9\begin{pmatrix}E_2 \\ O\end{pmatrix}\begin{pmatrix}E_2 & O\end{pmatrix},\]
            也就是 \[\begin{pmatrix}BA & O \\ O & O\end{pmatrix}=\begin{pmatrix}9E_2 & 0 \\ 0 & 0\end{pmatrix}.\]
            所以 $BA=9E_2$.
\end{solve}

\begin{exercise}\label{prob:L9-例-2}
求所有与$A$可交换的矩阵，其中
    \[A=\begin{pmatrix}
            1 & 0 & 0  \\
            0 & 1 & 2  \\
            0 & 1 & -2
        \end{pmatrix}.\]
\end{exercise}
\noindent 来源：\href{https://github.com/yhwu-is/Linear-Algebra-Left-Undone/blob/new/%E8%AE%B2%E4%B9%89/%E4%B8%93%E9%A2%98/9%20%E7%9F%A9%E9%98%B5%E8%BF%90%E7%AE%97%E8%BF%9B%E9%98%B6.tex}{L9-例-2}。
\begin{solve}
设$B$为与$A$可交换的矩阵，设
    \[B=\begin{pmatrix}
            a & b & c \\
            d & e & f \\
            g & h & i
        \end{pmatrix},\]
    则
    \[AB=\begin{pmatrix}
            a    & b    & c    \\
            d+2g & e+2h & f+2i \\
            d-2g & e-2h & f-2i
        \end{pmatrix}=BA=\begin{pmatrix}
            a & b+c & 2b-2c \\
            d & e+f & 2e-2f \\
            g & h+i & 2h-2i
        \end{pmatrix}.\]
    由此可得
    \[\begin{cases}
            b=b+c      \\
            c=2b-2c    \\
            d+2g=d     \\
            e+2h=e+f   \\
            f+2i=2e-2f \\
            d-2g=g     \\
            e-2h=h+i   \\
            f-2i=2h-2i
        \end{cases},\]
    很容易解得$b=c=d=g=0$，且$f=2h,\enspace e=3h+i$，因此
    \[B=\begin{pmatrix}
            a & 0    & 0  \\
            0 & 3h+i & 2h \\
            0 & h    & i
        \end{pmatrix}.\]
\end{solve}

\begin{exercise}\label{prob:L9-例-3}
求与矩阵$A=\begin{pmatrix}
            3  & 0  & 0 \\
            -1 & 3  & 0 \\
            0  & -1 & 3
        \end{pmatrix}$可交换的矩阵.
\end{exercise}
\noindent 来源：\href{https://github.com/yhwu-is/Linear-Algebra-Left-Undone/blob/new/%E8%AE%B2%E4%B9%89/%E4%B8%93%E9%A2%98/9%20%E7%9F%A9%E9%98%B5%E8%BF%90%E7%AE%97%E8%BF%9B%E9%98%B6.tex}{L9-例-3}。
\begin{solve}
由前述分析，取$t=3$可以使得$A-3E$对角线上元素全为0，便于计算，因此我们只需求与$A-3E$可交换的矩阵即可. 与原讲义“可交换矩阵1”处 同样的设未知数法，具体过程省略得到与$A$可交换的矩阵为
    \[B=\begin{pmatrix}
            a & 0 & 0 \\
            b & a & 0 \\
            c & b & a
        \end{pmatrix}.\]
\end{solve}

\begin{exercise}\label{prob:L9-例-10}
设 $A$ 为 $n$ 阶可逆矩阵，考虑 $B = A + XRY$，其中 $X$ 是 $n \times r$ 矩阵，$Y$ 是 $r \times n$ 矩阵，而 $R$ 是 $r \times r$ 的可逆矩阵. 若 $R^{-1} + YA^{-1}X$ 可逆，则
    \[B^{-1}=A^{-1}-A^{-1}X(R^{-1}+YA^{-1}X)^{-1}YA^{-1}.\]
\end{exercise}
\noindent 来源：\href{https://github.com/yhwu-is/Linear-Algebra-Left-Undone/blob/new/%E8%AE%B2%E4%B9%89/%E4%B8%93%E9%A2%98/9%20%E7%9F%A9%E9%98%B5%E8%BF%90%E7%AE%97%E8%BF%9B%E9%98%B6.tex}{L9-例-10}。
\begin{solve}
注意到 $B = A + XRY = A(E + A^{-1}XRY)$，将 $A^{-1}X$ 和 $RY$ 分别视为整体，因为 $E + RYA^{-1}X = R(R^{-1} + YA^{-1}X)$ 可逆，根据原讲义“凑因子3”处，则有 $E + A^{-1}XRY$ 可逆，并且
    \begin{align*}
        (E + A^{-1}XRY)^{-1} & = E - A^{-1}X(E + RYA^{-1}X)^{-1}RY \\
                             & = E - A^{-1}X(R^{-1} + YA^{-1}X)^{-1}Y.
    \end{align*}
    因此 $B^{-1} = (E + A^{-1}XRY)^{-1}A^{-1} = A^{-1} - A^{-1}X(R^{-1} + YA^{-1}X)^{-1}YA^{-1}$.
\end{solve}

\begin{exercise}\label{prob:L9-例-11}
求 $n$ 阶矩阵 
\[A = \begin{pmatrix}
        0 & 1 & 1 & \cdots & 1 \\
        1 & 0 & 1 & \cdots & 1 \\
        1 & 1 & 0 & \cdots & 1 \\
        \vdots & \vdots & \vdots & \ddots & \vdots \\
        1 & 1 & 1 & \cdots & 0
    \end{pmatrix}\]
 的逆矩阵.
\end{exercise}
\noindent 来源：\href{https://github.com/yhwu-is/Linear-Algebra-Left-Undone/blob/new/%E8%AE%B2%E4%B9%89/%E4%B8%93%E9%A2%98/9%20%E7%9F%A9%E9%98%B5%E8%BF%90%E7%AE%97%E8%BF%9B%E9%98%B6.tex}{L9-例-11}。
\begin{solve}
不难观察到这一矩阵与秩为 $1$ 的元素全为 $1$ 的矩阵之间只差一个单位矩阵，因此令 $\alpha = (1,1,\ldots,1)^\mathrm{T}$，则 $A = -E_n + \alpha\alpha^\mathrm{T}$，由 Sherman-Morrison 公式，我们有
    \begin{align*}
        A^{-1} &= -E_n - (1 + \alpha^\mathrm{T}(-E_n)\alpha)^{-1}(-E_n)\alpha\alpha^\mathrm{T}(-E_n) \\
                &= -E_n + \dfrac{1}{n-1}\alpha\alpha^\mathrm{T}.
    \end{align*}
\end{solve}

\begin{exercise}\label{prob:L9-例-12}
已知方阵$A$满足$A^k=O$，其中$k$是一个正整数，求$E-A$的逆.
\end{exercise}
\noindent 来源：\href{https://github.com/yhwu-is/Linear-Algebra-Left-Undone/blob/new/%E8%AE%B2%E4%B9%89/%E4%B8%93%E9%A2%98/9%20%E7%9F%A9%E9%98%B5%E8%BF%90%E7%AE%97%E8%BF%9B%E9%98%B6.tex}{L9-例-12}。
\begin{solve}
根据我们前面的分析，结合$A^k=O$，我们猜测
    \[(E-A)^{-1}=E+A+A^2+\cdots+A^{k-1}.\]
    事实上我们直接验证
    \[(E-A)(E+A+A^2+\cdots+A^{k-1})=E-A^k=E.\]
    因此$(E-A)^{-1}=E+A+A^2+\cdots+A^{k-1}$.
\end{solve}

\begin{exercise}\label{prob:L9-例-16}
设$A=\begin{pmatrix}0 & -1 & 0 \\ 1 & 0 & 0 \\ 0 & 0 & -1 \end{pmatrix},\enspace P^{-1}AP=B$，求$B^{2004}-2A^2$.
\end{exercise}
\noindent 来源：\href{https://github.com/yhwu-is/Linear-Algebra-Left-Undone/blob/new/%E8%AE%B2%E4%B9%89/%E4%B8%93%E9%A2%98/9%20%E7%9F%A9%E9%98%B5%E8%BF%90%E7%AE%97%E8%BF%9B%E9%98%B6.tex}{L9-例-16}。
\begin{solve}
事实上，$A^2=\begin{pmatrix}
                      -1 & 0 & 0 \\ 0 & -1 & 0 \\ 0 & 0 & 1
                  \end{pmatrix}$，因此$A^4=E$，于是$A^{2004}=(A^4)^{501}=E$，因此$B^{2004}-2A^2=E-2A^2=\begin{pmatrix}
                      3 & 0 & 0 \\ 0 & 3  & 0 \\ 0 & 0 & -1
                  \end{pmatrix}$.
\end{solve}

\begin{exercise}\label{prob:L9-例-21}
已知$M$是秩为 1 的矩阵，记$\operatorname{tr}(M)=b$，讨论$(aE+M)^n$的计算结果.
\end{exercise}
\noindent 来源：\href{https://github.com/yhwu-is/Linear-Algebra-Left-Undone/blob/new/%E8%AE%B2%E4%B9%89/%E4%B8%93%E9%A2%98/9%20%E7%9F%A9%E9%98%B5%E8%BF%90%E7%AE%97%E8%BF%9B%E9%98%B6.tex}{L9-例-21}。
\begin{solve}
事实上，由原讲义“相抵分解”处 可知，$M^k=b^{k-1}M$，因此
              \begin{enumerate}
                  \item 当$b=0$时，$M^k=O(k\geqslant 2)$，因此
                        \[(aE+M)^n=a^nE+na^{n-1}M.\]

                  \item 当$b\neq 0$时有
                        \begin{align*}
                            (aE+M)^n & =\sum\limits_{k=0}^nC_n^ka^{n-k}M^k                                    \\
                                     & =\sum\limits_{k=0}^nC_n^ka^{n-k}b^{k-1}M                               \\
                                     & =a^nE+\dfrac{1}{b}\sum\limits_{k=1}^nC_n^ka^{n-k}b^kM                  \\
                                     & =a^nE+\dfrac{1}{b}\left(\sum\limits_{k=0}^nC_n^ka^{n-k}b^k-a^n\right)M \\
                                     & =a^nE+\dfrac{(a+b)^n-a^n}{b}M
                        \end{align*}
              \end{enumerate}
\end{solve}

\begin{exercise}\label{prob:L9-例-22}
求矩阵的幂
\[A^n=\begin{pmatrix}
                      1  & -1 & -1 & -1 \\
                      -1 & 1  & -1 & -1 \\
                      -1 & -1 & 1  & -1 \\
                      -1 & -1 & -1 & 1
                  \end{pmatrix}^n\]
.
\end{exercise}
\noindent 来源：\href{https://github.com/yhwu-is/Linear-Algebra-Left-Undone/blob/new/%E8%AE%B2%E4%B9%89/%E4%B8%93%E9%A2%98/9%20%E7%9F%A9%E9%98%B5%E8%BF%90%E7%AE%97%E8%BF%9B%E9%98%B6.tex}{L9-例-22}。
\begin{solve}
很显然我们可以将$A$分为一个数量矩阵秩1矩阵的和，即\[A=2E+\begin{pmatrix}
                      -1 & -1 & -1 & -1 \\
                      -1 & -1 & -1 & -1 \\
                      -1 & -1 & -1 & -1 \\
                      -1 & -1 & -1 & -1
                  \end{pmatrix},\]
              然后我们利用上例的结论即可（或者不记得结论也可以将
\[\begin{pmatrix}
                      -1 & -1 & -1 & -1 \\
                      -1 & -1 & -1 & -1 \\
                      -1 & -1 & -1 & -1 \\
                      -1 & -1 & -1 & -1
                  \end{pmatrix}\]
分解成$\begin{pmatrix}
                      -1 \\ -1 \\ \vdots \\ -1
                  \end{pmatrix}\begin{pmatrix}
                      1 & 1 & \cdots & 1
                  \end{pmatrix}$然后现场推导）.
\end{solve}

\begin{exercise}\label{prob:L9-例-24}
设$A$为三阶矩阵，$P$为三阶可逆矩阵，$P^{-1}AP=B$，其中
              \[P=\begin{pmatrix}
                      0 & 2 & -1 \\ 1 & 1 & 2 \\ -1 & -1 & -1
                  \end{pmatrix},\enspace B=\begin{pmatrix}
                      0 & 0 & -1 \\ 0 & -1 & 0 \\ -1 & 0 & 0
                  \end{pmatrix},\]
              求$A^{2024}$.
\end{exercise}
\noindent 来源：\href{https://github.com/yhwu-is/Linear-Algebra-Left-Undone/blob/new/%E8%AE%B2%E4%B9%89/%E4%B8%93%E9%A2%98/9%20%E7%9F%A9%E9%98%B5%E8%BF%90%E7%AE%97%E8%BF%9B%E9%98%B6.tex}{L9-例-24}。
\begin{solve}
事实上$A=PBP^{-1}$，因此$A^{2024}=PB^{2024}P^{-1}$，由于$B^2=E$，因此$B^{2024}=(B^2)^{1012}=E$，因此$A^{2024}=PEP^{-1}=E$.
\end{solve}

\begin{exercise}\label{prob:L9-例-25}
当$D$和$A-BD^{-1}C$可逆时，通过分块矩阵初等变换求$P=\begin{pmatrix}A & B \\ C & D\end{pmatrix}$的逆矩阵.
\end{exercise}
\noindent 来源：\href{https://github.com/yhwu-is/Linear-Algebra-Left-Undone/blob/new/%E8%AE%B2%E4%B9%89/%E4%B8%93%E9%A2%98/9%20%E7%9F%A9%E9%98%B5%E8%BF%90%E7%AE%97%E8%BF%9B%E9%98%B6.tex}{L9-例-25}。
\begin{solve}
设$A$和$D$分别是$m$、$n$阶矩阵，根据前面的讨论，我们的想法就是对$P$和$E$同时做初等变换使得$P$变为单位矩阵，那么我们自然会想到打洞法，利用$D$可逆的条件消去左下角和右上角的矩阵. 我们首先对第二个分块行左乘$-BD^{-1}$加到第一个分块行，有
    \[\left(\begin{array}{cc:cc}
                A & B & E_m & O \\ C & D & O & E_n
            \end{array}\right)\rightarrow\left(\begin{array}{cc:cc}
                A-BD^{-1}C & O & E_m & -BD^{-1} \\ C & D & O & E_n
            \end{array}\right),\]
    这样消去了右上角的分块阵，接下来消去左下角的（当然也可以换个顺序消去，效果一致），我们此时利用$A-BD^{-1}C$可逆（继续利用$D$可逆得到的结果不具代表性，故不展开），将第一个分块行左乘$-C(A-BD^{-1}C)^{-1}$加到第二个分块行，得到
    \[\left(\begin{array}{cc:cc}
                A-BD^{-1}C & O & E_m & -BD^{-1} \\ O & D & -C(A-BD^{-1}C)^{-1} & E_n+C(A-BD^{-1}C)^{-1}BD^{-1}
            \end{array}\right),\]
    最后用$(A-BD^{-1}C)^{-1}$和$D^{-1}$分别左乘第一分块行和第二分块行，得到
    \[\left(\begin{array}{cc:cc}
                E_m & O & (A-BD^{-1}C)^{-1} & -(A-BD^{-1}C)^{-1}BD^{-1} \\ O & E_n & -D^{-1}C(A-BD^{-1}C)^{-1} & D^{-1}+D^{-1}C(A-BD^{-1}C)^{-1}BD^{-1}
            \end{array}\right),\]
    由此可得原矩阵的逆就是上述虚线右侧的
    \[\begin{pmatrix}
            (A-BD^{-1}C)^{-1} & -(A-BD^{-1}C)^{-1}BD^{-1} \\ -D^{-1}C(A-BD^{-1}C)^{-1} & D^{-1}+D^{-1}C(A-BD^{-1}C)^{-1}BD^{-1}
        \end{pmatrix}.\]
\end{solve}

\begin{exercise}\label{prob:L9-例-28}
若$n$阶矩阵$A$的各阶左上角子块矩阵都可逆，则存在上三角矩阵$U$和主对角元全为1的下三角矩阵$L$，使得$A=LU$（$LU$分解）.
\end{exercise}
\noindent 来源：\href{https://github.com/yhwu-is/Linear-Algebra-Left-Undone/blob/new/%E8%AE%B2%E4%B9%89/%E4%B8%93%E9%A2%98/9%20%E7%9F%A9%E9%98%B5%E8%BF%90%E7%AE%97%E8%BF%9B%E9%98%B6.tex}{L9-例-28}。
\begin{solve}
由于主对角元全为1的下三角矩阵可逆，其逆矩阵也是主对角元全为1的下三角矩阵，因此只要证明存在主对角元全为1的下三角矩阵$S$使得$SA=U$，我们可以利用数学归纳法来证明这一结论.

    当$n=1$时，结论显然成立. 假设命题对$n-1$阶矩阵成立. 对$n$阶矩阵$A$，将$A$分块为
    \[A=\begin{pmatrix}
            A_1 & \alpha_1 \\ \alpha_2 & a_{nn}
        \end{pmatrix},\]
    其中$A_1$为满足命题条件的$n-1$阶矩阵，根据归纳假设，存在上三角矩阵$U_1$和$n-1$阶主对角元全为1的下三角矩阵$S_1$使得$S_1A_1=U_1$. 根据我们前面的讨论，我们可以对$A$做分块矩阵初等变换将其化为上三角块矩阵，即
    \[PA=\begin{pmatrix}
            E_{n-1} & O \\ -\alpha_2A_1^{-1} & 1
        \end{pmatrix}\begin{pmatrix}
            A_1 & \alpha_1 \\ \alpha_2 & a_{nn}
        \end{pmatrix}=\begin{pmatrix}
            A_1 & \alpha_1 \\ O & a_{nn}-\alpha_2A_1^{-1}\alpha_1
        \end{pmatrix},\]
    再取$Q=\begin{pmatrix}
            S_1 & O \\ O & 1
        \end{pmatrix}$，即得
    \[QPA=\begin{pmatrix}
            U_1 & S_1\alpha_1 \\ O & a_{nn}-\alpha_2A_1^{-1}\alpha_1
        \end{pmatrix}=U.\]
    其中$U$为上三角矩阵，$S=QP$是主对角元全为1的下三角矩阵，故存在$L=S^{-1}$和$U$使得$A=LU$. 从而由归纳法得证.
\end{solve}

\begin{exercise}\label{prob:L9-例-29}
设矩阵$A=\begin{pmatrix}
            1 & 2 & -1 \\ 2 & 1 & -2 \\ -3 & 1 & 1
        \end{pmatrix}$，求$A$的$LU$分解，并利用$LU$分解求解方程组$Ax=b$，其中$b=(3,3,-6)^\mathrm{T}$.
\end{exercise}
\noindent 来源：\href{https://github.com/yhwu-is/Linear-Algebra-Left-Undone/blob/new/%E8%AE%B2%E4%B9%89/%E4%B8%93%E9%A2%98/9%20%E7%9F%A9%E9%98%B5%E8%BF%90%E7%AE%97%E8%BF%9B%E9%98%B6.tex}{L9-例-29}。
\begin{solve}
对于$A$，我们可以通过下述三个初等行变换化为上三角矩阵：
    \begin{enumerate}
        \item 第一行乘以$-2$加到第二行；
        \item 第一行乘以$3$加到第三行；
        \item 第二行乘以$\dfrac{7}{3}$加到第三行.
    \end{enumerate}
    记三个初等变换对应的矩阵依次为$P_1,P_2,P_3$，则上述初等变换过程可以写为
    \begin{equation}\label{eq:10:LU分解}
        P_3P_2P_1A=\begin{pmatrix}
            1 & 2 & -1 \\ 0 & -3 & 0 \\ 0 & 0 & -2
        \end{pmatrix}=U.
    \end{equation}
    即$A=(P_3P_2P_1)^{-1}U$，其中$U$为上三角矩阵，我们希望$L=(P_3P_2P_1)^{-1}$为下三角矩阵（事实上我们做的初等行变换保证了这一点），因此我们只需要求出$L$即可. 由于$P_1,P_2,P_3$均为初等矩阵，因此它们的逆矩阵也是初等矩阵，我们可以很容易地求出它们的逆矩阵为$L=(P_3P_2P_1)^{-1}=P_1^{-1}P_2^{-1}P_3^{-1}=\begin{pmatrix}
            1 & 0 & 0 \\ 2 & 1 & 0 \\ -3 & -\dfrac{7}{3} & 1
        \end{pmatrix}$. 因此根据式~\eqref{eq:10:LU分解} 可知我们得到了$A$的$LU$分解为
    \[A=\begin{pmatrix}
            1 & 0 & 0 \\ 2 & 1 & 0 \\ -3 & -\dfrac{7}{3} & 1
        \end{pmatrix}\begin{pmatrix}
            1 & 2 & -1 \\ 0 & -3 & 0 \\ 0 & 0 & -2
        \end{pmatrix}=LU.\]
    接下来我们利用分解求解方程组$Ax=b$. 我们首先令$c=Ux$，则方程组可以化为$Lc=b$，写成方程组的形式即为
    \[\begin{cases}
            c_1=3 \\ 2c_1+c_2=3 \\ -3c_1-\dfrac{7}{3}c_2+c_3=-6
        \end{cases},\]
    从第一个方程开始解很容易得到$c_1=3,c_2=-3,c_3=-4$，因此我们得到了$c=(3,-3,-4)^\mathrm{T}$. 最后我们利用$Ux=c$求解出$x$，即
    \[\begin{cases}
            x_1+2x_2-x_3=3 \\ -3x_2=-3 \\ -2x_3=-4
        \end{cases},\]
    从第三个方程开始解很容易得到$x_1=3,x_2=1,x_3=2$，因此我们得到了$x=(3,1,2)^\mathrm{T}$.
\end{solve}

\begin{exercise}\label{prob:L9-例-30}
设$A=\begin{pmatrix}1 & 0 & 0 \\ 1 & 1 & 0 \\ 1 & 1 & 1\end{pmatrix},\ B=\begin{pmatrix}0 & 1 & 1 \\ 1 & 0 & 1 \\ 1 & 1 & 0\end{pmatrix}$，求矩阵$X$满足：
    \[AXA+BXB=AXB+BXA+A(A-B)\]
\end{exercise}
\noindent 来源：\href{https://github.com/yhwu-is/Linear-Algebra-Left-Undone/blob/new/%E8%AE%B2%E4%B9%89/%E4%B8%93%E9%A2%98/9%20%E7%9F%A9%E9%98%B5%E8%BF%90%E7%AE%97%E8%BF%9B%E9%98%B6.tex}{L9-例-30}。
\begin{solve}
由$AXA+BXB=AXB+BXA+A(A-B)$得$(A-B)X(A-B)=A(A-B)$，又$A-B=\begin{pmatrix}
            1 & -1 & -1 \\ 0 & 1 & -1 \\ 0 & 0 & 1
        \end{pmatrix}$，很容易判断$A-B$可逆（可以参考原讲义“可逆等价条件”处），因此$(A-B)X=A$，即$X=(A-B)^{-1}A$.

    我们直接对$(A-B,A)$做初等行变换，有
    \[\left(\begin{array}{ccc:ccc}
                1 & -1 & -1 & 1 & 0 & 0 \\
                0 & 1  & -1 & 1 & 1 & 0 \\
                0 & 0  & 1  & 1 & 1 & 1
            \end{array}\right)\xrightarrow{\text{初等行变换}}\left(\begin{array}{ccc:ccc}
                1 & 0 & 0 & 4 & 3 & 2 \\
                0 & 1 & 0 & 2 & 2 & 1 \\
                0 & 0 & 1 & 1 & 1 & 1
            \end{array}\right),\]
    因此$X=\begin{pmatrix}
            4 & 3 & 2 \\ 2 & 2 & 1 \\ 1 & 1 & 1
        \end{pmatrix}$.
\end{solve}

\begin{exercise}\label{prob:L9-例-32}
若$A,B$为两个$n$阶矩阵，则
    \begin{enumerate}[label=\Alph*.]
        \item $r(A,B)=r(A^\mathrm{T},B^\mathrm{T})$

        \item $r(A,AB)=r(A)$

        \item $r(A,B)=\max\{r(A), r(B)\}$

        \item $r(A,BA)=r(A)$
    \end{enumerate}
\end{exercise}
\noindent 来源：\href{https://github.com/yhwu-is/Linear-Algebra-Left-Undone/blob/new/%E8%AE%B2%E4%B9%89/%E4%B8%93%E9%A2%98/9%20%E7%9F%A9%E9%98%B5%E8%BF%90%E7%AE%97%E8%BF%9B%E9%98%B6.tex}{L9-例-32}。
\begin{solve}
选择B项.
    \begin{enumerate}[label=\Alph*.]
        \item 要注意分块矩阵$(A\enspace B)$的转置是
              \[\begin{pmatrix}
                      A^\mathrm{T} \\ B^\mathrm{T}
                  \end{pmatrix}\]
              而非这里给出的$(A^\mathrm{T}, B^\mathrm{T})$. 因此不能应用``分块矩阵初等变换不改变矩阵的秩''，秩不一定相等. 事实上我们很好举出例子来说明这一点，例如
              \[A=\begin{pmatrix}
                      1 & 0 \\ 0 & 0
                  \end{pmatrix},\enspace B=\begin{pmatrix}
                      0 & 0 \\ 1 & 0
                  \end{pmatrix}.\]

        \item 我们在矩阵乘法一节中介绍过，左乘矩阵相当于做初等行变换，右乘矩阵相当于做初等列变换，因此$AB$相当于对$A$做了初等列变换，考察$(A\enspace AB)$的列向量组，显然$AB$的列向量组是$A$的列向量组的线性组合，因此$r(A,AB)=r(A)$.

              当然我们也可以有另一种证明方式. 事实上直接看列向量组，显然$(A,AB)$的列向量组包括了$A$的列向量组，因此$r(A,AB)\geqslant r(A)$. 反过来，$r(A,AB)=r(A(E,B))\leqslant r(A)$，其中$E$是$n$阶单位矩阵，综上$r(A,AB)=r(A)$.

        \item 显然$r(A,B)\geqslant \max\{r(A), r(B)\}$，这只需要看矩阵列向量组的秩即可轻松看出.

        \item 可以将$BA$理解为对$A$做初等行变换，但行变换不能像B选项那样直接得出秩相关的结论，因为$A$和$BA$的行混在一起，无法判断二者行秩的关系.

              另一方面，B选项的另一种证明也不可使用，因为下列写法根据分块矩阵乘法块的匹配规则是不合法的：$(A,BA)=(E,B)A$，因为$(E,B)$是$1\times 2$分块，后面应该乘以一个两行分块.

              我们直接举反例否定这一结论，例如
              \[A=\begin{pmatrix}
                      1 & 0 \\ 0 & 0
                  \end{pmatrix},\enspace B=\begin{pmatrix}
                      1 & 1 \\ 1 & 1
                  \end{pmatrix}.\]
    \end{enumerate}
\end{solve}

\section{行列式}\label{sec:topic10}
本节收录 79 道题。按题目结构选择展开、初等变换、分块或多线性方法；伴随矩阵中的代数余子式带有转置，阶数退化也须单独检查。
\begin{exercise}\label{prob:L10-A-2}
设$\alpha_1,\alpha_2,\alpha_3$为三维列向量，令$A=(\alpha_1,\alpha_2,\alpha_3)$，且$|A|=2$，求$|\alpha_1+\alpha_2+\alpha_3,\alpha_1+3\alpha_2+9\alpha_3,\alpha_1+4\alpha_2+16\alpha_3|$.
\end{exercise}
\noindent 来源：\href{https://github.com/yhwu-is/Linear-Algebra-Left-Undone/blob/new/%E8%AE%B2%E4%B9%89/%E4%B8%93%E9%A2%98/10%20%E8%A1%8C%E5%88%97%E5%BC%8F.tex}{L10-A-2}。
\begin{solve}
使用倍加列变换的性质，可得$|\alpha_1+\alpha_2+\alpha_3,\alpha_1+3\alpha_2+9\alpha_3,\alpha_1+4\alpha_2+16\alpha_3|=6|\alpha_1,\alpha_2,\alpha_3|=12$.
\end{solve}

\begin{exercise}\label{prob:L10-A-3}
求证以下命题：
        \begin{enumerate}
            \item 奇数阶反对称矩阵不可逆；

            \item 若$A$是$n$阶可逆对称矩阵，$B$是$n$阶反对称矩阵，则当$n$为奇数时，齐次线性方程组$(AB)X=O$有非零解.
        \end{enumerate}
\end{exercise}
\noindent 来源：\href{https://github.com/yhwu-is/Linear-Algebra-Left-Undone/blob/new/%E8%AE%B2%E4%B9%89/%E4%B8%93%E9%A2%98/10%20%E8%A1%8C%E5%88%97%E5%BC%8F.tex}{L10-A-3}。
\begin{solve}
\begin{enumerate}
                \item 首先由反对称矩阵的性质，有$A^T=-A$，其次由矩阵与行列式的性质，可得$|A^T|=|A|$，则可推出$|A|=|-A|=(-1)^n|A|$，又n为奇数，故$|A|=0$，矩阵$A$不可逆.

                \item 只需证明$AB$的秩小于$n$，即$|AB|=0$即可. $B$是奇数阶反对称矩阵，$|B|=0$,则$|AB|=|A||B|=0$得证.
            \end{enumerate}
\end{solve}

\begin{exercise}\label{prob:L10-A-4}
设 $A,B$ 分别为 $m,n$ 阶可逆矩阵，$a=\det A$、$b=\det B$。求分块矩阵
\[M=\begin{pmatrix}A&O\\O&B\end{pmatrix},\qquad N=\begin{pmatrix}O&A\\B&O\end{pmatrix}\]
的伴随矩阵。$N$ 中第一块列宽为 $n$，第二块列宽为 $m$。
\end{exercise}
\noindent 来源：\href{https://github.com/yhwu-is/Linear-Algebra-Left-Undone/blob/new/%E8%AE%B2%E4%B9%89/%E4%B8%93%E9%A2%98/10%20%E8%A1%8C%E5%88%97%E5%BC%8F.tex}{L10-A-4}。
\begin{solve}
由 $\det M=ab$ 及 $M^{-1}=\operatorname{diag}(A^{-1},B^{-1})$，得到
\[\operatorname{adj}(M)=\begin{pmatrix}b\operatorname{adj}(A)&O\\O&a\operatorname{adj}(B)\end{pmatrix}.\]
交换 $M$ 的两个块列，共需 $mn$ 次普通列交换，得到 $N$，故 $\det N=(-1)^{mn}ab$。又
\[N^{-1}=\begin{pmatrix}O&B^{-1}\\A^{-1}&O\end{pmatrix},\]
所以
\[\operatorname{adj}(N)=(-1)^{mn}\begin{pmatrix}O&a\operatorname{adj}(B)\\b\operatorname{adj}(A)&O\end{pmatrix}.\]
注意逆矩阵的块行高度依次为 $n,m$，与原矩阵的块列宽度对应。
\end{solve}

\begin{exercise}\label{prob:L10-A-5}
证明：上（下）三角矩阵的伴随矩阵是上（下）三角矩阵（对角矩阵为特例）.
\end{exercise}
\noindent 来源：\href{https://github.com/yhwu-is/Linear-Algebra-Left-Undone/blob/new/%E8%AE%B2%E4%B9%89/%E4%B8%93%E9%A2%98/10%20%E8%A1%8C%E5%88%97%E5%BC%8F.tex}{L10-A-5}。
\begin{solve}
先设 $A$ 可逆且为上三角矩阵，则 $A^{-1}$ 也是上三角矩阵，故 $\operatorname{adj}(A)=\det(A)A^{-1}$ 为上三角矩阵。若 $A$ 奇异，考察 $A+tI$：除有限多个 $t$ 外，它都可逆且上三角，所以 $\operatorname{adj}(A+tI)$ 的严格下三角元为零。每个这样的元素都是关于 $t$ 的多项式，在无限多个值上为零，因而恒为零；令 $t=0$，结论仍成立。下三角情形通过转置得到，因为 $\operatorname{adj}(A^{\mathrm T})=\operatorname{adj}(A)^{\mathrm T}$。
\end{solve}

\begin{exercise}\label{prob:L10-A-6}
设 $n\ge2$，$A$ 为 $n$ 阶方阵且 $\det A=0$。证明：代数余子式矩阵的任意两个行向量以及任意两个列向量均线性相关。若两个向量均非零，则它们互为非零标量倍。
\end{exercise}
\noindent 来源：\href{https://github.com/yhwu-is/Linear-Algebra-Left-Undone/blob/new/%E8%AE%B2%E4%B9%89/%E4%B8%93%E9%A2%98/10%20%E8%A1%8C%E5%88%97%E5%BC%8F.tex}{L10-A-6}。
\begin{solve}
若 $r(A)\le n-2$，则所有 $n-1$ 阶子式为零，伴随矩阵为零。若 $r(A)=n-1$，则伴随矩阵至少有一个非零元素；又 $A\operatorname{adj}(A)=O$，所以其每一列都属于一维空间 $\ker A$，从而 $r(\operatorname{adj}(A))=1$。代数余子式矩阵为伴随矩阵的转置，两者秩相同，均不超过 $1$，因此任意两个行向量或列向量线性相关。使用“线性相关”还涵盖一个向量为零、另一个非零的情形，避免“成比例”的歧义。
\end{solve}

\begin{exercise}\label{prob:L10-A-7}
设向量$\alpha_1,\alpha_2,\alpha_3$线性无关，讨论向量$\alpha_1-\alpha_2-2\alpha_3,\ 2\alpha_1+\alpha_2-\alpha_3,\ 3\alpha_1+\alpha_2+2\alpha_3$的线性相关性.
\end{exercise}
\noindent 来源：\href{https://github.com/yhwu-is/Linear-Algebra-Left-Undone/blob/new/%E8%AE%B2%E4%B9%89/%E4%B8%93%E9%A2%98/10%20%E8%A1%8C%E5%88%97%E5%BC%8F.tex}{L10-A-7}。
\begin{solve}
\[(\alpha_1-\alpha_2-2\alpha_3,2\alpha_1+\alpha_2-\alpha_3,3\alpha_1+\alpha_2+2\alpha_3)=(\alpha_1,\alpha_2,\alpha_3)\begin{pmatrix}
                1 & 2 & 3 \\ -1 & 1 & 1 \\ -2 & -1 & 2
            \end{pmatrix}\]
，又$\begin{vmatrix}
                1 & 2 & 3 \\ -1 & 1 & 1 \\ -2 & -1 & 2
            \end{vmatrix}=12\neq 0$，因此该矩阵可逆，则$(\alpha_1-\alpha_2-2\alpha_3,2\alpha_1+\alpha_2-\alpha_3,3\alpha_1+\alpha_2+2\alpha_3)$与$(\alpha_1,\alpha_2,\alpha_3)$秩相等，因此这三个向量线性无关.
\end{solve}

\begin{exercise}\label{prob:L10-A-8}
设$W=\operatorname{span}(\alpha_1,\alpha_2)$是$\mathbf{R}^4$的一个子空间，其中$\alpha_1=(1,2,1,-1)^\mathrm{T}$，$\alpha_2=(1,4,-1,-1)^\mathrm{T}$，试将$\alpha_1,\alpha_2$扩充为$\mathbf{R}^4$的基.
\end{exercise}
\noindent 来源：\href{https://github.com/yhwu-is/Linear-Algebra-Left-Undone/blob/new/%E8%AE%B2%E4%B9%89/%E4%B8%93%E9%A2%98/10%20%E8%A1%8C%E5%88%97%E5%BC%8F.tex}{L10-A-8}。
\begin{solve}
只要取$\alpha_3,\alpha_4 \in \mathbf{R}^4$，使得$D(\alpha_1,\alpha_2,\alpha_3,\alpha_4) \neq 0$即可. 易得$\begin{vmatrix}
                1  & 1  & 0 & 0 \\
                2  & 4  & 0 & 0 \\
                1  & -1 & 1 & 0 \\
                -1 & 1  & 0 & 1
            \end{vmatrix}=\begin{vmatrix}
                1 & 1 \\
                2 & 4
            \end{vmatrix} \cdot \begin{vmatrix}
                1 & 0 \\
                0 & 1
            \end{vmatrix}=2 \neq 0$. 于是可取$\alpha_3=(0,0,1,0)^T,\alpha_4=(0,0,0,1)^T$，能使得${\alpha_1,\alpha_2,\alpha_3,\alpha_4}$为$\mathbf{R}^4$的一组基.
\end{solve}

\begin{exercise}\label{prob:L10-A-9}
$ D_n=\begin{vmatrix}
                1      & 2      & 3      & 4      & \cdots & n-1    & n      \\
                x      & 1      & 2      & 3      & \cdots & n-2    & n-1    \\
                x      & x      & 1      & 2      & \cdots & n-3    & n-2    \\
                x      & x      & x      & 1      & \cdots & n-4    & n-3    \\
                \vdots & \vdots & \vdots & \vdots & \ddots & \vdots & \vdots \\
                x      & x      & x      & x      & \cdots & 1      & 2      \\
                x      & x      & x      & x      & \cdots & x      & 1      \\
            \end{vmatrix}$（提示：考虑滚动消去法）.
\end{exercise}
\noindent 来源：\href{https://github.com/yhwu-is/Linear-Algebra-Left-Undone/blob/new/%E8%AE%B2%E4%B9%89/%E4%B8%93%E9%A2%98/10%20%E8%A1%8C%E5%88%97%E5%BC%8F.tex}{L10-A-9}。
\begin{solve}
考虑滚动消去法，每行减下一行，得
            \[ \begin{vmatrix}
                    1-x    & 1      & 1      & 1      & \ldots & 1      & 1      \\
                    0      & 1-x    & 1      & 1      & \ldots & 1      & 1      \\
                    0      & 0      & 1-x    & 1      & \ldots & 1      & 1      \\
                    0      & 0      & 0      & 1-x    & \ldots & 1      & 1      \\
                    \vdots & \vdots & \vdots & \vdots & \ddots & \vdots & \vdots \\
                    0      & 0      & 0      & 0      & \ldots & 1-x    & 1      \\
                    x      & x      & x      & x      & \ldots & x      & 1      \\
                \end{vmatrix}. \]
            第一行乘$(-x)$到第$n$行，得
            \[\begin{vmatrix}
                    1-x    & 1      & 1      & 1      & \ldots & 1      & 1      \\
                    0      & 1-x    & 1      & 1      & \ldots & 1      & 1      \\
                    0      & 0      & 1-x    & 1      & \ldots & 1      & 1      \\
                    0      & 0      & 0      & 1-x    & \ldots & 1      & 1      \\
                    \vdots & \vdots & \vdots & \vdots & \ddots & \vdots & \vdots \\
                    0      & 0      & 0      & 0      & \ldots & 1-x    & 1      \\
                    x^2    & 0      & 0      & 0      & \ldots & 0      & 1-x    \\
                \end{vmatrix}.\]
            按第$n$行展开，得$D_n=(1-x)^n+(-1)^{n+1}x^2T_{n-1}$，其中
            \[T_{n-1}=\begin{vmatrix}
                    1      & 1      & 1      & \ldots & 1      & 1      \\
                    1-x    & 1      & 1      & \ldots & 1      & 1      \\
                    0      & 1-x    & 1      & \ldots & 1      & 1      \\
                    \vdots & \vdots & \vdots & \ddots & \vdots & \vdots \\
                    0      & 0      & 0      & \ldots & 1      & 1      \\
                    0      & 0      & 0      & \ldots & 1-x    & 1      \\
                \end{vmatrix}.\]
            按第$n-1$行展开$T_{n-1}$，有$T_{n-1}=T_{n-2}+(-1)(1-x)T_{n-2}=xT_{n-2}$. 由$T_2=\begin{vmatrix}
                    1&1\\1-x&1\end{vmatrix}=x$，递推得$T_{n-1}=x^{n-2}$. 代入可解得$D_n=(1-x)^n-(-x)^n$.
\end{solve}

\begin{exercise}\label{prob:L10-A-10}
$ D_n=\begin{vmatrix}
                a_1 & b_1 &     &        &         \\
                    & a_2 & b_2 &        &         \\
                    &     & a_3 & \ddots &         \\
                    &     &     & \ddots & b_{n-1} \\
                b_n &     &     &        & a_n
            \end{vmatrix}$.
\end{exercise}
\noindent 来源：\href{https://github.com/yhwu-is/Linear-Algebra-Left-Undone/blob/new/%E8%AE%B2%E4%B9%89/%E4%B8%93%E9%A2%98/10%20%E8%A1%8C%E5%88%97%E5%BC%8F.tex}{L10-A-10}。
\begin{solve}
按第一列展开，得$D_n=\displaystyle\prod_{i=1}^na_i+(-1)^{n+1}\prod_{i=1}^nb_i$.
\end{solve}

\begin{exercise}\label{prob:L10-A-11}
$D_n=\begin{vmatrix}
                a+b & ab  &        &     &     \\
                1   & a+b & \ddots &     &     \\
                    & 1   & \ddots & ab  &     \\
                    &     & \ddots & a+b & ab  \\
                    &     &        & 1   & a+b \\
            \end{vmatrix}$.
\end{exercise}
\noindent 来源：\href{https://github.com/yhwu-is/Linear-Algebra-Left-Undone/blob/new/%E8%AE%B2%E4%B9%89/%E4%B8%93%E9%A2%98/10%20%E8%A1%8C%E5%88%97%E5%BC%8F.tex}{L10-A-11}。
\begin{solve}
按第一列展开，得
            \begin{align*}
                D_{n} & =(a+b) D_{n-1}-\begin{vmatrix}
                                            a b    & 0      & 0      & 0   & \cdots & 0      & 0      \\
                                            1      & a+b    & a b    & 0   & \cdots & 0      & 0      \\
                                            0      & 1      & a+b    & a b & \cdots & 0      & 0      \\
                                            \vdots & \vdots & \vdots & a+b & \ddots & \vdots & \vdots \\
                                            0      & 0      & 0      & 0   & \cdots & a+b    & a b    \\
                                            0      & 0      & 0      & 0   & \cdots & 1      & a+b
                                        \end{vmatrix} \\
                        & =(a+b) D_{n-1}-a b D_{n-2}.
            \end{align*}
            \begin{enumerate}
                \item 方法一：特征方程法. 有特征方程$r^2-(a+b)r+ab=(r-a)(r-b)=0$，先考虑$a\neq b$，则$D_n=C_1a^n+C_2b^n$，$D_{1}=a+b$，$D_{2}=a^{2}+a b+b^{2}$，则有
                        \[\begin{cases}
                                C_1a+C_2b=a+b \\
                                C_1a^2+C_2b^2=a^2+ab+b^2
                            \end{cases},\]
                        考虑Cramer法则，首先判断系数行列式$\begin{vmatrix}
                                a   & b   \\
                                a^2 & b^2
                            \end{vmatrix}=ab(b-a)$，若$a=b=0$，原行列式为0；若$a=0,b\neq 0$，则有$C_2=1$，$D_n=b^n$；同理若$a=0,b\neq 0$，$D_n=a^n$. 考虑$a\neq 0,b\neq 0$，则系数行列式非0，根据Cramer法则可以解得
                        \[C_1=\frac{\begin{vmatrix}
                                    a+b        & b   \\
                                    a^2+ab+b^2 & b^2
                                \end{vmatrix}}{ab(b-a)}=\frac{a}{a-b},\quad
                            C_2=\frac{\begin{vmatrix}
                                    a   & a+b        \\
                                    a^2 & a^2+ab+b^2
                                \end{vmatrix}}{ab(b-a)}=-\frac{b}{a-b}.\]
                        故$D_n=\displaystyle\frac{a^{n+1}-b^{n+1}}{a-b}$. 考虑到$a$或$b$为0时也符合该公式，故合并为同一公式. 再考虑$a=b$，则$D_n=(C_1+C_2n)a^n$，有方程
                        \[\begin{cases}
                                (C_1+C_2)a=2a \\
                                (C_1+2C_2)a^2=3a^2
                            \end{cases},\]
                        $a=0$，则$D_n=0$；$a\neq 0$，可以解得$C_1=C_2=1$，有$D_n=(n+1)a^n$. 考虑到$a=0$时$D_n$也符合该式，故合并. 综上有
                        \[D_n=\begin{cases}
                                \displaystyle\frac{a^{n+1}-b^{n+1}}{a-b}, & a\neq b \\
                                (n+1)a^n,                                 & a=b
                            \end{cases}.\]

                \item 方法二：递推式变形法. 递推式变形, 得
                        \[\mathrm{D}_{n}-a D_{n-1}=b\left(D_{n-1}-a D_{n-2}\right),\]
                        由于 $D_{1}=a+b, D_{2}=a^{2}+a b+b^{2}$,从而利用上述递推公式得
                        \begin{align*}
                            \mathrm{D}_{n}-a D_{n-1} & =b\left(D_{n-1}-a D_{n-2}\right)                \\
                                                    & =b^{2}\left(D_{n-2}-a D_{n-3}\right)            \\
                                                    & =\cdots=b^{n-2}\left(D_{2}-a D_{1}\right)=b^{n}.
                        \end{align*}
                        故\begin{align*}
                            D_{n} & =a D_{n-1}+b^{n}=a\left(a D_{n-2}+b^{n-1}\right)+b^{n}     \\
                                  & =\cdots=a^{n-1} D_{1}+a^{n-2} b^{2}+\cdots+a b^{n-1}+b^{n} \\
                                  & =a^{n}+a^{n-1} b+\cdots+a b^{n-1}+b^{n}                    \\
                                  & =\begin{cases}
                                        \displaystyle\frac{a^{n+1}-b^{n+1}}{a-b}, & a\neq b \\
                                        (n+1)a^n,                                 & a=b
                                     \end{cases}.
                        \end{align*}
            \end{enumerate}
\end{solve}

\begin{exercise}\label{prob:L10-A-13}
解行列式 $ \begin{vmatrix}
                a^{2} & (a+1)^{2} & (a+2)^{2} & (a+3)^{2} \\
                b^{2} & (b+1)^{2} & (b+2)^{2} & (b+3)^{2} \\
                c^{2} & (c+1)^{2} & (c+2)^{2} & (c+3)^{2} \\
                d^{2} & (d+1)^{2} & (d+2)^{2} & (d+3)^{2}
            \end{vmatrix} $.
\end{exercise}
\noindent 来源：\href{https://github.com/yhwu-is/Linear-Algebra-Left-Undone/blob/new/%E8%AE%B2%E4%B9%89/%E4%B8%93%E9%A2%98/10%20%E8%A1%8C%E5%88%97%E5%BC%8F.tex}{L10-A-13}。
\begin{solve}
\begin{enumerate}
                \item 方法一：慢慢消去.
                      \begin{align*}
                          \begin{vmatrix}
                              a^{2} & (a+1)^{2} & (a+2)^{2} & (a+3)^{2} \\
                              b^{2} & (b+1)^{2} & (b+2)^{2} & (b+3)^{2} \\
                              c^{2} & (c+1)^{2} & (c+2)^{2} & (c+3)^{2} \\
                              d^{2} & (d+1)^{2} & (d+2)^{2} & (d+3)^{2}
                          \end{vmatrix}
                           & = \begin{vmatrix}
                                   a^{2} & 2a+1 & 4a+4 & 6a+9 \\
                                   b^{2} & 2b+1 & 4b+4 & 6b+9 \\
                                   c^{2} & 2c+1 & 4c+4 & 6c+9 \\
                                   d^{2} & 2d+1 & 4d+4 & 6d+9
                               \end{vmatrix} \\
                           & = \begin{vmatrix}
                                   a^{2} & 2a+1 & 4a+4 & 4 \\
                                   b^{2} & 2b+1 & 4b+4 & 4 \\
                                   c^{2} & 2c+1 & 4c+4 & 4 \\
                                   d^{2} & 2d+1 & 4d+4 & 4
                               \end{vmatrix}    \\
                           & = 4\begin{vmatrix}
                                    a^{2} & 2a & 4a & 1 \\
                                    b^{2} & 2b & 4b & 1 \\
                                    c^{2} & 2c & 4c & 1 \\
                                    d^{2} & 2d & 4d & 1
                                \end{vmatrix}=0.
                      \end{align*}

                \item 关注到行列式中只有$(a^2, b^2, c^2, d^2)^\mathrm{T}, (a, b, c, d)^\mathrm{T}$和$(1, 1, 1, 1)^\mathrm{T}$三个线性无关列向量，则一眼看出线性相关. 因此设原行列式为$|\alpha_1, \alpha_2, \alpha_3, \alpha_4|$，有$\alpha_4=\alpha_1-3\alpha_2+3\alpha_3$，因此原行列式值为0.
            \end{enumerate}
\end{solve}

\begin{exercise}\label{prob:L10-A-14}
设
        \[ D=\begin{vmatrix}
                1+a_1  & 1      & \cdots & 1      \\
                1      & 1+a_2  & \cdots & 1      \\
                \vdots & \vdots & \ddots & \vdots \\
                1      & 1      & \cdots & 1+a_n
            \end{vmatrix} \]
        \begin{enumerate}
            \item 用递推公式计算行列式$D$；

            \item 硬拆$D$为$2^n$个行列式，计算出结果.
        \end{enumerate}
\end{exercise}
\noindent 来源：\href{https://github.com/yhwu-is/Linear-Algebra-Left-Undone/blob/new/%E8%AE%B2%E4%B9%89/%E4%B8%93%E9%A2%98/10%20%E8%A1%8C%E5%88%97%E5%BC%8F.tex}{L10-A-14}。
\begin{solve}
\begin{enumerate}
                \item \begin{align*}
                          D_n & =\begin{vmatrix}
                                     D_{n-1} & a   \\
                                     b       & a_n
                                 \end{vmatrix},a=(0, \ldots, 0, -a_{n-1})^\mathrm{T}, b=(1, \ldots, 1) \\
                              & =a_nD_{n-1}+a_{n-1}\begin{vmatrix}
                                                       D_{n-2} & 1_{n-2} \\
                                                       1_{n-2} & 1
                                                   \end{vmatrix}
                          \quad \text{($1_{n-2}$表示由$(n-2)$个1构成的行/列向量)}                      \\
                              & =a_nD_{n-1}+a_{n-1}
                          \begin{vmatrix}
                              a_1    & 0      & 0      & \ldots & 0       & 1      \\
                              0      & a_2    & 0      & \ldots & 0       & 1      \\
                              0      & 0      & a_3    & \ldots & 0       & 1      \\
                              \vdots & \vdots & \vdots & \ddots & \vdots  & \vdots \\
                              0      & 0      & 0      & \ldots & a_{n-2} & 1      \\
                              0      & 0      & 0      & \ldots & 0       & 1      \\
                          \end{vmatrix} \quad \text{($1-n$列用第$n$列减)}                         \\
                              & =a_nD_{n-1}+\prod_{i=1}^{n-1}a_i.
                      \end{align*}
                      递推可得$\displaystyle D_n=\left(\prod_{k=1}^na_k\right)
                          \left(1+\sum_{i=1}^n\frac{1}{a_i}\right)$.

                \item 将 1 写成 $1+0$, 将 $D$ 硬拆成 $2^{n}$ 个行列式, 只有如下的 $n+1$ 个行列式非 0：
                      \begin{align*}
                          D_n = & \begin{vmatrix}
                                      1      &       &        &       \\
                                      1      & a_{2} &        &       \\
                                      \vdots &       & \ddots &       \\
                                      1      &       &        & a_{n}
                                  \end{vmatrix}
                          +\begin{vmatrix}
                               a_{1} & 1      &                \\
                                     & 1      &                \\
                                     & \vdots & \ddots &       \\
                                     & 1      &        & a_{n}
                           \end{vmatrix}+\cdots+                    \\
                                & \begin{vmatrix}
                                      a_{1} &       &        & 1      \\
                                            & a_{2} &        & 1      \\
                                            &       & \ddots & \vdots \\
                                            &       &        & 1
                                  \end{vmatrix}
                          +\begin{vmatrix}
                               a_{1} &                      \\
                                     & a_{2} &              \\
                                     &       & \ddots &     \\
                                     &       &        & a_n
                           \end{vmatrix}                       \\
                          =     & \left(\prod_{k=1}^na_k\right)
                          \left(1+\sum_{i=1}^n\frac{1}{a_i}\right).
                      \end{align*}
            \end{enumerate}
\end{solve}

\begin{exercise}\label{prob:L10-A-15}
解行列式
        \begin{enumerate}
            \item $D=\begin{vmatrix}
                          1      & 2      & \cdots & 2      & 2      \\
                          2      & 2      & \cdots & 2      & 2      \\
                          \vdots & \vdots & \ddots & \vdots & \vdots \\
                          2      & 2      & \cdots & n-1    & 2      \\
                          2      & 2      & \cdots & 2      & n\end{vmatrix}$；

            \item $D=\begin{vmatrix}
                          1      & 2      & \cdots & n-1    & n      \\
                          2      & 3      & \cdots & n      & 1      \\
                          3      & 4      & \cdots & 1      & 2      \\
                          \vdots & \vdots & \ddots & \vdots & \vdots \\
                          n      & 1      & \cdots & n-2    & n-1
                      \end{vmatrix}$.
        \end{enumerate}
\end{exercise}
\noindent 来源：\href{https://github.com/yhwu-is/Linear-Algebra-Left-Undone/blob/new/%E8%AE%B2%E4%B9%89/%E4%B8%93%E9%A2%98/10%20%E8%A1%8C%E5%88%97%E5%BC%8F.tex}{L10-A-15}。
\begin{solve}
\begin{enumerate}
                \item 行列式中2比较多，用全是2的第二行去消，然后按第2行展开.
                      \[D=\begin{vmatrix}
                              -1     & 0      & \cdots & 0      & 0      \\
                              2      & 2      & \cdots & 2      & 2      \\
                              \vdots & \vdots & \ddots & \vdots & \vdots \\
                              0      & 0      & \cdots & n-3    & 0      \\
                              0      & 0      & \cdots & 0      & n-2
                          \end{vmatrix}=2\cdot(-1)\cdot(n-2)!=-2(n-2)!.\]

                \item \begin{enumerate}
                          \item 方法一：考虑滚动消去法，有
                                \begin{align*}
                                    D & =\begin{vmatrix}
                                             1      & 2      & 3      & \cdots & n-1    & n      \\
                                             1      & 1      & 1      & \cdots & 1      & 1-n    \\
                                             1      & 1      & 1      & \cdots & 1-n    & 1      \\
                                             \vdots & \vdots & \vdots & \ddots & \vdots & \vdots \\
                                             1      & 1      & 1-n    & \cdots & 1      & 1      \\
                                             1      & 1-n    & 1      & \cdots & 1      & 1
                                         \end{vmatrix}                                    \\&=\begin{vmatrix}
                                        0      & n+1    & 2      & \cdots & n-2    & n-1    \\
                                        0      & n      & 0      & \cdots & 0      & -n     \\
                                        0      & n      & 0      & \cdots & -n     & 0      \\
                                        \vdots & \vdots & \vdots & \ddots & \vdots & \vdots \\
                                        0      & n      & -n     & \cdots & 0      & 0      \\
                                        1      & 1-n    & 1      & \cdots & 1      & 1
                                    \end{vmatrix}\\
                                      & =(-1)^{n+1}\begin{vmatrix}
                                                       n+1    & 2      & \cdots & n-2    & n-1    \\
                                                       n      & 0      & \cdots & 0      & -n     \\
                                                       n      & 0      & \cdots & -n     & 0      \\
                                                       \vdots & \vdots & \ddots & \vdots & \vdots \\
                                                       n      & -n     & \cdots & 0      & 0
                                                   \end{vmatrix}                                   \\
                                      & =(-1)^{n+1}(-1)^{\frac{(n-2)(n-3)}{2}}
                                    \begin{vmatrix}
                                        n+1    & n-1    & \cdots & 3      & 2      \\
                                        n      & -n     & \cdots & 0      & 0      \\
                                        n      & 0      & \cdots & 0      & 0      \\
                                        \vdots & \vdots & \ddots & \vdots & \vdots \\
                                        n      & 0      & \cdots & 0      & -n
                                    \end{vmatrix}                                                  \\
                                      & =(-1)^{\frac{n^2-3n+8}{2}}\begin{vmatrix}
                                                                      \frac{n(n+1)}{2} & n-1    & \cdots & 3      & 2      \\
                                                                      0                & -n     & \cdots & 0      & 0      \\
                                                                      0                & 0      & \cdots & 0      & 0      \\
                                                                      \vdots           & \vdots & \ddots & \vdots & \vdots \\
                                                                      0                & 0      & \cdots & 0      & -n
                                                                  \end{vmatrix}=(-1)^{\frac{n^2-n+12}{2}}\frac{n^{n-1}(n+1)}{2}.
                                \end{align*}

                          \item 方法二：考虑连加法. 将第 $2,3, \ldots, n$ 列都加到第 1 列, 提出公因子 $\dfrac{1}{2} n(n+1)$, 再依次将第 $n-1$ 行乘 $(-1)$ 加到第 $n$ 行, $\cdots$, 第 2 行乘 $(-1)$ 加到第 3 行, 第 1 行乘 $(-1)$ 加到第 2 行, 然后对第 1 列展开, 得到一个 $n-1$ 阶行列式, 它的副对角元为 $1-$ $n$, 其余元素均为 1 . 再把它的各列加到第 1 列, 并把它的第 1 行乘 $(-1)$ 加到其余各行, 得
                                \begin{align*}
                                    D & =\frac{n(n+1)}{2}\begin{vmatrix}
                                                             1      & 2      & \cdots & n-1    & n      \\
                                                             1      & 3      & \cdots & n      & 1      \\
                                                             1      & 4      & \cdots & 1      & 2      \\
                                                             \vdots & \vdots & \ddots & \vdots & \vdots \\
                                                             1      & 1      & \cdots & n-2    & n-1
                                                         \end{vmatrix} \\
                                      & =\frac{n(n+1)}{2}\begin{vmatrix}
                                                             1      & 2      & 3      & \cdots & n      \\
                                                             0      & 1      & 1      & \cdots & 1-n    \\
                                                             \vdots & \vdots & \vdots & \ddots & \vdots \\
                                                             0      & 1      & 1-n    & \cdots & 1      \\
                                                             0      & 1-n    & 1      & \cdots & 1
                                                         \end{vmatrix}_{n} \\
                                      & =\frac{n(n+1)}{2}
                                    \begin{vmatrix}
                                        1      & 1      & \cdots & 1      & 1-n \\
                                        \vdots & \vdots & \ddots & \vdots & 1   \\
                                        1      & 1-n    & \cdots & 1      & 1   \\
                                        1-n    & 1      & \cdots & 1      & 1
                                    \end{vmatrix}                         \\
                                      & =\frac{n(n+1)}{2}\begin{vmatrix}
                                                             -1     & 1      & \cdots & 1      & 1-n    \\
                                                             0      & 0      & \cdots & -n     & n      \\
                                                             \vdots & \vdots & \ddots & \vdots & \vdots \\
                                                             0      & -n     & \cdots & 0      & n      \\
                                                             0      & 0      & \cdots & 0      & n
                                                         \end{vmatrix}_{n-1}
                                \end{align*}
                                将上式先对第 1 列展开, 得到一个 $n-2$ 阶行列式, 再将它对最后一行展开, 得
                                \begin{align*}
                                    D & =\frac{-n(n+1)}{2} n
                                    \begin{vmatrix}
                                        0      & 0      & \cdots & 0      & -n     \\
                                        0      & 0      & \cdots & -n     & 0      \\
                                        \vdots & \vdots & \ddots & \vdots & \vdots \\
                                        0      & -n     & \cdots & 0      & 0      \\
                                        -n     & 0      & \cdots & 0      & 0
                                    \end{vmatrix}_{n-3}                      \\
                                      & =-\frac{n^{2}(n+1)}{2}(-1)^{\frac{(n-3)(n-4)}{2}}(-n)^{n-3} \\
                                      & =(-1)^{\frac{n(n-1)}{2} }\frac{(n+1) n^{n-1}}{2}.
                                \end{align*}
                                事实上，两种方法得到的答案是等价的.
                      \end{enumerate}
            \end{enumerate}
\end{solve}

\begin{exercise}\label{prob:L10-A-16}
记$a_{ij}=i-j$，且$A=(a_{ij})_{n\times n}$，计算$A$的行列式.
\end{exercise}
\noindent 来源：\href{https://github.com/yhwu-is/Linear-Algebra-Left-Undone/blob/new/%E8%AE%B2%E4%B9%89/%E4%B8%93%E9%A2%98/10%20%E8%A1%8C%E5%88%97%E5%BC%8F.tex}{L10-A-16}。
\begin{solve}
\begin{enumerate}
                \item 若 $n=1$，则 $|A| = |0| = 0$；

                \item 若 $n=2$, 则 $|A| = \begin{vmatrix} 0 & -1 \\ 1 & 0 \end{vmatrix} = 1$；

                \item 若 $n>2$，考虑使用滚动消去法：
                    \begin{align*}
                        |A| &= \begin{vmatrix}
                            0      & -1     & -2     & \cdots & -n+1   \\
                            1      & 0      & -1     & \cdots & -n+2   \\
                            2      & 1      & 0      & \cdots & -n+3   \\
                            \vdots & \vdots & \vdots & \ddots & \vdots \\
                            n-1    & n-2    & n-3    & \cdots & 0
                        \end{vmatrix} \\
                        &= \begin{vmatrix}
                            1 & 1 & \cdots & 1 & -n+1 \\
                            1 & 1 & \cdots & 1 & -n+2 \\
                            \vdots & \vdots & \ddots & \vdots & \vdots \\
                            1 & 1 & \cdots & 1 & 0
                        \end{vmatrix} \\
                        &= \begin{vmatrix}
                            1 & 0 & \cdots & 0 & -n+1 \\
                            1 & 0 & \cdots & 0 & -n+2 \\
                            \vdots & \vdots & \ddots & \vdots & \vdots \\
                            1 & 0 & \cdots & 0 & 0
                        \end{vmatrix}
                        = 0.
                    \end{align*}
            \end{enumerate}
\end{solve}

\begin{exercise}\label{prob:L10-A-17}
计算行列式$D=\begin{vmatrix}
                1 & 0 & 2 & 0 & x   \\
                0 & 2 & 0 & 1 & 0   \\
                1 & 0 & 4 & 0 & x^2 \\
                0 & 5 & 0 & 3 & 0   \\
                1 & 0 & 8 & 0 & x^3
            \end{vmatrix}$.
\end{exercise}
\noindent 来源：\href{https://github.com/yhwu-is/Linear-Algebra-Left-Undone/blob/new/%E8%AE%B2%E4%B9%89/%E4%B8%93%E9%A2%98/10%20%E8%A1%8C%E5%88%97%E5%BC%8F.tex}{L10-A-17}。
\begin{solve}
我们可以先构造分块对角矩阵，然后分别处理两个低阶的行列式：
            \[
                D = (-1)^{3+3} \begin{vmatrix}
                    1 & 2 & x   & 0 & 0 \\
                    1 & 4 & x^2 & 0 & 0 \\
                    1 & 8 & x^3 & 0 & 0 \\
                    0 & 0 & 0   & 2 & 1 \\
                    0 & 0 & 0   & 5 & 3
                \end{vmatrix} = (-1)^{3+3} \begin{vmatrix}
                    1 & 2 & x   \\
                    1 & 4 & x^2 \\
                    1 & 8 & x^3
                \end{vmatrix} \begin{vmatrix}
                    2 & 1 \\
                    5 & 3
                \end{vmatrix} = \begin{vmatrix}
                    1 & 2 & x   \\
                    1 & 4 & x^2 \\
                    1 & 8 & x^3
                \end{vmatrix}.
            \]

            该行列式可以通过加边法构造范德蒙行列式：
            \[
                D = \begin{vmatrix}
                    1 & 2 & x   \\
                    1 & 4 & x^2 \\
                    1 & 8 & x^3
                \end{vmatrix} = \begin{vmatrix}
                    1 & 1 & 1 & 1   \\
                    0 & 1 & 2 & x   \\
                    0 & 1 & 4 & x^2 \\
                    0 & 1 & 8 & x^3
                \end{vmatrix} = 2x(x-1)(x-2).
            \]
\end{solve}

\begin{exercise}\label{prob:L10-B-1}
设$D=\begin{vmatrix}
                3 & 0 & 4 & 1 \\ 2 & 3 & 1 & 4 \\ 0 & -7 & 8 & 3 \\ 5 & 3 & -2 & 2
            \end{vmatrix}$，求
        \begin{enumerate}
            \item $A_{21}+A_{22}+A_{23}+A_{24}$；

            \item $A_{31}+A_{33}$；

            \item $M_{41}+M_{42}+M_{43}+M_{44}$.
        \end{enumerate}
\end{exercise}
\noindent 来源：\href{https://github.com/yhwu-is/Linear-Algebra-Left-Undone/blob/new/%E8%AE%B2%E4%B9%89/%E4%B8%93%E9%A2%98/10%20%E8%A1%8C%E5%88%97%E5%BC%8F.tex}{L10-B-1}。
\begin{solve}
这三小问实际上都是对原行列式中的某一行进行了替换进行计算.
            \begin{enumerate}
                \item \begin{align*}
                            A_{21}+A_{22}+A_{23}+A_{24}
                            & = \begin{vmatrix}
                                    3 & 0  & 4  & 1 \\
                                    1 & 1  & 1  & 1 \\
                                    0 & -7 & 8  & 3 \\
                                    5 & 3  & -2 & 2
                                \end{vmatrix}
                            = \begin{vmatrix}
                                3 & -3 & 1  & -2 \\
                                1 & 0  & 0  & 0  \\
                                0 & -7 & 8  & 3  \\
                                5 & -2 & -7 & -3
                            \end{vmatrix}             \\
                            & = (-1)^{2+1} \begin{vmatrix}
                                                -3 & 1  & -1 \\
                                                -7 & 8  & 3  \\
                                                -2 & -7 & -3
                                            \end{vmatrix} \\
                            & = 148.
                        \end{align*}

                \item \begin{align*}
                            A_{31}+A_{33}
                            & = 1A_{31}+0A_{32}+1A_{33}+0A_{34}
                            = \begin{vmatrix}
                                3 & 0 & 4  & 1 \\
                                2 & 3 & 1  & 4 \\
                                1 & 0 & 1  & 0 \\
                                5 & 3 & -2 & 2
                            \end{vmatrix}
                            = 3 \begin{vmatrix}
                                    3 & 0 & 4  & 1 \\
                                    2 & 1 & 1  & 4 \\
                                    1 & 0 & 1  & 0 \\
                                    5 & 1 & -2 & 2
                                \end{vmatrix}                   \\
                            & = 3 \begin{vmatrix}
                                    3 & 0 & 1  & 1 \\
                                    2 & 1 & -1 & 4 \\
                                    1 & 0 & 0  & 0 \\
                                    5 & 3 & -7 & 2
                                \end{vmatrix}
                            = 3 \cdot (-1)^{3+1} \begin{vmatrix}
                                                    0 & 1  & 1 \\
                                                    1 & -1 & 4 \\
                                                    3 & -7 & 2
                                                \end{vmatrix}  \\
                            & = -12.
                        \end{align*}

                \item \begin{align*}
                            M_{41}+M_{42}+M_{43}+M_{44}
                            & = -A_{41}+A_{42}-A_{43}+A_{44}
                            = \begin{vmatrix}
                                3  & 0  & 4  & 1 \\
                                2  & 3  & 1  & 4 \\
                                0  & -7 & 8  & 3 \\
                                -1 & 1  & -1 & 1
                            \end{vmatrix}                     \\
                            & = \begin{vmatrix}
                                    3  & 3  & 1  & 4 \\
                                    2  & 5  & -1 & 6 \\
                                    0  & -7 & 8  & 3 \\
                                    -1 & 0  & 0  & 0
                                \end{vmatrix}
                            = {-1}^{4+1} \cdot (-1) \begin{vmatrix}
                                                        3  & 1  & 4 \\
                                                        5  & -1 & 6 \\
                                                        -7 & 8  & 3
                                                    \end{vmatrix} \\
                            & = -78.
                        \end{align*}
            \end{enumerate}
\end{solve}

\begin{exercise}\label{prob:L10-B-3}
设$A,B$为三阶矩阵，且$|A|=3,|B|=2$，且$|A^{-1}+B|=2$，求$|A+B^{-1}|$.
\end{exercise}
\noindent 来源：\href{https://github.com/yhwu-is/Linear-Algebra-Left-Undone/blob/new/%E8%AE%B2%E4%B9%89/%E4%B8%93%E9%A2%98/10%20%E8%A1%8C%E5%88%97%E5%BC%8F.tex}{L10-B-3}。
\begin{solve}
\begin{align*}
                \lvert A+B^{-1} \rvert ={} & \lvert B^{-1}BA+B^{-1}E \rvert = \lvert B^{-1} \rvert \cdot \lvert BA+E \rvert                  \\
                ={}                        & \frac{1}{2} \lvert BA+A^{-1}A \rvert = \frac{1}{2} \lvert B+A^{-1} \rvert \cdot \lvert A \rvert \\
                ={}                        & \frac{3}{2} \lvert A^{-1}+B \rvert = 3.
            \end{align*}
\end{solve}

\begin{exercise}\label{prob:L10-B-4}
设$A$为$n$阶正交矩阵，即$AA^\mathrm{T}=A^\mathrm{T}A=E$，且$|A|<0$，证明：$|E+A|=0$.
\end{exercise}
\noindent 来源：\href{https://github.com/yhwu-is/Linear-Algebra-Left-Undone/blob/new/%E8%AE%B2%E4%B9%89/%E4%B8%93%E9%A2%98/10%20%E8%A1%8C%E5%88%97%E5%BC%8F.tex}{L10-B-4}。
\begin{solve}
正交矩阵满足 $AA^{\mathrm{T}} = A^{\mathrm{T}}A = E$，所以 $\lvert AA^{\mathrm{T}} \rvert = \lvert A \rvert^2 = \lvert E \rvert = 1$. 而 $\lvert A \rvert < 0$，所以 $\lvert A \rvert = -1$.
            \[\lvert E+A \rvert = \lvert AA^{\mathrm{T}}+A \rvert = \lvert A \rvert \cdot \lvert A^{\mathrm{T}}+E \rvert = -\lvert (A+E)^{\mathrm{T}} \rvert = -\lvert E+A \rvert.\]
            故 $\lvert E+A \rvert = 0$.

            以上两道题都多次运用了 $E = AA^{-1} = A^{-1}A$ 的技巧，请大家留意.
\end{solve}

\begin{exercise}\label{prob:L10-B-5}
设实方阵$A$的伴随矩阵$\operatorname{adj}(A)=\begin{pmatrix}
                1 & 0 & 0 \\ 0 & 1 & 0 \\ 1 & 0 & 1
            \end{pmatrix}$，且$|A|>0$，已知矩阵$B$满足$AB=E+3A$，求矩阵$B$.
\end{exercise}
\noindent 来源：\href{https://github.com/yhwu-is/Linear-Algebra-Left-Undone/blob/new/%E8%AE%B2%E4%B9%89/%E4%B8%93%E9%A2%98/10%20%E8%A1%8C%E5%88%97%E5%BC%8F.tex}{L10-B-5}。
\begin{solve}
由 $A$ 为三阶矩阵，可得 $|A|^2 = |\operatorname{adj}(A)| = 1$，而 $|A|>0$，所以 $|A|=1$. $AB=E+3A$ 经变形可得 $A(B-3E) = E$，故有
            \[
                B = 3E + A^{-1} = 3E + \frac{\operatorname{adj}(A)}{|A|} = \begin{pmatrix}
                    4 & 0 & 0 \\ 0 & 4 & 0 \\ 1 & 0 & 4
                \end{pmatrix}.
            \]
\end{solve}

\begin{exercise}\label{prob:L10-B-6}
设$B=\begin{pmatrix}
                1 & 1 & 1 \\ 0 & 1 & 1 \\ 0 & 0 & 1
            \end{pmatrix}$，矩阵$A$满足方程$BA=B^\mathrm{T}B$，求$|\operatorname{adj}(A)-2E|$.
\end{exercise}
\noindent 来源：\href{https://github.com/yhwu-is/Linear-Algebra-Left-Undone/blob/new/%E8%AE%B2%E4%B9%89/%E4%B8%93%E9%A2%98/10%20%E8%A1%8C%E5%88%97%E5%BC%8F.tex}{L10-B-6}。
\begin{solve}
$|B|=1$，对 $BA=B^\mathrm{T}B$ 两边同时取行列式有
            \[
                |B||A| = |B^\mathrm{T}||B| = |B|^2,
            \]
            故 $|A| = 1$. 从而有
            \[
                \operatorname{adj}(A) = |A|A^{-1} = B^{-1} \left(B^\mathrm{T}\right)^{-1} B.
            \]
            $B^{-1} = \begin{pmatrix}
                1 & 0 & 0 \\ -1 & 1 & 0 \\ 0 & -1 & 1
            \end{pmatrix}$，于是有
            \[
                |\operatorname{adj}(A)-2E| = \lvert B^{-1} \left(B^\mathrm{T}\right)^{-1} B - 2 B^{-1} B \rvert
                         = \lvert B^{-1} \rvert \lvert \left(B^\mathrm{T}\right)^{-1} - 2E \rvert \lvert B \rvert
                         = \begin{vmatrix}
                            -1 & 0 & 0 \\ -1 & -1 & 0 \\ 0 & -1 & -1
                         \end{vmatrix}
                         = -1.
            \]
\end{solve}

\begin{exercise}\label{prob:L10-B-7}
已知齐次线性方程组
        \[\begin{cases} \begin{aligned}
                    a_{11}x_1+a_{12}x_2+\cdots+a_{1n}x_n          & = 0             \\
                    a_{21}x_1+a_{22}x_2+\cdots+a_{2n}x_n          & = 0             \\
                                                                  & \vdotswithin{=} \\
                    a_{n-1,1}x_1+a_{n-1,2}x_2+\cdots+a_{n-1,n}x_n & = 0
                \end{aligned} \end{cases}\]
        设$M_j\enspace(j=1,2,\ldots,n)$表示$A=(a_{ij})_{n-1 \times n}$划掉第$j$列所得的$n-1$阶子式，证明：
        \begin{enumerate}
            \item $(M_1,-M_2,\ldots,(-1)^{n-1}M_n)$是方程组的一个解；

            \item 若$r(A)=n-1$，则方程组的解全是$(M_1,-M_2,\ldots,(-1)^{n-1}M_n)$的倍数.
        \end{enumerate}
\end{exercise}
\noindent 来源：\href{https://github.com/yhwu-is/Linear-Algebra-Left-Undone/blob/new/%E8%AE%B2%E4%B9%89/%E4%B8%93%E9%A2%98/10%20%E8%A1%8C%E5%88%97%E5%BC%8F.tex}{L10-B-7}。
\begin{solve}
\begin{enumerate}
                \item 由于行列式
                      \[ D = \begin{vmatrix}
                              a_{11}     & a_{12}     & a_{13}     & \cdots & a_{1n}     \\
                              a_{11}     & a_{12}     & a_{13}     & \cdots & a_{1n}     \\
                              a_{21}     & a_{22}     & a_{23}     & \cdots & a_{2n}     \\
                              \vdots     & \vdots     & \vdots     & \ddots & \vdots     \\
                              a_{n-1, 1} & a_{n-1, 2} & a_{n-1, 3} & \cdots & a_{n-1, n}
                          \end{vmatrix} = 0,\]
                      而 $M_1, -M_2, \ldots, (-1)^{n-1}M_n$ 恰是 $D$ 的第一行元素的代数余子式，所以将 $D$ 按第一行展开，可知\[a_{11}M_1+a_{12}(-M_2)+\cdots+a_{1n}(-1)^{n-1}M_n=0.\]
                      而 $D$ 的其他行元素与第一行元素的代数余子式乘积之和为 0，于是结论成立.

                \item 因为 $r(A) = n-1$，所以 $M_1, -M_2, \ldots, (-1)^{n-1}M_n$，不全为 0. 且该方程组解空间维数为 1，$M_1, -M_2, \ldots, (-1)^{n-1}M_n$ 正是该方程组的非零解，结论成立.
            \end{enumerate}
\end{solve}

\begin{exercise}\label{prob:L10-B-8}
设$A,B$均为$n$阶矩阵，且$|A|=2,|B|=1$，求$|2\operatorname{adj}(A)\operatorname{adj}(B)-A^{-1}B^{-1}|$.
\end{exercise}
\noindent 来源：\href{https://github.com/yhwu-is/Linear-Algebra-Left-Undone/blob/new/%E8%AE%B2%E4%B9%89/%E4%B8%93%E9%A2%98/10%20%E8%A1%8C%E5%88%97%E5%BC%8F.tex}{L10-B-8}。
\begin{solve}
$\operatorname{adj}(A) = \lvert A \rvert A^{-1} = 2A^{-1}, \operatorname{adj}(B) = \lvert B \rvert B^{-1} = B^{-1}$，故 \[\lvert 2\operatorname{adj}(A)\operatorname{adj}(B)-A^{-1}B^{-1} \rvert = \lvert 4A^{-1}B^{-1}-A^{-1}B^{-1} \rvert = \lvert 3A^{-1}B^{-1} \rvert = \dfrac{3^n}{\lvert A \rvert \lvert B \rvert} = \dfrac{3^n}{2}.\]
\end{solve}

\begin{exercise}\label{prob:L10-B-9}
设 $n\ge3$，$A$ 是非零实 $n$ 阶矩阵，且 $A^{\mathrm T}=\operatorname{adj}(A)$。证明：
\begin{enumerate}
\item $\det A>0$；
\item $\det A=1$；若第一行的全部元素相等，求其公共值；
\item $A$ 是正交矩阵，即 $AA^{\mathrm T}=A^{\mathrm T}A=I_n$；
\item 若 $n$ 为奇数，则 $\det(I_n-A)=0$。
\end{enumerate}
\end{exercise}
\noindent 来源：\href{https://github.com/yhwu-is/Linear-Algebra-Left-Undone/blob/new/%E8%AE%B2%E4%B9%89/%E4%B8%93%E9%A2%98/10%20%E8%A1%8C%E5%88%97%E5%BC%8F.tex}{L10-B-9}。
\begin{solve}
\begin{enumerate}
\item 由 $AA^{\mathrm T}=A\operatorname{adj}(A)=\det(A)I_n$，任取 $A$ 的一个非零行，其元素平方和等于 $\det A$，所以 $\det A>0$。实数域条件保证该平方和严格为正。
\item 对上述矩阵恒等式取行列式，得 $(\det A)^2=(\det A)^n$。因 $n\ge3$ 且 $\det A>0$，有 $\det A=1$。若第一行元素均为 $c$，由该行长度平方为 $1$ 得 $nc^2=1$，故 $c=\pm1/\sqrt n$。
\item 代回得 $AA^{\mathrm T}=I_n$，故 $A^{-1}=A^{\mathrm T}$，也有 $A^{\mathrm T}A=I_n$。
\item 因 $\det A=1$，有
\[\det(I_n-A)=\det\bigl(A(A^{\mathrm T}-I_n)\bigr)=\det(A-I_n)=(-1)^n\det(I_n-A).\]
奇数 $n$ 给出所求行列式为零。
\end{enumerate}
阶数条件不能删去：二阶矩阵 $A=2I_2$ 满足 $A^{\mathrm T}=\operatorname{adj}(A)$，但行列式为 $4$，也不是正交矩阵。
\end{solve}

\begin{exercise}\label{prob:L10-B-10}
已知$A$是一个秩为$n-1$的$n\enspace(n \geqslant 2)$阶方阵，且已知某个元素$a_{ij}$的代数余子式$A_{ij} \neq 0$，求方程组$AX=0$的基础解系.
\end{exercise}
\noindent 来源：\href{https://github.com/yhwu-is/Linear-Algebra-Left-Undone/blob/new/%E8%AE%B2%E4%B9%89/%E4%B8%93%E9%A2%98/10%20%E8%A1%8C%E5%88%97%E5%BC%8F.tex}{L10-B-10}。
\begin{solve}
$r(A) = n-1$，则 $\lvert A \rvert = 0$ 且 $AX = 0$ 的解空间的维数为 1. 而考虑 $A\operatorname{adj}(A) = \lvert A \rvert E = 0$，且 $\exists A_{ij} \neq 0$，所以 $(A_{i1}, A_{i2}, \ldots , A_{in})^{\mathrm{T}}$ 是所求的基础解系.
\end{solve}

\begin{exercise}\label{prob:L10-B-11}
设$D=|a_{ij}|_{n \times n}$，$A_{ij}$是$a_{ij}$的代数余子式. 求证：
        \[\begin{vmatrix}
                A_{11}    & A_{12}    & \cdots & A_{1,n-1}   \\
                A_{21}    & A_{22}    & \cdots & A_{2,n-1}   \\
                \vdots    & \vdots    & \ddots & \vdots      \\
                A_{n-1,1} & A_{n-1,2} & \cdots & A_{n-1,n-1}
            \end{vmatrix}=a_{nn}D^{n-2}.\]
\end{exercise}
\noindent 来源：\href{https://github.com/yhwu-is/Linear-Algebra-Left-Undone/blob/new/%E8%AE%B2%E4%B9%89/%E4%B8%93%E9%A2%98/10%20%E8%A1%8C%E5%88%97%E5%BC%8F.tex}{L10-B-11}。
\begin{solve}
设 \[A = \begin{pmatrix}
                a_{11} & a_{12} & \cdots & a_{1n} \\
                a_{21} & a_{22} & \cdots & a_{2n} \\
                \vdots & \vdots & \ddots & \vdots \\
                a_{n1} & a_{n2} & \cdots & a_{nn} \\
            \end{pmatrix},\]
            则 \[\operatorname{adj}(A) = \begin{pmatrix}
                    A_{11}     & A_{21}     & \cdots & A_{n-1, 1}   & A_{n1}     \\
                    A_{12}     & A_{22}     & \cdots & A_{n-1, 2}   & A_{n2}     \\
                    \vdots     & \vdots     & \ddots & \vdots       & \vdots     \\
                    A_{1, n-1} & A_{2, n-1} & \cdots & A_{n-1, n-1} & A_{n, n-1} \\
                    A_{1n}     & A_{2n}     & \cdots & A_{n-1, n}   & A_{nn}     \\
                \end{pmatrix}.\]
            注意到目标行列式是 $\operatorname{adj}(A)$ 中元素 $A_{nn}$ 的代数余子式，也就是 $\operatorname{adj}(\operatorname{adj}(A))$ 中 $(n, n)$ 位置的元素. 由 {例10.16(4)} $\operatorname{adj}(\operatorname{adj}(A)) = \lvert A \rvert^{n-2}A$ 可知结论成立.
\end{solve}

\begin{exercise}\label{prob:L10-B-12}
设 $a_1,a_2,\ldots,a_n$ 是 $n$ 个两两互异的实数，$f_1(x),f_2(x),\ldots,f_n(x)$ 是 $n$ 个次数不大于 $n-2$ 的实系数多项式，证明：

        \[\begin{vmatrix}
                f_1(a_1) & f_1(a_2) & \cdots & f_1(a_n) \\
                f_2(a_1) & f_2(a_2) & \cdots & f_2(a_n) \\
                \vdots   & \vdots   & \ddots & \vdots   \\
                f_n(a_1) & f_n(a_2) & \cdots & f_n(a_n)
            \end{vmatrix}=0.\]
\end{exercise}
\noindent 来源：\href{https://github.com/yhwu-is/Linear-Algebra-Left-Undone/blob/new/%E8%AE%B2%E4%B9%89/%E4%B8%93%E9%A2%98/10%20%E8%A1%8C%E5%88%97%E5%BC%8F.tex}{L10-B-12}。
\begin{solve}
设 $f_i(x) = c_0 + c_1 x + c_2 x^2 + \cdots + c_{n-2} x^{n-2} \in R[x]_{n-1}$，由于 $n-1$ 维线性空间中的 $n$ 个向量必线性相关，我们可以得到 $\left\{ f_1(x), f_2(x), \ldots, f_n(x) \right\}$ 线性相关. 即存在不全为零的 $k_1, k_2, \ldots, k_n$，使得 $\forall x \in \mathbb{R}, k_1 f_1(x) + k_2 f_2(x) + \cdots + k_n f_n(x) = 0$. 设原行列式的行向量分别为 $\alpha_1, \alpha_2, \ldots, \alpha_n$，则有
            \[
                k_1 \alpha_1 + k_2 \alpha_2 + \cdots + k_n \alpha_n = \mathbf{0}.
            \]
            原行列式的行向量组线性相关，因此其值为 $0$.
\end{solve}

\begin{exercise}\label{prob:L10-B-13}
证明：$n$维向量组$\alpha_1,\alpha_2,\ldots,\alpha_n$线性无关的充要条件是
        \[\begin{vmatrix}
                \alpha_1^\mathrm{T}\alpha_1 & \alpha_1^\mathrm{T}\alpha_2 & \cdots & \alpha_1^\mathrm{T}\alpha_n \\
                \alpha_2^\mathrm{T}\alpha_1 & \alpha_2^\mathrm{T}\alpha_2 & \cdots & \alpha_2^\mathrm{T}\alpha_n \\
                \vdots                      & \vdots                      & \ddots & \vdots                      \\
                \alpha_n^\mathrm{T}\alpha_1 & \alpha_n^\mathrm{T}\alpha_2 & \cdots & \alpha_n^\mathrm{T}\alpha_n
            \end{vmatrix}\neq 0.\]
\end{exercise}
\noindent 来源：\href{https://github.com/yhwu-is/Linear-Algebra-Left-Undone/blob/new/%E8%AE%B2%E4%B9%89/%E4%B8%93%E9%A2%98/10%20%E8%A1%8C%E5%88%97%E5%BC%8F.tex}{L10-B-13}。
\begin{solve}
令 
\[A = (\alpha_1, \alpha_2, \ldots, \alpha_n), B = \begin{pmatrix}
                \alpha_1^{\mathrm{T}}\alpha_1 & \alpha_1^{\mathrm{T}}\alpha_2 & \cdots & \alpha_1^{\mathrm{T}}\alpha_n \\
                \alpha_2^{\mathrm{T}}\alpha_1 & \alpha_2^{\mathrm{T}}\alpha_2 & \cdots & \alpha_2^{\mathrm{T}}\alpha_n \\
                \vdots                        & \vdots                        & \ddots & \vdots                        \\
                \alpha_n^{\mathrm{T}}\alpha_1 & \alpha_n^{\mathrm{T}}\alpha_2 & \cdots & \alpha_n^{\mathrm{T}}\alpha_n \\
            \end{pmatrix}\]
，显然 $B = A^{\mathrm{T}}A$，由 {11.3 秩不等式第 4 个} 有 $r(B) = r(A)$. 进而 $n$ 维向量组 $(\alpha_1, \alpha_2, \ldots, \alpha_n)$ 线性无关等价于 $r(A) = n$，也就等价于 $r(B) = n$，进而等价于 $\lvert B \rvert \neq 0$，命题得证.
\end{solve}

\begin{exercise}\label{prob:L10-B-14}
设$a_1,\ldots,a_n$为$n$个$n$维向量，证明：向量组$a_1,\ldots,a_n$线性无关的充要条件是任一个$n$维向量都可以由其线性表示（不使用线性空间维数的方式完成）.
\end{exercise}
\noindent 来源：\href{https://github.com/yhwu-is/Linear-Algebra-Left-Undone/blob/new/%E8%AE%B2%E4%B9%89/%E4%B8%93%E9%A2%98/10%20%E8%A1%8C%E5%88%97%E5%BC%8F.tex}{L10-B-14}。
\begin{solve}
$(\implies)$ 记 $A = (\alpha_1, \alpha_2, \ldots, \alpha_n), \varepsilon_i = (0, \ldots, 1, 0, \ldots, 0)^{\mathrm{T}}$(第 $i$ 个为 1)，$E = (\varepsilon_1, \varepsilon_2, \ldots, \varepsilon_n)$. 因为 $\alpha_1, \alpha_2, \ldots, \alpha_n$ 是 $n$ 维线性无关的向量组，故 $\lvert A \rvert \neq 0$，即 $A$ 可逆，$A = EA, AA^{-1} = E$,\[(\alpha_1, \alpha_2, \ldots, \alpha_n) = (\varepsilon_1, \varepsilon_2, \ldots, \varepsilon_n)A, (\varepsilon_1, \varepsilon_2, \ldots, \varepsilon_n) = (\alpha_1, \alpha_2, \ldots, \alpha_n)A^{-1}.\] 即 $\alpha_1, \alpha_2, \ldots, \alpha_n$ 与 $\varepsilon_1, \varepsilon_2, \ldots, \varepsilon_n$ 等价. $\varepsilon_1, \varepsilon_2, \ldots, \varepsilon_n$ 为 $\mathbf{R}^n$ 的一个基，能表出任一 $n$ 维向量，故 $\alpha_1, \alpha_2, \ldots, \alpha_n$ 能表出 $\mathbf{R}^n$ 中任一 $n$ 维向量.

            $(\impliedby)$ 若 $\alpha_1, \alpha_2, \ldots, \alpha_n$ 能表出任一 $n$ 维向量，则 $\varepsilon_i = \lambda_{i1}\alpha_1+\cdots+\lambda_{in}\alpha_n, i = 1, 2, \ldots, n$. 即
            \[E = (\varepsilon_1, \varepsilon_2, \ldots, \varepsilon_n) = (\alpha_1, \alpha_2, \ldots, \alpha_n)\begin{pmatrix}
                    \lambda_{11} & \lambda_{21} & \cdots & \lambda_{n1} \\
                    \lambda_{11} & \lambda_{21} & \cdots & \lambda_{n1} \\
                    \vdots       & \vdots       & \ddots & \vdots       \\
                    \lambda_{11} & \lambda_{21} & \cdots & \lambda_{n1} \\
                \end{pmatrix}.\]
            进而 $\lvert \alpha_1, \alpha_2, \ldots, \alpha_n \rvert \begin{vmatrix}
                    \lambda_{11} & \lambda_{21} & \cdots & \lambda_{n1} \\
                    \lambda_{11} & \lambda_{21} & \cdots & \lambda_{n1} \\
                    \vdots       & \vdots       & \ddots & \vdots       \\
                    \lambda_{11} & \lambda_{21} & \cdots & \lambda_{n1} \\
                \end{vmatrix} = 1$，所以 $\lvert \alpha_1, \alpha_2, \ldots, \alpha_n \rvert \neq 0$，$\alpha_1, \alpha_2, \ldots, \alpha_n$ 线性无关.
\end{solve}

\begin{exercise}\label{prob:L10-B-15}
设$s \times n\enspace(s\leqslant n)$矩阵为
        \[\begin{pmatrix}
                1      & a      & a^2    & \cdots & a^{n-1}    \\
                1      & a^2    & a^4    & \cdots & a^{2(n-1)} \\
                \vdots & \vdots & \vdots & \ddots & \vdots     \\
                1      & a^s    & a^{2s} & \cdots & a^{s(n-1)}
            \end{pmatrix}\]
        且$a^r\neq 1\enspace(0<r<n)$，求$A$的秩和它的列向量组的一个极大线性无关组.
\end{exercise}
\noindent 来源：\href{https://github.com/yhwu-is/Linear-Algebra-Left-Undone/blob/new/%E8%AE%B2%E4%B9%89/%E4%B8%93%E9%A2%98/10%20%E8%A1%8C%E5%88%97%E5%BC%8F.tex}{L10-B-15}。
\begin{solve}
$A$ 的前 $s$ 列组成的 $s$ 阶子式为Vandermonde行列式
            \[D = \begin{vmatrix}
                    1      & a      & a^2    & \cdots & a^{s-1}    \\
                    1      & a^2    & a^4    & \cdots & a^{2(s-1)} \\
                    \vdots & \vdots & \vdots & \ddots & \vdots     \\
                    1      & a^s    & a^{2s} & \cdots & a^{s(s-1)} \\
                \end{vmatrix}.\]
            由于当 $0 < r < n$ 时，$a^r \neq 1$，因此 $a, a^2, \ldots, a^s$ 两两不同，进而 $D \neq 0$，于是 $r(A) \geqslant s$. 又因为 $A$ 的行数是 $s$，所以 $r(A) \leqslant s$. 从而 $r(A) = s$.

            由于 $D \neq 0$，故 $D$ 的列向量线性无关，因此 $A$ 的前 $s$ 个列向量即为其列向量组的一个极大线性无关组.
\end{solve}

\begin{exercise}\label{prob:L10-B-16}
设$A,B,C,D \in \mathbf{F}^{n \times n}$，定义变换 $ T : \mathbf{F}^{n \times n} \to \mathbf{F}^{n \times n}$ 为
        \[ T(X) = AXB+CX+XD \]
        证明：
        \begin{enumerate}
            \item $T$为$\mathbf{F}^{n \times n}$上的线性变换；

            \item 当$C=D=O$时，$T$可逆的充要条件是$|AB| \neq 0$.
        \end{enumerate}
\end{exercise}
\noindent 来源：\href{https://github.com/yhwu-is/Linear-Algebra-Left-Undone/blob/new/%E8%AE%B2%E4%B9%89/%E4%B8%93%E9%A2%98/10%20%E8%A1%8C%E5%88%97%E5%BC%8F.tex}{L10-B-16}。
\begin{solve}
\begin{enumerate}
\item 对 $X,Y\in\mathbf F^{n\times n}$ 及 $a,b\in\mathbf F$，逐项使用矩阵乘法分配律可得 $T(aX+bY)=aT(X)+bT(Y)$，故 $T$ 为线性变换。
\item 此时 $T(X)=AXB$。若 $A,B$ 均可逆，则逆映射为 $X\mapsto A^{-1}XB^{-1}$。反之，若 $A$ 奇异，取非零列向量 $u\in\ker A$ 和非零列向量 $v$，则 $X=uv^{\mathrm T}\ne O$ 且 $AXB=O$，故 $T$ 不单射。若 $B$ 奇异，取非零 $v\in\ker B^{\mathrm T}$ 与非零 $u$，同样有 $uv^{\mathrm T}B=O$。因此 $T$ 可逆当且仅当 $A,B$ 都可逆，也即 $|AB|\ne0$。\end{enumerate}
\end{solve}

\begin{exercise}\label{prob:L10-B-17}
设$A$为$n$阶矩阵，且$r(A) < n$，又$A_{11} \neq 0$，证明：存在常数$k$，使得$(\operatorname{adj}(A))^2=k\operatorname{adj}(A)$.
\end{exercise}
\noindent 来源：\href{https://github.com/yhwu-is/Linear-Algebra-Left-Undone/blob/new/%E8%AE%B2%E4%B9%89/%E4%B8%93%E9%A2%98/10%20%E8%A1%8C%E5%88%97%E5%BC%8F.tex}{L10-B-17}。
\begin{solve}
因为 $r(A) < n, A_{11} \neq 0$，所以 $r(A) = n-1$，进而由 {例10.16 (6)} 可知 $r(\operatorname{adj}(A)) = 1$，所以有 
            \[\operatorname{adj}(A) = \begin{pmatrix} a_1 \\ a_2 \\ \vdots \\ a_n \end{pmatrix} (\lambda_1, \lambda_2, \ldots, \lambda_n).\]
            设 $(\lambda_1, \lambda_2, \ldots, \lambda_n) \begin{pmatrix} a_1 \\ a_2 \\ \vdots \\ a_n \end{pmatrix} = k$，则
            \begin{align*}
                (\operatorname{adj}(A))^2 & = \begin{pmatrix}a_1 \\ a_2 \\ \vdots \\ a_n \end{pmatrix}
                (\lambda_1, \lambda_2, \ldots, \lambda_n)
                \begin{pmatrix} a_1 \\ a_2 \\ \vdots \\ a_n \end{pmatrix}
                (\lambda_1, \lambda_2, \ldots, \lambda_n)                                     \\
                        & = \begin{pmatrix} a_1 \\ a_2 \\ \vdots \\ a_n \end{pmatrix} \cdot k
                \cdot (\lambda_1, \lambda_2, \ldots, \lambda_n) = k
                \begin{pmatrix} a_1 \\ a_2 \\ \vdots \\ a_n \end{pmatrix}
                (\lambda_1, \lambda_2, \ldots, \lambda_n) = k\operatorname{adj}(A).
            \end{align*}
\end{solve}

\begin{exercise}\label{prob:L10-B-18}
设$V$是一个$n$维实线性空间，证明：存在$V$中的一个由可列无穷多个向量组成的向量组$\{\alpha_i \mid i\in\mathbf{Z}_+\}$，使得其中任意$n$个向量组成的向量组都是$V$的一组基.
\end{exercise}
\noindent 来源：\href{https://github.com/yhwu-is/Linear-Algebra-Left-Undone/blob/new/%E8%AE%B2%E4%B9%89/%E4%B8%93%E9%A2%98/10%20%E8%A1%8C%E5%88%97%E5%BC%8F.tex}{L10-B-18}。
\begin{solve}
$\forall a_i \in \mathbf{R}, i \in \mathbf{Z}_{+}$ 满足：若 $i \neq j$，则 $a_i \neq a_j$，考虑以Vandermonde行列式的形式进行排布. 设
            \[ A = \begin{pmatrix}
                    1         & 1         & \cdots & 1         & \cdots \\
                    a_1       & a_2       & \cdots & a_k       & \cdots \\
                    a_1^2     & a_2^2     & \cdots & a_k^2     & \cdots \\
                    \vdots    & \vdots    & \ddots & \vdots    &        \\
                    a_1^{n-1} & a_2^{n-1} & \cdots & a_k^{n-1} & \cdots
                \end{pmatrix}.\]
            考虑 $A$ 的任意 $n$ 阶主子式 $D_n$，其均构成 Vandermonde 行列式，又 $i \neq j$ 有 $a_i \neq a_j$，所以值均不为 0，也就是说任意 $n$ 个向量都线性无关，其个数也恰好为 $n$，构成 $V$ 的一组基，命题得证.
\end{solve}

\begin{exercise}\label{prob:L10-B-20}
解行列式
        \begin{enumerate} 
                \item $ D_1=\begin{vmatrix}
                        1 & 2 & 3 & 4 \\2&3&4&1\\3&4&1&2\\4&1&2&3
                    \end{vmatrix} $

                \item $ D_2=\begin{vmatrix}
                        \lambda+2 & -1        & -1        & -1        \\
                        -1        & \lambda+2 & -1        & -1        \\
                        -1        & -1        & \lambda+2 & -1        \\
                        -1        & -1        & -1        & \lambda+2
                    \end{vmatrix} $

              

                \item $ D_3=\begin{vmatrix}
                        1^{2} & 2^{2} & 3^{2} & 4^{2} \\
                        2^{2} & 3^{2} & 4^{2} & 5^{2} \\
                        3^{2} & 4^{2} & 5^{2} & 6^{2} \\
                        4^{2} & 5^{2} & 6^{2} & 7^{2}
                    \end{vmatrix} $

                \item $ D_4=\begin{vmatrix}
                        3 & 2 & 0 & 0 \\
                        1 & 3 & 2 & 0 \\
                        0 & 1 & 3 & 2 \\
                        0 & 0 & 1 & 3
                    \end{vmatrix} $
             \end{enumerate}
\end{exercise}
\noindent 来源：\href{https://github.com/yhwu-is/Linear-Algebra-Left-Undone/blob/new/%E8%AE%B2%E4%B9%89/%E4%B8%93%E9%A2%98/10%20%E8%A1%8C%E5%88%97%E5%BC%8F.tex}{L10-B-20}。
\begin{solve}
\begin{enumerate}
                \item 连加法得结果为160；

                \item 连加法得结果为$(\lambda-1)(\lambda+3)^3$；

                \item 如果对数据敏感能看出线性相关，可以直接由$\alpha_4-\alpha_1=3(\alpha_3-\alpha_2)$证得其线性相关，则为0. 也可以使用滚动消去法，如\[
                          \begin{vmatrix}
                              1 & 4 & 9  & 16 \\
                              3 & 5 & 7  & 9  \\
                              5 & 7 & 9  & 11 \\
                              7 & 9 & 11 & 13
                          \end{vmatrix}=\begin{vmatrix}
                              1 & 4 & 9 & 16 \\
                              3 & 5 & 7 & 9  \\
                              2 & 2 & 2 & 2  \\
                              2 & 2 & 2 & 2
                          \end{vmatrix}=0.\]

                \item 当然是递推法的常见形式，但是只有4阶，所以可以暴算，转化为上三角行列式，如\[
                          \begin{vmatrix}
                              3 & 2           & 0            & 0             \\
                              0 & \frac{7}{3} & 2            & 0             \\
                              0 & 0           & \frac{15}{7} & 2             \\
                              0 & 0           & 0            & \frac{31}{15}
                          \end{vmatrix}=31.\]
            \end{enumerate}
\end{solve}

\begin{exercise}\label{prob:L10-B-21}
解行列式 $\begin{vmatrix}
                a_{1}+a_{2}         & a_{2}+a_{3}         & \cdots & a_{n-1}+a_n         & a_n+a_{1}         \\
                a_{1}^{2}+a_{2}^{2} & a_{2}^{2}+a_{3}^{2} & \cdots & a_{n-1}^{2}+a_n^{2} & a_n^{2}+a_{1}^{2} \\
                \vdots              & \vdots              & \ddots & \vdots              & \vdots            \\
                a_{1}^{n}+a_{2}^{n} & a_{2}^{n}+a_{3}^{n} & \cdots & a_{n-1}^{n}+a_n^{n} & a_n^{n}+a_{1}^{n}
            \end{vmatrix}$.
\end{exercise}
\noindent 来源：\href{https://github.com/yhwu-is/Linear-Algebra-Left-Undone/blob/new/%E8%AE%B2%E4%B9%89/%E4%B8%93%E9%A2%98/10%20%E8%A1%8C%E5%88%97%E5%BC%8F.tex}{L10-B-21}。
\begin{solve}
暴力拆成$2^n$个行列式，其中只有2个每列各不相同，其他都因为有相同列而为0. 因此有
            \begin{align*}
                \text{原式} & =\begin{vmatrix}
                                    a_{1}     & a_{2}     & \cdots & a_{n-1}     & a_{n}     \\
                                    a_{1}^{2} & a_{2}^{2} & \cdots & a_{n-1}^{2} & a_{n}^{2} \\
                                    \vdots    & \vdots    & \ddots & \vdots      & \vdots    \\
                                    a_{1}^{n} & a_{2}^{n} & \cdots & a_{n-1}^{n} & a_{n}^{n}
                                \end{vmatrix}+
                \begin{vmatrix}
                    a_{2}     & a_{3}     & \cdots & a_{n}     & a_{1}     \\
                    a_{2}^{2} & a_{3}^{2} & \cdots & a_{n}^{2} & a_{1}^{2} \\
                    \vdots    & \vdots    & \ddots & \vdots    & \vdots    \\
                    a_{2}^{n} & a_{3}^{n} & \cdots & a_{n}^{n} & a_{1}^{n}
                \end{vmatrix}                                               \\
                            & =(1+(-1)^{n-1})\begin{vmatrix}
                                                a_{1}     & a_{2}     & \cdots & a_{n-1}     & a_{n}     \\
                                                a_{1}^{2} & a_{2}^{2} & \cdots & a_{n-1}^{2} & a_{n}^{2} \\
                                                \vdots    & \vdots    & \ddots & \vdots      & \vdots    \\
                                                a_{1}^{n} & a_{2}^{n} & \cdots & a_{n-1}^{n} & a_{n}^{n}
                                            \end{vmatrix}                \\
                            & =(1+(-1)^{n-1})\prod_{i=1}^na_i
                \begin{vmatrix}
                    1           & 1           & \cdots & 1             & 1           \\
                    a_{1}       & a_{2}       & \cdots & a_{n-1}       & a_{n}       \\
                    \vdots      & \vdots      & \ddots & \vdots        & \vdots      \\
                    a_{1}^{n-1} & a_{2}^{n-1} & \cdots & a_{n-1}^{n-1} & a_{n}^{n-1}
                \end{vmatrix}                                     \\
                            & =(1+(-1)^{n-1})\left(\prod_{i=1}^na_i\right)\prod_{1\leqslant i<j\leqslant n}(a_j-a_i). \\
            \end{align*}
\end{solve}

\begin{exercise}\label{prob:L10-B-22}
解行列式
         \begin{enumerate}
                \item $D=\begin{vmatrix}
                              ax+by & ay+bz & az+bx \\
                              ay+bz & az+bx & ax+by \\
                              az+bx & ax+by & ay+bz
                          \end{vmatrix}$

                \item $D=\begin{vmatrix}
                              x^2+1 & xy    & xz    \\
                              xy    & y^2+1 & yz    \\
                              xz    & yz    & z^2+1
                          \end{vmatrix}$
            \end{enumerate}
\end{exercise}
\noindent 来源：\href{https://github.com/yhwu-is/Linear-Algebra-Left-Undone/blob/new/%E8%AE%B2%E4%B9%89/%E4%B8%93%E9%A2%98/10%20%E8%A1%8C%E5%88%97%E5%BC%8F.tex}{L10-B-22}。
\begin{solve}
\begin{enumerate}
\item 写成矩阵乘积：
\[M=\begin{pmatrix}a&b&0\\0&a&b\\b&0&a\end{pmatrix}\begin{pmatrix}x&y&z\\y&z&x\\z&x&y\end{pmatrix}.\]
第一个因子的行列式为 $a^3+b^3$，第二个为 $3xyz-x^3-y^3-z^3$，所以
\[D=(a^3+b^3)(3xyz-x^3-y^3-z^3).\]
这里不能把三阶行列式中的标量因子误写成一次项，也不能把立方和写成平方和。
\item 令 $u=(x,y,z)^{\mathrm T}$，则矩阵为 $I_3+uu^{\mathrm T}$。由恒等式 $\det(I+uv^{\mathrm T})=1+v^{\mathrm T}u$，得到 $D=1+x^2+y^2+z^2$。该恒等式也可由行列式按列的多线性直接证明：取两列以上的秩一部分时这些列成比例，相应行列式为零，剩下常数项和三个平方项。\end{enumerate}
\end{solve}

\begin{exercise}\label{prob:L10-B-23}
计算行列式$|2E-\alpha_1^\mathrm{T}\beta_1-\alpha_2^\mathrm{T}\beta_2|$，其中$\alpha_1=(a_1,a_2,\ldots,a_n),\enspace \beta_1=(b_1,b_2,\ldots,b_n),\enspace \alpha_2=(c_1,c_2,\ldots,c_n),\enspace \beta_2 = (d_1,d_2,\ldots,d_n)$.（提示：利用$|\lambda E_m-AB|=\lambda^{m-n}|\lambda E_n-BA|$）
\end{exercise}
\noindent 来源：\href{https://github.com/yhwu-is/Linear-Algebra-Left-Undone/blob/new/%E8%AE%B2%E4%B9%89/%E4%B8%93%E9%A2%98/10%20%E8%A1%8C%E5%88%97%E5%BC%8F.tex}{L10-B-23}。
\begin{solve}
利用$|\lambda E_m-AB|=\lambda^{m-n}|\lambda E_n-BA|$
            \begin{align*}
                |2E-\alpha_1^\mathrm{T}\beta_1-\alpha_2^\mathrm{T}\beta_2|
                & =\begin{vmatrix}
                        2E-\begin{pmatrix}
                        \alpha_1^\mathrm{T} & \alpha_2^\mathrm{T}
                    \end{pmatrix}\begin{pmatrix}
                                        \beta_1 \\\beta_2
                                    \end{pmatrix}
                    \end{vmatrix}
                =2^{n-2}\begin{vmatrix}
                            2E-\begin{pmatrix}
                        \beta_1 \\\beta_2
                    \end{pmatrix}\begin{pmatrix}
                                        \alpha_1^\mathrm{T} & \alpha_2^\mathrm{T}
                                    \end{pmatrix}
                        \end{vmatrix}         \\
                & =2^{n-2}\begin{vmatrix}
                            2-\beta_1\alpha_1^\mathrm{T} & -\beta_1\alpha_2^\mathrm{T}  \\
                            -\beta_2\alpha_1^\mathrm{T}  & 2-\beta_2\alpha_2^\mathrm{T}
                        \end{vmatrix} \\
                & =2^{n-2}
                \left(\left(2-\sum_{i=1}^na_ib_i\right)
                \left(2-\sum_{i=1}^nc_id_i\right)
                -\left(\sum_{i=1}^na_id_i\right)
                \left(\sum_{i=1}^nb_ic_i\right)\right)
            \end{align*}
\end{solve}

\begin{exercise}\label{prob:L10-B-24}
求$2n$阶行列式的值（空缺处都是0）：
        \[\begin{vmatrix}
                a &         &   &   &         & b \\
                  & \ddots  &   &   & \iddots &   \\
                  &         & a & b &         &   \\
                  &         & b & a &         &   \\
                  & \iddots &   &   & \ddots  &   \\
                b &         &   &   &         & a
            \end{vmatrix}.\]
\end{exercise}
\noindent 来源：\href{https://github.com/yhwu-is/Linear-Algebra-Left-Undone/blob/new/%E8%AE%B2%E4%B9%89/%E4%B8%93%E9%A2%98/10%20%E8%A1%8C%E5%88%97%E5%BC%8F.tex}{L10-B-24}。
\begin{solve}
设所求 $2n$ 阶行列式为 $D$.
            \begin{enumerate}
                \item 若 $a=0$，则
                    \[
                        D = \begin{vmatrix}
                            0 &         &   &   &         & b \\
                              & \ddots  &   &   &         &   \\
                              &         & 0 & b &         &   \\
                              &         & b & 0 &         &   \\
                              &         &   &   & \ddots  &   \\
                            b &         &   &   &         & 0
                        \end{vmatrix}
                        = (-1)^{\frac{2n(2n-1)}{2}} \begin{vmatrix}
                            b &         &   &   &         & 0 \\
                              & \ddots  &   &   &         &   \\
                              &         & b & 0 &         &   \\
                              &         & 0 & b &         &   \\
                              &         &   &   & \ddots  &   \\
                            0 &         &   &   &         & b
                        \end{vmatrix}
                        = b^{2n}.
                    \]
                \item 若 $a \neq 0$，我们将第 $i$ 列乘以 $-\dfrac{b}{a}$ 后加到第 $2n-i+1$ 列，化为下三角矩阵：
                    \[
                        D = \begin{vmatrix}
                            a &         &   &                   &         & 0                 \\
                              & \ddots  &   &                   &         &                   \\
                              &         & a & 0                 &         &                   \\
                              &         & b & a - \frac{b^2}{a} &         &                   \\
                              &         &   &                   & \ddots  &                   \\
                            b &         &   &                   &         & a - \frac{b^2}{a}
                        \end{vmatrix}
                        = a^n (a - \frac{b^2}{a})^n
                        = (a^2 - b^2)^n.
                    \]
            \end{enumerate}
\end{solve}

\begin{exercise}\label{prob:L10-C-3}
设 $a_i, b_i$ 都是实数，利用 Cauchy-Binet 公式证明 Cauchy-Schwarz 不等式：
        \[
            \left( \sum_{i=1}^{n}\limits a_i^2 \right) \left( \sum_{i=1}^{n}\limits b_i^2 \right) \geqslant \left( \sum_{i=1}^{n}\limits a_i b_i \right)^2.
        \]
\end{exercise}
\noindent 来源：\href{https://github.com/yhwu-is/Linear-Algebra-Left-Undone/blob/new/%E8%AE%B2%E4%B9%89/%E4%B8%93%E9%A2%98/10%20%E8%A1%8C%E5%88%97%E5%BC%8F.tex}{L10-C-3}。
\begin{solve}
设

            \[
                A = \begin{pmatrix}
                    a_1 & a_2 & \cdots & a_n \\
                    b_1 & b_2 & \cdots & b_n
                \end{pmatrix}.
            \]

            由 Cauchy-Binet 公式的推论，$A A^{\mathrm{T}}$ 的任一主子式非负，故 $|AA^\mathrm{T}|$ 非负，即

            \[
                \begin{vmatrix}
                    \sum\limits_{i=1}^{n} a_i^2 & \sum\limits_{i=1}^{n} a_ib_i \\
                    \sum\limits_{i=1}^{n} a_ib_i & \sum\limits_{i=1}^{n} b_i^2
                \end{vmatrix} = |AA^\mathrm{T}| \geqslant 0.
            \]

            将上述行列式展开即有

            \[
                \left( \sum\limits_{i=1}^{n} a_i^2 \right) \left( \sum\limits_{i=1}^{n} b_i^2 \right) - \left( \sum\limits_{i=1}^{n} a_i b_i \right)^2 \geqslant 0.
            \]

            故 Cauchy-Schwarz 不等式成立.
\end{solve}

\begin{exercise}\label{prob:L10-C-4}
设$A=(a_{ij})$是$n\enspace(n\geqslant 2)$阶整数方阵，满足对任意的$i,j$，$|A|$均可整除$a_{ij}$，证明：$|A|=\pm 1$.
\end{exercise}
\noindent 来源：\href{https://github.com/yhwu-is/Linear-Algebra-Left-Undone/blob/new/%E8%AE%B2%E4%B9%89/%E4%B8%93%E9%A2%98/10%20%E8%A1%8C%E5%88%97%E5%BC%8F.tex}{L10-C-4}。
\begin{solve}
设 $a_{ij} = |A| b_{ij}$，其中 $b_{ij}$ 为整数，则
            \[
                |A| = \begin{vmatrix}
                    |A| b_{11} & |A| b_{12} & \cdots & |A| b_{1n} \\
                    |A| b_{21} & |A| b_{22} & \cdots & |A| b_{2n} \\
                    \vdots     & \vdots     & \ddots & \vdots     \\
                    |A| b_{n1} & |A| b_{n2} & \cdots & |A| b_{nn}
                \end{vmatrix}
                = |A|^n \begin{vmatrix}
                    b_{11} & b_{12} & \cdots & b_{1n} \\
                    b_{21} & b_{22} & \cdots & b_{2n} \\
                    \vdots & \vdots & \ddots & \vdots \\
                    b_{n1} & b_{n2} & \cdots & b_{nn}
                \end{vmatrix}.
            \]

            由此可得
            \[
                |A|^{n-1} \begin{vmatrix}
                    b_{11} & b_{12} & \cdots & b_{1n} \\
                    b_{21} & b_{22} & \cdots & b_{2n} \\
                    \vdots & \vdots & \ddots & \vdots \\
                    b_{n1} & b_{n2} & \cdots & b_{nn}
                \end{vmatrix} = 1.
            \]

            由于 $a_{ij}$ 与 $b_{ij}$ 均为整数，所以 $|A|$ 与 $|b_{ij}|_{n \times n}$ 也为整数，且 $|A| \neq 0$，所以 $|A| = \pm 1$.
\end{solve}

\begin{exercise}\label{prob:L10-C-5}
设$A,B,C,D$均为$n$阶方阵，$\lvert A \rvert \neq 0$且$AC=CA$. 证明：
        \[\begin{vmatrix}
                A & B \\ C & D
            \end{vmatrix} = |AD-CB|.\]
\end{exercise}
\noindent 来源：\href{https://github.com/yhwu-is/Linear-Algebra-Left-Undone/blob/new/%E8%AE%B2%E4%B9%89/%E4%B8%93%E9%A2%98/10%20%E8%A1%8C%E5%88%97%E5%BC%8F.tex}{L10-C-5}。
\begin{solve}
对原矩阵进行分块初等变换化为上三角块矩阵后进行计算. 因为 $\lvert A \rvert \neq 0$，所以 $A$ 可逆，进而有如下变换：
            \[\begin{pmatrix}
                    E        & O \\
                    -CA^{-1} & E
                \end{pmatrix} \begin{pmatrix}
                    A & B \\
                    C & D \\
                \end{pmatrix} = \begin{pmatrix}
                    A & B          \\
                    O & D-CA^{-1}B
                \end{pmatrix},\]
            所以
            \begin{align*}
                    & \begin{vmatrix}
                            A & B \\
                            C & D \\
                        \end{vmatrix}
                = \begin{vmatrix}
                        E        & O \\
                        -CA^{-1} & E \\
                    \end{vmatrix}
                \begin{vmatrix}
                    A & B \\
                    C & D \\
                \end{vmatrix}       \\
                ={} & \begin{vmatrix}
                            A & B          \\
                            O & D-CA^{-1}B
                        \end{vmatrix}
                = \lvert A \rvert \lvert D-CA^{-1}B \rvert = \lvert AD-ACA^{-1}B \rvert.
            \end{align*} 由于 $AC = CA$，所以有 $ACA^{-1} = CAA^{-1} = C$，所以
            \[\begin{vmatrix}
                    A & B \\
                    C & D \\
                \end{vmatrix} = \lvert AD-CB \rvert.\]
\end{solve}

\begin{exercise}\label{prob:L10-C-6}
设$A$为$n$阶可逆矩阵，$\alpha,\beta$为$n$维列向量，证明：
        \[|A+\alpha\beta^{\mathrm{T}}|=|A|(1+\beta^\mathrm{T}A^{-1}\alpha).\]
\end{exercise}
\noindent 来源：\href{https://github.com/yhwu-is/Linear-Algebra-Left-Undone/blob/new/%E8%AE%B2%E4%B9%89/%E4%B8%93%E9%A2%98/10%20%E8%A1%8C%E5%88%97%E5%BC%8F.tex}{L10-C-6}。
\begin{solve}
这道题目我们利用了一个分块矩阵作为中间``桥梁''使得其通过分块初等变换之后能分别得到两个方向上的结果. 考虑矩阵 $\begin{pmatrix}
                A                   & \alpha \\
                -\beta^{\mathrm{T}} & 1
            \end{pmatrix}$，有
            \[\begin{pmatrix}
                    A                   & \alpha \\
                    -\beta^{\mathrm{T}} & 1
                \end{pmatrix} \begin{pmatrix}
                    E                  & O \\
                    \beta^{\mathrm{T}} & E
                \end{pmatrix} = \begin{pmatrix}
                    A+\alpha \beta^{\mathrm{T}} & \alpha \\
                    O                           & 1
                \end{pmatrix},\]
            所以
            \[\begin{vmatrix}
                    A                   & \alpha \\
                    -\beta^{\mathrm{T}} & 1
                \end{vmatrix} = \begin{vmatrix}
                    A                   & \alpha \\
                    -\beta^{\mathrm{T}} & 1
                \end{vmatrix} \begin{vmatrix}
                    E                  & O \\
                    \beta^{\mathrm{T}} & E
                \end{vmatrix} = \begin{vmatrix}
                    A+\alpha \beta^{\mathrm{T}} & \alpha \\
                    O                           & 1
                \end{vmatrix} = \lvert A+\alpha \beta^{\mathrm{T}} \rvert.\]
            另一方面，
            \[\begin{pmatrix}
                    E                        & O \\
                    \beta^{\mathrm{T}}A^{-1} & 1
                \end{pmatrix} \begin{pmatrix}
                    A                   & \alpha \\
                    -\beta^{\mathrm{T}} & 1
                \end{pmatrix} = \begin{pmatrix}
                    A & \alpha                           \\
                    O & 1+\beta^{\mathrm{T}}A^{-1}\alpha
                \end{pmatrix}.\]
            注意 $\beta^{\mathrm{T}}A^{-1}\alpha$ 的最终结果是一个数. 进而
            \[\begin{vmatrix}
                    A                   & \alpha \\
                    -\beta^{\mathrm{T}} & 1
                \end{vmatrix} = \begin{vmatrix}
                    E                        & O \\
                    \beta^{\mathrm{T}}A^{-1} & 1
                \end{vmatrix} \begin{vmatrix}
                    A                   & \alpha \\
                    -\beta^{\mathrm{T}} & 1
                \end{vmatrix} = \begin{vmatrix}
                    A & \alpha                           \\
                    O & 1+\beta^{\mathrm{T}}A^{-1}\alpha
                \end{vmatrix} = \lvert A \rvert(1+\beta^{\mathrm{T}}A^{-1}\alpha).\]
            所以有
            \[\lvert A+\alpha \beta^{\mathrm{T}} \rvert = \lvert A \rvert(1+\beta^{\mathrm{T}}A^{-1}\alpha).\]
\end{solve}

\begin{exercise}\label{prob:L10-C-7}
设$A,B$均为$n$阶方阵，证明：
        \[\begin{vmatrix}
                A & B \\ B & A
            \end{vmatrix} = |A+B||A-B|.\]
\end{exercise}
\noindent 来源：\href{https://github.com/yhwu-is/Linear-Algebra-Left-Undone/blob/new/%E8%AE%B2%E4%B9%89/%E4%B8%93%E9%A2%98/10%20%E8%A1%8C%E5%88%97%E5%BC%8F.tex}{L10-C-7}。
\begin{solve}
依旧是对分块矩阵做初等分块变换. 考虑到
            \begin{align*}
                \left(\begin{pmatrix}
                            E & E \\
                            O & E
                        \end{pmatrix}
                \begin{pmatrix}
                    A & B \\
                    B & A
                \end{pmatrix} \right)
                \begin{pmatrix}
                    E & -E \\
                    O & E
                \end{pmatrix}
                & = \begin{pmatrix}
                        A+B & A+B \\
                        B   & A
                    \end{pmatrix}
                \begin{pmatrix}
                    E & -E \\
                    O & E
                \end{pmatrix}      \\
                & = \begin{pmatrix}
                        A+B & O   \\
                        B   & A-B
                    \end{pmatrix},
            \end{align*}
            所以有 \begin{align*}
                \begin{vmatrix}
                    A & B \\
                    B & A
                \end{vmatrix}
                & = \begin{vmatrix}
                        E & E \\
                        O & E
                    \end{vmatrix}
                \begin{vmatrix}
                    A & B \\
                    B & A
                \end{vmatrix}
                \begin{vmatrix}
                    E & -E \\
                    O & E
                \end{vmatrix}
                = \begin{vmatrix}
                        A+B & O   \\
                        B   & A-B
                    \end{vmatrix}    \\
                & =
                \lvert (A+B)(A-B) \rvert = \lvert A+B \rvert \lvert A-B \rvert
            \end{align*}
\end{solve}

\begin{exercise}\label{prob:L10-C-8}
设$A,B,C,D$均为$n$阶方阵，且$r\begin{pmatrix}
                A & B \\ C & D
            \end{pmatrix}=n$，证明：
        \[\begin{vmatrix}
                |A| & |B| \\ |C| & |D|
            \end{vmatrix} = 0.\]
\end{exercise}
\noindent 来源：\href{https://github.com/yhwu-is/Linear-Algebra-Left-Undone/blob/new/%E8%AE%B2%E4%B9%89/%E4%B8%93%E9%A2%98/10%20%E8%A1%8C%E5%88%97%E5%BC%8F.tex}{L10-C-8}。
\begin{solve}
首先有 $\begin{vmatrix}
                \lvert A \rvert & \lvert B \rvert \\
                \lvert C \rvert & \lvert D \rvert
            \end{vmatrix} = \lvert A \rvert \lvert D \rvert-\lvert B \rvert \lvert C \rvert$.
            \begin{enumerate}
                \item 若 $\lvert A \rvert \neq 0$，即 $A$ 可逆. 因为矩阵的初等变换不改变矩阵的秩，所以由
                    \[\begin{pmatrix}
                            E        & O \\
                            -CA^{-1} & E
                        \end{pmatrix} \begin{pmatrix}
                            A & B \\
                            C & D
                        \end{pmatrix} \begin{pmatrix}
                            E & -A^{-1}B \\
                            O & E        \\
                        \end{pmatrix} = \begin{pmatrix}
                            A & O          \\
                            O & D-CA^{-1}B
                        \end{pmatrix},\]
                    条件 $r\left(\begin{pmatrix}
                                A & B \\
                                C & D
                            \end{pmatrix}\right) = n$ 以及 $A$ 可逆，可以得到
                    \[D-CA^{-1}B = O.\]
                    即若 $A$ 可逆，则 $D = CA^{-1}B$，并且
                    \[\begin{vmatrix}
                            \lvert A \rvert & \lvert B \rvert \\
                            \lvert C \rvert & \lvert D \rvert
                        \end{vmatrix} = \lvert A \rvert \lvert CA^{-1}B \rvert-\lvert B \rvert \lvert C \rvert = \lvert A \rvert \lvert C \rvert \lvert A^{-1} \rvert \lvert B \rvert-\lvert B \rvert \lvert C \rvert = 0.\]

                \item 若 $\lvert A \rvert = 0$，只需证 $\lvert B \rvert \lvert C \rvert = 0$. 若 $\lvert B \rvert \neq 0$，则由
                    \[\begin{pmatrix}
                            E        & O \\
                            -DB^{-1} & E
                        \end{pmatrix} \begin{pmatrix}
                            A & B \\
                            C & D
                        \end{pmatrix} \begin{pmatrix}
                            E        & O \\
                            -B^{-1}A & E
                        \end{pmatrix} = \begin{pmatrix}
                            O          & B \\
                            C-DB^{-1}A & O
                        \end{pmatrix},\]
                    有 $C-DB^{-1}A = O$. 注意到 $\lvert A \rvert = 0$，故 \[\lvert C \rvert = \lvert DB^{-1}A \rvert = \lvert D \rvert \lvert B^{-1} \rvert \lvert A \rvert = 0.\] 同理可证若 $\lvert C \rvert \neq 0$，则 $\lvert B \rvert = 0$.
            \end{enumerate}
            综上，结论成立.
\end{solve}

\begin{exercise}\label{prob:L10-C-9}
设$A$为$n$阶实对称矩阵（$n>1$），$|A|=0$，证明：$A_{ii}A_{jj}=(A_{ij})^2(i,j=1,\ldots,n)$.
\end{exercise}
\noindent 来源：\href{https://github.com/yhwu-is/Linear-Algebra-Left-Undone/blob/new/%E8%AE%B2%E4%B9%89/%E4%B8%93%E9%A2%98/10%20%E8%A1%8C%E5%88%97%E5%BC%8F.tex}{L10-C-9}。
\begin{solve}
由 $|A| = 0$ 可知 $r(A) < n$. 利用结论
            \[
                r(\operatorname{adj}(A))=\begin{cases}
                    n & r(A)=n     \\
                    1 & r(A)=n-1   \\
                    0 & r(A) < n-1
                \end{cases}
            \]
            可知 $r(\operatorname{adj}(A)) \leqslant 1$.

            若 $r(\operatorname{adj}(A)) = 0$，则 $\operatorname{adj}(A) = O$，故 $A_{ii} = A_{jj} = A_{ij} = 0$，自然有 $A_{ii}A_{jj}=(A_{ij})^2$；

            若 $r(\operatorname{adj}(A)) = 1$，则 $\operatorname{adj}(A)$ 的任意两个列向量线性相关，可得 $A_{ii}A_{jj} = A_{ij}A_{ji}$. 又由 $A$ 是实对称矩阵，有 $A_{ij} = A_{ji}$，因此 $A_{ii}A_{jj} = (A_{ij})^2$.
\end{solve}

\begin{exercise}\label{prob:L10-C-10}
分别求下列分块矩阵的伴随矩阵，所有块的尺寸均相容：
\begin{enumerate}
\item $M=\begin{pmatrix}A&C\\O&B\end{pmatrix}$，其中 $A,B$ 为方阵，允许奇异；
\item $N=\begin{pmatrix}A&B\\C&D\end{pmatrix}$，其中 $A,D$ 为方阵，且 $A$ 可逆。
\end{enumerate}
\end{exercise}
\noindent 来源：\href{https://github.com/yhwu-is/Linear-Algebra-Left-Undone/blob/new/%E8%AE%B2%E4%B9%89/%E4%B8%93%E9%A2%98/10%20%E8%A1%8C%E5%88%97%E5%BC%8F.tex}{L10-C-10}。
\begin{solve}
\begin{enumerate}
\item 当 $A,B$ 均可逆时，利用分块求逆及 $\det M=\det A\det B$ 得
\[\operatorname{adj}(M)=\begin{pmatrix}\det(B)\operatorname{adj}(A)&-\operatorname{adj}(A)C\operatorname{adj}(B)\\O&\det(A)\operatorname{adj}(B)\end{pmatrix}.\]
公式对奇异块仍成立：将 $A,B$ 同时换成 $A+tI,B+tI$ 后，除有限多个 $t$ 外都可逆；等式两边的每个元素都是 $t$ 的多项式，在无限多个值上相等，故恒等。取 $t=0$ 即得一般结论。
\item 记 Schur 补 $S=D-CA^{-1}B$。消元分解给出 $\det N=\det A\det S$。当 $S$ 可逆时，分块求逆后乘行列式得到
\[\operatorname{adj}(N)=\det(A)\begin{pmatrix}
\det(S)A^{-1}+A^{-1}B\operatorname{adj}(S)CA^{-1}&-A^{-1}B\operatorname{adj}(S)\\
-\operatorname{adj}(S)CA^{-1}&\operatorname{adj}(S)
\end{pmatrix}.\]
右端不含 $S^{-1}$，所以也适用于奇异的 $S$。严格地说，固定 $A,B,C$ 并将 $D$ 换成 $D+tI$，则 $S$ 换成 $S+tI$；两边再次是关于 $t$ 的多项式恒等式，令 $t=0$ 即可。题设仅要求 $A$ 可逆，无须额外假定整个 $N$ 可逆。
\end{enumerate}
\end{solve}

\begin{exercise}\label{prob:L10-C-11}
以下矩阵均在复数域上考虑，且 $n\ge2$。下面三个小问探讨伴随矩阵的反问题，即对任意给定的$n$阶方阵$B$，是否存在$n$阶方阵$A$使得$\operatorname{adj}(A)=B$.
        \begin{enumerate}
            \item 证明：若$n=2$，则存在唯一的2阶方阵$A$使得$\operatorname{adj}(A)=B$；

            \item 证明：若$n > 2$，则存在$n$阶方阵$A$使得$\operatorname{adj}(A)=B$的充要条件为$r(B) \in \{0,1,n\}$，并且
                  \begin{enumerate}
                      \item $r(B)=n$时，$A=\sqrt[n-1]{|B|}B^{-1}$；

                      \item $r(B)=1$时，$A=Q^{-1}\begin{pmatrix}
                                    0 & O \\ O & X_{n-1}
                                \end{pmatrix}P^{-1}$，且$|X_{n-1}|=|PQ|$，$B=P\begin{pmatrix}
                                    1 & O \\ O & O
                                \end{pmatrix}Q$；
                  \end{enumerate}

            \item 设$A=\begin{pmatrix}
                          1 & 1 & 1 \\ 1 & 1 & 1 \\ 1 & 1 & 1
                      \end{pmatrix}$，求矩阵$B$使得$\operatorname{adj}(B)=A$.
        \end{enumerate}
\end{exercise}
\noindent 来源：\href{https://github.com/yhwu-is/Linear-Algebra-Left-Undone/blob/new/%E8%AE%B2%E4%B9%89/%E4%B8%93%E9%A2%98/10%20%E8%A1%8C%E5%88%97%E5%BC%8F.tex}{L10-C-11}。
\begin{solve}
\begin{enumerate}
\item 二阶矩阵满足 $\operatorname{adj}(\operatorname{adj}A)=A$，因此唯一解为 $A=\operatorname{adj}B$。
\item 对 $n>2$，若 $r(A)=n$，则其伴随矩阵满秩；若 $r(A)=n-1$，则至少一个 $n-1$ 阶子式非零，且 $A\operatorname{adj}A=O$，故伴随矩阵秩为 $1$；若 $r(A)\le n-2$，则全部代数余子式为零。这证明了秩条件的必要性。
若 $B$ 满秩，所有解恰为 $A=cB^{-1}$，其中 $c^{n-1}=\det B$。必要性由 $\operatorname{adj}A=\det(A)A^{-1}$ 得到；充分性直接代入即可。在复数域上，该方程有 $n-1$ 个不同的非零根，故存在解。数域不能省略：例如实数域上偶次根未必存在。
若 $r(B)=1$，写 $B=P\operatorname{diag}(1,O)Q$，其中 $P,Q$ 可逆。取任意可逆的 $n-1$ 阶矩阵 $X$ 使 $\det X=\det(PQ)$，令 $A=Q^{-1}\operatorname{diag}(0,X)P^{-1}$。利用 $\operatorname{adj}(UV)=\operatorname{adj}(V)\operatorname{adj}(U)$ 可得
\[\operatorname{adj}A=\frac{\det X}{\det(PQ)}P\operatorname{diag}(1,O)Q=B.\]
若 $B=O$，取任何秩不超过 $n-2$ 的 $A$ 即可。三种情形都可构造所需矩阵。
\item 令 $e=(1,1,1)^{\mathrm T}$，目标是 $\operatorname{adj}B=ee^{\mathrm T}$。取
\[Q=\begin{pmatrix}1/\sqrt3&1/\sqrt2&1/\sqrt6\\1/\sqrt3&-1/\sqrt2&1/\sqrt6\\1/\sqrt3&0&-2/\sqrt6\end{pmatrix}.\]
则全部解可写为 $B=Q\operatorname{diag}(0,C)Q^{\mathrm T}$，其中 $C$ 为任意满足 $\det C=3$ 的二阶矩阵。事实上，相应伴随矩阵为 $Q\operatorname{diag}(3,0,0)Q^{\mathrm T}=ee^{\mathrm T}$；反之，由秩为 $2$ 及 $Be=0,e^{\mathrm T}B=0$ 可知在这个正交基下第一行和第一列都为零，剩余二阶块必须满足此行列式条件。例如可取 $C=\operatorname{diag}(1,3)$。\end{enumerate}
\end{solve}

\begin{exercise}\label{prob:L10-C-12}
证明$\begin{vmatrix}
                a      & c      & c      & \cdots & c      \\
                b      & a      & c      & \cdots & c      \\
                b      & b      & a      & \cdots & c      \\
                \vdots & \vdots & \vdots & \ddots & \vdots \\
                b      & b      & b      & \cdots & a
            \end{vmatrix}=\dfrac{b(a-c)^{n}-c(a-b)^{n}}{b-c}$，其中 $b \neq c$, 等式左端是 $n$ 阶行列式.
\end{exercise}
\noindent 来源：\href{https://github.com/yhwu-is/Linear-Algebra-Left-Undone/blob/new/%E8%AE%B2%E4%B9%89/%E4%B8%93%E9%A2%98/10%20%E8%A1%8C%E5%88%97%E5%BC%8F.tex}{L10-C-12}。
\begin{solve}
我们将第 $i$ 行乘 $(-1)$ 加到第 $i-1$ 行（$i$ 依次取 $n, n-1, \ldots, 2$），再对第 1 列展开,
            \[D_{n}=\begin{vmatrix}
                    a-b    & c-a    & 0      & \cdots & 0      \\
                    0      & a-b    & c-a    & \cdots & 0      \\
                    \vdots & \vdots & \vdots & \ddots & \vdots \\
                    0      & 0      & 0      & \cdots & c-a    \\
                    b      & b      & b      & \cdots & a
                \end{vmatrix}=(a-b)D_{n-1}+(-1)^{n+1}b(c-a)^{n-1},\]
            有了递推式，但是递推式中有$(-1)^{n+1}$，递推比较麻烦可能还需要处理. 如果用数归证明的话应该可以直接下手了，但是有更聪明的办法：取转置，则行列式的值不变，但是 $b$ 与 $c$ 的位置交换，由此可以写出 $D_n$ 的第二个递推式：$D_n=(a-c)D_{n-1}+(-1)^{n+1}c(b-a)^{n-1}$. 两个递推式分别乘以 $a-c$ 与 $a-b$ 后相减即有
            \[
                D_n=\dfrac{b(a-c)^{n}-c(a-b)^{n}}{b-c}.
            \]
\end{solve}

\begin{exercise}\label{prob:L10-C-13}
计算$n$阶行列式
        \[|A|=\begin{vmatrix}
                a_1b_1 & a_1b_2 & a_1b_3 & \cdots & a_1b_n \\
                a_2b_1 & a_2b_2 & a_2b_3 & \cdots & a_2b_n \\
                a_3b_1 & a_3b_2 & a_3b_3 & \cdots & a_3b_n \\
                \vdots & \vdots & \vdots & \ddots & \vdots \\
                a_nb_1 & a_nb_2 & a_nb_3 & \cdots & a_nb_n
            \end{vmatrix}.\]
\end{exercise}
\noindent 来源：\href{https://github.com/yhwu-is/Linear-Algebra-Left-Undone/blob/new/%E8%AE%B2%E4%B9%89/%E4%B8%93%E9%A2%98/10%20%E8%A1%8C%E5%88%97%E5%BC%8F.tex}{L10-C-13}。
\begin{solve}
\begin{enumerate}
                \item $n=1$ 时，显然有 $|A|=a_1b_1$.

                \item $n \geqslant 2$ 时，
                    \[
                        A = \begin{pmatrix}
                            a_1 \\ a_2 \\ \vdots \\ a_n
                        \end{pmatrix} \begin{pmatrix}
                            b_1 & b_2 & \cdots & b_n
                        \end{pmatrix},
                    \]
                    故 $r(A) = 1 < n$，$A$ 不可逆，因此 $|A| = 0$.
            \end{enumerate}
\end{solve}

\begin{exercise}\label{prob:L10-C-14}
计算$n$阶行列式
        \[|A|=\begin{vmatrix}
                1+a_1^2 & a_1a_2  & \cdots & a_1a_n  \\
                a_2a_1  & 1+a_2^2 & \cdots & a_2a_n  \\
                \vdots  & \vdots  & \ddots & \vdots  \\
                a_na_1  & a_na_2  & \cdots & 1+a_n^2
            \end{vmatrix}.\]
\end{exercise}
\noindent 来源：\href{https://github.com/yhwu-is/Linear-Algebra-Left-Undone/blob/new/%E8%AE%B2%E4%B9%89/%E4%B8%93%E9%A2%98/10%20%E8%A1%8C%E5%88%97%E5%BC%8F.tex}{L10-C-14}。
\begin{solve}
设
            \[
                D_n = \begin{vmatrix}
                    1+a_1^2 & a_1a_2  & \cdots & a_1a_n  \\
                    a_2a_1  & 1+a_2^2 & \cdots & a_2a_n  \\
                    \vdots  & \vdots  & \ddots & \vdots  \\
                    a_na_1  & a_na_2  & \cdots & 1+a_n^2
                \end{vmatrix},
            \]

            利用行列式公理化定义中的线性性，有
            \begin{align*}
                D_n &= \begin{vmatrix}
                    1+a_1^2 & a_1a_2  & \cdots & a_1a_n  \\
                    a_2a_1  & 1+a_2^2 & \cdots & a_2a_n  \\
                    \vdots  & \vdots  & \ddots & \vdots  \\
                    a_na_1  & a_na_2  & \cdots & a_n^2
                \end{vmatrix} + \begin{vmatrix}
                    1+a_1^2 & a_1a_2  & \cdots & 0       \\
                    a_2a_1  & 1+a_2^2 & \cdots & 0       \\
                    \vdots  & \vdots  & \ddots & \vdots  \\
                    a_na_1  & a_na_2  & \cdots & 1
                \end{vmatrix} \\
                &= a_n \begin{vmatrix}
                    1+a_1^2 & a_1a_2  & \cdots & a_1     \\
                    a_2a_1  & 1+a_2^2 & \cdots & a_2     \\
                    \vdots  & \vdots  & \ddots & \vdots  \\
                    a_na_1  & a_na_2  & \cdots & a_n
                \end{vmatrix} + D_{n-1} \\
                &= a_n \begin{vmatrix}
                    1      & 0      & \cdots & a_1     \\
                    0      & 1      & \cdots & a_2     \\
                    \vdots & \vdots & \ddots & \vdots  \\
                    0      & 0      & \cdots & a_n
                \end{vmatrix} \\
                &= D_{n-1} + a_n^2.
            \end{align*}
            而 $D_1 = 1 + a_1^2$, 因此
            \[
                |A| = D_n = 1 + \sum_{i=1}^n a_i^2.
            \]
\end{solve}

\begin{exercise}\label{prob:L10-C-15}
计算$n$阶行列式
        \[|A|=\begin{vmatrix}
                1+x_1  & 1+x_1^2 & \cdots & 1+x_1^n \\
                1+x_2  & 1+x_2^2 & \cdots & 1+x_2^n \\
                \vdots & \vdots  & \ddots & \vdots  \\
                1+x_n  & 1+x_n^2 & \cdots & 1+x_n^n
            \end{vmatrix}.\]
\end{exercise}
\noindent 来源：\href{https://github.com/yhwu-is/Linear-Algebra-Left-Undone/blob/new/%E8%AE%B2%E4%B9%89/%E4%B8%93%E9%A2%98/10%20%E8%A1%8C%E5%88%97%E5%BC%8F.tex}{L10-C-15}。
\begin{solve}
使用硬拆法，将 $A$ 的第 $i$ 行分解为 $(1, 1, \ldots, 1) + (x_i, x_i^2, \ldots, x_i^n)$. 考虑到拆出的行列式中包含不少于两列为 $(1, 1, \ldots, 1)$ 的项均为 $0$，因此可将 $|A|$ 化简为
            \[
                \begin{vmatrix}
                    x_1    & x_1^2  & \cdots & x_1^n  \\
                    x_2    & x_2^2  & \cdots & x_2^n  \\
                    \vdots & \vdots & \ddots & \vdots \\
                    x_n    & x_n^2  & \cdots & x_n^n
                \end{vmatrix} + \sum_{i=1}^n \begin{vmatrix}
                    x_1     & x_1^2     & \cdots & x_1^n     \\
                    \vdots  & \vdots    &        & \vdots    \\
                    x_{i-1} & x_{i-1}^2 & \cdots & x_{i-1}^n \\
                    1       & 1         & \cdots & 1         \\
                    x_{i+1} & x_{i+1}^2 & \cdots & x_{i+1}^n \\
                    \vdots  & \vdots    &        & \vdots    \\
                    x_n     & x_n^2     & \cdots & x_n^n
                \end{vmatrix}.
            \]

            上式中第一项可利用 Vandermonde 行列式化简：
            \[
                \begin{vmatrix}
                    x_1    & x_1^2  & \cdots & x_1^n  \\
                    x_2    & x_2^2  & \cdots & x_2^n  \\
                    \vdots & \vdots & \ddots & \vdots \\
                    x_n    & x_n^2  & \cdots & x_n^n
                \end{vmatrix} = \prod_{i=1}^n x_i \begin{vmatrix}
                    1      & x_1    & x_1^2  & \cdots & x_1^n  \\
                    1      & x_2    & x_2^2  & \cdots & x_2^n  \\
                    \vdots & \vdots & \vdots & \ddots & \vdots \\
                    1      & x_n    & x_n^2  & \cdots & x_n^n
                \end{vmatrix} = \left(\prod_{i=1}^n x_i\right) \left(\prod_{1 \leqslant i < j \leqslant n} (x_j - x_i)\right).
            \]

            对于第二项，
            \begin{align*}
                \begin{vmatrix}
                    x_1     & x_1^2     & \cdots & x_1^n     \\
                    \vdots  & \vdots    &        & \vdots    \\
                    x_{i-1} & x_{i-1}^2 & \cdots & x_{i-1}^n \\
                    1       & 1         & \cdots & 1         \\
                    x_{i+1} & x_{i+1}^2 & \cdots & x_{i+1}^n \\
                    \vdots  & \vdots    &        & \vdots    \\
                    x_n     & x_n^2     & \cdots & x_n^n
                \end{vmatrix}
                &= \prod_{j \neq i} x_j \begin{vmatrix}
                    1      & x_1     & \cdots & x_1^{n-1}     \\
                    \vdots & \vdots  &        & \vdots        \\
                    1      & x_{i-1} & \cdots & x_{i-1}^{n-1} \\
                    1      & 1^1     & \cdots & 1^{n-1}       \\
                    1      & x_{i+1} & \cdots & x_{i-1}^{n+1} \\
                    \vdots & \vdots  &        & \vdots        \\
                    1      & x_n     & \cdots & x_n^{n-1}
                \end{vmatrix} \\
                &= \prod_{j \neq i} x_j \cdot \prod_{1 \leqslant k < j \leqslant n} (x_j - x_k) \cdot \prod_{1 \leqslant j < i} \frac{1 - x_j}{x_i - x_j} \cdot \prod_{i < j \leqslant n} \frac{x_j - 1}{x_j - x_i}  \\
                &= \prod_{j \neq i} x_j \cdot \prod_{j \neq i} \frac{x_j - 1}{x_j - x_i} \cdot \prod_{1 \leqslant k < j \leqslant n} (x_j - x_k).
            \end{align*}

            综上可得，
            \[
                |A| = \left( 1 + \sum_{i=1}^n \frac{1}{x_i} \prod_{j \neq i} \frac{x_j - 1}{x_j - x_i} \right) \prod_{i=1}^n x_i \cdot \prod_{1 \leqslant k < j \leqslant n} (x_j - x_k).
            \]
\end{solve}

\begin{exercise}\label{prob:L10-C-16}
计算$n$阶行列式，其中$x,y,z$为任意实常数：
        \[|A|=\begin{vmatrix}
                x_1    & y      & y      & \cdots & y      \\
                z      & x_2    & y      & \cdots & y      \\
                z      & z      & x_3    & \cdots & y      \\
                \vdots & \vdots & \vdots & \ddots & \vdots \\
                z      & z      & z      & \cdots & x_n
            \end{vmatrix}.\]
\end{exercise}
\noindent 来源：\href{https://github.com/yhwu-is/Linear-Algebra-Left-Undone/blob/new/%E8%AE%B2%E4%B9%89/%E4%B8%93%E9%A2%98/10%20%E8%A1%8C%E5%88%97%E5%BC%8F.tex}{L10-C-16}。
\begin{solve}
\begin{enumerate}
                \item 若 $y = z$，
                    则
                    \[
                        |A| = \begin{vmatrix}
                            x_1    & y      & y      & \cdots & y      \\
                            y      & x_2    & y      & \cdots & y      \\
                            y      & y      & x_3    & \cdots & y      \\
                            \vdots & \vdots & \vdots & \ddots & \vdots \\
                            y      & y      & y      & \cdots & x_n
                        \end{vmatrix}.
                    \]

                    \begin{enumerate}
                        \item 若存在 $i \neq j$，使得 $x_i = x_j = y$，则 $A$ 中有两列元素全为 $y$，因此 $|A|$ 为 $0$.

                        \item 若存在唯一 $i$，使得 $x_i = y$，则
                            \begin{align*}
                                |A| &= (-1)^{i - 1 + i - 1} \begin{vmatrix}
                                    x_i    & y      & y      & \cdots & y      \\
                                    y      & x_1    & y      & \cdots & y      \\
                                    y      & y      & x_2    & \cdots & y      \\
                                    \vdots & \vdots & \vdots & \ddots & \vdots \\
                                    y      & y      & y      & \cdots & x_n
                                \end{vmatrix}
                                = \begin{vmatrix}
                                    y      & y       & y       & \cdots & y       \\
                                    0      & x_1 - y & 0       & \cdots & 0       \\
                                    0      & 0       & x_2 - y & \cdots & 0       \\
                                    \vdots & \vdots  & \vdots  & \ddots & \vdots  \\
                                    0      & 0       & 0       & \cdots & x_n - y
                                \end{vmatrix} \\
                                &= y \prod_{j \neq i} (x_j - y).
                            \end{align*}

                            \item 若不存在 $i$，使得 $x_i = y$，则
                                \begin{align*}
                                    |A| &= \begin{vmatrix}
                                        1      & y      & y      & y      & \cdots & y      \\
                                        0      & x_1    & y      & y      & \cdots & y      \\
                                        0      & y      & x_2    & y      & \cdots & y      \\
                                        0      & y      & y      & x_3    & \cdots & y      \\
                                        \vdots & \vdots & \vdots & \vdots & \ddots & \vdots \\
                                        0      & y      & y      & y      & \cdots & x_n
                                    \end{vmatrix}_{n+1} \\
                                    &= \begin{vmatrix}
                                        1      & y      & y      & y      & \cdots & y      \\
                                        -1     & x_1-y  & 0      & 0      & \cdots & 0      \\
                                        -1     & 0      & x_2-y  & 0      & \cdots & 0      \\
                                        -1     & 0      & 0      & x_3-y  & \cdots & 0      \\
                                        \vdots & \vdots & \vdots & \vdots & \ddots & \vdots \\
                                        -1     & 0      & 0      & 0      & \cdots & x_n-y
                                    \end{vmatrix}_{n+1} \\
                                    &= \begin{vmatrix}
                                        1 + \sum_{j=1}^n \frac{y}{x_j-y} & y      & y      & y      & \cdots & y      \\
                                        0                                & x_1-y  & 0      & 0      & \cdots & 0      \\
                                        0                                & 0      & x_2-y  & 0      & \cdots & 0      \\
                                        0                                & 0      & 0      & x_3-y  & \cdots & 0      \\
                                        \vdots                           & \vdots & \vdots & \vdots & \ddots & \vdots \\
                                        0                                & 0      & 0      & 0      & \cdots & x_n-y
                                    \end{vmatrix} \\
                                    &= \prod_{j=1}^n (x_j - y) + y\sum_{i=1}^n \prod_{j \neq i} (x_j - y).
                                \end{align*}
                    \end{enumerate}
                    综上，当 $y = z$ 时，
                    \[
                        |A| = \prod_{j=1}^n (x_j - y) + y\sum_{i=1}^n \prod_{j \neq i} (x_j - y).
                    \]

                \item 若 $y \neq z$，则可按照与第 12 题完全相同的方法进行处理，最后得到
                    \[
                        |A| = \frac{z \prod_{i=1}^n (x_i - y) - y \prod_{i=1}^n (x_i - z)}{z - y}.
                    \]
            \end{enumerate}
\end{solve}

\begin{exercise}\label{prob:L10-例-2}
计算行列式 $D = \begin{vmatrix}
        1 & 2 \\
        2 & 3
    \end{vmatrix}$.
\end{exercise}
\noindent 来源：\href{https://github.com/yhwu-is/Linear-Algebra-Left-Undone/blob/new/%E8%AE%B2%E4%B9%89/%E4%B8%93%E9%A2%98/10%20%E8%A1%8C%E5%88%97%E5%BC%8F.tex}{L10-例-2}。
\begin{solve}
设 $\alpha_1 = (1, 2)$, $\alpha_2 = (2, 3)$，将行列式 $D$ 写成向量形式并用公理化定义展开可得：
    \begin{align*}
        D &= D(\alpha_1, \alpha_2) = D(e_1 + 2e_2, 2e_1 + 3e_2) \\
          &= D(e_1, 2e_1 + 3e_2) + 2 D(e_2, 2e_1 + 3e_2) \\
          &= 2D(e_1, e_1) + 3D(e_1, e_2) + 4D(e_2, e_1) + 6D(e_2, e_2) \\
          &= 3 - 4 = -1.
    \end{align*}
\end{solve}

\begin{exercise}\label{prob:L10-例-3}
设向量$\alpha_1,\alpha_2,\beta_1,\beta_2$为三维列向量，又$A=(\alpha_1,\alpha_2,\beta_1),B=(\alpha_1,\alpha_2,\beta_2)$，且$|A|=3$，$|B|=2$，求$|2A+3B|$.
\end{exercise}
\noindent 来源：\href{https://github.com/yhwu-is/Linear-Algebra-Left-Undone/blob/new/%E8%AE%B2%E4%B9%89/%E4%B8%93%E9%A2%98/10%20%E8%A1%8C%E5%88%97%E5%BC%8F.tex}{L10-例-3}。
\begin{solve}
$2A+3B=(2\alpha_1+3\alpha_1,2\alpha_2+3\alpha_2,2\beta_1+3\beta_2)=(5\alpha_1,5\alpha_2,2\beta_1+3\beta_2)$，因此
    \begin{align*}
        |2A+3B| & =D(5\alpha_1,5\alpha_2,2\beta_1+3\beta_2)                        \\
                & =25D(\alpha_1,\alpha_2,2\beta_1+3\beta_2)                        \\
                & =25(2D(\alpha_1,\alpha_2,\beta_1)+3D(\alpha_1,\alpha_2,\beta_2)) \\
                & =25(2|A|+3|B|)=300.
    \end{align*}
\end{solve}

\begin{exercise}\label{prob:L10-例-5}
求排列 $(3,1,4,2)$ 的逆序数.
\end{exercise}
\noindent 来源：\href{https://github.com/yhwu-is/Linear-Algebra-Left-Undone/blob/new/%E8%AE%B2%E4%B9%89/%E4%B8%93%E9%A2%98/10%20%E8%A1%8C%E5%88%97%E5%BC%8F.tex}{L10-例-5}。
\begin{solve}
我们可以直接计算：$3$ 后面有 $2$ 个数比 $3$ 小，$1$ 后面没有数比 $1$ 小，$4$ 后面有 $1$ 个数比 $4$ 小，$2$ 后面没有数比 $2$ 小，因此逆序数为 $2+0+1+0=3$.
\end{solve}

\begin{exercise}\label{prob:L10-例-11}
求证以下命题：
    \begin{enumerate}
        \item 奇数阶反对称矩阵不可逆；

        \item 若$A$是$n$阶可逆对称矩阵，$B$是$n$阶反对称矩阵，则当$n$为奇数时，齐次线性方程组$(AB)X=O$有非零解.
    \end{enumerate}
\end{exercise}
\noindent 来源：\href{https://github.com/yhwu-is/Linear-Algebra-Left-Undone/blob/new/%E8%AE%B2%E4%B9%89/%E4%B8%93%E9%A2%98/10%20%E8%A1%8C%E5%88%97%E5%BC%8F.tex}{L10-例-11}。
\begin{solve}
\begin{enumerate}
        \item 设$A$为$n$阶反对称矩阵，则$A^\mathrm{T}=-A$，因此$|A|=|A^\mathrm{T}|=|-A|=(-1)^n|A|=-|A|$，从而$|A|=0$，故$A$不可逆.
        \item 反证法，则$(AB)X=O$只有零解，由原讲义“可逆等价条件”处 的 原讲义“item:11:可逆等价条件:4”处 可知$AB$可逆，故$|AB|=|A||B|\neq 0$，这说明$|B|\neq 0$，即$B$可逆. 但由于$B$为奇数阶反对称矩阵，因此$B$不可逆，矛盾. 因此原方程组有非零解.
    \end{enumerate}
\end{solve}

\begin{exercise}\label{prob:L10-例-12}
设$A$为$n$阶正交矩阵，即$AA^\mathrm{T}=A^\mathrm{T}A=E$，且$|A|<0$，证明：$|E+A|=0$.
\end{exercise}
\noindent 来源：\href{https://github.com/yhwu-is/Linear-Algebra-Left-Undone/blob/new/%E8%AE%B2%E4%B9%89/%E4%B8%93%E9%A2%98/10%20%E8%A1%8C%E5%88%97%E5%BC%8F.tex}{L10-例-12}。
\begin{solve}
正交矩阵满足 $AA^{\mathrm{T}} = A^{\mathrm{T}}A = E$，所以 $\lvert AA^{\mathrm{T}} \rvert = \lvert A \rvert^2 = \lvert E \rvert = 1$. 而 $\lvert A \rvert < 0$，所以 $\lvert A \rvert = -1$.
    \[
        \lvert E+A \rvert = \lvert AA^{\mathrm{T}}+A \rvert = \lvert A \rvert \cdot \lvert A^{\mathrm{T}}+E \rvert = -\lvert (A+E)^{\mathrm{T}} \rvert = -\lvert E+A \rvert.
    \]

    故 $\lvert E+A \rvert = 0$.
\end{solve}

\begin{exercise}\label{prob:L10-例-17}
判断矩阵$\begin{pmatrix}
            1 & 2 & 3 \\ 2 & 1 & 2 \\ 1 & 3 & 3
        \end{pmatrix}$是否可逆. 若可逆，利用伴随矩阵求其逆矩阵.
\end{exercise}
\noindent 来源：\href{https://github.com/yhwu-is/Linear-Algebra-Left-Undone/blob/new/%E8%AE%B2%E4%B9%89/%E4%B8%93%E9%A2%98/10%20%E8%A1%8C%E5%88%97%E5%BC%8F.tex}{L10-例-17}。
\begin{solve}
我们使用行列式来判断矩阵是否可逆：设 $A=\begin{pmatrix}
        1 & 2 & 3 \\ 2 & 1 & 2 \\ 1 & 3 & 3
    \end{pmatrix}$，则有 $|A| = 4$，因此 $A$ 可逆.

    接着我们使用伴随矩阵求 $A$ 的逆矩阵. 首先求 $|A|$ 中各元素的代数余子式：

    \[
        A_{11} = \begin{vmatrix}
            1 & 2 \\ 3 & 3
        \end{vmatrix} = -3, \quad
        A_{12} = - \begin{vmatrix}
            2 & 2 \\ 1 & 3
        \end{vmatrix} = -4, \quad
        A_{13} = \begin{vmatrix}
            2 & 1 \\ 1 & 3
        \end{vmatrix} = 5.
    \]

    类似地，我们有
    \begin{gather*}
        A_{21} = 3, \quad A_{22} = 0, \quad A_{23} = -1, \\
        A_{31} = 1, \quad A_{32} = 4, \quad A_{33} = -3,
    \end{gather*}

    由此可得 $A$ 的伴随矩阵 
\[\operatorname{adj}(A) = \begin{pmatrix}
            A_{11} & A_{21} & A_{31} \\
            A_{12} & A_{22} & A_{32} \\
            A_{13} & A_{23} & A_{33}
        \end{pmatrix} = \begin{pmatrix}
            -3 & 3 & 1 \\
            -4 & 0 & 4 \\
            5 & -1 & -3
        \end{pmatrix}\]
.

    因此，

    \[
        A^{-1} = \dfrac{\operatorname{adj}(A)}{|A|} = \dfrac{1}{4} \begin{pmatrix}
            -3 & 3 & 1 \\
            -4 & 0 & 4 \\
            5 & -1 & -3
            \end{pmatrix}.
    \]
\end{solve}

\begin{exercise}\label{prob:L10-例-18}
已知 $\operatorname{adj}(A)=\begin{pmatrix}
            1 & -2 & 1 \\ 0 & 2 & -2 \\ -1 & 2 & 1
        \end{pmatrix}$，求 $A$.
\end{exercise}
\noindent 来源：\href{https://github.com/yhwu-is/Linear-Algebra-Left-Undone/blob/new/%E8%AE%B2%E4%B9%89/%E4%B8%93%E9%A2%98/10%20%E8%A1%8C%E5%88%97%E5%BC%8F.tex}{L10-例-18}。
\begin{solve}
不难求得 $|\operatorname{adj}(A)| = 4$，根据伴随矩阵行列式的性质有，$|\operatorname{adj}(A)| = |A|^2$，因此 $|A| = \pm 2$. 若 $|A| = 2$，则根据 $A\operatorname{adj}(A) = |A|E$，有

    \[A^{-1} = |A|^{-1}\operatorname{adj}(A) = \frac{1}{2}\begin{pmatrix}
            1 & -2 & 1 \\ 0 & 2 & -2 \\ -1 & 2 & 1
        \end{pmatrix} = \begin{pmatrix}
            \dfrac{1}{2} & -1 & \dfrac{1}{2} \\ 0 & 1 & -1 \\ -\dfrac{1}{2} & 1 & \dfrac{1}{2}
        \end{pmatrix}.\]

    因此 $A = (A^{-1})^{-1} = \begin{pmatrix}
            3 & 2 & 1 \\ 1 & 1 & 1 \\ 1 & 0 & 1
        \end{pmatrix}$.

    若 $|A| = -2$，则同理可以求得 $A = -\begin{pmatrix}
            3 & 2 & 1 \\ 1 & 1 & 1 \\ 1 & 0 & 1
        \end{pmatrix}$.
\end{solve}

\begin{exercise}\label{prob:L10-例-19}
写出矩阵$\begin{pmatrix}
            1 & 5 & -2 \\ 2 & 3 & 4 \\ -1 & -3 & 0
        \end{pmatrix}$的所有一阶、二阶子式、主子式和顺序主子式.
\end{exercise}
\noindent 来源：\href{https://github.com/yhwu-is/Linear-Algebra-Left-Undone/blob/new/%E8%AE%B2%E4%B9%89/%E4%B8%93%E9%A2%98/10%20%E8%A1%8C%E5%88%97%E5%BC%8F.tex}{L10-例-19}。
\begin{solve}
\begin{enumerate}
        \item 一阶子式：根据子式定义，任意1行1列交点组成的1个元素就是一阶子式，即所有的元素本身都是一阶子式；

        \item 二阶子式：根据子式定义，任意2行2列交点组成的4个元素排成的行列式就是二阶子式，即
              \[\begin{vmatrix}
                      1 & 5 \\
                      2 & 3
                  \end{vmatrix},\begin{vmatrix}
                      1 & -2 \\
                      2 & 4
                  \end{vmatrix},\begin{vmatrix}
                      5 & -2 \\
                      3 & 4
                  \end{vmatrix},\begin{vmatrix}
                      1  & 5  \\
                      -1 & -3
                  \end{vmatrix},\begin{vmatrix}
                      1  & -2 \\
                      -1 & 0
                  \end{vmatrix},\]\[\begin{vmatrix}
                      5  & -2 \\
                      -3 & 0
                  \end{vmatrix},\begin{vmatrix}
                      2  & 3  \\
                      -1 & -3
                  \end{vmatrix},\begin{vmatrix}
                      2  & 4 \\
                      -1 & 0
                  \end{vmatrix},\begin{vmatrix}
                      3  & 4 \\
                      -3 & 0
                  \end{vmatrix};\]

        \item 主子式：根据主子式定义，要求取的行列号相同，故一阶主子式就是1行1列、2行2列、3行3列的元素，二阶主子式就是选取1、2/1、3/2、3行与列构成的子式，即
              \[\begin{vmatrix}
                      1 & 5 \\ 2 & 3
                  \end{vmatrix},\begin{vmatrix}
                      1 & -2 \\ -1 & 0
                  \end{vmatrix},\begin{vmatrix}
                      3 & 4 \\ -3 & 0
                  \end{vmatrix};\]
              三阶主子式就是矩阵本身对应的行列式，不再赘述.

        \item 顺序主子式根据定义，一阶就是第一行第一列的元素，二阶就是前两行前两节元素构成的子式，三阶就是本身的行列式.
    \end{enumerate}
\end{solve}

\begin{exercise}\label{prob:L10-例-20}
利用定义求矩阵$\begin{pmatrix}
            1 & 1 & -1 & 3 \\ 1 & 2 & 1 & 1 \\ 3 & 3 & 0 & 4
        \end{pmatrix}$的行列式秩.
\end{exercise}
\noindent 来源：\href{https://github.com/yhwu-is/Linear-Algebra-Left-Undone/blob/new/%E8%AE%B2%E4%B9%89/%E4%B8%93%E9%A2%98/10%20%E8%A1%8C%E5%88%97%E5%BC%8F.tex}{L10-例-20}。
\begin{solve}
记该矩阵为$A$，由于$A$为3行4列矩阵，因此$r(A)\leqslant 3$. 又我们可以发现其三阶子式
    \[\begin{vmatrix}
            1 & -1 & 3 \\ 1 & 1 & 1 \\ 3 & 0 & 4
        \end{vmatrix}=-4\neq 0,\]
    故$r(A)\geqslant 3$，因此$r(A)=3$.
\end{solve}

\begin{exercise}\label{prob:L10-例-21}
求多项式$f(x)=\begin{vmatrix}
            1 & a_1   & a_2     & a_3     \\
            1 & a_1+x & a_2     & a_3     \\
            1 & a_1   & a_2+x+1 & a_3     \\
            1 & a_1   & a_2     & a_3+x+2
        \end{vmatrix}$的所有零点.
\end{exercise}
\noindent 来源：\href{https://github.com/yhwu-is/Linear-Algebra-Left-Undone/blob/new/%E8%AE%B2%E4%B9%89/%E4%B8%93%E9%A2%98/10%20%E8%A1%8C%E5%88%97%E5%BC%8F.tex}{L10-例-21}。
\begin{solve}
事实上，本题可以直接展开求出四阶行列式的值然后解方程$f(x)=0$即可，但我们这里使用更为本质的方法. $f(x)=0$实际上就是行列式等于零，即此时$x$使得行列式中两列（或两行）线性相关了. 事实上，我们很容易发现$x=0$时，第一列与第二列成比例，故此时$f(x)=0$成立. 同理，$x=-1$和$x=-2$也是$f(x)=0$的解.

    事实上，我们知道这一行列式展开后是一个次数最高为3的多项式，因此$0,-1,-2$就是$f(x)$的所有零点.
\end{solve}

\begin{exercise}\label{prob:L10-例-22}
计算行列式$D_{n+1}=\begin{vmatrix}
            1      & a_{1}       & a_{2}       & \cdots & a_n       \\
            1      & a_{1}+b_{1} & a_{2}       & \cdots & a_n       \\
            1      & a_{1}       & a_{2}+b_{2} & \cdots & a_n       \\
            \vdots & \vdots      & \vdots      & \ddots & \vdots    \\
            1      & a_{1}       & a_{2}       & \cdots & a_n+b_{n}
        \end{vmatrix}$.
\end{exercise}
\noindent 来源：\href{https://github.com/yhwu-is/Linear-Algebra-Left-Undone/blob/new/%E8%AE%B2%E4%B9%89/%E4%B8%93%E9%A2%98/10%20%E8%A1%8C%E5%88%97%E5%BC%8F.tex}{L10-例-22}。
\begin{solve}
观察这个行列式的特点，主对角线下方的元素与第 1 行元素对应相同，故用第 1 行的 $-1$ 倍加到下面各行便可使主对角线下方的元素全部变为0，即化为上三角形.
    \[ D_{n+1}=\begin{vmatrix}
            1 & a_{1} & a_{2} & \ldots & a_n   \\
              & b_{1} &       &        &       \\
              &       & b_{2} &        &       \\
              &       &       & \ddots &       \\
              &       &       &        & b_{n}
        \end{vmatrix}=\prod_{i=1}^n b_i \]
\end{solve}

\begin{exercise}\label{prob:L10-例-23}
计算$n$阶行列式，其中$a_i\neq 0\enspace(i=2,\ldots,n)$：
    \[|A|=\begin{vmatrix}
            a_1    & b_2    & b_3    & \cdots & b_n    \\
            c_2    & a_2    & 0      & \cdots & 0      \\
            c_3    & 0      & a_3    & \cdots & 0      \\
            \vdots & \vdots & \vdots & \ddots & \vdots \\
            c_n    & 0      & 0      & \cdots & a_n
        \end{vmatrix}.\]
\end{exercise}
\noindent 来源：\href{https://github.com/yhwu-is/Linear-Algebra-Left-Undone/blob/new/%E8%AE%B2%E4%B9%89/%E4%B8%93%E9%A2%98/10%20%E8%A1%8C%E5%88%97%E5%BC%8F.tex}{L10-例-23}。
\begin{solve}
这一行列式因为形似箭头而被称为箭形行列式，将第$i\enspace(2\leqslant i\leqslant n)$列乘以$-c_i/a_i$加到第一列上可以消去第一列所有的$c_i$得到上三角矩阵，最后答案为
    \[
    (a_1-\sum\limits_{i=2}^n\dfrac{b_ic_i}{a_i})a_2a_3\cdots a_n.
    \]
\end{solve}

\begin{exercise}\label{prob:L10-例-24}
计算$n$阶行列式，其中$a_i\neq 0,\enspace i=1,\ldots,n$：
    \[|A|=\begin{vmatrix}
            x_1-a_1 & x_2     & x_3     & \cdots & x_n     \\
            x_1     & x_2-a_2 & x_3     & \cdots & x_n     \\
            x_1     & x_2     & x_3-a_3 & \cdots & x_n     \\
            \vdots  & \vdots  & \vdots  & \ddots & \vdots  \\
            x_1     & x_2     & x_3     & \cdots & x_n-a_n
        \end{vmatrix}.\]
\end{exercise}
\noindent 来源：\href{https://github.com/yhwu-is/Linear-Algebra-Left-Undone/blob/new/%E8%AE%B2%E4%B9%89/%E4%B8%93%E9%A2%98/10%20%E8%A1%8C%E5%88%97%E5%BC%8F.tex}{L10-例-24}。
\begin{solve}
第一行乘以-1依次加到其它行上去可以得到一个箭形行列式，使用上一例题的结论可知答案为\[(-1)^{n-1}a_1a_2\cdots a_n(\sum\limits_{i=1}^n\dfrac{x_i}{a_i}-1).\]
\end{solve}

\begin{exercise}\label{prob:L10-例-26}
计算 $n$ 阶行列式：
    \[
        |A| = \begin{vmatrix}
            \dfrac{1}{x_1 - 1} & 1      & x_1    & x_1^2  & \cdots & x_1^{n-2} \\[1.2em]
            \dfrac{1}{x_2 - 1} & 1      & x_2    & x_2^2  & \cdots & x_2^{n-2} \\[1.2em]
            \vdots             & \vdots & \vdots & \vdots & \ddots & \vdots    \\[1.2em]
            \dfrac{1}{x_n - 1} & 1      & x_n    & x_n^2  & \cdots & x_n^{n-2}
        \end{vmatrix}
    \]
\end{exercise}
\noindent 来源：\href{https://github.com/yhwu-is/Linear-Algebra-Left-Undone/blob/new/%E8%AE%B2%E4%B9%89/%E4%B8%93%E9%A2%98/10%20%E8%A1%8C%E5%88%97%E5%BC%8F.tex}{L10-例-26}。
\begin{solve}
对第 $i$ 行提出因子 $\dfrac{1}{x_i-1}$，则有

    \[
        |A| = \left(\displaystyle\prod_{i=1}^n \frac{1}{x_i-1} \right) \begin{vmatrix}
            1 & x_1 - 1 & x_1^2 - x_1 & \cdots & x_1^{n-1} - x_1^{n-2} \\
            1 & x_2 - 1 & x_2^2 - x_2 & \cdots & x_2^{n-1} - x_2^{n-2} \\
            1 & x_3 - 1 & x_3^2 - x_3 & \cdots & x_3^{n-1} - x_3^{n-2} \\
            \vdots & \vdots & \vdots & \ddots & \vdots \\
            1 & x_n - 1 & x_n^2 - x_n & \cdots & x_n^{n-1} - x_n^{n-2}
        \end{vmatrix}.
    \]

    通过初等列变换将 $A$ 化为范德蒙行列式的形式：

    \begin{align*}
        |A| &= \left(\displaystyle\prod_{i=1}^n \frac{1}{x_i-1} \right) \begin{vmatrix}
            1 & x_1 & x_1^2 & \cdots & x_1^{n-1} \\
            1 & x_2 & x_2^2 & \cdots & x_2^{n-1} \\
            1 & x_3 & x_3^2 & \cdots & x_3^{n-1} \\
            \vdots & \vdots & \vdots & \ddots & \vdots \\
            1 & x_n & x_n^2 & \cdots & x_n^{n-1}
        \end{vmatrix} \\
        &= \left(\displaystyle\prod_{i=1}^n \frac{1}{x_i-1} \right)
           \left(\displaystyle\prod_{1 \leqslant i < j \leqslant n} (x_j - x_i)\right).
    \end{align*}
\end{solve}

\begin{exercise}\label{prob:L10-例-29}
计算行列式$D_n=\begin{vmatrix}
            1      & 2      & 3      & \cdots & n-1    & n      \\
            2      & 1      & 2      & \cdots & n-2    & n-1    \\
            3      & 2      & 1      & \cdots & n-3    & n-2    \\
            \vdots & \vdots & \vdots & \ddots & \vdots & \vdots \\
            n-1    & n-2    & n-3    & \cdots & 1      & 2      \\
            n      & n-1    & n-2    & \cdots & 2      & 1
        \end{vmatrix},\enspace n \geqslant 2$.
\end{exercise}
\noindent 来源：\href{https://github.com/yhwu-is/Linear-Algebra-Left-Undone/blob/new/%E8%AE%B2%E4%B9%89/%E4%B8%93%E9%A2%98/10%20%E8%A1%8C%E5%88%97%E5%BC%8F.tex}{L10-例-29}。
\begin{solve}
这是一个循环行列式，这种行列式的特点是相邻的行或列之间有共同的变化规律，所以我们可以使用``差分''的思想，从最后一行开始每行减去上一行
    \begin{align*}
        D_n & =\begin{vmatrix}
                   1      & 2      & 3      & \cdots & n-1    & n      \\
                   1      & -1     & -1     & \cdots & -1     & -1     \\
                   1      & 1      & -1     & \cdots & -1     & -1     \\
                   \vdots & \vdots & \vdots & \ddots & \vdots & \vdots \\
                   1      & 1      & 1      & \cdots & -1     & -1     \\
                   1      & 1      & 1      & \cdots & 1      & -1
               \end{vmatrix}=\begin{vmatrix}
                                 1      & 2      & 3      & \cdots & n-1    & n      \\
                                 2      & 0      & 0      & \cdots & 0      & -2     \\
                                 2      & 2      & 0      & \cdots & 0      & -2     \\
                                 \vdots & \vdots & \vdots & \ddots & \vdots & \vdots \\
                                 2      & 2      & 2      & \cdots & 0      & -2     \\
                                 1      & 1      & 1      & \cdots & 1      & -1
                             \end{vmatrix} \\
            & =\begin{vmatrix}
                   1      & 2      & 3      & \cdots & n-1    & n+1    \\
                   2      & 0      & 0      & \cdots & 0      & 0      \\
                   2      & 2      & 0      & \cdots & 0      & 0      \\
                   \vdots & \vdots & \vdots & \ddots & \vdots & \vdots \\
                   2      & 2      & 2      & \cdots & 0      & 0      \\
                   1      & 1      & 1      & \cdots & 1      & 0
               \end{vmatrix}=2^{n-2}
        \begin{vmatrix}
            1      & 2      & 3      & \cdots & n-1    & n+1    \\
            1      & 0      & 0      & \cdots & 0      & 0      \\
            1      & 1      & 0      & \cdots & 0      & 0      \\
            \vdots & \vdots & \vdots & \ddots & \vdots & \vdots \\
            1      & 1      & 1      & \cdots & 0      & 0      \\
            1      & 1      & 1      & \cdots & 1      & 0
        \end{vmatrix}                      \\
            & =(-1)^{n+1}(n+1) 2^{n-2}.
    \end{align*}
\end{solve}

\begin{exercise}\label{prob:L10-例-30}
解行列式 $D=\begin{vmatrix}
            0      & 1      & 1      & \cdots & 1      & 1      \\
            1      & 0      & 1      & \cdots & 1      & 1      \\
            1      & 1      & 0      & \cdots & 1      & 1      \\
            \vdots & \vdots & \vdots & \ddots & \vdots & \vdots \\
            1      & 1      & 1      & \cdots & 0      & 1      \\
            1      & 1      & 1      & \cdots & 1      & 0
        \end{vmatrix}$.
\end{exercise}
\noindent 来源：\href{https://github.com/yhwu-is/Linear-Algebra-Left-Undone/blob/new/%E8%AE%B2%E4%B9%89/%E4%B8%93%E9%A2%98/10%20%E8%A1%8C%E5%88%97%E5%BC%8F.tex}{L10-例-30}。
\begin{solve}
使行列式 $D$ 变成 $n+1$ 阶行列式, 即
    \[ D=\begin{vmatrix}
            1      & 1      & 1      & \cdots & 1      & 1      \\
            0      & 0      & 1      & \cdots & 1      & 1      \\
            0      & 1      & 0      & \cdots & 1      & 1      \\
            \vdots & \vdots & \vdots & \ddots & \vdots & \vdots \\
            0      & 1      & 1      & \cdots & 0      & 1      \\
            0      & 1      & 1      & \cdots & 1      & 0
        \end{vmatrix}. \]

    再将第一行的 $-1$ 倍加到其他各行, 得:
    \[ D =\begin{vmatrix}
            1      & 1  & \cdots & 1  \\
            -1     & -1 &        &    \\
            \vdots &    & \ddots &    \\
            -1     &    &        & -1
        \end{vmatrix}. \]

    从第二列开始, 每列乘以 $-1$ 加到第一列, 得:
    \begin{align*}
        D & =\begin{vmatrix}
                 -(n-1) & 1  & \cdots & 1  \\
                        & -1 &        &    \\
                        &    & \ddots &    \\
                        &    &        & -1
             \end{vmatrix} \\
          & =(-1)^{n+1}(n-1).
    \end{align*}
\end{solve}

\begin{exercise}\label{prob:L10-例-31}
求行列式 $D_n=\begin{vmatrix}
            1         & 1         & \cdots & 1         \\
            x_1       & x_2       & \cdots & x_n       \\
            x_1^{2}   & x_2^{2}   & \cdots & x_n^{2}   \\
            \vdots    & \vdots    & \ddots & \vdots    \\
            x_1^{n-2} & x_2^{n-2} & \cdots & x_n^{n-2} \\
            x_1^{n}   & x_2^{n}   & \cdots & x_n^{n}
        \end{vmatrix}$.
\end{exercise}
\noindent 来源：\href{https://github.com/yhwu-is/Linear-Algebra-Left-Undone/blob/new/%E8%AE%B2%E4%B9%89/%E4%B8%93%E9%A2%98/10%20%E8%A1%8C%E5%88%97%E5%BC%8F.tex}{L10-例-31}。
\begin{solve}
观察到这个行列式与范德蒙行列式十分类似，但是把 $n-1$ 次幂换成了 $n$ 次幂，因此我们可以考虑将其用加边法补成一个 $n+1$ 阶的范德蒙行列式. 普通的加边法套路是加一列只有一个元素非零的列，同时在非零元素所在行补上任意元素以化简. 如果我们在这道题中使用此方法，那么我们需要在倒数第二行补 $x_i^{n-1}$，这样的话就必须补上一个 $n+1$ 维列向量 $(0, 0, \ldots, 0, 1, 0)^{\mathrm{T}}$，无法构成范德蒙行列式.

    

    因此，我们必须另辟蹊径：假设我们已经构造出了一个 $n+1$ 阶的范德蒙行列式 $V_{n+1}$，加的行在倒数第二行，加的列在最后一列，那么我们其实只需要求这个行列式对加的行列的公共元素的余子式 $M_{n,n+1}$ 即可. 求余子式最直接的方法就是将 $n+1$ 阶的行列式按照加的列展开，我们已经知道了整个行列式的值，那么只需要将待求余子式 $M_{n,n+1}$ 与其他部分分离. 由于范德蒙行列式的结构特征，加的列必然是形如 $(1, x, x^2, \ldots, x^n)$ 的列向量，我们发现如果将 $x$ 视为主元，按该列展开后便形成了一个 $n$ 次的多项式，其中 $x^{n-1}$ 的系数正是我们需要的代数余子式 $(-1)^{2n+1} M_{n,n+1}$. 因此就有

    \begin{align*}
        V_{n+1} &= \begin{vmatrix}
            1      & 1      & \cdots & 1      & 1      \\
            x_1    & x_2    & \cdots & x_n    & x      \\
            x_1^2  & x_2^2  & \cdots & x_n^2  & x^2    \\
            \vdots & \vdots & \ddots & \vdots & \vdots \\
            x_1^{n-1} & x_2^{n-1} & \cdots & x_n^{n-1} & x^{n-1} \\
            x_1^n  & x_2^n  & \cdots & x_n^n  & x^n
        \end{vmatrix} = \sum\limits_{i=0}^n (-1)^{i+n+1} M_{i, n+1} x^i. \\
        V_{n+1} &= \left( \prod\limits_{1 \leqslant i < j \leqslant n} (x_j - x_i) \right) \left(\prod\limits_{i=1}^n (x - x_i) \right).
    \end{align*}

    上下两个多项式对应系数相等，故有

    \[
        - M_{n,n+1} = \left( \prod\limits_{1 \leqslant i < j \leqslant n} (x_j - x_i) \right) \sum\limits_{i=1}^n -x_i.
    \]

    因此，$D_n = M_{n, n+1} = \left( \displaystyle\prod_{1 \leqslant i < j \leqslant n} (x_j - x_i) \right) \displaystyle\sum_{i=1}^n x_i$.
\end{solve}

\begin{exercise}\label{prob:L10-例-32}
计算行列式$D_n=\begin{vmatrix}
            1-a_{1} & a_{2}   &        &           &         \\
            -1      & 1-a_{2} & \ddots &           &         \\
                    & -1      & \ddots & a_{n-1}   &         \\
                    &         & \ddots & 1-a_{n-1} & a_{n}   \\
                    &         &        & -1        & 1-a_{n}
        \end{vmatrix}$.
\end{exercise}
\noindent 来源：\href{https://github.com/yhwu-is/Linear-Algebra-Left-Undone/blob/new/%E8%AE%B2%E4%B9%89/%E4%B8%93%E9%A2%98/10%20%E8%A1%8C%E5%88%97%E5%BC%8F.tex}{L10-例-32}。
\begin{solve}
把第一列的元素看成两项的和进行拆列, 得
    \begin{align*}
        D_n= & \begin{vmatrix}
                   1-a_{1} & a_{2}   &        &           &         \\
                   -1      & 1-a_{2} & \ddots &           &         \\
                   0 + 0   & -1      & \ddots & a_{n-1}   &         \\
                   0 + 0   &         & \ddots & 1-a_{n-1} & a_{n}   \\
                   0 + 0   &         &        & -1        & 1-a_{n}
               \end{vmatrix}                \\
        =    & \begin{vmatrix}
                   1  & a_{2}   &        &           &         \\
                   -1 & 1-a_{2} & \ddots &           &         \\
                      & -1      & \ddots & a_{n-1}   &         \\
                      &         & \ddots & 1-a_{n-1} & a_{n}   \\
                      &         &        & -1        & 1-a_{n}
               \end{vmatrix} + \begin{vmatrix}
                                   -a_{1} & a_{2}   &        &           &         \\
                                          & 1-a_{2} & \ddots &           &         \\
                                          & -1      & \ddots & a_{n-1}   &         \\
                                          &         & \ddots & 1-a_{n-1} & a_{n}   \\
                                          &         &        & -1        & 1-a_{n}
                               \end{vmatrix}.
    \end{align*}

    上面第一个行列式的值为 1（从第 1 行开始，每一行依次加到下一行），所以
    \[ \begin{aligned}
            D_n & =1-a_{1}\begin{vmatrix}
                              1-a_{2} & a_{3}   &        &           &         \\
                              -1      & 1-a_{3} & \ddots &           &         \\
                                      & -1      & \ddots & a_{n-1}   &         \\
                                      &         & \ddots & 1-a_{n-1} & a_{n}   \\
                                      &         &        & -1        & 1-a_{n}
                          \end{vmatrix} \\
                & =1-a_{1} D_{n-1} \cdot
        \end{aligned} \]

    这个式子对任何 $n \geqslant 2$ 都成立, 因此有
    \begin{align*}
        D_n & =1-a_{1} D_{n-1}                                          \\
            & =1-a_{1}(1-a_{2} D_{n-2})                                 \\
            & =\cdots                                                   \\
            & =1-a_{1}+a_{1} a_{2}+\cdots+(-1)^n a_{1} a_{2} \cdots a_n \\
            & =1+\sum_{i=1}^{n}(-1)^i \prod_{j=1}^i a_j.
    \end{align*}
\end{solve}

\begin{exercise}\label{prob:L10-例-33}
定义
    \[
        f(x) = \begin{vmatrix}
            x - a_{11} & - a_{12} & \cdots & - a_{1n} \\
            - a_{21} & x - a_{22} & \cdots & - a_{2n} \\
            \vdots   & \vdots   & \ddots & \vdots   \\
            - a_{n1} & - a_{n2} & \cdots & x - a_{nn}
        \end{vmatrix},
    \]
    试使用硬拆法求
    \begin{enumerate}
        \item $f(x)$ 中 $n-1$ 次项系数 $c_{n-1}$ 的值.
        \item $f(x)$ 中常数项系数 $c_0$ 的值.
    \end{enumerate}
\end{exercise}
\noindent 来源：\href{https://github.com/yhwu-is/Linear-Algebra-Left-Undone/blob/new/%E8%AE%B2%E4%B9%89/%E4%B8%93%E9%A2%98/10%20%E8%A1%8C%E5%88%97%E5%BC%8F.tex}{L10-例-33}。
\begin{solve}
\begin{enumerate}
        \item 设列向量组 $\alpha_i = (-a_{1i}, -a_{2i}, \ldots, -a_{ni})^{\mathrm{T}}$，将行列式的第 $i$ 列分解为 $xe_i + \alpha_i$，即

            \[
                f(x) = D(xe_1 + \alpha_1, xe_2 + \alpha_2, \ldots, xe_n + \alpha_n).
            \]

            逐列进行分解，我们可以得到 $2^n$ 个 $n$ 阶行列式. 设某个行列式中有 $p$ 列为 $xe_i$，$n-p$ 列为 $\alpha_i$，对所有取 $xe_i$ 的列展开，就可以得到一个形如 $cx^p$ 的多项式. 因此，我们只需要考虑拆出的 $2^n$ 个行列式中有 $n-1$ 列为 $xe_i$ 的行列式即可. 因此有

            \[
                c_{n-1}x^{n-1} = D(\alpha_1, xe_2, \ldots, xe_n) + D(xe_1, \alpha_2, \ldots, xe_n) + \cdots + D(xe_1, xe_2, \ldots, \alpha_n).
            \]

            化简后便有 $c_{n-1} = -\displaystyle\sum_{i=1}^n a_{ii}$.

        \item 使用硬拆法，拆出的 $2^n$ 个行列式中只有 $D(\alpha_1, \alpha_2, \ldots, \alpha_n)$ 是常数，因此 $c_0 = D(\alpha_1, \alpha_2, \ldots, \alpha_n)$.
    \end{enumerate}
\end{solve}

\begin{exercise}\label{prob:L10-例-34}
解行列式$D_n=\begin{vmatrix}
            x     & -1    &        &         &         \\
                  & x     & \ddots &         &         \\
                  &       & \ddots & -1      &         \\
                  &       &        & x       & -1      \\
            a_{0} & a_{1} & \cdots & a_{n-2} & a_{n-1}
        \end{vmatrix}$.
\end{exercise}
\noindent 来源：\href{https://github.com/yhwu-is/Linear-Algebra-Left-Undone/blob/new/%E8%AE%B2%E4%B9%89/%E4%B8%93%E9%A2%98/10%20%E8%A1%8C%E5%88%97%E5%BC%8F.tex}{L10-例-34}。
\begin{solve}
按最后一行展开，得
    \begin{align*}
        D_n & =\sum_{i=0}^{n-1}(-1)^{n+i+1}a_i\begin{vmatrix} A & O \\ O & B \end{vmatrix}
        =\sum_{i=0}^{n-1}(-1)^{n+i+1}a_i|A||B|                                                        \\
            & =\sum_{i=0}^{n-1}(-1)^{n+i+1}a_i(x^i)\left((-1)^{n-1-i}\right) =\sum_{i=0}^{n-1}a_i x^i,
    \end{align*}

    其中$A\in \mathbf{M}_i(\mathbf{R}),\enspace B\in \mathbf{M}_{n-1-i}(\mathbf{R})$，
    \[ A=\begin{pmatrix}
            x & -1 &        &    &    \\
              & x  & \ddots &    &    \\
              &    & \ddots & -1 &    \\
              &    &        & x  & -1 \\
              &    &        &    & x
        \end{pmatrix},\enspace
        B=\begin{pmatrix}
            -1 &    &        &    &    \\
            x  & -1 &        &    &    \\
               & x  & \ddots &    &    \\
               &    & \ddots & -1 &    \\
               &    &        & x  & -1
        \end{pmatrix}. \]
\end{solve}

\begin{exercise}\label{prob:L10-例-36}
计算行列式$\begin{vmatrix}
            -1 & 1  & 1 & 2 & -1 \\
            0  & -1 & 0 & 1 & 2  \\
            2  & 1  & 1 & 3 & -1 \\
            1  & 2  & 2 & 1 & 0  \\
            0  & 3  & 0 & 1 & 3
        \end{vmatrix}$.
\end{exercise}
\noindent 来源：\href{https://github.com/yhwu-is/Linear-Algebra-Left-Undone/blob/new/%E8%AE%B2%E4%B9%89/%E4%B8%93%E9%A2%98/10%20%E8%A1%8C%E5%88%97%E5%BC%8F.tex}{L10-例-36}。
\begin{solve}
因为第一、第三列含有较多的 $0$，因此在这两列上做 Laplace 展开，得
    \begin{align*}
        |A| & = \begin{vmatrix}
                    -1 & 1 \\ 2 & 1
                \end{vmatrix} \cdot (-1)^{1+3+1+3}
        \begin{vmatrix}
            -1 & 1 & 2 \\ 2 & 1 & 0 \\ 3 & 1 & 3
        \end{vmatrix}                                 \\
            & + \begin{vmatrix}
                    -1 & 1 \\ 1 & 2
                \end{vmatrix} \cdot (-1)^{1+4+1+3}
        \begin{vmatrix}
            -1 & 1 & 2 \\ 1 & 3 & -1 \\ 3 & 1 & 3
        \end{vmatrix}                                \\
            & + \begin{vmatrix}
                    2 & 1 \\ 1 & 1
                \end{vmatrix} \cdot (-1)^{3+4+1+3}
        \begin{vmatrix}
            1 & 2 & -1 \\ -1 & 1 & 2 \\ 3 & 1 & 3
        \end{vmatrix}                                \\
            & = (-3) \times (-11) + (-3) \times 32 + 3 \times (-23) = -132.
    \end{align*}
\end{solve}

\begin{exercise}\label{prob:L10-例-37}
计算行列式$D_n=\begin{vmatrix}
            \cos \beta & 1            &        &              &              \\
            1          & 2 \cos \beta & \ddots &              &              \\
                       & 1            & \ddots & 1            &              \\
                       &              & \ddots & 2 \cos \beta & 1            \\
                       &              &        & 1            & 2 \cos \beta
        \end{vmatrix}$.
\end{exercise}
\noindent 来源：\href{https://github.com/yhwu-is/Linear-Algebra-Left-Undone/blob/new/%E8%AE%B2%E4%B9%89/%E4%B8%93%E9%A2%98/10%20%E8%A1%8C%E5%88%97%E5%BC%8F.tex}{L10-例-37}。
\begin{solve}
\begin{enumerate}
        \item 法一：归纳法
            \begin{align*}
                D_1 & =\cos\beta,                               \\
                D_2 & =\begin{vmatrix}
                            \cos\beta & 1          \\
                            1         & 2\cos\beta
                        \end{vmatrix}=2\cos^2\beta-1=\cos2\beta.
            \end{align*}

            猜想$D_n=\cos n\beta$. 数学归纳证明：

            假设当 $n=k$ 时，结论成立，即 $D_{k}=\cos k \beta$. 现证当 $n=k+1$ 时，结论也成立. $ n=k+1 $ 时，
            \[ D_{k+1} = \begin{vmatrix}
                    \cos \beta & 1            &        &              &              \\
                    1          & 2 \cos \beta & \ddots &              &              \\
                                & 1            & \ddots & 1            &              \\
                                &              & \ddots & 2 \cos \beta & 1            \\
                                &              &        & 1            & 2 \cos \beta
                \end{vmatrix}. \]

            将 $D_{k+1}$ 按最后一行展开, 得
            \begin{align*}
                D_{k+1}={} & (-1)^{k+1+k+1} \cdot 2 \cos \beta
                \begin{vmatrix}
                    \cos \beta & 1            &        &              &              \\
                    1          & 2 \cos \beta & \ddots &              &              \\
                                & 1            & \ddots & 1            &              \\
                                &              & \ddots & 2 \cos \beta & 1            \\
                                &              &        & 1            & 2 \cos \beta
                \end{vmatrix} \\
                            & +(-1)^{k+1+k}
                \begin{vmatrix}
                    \cos \beta & 1            & 0            & \cdots & 0      \\
                    1          & 2 \cos \beta & 1            & \cdots & 0      \\
                    0          & 1            & 2 \cos \beta & \cdots & 0      \\
                    \vdots     & \vdots       & \vdots       & \ddots & \vdots \\
                    0          & 0            & 0            & \cdots & 1
                \end{vmatrix}       \\
                ={}        & 2\cos\beta D_k-D_{k-1}.
            \end{align*}

            而$D_{k}=\cos k \beta,\enspace
                D_{k-1}=\cos (k-1)\beta = \cos (k \beta-\beta) = \cos k \beta \cos \beta+\sin k \beta \sin \beta$. 所以有
            \begin{align*}
                D_{k+1} & = 2 \cos \beta D_{k}-D_{k-1}                                               \\
                        & =2 \cos \beta \cos k \beta-\cos k \beta \cos \beta-\sin k \beta \sin \beta \\
                        & =\cos k \beta \cos \beta-\sin k \beta \sin \beta                           \\
                        & =\cos (k+1) \beta.
            \end{align*}

            则证得$D_n=\cos n\beta,\enspace n\in \mathbf{N}$.

        \item 法二：递推法
            先得到 $D_n$ 的递推式：
            \[
                D_n - 2 \cos \beta D_{n-1} + D_{n-2} = 0,
            \]
            其中 $D_1 = \cos \beta, \enspace D_2 = \cos 2\beta$.

            现求其特征方程
            \[
                \lambda^2 - 2 \lambda \cos \beta + 1 = 0
            \]
            的根，有
            \[
                \lambda = 2 \cos \beta \pm 2 \mathrm{i} \sin \beta.
            \]
            故令
            \[
                D_n = C_1 \cos n\beta + C_2 \sin n\beta,
            \]
            代入初值，有
            \[
                \begin{cases}
                    C_1 \cos \beta + C_2 \sin \beta = \cos \beta, \\
                    C_1 \cos 2\beta + C_2 \sin 2\beta = \cos 2\beta.
                \end{cases}
            \]
            易知 $C_1 = 1, C_2 = 0$ 为其解，故 $D_n = \cos n\beta$.
    \end{enumerate}
\end{solve}

\begin{exercise}\label{prob:L10-例-38}
计算行列式$D_n=
        \begin{vmatrix}
            9 & 5 &        &   &   \\
            4 & 9 & \ddots &   &   \\
              & 4 & \ddots & 5 &   \\
              &   & \ddots & 9 & 5 \\
              &   &        & 4 & 9
        \end{vmatrix}$.
\end{exercise}
\noindent 来源：\href{https://github.com/yhwu-is/Linear-Algebra-Left-Undone/blob/new/%E8%AE%B2%E4%B9%89/%E4%B8%93%E9%A2%98/10%20%E8%A1%8C%E5%88%97%E5%BC%8F.tex}{L10-例-38}。
\begin{solve}
按第一列展开, 得
    \[ D_n=9 D_{n-1}-20 D_{n-2}, \]
    即 $ D_n-9 D_{n-1}+20 D_{n-2}=0 $.
    作特征方程
    \[ x^{2}-9 x+20=0, \]
    解得 $ x_1=4,\enspace x_2=5 $. 则
    \[ D_n=A \cdot 4^n+B \cdot 5^n, \]

    当 $n=1$ 时, $9=4A+5B$；

    当 $n=2$ 时，$61=16A+25B$.

    解得$A=-4,\enspace B=5$，所以
    \[ D_n=5^{n+1}-4^{n+1}. \]
\end{solve}

\begin{exercise}\label{prob:L10-例-39}
计算行列式 $D_n=\begin{vmatrix}
            (a_1+b_1)^{-1} & (a_1+b_2)^{-1} & \cdots & (a_1+b_n)^{-1} \\
            (a_2+b_1)^{-1} & (a_2+b_2)^{-1} & \cdots & (a_2+b_n)^{-1} \\
            \vdots         & \vdots         & \ddots & \vdots         \\
            (a_n+b_1)^{-1} & (a_n+b_2)^{-1} & \cdots & (a_n+b_n)^{-1}
        \end{vmatrix}$.
\end{exercise}
\noindent 来源：\href{https://github.com/yhwu-is/Linear-Algebra-Left-Undone/blob/new/%E8%AE%B2%E4%B9%89/%E4%B8%93%E9%A2%98/10%20%E8%A1%8C%E5%88%97%E5%BC%8F.tex}{L10-例-39}。
\begin{solve}
我们尝试求解 $D_n$ 和 $D_{n-1}$ 之间的递推式. 首先将行列式的前 $n-1$ 列每列都减去第 $n$ 列：

    \begin{align*}
        D_n & = \begin{vmatrix}
                    \dfrac{1}{a_1+b_1}     & \cdots & \dfrac{1}{a_1+b_{n-1}}     & \dfrac{1}{a_1+b_n}     \\[2ex]
                    \vdots                 & \ddots & \vdots                     & \vdots                 \\[2ex]
                    \dfrac{1}{a_{n-1}+b_1} & \cdots & \dfrac{1}{a_{n-1}+b_{n-1}} & \dfrac{1}{a_{n-1}+b_n} \\[2ex]
                    \dfrac{1}{a_n+b_1}     & \cdots & \dfrac{1}{a_n+b_{n-1}}     & \dfrac{1}{a_n+b_n}
                \end{vmatrix}                                                                       \\[2ex]
            & = \begin{vmatrix}
                    \dfrac{b_n-b_1}{(a_1+b_1)(a_1+b_n)}         & \cdots & \dfrac{b_n-b_{n-1}}{(a_1+b_{n-1})(a_1+b_n)}         & \dfrac{1}{a_1+b_n}     \\[2ex]
                    \vdots                                      & \ddots & \vdots                                              & \vdots                 \\[2ex]
                    \dfrac{b_n-b_1}{(a_{n-1}+b_1)(a_{n-1}+b_n)} & \cdots & \dfrac{b_n-b_{n-1}}{(a_{n-1}+b_{n-1})(a_{n-1}+b_n)} & \dfrac{1}{a_{n-1}+b_n} \\[2ex]
                    \dfrac{b_n-b_1}{(a_n+b_1)(a_n+b_n)}         & \cdots & \dfrac{b_n-b_{n-1}}{(a_n+b_{n-1})(a_n+b_n)}         & \dfrac{1}{a_n+b_n}
                \end{vmatrix}.
    \end{align*}
    第二步是提取公因式：
    \[ \text{上式} = \dfrac{\prod\limits_{i=1}^{n-1}(b_n-b_i)}{\prod\limits_{j=1}^{n}(a_j+b_n)}
        \begin{vmatrix}
            \dfrac{1}{a_1+b_1}     & \cdots & \dfrac{1}{a_1+b_{n-1}}     & 1      \\[2ex]
            \vdots                 & \ddots & \vdots                     & \vdots \\[2ex]
            \dfrac{1}{a_{n-1}+b_1} & \cdots & \dfrac{1}{a_{n-1}+b_{n-1}} & 1      \\[2ex]
            \dfrac{1}{a_n+b_1}     & \cdots & \dfrac{1}{a_n+b_{n-1}}     & 1
        \end{vmatrix}. \]
    第三步是将行列式的前 $n-1$ 行每行都减去第 $n$ 行：
    \[ \text{上式} = \dfrac{\prod\limits_{i=1}^{n-1}(b_n-b_i)}{\prod\limits_{j=1}^{n}(a_j+b_n)}
        \begin{vmatrix}
            \dfrac{a_n-a_1}{(a_1+b_1)(a_n+b_1)}         & \cdots & \dfrac{a_n-a_1}{(a_1+b_{n-1})(a_n+b_{n-1})}         & 0      \\[2ex]
            \vdots                                      & \ddots & \vdots                                              & \vdots \\[2ex]
            \dfrac{a_n-a_{n-1}}{(a_{n-1}+b_1)(a_n+b_1)} & \cdots & \dfrac{a_n-a_{n-1}}{(a_{n-1}+b_{n-1})(a_n+b_{n-1})} & 0      \\[2ex]
            \dfrac{1}{a_n+b_1}                          & \cdots & \dfrac{1}{a_n+b_{n-1}}                              & 1
        \end{vmatrix}. \]
    然后再次提取公因式：
    \begin{align*}
        \text{上式} & = \dfrac{\prod\limits_{i=1}^{n-1}(a_n-a_i)(b_n-b_i)}{\prod\limits_{j=1}^{n}(a_j+b_n)\prod\limits_{k=1}^{n-1}(a_n+b_k)}
        \begin{vmatrix}
            \dfrac{1}{a_1+b_1}     & \cdots & \dfrac{1}{a_1+b_{n-1}}     \\[2ex]
            \vdots                 & \ddots & \vdots                     \\[2ex]
            \dfrac{1}{a_{n-1}+b_1} & \cdots & \dfrac{1}{a_{n-1}+b_{n-1}}
        \end{vmatrix}                                                         \\[2ex]
                    & = \dfrac{\prod\limits_{i=1}^{n-1}(a_n-a_i)(b_n-b_i)}{\prod\limits_{j=1}^{n}(a_j+b_n)\prod\limits_{k=1}^{n-1}(a_n+b_k)} D_{n-1}.
    \end{align*}
    不断递推下去即得

    \[|A| = \dfrac{\prod\limits_{1 \leqslant i < j \leqslant n}(a_j - a_i)(b_j - b_i)}{\prod\limits_{i,j=1}^{n}(a_i + b_j)}.\]
\end{solve}

\begin{exercise}\label{prob:L10-例-40}
求下列矩阵行列式的值：
            \[A=\begin{pmatrix}
                    a_1^2    & 1+a_1a_2   & \cdots & 1+a_1a_n \\
                    1+a_2a_1 & a_2^2    & \cdots & 1+a_2a_n \\
                    \vdots   & \vdots   & \ddots & \vdots   \\
                    1+a_na_1 & 1+a_na_2 & \cdots & a_n^2
                \end{pmatrix}.\]
\end{exercise}
\noindent 来源：\href{https://github.com/yhwu-is/Linear-Algebra-Left-Undone/blob/new/%E8%AE%B2%E4%B9%89/%E4%B8%93%E9%A2%98/10%20%E8%A1%8C%E5%88%97%E5%BC%8F.tex}{L10-例-40}。
\begin{solve}
我们可以将 $A$ 化为

            \[ A = -E_n + \begin{pmatrix}
                    a_1 & 1 \\ a_2 & 1 \\ \vdots & \vdots \\ a_n & 1
                \end{pmatrix} E_2^{-1} \begin{pmatrix}
                    a_1 & a_2 & \cdots & a_n \\ 1 & 1 & \cdots & 1
                \end{pmatrix}. \]

            故根据降阶公式有

            \begin{align*}
                |A| & = \left\vert-E_n +
                \begin{pmatrix}
                    a_1 & 1 \\ a_2 & 1 \\ \vdots & \vdots \\ a_n & 1
                \end{pmatrix} E_2^{-1}
                \begin{pmatrix}
                    a_1 & a_2 & \cdots & a_n \\ 1 & 1 & \cdots & 1
                \end{pmatrix}\right\vert                                                                       \\
                    & = (-1)^n \left|E_2 -
                \begin{pmatrix}
                    \sum\limits_{i=1}^n a_i^2 & \sum\limits_{i=1}^n a_i \\[2ex]
                    \sum\limits_{i=1}^n a_i   & n
                \end{pmatrix}\right|                                                          \\
                    & = (-1)^n \left((1-n)\left(1-\sum\limits_{i=1}^n a_i^2\right) - \left(\sum\limits_{i=1}^n a_i\right)^2\right).
            \end{align*}
\end{solve}

\begin{exercise}\label{prob:L10-例-41}
求行列式 $\begin{vmatrix}
                    0      & 2a_1   & 3a_1   & \cdots & na_1     \\
                    a_2    & a_2    & 3a_2   & \cdots & na_2     \\
                    a_3    & 2a_3   & 2a_3   & \cdots & na_3     \\
                    \vdots & \vdots & \vdots & \ddots & \vdots   \\
                    a_n    & 2a_n   & 3a_n   & \cdots & (n-1)a_n \\
                \end{vmatrix}$.
\end{exercise}
\noindent 来源：\href{https://github.com/yhwu-is/Linear-Algebra-Left-Undone/blob/new/%E8%AE%B2%E4%B9%89/%E4%B8%93%E9%A2%98/10%20%E8%A1%8C%E5%88%97%E5%BC%8F.tex}{L10-例-41}。
\begin{solve}
\[ \text{原式}=\prod_{i=1}^na_i \begin{vmatrix}
                0      & 2      & 3      & \cdots & n      \\
                1      & 1      & 3      & \cdots & n      \\
                1      & 2      & 2      & \cdots & n      \\
                \vdots & \vdots & \vdots & \ddots & \vdots \\
                1      & 2      & 3      & \cdots & n-1    \\
            \end{vmatrix}. \]

            注意到这个行列式对应的矩阵可以加上一个单位矩阵构成一个每行都相等的矩阵，所以我们可以利用$|E_m-AB|=|E_n-BA|$来求解这个行列式.
            \begin{align*}
                \begin{vmatrix}
                    0      & 2      & 3      & \cdots & n      \\
                    1      & 1      & 3      & \cdots & n      \\
                    1      & 2      & 2      & \cdots & n      \\
                    \vdots & \vdots & \vdots & \ddots & \vdots \\
                    1      & 2      & 3      & \cdots & n-1
                \end{vmatrix}
                & = \begin{vmatrix}(-1)\left(E_n-\begin{pmatrix}
                        1      & 2      & 3      & \cdots & n      \\
                        1      & 2      & 3      & \cdots & n      \\
                        1      & 2      & 3      & \cdots & n      \\
                        \vdots & \vdots & \vdots & \ddots & \vdots \\
                        1      & 2      & 3      & \cdots & n
                    \end{pmatrix}\right)\end{vmatrix} \\
                & =(-1)^n\begin{vmatrix}E_n-
                            \begin{pmatrix}
                        1 \\1\\1\\\vdots\\1
                    \end{pmatrix}\begin{pmatrix}1 & 2 & 3 & \cdots & n\end{pmatrix}\end{vmatrix}.
            \end{align*}

            而 \[ \begin{vmatrix}E_n-\begin{pmatrix}
                        1 \\1\\1\\\vdots\\1
                    \end{pmatrix}\begin{pmatrix}1 & 2 & 3 & \cdots & n\end{pmatrix}\end{vmatrix}
                =1-\begin{pmatrix}1 & 2 & 3 & \cdots & n\end{pmatrix}
                \begin{pmatrix}1 \\ 1 \\ 1 \\ \vdots \\ 1\end{pmatrix}
                =-\frac{n^2+n-2}{2}. \]

            所以原式$\displaystyle =(-1)^{n+1} \cdot \frac{n^2+n-2}{2} \left(\prod_{i=1}^na_i \right)$.
\end{solve}

\begin{exercise}\label{prob:L10-例-42}
解行列式$D_n=\begin{vmatrix}
                    \lambda & a      & a      & a      & \cdots & a      \\
                    b       & \gamma & \beta  & \beta  & \cdots & \beta  \\
                    b       & \beta  & \gamma & \beta  & \cdots & \beta  \\
                    \vdots  & \vdots & \vdots & \vdots & \ddots & \vdots \\
                    b       & \beta  & \beta  & \beta  & \cdots & \gamma
                \end{vmatrix}$.
\end{exercise}
\noindent 来源：\href{https://github.com/yhwu-is/Linear-Algebra-Left-Undone/blob/new/%E8%AE%B2%E4%B9%89/%E4%B8%93%E9%A2%98/10%20%E8%A1%8C%E5%88%97%E5%BC%8F.tex}{L10-例-42}。
\begin{solve}
从第$n$行到第3行，每行都减去上一行；再从第3列到第$n$列，每列都加到第2列，得

            \begin{align*}
                D_n & =
                \begin{vmatrix}
                    \lambda & a            & a            & a            & \cdots & a            \\
                    b       & \gamma       & \beta        & \beta        & \cdots & \beta        \\
                    0       & \beta-\gamma & \gamma-\beta & 0            & \cdots & 0            \\
                    0       & 0            & \beta-\gamma & \gamma-\beta & \cdots & 0            \\
                    \vdots  & \vdots       & \vdots       & \vdots       & \ddots & \vdots       \\
                    0       & 0            & 0            & 0            & \cdots & \gamma-\beta
                \end{vmatrix}             \\
                    & =\begin{vmatrix}
                           \lambda & (n-1)a            & a            & a            & \cdots & a            \\
                           b       & \gamma+(n-2)\beta & \beta        & \beta        & \cdots & \beta        \\
                           0       & 0                 & \gamma-\beta & 0            & \cdots & 0            \\
                           0       & 0                 & \beta-\gamma & \gamma-\beta & \cdots & 0            \\
                           \vdots  & \vdots            & \vdots       & \vdots       & \ddots & \vdots       \\
                           0       & 0                 & 0            & 0            & \cdots & \gamma-\beta
                       \end{vmatrix} \\
                    & =\begin{vmatrix}
                           \lambda & (n-1)a            \\
                           b       & \gamma+(n-2)\beta
                       \end{vmatrix} \cdot \begin{vmatrix}
                                               \gamma-\beta & 0            & \cdots & 0            \\
                                               \beta-\gamma & \gamma-\beta & \cdots & 0            \\
                                               \vdots       & \vdots       & \ddots & \vdots       \\
                                               0            & 0            & \cdots & \gamma-\beta
                                           \end{vmatrix}           \\
                    & =(\lambda \gamma+\lambda(n-2)\beta-(n-1)ab)(\gamma-\beta)^{n-2}.
            \end{align*}
\end{solve}

\begin{exercise}\label{prob:L10-例-43}
证明 Lagrange 恒等式 $(n \geqslant 2)$：
    \[ \left( \sum_{i=1}^{n}\limits a_i^2 \right) \left( \sum_{i=1}^{n}\limits b_i^2 \right) - \left( \sum_{i=1}^{n}\limits a_i b_i \right)^2 = \sum_{1 \leqslant i < j \leqslant n} (a_i b_j - a_j b_i)^2. \]
\end{exercise}
\noindent 来源：\href{https://github.com/yhwu-is/Linear-Algebra-Left-Undone/blob/new/%E8%AE%B2%E4%B9%89/%E4%B8%93%E9%A2%98/10%20%E8%A1%8C%E5%88%97%E5%BC%8F.tex}{L10-例-43}。
\begin{solve}
左边的式子等于
    \[ \begin{vmatrix}
        \sum_{i=1}^{n}\limits a_i^2 & \sum_{i=1}^{n}\limits a_i b_i \\
        \sum_{i=1}^{n}\limits a_i b_i & \sum_{i=1}^{n}\limits b_i^2
    \end{vmatrix}, \]

    这个行列式对应的矩阵可化为：
    \[ \begin{pmatrix}
        a_1 & a_2 & \cdots & a_n \\
        b_1 & b_2 & \cdots & b_n
    \end{pmatrix}
    \begin{pmatrix}
        a_1 & b_1 \\
        a_2 & b_2 \\
        \vdots & \vdots \\
        a_n & b_n
    \end{pmatrix}. \]
    用 Cauchy-Binet 公式得

    \[
    \begin{vmatrix}
        \sum_{i=1}^{n}\limits a_i^2 & \sum_{i=1}^{n}\limits a_i b_i \\
        \sum_{i=1}^{n}\limits a_i b_i & \sum_{i=1}^{n}\limits b_i^2
    \end{vmatrix}
    = \sum_{1 \leqslant i < j \leqslant n}
    \begin{vmatrix}
        a_i & a_j \\
        b_i & b_j
    \end{vmatrix}^2
    = \sum_{1 \leqslant i < j \leqslant n} (a_i b_j - a_j b_i)^2.
    \]
\end{solve}

\section{线性方程组与矩阵综合}\label{sec:topic11}
本节收录 56 道题，其中 12 道为编者补充计算题。先用系数矩阵和增广矩阵的秩判定有解性，再把非齐次解集写成特解加齐次解空间，最后验证所得基础解系。
\begin{exercise}\label{prob:L11-A-2}
设$A$为四阶矩阵，$r(A)<4$，且$A_{21}\neq 0$，求方程组$AX=\vec{0}$的通解.
\end{exercise}
\noindent 来源：\href{https://github.com/yhwu-is/Linear-Algebra-Left-Undone/blob/new/%E8%AE%B2%E4%B9%89/%E4%B8%93%E9%A2%98/11%20%E6%9C%9D%E8%8A%B1%E5%A4%95%E6%8B%BE.tex}{L11-A-2}。
\begin{solve}
$r(\operatorname{adj}(A))=
              \begin{cases}
                  1 & r(A)=3 \\
                  0 & r(A)<3
              \end{cases}$，而已知$A_{21}\neq 0$，故$r(\operatorname{adj}(A))=1$，故$r(A)=3$，故$AX = 0$解空间维数为$4-r(A)=1$. 而$|A|=0$，由行列式展开可知
          \[ a_{i1}A_{21}+a_{i2}A_{22}+a_{i3}A_{23}+a_{i4}A_{24}=0 \qquad i=1,2,3,4 \]
          故解为$k(A_{21},A_{22},A_{23},A_{24})^\mathrm{T},\enspace k\in \mathbf{R}$.
\end{solve}

\begin{exercise}\label{prob:L11-A-3}
设$A=(\alpha_1,\alpha_2,\alpha_3,\alpha_4)$为四阶矩阵，方程组$AX=\vec{0}$的通解为$X=k(1,0,-4,0)^\mathrm{T}$，求$\operatorname{adj}(A)X=0$的基础解系.
\end{exercise}
\noindent 来源：\href{https://github.com/yhwu-is/Linear-Algebra-Left-Undone/blob/new/%E8%AE%B2%E4%B9%89/%E4%B8%93%E9%A2%98/11%20%E6%9C%9D%E8%8A%B1%E5%A4%95%E6%8B%BE.tex}{L11-A-3}。
\begin{solve}
由题意得$r(A)=3<4$且$\alpha_1-4\alpha_3=0$，或$\alpha_1=4\alpha_3$. 由$r(A)=3$得$r(\operatorname{adj}(A))=1$，$\operatorname{adj}(A)X=0$的基础解系含3个线性无关的解向量. 由$\operatorname{adj}(A)A=|A|E=O$得$\alpha_1,\alpha_2,\alpha_3,\alpha_4$为$\operatorname{adj}(A)X=0$的解，从而$\alpha_1,\alpha_2,\alpha_4$或$\alpha_2,\alpha_3,\alpha_4$为$\operatorname{adj}(A)X=0$的基础解系.
\end{solve}

\begin{exercise}\label{prob:L11-A-4}
设$A$为$n$阶实矩阵，$W=\{\beta\in\mathbf{R}^n \mid \alpha^\mathrm{T}A\beta=0,\enspace \forall \alpha\in\mathbf{R}^n\}$，证明：
        \begin{enumerate}
            \item $\dim W+r(A)=n$；

            \item $W$为$\mathbf{R}^n$的子空间.
        \end{enumerate}
\end{exercise}
\noindent 来源：\href{https://github.com/yhwu-is/Linear-Algebra-Left-Undone/blob/new/%E8%AE%B2%E4%B9%89/%E4%B8%93%E9%A2%98/11%20%E6%9C%9D%E8%8A%B1%E5%A4%95%E6%8B%BE.tex}{L11-A-4}。
\begin{solve}
$\forall \alpha \in \mathbf{R}^n, \alpha^\mathrm{T}A\beta =0\implies A\beta =0$.
\end{solve}

\begin{exercise}\label{prob:L11-A-5}
已知4级方阵$A=(\alpha_1,\alpha_2,\alpha_3,\alpha_4)$的列向量$\alpha_1,\alpha_2,\alpha_4$线性无关，且$\alpha_1=2\alpha_2-\alpha_3$，若$\beta=\alpha_1-\alpha_2+3\alpha_4$，求方程组$AX=\beta$的通解.
\end{exercise}
\noindent 来源：\href{https://github.com/yhwu-is/Linear-Algebra-Left-Undone/blob/new/%E8%AE%B2%E4%B9%89/%E4%B8%93%E9%A2%98/11%20%E6%9C%9D%E8%8A%B1%E5%A4%95%E6%8B%BE.tex}{L11-A-5}。
\begin{solve}
$r(A)=3\implies \dim{N(A)}=4-r(A)=1$.

          $\beta = x_1\alpha_1+x_2\alpha_2+x_3\alpha_3+x_4\alpha_4\implies$特解为$(1,-1,0,3)^\mathrm{T}$

          $\alpha_1-2\alpha_2+\alpha_3=0\implies$导出组基础解系为$(1,-2,1,0)^\mathrm{T}$.

          故通解为$k(1,-2,1,0)^\mathrm{T}+(1,-1,0,3)^\mathrm{T},\enspace k\in\mathbf{R}$
\end{solve}

\begin{exercise}\label{prob:L11-A-6}
设四元非齐次线性方程组的系数矩阵的秩为3，已知$\eta_1,\eta_2,\eta_3$是它的三个解向量，且$\eta_1=(2,3,4,5)^\mathrm{T},\eta_2+\eta_3=(1,2,3,4)^\mathrm{T}$，求该方程组的通解.
\end{exercise}
\noindent 来源：\href{https://github.com/yhwu-is/Linear-Algebra-Left-Undone/blob/new/%E8%AE%B2%E4%B9%89/%E4%B8%93%E9%A2%98/11%20%E6%9C%9D%E8%8A%B1%E5%A4%95%E6%8B%BE.tex}{L11-A-6}。
\begin{solve}
设方程组$Ax=b$，由$r(A)=3$知其导出组$Ax=0$的基础解系只含一个解向量.

          而$A\eta_1=b,A(\eta_2+\eta_3)=2b$，故$2\eta_1-(\eta_2+\eta_3)=(3,4,5,6)^\mathrm{T}$为$Ax=0$的基础解系，从而所求通解为
          \[ x = \eta_1+c(3,4,5,6)^\mathrm{T} = (2,3,4,5)^\mathrm{T} + c(3,4,5,6)^\mathrm{T} \]
          其中$c$为任意常数.
\end{solve}

\begin{exercise}\label{prob:L11-A-7}
设$\beta_1,\beta_2,\beta_3$是$n$元非齐次线性方程组$AX=\vec{b}$的三个线性无关的解，且$r(A)=n-2$，求：
        \begin{enumerate}
            \item 导出组$AX=\vec{0}$的一个基础解系；

            \item $AX=2\vec{b}$的一般解.
        \end{enumerate}
\end{exercise}
\noindent 来源：\href{https://github.com/yhwu-is/Linear-Algebra-Left-Undone/blob/new/%E8%AE%B2%E4%B9%89/%E4%B8%93%E9%A2%98/11%20%E6%9C%9D%E8%8A%B1%E5%A4%95%E6%8B%BE.tex}{L11-A-7}。
\begin{solve}
\begin{enumerate}
                \item 取$\beta_1-\beta_2,\beta_1-\beta_3$（取法不唯一），有$A(\beta_1-\beta_2)=A(\beta_1-\beta_3)=0$且$\beta_1-\beta_2,\beta_1-\beta_3$线性无关，否则$\beta_1,\beta_2,\beta_3$线性相关.

                \item 取$2\beta_1+k_1(\beta_1+\beta_2)+k_2(\beta_2+\beta_3)$可行（取法不唯一）.
            \end{enumerate}
\end{solve}

\begin{exercise}\label{prob:L11-A-8}
已知$A$是一个$s\times n$矩阵，证明：线性方程组$AX=\vec{b}$对任意列向量$\vec{b}_{s\times 1}$都有解的充要条件是$A$行满秩.
\end{exercise}
\noindent 来源：\href{https://github.com/yhwu-is/Linear-Algebra-Left-Undone/blob/new/%E8%AE%B2%E4%B9%89/%E4%B8%93%E9%A2%98/11%20%E6%9C%9D%E8%8A%B1%E5%A4%95%E6%8B%BE.tex}{L11-A-8}。
\begin{solve}
对于任意的$b_{s\times 1}$，$AX=b$都有解，说明$A$的列向量可以线性表出出任意$s$维向量$b$，从而$r(A)\geqslant s$，而$r(A)\leqslant s$，所以$r(A)=s$.

          反之，如果$r(A)=s$，则$A$有$s$个线性无关的列向量，它们就是$s$维空间$V$的一组基，对于$V$中任意向量$b$，都可以由$A$的列向量线性表出，从而$AX=b$有解.

          当然，也可以直接用秩不等式：$s=r(A)\leqslant r(A,b)\leqslant s$得到$r(A,b)=s$，所以$AX=b$有解.
\end{solve}

\begin{exercise}\label{prob:L11-A-9}
设$A,B$分别是$m \times n$和$n \times s$矩阵，且$r(B)=n$，证明：若$AB=O$，则$A=O$.
\end{exercise}
\noindent 来源：\href{https://github.com/yhwu-is/Linear-Algebra-Left-Undone/blob/new/%E8%AE%B2%E4%B9%89/%E4%B8%93%E9%A2%98/11%20%E6%9C%9D%E8%8A%B1%E5%A4%95%E6%8B%BE.tex}{L11-A-9}。
\begin{solve}
$AB=0\implies r(A)+r(B)\leqslant n,r(B)=n\implies r(A)\leqslant 0\implies A=O$.
\end{solve}

\begin{exercise}\label{prob:L11-A-10}
设$A \in \mathbf{F}^{m \times n},B \in \mathbf{F}^{(n-m) \times n}\enspace(m<n)$，$V_1,V_2$分别为齐次线性方程组$AX=\vec{0}$和$BX=\vec{0}$的解空间，证明：$\mathbf{F}^n=V_1\oplus V_2$的充要条件是$\begin{pmatrix} A \\ B \end{pmatrix}X=\vec{0}$只有零解.
\end{exercise}
\noindent 来源：\href{https://github.com/yhwu-is/Linear-Algebra-Left-Undone/blob/new/%E8%AE%B2%E4%B9%89/%E4%B8%93%E9%A2%98/11%20%E6%9C%9D%E8%8A%B1%E5%A4%95%E6%8B%BE.tex}{L11-A-10}。
\begin{solve}
必要性. 有条件可知$V_1\cap V_2=\{0\}$，$\begin{pmatrix}A\\B\end{pmatrix}\in\mathbf{F}^{n\times n}$，而$\begin{pmatrix}A\\B\end{pmatrix}x=0$只有零解，故$\begin{pmatrix}A\\B\end{pmatrix}$可逆，从而$r(A)=m,r(B)=n-m$，于是
          \begin{align*}
              \dim V_1 & =n-r(A)=n-m,       \\
              \dim V_2 & =n-r(B)=n-(n-m)=m.
          \end{align*}
          又$V_1+V_2$是$\mathbf{F}^n$的子空间，且$\dim(V_1+V_2)=\dim V_1+\dim V_2-\dim(V_1\cap V_2)=\dim\mathbf{F}^n$，故$\mathbf{F}^n=V_1\oplus V_2$.

          充分性. 若$\begin{pmatrix}A\\B\end{pmatrix}x=0$有非零解$x_1$，则$x_1\in V_1\cap V_2$. 这与$\mathbf{F}^n=V_1\oplus V_2$矛盾.
\end{solve}

\begin{exercise}\label{prob:L11-A-11}
齐次线性方程组\[\begin{cases}
                x_2+ax_3+bx_4=0  \\
                -x_1+cx_3+dx_4=0 \\
                ax_1+cx_2-ex_4=0 \\
                bx_1+dx_2+ex_3=0
            \end{cases}\]的一般解以$x_3,x_4$作为自由未知量.
        \begin{enumerate}
            \item 求$a,b,c,d,e$满足的的条件；

            \item 求齐次线性方程组的基础解系.
        \end{enumerate}
\end{exercise}
\noindent 来源：\href{https://github.com/yhwu-is/Linear-Algebra-Left-Undone/blob/new/%E8%AE%B2%E4%B9%89/%E4%B8%93%E9%A2%98/11%20%E6%9C%9D%E8%8A%B1%E5%A4%95%E6%8B%BE.tex}{L11-A-11}。
\begin{solve}
\begin{enumerate}
                \item 由条件知，方程组系数矩阵的为2，系数矩阵
                      \[ A=\begin{pmatrix}
                              0  & 1 & a & b  \\
                              -1 & 0 & c & d  \\
                              a  & c & 0 & -e \\
                              b  & d & e & 0
                          \end{pmatrix}\rightarrow
                          \begin{pmatrix}
                              0  & 1 & a          & b       \\
                              -1 & 0 & c          & d       \\
                              0  & 0 & 0          & ad-e-bc \\
                              0  & 0 & -(ad-e-bc) & 0
                          \end{pmatrix} \]
                      故$ad-e-bc=0$.

                \item 易求得基础解系为$(c,-a,1,0)^\mathrm{T},(d,-b,0,1)^\mathrm{T}$.
            \end{enumerate}
\end{solve}

\begin{exercise}\label{prob:L11-B-1}
设$A$为方阵，证明：$A^2=E \iff r(A+E)+r(A-E)=n$.
\end{exercise}
\noindent 来源：\href{https://github.com/yhwu-is/Linear-Algebra-Left-Undone/blob/new/%E8%AE%B2%E4%B9%89/%E4%B8%93%E9%A2%98/11%20%E6%9C%9D%E8%8A%B1%E5%A4%95%E6%8B%BE.tex}{L11-B-1}。
\begin{solve}
\begin{enumerate}
            \item 充分性：有 $(A+E)(A-E) = A^2 - E = 0$, 由定理 $BC=0 \implies r(B)+r(C) \leqslant n$，得 $r(A+E)+r(A-E) \leqslant n$；而 $r(A+E)+r(A-E) \geqslant r(A+E+A-E) = n$，故 $r(A+E)+r(A-E)=n$.
            \item 必要性：由$r(A+E)+r(A-E)=n$可得出 $A$ 的特征值仅可能为 1 或者 -1，并且可以对角化.

                因此，记 $A = PBP^{-1}$，其中 $B$ 为对角矩阵，则 $(A+E)(A-E)=P(B-E)(B+E)P^{-1}=0$，故命题得证.
        \end{enumerate}
\end{solve}

\begin{exercise}\label{prob:L11-B-2}
设 $n\ge1$，$a_1,\ldots,a_n$ 为互异实数，$b_1,\ldots,b_n$ 为任意实数。证明：存在唯一的次数不超过 $n-1$ 的实多项式 $f$（包括零多项式），满足 $f(a_i)=b_i$（$1\le i\le n$）。
\end{exercise}
\noindent 来源：\href{https://github.com/yhwu-is/Linear-Algebra-Left-Undone/blob/new/%E8%AE%B2%E4%B9%89/%E4%B8%93%E9%A2%98/11%20%E6%9C%9D%E8%8A%B1%E5%A4%95%E6%8B%BE.tex}{L11-B-2}。
\begin{solve}
设 $f(x) = c_0+c_1x+c_2x^2+\cdots+c_{n-1}x^{n-1}$，则有
            \[\begin{cases} \begin{aligned}
                        c_0+c_1a_1+c_2a_1^2+\cdots+c_{n-1}a_1^{n-1} & = b_1,            \\
                        c_0+c_1a_2+c_2a_2^2+\cdots+c_{n-1}a_2^{n-1} & = b_2,            \\
                                                                    & \vdotswithin{ = } \\
                        c_0+c_1a_n+c_2a_n^2+\cdots+c_{n-1}a_n^{n-1} & = b_n.            \\
                    \end{aligned} \end{cases}\]
            注意此处我们研究的对象是 $(c_0, c_1, \ldots, c_{n-1})^{\mathrm{T}}$. 因为系数行列式
            \[D = \begin{vmatrix}
                    1      & a_1    & \cdots & a_1^{n-1} \\
                    1      & a_2    & \cdots & a_2^{n-1} \\
                    \vdots & \vdots & \ddots & \vdots    \\
                    1      & a_n    & \cdots & a_n^{n-1} \\
                \end{vmatrix} = \prod_{1 \leqslant j < i \leqslant n} (a_i-a_j) \neq 0.\]
            所以由 Cramer 法则，上述关于 $(c_0, c_1, \ldots, c_{n-1})^{\mathrm{T}}$ 的方程组有唯一解，所以满足条件的多项式函数 $f$ 是唯一存在的.
\end{solve}

\begin{exercise}\label{prob:L11-B-3}
设$A$为$m \times n$矩阵，$r(A)=m$，$B$是$m$阶可逆矩阵，已知$A$的行空间$R(A^\mathrm{T})$是方程组$CX=\vec{0}$的解空间，证明：$BA$的行向量也是$CX=\vec{0}$的一个基础解系.
\end{exercise}
\noindent 来源：\href{https://github.com/yhwu-is/Linear-Algebra-Left-Undone/blob/new/%E8%AE%B2%E4%B9%89/%E4%B8%93%E9%A2%98/11%20%E6%9C%9D%E8%8A%B1%E5%A4%95%E6%8B%BE.tex}{L11-B-3}。
\begin{solve}
因 $r(A)=m$，$A$ 的 $m$ 个行向量线性无关。左乘可逆矩阵 $B$ 保持行空间与秩，所以 $BA$ 的 $m$ 个行向量仍线性无关，并张成同一行空间。把行向量转置为列向量后，有
\[C(BA)^{\mathrm T}=CA^{\mathrm T}B^{\mathrm T}=O.\]
故它们都是方程 $CX=0$ 的解；由已知该解空间就是 $A$ 的行空间，$BA$ 的各行转置后构成所求基础解系。
\end{solve}

\begin{exercise}\label{prob:L11-B-4}
设$A$是数域$\mathbf{F}$上的一个$n$阶可逆方阵，$A$的前$r$个行向量组成的矩阵为$B$，后$n-r$个行向量组成的矩阵为$C$，$n$元线性方程组$BX=0$与$CX=0$的解空间分别为$V_1,V_2$. 证明：$\mathbf{F}^n=V_1\oplus V_2$.
\end{exercise}
\noindent 来源：\href{https://github.com/yhwu-is/Linear-Algebra-Left-Undone/blob/new/%E8%AE%B2%E4%B9%89/%E4%B8%93%E9%A2%98/11%20%E6%9C%9D%E8%8A%B1%E5%A4%95%E6%8B%BE.tex}{L11-B-4}。
\begin{solve}
先记$W=V_1+V_2$，若$\alpha\in V_1\cap V_2$，则$B\alpha=C\alpha=0$，所以
            \[A\alpha=\begin{pmatrix}
                    B \\
                    C
                \end{pmatrix}\alpha=0,\]
            由于$A$可逆，因此$\alpha=0$，即$V_1\cap V_2=\{0\}$，因此$V_1+V_2$是直和，因此只需证$W=\mathbf{F}^n$即可. 事实上，我们知道$r(B)=r,r(C)=n-r$，因此$\dim V_1=n-r,\enspace \dim V_2=n-(n-r)=r$，所以
            \[\dim W=\dim V_1+\dim V_2=n-r+r=n=\dim \mathbf{F}^n,\]
            又$W=V_1\oplus V_2\subseteq \mathbf{F}^n$，因此$W=\mathbf{F}^n$，故得证.
\end{solve}

\begin{exercise}\label{prob:L11-B-5}
设$A$是$n$阶矩阵，且$A_{11}\neq 0$，证明：方程组$AX=\vec{b}$（$\vec{b}$为非零向量）有无穷多解的充要条件为$\operatorname{adj}(A)\vec{b}=\vec{0}$.
\end{exercise}
\noindent 来源：\href{https://github.com/yhwu-is/Linear-Algebra-Left-Undone/blob/new/%E8%AE%B2%E4%B9%89/%E4%B8%93%E9%A2%98/11%20%E6%9C%9D%E8%8A%B1%E5%A4%95%E6%8B%BE.tex}{L11-B-5}。
\begin{solve}
必要性：设$AX=b$有无穷多解，则$r(A)<n$，从而$|A|=0$，于是$\operatorname{adj}(A)b=\operatorname{adj}(A)AX=|A|X=0$.

          充分性：设$\operatorname{adj}(A)b=0$，即方程组$\operatorname{adj}(A)X=0$有非零解，则$r(\operatorname{adj}(A))<n$，又$A_{11}\neq 0$，所以$r(A)=n-1$.

          令$A=(\alpha_1,\alpha_2,\ldots,\alpha_n)$，因为$\operatorname{adj}(A)A=|A|E=O$，所以$\alpha_1,\alpha_2,\ldots,\alpha_n$为方程组$\operatorname{adj}(A)X=0$的解，又因为$A_{11}\neq 0$，所以$\alpha_2,\alpha_3,\ldots,\alpha_n$线性无关.

          由$r(\operatorname{adj}(A))=1$，得方程组$\operatorname{adj}(A)X=0$的基础解系含有$n-1$个线性无关的解向量，所以$\alpha_2,\alpha_3,\ldots,\alpha_n$为方程组$\operatorname{adj}(A)X=0$的一个基础解系.

          因为$\operatorname{adj}(A)b=0$，所以$b$为方程组$\operatorname{adj}(A)X=0$的一个解，从而$b$可由$\alpha_2,\alpha_3,\ldots,\alpha_n$线性表示，$b$也可由$\alpha_1,\alpha_2,\ldots,\alpha_n$线性表示，于是$r(A) = r\begin{pmatrix}A & b\end{pmatrix}=n-1<n$，故方程组$AX=b$有无穷多解.
\end{solve}

\begin{exercise}\label{prob:L11-B-6}
若$n$阶矩阵$A$的各行、各列元素之和都为0，证明：$|A|$的所有元素的代数余子式都相等.
\end{exercise}
\noindent 来源：\href{https://github.com/yhwu-is/Linear-Algebra-Left-Undone/blob/new/%E8%AE%B2%E4%B9%89/%E4%B8%93%E9%A2%98/11%20%E6%9C%9D%E8%8A%B1%E5%A4%95%E6%8B%BE.tex}{L11-B-6}。
\begin{solve}
当 $n=1$ 时只有一个代数余子式，结论成立。设 $n\ge2$，令 $e=(1,\ldots,1)^{\mathrm T}$，则 $Ae=0$、$A^{\mathrm T}e=0$，故 $A$ 奇异。若 $r(A)\le n-2$，所有代数余子式都是零。若 $r(A)=n-1$，则 $\ker A=\ker A^{\mathrm T}=\operatorname{span}(e)$。由 $A\operatorname{adj}(A)=O$，伴随矩阵各列都是 $e$ 的倍数；再由 $\operatorname{adj}(A)A=O$，其各行也都是 $e^{\mathrm T}$ 的倍数。因此存在标量 $c$ 使 $\operatorname{adj}(A)=cee^{\mathrm T}$，所有代数余子式均等于 $c$。
\end{solve}

\begin{exercise}\label{prob:L11-B-7}
已知 $\alpha_1,\alpha_2,\ldots,\alpha_s$ 是齐次线性方程组 $AX=\vec{0}$ 的一组基础解系，向量组
        \[\beta_1=t_1\alpha_1+t_2\alpha_2,\ \beta_2=t_1\alpha_2+t_2\alpha_3,\ \ldots,\ \beta_{s-1}=t_1\alpha_{s-1}+t_2\alpha_s\]
        试问当实数 $t_1,t_2$ 满足何条件时，$AX=\vec{0}$ 有基础解系包含向量 $\beta_1,\beta_2,\ldots,\beta_{s-1}$，并写出该基础解系中的其余向量.
\end{exercise}
\noindent 来源：\href{https://github.com/yhwu-is/Linear-Algebra-Left-Undone/blob/new/%E8%AE%B2%E4%B9%89/%E4%B8%93%E9%A2%98/11%20%E6%9C%9D%E8%8A%B1%E5%A4%95%E6%8B%BE.tex}{L11-B-7}。
\begin{solve}
设基础解系含 $s\ge2$ 个向量。所求条件为 $(t_1,t_2)\ne(0,0)$。
若 $t_1\ne0$，取补充向量 $\alpha_s$。这 $s$ 个向量相对于原基 $\alpha_1,\ldots,\alpha_s$ 的坐标矩阵为下三角矩阵，对角元依次为 $t_1,\ldots,t_1,1$，行列式为 $t_1^{s-1}\ne0$；因此 $\beta_1,\ldots,\beta_{s-1},\alpha_s$ 是一组基础解系。
若 $t_1=0$、$t_2\ne0$，则 $\beta_j=t_2\alpha_{j+1}$，补入 $\alpha_1$ 即可。若 $t_1=t_2=0$，则题给全部 $\beta_j$ 为零向量，不能包含在基中。若允许 $s=1$，题给向量组为空，任意 $t_1,t_2$ 都可以，直接取基础解系 $\{\alpha_1\}$。
\end{solve}

\begin{exercise}\label{prob:L11-B-8}
已知线性方程组
        \[\begin{cases} \begin{aligned}
                    a_{11}x_1+a_{12}x_2+\cdots+a_{1,2n}x_{2n} & =0              \\
                    a_{21}x_1+a_{22}x_2+\cdots+a_{2,2n}x_{2n} & =0              \\
                                                              & \vdotswithin{=} \\
                    a_{n1}x_1+a_{n2}x_2+\cdots+a_{n,2n}x_{2n} & =0              \\
                \end{aligned}\end{cases}\]
        的一个基础解系为$(b_{11},b_{12},\ldots,b_{1,2n})^\mathrm{T},(b_{21},b_{22},\ldots,b_{2,2n})^\mathrm{T},(b_{n1},b_{n2},\ldots,b_{n,2n})^\mathrm{T}$，求解线性方程组
        \[\begin{cases} \begin{aligned}
                    b_{11}x_1+b_{12}x_2+\cdots+b_{1,2n}x_{2n} & =0              \\
                    b_{21}x_1+b_{22}x_2+\cdots+b_{2,2n}x_{2n} & =0              \\
                                                              & \vdotswithin{=} \\
                    b_{n1}x_1+b_{n2}x_2+\cdots+b_{n,2n}x_{2n} & =0              \\
                \end{aligned} \end{cases}.\]
        （注：本题的一般形式：设$A=(a_{ij})_{m\times n},m<n,r(A)=m$；$B=(b_{ij})_{n\times(n-m)},m<n,r(B)=n-m$.已知齐次线性方程组$A\vec{X}=\vec{0}$的解空间$N(A)=R(B)$（矩阵$B$的列空间），求齐次线性方程组$\displaystyle\sum_{i=1}^{n}b_{ij}y_i=0,j=1,2,\ldots,n-m$的一个基础解系.）
\end{exercise}
\noindent 来源：\href{https://github.com/yhwu-is/Linear-Algebra-Left-Undone/blob/new/%E8%AE%B2%E4%B9%89/%E4%B8%93%E9%A2%98/11%20%E6%9C%9D%E8%8A%B1%E5%A4%95%E6%8B%BE.tex}{L11-B-8}。
\begin{solve}
令原方程组$AX=0$，解系组成矩阵$B$. 则新方程组为$B^\mathrm{T}X=0$. 而$B^\mathrm{T}A^\mathrm{T}=(AB)^\mathrm{T}=0$，推测$A^\mathrm{T}$的列向量为基础解系. 而由维数公式，
          \[ r(A)=\dim V-\dim\ker(A) = 2n-n=n. \]
          故$A^\mathrm{T}$列空间维数为$n$，即$r_c(A^\mathrm{T})=n$. 而由题意$r(B)=n$（否则不是基础解系），故
          \[ \dim{\ker(B)}=\dim{V^{\prime}}-r(B)=2n-n=n=r_c(A^\mathrm{T}). \]
          猜想成立.
\end{solve}

\begin{exercise}\label{prob:L11-B-11}
设$A$是元素全为1的$n$阶方阵.
        \begin{enumerate}
            \item 求行列式$|aE+bA|$，其中$a,b$为实常数；

            \item 已知$1<r(aE+bA)<n$，试确定$a,b$满足的条件，并求下列线性子空间的维数：
                  \[W=\{X \mid (aE+bA)X=0,\enspace X\in\mathbf{R}^n\}.\]
        \end{enumerate}
\end{exercise}
\noindent 来源：\href{https://github.com/yhwu-is/Linear-Algebra-Left-Undone/blob/new/%E8%AE%B2%E4%B9%89/%E4%B8%93%E9%A2%98/11%20%E6%9C%9D%E8%8A%B1%E5%A4%95%E6%8B%BE.tex}{L11-B-11}。
\begin{solve}
\begin{enumerate}
                \item 容易计算$|aE=bA|=(a+nb)a^{n-1}$.

                \item 由$1<r(aE+bA)<n$知$|aE+bA|=0$. 故$a\neq 0$，且$a+nb=0$，此时$aE+bA$左上角的$n-1$阶子式
                      \[ \begin{vmatrix}
                              a+b    & b      & \cdots & b      \\
                              b      & a+b    & \cdots & b      \\
                              \vdots & \vdots & \ddots & \vdots \\
                              b      & b      & \cdots & a+b
                          \end{vmatrix}=(a+(n-1)b)a^{n-2}=\frac{a^{n-1}}{n}\neq 0, \]
                      故$\dim{W}=n-r(aE+bA)=n-(n-1)=1$.
            \end{enumerate}
\end{solve}

\begin{exercise}\label{prob:L11-B-13}
设$A=(a_{ij})_{m\times n}$，$\vec{b}$和$X$为$m$元列向量，$Y$为$n$元列向量，证明：
        \begin{enumerate}
            \item 若$AY=\vec{b}$有解，则$A^\mathrm{T}X=\vec{0}$的任一组解都满足$\vec{b}^\mathrm{T}X=0$；

            \item 方程组$AY=\vec{b}$有解的充要条件是方程组
\[\begin{pmatrix}
                          A^\mathrm{T} \\ \vec{b}^\mathrm{T}
                      \end{pmatrix}X=\begin{pmatrix}
                          \vec{0} \\ 1
                      \end{pmatrix}\]
无解（其中$\vec{0}$是$n$元零向量）.
        \end{enumerate}
\end{exercise}
\noindent 来源：\href{https://github.com/yhwu-is/Linear-Algebra-Left-Undone/blob/new/%E8%AE%B2%E4%B9%89/%E4%B8%93%E9%A2%98/11%20%E6%9C%9D%E8%8A%B1%E5%A4%95%E6%8B%BE.tex}{L11-B-13}。
\begin{solve}
\begin{enumerate}
                \item 由条件 $(A^\mathrm{T}X)^\mathrm{T}Y=X^\mathrm{T}(AY)=X^\mathrm{T}\vec{b}=0$，故 $\vec{b}^\mathrm{T}X=0$.
                \item 必要性由上一问显然成立.

                充分性：由条件可知，若 $X_0$ 使得 $A^\mathrm{T}X_0=\vec{0}$，则 $\vec{b}^\mathrm{T}X_0=0$，否则令 $X_1 =\dfrac{X_0}{\vec{b}^\mathrm{T}X_0}$，则 $X_1$ 是方程组 $\begin{pmatrix} A^\mathrm{T} \\ \vec{b}^\mathrm{T} \end{pmatrix}$ 的解，与条件矛盾.

                也就是说 $A^\mathrm{T}X=\vec{0}$ 与 $\begin{pmatrix} A^\mathrm{T} \\ \vec{b}^\mathrm{T} \end{pmatrix} X = \vec{0}$ 同解，从而 $r(A^\mathrm{T}) = r\begin{pmatrix} A^\mathrm{T} \\ \vec{b}^\mathrm{T} \end{pmatrix}$，这表明 $\vec{b}^\mathrm{T}$ 可由 $A^\mathrm{T}$ 的行向量（即 $A$ 的列向量）线性表示，因此 $AY=\vec{b}$ 有解.
            \end{enumerate}
\end{solve}

\begin{exercise}\label{prob:L11-B-14}
判断：设 $A$ 是复数域上 $m \times n$ 阶矩阵，则矩阵秩 $r\left(A^T A\right)=r(A)$.
\end{exercise}
\noindent 来源：\href{https://github.com/yhwu-is/Linear-Algebra-Left-Undone/blob/new/%E8%AE%B2%E4%B9%89/%E4%B8%93%E9%A2%98/11%20%E6%9C%9D%E8%8A%B1%E5%A4%95%E6%8B%BE.tex}{L11-B-14}。
\begin{solve}
假. 取$A=\begin{pmatrix}
                1-i & 1+i  \\
                1+i & -1+i
            \end{pmatrix}$，有$A^\mathrm{T}A=0$.
\end{solve}

\begin{exercise}\label{prob:L11-B-15}
证明：对于$m \times n$实矩阵$A$，方程$A^\mathrm{T}AX = A^\mathrm{T}\vec{b}$总是有解，且$A$为方阵时，$A^\mathrm{T}AX = \vec{0}$和$AX=\vec{0}$同解. 当$r(A)=n$时求其解，并证明$A(A^\mathrm{T}A)^{-1}A^\mathrm{T}$是幂等的对称矩阵.
\end{exercise}
\noindent 来源：\href{https://github.com/yhwu-is/Linear-Algebra-Left-Undone/blob/new/%E8%AE%B2%E4%B9%89/%E4%B8%93%E9%A2%98/11%20%E6%9C%9D%E8%8A%B1%E5%A4%95%E6%8B%BE.tex}{L11-B-15}。
\begin{solve}
由于 $\vec{b}$ 的任意性，$A^\mathrm{T}\vec{b}$ 可以取遍 $R(A^\mathrm{T})$ 中的任意元素，这里 $R(A^\mathrm{T})$ 表示 $A^\mathrm{T}$ 的列向量张成的空间. 因此我们只需证明 $R(A^\mathrm{T}) = R(A^\mathrm{T}A)$.

            $A^\mathrm{T}A$ 的每一列都是 $A^\mathrm{T}$ 的列向量的线性组合，因此 $R(A^\mathrm{T}A) \subseteq R(A^\mathrm{T})$. 而根据原讲义“线性方程组理论与秩不等式”处 的结论，有 $r(A^\mathrm{T}A) = r(A)$，因此 $A^\mathrm{T}A$ 的列向量组和 $A^\mathrm{T}$ 的列秩相等，即 $\dim R(A^\mathrm{T}) = \dim R(A^\mathrm{T}A)$. 从而 $R(A^\mathrm{T}) = R(A^\mathrm{T}A)$，命题得证.
\end{solve}

\begin{exercise}\label{prob:L11-B-16}
设$A,B,C$为$n$阶实方阵，且$BAA^\mathrm{T}=CAA^\mathrm{T}$，证明：$BA=CA$.
\end{exercise}
\noindent 来源：\href{https://github.com/yhwu-is/Linear-Algebra-Left-Undone/blob/new/%E8%AE%B2%E4%B9%89/%E4%B8%93%E9%A2%98/11%20%E6%9C%9D%E8%8A%B1%E5%A4%95%E6%8B%BE.tex}{L11-B-16}。
\begin{solve}
记 $D=B-C$，由题设 $DAA^{\mathrm T}=O$，右乘 $D^{\mathrm T}$ 得
\[(DA)(DA)^{\mathrm T}=O.\]
左端第 $i$ 个对角元是实矩阵 $DA$ 第 $i$ 行元素的平方和，故各行都为零，即 $DA=O$。因此 $BA=CA$。
\end{solve}

\begin{exercise}\label{prob:L11-B-17}
设$A$为数域$\mathbf{F}$上的$n$阶方阵，又设线性空间$\mathbf{F^n}$的两个子空间为$W_1=\{X\in\mathbf{F}^n \mid AX=\vec{0}\}$，$W_2=\{X\in\mathbf{F}^n \mid (A-E)X=\vec{0}\}$. 证明：$A^2=A \iff \mathbf{F}^n=W_1\oplus W_2$.
\end{exercise}
\noindent 来源：\href{https://github.com/yhwu-is/Linear-Algebra-Left-Undone/blob/new/%E8%AE%B2%E4%B9%89/%E4%B8%93%E9%A2%98/11%20%E6%9C%9D%E8%8A%B1%E5%A4%95%E6%8B%BE.tex}{L11-B-17}。
\begin{solve}
若 $A^2=A$，则任意 $x\in\mathbf F^n$ 可写成 $x=(x-Ax)+Ax$。其中 $A(x-Ax)=0$，而 $(A-I)Ax=0$，故两项分别属于 $W_1,W_2$。若 $z\in W_1\cap W_2$，则同时有 $Az=0$、$Az=z$，所以 $z=0$，从而是直和。
反过来，若 $\mathbf F^n=W_1\oplus W_2$，写 $x=u+v$，其中 $Au=0$、$Av=v$，则 $A^2x=Av=v=Ax$。因为这对任意 $x$ 成立，故 $A^2=A$。
\end{solve}

\begin{exercise}\label{prob:L11-B-18}
$n$阶方阵$A,B$满足$AB=BA$，证明：$r(AB)+r(A+B) \leqslant r(A)+r(B)$.
\end{exercise}
\noindent 来源：\href{https://github.com/yhwu-is/Linear-Algebra-Left-Undone/blob/new/%E8%AE%B2%E4%B9%89/%E4%B8%93%E9%A2%98/11%20%E6%9C%9D%E8%8A%B1%E5%A4%95%E6%8B%BE.tex}{L11-B-18}。
\begin{solve}
方法一. 用分块矩阵的方法，我们知道
          \[ \begin{pmatrix}
                  A & O \\
                  O & B
              \end{pmatrix}
              \rightarrow
              \begin{pmatrix}
                  A & O \\
                  A & B
              \end{pmatrix}
              \rightarrow
              \begin{pmatrix}
                  A & A   \\
                  A & A+B
              \end{pmatrix}. \]
          结合$AB=BA$，我们知道
          \[ \begin{pmatrix}
                  A & A   \\
                  A & A+B
              \end{pmatrix}
              \underbrace{
                  \begin{pmatrix}
                      A+B & O \\
                      -A  & E
                  \end{pmatrix}
              }_{\text{非广义初等变换，难以想到}}
              =
              \begin{pmatrix}
                  AB & A   \\
                  O  & A+B
              \end{pmatrix}. \]
          于是
          \[ r(A)+r(B)=r
              \begin{pmatrix}
                  A & O \\
                  O & B
              \end{pmatrix}=
              \begin{pmatrix}
                  A & A   \\
                  A & A+B
              \end{pmatrix}\geqslant r
              \begin{pmatrix}
                  AB & A   \\
                  O  & A+B
              \end{pmatrix}\geqslant
              r(AB)+r(A+B). \]
          方法二. 设方程组$AX=0$与$BX=0$的解空间分别是$V_1, V_2$，方程组$ABX=BAX=0$与$(A+B)X=0$的解空间分别为$W_1, W_2$，则$V_1\subseteq W_1, V_2\subseteq W_1$，从而$V_1+V_2\subseteq W_1$，同时$V_1\cap V_2\subseteq W_1$，同时$V_1\cap V_2\subseteq W_2$，利用维数公式就有
          \[ \dim V_1+\dim V_2=\dim(V_1+V_2)+\dim(V_1\cap V_2)\leqslant \dim W_1+\dim W_2. \]
          即
          \[ (n-r(A))+(n-r(B))\leqslant (n-r(AB))+(n-r(A+B)). \]
          化简便知$r(A)+r(B)\geqslant r(AB)+r(A+B)$.
\end{solve}

\begin{exercise}\label{prob:L11-B-19}
请按序证明以下结论：
        \begin{enumerate}
            \item $A,B$分别是$s \times m,m \times n$矩阵，则$ABX=\vec{0}$与$BX=\vec{0}$同解的充要条件是$r(AB)=r(B)$；

            \item $A,B$分别是$s \times m,m \times n$矩阵，且$r(AB)=r(B)$，则对任意的$n \times t$矩阵都有$r(ABC)=r(BC)$；

            \item 设$A$是$n$阶方阵，则存在正整数$k$使得$r(A^k)=r(A^{k+1})=r(A^{k+2})=\cdots$，且对任意正整数$m$，有$r(A^n)=r(A^{n+m})$.
        \end{enumerate}
\end{exercise}
\noindent 来源：\href{https://github.com/yhwu-is/Linear-Algebra-Left-Undone/blob/new/%E8%AE%B2%E4%B9%89/%E4%B8%93%E9%A2%98/11%20%E6%9C%9D%E8%8A%B1%E5%A4%95%E6%8B%BE.tex}{L11-B-19}。
\begin{solve}
\begin{enumerate}
                \item 必要性：$ABX=0$与$BX=0$同解可知它们基础解系所含向量个数相同，即
                      \[ n-r(AB)=n-r(B)\implies r(AB)=r(B). \]
                      充分性：由必要性，当$r(AB)=r(B)$时，$ABX=0$与$BX=0$的基础解系所含向量个数相同，而$BX=0$的解都是$ABX=0$的解，所以$ABX=0$与$BX=0$同解.

                \item $r(AB)=r(B)$说明$ABX=0$与$BX=0$同解，用$CX$代替$X$就得到$ABCX=0$与$BCX=0$同解，从而有$r(ABC)=r(BC)$.\\
                      \textbf{推论.}设$A$是一个方阵，且存在正整数$k$使得$r(A^{k+1})=r(A^k)$，递推就有
                      \[ r(A^k)=r(A^{k+1})=r(A^{k+2})=\cdots. \]

                \item 当$A$可逆时，结论是显然的. 当$A$不可逆时，有$r(A)\leqslant n-1$，现在考虑$n+1$个矩阵$A,A^2,\ldots,A^{n+1}$，有
                      \[ n-1\geqslant r(A)\geqslant r(A^2)\geqslant\cdots\geqslant r(A^n)\geqslant r(A^{n+1})\geqslant 0. \]
                      这$n+1$个矩阵的秩只能从$0,1,\ldots,n-1$这$n$个数中取，所以必有两个矩阵的秩相同，即存在$m(1\leqslant m\leqslant n)$使得$r(A^m)=r(A^{m+1})$，由上面的推论可得：
                      \[ r(A^m)=r(A^{m+1})=\cdots=r(A^{n})=r(A^{n+1})=\cdots. \]
                      特别的，对于任意的正整数$k$有$r(A^n)=r(A^{n+k})$.
            \end{enumerate}
\end{solve}

\begin{exercise}\label{prob:L11-B-20}
如果齐次线性方程组\[\begin{cases}
                x_1+x_2+bx_3-x_4+x_5=0   \\
                2x_1+3x_2+x_3+x_4-2x_5=0 \\
                x_2+ax_3+3x_4-4x_5=0     \\
                -3x_1-3x_2-3bx_3+bx_4+(a+2)x_5=0
            \end{cases}\]的解空间是3维的，试求$a,b$的值与解空间的基. 解空间可能为2维吗？
\end{exercise}
\noindent 来源：\href{https://github.com/yhwu-is/Linear-Algebra-Left-Undone/blob/new/%E8%AE%B2%E4%B9%89/%E4%B8%93%E9%A2%98/11%20%E6%9C%9D%E8%8A%B1%E5%A4%95%E6%8B%BE.tex}{L11-B-20}。
\begin{solve}
对系数矩阵施行 $R_2\leftarrow R_2-2R_1$、$R_3\leftarrow R_3-R_2$（第二步使用更新后的 $R_2$）、$R_4\leftarrow R_4+3R_1$，得到
\[\begin{pmatrix}1&1&b&-1&1\\0&1&1-2b&3&-4\\0&0&a+2b-1&0&0\\0&0&0&b-3&a+5\end{pmatrix}.\]
前两行恒线性无关。因此解空间三维等价于后两行同时为零，即 $a=-5,b=3$。此时令自由变量为 $x_3,x_4,x_5$，得到
\[x_2=5x_3-3x_4+4x_5,\qquad x_1=-8x_3+4x_4-5x_5.\]
一组基础解系为
\[(-8,5,1,0,0)^{\mathrm T},\quad(4,-3,0,1,0)^{\mathrm T},\quad(-5,4,0,0,1)^{\mathrm T}.\]
解空间也可能为二维。最后两行若同时非零则相互独立；而第四行全零时必有 $a=-5,b=3$，第三行也随之为零。因此秩恰为 $3$ 的充要条件是 $a+2b=1$ 且 $b\ne3$。例如 $a=1,b=0$ 即满足要求。
\end{solve}

\begin{exercise}\label{prob:L11-B-21}
设$W_1,W_2$分别为$n$元齐次线性方程组$AX=\vec{0}$和$BX=\vec{0}$的解空间，试构造两个$n$元齐次线性方程组，使它们的解空间分别为$W_1 \cap W_2$和$W_1+W_2$.
\end{exercise}
\noindent 来源：\href{https://github.com/yhwu-is/Linear-Algebra-Left-Undone/blob/new/%E8%AE%B2%E4%B9%89/%E4%B8%93%E9%A2%98/11%20%E6%9C%9D%E8%8A%B1%E5%A4%95%E6%8B%BE.tex}{L11-B-21}。
\begin{solve}
$W_1\cap W_2:\begin{pmatrix}
                  A \\ B
              \end{pmatrix}X=0$. \qquad
          $W_1+ W_2:\begin{cases}
                  AX = 0 \\
                  BX = 0
              \end{cases}$.
\end{solve}

\begin{exercise}\label{prob:L11-B-22}
已知方程组$\begin{cases}
                x_1+x_2+ax_3+x_4=1 \\ -x_1+x_2-x_3+bx_4=2 \\ 2x_1+x_2+x_3+x_4=c
            \end{cases}$与$\begin{cases}
                x_1+x_4=-1 \\ x_2-2x_4=d \\ x_3+x_4=e
            \end{cases}$同解，求$a,b,c,d,e$.
\end{exercise}
\noindent 来源：\href{https://github.com/yhwu-is/Linear-Algebra-Left-Undone/blob/new/%E8%AE%B2%E4%B9%89/%E4%B8%93%E9%A2%98/11%20%E6%9C%9D%E8%8A%B1%E5%A4%95%E6%8B%BE.tex}{L11-B-22}。
\begin{solve}
方法一. 令$x_4=t$，则方程组$\begin{cases}
                x_1+x_4=-1 \\
                x_2-2x_4=d \\
                x_3+x_4=e
            \end{cases}$的一般解为$\begin{cases}
                x_1=-1-t \\
                x_2=d+2t \\
                x_3=e-t  \\
                x_4=t
            \end{cases}$.

        代入方程组$\begin{cases}
                x_1+x_2+ax_3+x_4=1  \\
                -x_1+x_2-x_3+bx_4=2 \\
                2x_1+x_2+x_3+x_4=c
            \end{cases}$可得$\begin{cases}
                (2-a)t=2-d-ae \\
                (b+4)t=1-d+e  \\
                0=c-d-e+2
            \end{cases}$. 由$t$的任意性，可得$a=2,b=-4$. 从而$\begin{cases}
                0=2-d-2e \\
                0=1-d+e  \\
                0=c-d-e+2
            \end{cases}$，解得$d=\dfrac{4}{3},e=\dfrac{1}{3},c=-\dfrac{1}{3}$.

        方法二. 由于
        \[\begin{pmatrix}
                A & b \\
                B & d
            \end{pmatrix}=
            \begin{pmatrix}
                1  & 1 & a  & 1  & 1  \\
                -1 & 1 & -1 & b  & 2  \\
                2  & 1 & 1  & 1  & c  \\
                1  & 0 & 0  & 1  & -1 \\
                0  & 1 & 0  & -2 & d  \\
                0  & 0 & 1  & 1  & e
            \end{pmatrix}\rightarrow
            \begin{pmatrix}
                1 & 0 & 0 & 0     & 0            \\
                0 & 0 & 0 & b-a+6 & 3-2d+(1-a)e  \\
                0 & 0 & 0 & 2a-4  & c-2+d(2a-1)e \\
                0 & 0 & 0 & a-2   & d-2+ae       \\
                0 & 1 & 0 & 0     & 0            \\
                0 & 0 & 1 & 0     & 0
            \end{pmatrix}.\]
        易知$r(B)=3$，由$r\begin{pmatrix}
                A & b \\
                B & d
            \end{pmatrix}=r(B)$可得$\begin{cases}
                a-2=0          \\
                2a-4=0         \\
                b-a+6=0        \\
                3-2d+(1-a)e=0  \\
                c-2+d(2a-1)e=0 \\
                d-2+ae=0
            \end{cases}$，解得$\begin{cases}
                a=2          \\
                b=-4         \\
                c=-\dfrac 13 \\[1ex]
                d=\dfrac 43  \\[1ex]
                e=\dfrac 13
            \end{cases}$.
\end{solve}

\begin{exercise}\label{prob:L11-B-23}
设有两个非齐次线性方程组 (1) 和 (2)，它们的通解分别是$X=\gamma+t_1\eta_1+t_2\eta_2, X=\delta+k_1\xi_1+k_2\xi_2$. 其中$\gamma=(5,-3,0,0)^\mathrm{T},\eta_1=(-6,5,1,0)^\mathrm{T},\eta_2=(-5,4,0,1)^\mathrm{T},\delta=(-11,3,0,0)^\mathrm{T},\xi_1=(8,-1,1,0)^\mathrm{T},\xi_2=(10,-2,0,1)^\mathrm{T}$，求这两个方程组的公共解.
\end{exercise}
\noindent 来源：\href{https://github.com/yhwu-is/Linear-Algebra-Left-Undone/blob/new/%E8%AE%B2%E4%B9%89/%E4%B8%93%E9%A2%98/11%20%E6%9C%9D%E8%8A%B1%E5%A4%95%E6%8B%BE.tex}{L11-B-23}。
\begin{solve}
注意如果$X$是两个方程组的公共解，这等价于存在$t_1,t_2,k_1,k_2$使得
          \[ X=\gamma+t_1\eta_1+t_2\eta_2=\delta+k_1\xi_1+k_2\xi_2. \]
          从而对应有
          \[ t_1\eta_1+t_2\eta_2-k_1\xi_1-k_2\xi_2=\gamma-\delta. \]
          将$t_1,t_2,k_1,k_2$设为未知量，对上述方程组的增广矩阵进行初等行变换，可得
          \[ (\eta_1,\eta_2,-\xi_1,-\xi_2,\delta-\gamma)=
              \begin{pmatrix}
                  -6 & -5 & -8 & -10 & -16 \\
                  5  & 4  & 1  & 2   & 6   \\
                  1  & 0  & -1 & 0   & 0   \\
                  0  & 1  & 0  & -1  & 0
              \end{pmatrix}\rightarrow
              \begin{pmatrix}
                  1 & 0 & -1 & 0  & 0  \\
                  0 & 1 & 0  & -1 & 0  \\
                  0 & 0 & 1  & 1  & 1  \\
                  0 & 0 & 0  & 1  & 2.
              \end{pmatrix} \]
          故$(t_1,t_2,k_1,k_2)^\mathrm{T}$有唯一解$(-1,2,-1,2)^\mathrm{T}$. 因此公共解为
          \[ X=\gamma-\eta_1+2\eta_2=\delta-\xi_1+2\xi_2=(1,0,-1,2)\mathrm{T}. \]
\end{solve}

\begin{exercise}\label{prob:L11-B-24}
若方程组$\begin{cases}
                x_1+x_2+x_3=0   \\
                x_1+2x_2+ax_3=0 \\
                x_1+4x_2+a^2x_3=0
            \end{cases}$与$x_1+2x_2+x_3=a-1$有公共解，求$a$的值及所有公共解.
\end{exercise}
\noindent 来源：\href{https://github.com/yhwu-is/Linear-Algebra-Left-Undone/blob/new/%E8%AE%B2%E4%B9%89/%E4%B8%93%E9%A2%98/11%20%E6%9C%9D%E8%8A%B1%E5%A4%95%E6%8B%BE.tex}{L11-B-24}。
\begin{solve}
方程组 (1)、(2) 有公共解，即方程组$\begin{cases}
                x_1+x_2+x_3=0,     \\
                x_1+2x_2+ax_3=0,   \\
                x_1+4x_2+a^2x_3=0, \\
                x_1+2x_2+x_3=a-1
            \end{cases}$有解，对其增广矩阵进行初等行变换，
        \begin{align*}
            (A,\vec{b})= & \begin{pmatrix}
                           1 & 1 & 1   & 0   \\
                           1 & 2 & a   & 0   \\
                           1 & 4 & a^2 & 0   \\
                           1 & 2 & 1   & a-1
                       \end{pmatrix}\rightarrow
            \begin{pmatrix}
                1 & 1 & 1     & 0   \\
                0 & 1 & a-1   & 0   \\
                0 & 3 & a^2-1 & 0   \\
                0 & 1 & 0     & a-1
            \end{pmatrix}\rightarrow            \\
                     & \begin{pmatrix}
                           1 & 1 & 1        & 0   \\
                           0 & 1 & a-1      & 0   \\
                           0 & 0 & a^2-3a+2 & 0   \\
                           0 & 0 & 1-a      & a-1
                       \end{pmatrix}\rightarrow
            \begin{pmatrix}
                1 & 1 & 1   & 0          \\
                0 & 1 & a-1 & 0          \\
                0 & 0 & 1-a & a-1        \\
                0 & 0 & 0   & (a-1)(a-2)
            \end{pmatrix}.
        \end{align*}
        当$a\neq 1$且$a\neq 2$时，方程组 (1)、(2) 没有公共解.

        当$a=1$时，
\[(A,\vec{b})\rightarrow\begin{pmatrix}
                1 & 1 & 1 & 0 \\
                0 & 1 & 0 & 0 \\
                0 & 0 & 0 & 0 \\
                0 & 0 & 0 & 0
            \end{pmatrix}\rightarrow\begin{pmatrix}
                1 & 0 & 1 & 0 \\
                0 & 1 & 0 & 0 \\
                0 & 0 & 0 & 0 \\
                0 & 0 & 0 & 0
            \end{pmatrix}\]
，因为$r(A)=r(A,\vec{b})=2$，所以方程组 (1)、(2) 有公共解，公共解为$X=k\begin{pmatrix} -1 \\ 0  \\ 1 \end{pmatrix}$（其中$k$为任意常数）.

        当$a=2$时，
\[(A,\vec{b})\rightarrow\begin{pmatrix}
                1 & 1 & 1  & 0 \\
                0 & 1 & 1  & 0 \\
                0 & 0 & -1 & 1 \\
                0 & 0 & 0  & 0
            \end{pmatrix}\rightarrow\begin{pmatrix}
                1 & 0 & 0 & 0  \\
                0 & 1 & 0 & 1  \\
                0 & 0 & 1 & -1 \\
                0 & 0 & 0 & 0
            \end{pmatrix}\]
，方程组 (1)、(2) 有唯一的公共解为$X=\begin{pmatrix}
                0 \\
                1 \\
                -1
            \end{pmatrix}$.
\end{solve}

\begin{exercise}\label{prob:L11-C-1}
用方程组的理论证明：一个$n$次多项式不可能有多于$n$个不同的根.
\end{exercise}
\noindent 来源：\href{https://github.com/yhwu-is/Linear-Algebra-Left-Undone/blob/new/%E8%AE%B2%E4%B9%89/%E4%B8%93%E9%A2%98/11%20%E6%9C%9D%E8%8A%B1%E5%A4%95%E6%8B%BE.tex}{L11-C-1}。
\begin{solve}
令$f(x)=a_nx^n+\cdots+a_1x+a_0$为一个$n$次多项式，设$\lambda_1,\ldots,\lambda_{n+1}$为$f(x)$的$n+1$个不同的根，则有齐次线性方程组$\begin{cases} \begin{aligned}
                \lambda_1^n a_n+\cdots+\lambda_1a_1+a_0         & =0,               \\
                \lambda_2^n a_n+\cdots+\lambda_2a_1+a_0         & =0,               \\
                                                                & \vdotswithin{ = } \\
                \lambda_{n+1}^n a_n+\cdots+\lambda_{n+1}a_1+a_0 & =0.
            \end{aligned} \end{cases}$. 其系数矩阵为
\[A=\begin{pmatrix}
            \lambda_1^n     & \lambda_1^{n-1}     & \cdots & 1      \\
            \lambda_2^n     & \lambda_2^{n-1}     & \cdots & 1      \\
            \vdots          & \vdots              & \ddots & \vdots \\
            \lambda_{n+1}^n & \lambda_{n+1}^{n-1} & \cdots & 1
        \end{pmatrix}\]
. 显然$|A|\neq 0$，从而$a_n=a_{n-1}=\cdots=a_1=a_0=0$，这与$f(x)$是$n$次多项式矛盾.
\end{solve}

\begin{exercise}\label{prob:L11-C-3}
已知$A$是$n$阶对称矩阵，$\beta$为$n$元非零列向量，$B=\begin{pmatrix}
                A & \beta \\ \beta^\mathrm{T} & 0
            \end{pmatrix}$，证明：
        \begin{enumerate}
            \item 若$r(A)=n$，则$B$可逆的充要条件是$\beta^\mathrm{T}A^{-1}\beta \neq \vec{0}$；

            \item 若$r(A)=r$，则$r(B)=r$的充要条件是方程组$\begin{cases}
                          AX=\beta \\ \beta^\mathrm{T}X=\vec{0}
                      \end{cases}$有解；

            \item 若$r(A)=n-1$，则$B$可逆的充要条件是$AX=\beta$无解.
        \end{enumerate}
\end{exercise}
\noindent 来源：\href{https://github.com/yhwu-is/Linear-Algebra-Left-Undone/blob/new/%E8%AE%B2%E4%B9%89/%E4%B8%93%E9%A2%98/11%20%E6%9C%9D%E8%8A%B1%E5%A4%95%E6%8B%BE.tex}{L11-C-3}。
\begin{solve}
\begin{enumerate}
                \item 由于$r(A)=n$，所以$A$可逆，由打洞原理可知$|B|=|A|(0-\beta' A^{-1}\beta)=-|A|\beta' A^{-1}\beta$，从而$B$可逆的充要条件为$\beta' A^{-1}\beta\neq 0$.

                \item 必要性. 由于
                      \[ r=r(A)\leqslant r(A,\beta)\leqslant r\begin{pmatrix} A & \beta \\ \beta' & 0 \end{pmatrix}=r(B)=r \]
                      再结合$A'=A$，可知
                      \[ r\begin{pmatrix} A & \beta \\ \beta' & 0 \end{pmatrix} =r(A,\beta)=r\begin{pmatrix} A \\ \beta' \end{pmatrix} \]
                      于是由原讲义定理 15.1可知方程组$\begin{pmatrix} A \\ \beta' \end{pmatrix}X=\begin{pmatrix} \beta \\ 0 \end{pmatrix}$有解，即$\begin{cases} AX=\beta, \\ \beta' X=0. \end{cases}$有解.

                      充分性. 由于$\begin{cases} AX=\beta \\ \beta' X=0 \end{cases}$有解，从而$AX=\beta$也有解，即有$r(A)=r(A,\beta)$. 另外$\begin{cases} AX=\beta \\ \beta' X=0 \end{cases}$有解也说明$\begin{pmatrix} A \\ \beta' \end{pmatrix}X=\begin{pmatrix} \beta \\ 0 \end{pmatrix}$有解，于是结合原讲义定理 15.1可知$r\begin{pmatrix} A & \beta \\ \beta' & 0 \end{pmatrix}=r\begin{pmatrix} A \\ \beta' \end{pmatrix}$. 而显然$r\begin{pmatrix} A \\ \beta' \end{pmatrix}=r(A,\beta)=r(A)$，于是$r(B)=r\begin{pmatrix} A & \beta \\ \beta' & 0 \end{pmatrix}=r$.

                \item 必要性. 若$B$可逆，则$B$的行向量组线性无关，从而$r(A,\beta)=n$，又由于$r(A)=n-1$，所以$r(A,\beta)=r(A)+1$，从而由原讲义定理 15.1知$AX=\beta$无解.

                      充分性. 由于$r(A)=n-1$，若$AX=\beta$无解，则由原讲义定理 15.1可知$r(A,\beta)=r(A)+1=n$，于是
                      \[ r(B)=r\begin{pmatrix}
                              A      & \beta \\
                              \beta' & 0
                          \end{pmatrix}\geqslant r(A,\beta)=n \]
                      若$r(B)=n$，则$r(B)=r\begin{pmatrix}
                              A      & \beta \\
                              \beta' & 0
                          \end{pmatrix}=r(A,\beta)$，取转置结合$A'=A$有
\[r\begin{pmatrix}
                              A      & \beta \\
                              \beta' & 0
                          \end{pmatrix}=r\begin{pmatrix}
                              A \\
                              \beta'
                          \end{pmatrix}\]
，再次结合原讲义定理 15.1可知方程组$\begin{cases}
                              AX=\beta, \\
                              \beta' X=0.
                          \end{cases}$有解，特别地，$AX=\beta$也有解，这就与已知产生了矛盾. 所以$r(B)=n+1$，即$B$可逆.
            \end{enumerate}
\end{solve}

\begin{exercise}\label{prob:L11-C-4}
讨论$b_1,b_2,\ldots,b_n\enspace(n \geqslant 2)$满足什么条件时，下列方程组
        \[\begin{cases} \begin{aligned}
                    x_1+x_2     & =b_1            \\
                    x_2+x_3     & =b_2            \\
                                & \vdotswithin{=} \\
                    x_{n-1}+x_n & =b_{n-1}        \\
                    x_n+x_1     & =b_n
                \end{aligned} \end{cases}\]有解，并求解.
\end{exercise}
\noindent 来源：\href{https://github.com/yhwu-is/Linear-Algebra-Left-Undone/blob/new/%E8%AE%B2%E4%B9%89/%E4%B8%93%E9%A2%98/11%20%E6%9C%9D%E8%8A%B1%E5%A4%95%E6%8B%BE.tex}{L11-C-4}。
\begin{solve}
\begin{enumerate}
                \item 当$n$为偶数时，将增广矩阵$\bar{A}$的第$i\enspace(i=1,2,\ldots,n-1)$行乘以$(-1)^i$加到最后一行，得
                      \[ \bar{A}\rightarrow\begin{pmatrix}
                              1      & 1      & \cdots & 0      & 0      & b_1                                  \\
                              0      & 1      & \cdots & 0      & 0      & b_2                                  \\
                              0      & 0      & \cdots & 0      & 0      & b_3                                  \\
                              \vdots & \vdots & \ddots & \vdots & \vdots & \vdots                               \\
                              0      & 0      & \cdots & 1      & 1      & b_{n-1}                              \\
                              0      & 0      & \cdots & 0      & 0      & \displaystyle\sum_{i=1}^{n}(-1)^ib_i
                          \end{pmatrix}. \]
                      故当$\displaystyle\sum_{i=1}^{n}(-1)^ib_i=0$时，方程组有无穷多解，一般解为
                      \[\begin{cases} \begin{aligned}
                                  x_1     & =\displaystyle\sum_{i=1}^{n-1}(-1)^{i-1}b_i+(-1)^1x_n, \\
                                  x_2     & =\displaystyle\sum_{i=2}^{n-1}(-1)^{i-2}b_i+(-1)^2x_n, \\
                                          & \vdotswithin{ = }                                      \\
                                  x_{n-2} & =b_{n-2}-b_{n-1}+(-1)^{n-2}x_n,                        \\
                                  x_{n-1} & =b_{n-1}+(-1)^{n-1}x_n,                                \\
                              \end{aligned} \end{cases}\]
                      其中$x_n$为自由未知量.

                \item 当$n$为奇数时，有
                      \[ \bar{A}\rightarrow\begin{pmatrix}
                              1      & 1      & \cdots & 0      & 0      & b_1                                        \\
                              0      & 1      & \cdots & 0      & 0      & b_2                                        \\
                              0      & 0      & \cdots & 0      & 0      & b_3                                        \\
                              \vdots & \vdots & \ddots & \vdots & \vdots & \vdots                                     \\
                              0      & 0      & \cdots & 1      & 1      & b_{n-1}                                    \\
                              0      & 0      & \cdots & 0      & 2      & b_n+\displaystyle\sum_{i=1}^{n-1}(-1)^ib_i
                          \end{pmatrix}. \]
                      此时无论$b_1,b_2,\ldots,b_n\enspace(n\geqslant 2)$取何值，方程组都有唯一解为
                      \[\begin{cases} \begin{aligned}
                                  x_1     & =\displaystyle\sum_{i=1}^{n-1}(-1)^{i-1}b_i+(-1)^{n-1}x_n, \\
                                  x_2     & =\displaystyle\sum_{i=2}^{n-1}(-1)^{i-2}b_i+(-1)^{n-2}x_n, \\
                                          & \vdotswithin{ = }                                          \\
                                  x_{n-2} & =b_{n-2}-b_{n-1}+(-1)^2x_n,                                \\
                                  x_{n-1} & =b_{n-1}+(-1)^1x_n,                                        \\
                                  x_n     & =\frac 12b_n+\displaystyle\sum_{i=1}^{n-1}(-1)^ib_i.
                              \end{aligned} \end{cases}\]
            \end{enumerate}
\end{solve}

\begin{exercise}\label{prob:L11-C-5}
已知$A$是$s \times n$矩阵，$B$是$m \times n$矩阵，$X,\vec{a},\vec{b}$分别是$n,s,m$元列向量，证明：
        \begin{enumerate}
            \item 齐次线性方程组$AX=\vec{0}$和$BX=\vec{0}$同解的充要条件是$A$与$B$行向量等价（列向量不一定）；

            \item 齐次线性方程组$AX=\vec{0}$和$BX=\vec{0}$解空间分别为$V_1,V_2$，证明：$V_1 \subseteq V_2$的充要条件是存在$m \times s$矩阵$C$使得$B=CA$；

            \item 线性方程组$AX=\vec{a}$的解都是$BX=\vec{b}$的解的充要条件是增广矩阵$(B,\vec{b})$的每个行向量都可以被$(A,\vec{a})$的行向量线性表示；

            \item 线性方程组$AX=\vec{a}$与$BX=\vec{b}$同解的充要条件是$(A,\vec{a})$与$(B,\vec{b})$行向量等价.
        \end{enumerate}
\end{exercise}
\noindent 来源：\href{https://github.com/yhwu-is/Linear-Algebra-Left-Undone/blob/new/%E8%AE%B2%E4%B9%89/%E4%B8%93%E9%A2%98/11%20%E6%9C%9D%E8%8A%B1%E5%A4%95%E6%8B%BE.tex}{L11-C-5}。
\begin{solve}
\begin{enumerate}
                \item 必要性. 设$A$的行向量为$\alpha_1,\alpha_2,\ldots,\alpha_s$，$B$的行向量为$\beta_1,\beta_2,\ldots,\beta_m$，由$AX=0$与$BX=0$同解得$AX=0$与$\begin{pmatrix}A \\ B\end{pmatrix}=0$同解，从而系数矩阵的秩相同，即行向量$\alpha_1,\alpha_2,\ldots,\alpha_s$与$\alpha_1,\alpha_2,\ldots,\alpha_s,\beta_1,\beta_2,\ldots,\beta_m$秩相同. 

                      充分性可由必要性的证明得到. 但是由于$A,B$列向量的维数都可能不同，所以不存在等价关系.

                \item 提示. $V_1\subseteq V_2$等价于$AX=0$与$\begin{pmatrix}
                              A \\
                              B
                          \end{pmatrix}X=0$同解.

                \item 必要性. 由题意可知$AX=a$与$\begin{pmatrix} A \\ B \end{pmatrix}X=\begin{pmatrix} a \\ b \end{pmatrix}$同解，从而它们的增广矩阵秩相同，即$r(A,a)=r\begin{pmatrix}
                              A & a \\
                              B & b
                          \end{pmatrix}$，可知$(A,a)$与$\begin{pmatrix}
                              A & a \\
                              B & b
                          \end{pmatrix}$的行向量组等价，从而$(B,b)$的每一个行向量都可以由$(A,a)$的行向量线性表出.\\
                      充分性可由必要性得到.

                \item 证明留给读者.
            \end{enumerate}
\end{solve}

\begin{exercise}\label{prob:L11-C-6}
设$A=(a_{ij})_{n \times n}$是一个$n$级矩阵，证明：
        \begin{enumerate}
            \item 若$A$为复矩阵，且$|a_{ii}|>\displaystyle\sum_{j \neq i}|a_{ij}|$，那么$|A|\neq 0$；

            \item 若$A$为实矩阵，且$a_{ii}>\displaystyle\sum_{j \neq i}|a_{ij}|$，那么$|A|>0$；

            \item （推论）存在充分大的实数$M$，使得$t>M$时，$tE+A$可逆.
        \end{enumerate}
\end{exercise}
\noindent 来源：\href{https://github.com/yhwu-is/Linear-Algebra-Left-Undone/blob/new/%E8%AE%B2%E4%B9%89/%E4%B8%93%E9%A2%98/11%20%E6%9C%9D%E8%8A%B1%E5%A4%95%E6%8B%BE.tex}{L11-C-6}。
\begin{solve}
\begin{enumerate}
                \item \label{item:15:B:5:1}
                      采用反证法. 设 $\lvert A \rvert = 0$，则线性方程组 $AX = 0$ 有非零解，设为 $X_0 = (x_1, x_2, \ldots, x_n)^{\mathrm{T}}$，记
                      \[\lvert x_k \rvert = \max \{\lvert x_1 \rvert, \lvert x_2 \rvert, \ldots, \lvert x_n \rvert\}.\]
                      由 $X_0 \neq 0$ 可知 $\lvert x_k \rvert > 0$，考虑 $AX = 0$ 的第 $k$ 个方程，有 $\displaystyle\sum_{j=1}^n a_{kj}x_j = 0$，于是
                      \[\lvert a_{kk} \rvert \lvert x_k \rvert = \lvert -\displaystyle\sum_{j \neq k}a_{kj}x_j \rvert \leqslant \sum_{j \neq k}\lvert a_{kj} \rvert \lvert x_k \rvert.\]
                      约去 $\lvert x_k \rvert$ 后可得 $\lvert a_{kk} \rvert \leqslant \displaystyle\sum_{j \neq k} \lvert a_{kj} \rvert$，这与条件矛盾. 所以 $\lvert A \rvert \neq 0$.

                \item 构造实函数
                      \[f(t) = \begin{vmatrix}
                              a_{11}  & ta_{12} & ta_{13} & \cdots & ta_{1n} \\
                              ta_{21} & a_{22}  & ta_{23} & \cdots & ta_{2n} \\
                              ta_{31} & ta_{32} & a_{33}  & \cdots & ta_{3n} \\
                              \vdots  & \vdots  & \vdots  & \ddots & \vdots  \\
                              ta_{n1} & ta_{n2} & ta_{n3} & \cdots & a_{nn}  \\
                          \end{vmatrix}\]
                      由于 $A$ 是实矩阵，所以 $f(t)$ 是关于 $t$ 的一个实系数多项式（连续）函数，同时
                      \[f(0) = a_{11}a_{22}\cdots a_{nn} > 0.\]
                      当$t \in [0, 1]$ 时，还有
                      \[a_{ii} > \sum_{j \neq i} \lvert a_{ij} \rvert \geqslant \sum_{j \neq i} \lvert ta_{ij} \rvert.\]
                      由 \eqref{item:15:B:5:1} 可知 $f(t)$ 在 $[0, 1]$ 上非零，由连续函数的介值定理可知 $f(1) > 0$，即 $\lvert A \rvert > 0$.

                \item 此为直接推论不再赘述.
            \end{enumerate}
\end{solve}

\begin{exercise}\label{prob:L11-例-1}
设$n$阶矩阵$A$的行列式$|A|\neq 0$，记$A$的前$n-1$列形成的矩阵为$A_1$，$A$的第$n$列为$\vec{b}$，问：线性方程组$A_1X=\vec{b}$是否有解？
\end{exercise}
\noindent 来源：\href{https://github.com/yhwu-is/Linear-Algebra-Left-Undone/blob/new/%E8%AE%B2%E4%B9%89/%E4%B8%93%E9%A2%98/11%20%E6%9C%9D%E8%8A%B1%E5%A4%95%E6%8B%BE.tex}{L11-例-1}。
\begin{solve}
无解；$|A|\neq 0$可知$r(A)=n$，$r(A_1)=n-1$，因此系数矩阵$A_1$的秩小于增广矩阵$A$的秩$n$，故无解.
\end{solve}

\begin{exercise}\label{prob:L11-例-2}
求解方程组$\begin{cases}
            x_1+x_2+x_3=1          \\
            a_1x_1+a_2x_2+a_3x_3=0 \\
            a_1^2x_1+a_2^2x_2+a_3^2x_3=0
        \end{cases}$，其中$a_1,a_2,a_3$两两不等.
\end{exercise}
\noindent 来源：\href{https://github.com/yhwu-is/Linear-Algebra-Left-Undone/blob/new/%E8%AE%B2%E4%B9%89/%E4%B8%93%E9%A2%98/11%20%E6%9C%9D%E8%8A%B1%E5%A4%95%E6%8B%BE.tex}{L11-例-2}。
\begin{solve}
$a_1,a_2,a_3$两两不等时，我们有
    \[D=\begin{vmatrix}
            1     & 1     & 1     \\
            a_1   & a_2   & a_3   \\
            a_1^2 & a_2^2 & a_3^2
        \end{vmatrix}=(a_2-a_1)(a_3-a_1)(a_3-a_2)\neq 0,\]
    根据Cramer法则，方程组有唯一解，且
    \[D_1=\begin{vmatrix}
            1 & 1     & 1     \\
            0 & a_2   & a_3   \\
            0 & a_2^2 & a_3^2
        \end{vmatrix}=(a_3-a_2)a_2a_3,\]
    \[D_2=\begin{vmatrix}
            1     & 1 & 1     \\
            a_1   & 0 & a_3   \\
            a_1^2 & 0 & a_3^2
        \end{vmatrix}=(a_1-a_3)a_1a_3,\]
    \[D_3=\begin{vmatrix}
            1     & 1     & 1 \\
            a_1   & a_2   & 0 \\
            a_1^2 & a_2^2 & 0
        \end{vmatrix}=(a_2-a_1)a_1a_2,\]
    因此
    \[x_1=\dfrac{D_1}{D}=\dfrac{a_2a_3}{(a_2-a_1)(a_3-a_1)},\]\[x_2=\dfrac{D_2}{D}=\dfrac{a_1a_3}{(a_1-a_2)(a_3-a_2)},\]\[x_3=\dfrac{D_3}{D}=\dfrac{a_1a_2}{(a_1-a_3)(a_2-a_3)}.\]
\end{solve}

\begin{exercise}\label{prob:L11-例-9}
当$k$取何值时，方程组：
    \[\begin{cases}
            x_1+x_2+kx_3=4 \\ -x_1+kx_2+x_3=k^2 \\ x_1-x_2+2x_3=-4
        \end{cases}\]
    有唯一解、无解、有无穷多解？在有解的情况下，求出方程组的全部解.
\end{exercise}
\noindent 来源：\href{https://github.com/yhwu-is/Linear-Algebra-Left-Undone/blob/new/%E8%AE%B2%E4%B9%89/%E4%B8%93%E9%A2%98/11%20%E6%9C%9D%E8%8A%B1%E5%A4%95%E6%8B%BE.tex}{L11-例-9}。
\begin{solve}
对于有无解的区别，我们一般都考虑直接使用高斯-若当消元法. 由高斯-若当消元法有（省略中间步骤直接得到阶梯矩阵）：
    \[\begin{pmatrix}
            1  & 1  & k & 4   \\
            -1 & k  & 1 & k^2 \\
            1  & -1 & 2 & -4
        \end{pmatrix}\to\begin{pmatrix}
            1 & 1 & k                     & 4      \\
            0 & 2 & k-2                   & 8      \\
            0 & 0 & \dfrac{(k+1)(4-k)}{2} & k(k-4)
        \end{pmatrix}.\]
    \begin{enumerate}
        \item 当$k=-1$时，增广矩阵秩为3，系数矩阵秩为2（或者说最后一行出现矛盾方程），无解；

        \item 当$k=4$时，增广矩阵和系数矩阵秩均为2，方程有无穷多解，解得同解为$k_1(-3,-1,1)^{\mathrm{T}}+(0,4,0)^{\mathrm{T}}(k_1\in\mathbf{R})$；

        \item 当$k\neq-1,4$时，增广矩阵和系数矩阵秩均为3，方程有唯一解，解为
              \[(\dfrac{k^2+2k}{k+1},\dfrac{k^2+2k+4}{k+1},-\dfrac{2k}{k+1})^{\mathrm{T}}.\]
    \end{enumerate}

    事实上，本题方程个数与未知数个数相等，因此可以运用 原讲义“Cramer”处解决. 首先求解系数矩阵$A$的行列式为
    \[|A|=\begin{vmatrix}
            1  & 1  & k \\
            -1 & k  & 1 \\
            1  & -1 & 2
        \end{vmatrix}=(k+1)(4-k),\]
    由Cramer法则，当$|A|\neq 0$时，方程组有唯一解，当$|A|=0$时，方程组无解或有无穷多解，其余关于有无解的讨论与上面一致.
\end{solve}

\begin{exercise}\label{prob:L11-例-10}
已知方程组\begin{gather*}
        \begin{cases}
            x_1+2x_2+3x_3=0  \\
            2x_1+3x_2+5x_3=0 \\
            x_1+x_2+ax_3=0
        \end{cases} \\
        \begin{cases}
            x_1+bx_2+cx_3=0 \\
            2x_1+b^2x_2+(c+1)x_3=0
        \end{cases}
    \end{gather*}
    同解，求$a,b,c$的值.
\end{exercise}
\noindent 来源：\href{https://github.com/yhwu-is/Linear-Algebra-Left-Undone/blob/new/%E8%AE%B2%E4%B9%89/%E4%B8%93%E9%A2%98/11%20%E6%9C%9D%E8%8A%B1%E5%A4%95%E6%8B%BE.tex}{L11-例-10}。
\begin{solve}
设第一个方程系数矩阵为$A$，第二个方程系数矩阵为$B$，则
    \[\begin{pmatrix}
            A \\ B
        \end{pmatrix}=\begin{pmatrix}
            1 & 2 & 3 \\ 2 & 3 & 5 \\ 1 & 1 & a \\ 1 & b & c \\ 2 & b^2 & c+1
        \end{pmatrix}\to\begin{pmatrix}
            1 & 0 & 1 \\ 0 & 1 & 1 \\ 0 & 0 & a-2 \\ 0 & 0 & c-b-1 \\ 0 & 0 & c-b^2-1
        \end{pmatrix},\]
    因此由同解条件$r\begin{pmatrix}
            A \\ B
        \end{pmatrix}=r(A)=r(B)$可知必有$r\begin{pmatrix}
            A \\ B
        \end{pmatrix}=r(A)=r(B)=2$（因为$r(B)\leqslant 2$，$r\begin{pmatrix}
            A \\ B
        \end{pmatrix}\geqslant 2$）.

    从而有$a-2=c-b-1=c-b^2-1=0$. 因此$a=2,b=0,c=1$或$a=2,b=1,c=2$. 当$a=2,b=0,c=1$时，$r(B)=1\neq r(A)=2$，舍去. 故$a=2,b=1,c=2$.
\end{solve}

\begin{exercise}\label{prob:L11-例-11}
设四元齐次线性方程组(1) 为\[\begin{cases}
            2x_1+3x_2-x_3=0 \\ x_1+2x_2+x_3-x_4=0
        \end{cases}\]已知另一个四元齐次线性方程组 (2) 的基础解系为
    \[\alpha_1=(2,-1,a+2,1)^\mathrm{T},\enspace\alpha_2=(-1,2,4,a+8)^\mathrm{T}\]
    \begin{enumerate}
        \item 求方程组 (1) 的一个基础解系；

        \item 当$a$为何值时，方程组 (1) 和 (2) 有非零公共解，并求出非零公共解.
    \end{enumerate}
\end{exercise}
\noindent 来源：\href{https://github.com/yhwu-is/Linear-Algebra-Left-Undone/blob/new/%E8%AE%B2%E4%B9%89/%E4%B8%93%E9%A2%98/11%20%E6%9C%9D%E8%8A%B1%E5%A4%95%E6%8B%BE.tex}{L11-例-11}。
\begin{solve}
\begin{enumerate}
        \item 直接给出结论：方程组 (1) 的一个基础解系为
              \[\beta_1=(5,-3,1,0)^{\mathrm{T}},\enspace\beta_2=(-3,2,0,1)^{\mathrm{T}}.\]

        \item 根据前述定理，两方程有非零公共解当且仅当$\alpha_1,\alpha_2,\beta_1,\beta_2$线性相关，可以将这$\beta_1,\beta_2,\alpha_1,\alpha_2$按列排成矩阵
              \[A = \begin{pmatrix}
                      5 & -3 & 2 & -1 \\ -3 & 2 & -1 & 2 \\ 1 & 0 & a+2 & 4 \\ 0 & 1 & 1 & a+8
                  \end{pmatrix},\]
              利用求解极大线性无关组的思想，对 $A$ 进行初等行变换得到阶梯矩阵
              \[\begin{pmatrix}
                      1 & 0 & a+2 & 4 \\ 0 & 1 & 1 & a+8 \\ 0 & 0 & 3a+3 & -2a-2 \\ 0 & 0 & 0 & -\frac{1}{3}a-\frac{1}{3}
                  \end{pmatrix},\]
              使得 $\alpha_1,\alpha_2,\beta_1,\beta_2$ 线性相关，则必有 $r(A) < 4$，即 $|A| = 0$，不难解得 $a = -1$. 此时阶梯矩阵为
              \[\begin{pmatrix}
                      1 & 0 & 1 & 4 \\ 0 & 1 & 1 & 7 \\ 0 & 0 & 0 & 0 \\ 0 & 0 & 0 & 0
                  \end{pmatrix},\]
              不难看出 $\beta_1,\beta_2$ 是向量组 $\beta_1,\beta_2,\alpha_1,\alpha_2$ 的极大线性无关组，因此 $\alpha_1,\alpha_2$ 可以完全由 $\beta_1,\beta_2$ 线性表示，故 $W_1 \cap W_2 = \operatorname{span}(\beta_1,\beta_2)$. 而求两方程组公共解就是求两子空间 $W_1$ 和 $W_2$ 的交集，故公共解为 $\beta_1,\beta_2$ 的线性组合，即
              \[k_1\beta_1 + k_2\beta_2 = k_1(5,-3,1,0)^{\mathrm{T}} + k_2(-3,2,0,1)^{\mathrm{T}}.\]
              其中 $k_1,k_2$ 为任意常数.
    \end{enumerate}
\end{solve}

\begin{exercise}\label{prob:L11-例-12}
已知$n$维列向量组$\alpha_1,\ldots,\alpha_s$线性无关，求一齐次线性方程组以$\alpha_1,\ldots,\alpha_s$为基础解系.
\end{exercise}
\noindent 来源：\href{https://github.com/yhwu-is/Linear-Algebra-Left-Undone/blob/new/%E8%AE%B2%E4%B9%89/%E4%B8%93%E9%A2%98/11%20%E6%9C%9D%E8%8A%B1%E5%A4%95%E6%8B%BE.tex}{L11-例-12}。
\begin{solve}
设所求的齐次线性方程组为$AX=0$，令$B=(\alpha_1,\ldots,\alpha_s)$，则$B$为$n\times s$矩阵且$AB=O$，于是$B^\mathrm{T}A^\mathrm{T}=O$，解线性方程组$B^\mathrm{T}X=0$，得到其基础解系为
    \[\beta_1,\beta_2,\ldots,\beta_{n-s},\]
    令$A^\mathrm{T}=(\beta_1,\beta_2,\ldots,\beta_{n-s})$即可，因为此时$B^\mathrm{T}A^\mathrm{T}=O$，故$AB=O$.
\end{solve}

\begin{exercise}\label{prob:L11-例-13}
设向量组$\alpha_1,\ldots,\alpha_s$线性无关，求一非齐次线性方程组$AX=\vec{b}$，使其解集以$\alpha_1,\ldots,\alpha_s$为极大线性无关组.
\end{exercise}
\noindent 来源：\href{https://github.com/yhwu-is/Linear-Algebra-Left-Undone/blob/new/%E8%AE%B2%E4%B9%89/%E4%B8%93%E9%A2%98/11%20%E6%9C%9D%E8%8A%B1%E5%A4%95%E6%8B%BE.tex}{L11-例-13}。
\begin{solve}
根据原讲义“非齐次线性方程组公共解”处，$\alpha_2-\alpha_1,\ldots,\alpha_s-\alpha_1$是$AX=\vec{0}$的基础解系，因此根据齐次线性方程组反问题的解法可以得到$A$，然后令$b=A\alpha_1$即可符合题意.
\end{solve}

\begin{exercise}\label{prob:L11-例-14}
已知$\alpha_1=(1,2,-1,0,4)^\mathrm{T},\enspace\alpha_2=(-1,3,2,4,1)^\mathrm{T},\enspace\alpha_3=(2,9,-1,4,13)^\mathrm{T}$，且有$W=\operatorname{span}(\alpha_1,\alpha_2,\alpha_3)$.
    \begin{enumerate}
        \item 求以$W$为解空间的一个齐次线性方程组；

        \item 求以$W'=\{\eta+\alpha \mid \alpha\in W\}$为解集的一个非齐次线性方程组，其中$\eta=(1,2,1,2,1)^\mathrm{T}$.
    \end{enumerate}
\end{exercise}
\noindent 来源：\href{https://github.com/yhwu-is/Linear-Algebra-Left-Undone/blob/new/%E8%AE%B2%E4%B9%89/%E4%B8%93%E9%A2%98/11%20%E6%9C%9D%E8%8A%B1%E5%A4%95%E6%8B%BE.tex}{L11-例-14}。
\begin{solve}
首先我们通过求解极大线性无关组的方法可以得到，$\alpha_1,\alpha_2,\alpha_3$的极大线性无关组是$\alpha_1,\alpha_2$. 然后我们基于此根据前面两个例题介绍的方法，我们有如下求解过程：
    \begin{enumerate}
        \item 设所求的齐次线性方程组为$AX=0$，令$B=(\alpha_1,\alpha_2)$，解线性方程组$B^\mathrm{T}X=0$，得到其基础解系为
              \[\beta_1=(7,-1,5,0,0)^\mathrm{T},\enspace\beta_2=(8,-4,0,5,0)^\mathrm{T},\enspace\beta_3=(-2,-1,0,0,1)^\mathrm{T},\]
              令$A^\mathrm{T}=(\beta_1,\beta_2,\beta_3)$即可，即
              \[A=\begin{pmatrix}
                      7 & -1 & 5 & 0 & 0 \\ 8 & -4 & 0 & 5 & 0 \\ -2 & -1 & 0 & 0 & 1
                  \end{pmatrix}.\]
              故所求的线性方程组为$\begin{cases}
                      7x_1-x_2+5x_3=0 \\ 8x_1-4x_2+5x_4=0 \\ -2x_1-x_2+x_5=0
                  \end{cases}$.

        \item 这一问可以利用前面例子中给出的方法，具体步骤不在此赘述（注意因为$\eta$不在$W$中，因此$W'$中可以有三个线性无关向量，不可想当然认为只有$\alpha_2-\alpha_1$是$AX=0$（设$A$为本题要求的方程组的系数矩阵）的解）. 我们这里给出更简单的方法，读者只需回顾根据原讲义“通解加特解”处 即可发现奥秘. 事实上，将$x_1=1,x_2=2,x_3=1,x_4=2,x_5=1$代入有$7x_1-x_2+5x_3=10$，$8x_1-4x_2+5x_4=10$，$-2x_1-x_2+x_5=-3$，因此$\eta$就是方程组$\begin{cases}
                      7x_1-x_2+5x_3=10 \\ 8x_1-4x_2+5x_4=10 \\ -2x_1-x_2+x_5=-3
                  \end{cases}$的一个特解，从而根据原讲义“通解加特解”处 可以知道上述方程组符合题意.
    \end{enumerate}
\end{solve}

\begin{exercise}\label{prob:S-SYS-1}
设 $A=\begin{pmatrix}1&-6&0\\0&2&0\\1&0&3\end{pmatrix}$，$b=\begin{pmatrix}-7\\2\\8\end{pmatrix}$。求 $\det A$、$A^{-1}$，解 $Ax=b$，并利用消元结果说明解的唯一性。
\end{exercise}
\noindent 编者补充：S-SYS-1；相关主题见原讲义对应讲，题目数据与解答由本选编另行构造。
\begin{solve}
矩阵的行列式为 $\det A=6\ne0$。对增广矩阵 $(A\mid I_3)$ 作行初等变换，得到
\[(A\mid I_3)\sim(I_3\mid A^{-1}),\qquad A^{-1}=\begin{pmatrix}1&3&0\\0&\frac{1}{2}&0\\-\frac{1}{3}&-1&\frac{1}{3}\end{pmatrix}.\]
因此 $x=A^{-1}b=\begin{pmatrix}-1\\1\\3\end{pmatrix}$。也可以直接化简含常数列的增广矩阵：
\[(A\mid b)\sim \begin{pmatrix}1&0&0&-1\\0&1&0&1\\0&0&1&3\end{pmatrix}.\]
三个未知量所在的列都有主元，没有自由变量，所以此解唯一。检验时将所求向量代回 $Ax$，得到的三维列向量恰为题给 $b$；求逆过程中的等式可通过 $AA^{-1}=I_3$ 复核。
\end{solve}

\begin{exercise}\label{prob:S-NULL-1}
给定 $A=\begin{pmatrix}1&-3&-2&-4&-1\\0&1&1&1&1\\1&0&1&-1&2\end{pmatrix}$ 和 $b=\begin{pmatrix}-2\\1\\1\end{pmatrix}$。求 $r(A)$、$\ker A$ 的一组基及 $Ax=b$ 的全部解；再说明其解集为什么通常不是线性子空间。
\end{exercise}
\noindent 编者补充：S-NULL-1；相关主题见原讲义对应讲，题目数据与解答由本选编另行构造。
\begin{solve}
行化简得到
\[(A\mid b)\sim \begin{pmatrix}1&0&1&-1&2&1\\0&1&1&1&1&1\\0&0&0&0&0&0\end{pmatrix}.\]
前两列是主元列，所以 $r(A)=2$，核的维数为 $5-2=3$。把 $x_3,x_4,x_5$ 分别作为自由变量，可得核的一组基为下列矩阵 $N$ 的三列：
\[N=\begin{pmatrix}-1&1&-2\\-1&-1&-1\\1&0&0\\0&1&0\\0&0&1\end{pmatrix}.\]
取 $x_0=\begin{pmatrix}1\\1\\0\\0\\0\end{pmatrix}$，所有解可统一写成 $x=x_0+Nt$，其中 $t\in\mathbf R^3$ 任意。直接相乘得到 $AN=O$ 和 $Ax_0=b$，验证了这一表示；反过来，两特解之差在 $\ker A$ 中，说明没有遗漏任何解。由于此处 $b\ne0$，零向量不是解，故解集是三维仿射子空间而不是线性子空间。
\end{solve}

\begin{exercise}\label{prob:S-SYS-2}
设 $A=\begin{pmatrix}1&-2&-8\\0&2&8\\0&0&4\end{pmatrix}$，$b=\begin{pmatrix}-24\\24\\12\end{pmatrix}$。求 $\det A$、$A^{-1}$，解 $Ax=b$，并利用消元结果说明解的唯一性。
\end{exercise}
\noindent 编者补充：S-SYS-2；相关主题见原讲义对应讲，题目数据与解答由本选编另行构造。
\begin{solve}
矩阵的行列式为 $\det A=8\ne0$。对增广矩阵 $(A\mid I_3)$ 作行初等变换，得到
\[(A\mid I_3)\sim(I_3\mid A^{-1}),\qquad A^{-1}=\begin{pmatrix}1&1&0\\0&\frac{1}{2}&-1\\0&0&\frac{1}{4}\end{pmatrix}.\]
因此 $x=A^{-1}b=\begin{pmatrix}0\\0\\3\end{pmatrix}$。也可以直接化简含常数列的增广矩阵：
\[(A\mid b)\sim \begin{pmatrix}1&0&0&0\\0&1&0&0\\0&0&1&3\end{pmatrix}.\]
三个未知量所在的列都有主元，没有自由变量，所以此解唯一。检验时将所求向量代回 $Ax$，得到的三维列向量恰为题给 $b$；求逆过程中的等式可通过 $AA^{-1}=I_3$ 复核。
\end{solve}

\begin{exercise}\label{prob:S-NULL-2}
给定 $A=\begin{pmatrix}1&-1&1&-3&1\\0&1&1&2&1\\0&0&0&0&0\end{pmatrix}$ 和 $b=\begin{pmatrix}-1\\2\\0\end{pmatrix}$。求 $r(A)$、$\ker A$ 的一组基及 $Ax=b$ 的全部解；再说明其解集为什么通常不是线性子空间。
\end{exercise}
\noindent 编者补充：S-NULL-2；相关主题见原讲义对应讲，题目数据与解答由本选编另行构造。
\begin{solve}
行化简得到
\[(A\mid b)\sim \begin{pmatrix}1&0&2&-1&2&1\\0&1&1&2&1&2\\0&0&0&0&0&0\end{pmatrix}.\]
前两列是主元列，所以 $r(A)=2$，核的维数为 $5-2=3$。把 $x_3,x_4,x_5$ 分别作为自由变量，可得核的一组基为下列矩阵 $N$ 的三列：
\[N=\begin{pmatrix}-2&1&-2\\-1&-2&-1\\1&0&0\\0&1&0\\0&0&1\end{pmatrix}.\]
取 $x_0=\begin{pmatrix}1\\2\\0\\0\\0\end{pmatrix}$，所有解可统一写成 $x=x_0+Nt$，其中 $t\in\mathbf R^3$ 任意。直接相乘得到 $AN=O$ 和 $Ax_0=b$，验证了这一表示；反过来，两特解之差在 $\ker A$ 中，说明没有遗漏任何解。由于此处 $b\ne0$，零向量不是解，故解集是三维仿射子空间而不是线性子空间。
\end{solve}

\begin{exercise}\label{prob:S-SYS-3}
设 $A=\begin{pmatrix}1&-4&-20\\1&-2&-10\\0&0&5\end{pmatrix}$，$b=\begin{pmatrix}-55\\-27\\15\end{pmatrix}$。求 $\det A$、$A^{-1}$，解 $Ax=b$，并利用消元结果说明解的唯一性。
\end{exercise}
\noindent 编者补充：S-SYS-3；相关主题见原讲义对应讲，题目数据与解答由本选编另行构造。
\begin{solve}
矩阵的行列式为 $\det A=10\ne0$。对增广矩阵 $(A\mid I_3)$ 作行初等变换，得到
\[(A\mid I_3)\sim(I_3\mid A^{-1}),\qquad A^{-1}=\begin{pmatrix}-1&2&0\\-\frac{1}{2}&\frac{1}{2}&-1\\0&0&\frac{1}{5}\end{pmatrix}.\]
因此 $x=A^{-1}b=\begin{pmatrix}1\\-1\\3\end{pmatrix}$。也可以直接化简含常数列的增广矩阵：
\[(A\mid b)\sim \begin{pmatrix}1&0&0&1\\0&1&0&-1\\0&0&1&3\end{pmatrix}.\]
三个未知量所在的列都有主元，没有自由变量，所以此解唯一。检验时将所求向量代回 $Ax$，得到的三维列向量恰为题给 $b$；求逆过程中的等式可通过 $AA^{-1}=I_3$ 复核。
\end{solve}

\begin{exercise}\label{prob:S-NULL-3}
给定 $A=\begin{pmatrix}1&-2&1&-7&0\\1&-1&2&-4&1\\0&0&0&0&0\end{pmatrix}$ 和 $b=\begin{pmatrix}-5\\-2\\0\end{pmatrix}$。求 $r(A)$、$\ker A$ 的一组基及 $Ax=b$ 的全部解；再说明其解集为什么通常不是线性子空间。
\end{exercise}
\noindent 编者补充：S-NULL-3；相关主题见原讲义对应讲，题目数据与解答由本选编另行构造。
\begin{solve}
行化简得到
\[(A\mid b)\sim \begin{pmatrix}1&0&3&-1&2&1\\0&1&1&3&1&3\\0&0&0&0&0&0\end{pmatrix}.\]
前两列是主元列，所以 $r(A)=2$，核的维数为 $5-2=3$。把 $x_3,x_4,x_5$ 分别作为自由变量，可得核的一组基为下列矩阵 $N$ 的三列：
\[N=\begin{pmatrix}-3&1&-2\\-1&-3&-1\\1&0&0\\0&1&0\\0&0&1\end{pmatrix}.\]
取 $x_0=\begin{pmatrix}1\\3\\0\\0\\0\end{pmatrix}$，所有解可统一写成 $x=x_0+Nt$，其中 $t\in\mathbf R^3$ 任意。直接相乘得到 $AN=O$ 和 $Ax_0=b$，验证了这一表示；反过来，两特解之差在 $\ker A$ 中，说明没有遗漏任何解。由于此处 $b\ne0$，零向量不是解，故解集是三维仿射子空间而不是线性子空间。
\end{solve}

\begin{exercise}\label{prob:S-SYS-4}
设 $A=\begin{pmatrix}1&0&-12\\0&2&-12\\0&0&6\end{pmatrix}$，$b=\begin{pmatrix}-34\\-40\\18\end{pmatrix}$。求 $\det A$、$A^{-1}$，解 $Ax=b$，并利用消元结果说明解的唯一性。
\end{exercise}
\noindent 编者补充：S-SYS-4；相关主题见原讲义对应讲，题目数据与解答由本选编另行构造。
\begin{solve}
矩阵的行列式为 $\det A=12\ne0$。对增广矩阵 $(A\mid I_3)$ 作行初等变换，得到
\[(A\mid I_3)\sim(I_3\mid A^{-1}),\qquad A^{-1}=\begin{pmatrix}1&0&2\\0&\frac{1}{2}&1\\0&0&\frac{1}{6}\end{pmatrix}.\]
因此 $x=A^{-1}b=\begin{pmatrix}2\\-2\\3\end{pmatrix}$。也可以直接化简含常数列的增广矩阵：
\[(A\mid b)\sim \begin{pmatrix}1&0&0&2\\0&1&0&-2\\0&0&1&3\end{pmatrix}.\]
三个未知量所在的列都有主元，没有自由变量，所以此解唯一。检验时将所求向量代回 $Ax$，得到的三维列向量恰为题给 $b$；求逆过程中的等式可通过 $AA^{-1}=I_3$ 复核。
\end{solve}

\begin{exercise}\label{prob:S-NULL-4}
给定 $A=\begin{pmatrix}1&0&4&-1&2\\0&1&1&4&1\\0&0&0&0&0\end{pmatrix}$ 和 $b=\begin{pmatrix}1\\4\\0\end{pmatrix}$。求 $r(A)$、$\ker A$ 的一组基及 $Ax=b$ 的全部解；再说明其解集为什么通常不是线性子空间。
\end{exercise}
\noindent 编者补充：S-NULL-4；相关主题见原讲义对应讲，题目数据与解答由本选编另行构造。
\begin{solve}
行化简得到
\[(A\mid b)\sim \begin{pmatrix}1&0&4&-1&2&1\\0&1&1&4&1&4\\0&0&0&0&0&0\end{pmatrix}.\]
前两列是主元列，所以 $r(A)=2$，核的维数为 $5-2=3$。把 $x_3,x_4,x_5$ 分别作为自由变量，可得核的一组基为下列矩阵 $N$ 的三列：
\[N=\begin{pmatrix}-4&1&-2\\-1&-4&-1\\1&0&0\\0&1&0\\0&0&1\end{pmatrix}.\]
取 $x_0=\begin{pmatrix}1\\4\\0\\0\\0\end{pmatrix}$，所有解可统一写成 $x=x_0+Nt$，其中 $t\in\mathbf R^3$ 任意。直接相乘得到 $AN=O$ 和 $Ax_0=b$，验证了这一表示；反过来，两特解之差在 $\ker A$ 中，说明没有遗漏任何解。由于此处 $b\ne0$，零向量不是解，故解集是三维仿射子空间而不是线性子空间。
\end{solve}

\begin{exercise}\label{prob:S-SYS-5}
设 $A=\begin{pmatrix}1&8&21\\0&2&0\\0&4&7\end{pmatrix}$，$b=\begin{pmatrix}42\\-6\\9\end{pmatrix}$。求 $\det A$、$A^{-1}$，解 $Ax=b$，并利用消元结果说明解的唯一性。
\end{exercise}
\noindent 编者补充：S-SYS-5；相关主题见原讲义对应讲，题目数据与解答由本选编另行构造。
\begin{solve}
矩阵的行列式为 $\det A=14\ne0$。对增广矩阵 $(A\mid I_3)$ 作行初等变换，得到
\[(A\mid I_3)\sim(I_3\mid A^{-1}),\qquad A^{-1}=\begin{pmatrix}1&2&-3\\0&\frac{1}{2}&0\\0&-\frac{2}{7}&\frac{1}{7}\end{pmatrix}.\]
因此 $x=A^{-1}b=\begin{pmatrix}3\\-3\\3\end{pmatrix}$。也可以直接化简含常数列的增广矩阵：
\[(A\mid b)\sim \begin{pmatrix}1&0&0&3\\0&1&0&-3\\0&0&1&3\end{pmatrix}.\]
三个未知量所在的列都有主元，没有自由变量，所以此解唯一。检验时将所求向量代回 $Ax$，得到的三维列向量恰为题给 $b$；求逆过程中的等式可通过 $AA^{-1}=I_3$ 复核。
\end{solve}

\begin{exercise}\label{prob:S-NULL-5}
给定 $A=\begin{pmatrix}1&4&9&19&6\\0&1&1&5&1\\0&2&2&10&2\end{pmatrix}$ 和 $b=\begin{pmatrix}21\\5\\10\end{pmatrix}$。求 $r(A)$、$\ker A$ 的一组基及 $Ax=b$ 的全部解；再说明其解集为什么通常不是线性子空间。
\end{exercise}
\noindent 编者补充：S-NULL-5；相关主题见原讲义对应讲，题目数据与解答由本选编另行构造。
\begin{solve}
行化简得到
\[(A\mid b)\sim \begin{pmatrix}1&0&5&-1&2&1\\0&1&1&5&1&5\\0&0&0&0&0&0\end{pmatrix}.\]
前两列是主元列，所以 $r(A)=2$，核的维数为 $5-2=3$。把 $x_3,x_4,x_5$ 分别作为自由变量，可得核的一组基为下列矩阵 $N$ 的三列：
\[N=\begin{pmatrix}-5&1&-2\\-1&-5&-1\\1&0&0\\0&1&0\\0&0&1\end{pmatrix}.\]
取 $x_0=\begin{pmatrix}1\\5\\0\\0\\0\end{pmatrix}$，所有解可统一写成 $x=x_0+Nt$，其中 $t\in\mathbf R^3$ 任意。直接相乘得到 $AN=O$ 和 $Ax_0=b$，验证了这一表示；反过来，两特解之差在 $\ker A$ 中，说明没有遗漏任何解。由于此处 $b\ne0$，零向量不是解，故解集是三维仿射子空间而不是线性子空间。
\end{solve}

\begin{exercise}\label{prob:S-SYS-6}
设 $A=\begin{pmatrix}1&0&0\\0&2&0\\-1&-2&8\end{pmatrix}$，$b=\begin{pmatrix}4\\-8\\28\end{pmatrix}$。求 $\det A$、$A^{-1}$，解 $Ax=b$，并利用消元结果说明解的唯一性。
\end{exercise}
\noindent 编者补充：S-SYS-6；相关主题见原讲义对应讲，题目数据与解答由本选编另行构造。
\begin{solve}
矩阵的行列式为 $\det A=16\ne0$。对增广矩阵 $(A\mid I_3)$ 作行初等变换，得到
\[(A\mid I_3)\sim(I_3\mid A^{-1}),\qquad A^{-1}=\begin{pmatrix}1&0&0\\0&\frac{1}{2}&0\\\frac{1}{8}&\frac{1}{8}&\frac{1}{8}\end{pmatrix}.\]
因此 $x=A^{-1}b=\begin{pmatrix}4\\-4\\3\end{pmatrix}$。也可以直接化简含常数列的增广矩阵：
\[(A\mid b)\sim \begin{pmatrix}1&0&0&4\\0&1&0&-4\\0&0&1&3\end{pmatrix}.\]
三个未知量所在的列都有主元，没有自由变量，所以此解唯一。检验时将所求向量代回 $Ax$，得到的三维列向量恰为题给 $b$；求逆过程中的等式可通过 $AA^{-1}=I_3$ 复核。
\end{solve}

\begin{exercise}\label{prob:S-NULL-6}
给定 $A=\begin{pmatrix}1&0&6&-1&2\\0&1&1&6&1\\-1&-1&-7&-5&-3\end{pmatrix}$ 和 $b=\begin{pmatrix}1\\6\\-7\end{pmatrix}$。求 $r(A)$、$\ker A$ 的一组基及 $Ax=b$ 的全部解；再说明其解集为什么通常不是线性子空间。
\end{exercise}
\noindent 编者补充：S-NULL-6；相关主题见原讲义对应讲，题目数据与解答由本选编另行构造。
\begin{solve}
行化简得到
\[(A\mid b)\sim \begin{pmatrix}1&0&6&-1&2&1\\0&1&1&6&1&6\\0&0&0&0&0&0\end{pmatrix}.\]
前两列是主元列，所以 $r(A)=2$，核的维数为 $5-2=3$。把 $x_3,x_4,x_5$ 分别作为自由变量，可得核的一组基为下列矩阵 $N$ 的三列：
\[N=\begin{pmatrix}-6&1&-2\\-1&-6&-1\\1&0&0\\0&1&0\\0&0&1\end{pmatrix}.\]
取 $x_0=\begin{pmatrix}1\\6\\0\\0\\0\end{pmatrix}$，所有解可统一写成 $x=x_0+Nt$，其中 $t\in\mathbf R^3$ 任意。直接相乘得到 $AN=O$ 和 $Ax_0=b$，验证了这一表示；反过来，两特解之差在 $\ker A$ 中，说明没有遗漏任何解。由于此处 $b\ne0$，零向量不是解，故解集是三维仿射子空间而不是线性子空间。
\end{solve}

\section{一元多项式}\label{sec:topic13}
本节收录 23 道题，其中 12 道为编者补充计算题。区分多项式恒等式与有限个点的取值关系，注意零多项式及次数上界；带余除法、插值和互素性是后续标准形的基础。
\begin{exercise}\label{prob:L13-A-1}
设 $\mathbf F$ 是域，证明：$\mathbf F[x]$ 在通常的多项式加法和乘法下构成含幺交换环。
\end{exercise}
\noindent 来源：\href{https://github.com/yhwu-is/Linear-Algebra-Left-Undone/blob/new/%E8%AE%B2%E4%B9%89/%E4%B8%93%E9%A2%98/13%20%E5%A4%9A%E9%A1%B9%E5%BC%8F.tex}{L13-A-1}。
\begin{solve}
将多项式看成仅有有限个非零项的系数序列，规定
\[(f+g)_k=a_k+b_k,\qquad(fg)_k=\sum_{i+j=k}a_ib_j,\]
其中 $f=\sum_ia_ix^i$、$g=\sum_jb_jx^j$，未出现的系数取零。加法逐项继承域中的交换律、结合律；零多项式为加法单位元，$-f=\sum_i(-a_i)x^i$ 为加法逆元。有限支撑保证两种运算的结果仍是多项式。
乘法交换律由交换 $i,j$ 得到；若 $h=\sum_kc_kx^k$，则 $(fg)h$ 和 $f(gh)$ 的第 $r$ 个系数都为 $\sum_{i+j+k=r}a_ib_jc_k$，所以乘法结合。各系数的分配律给出 $f(g+h)=fg+fh$，另一侧同理。常数多项式 $1$ 为乘法单位元，故全部环公理成立。
\end{solve}

\begin{exercise}\label{prob:L13-A-2}
设$p(x),q(x),r(x)\in\mathbf{F}[x]$，则
        \begin{enumerate}
            \item 若$p(x),q(x)\neq 0$，则$p(x)q(x)\neq 0$；
            \item 若$p(x)\neq 0$，且$p(x)q(x)=p(x)r(x)$，则$q(x)=r(x)$.
        \end{enumerate}
\end{exercise}
\noindent 来源：\href{https://github.com/yhwu-is/Linear-Algebra-Left-Undone/blob/new/%E8%AE%B2%E4%B9%89/%E4%B8%93%E9%A2%98/13%20%E5%A4%9A%E9%A1%B9%E5%BC%8F.tex}{L13-A-2}。
\begin{solve}
\begin{enumerate}
\item 若 $p,q$ 非零，设其次数为 $m,n$，首项系数为 $a_m,b_n\ne0$。乘积中 $x^{m+n}$ 的系数为 $a_mb_n\ne0$，故 $pq\ne0$，并有 $\deg(pq)=m+n$。
\item 由 $pq=pr$ 得 $p(q-r)=0$。若 $q-r\ne0$，第 (1) 项将给出乘积非零，矛盾。因此 $q=r$。
\end{enumerate}
此证明对任意系数域成立，无需通过选取多项式的非零函数值来论证。
\end{solve}

\begin{exercise}\label{prob:L13-A-4}
证明：每个奇数次的实系数多项式都有实的零点.
\end{exercise}
\noindent 来源：\href{https://github.com/yhwu-is/Linear-Algebra-Left-Undone/blob/new/%E8%AE%B2%E4%B9%89/%E4%B8%93%E9%A2%98/13%20%E5%A4%9A%E9%A1%B9%E5%BC%8F.tex}{L13-A-4}。
\begin{solve}
设 $f(x)=a_nx^n+\cdots+a_0$，其中 $n$ 为奇数且 $a_n\ne0$。必要时用 $-f$ 代替 $f$，可设 $a_n>0$。由
\[\frac{f(x)}{x^n}=a_n+\frac{a_{n-1}}x+\cdots+\frac{a_0}{x^n}\longrightarrow a_n\quad(|x|\to\infty),\]
可取充分大的 $R>0$，使 $f(R)>0$ 而 $f(-R)<0$。多项式在 $[-R,R]$ 上连续，由介值定理存在 $c\in(-R,R)$ 使 $f(c)=0$。
\end{solve}

\begin{exercise}\label{prob:L13-B-1}
设多项式$f(x)$被$(x-1),(x-2),(x-3)$除后，余式分别为$4,8,16$. 求$f(x)$被$(x-1)(x-2)(x-3)$除后的余式.
\end{exercise}
\noindent 来源：\href{https://github.com/yhwu-is/Linear-Algebra-Left-Undone/blob/new/%E8%AE%B2%E4%B9%89/%E4%B8%93%E9%A2%98/13%20%E5%A4%9A%E9%A1%B9%E5%BC%8F.tex}{L13-B-1}。
\begin{solve}
$1=\dfrac{1}{2}(x-2)(x-3)-\dfrac{1}{2}(x-4)(x-1)$，故$\dfrac{1}{2}(x-2)(x-3) \equiv1\bmod{(x-1)}$.

            同理有$-(x-1)(x-3)\equiv1\bmod{(x-2)}$，$\dfrac{1}{2}(x-1)(x-2)\equiv1\bmod{(x-3)}$.

            从而$f(x)\equiv\dfrac{1}{2}\times 4 \times(x-2)(x-3)-1\times 8\times(x-1)(x-3)+\dfrac{1}{2}\times 16\times(x-1)(x-2)\equiv 2x^2-2x+4\bmod{(x-1)(x-2)(x-3)}$.
\end{solve}

\begin{exercise}\label{prob:L13-B-2}
设$p,q\in\mathbf{F}[x]$，证明：$p^2 \mid q^2\iff p \mid q$.
\end{exercise}
\noindent 来源：\href{https://github.com/yhwu-is/Linear-Algebra-Left-Undone/blob/new/%E8%AE%B2%E4%B9%89/%E4%B8%93%E9%A2%98/13%20%E5%A4%9A%E9%A1%B9%E5%BC%8F.tex}{L13-B-2}。
\begin{solve}
若 $p\mid q$，写 $q=ph$，则 $q^2=p^2h^2$，故 $p^2\mid q^2$。反过来，若 $q=0$，结论直接成立；若 $q\ne0$ 且 $p^2\mid q^2$，则 $p\ne0$。利用域上多项式的唯一因式分解，对每个首一不可约多项式 $\pi$，记 $\nu_\pi(p)$、$\nu_\pi(q)$ 为其在分解中的指数。整除关系给出 $2\nu_\pi(p)\le2\nu_\pi(q)$，所以 $\nu_\pi(p)\le\nu_\pi(q)$ 对所有 $\pi$ 成立，即 $p\mid q$。若 $p=0$，按 $p\mid q\iff q=ph$ 的定义，两侧都等价于 $q=0$。
\end{solve}

\begin{exercise}\label{prob:L13-B-3}
证明：$\mathbf{F}[x]$中两个次数大于0的多项式没有公共复根的充要条件是它们互素.
\end{exercise}
\noindent 来源：\href{https://github.com/yhwu-is/Linear-Algebra-Left-Undone/blob/new/%E8%AE%B2%E4%B9%89/%E4%B8%93%E9%A2%98/13%20%E5%A4%9A%E9%A1%B9%E5%BC%8F.tex}{L13-B-3}。
\begin{solve}
\begin{enumerate}
                \item 充分性：假设多项式$p,q$有公共复根$\lambda$，则$(x-\lambda)$为$p$和$q$的公因式，这与$p$和$q$互素矛盾.
                \item 必要性：设$d(x)$是$p(x)$和$q(x)$的公因式，则$\deg d>1$，由代数基本定理知$d(x)$在复数域上有根，这表明$p$和$q$有公共复根，矛盾.从而$p$和$q$互素.
            \end{enumerate}
\end{solve}

\begin{exercise}\label{prob:L13-B-4}
设 $\mathbf F$ 为数域，$0\ne f\in\mathbf F[x]$，$c\in\mathbf F$，$k\ge1$。证明：$c$ 恰为 $f$ 的 $k$ 重根，当且仅当
\[f(c)=f'(c)=\cdots=f^{(k-1)}(c)=0,\qquad f^{(k)}(c)\ne0.\]
\end{exercise}
\noindent 来源：\href{https://github.com/yhwu-is/Linear-Algebra-Left-Undone/blob/new/%E8%AE%B2%E4%B9%89/%E4%B8%93%E9%A2%98/13%20%E5%A4%9A%E9%A1%B9%E5%BC%8F.tex}{L13-B-4}。
\begin{solve}
在点 $c$ 处作多项式的有限 Taylor 展开：
\[f(x)=\sum_{j=0}^{\deg f}\frac{f^{(j)}(c)}{j!}(x-c)^j.\]
数域的特征为零，所以每个 $j!$ 均可逆。根的重数正是展开中第一个非零项的次数；因此重数恰为 $k$ 与题给导数条件等价。也可将 $f=(x-c)^kg$ 且 $g(c)\ne0$ 代入逐次求导，得到 $f^{(k)}(c)=k!g(c)\ne0$。
\end{solve}

\begin{exercise}\label{prob:L13-B-5}
设$p(x)$和$q(x)$是$\mathbf{F}[x]$中的两个次数不超过$n$的多项式，若存在$\mathbf{F}$上$n+1$个不同的数$\lambda_0,\lambda_1,\ldots,\lambda_n$使得$p(\lambda_i)=q(\lambda_i),\enspace i=0,1,\ldots,n$，则$p(x)=q(x)$.
\end{exercise}
\noindent 来源：\href{https://github.com/yhwu-is/Linear-Algebra-Left-Undone/blob/new/%E8%AE%B2%E4%B9%89/%E4%B8%93%E9%A2%98/13%20%E5%A4%9A%E9%A1%B9%E5%BC%8F.tex}{L13-B-5}。
\begin{solve}
考虑多项式$f(x)=p(x)-q(x)$，$\deg f\leqslant \max\{\deg p,\deg q\}=n$，而$m\enspace(m\geqslant 0)$次多项式至多有$m$个根，这表明$\deg f=-\infty$，只能有$f(x)\equiv 0$，从而$p(x)=q(x)$.
\end{solve}

\begin{exercise}\label{prob:L13-C-1}
设 $p,s\in\mathbf F[x]$ 且 $s\ne0$，其中 $\mathbf F$ 为域。证明存在唯一的 $q,r\in\mathbf F[x]$，使 $p=sq+r$ 且 $r=0$ 或 $\deg r<\deg s$。请用余式次数最小的方法证明存在性。
\end{exercise}
\noindent 来源：\href{https://github.com/yhwu-is/Linear-Algebra-Left-Undone/blob/new/%E8%AE%B2%E4%B9%89/%E4%B8%93%E9%A2%98/13%20%E5%A4%9A%E9%A1%B9%E5%BC%8F.tex}{L13-C-1}。
\begin{solve}
集合 $\{p-sq:q\in\mathbf F[x]\}$ 非空，因为取 $q=0$ 即含 $p$。若它含零，存在性成立；否则在所有元素中选次数最小者 $r=p-sq$。设 $\deg r=m\ge d=\deg s$，两多项式的首项系数分别为 $a,b$。则
\[r_1=r-\frac abx^{m-d}s=p-s\left(q+\frac abx^{m-d}\right)\]
仍属于该集合，而首项抵消，故 $r_1=0$ 或 $\deg r_1<m$。前一种情形与集合不含零矛盾，后一种与最小性矛盾。因此 $\deg r<d$。若另有 $p=sq_1+r_1=sq_2+r_2$ 满足要求，则 $s(q_1-q_2)=r_2-r_1$；若 $q_1\ne q_2$，左端次数至少为 $\deg s$，右端次数却更小，矛盾。故 $q_1=q_2,r_1=r_2$。
\end{solve}

\begin{exercise}\label{prob:L13-C-2}
设 $p,q\in\mathbf F[x]$，$\deg p=n\ge1$、$\deg q=m\ge1$。用有限维线性映射证明：$p,q$ 互素当且仅当存在多项式 $u,v$ 使 $up+vq=1$。证明互素时可以选择 $\deg u<m$、$\deg v<n$。
\end{exercise}
\noindent 来源：\href{https://github.com/yhwu-is/Linear-Algebra-Left-Undone/blob/new/%E8%AE%B2%E4%B9%89/%E4%B8%93%E9%A2%98/13%20%E5%A4%9A%E9%A1%B9%E5%BC%8F.tex}{L13-C-2}。
\begin{solve}
若互素，考察线性映射
\[T:\mathbf F[x]_m\times\mathbf F[x]_n\longrightarrow\mathbf F[x]_{m+n},\qquad T(u,v)=up+vq.\]
设 $up+vq=0$，则 $p\mid vq$；由互素性和唯一因子分解可知 $p\mid v$。但 $\deg v<n=\deg p$，故 $v=0$，进而 $u=0$。所以 $T$ 单射。定义域与陪域维数都为 $m+n$，故它也是满射；特别地常值多项式 $1$ 有原像，从而得到所需次数界。反之，若 $up+vq=1$，任何公共因子都整除 $1$，所以 $p,q$ 的首一最大公因式为 $1$。
\end{solve}

\begin{exercise}\label{prob:L13-例-2}
设$p(x)=x^4-x^3-x^2+2x-1$，$q(x)=x^3-2x+1$，求$(p(x),q(x))$，以及$u(x),v(x)$使得$(p(x),q(x))=u(x)p(x)+v(x)q(x)$.
\end{exercise}
\noindent 来源：\href{https://github.com/yhwu-is/Linear-Algebra-Left-Undone/blob/new/%E8%AE%B2%E4%B9%89/%E4%B8%93%E9%A2%98/13%20%E5%A4%9A%E9%A1%B9%E5%BC%8F.tex}{L13-例-2}。
\begin{solve}
我们可以利用带余除法进行计算：
    \begin{align*}
        p(x)   & =q(x)(x-1)+x^2-x=q(x)s_1(x)+r_1(x),   \\
        q(x)   & =r_1(x)(x+1)-x+1=r_1(x)s_2(x)+r_2(x), \\
        r_1(x) & =r_2(x)(-x)=r_2(x)s_3(x),
    \end{align*}
    因此最大公因式$(p(x),q(x))=r_2(x)=-x+1$. 接下来开始回推，我们有
    \begin{align*}
        r_2(x) & =q(x)-r_1(x)s_2(x)                   \\
               & =q(x)-(p(x)-q(x)s_1(x))s_2(x)        \\
               & =p(x)(-s_2(x))+q(x)(1+s_1(x)s_2(x)),
    \end{align*}
    所以$u(x)=-s_2(x)=-x-1$，$v(x)=1+s_1(x)s_2(x)=x^2$.
\end{solve}

\begin{exercise}\label{prob:S-GCD-1}
设 $f(x)=x^{2}-3x+2$，$g(x)=x^{2}-5x+4$。求首一最大公因式、首一最小公倍式，并找出多项式 $u,v$ 使 $uf+vg=\gcd(f,g)$。进一步求 $f+g$ 与 $f-g$ 的首一最大公因式。
\end{exercise}
\noindent 编者补充：S-GCD-1；相关主题见原讲义对应讲，题目数据与解答由本选编另行构造。
\begin{solve}
分解得到 $f=(x-1)(x-2)$、$g=(x-1)(x-4)$。三个根两两不同，所以
\[\gcd(f,g)=x-1,\qquad\operatorname{lcm}(f,g)=x^{3}-7x^{2}+14x-8.\]
相减得 $f-g=(2)(x-1)$，从而可以取常数多项式
\[u=\frac{1}{2},\qquad v=-\frac{1}{2}.\]
这既给出了裴蜀等式，也证明任何公共因子都整除 $x-1$。最后，由 $2f=(f+g)+(f-g)$、$2g=(f+g)-(f-g)$，且标量域为实数域，$2$ 可逆，故两对多项式的公因子相同，最后一问的答案仍为 $x-1$。
\end{solve}

\begin{exercise}\label{prob:S-INTERP-1}
求次数不超过 $2$ 的实多项式 $p$，使 $p(-1)=5$、$p(0)=1$、$p(2)=5$。用待定系数法求解并说明唯一性；若任意多项式 $f$ 具有相同的三个函数值，求它除以 $(x+1)x(x-2)$ 的余式。
\end{exercise}
\noindent 编者补充：S-INTERP-1；相关主题见原讲义对应讲，题目数据与解答由本选编另行构造。
\begin{solve}
写 $p=c_0+c_1x+c_2x^2$，得到范德蒙方程组
\[\begin{pmatrix}1&-1&1\\1&0&0\\1&2&4\end{pmatrix}\begin{pmatrix}c_0\\c_1\\c_2\end{pmatrix}=\begin{pmatrix}5\\1\\5\end{pmatrix}.\]
系数行列式为 $6\ne0$，解为 $(c_0,c_1,c_2)^{\mathrm T}=\begin{pmatrix}1\\-2\\2\end{pmatrix}$，即 $p(x)=2x^{2}-2x+1$。若还有另一多项式满足条件，两者之差次数至多 $2$，却有三个不同零点，只能为零，因此解唯一。作带余除法 $f=q(x)(x+1)x(x-2)+r$，则 $\deg r<3$ 且 $r$ 与 $f$ 在三个插值点的值相同，故余式就是所求 $p$。
\end{solve}

\begin{exercise}\label{prob:S-GCD-2}
设 $f(x)=x^{2}-5x+6$，$g(x)=x^{2}-7x+10$。求首一最大公因式、首一最小公倍式，并找出多项式 $u,v$ 使 $uf+vg=\gcd(f,g)$。进一步求 $f+g$ 与 $f-g$ 的首一最大公因式。
\end{exercise}
\noindent 编者补充：S-GCD-2；相关主题见原讲义对应讲，题目数据与解答由本选编另行构造。
\begin{solve}
分解得到 $f=(x-2)(x-3)$、$g=(x-2)(x-5)$。三个根两两不同，所以
\[\gcd(f,g)=x-2,\qquad\operatorname{lcm}(f,g)=x^{3}-10x^{2}+31x-30.\]
相减得 $f-g=(2)(x-2)$，从而可以取常数多项式
\[u=\frac{1}{2},\qquad v=-\frac{1}{2}.\]
这既给出了裴蜀等式，也证明任何公共因子都整除 $x-2$。最后，由 $2f=(f+g)+(f-g)$、$2g=(f+g)-(f-g)$，且标量域为实数域，$2$ 可逆，故两对多项式的公因子相同，最后一问的答案仍为 $x-2$。
\end{solve}

\begin{exercise}\label{prob:S-INTERP-2}
求次数不超过 $2$ 的实多项式 $p$，使 $p(-1)=7$、$p(0)=2$、$p(3)=23$。用待定系数法求解并说明唯一性；若任意多项式 $f$ 具有相同的三个函数值，求它除以 $(x+1)x(x-3)$ 的余式。
\end{exercise}
\noindent 编者补充：S-INTERP-2；相关主题见原讲义对应讲，题目数据与解答由本选编另行构造。
\begin{solve}
写 $p=c_0+c_1x+c_2x^2$，得到范德蒙方程组
\[\begin{pmatrix}1&-1&1\\1&0&0\\1&3&9\end{pmatrix}\begin{pmatrix}c_0\\c_1\\c_2\end{pmatrix}=\begin{pmatrix}7\\2\\23\end{pmatrix}.\]
系数行列式为 $12\ne0$，解为 $(c_0,c_1,c_2)^{\mathrm T}=\begin{pmatrix}2\\-2\\3\end{pmatrix}$，即 $p(x)=3x^{2}-2x+2$。若还有另一多项式满足条件，两者之差次数至多 $2$，却有三个不同零点，只能为零，因此解唯一。作带余除法 $f=q(x)(x+1)x(x-3)+r$，则 $\deg r<3$ 且 $r$ 与 $f$ 在三个插值点的值相同，故余式就是所求 $p$。
\end{solve}

\begin{exercise}\label{prob:S-GCD-3}
设 $f(x)=x^{2}-7x+12$，$g(x)=x^{2}-9x+18$。求首一最大公因式、首一最小公倍式，并找出多项式 $u,v$ 使 $uf+vg=\gcd(f,g)$。进一步求 $f+g$ 与 $f-g$ 的首一最大公因式。
\end{exercise}
\noindent 编者补充：S-GCD-3；相关主题见原讲义对应讲，题目数据与解答由本选编另行构造。
\begin{solve}
分解得到 $f=(x-3)(x-4)$、$g=(x-3)(x-6)$。三个根两两不同，所以
\[\gcd(f,g)=x-3,\qquad\operatorname{lcm}(f,g)=x^{3}-13x^{2}+54x-72.\]
相减得 $f-g=(2)(x-3)$，从而可以取常数多项式
\[u=\frac{1}{2},\qquad v=-\frac{1}{2}.\]
这既给出了裴蜀等式，也证明任何公共因子都整除 $x-3$。最后，由 $2f=(f+g)+(f-g)$、$2g=(f+g)-(f-g)$，且标量域为实数域，$2$ 可逆，故两对多项式的公因子相同，最后一问的答案仍为 $x-3$。
\end{solve}

\begin{exercise}\label{prob:S-INTERP-3}
求次数不超过 $2$ 的实多项式 $p$，使 $p(-1)=9$、$p(0)=3$、$p(4)=59$。用待定系数法求解并说明唯一性；若任意多项式 $f$ 具有相同的三个函数值，求它除以 $(x+1)x(x-4)$ 的余式。
\end{exercise}
\noindent 编者补充：S-INTERP-3；相关主题见原讲义对应讲，题目数据与解答由本选编另行构造。
\begin{solve}
写 $p=c_0+c_1x+c_2x^2$，得到范德蒙方程组
\[\begin{pmatrix}1&-1&1\\1&0&0\\1&4&16\end{pmatrix}\begin{pmatrix}c_0\\c_1\\c_2\end{pmatrix}=\begin{pmatrix}9\\3\\59\end{pmatrix}.\]
系数行列式为 $20\ne0$，解为 $(c_0,c_1,c_2)^{\mathrm T}=\begin{pmatrix}3\\-2\\4\end{pmatrix}$，即 $p(x)=4x^{2}-2x+3$。若还有另一多项式满足条件，两者之差次数至多 $2$，却有三个不同零点，只能为零，因此解唯一。作带余除法 $f=q(x)(x+1)x(x-4)+r$，则 $\deg r<3$ 且 $r$ 与 $f$ 在三个插值点的值相同，故余式就是所求 $p$。
\end{solve}

\begin{exercise}\label{prob:S-GCD-4}
设 $f(x)=x^{2}-9x+20$，$g(x)=x^{2}-11x+28$。求首一最大公因式、首一最小公倍式，并找出多项式 $u,v$ 使 $uf+vg=\gcd(f,g)$。进一步求 $f+g$ 与 $f-g$ 的首一最大公因式。
\end{exercise}
\noindent 编者补充：S-GCD-4；相关主题见原讲义对应讲，题目数据与解答由本选编另行构造。
\begin{solve}
分解得到 $f=(x-4)(x-5)$、$g=(x-4)(x-7)$。三个根两两不同，所以
\[\gcd(f,g)=x-4,\qquad\operatorname{lcm}(f,g)=x^{3}-16x^{2}+83x-140.\]
相减得 $f-g=(2)(x-4)$，从而可以取常数多项式
\[u=\frac{1}{2},\qquad v=-\frac{1}{2}.\]
这既给出了裴蜀等式，也证明任何公共因子都整除 $x-4$。最后，由 $2f=(f+g)+(f-g)$、$2g=(f+g)-(f-g)$，且标量域为实数域，$2$ 可逆，故两对多项式的公因子相同，最后一问的答案仍为 $x-4$。
\end{solve}

\begin{exercise}\label{prob:S-INTERP-4}
求次数不超过 $2$ 的实多项式 $p$，使 $p(-1)=11$、$p(0)=4$、$p(5)=119$。用待定系数法求解并说明唯一性；若任意多项式 $f$ 具有相同的三个函数值，求它除以 $(x+1)x(x-5)$ 的余式。
\end{exercise}
\noindent 编者补充：S-INTERP-4；相关主题见原讲义对应讲，题目数据与解答由本选编另行构造。
\begin{solve}
写 $p=c_0+c_1x+c_2x^2$，得到范德蒙方程组
\[\begin{pmatrix}1&-1&1\\1&0&0\\1&5&25\end{pmatrix}\begin{pmatrix}c_0\\c_1\\c_2\end{pmatrix}=\begin{pmatrix}11\\4\\119\end{pmatrix}.\]
系数行列式为 $30\ne0$，解为 $(c_0,c_1,c_2)^{\mathrm T}=\begin{pmatrix}4\\-2\\5\end{pmatrix}$，即 $p(x)=5x^{2}-2x+4$。若还有另一多项式满足条件，两者之差次数至多 $2$，却有三个不同零点，只能为零，因此解唯一。作带余除法 $f=q(x)(x+1)x(x-5)+r$，则 $\deg r<3$ 且 $r$ 与 $f$ 在三个插值点的值相同，故余式就是所求 $p$。
\end{solve}

\begin{exercise}\label{prob:S-GCD-5}
设 $f(x)=x^{2}-11x+30$，$g(x)=x^{2}-13x+40$。求首一最大公因式、首一最小公倍式，并找出多项式 $u,v$ 使 $uf+vg=\gcd(f,g)$。进一步求 $f+g$ 与 $f-g$ 的首一最大公因式。
\end{exercise}
\noindent 编者补充：S-GCD-5；相关主题见原讲义对应讲，题目数据与解答由本选编另行构造。
\begin{solve}
分解得到 $f=(x-5)(x-6)$、$g=(x-5)(x-8)$。三个根两两不同，所以
\[\gcd(f,g)=x-5,\qquad\operatorname{lcm}(f,g)=x^{3}-19x^{2}+118x-240.\]
相减得 $f-g=(2)(x-5)$，从而可以取常数多项式
\[u=\frac{1}{2},\qquad v=-\frac{1}{2}.\]
这既给出了裴蜀等式，也证明任何公共因子都整除 $x-5$。最后，由 $2f=(f+g)+(f-g)$、$2g=(f+g)-(f-g)$，且标量域为实数域，$2$ 可逆，故两对多项式的公因子相同，最后一问的答案仍为 $x-5$。
\end{solve}

\begin{exercise}\label{prob:S-INTERP-5}
求次数不超过 $2$ 的实多项式 $p$，使 $p(-1)=13$、$p(0)=5$、$p(6)=209$。用待定系数法求解并说明唯一性；若任意多项式 $f$ 具有相同的三个函数值，求它除以 $(x+1)x(x-6)$ 的余式。
\end{exercise}
\noindent 编者补充：S-INTERP-5；相关主题见原讲义对应讲，题目数据与解答由本选编另行构造。
\begin{solve}
写 $p=c_0+c_1x+c_2x^2$，得到范德蒙方程组
\[\begin{pmatrix}1&-1&1\\1&0&0\\1&6&36\end{pmatrix}\begin{pmatrix}c_0\\c_1\\c_2\end{pmatrix}=\begin{pmatrix}13\\5\\209\end{pmatrix}.\]
系数行列式为 $42\ne0$，解为 $(c_0,c_1,c_2)^{\mathrm T}=\begin{pmatrix}5\\-2\\6\end{pmatrix}$，即 $p(x)=6x^{2}-2x+5$。若还有另一多项式满足条件，两者之差次数至多 $2$，却有三个不同零点，只能为零，因此解唯一。作带余除法 $f=q(x)(x+1)x(x-6)+r$，则 $\deg r<3$ 且 $r$ 与 $f$ 在三个插值点的值相同，故余式就是所求 $p$。
\end{solve}

\begin{exercise}\label{prob:S-GCD-6}
设 $f(x)=x^{2}-13x+42$，$g(x)=x^{2}-15x+54$。求首一最大公因式、首一最小公倍式，并找出多项式 $u,v$ 使 $uf+vg=\gcd(f,g)$。进一步求 $f+g$ 与 $f-g$ 的首一最大公因式。
\end{exercise}
\noindent 编者补充：S-GCD-6；相关主题见原讲义对应讲，题目数据与解答由本选编另行构造。
\begin{solve}
分解得到 $f=(x-6)(x-7)$、$g=(x-6)(x-9)$。三个根两两不同，所以
\[\gcd(f,g)=x-6,\qquad\operatorname{lcm}(f,g)=x^{3}-22x^{2}+159x-378.\]
相减得 $f-g=(2)(x-6)$，从而可以取常数多项式
\[u=\frac{1}{2},\qquad v=-\frac{1}{2}.\]
这既给出了裴蜀等式，也证明任何公共因子都整除 $x-6$。最后，由 $2f=(f+g)+(f-g)$、$2g=(f+g)-(f-g)$，且标量域为实数域，$2$ 可逆，故两对多项式的公因子相同，最后一问的答案仍为 $x-6$。
\end{solve}

\begin{exercise}\label{prob:S-INTERP-6}
求次数不超过 $2$ 的实多项式 $p$，使 $p(-1)=15$、$p(0)=6$、$p(7)=335$。用待定系数法求解并说明唯一性；若任意多项式 $f$ 具有相同的三个函数值，求它除以 $(x+1)x(x-7)$ 的余式。
\end{exercise}
\noindent 编者补充：S-INTERP-6；相关主题见原讲义对应讲，题目数据与解答由本选编另行构造。
\begin{solve}
写 $p=c_0+c_1x+c_2x^2$，得到范德蒙方程组
\[\begin{pmatrix}1&-1&1\\1&0&0\\1&7&49\end{pmatrix}\begin{pmatrix}c_0\\c_1\\c_2\end{pmatrix}=\begin{pmatrix}15\\6\\335\end{pmatrix}.\]
系数行列式为 $56\ne0$，解为 $(c_0,c_1,c_2)^{\mathrm T}=\begin{pmatrix}6\\-2\\7\end{pmatrix}$，即 $p(x)=7x^{2}-2x+6$。若还有另一多项式满足条件，两者之差次数至多 $2$，却有三个不同零点，只能为零，因此解唯一。作带余除法 $f=q(x)(x+1)x(x-7)+r$，则 $\deg r<3$ 且 $r$ 与 $f$ 在三个插值点的值相同，故余式就是所求 $p$。
\end{solve}

\section{特征值与特征向量}\label{sec:topic14}
本节收录 24 道题。特征向量必须非零。求得特征值后分别解齐次方程，核对代数重数与几何重数，不能把二者默认等同。
\begin{exercise}\label{prob:L14-A-1}
设$S,T\in \mathcal{L}(V)$满足$ST=TS$，证明：$\ker S$和$\operatorname{im} S$都在$T$下不变.
\end{exercise}
\noindent 来源：\href{https://github.com/yhwu-is/Linear-Algebra-Left-Undone/blob/new/%E8%AE%B2%E4%B9%89/%E4%B8%93%E9%A2%98/14%20%E7%9B%B8%E4%BC%BC%E6%A0%87%E5%87%86%E5%BD%A2%EF%BC%9A%E5%8A%A8%E6%9C%BA%E4%B8%8E%E5%9F%BA%E7%A1%80.tex}{L14-A-1}。
\begin{solve}
对于$\forall v\in\ker S,Sv=0\implies STv=TSv=0$，这表明$Tv\in \ker S$，$\ker S$在$T$下不变.
            对于$\forall v \in \operatorname{im} S,v=Su\implies Tv=TSu=S(Tu)$，这表明$Tv \in \operatorname{im} S$，$\operatorname{im} S$在$T$下不变.
\end{solve}

\begin{exercise}\label{prob:L14-A-2}
已知$\mathbf{R}^2$上线性变换$T$在基$e_1=(1,0),e_2=(0,1)$下的矩阵为$\begin{pmatrix}2 & 1 \\ 0 & 2\end{pmatrix}$. 证明：
        \begin{enumerate}
            \item 设$W_1$为由$e_1$张成的子空间，则$W_1$为$T$的不变子空间；

            \item $\mathbf{R}^2$不能表示为$T$的任何不变子空间$W_2$与$W_1$的直和.
        \end{enumerate}
\end{exercise}
\noindent 来源：\href{https://github.com/yhwu-is/Linear-Algebra-Left-Undone/blob/new/%E8%AE%B2%E4%B9%89/%E4%B8%93%E9%A2%98/14%20%E7%9B%B8%E4%BC%BC%E6%A0%87%E5%87%86%E5%BD%A2%EF%BC%9A%E5%8A%A8%E6%9C%BA%E4%B8%8E%E5%9F%BA%E7%A1%80.tex}{L14-A-2}。
\begin{solve}
\begin{enumerate}
                \item 对于$\forall v\in W_1$，有$v=ke_1,Tv=T(ke_1)=kT(e_1)=2ke_1\in W_1$，从而$W_1$为$T$的不变子空间.
                \item 容易看出$T$在自然基下的矩阵不可对角化，从而$\mathbf{R}^2$不能表示为$T$的任何不变子空间$W_2$与$W_1$的直和.
            \end{enumerate}
\end{solve}

\begin{exercise}\label{prob:L14-A-3}
定义线性变换$T\in \mathcal{L}(\mathbf{F}^2)$为$T(x,y)=(y,0)$. 令$U=\{(x,0) \mid x\in\mathbf{F}\}$. 证明：
        \begin{enumerate}
            \item $U$在$T$下不变，且$T|_{U}$是$U$上的零线性变换；

            \item 不存在$\mathbf{F}^2$的在$T$下不变的子空间$W$使得$\mathbf{F}^2=U\oplus W$；

            \item $T/U$是$\mathbf{F}^2/U$上的0线性变换.
        \end{enumerate}
\end{exercise}
\noindent 来源：\href{https://github.com/yhwu-is/Linear-Algebra-Left-Undone/blob/new/%E8%AE%B2%E4%B9%89/%E4%B8%93%E9%A2%98/14%20%E7%9B%B8%E4%BC%BC%E6%A0%87%E5%87%86%E5%BD%A2%EF%BC%9A%E5%8A%A8%E6%9C%BA%E4%B8%8E%E5%9F%BA%E7%A1%80.tex}{L14-A-3}。
\begin{solve}
\begin{enumerate}
                \item 对于$\forall u=(x,0)\in U$，$Tu=T(x,0)=(0,0)\in U$，因此$U$在$T$下不变，且$T|_{U}$是$U$上的零线性变换.
                \item $T$在自然基下的矩阵为$\begin{pmatrix}
                    0 & 1 \\ 0 & 0
                \end{pmatrix}$，它不可对角化.
                \item 对于$\forall v+U=(x,y)+U\in \mathbf{F}^2/U$，$(T/U)(v+U)=Tv+U=(y,0)+U=0+U$，从而$T/U$是$\mathbf{F}^2/U$上的$0$线性变换.
            \end{enumerate}
\end{solve}

\begin{exercise}\label{prob:L14-A-4}
设$V$是有限维的且$T\in \mathcal{L}(V)$，设$\lambda_1,\ldots,\lambda_m$是非零互异特征值，证明：
        \[ \dim E(\lambda_1,T)+\cdots+\dim E(\lambda_m,T)\leqslant\dim\operatorname{im} T. \]
\end{exercise}
\noindent 来源：\href{https://github.com/yhwu-is/Linear-Algebra-Left-Undone/blob/new/%E8%AE%B2%E4%B9%89/%E4%B8%93%E9%A2%98/14%20%E7%9B%B8%E4%BC%BC%E6%A0%87%E5%87%86%E5%BD%A2%EF%BC%9A%E5%8A%A8%E6%9C%BA%E4%B8%8E%E5%9F%BA%E7%A1%80.tex}{L14-A-4}。
\begin{solve}
考虑$U=E(\lambda_1,T)\oplus E(\lambda_2,T) \oplus \cdots \oplus E(\lambda_m,T)$，对于$\forall u\in U$，设$u=v_1+v_2+\cdots+v_m$，其中$v_i\in E(\lambda_i,T)$.则$v_i=T(\dfrac{v_i}{\lambda_i})$，从而$u=T(\sum\limits_{i=1}^{m}\dfrac{v_i}{\lambda_i}) \in \operatorname{im} T$.因此$U$是$\operatorname{im} T$的子空间.进而有$\dim U=\dim E(\lambda_1,T)+\cdots+\dim E(\lambda_m,T) \leqslant \dim \operatorname{im} T$.
\end{solve}

\begin{exercise}\label{prob:L14-A-5}
设$T\in \mathcal{L}(V)$且$\dim\operatorname{im} T=k$. 证明$T$至多有$k+1$个特征值.
\end{exercise}
\noindent 来源：\href{https://github.com/yhwu-is/Linear-Algebra-Left-Undone/blob/new/%E8%AE%B2%E4%B9%89/%E4%B8%93%E9%A2%98/14%20%E7%9B%B8%E4%BC%BC%E6%A0%87%E5%87%86%E5%BD%A2%EF%BC%9A%E5%8A%A8%E6%9C%BA%E4%B8%8E%E5%9F%BA%E7%A1%80.tex}{L14-A-5}。
\begin{solve}
设$T$的非零特征值为$\lambda_1,\ldots,\lambda_m$，则$m\leqslant \dim E(\lambda_1,T)+\cdots+\dim E(\lambda_m,T)\leqslant \dim\operatorname{im} T=k$.若$0$是$T$的非零特征值，则$T$至多有$k+1$个特征值.否则$T$至多有$k$个特征值.综上可知$T$至多有$k+1$个特征值.
\end{solve}

\begin{exercise}\label{prob:L14-A-6}
设$\sigma$是线性空间$\mathbf{R}[x]_3$上的线性变换，它在基$1,x,x^2$下的矩阵为
        \[ A=\begin{pmatrix}
                1 & 2 & 2 \\ 2 & 1 & 2 \\ 2 & 2 & 1
            \end{pmatrix}\]
        求$\sigma$的特征值与特征子空间.
\end{exercise}
\noindent 来源：\href{https://github.com/yhwu-is/Linear-Algebra-Left-Undone/blob/new/%E8%AE%B2%E4%B9%89/%E4%B8%93%E9%A2%98/14%20%E7%9B%B8%E4%BC%BC%E6%A0%87%E5%87%86%E5%BD%A2%EF%BC%9A%E5%8A%A8%E6%9C%BA%E4%B8%8E%E5%9F%BA%E7%A1%80.tex}{L14-A-6}。
\begin{solve}
令$|\lambda E-A|=0$，解得$\lambda_1=5,\lambda_2=-1$.然后解$(5E-A)X=0$，解得$X=k(1,1,1)^\mathrm{T}$.解$(-E-A)X=0$，解得$X=k_1(1,-1,0)^\mathrm{T}+k_2(1,0,-1)^\mathrm{T}$.因此$\lambda_1=5$对应的特征子空间为$\operatorname{span}(1+x+x^2)$，$\lambda_2=-1$对应的特征子空间为$\operatorname{span}(1-x,1-x^2)$.
\end{solve}

\begin{exercise}\label{prob:L14-A-7}
设$A,P$都是3阶方阵，$P$可逆，已知$A$的特征值$\lambda_1=1,\lambda_2=-1,\lambda_3=2$，$B=A^3-5A^2$，求$|B|$，$|A+5E|$，$|5E+P^{-1}AP|$.
\end{exercise}
\noindent 来源：\href{https://github.com/yhwu-is/Linear-Algebra-Left-Undone/blob/new/%E8%AE%B2%E4%B9%89/%E4%B8%93%E9%A2%98/14%20%E7%9B%B8%E4%BC%BC%E6%A0%87%E5%87%86%E5%BD%A2%EF%BC%9A%E5%8A%A8%E6%9C%BA%E4%B8%8E%E5%9F%BA%E7%A1%80.tex}{L14-A-7}。
\begin{solve}
矩阵的行列式等于其特征值之积，知$|A|=\lambda_1\lambda_2\lambda_3=-2$.
            进一步地，$A-5E$的特征值为$-4,-6,-3$，$A+5E$和$5E+P^{-1}AP$的特征值均为$6,4,7$.从而$|A-5E|=-72,|A+5E|=|5E+P^{-1}AP|=168$.
            $|B|=|A^2(A-5E)|=|A||A||A-5E|=-288$.
\end{solve}

\begin{exercise}\label{prob:L14-A-8}
设$A=\begin{pmatrix}
                a & -1 & c \\ 5 & b & 3 \\ 1-c & 0 & -a
            \end{pmatrix}$，$|A|=-1$，$\alpha=(-1,-1,1)^\mathrm{T}$为$\operatorname{adj}(A)$的特征向量，求$\operatorname{adj}(A)$对应于$\alpha$的特征值及$a,b,c$和$A$对应于$\alpha$的特征值$\mu$.
\end{exercise}
\noindent 来源：\href{https://github.com/yhwu-is/Linear-Algebra-Left-Undone/blob/new/%E8%AE%B2%E4%B9%89/%E4%B8%93%E9%A2%98/14%20%E7%9B%B8%E4%BC%BC%E6%A0%87%E5%87%86%E5%BD%A2%EF%BC%9A%E5%8A%A8%E6%9C%BA%E4%B8%8E%E5%9F%BA%E7%A1%80.tex}{L14-A-8}。
\begin{solve}
设$\operatorname{adj}(A)\alpha=\lambda \alpha$，则$A\operatorname{adj}(A)\alpha=\lambda A\alpha$.从而$\lambda A\alpha=|A|\alpha=-\alpha\implies A\alpha=-\dfrac{1}{\lambda}\alpha$，故$\alpha$为$A$的特征向量.

            由
\[A\alpha=\begin{pmatrix}
                a & -1 & c \\ 5 & b & 3 \\ 1-c & 0 & -a
            \end{pmatrix}\cdot\begin{pmatrix}
                -1 \\ -1 \\ 1
            \end{pmatrix}=\begin{pmatrix}
                -a+1+c \\ -2-b \\ c-1-a
            \end{pmatrix}=-\dfrac{1}{\lambda}\alpha\]
可知：
            $(-a+1+c)+(c-1-a)=0$，从而$a=c$，代入可知$b=-3,-\dfrac{1}{\lambda}=-1,\lambda=1$.
            此外有$A=\begin{pmatrix}
                a & -1 & a \\ 5 & -3 & 3 \\ 1-a & 0 & -a
            \end{pmatrix}$，由$|A|=-1$解得$a=2$，因此$c=a=2$.
\end{solve}

\begin{exercise}\label{prob:L14-A-9}
设$A,B\in \mathbf{M}_n(\mathbf{F})$，$AB=BA$，证明：若$X$是矩阵$A$属于特征值$\lambda_0$的特征向量，则$BX\in V_{\lambda_0}$（注：本题是解决很多$AB=BA$类问题的基础）.
\end{exercise}
\noindent 来源：\href{https://github.com/yhwu-is/Linear-Algebra-Left-Undone/blob/new/%E8%AE%B2%E4%B9%89/%E4%B8%93%E9%A2%98/14%20%E7%9B%B8%E4%BC%BC%E6%A0%87%E5%87%86%E5%BD%A2%EF%BC%9A%E5%8A%A8%E6%9C%BA%E4%B8%8E%E5%9F%BA%E7%A1%80.tex}{L14-A-9}。
\begin{solve}
由于$AX=\lambda_0 X$，故$A(BX)=(AB)X=(BA)X=B(AX)=B(\lambda_0X)=\lambda_0(BX)$，从而$BX\in V_{\lambda_0}$.
\end{solve}

\begin{exercise}\label{prob:L14-B-1}
设$\sigma\in \mathcal{L}(V,W)$，定义$\widetilde{\sigma}:(V/(\ker \sigma))\to W$如下：
        \[\widetilde{\sigma}(v+\ker\sigma)=\sigma(v).\]
        \begin{enumerate}
            \item $\widetilde{\sigma}$是良定义的，且是$(V/(\ker \sigma))$到$W$上的线性映射；

            \item $\widetilde{\sigma}$是单射；

            \item $\operatorname{im} \widetilde{\sigma}=\operatorname{im} \sigma$；

            \item $V/(\ker \sigma)$同构于$\operatorname{im} \sigma$.
        \end{enumerate}
\end{exercise}
\noindent 来源：\href{https://github.com/yhwu-is/Linear-Algebra-Left-Undone/blob/new/%E8%AE%B2%E4%B9%89/%E4%B8%93%E9%A2%98/14%20%E7%9B%B8%E4%BC%BC%E6%A0%87%E5%87%86%E5%BD%A2%EF%BC%9A%E5%8A%A8%E6%9C%BA%E4%B8%8E%E5%9F%BA%E7%A1%80.tex}{L14-B-1}。
\begin{solve}
\begin{enumerate}
                \item 对于$v_1+\ker\sigma=v_2+\ker\sigma$，有$v_1=v_2+u$，其中$u\in\ker\sigma$.因此$\widetilde{\sigma}(v_1+\ker\sigma)=\sigma(v_1)=\sigma(v_1+u)=\sigma(v_2)=\widetilde{\sigma}(v_2+\ker\sigma)$，从而$\widetilde{\sigma}$是良定义的.

                对于$\forall v_1+\ker\sigma,v_2+\ker\sigma \in V/\ker\sigma$和$l,h\in \mathbf{F}$，都有$\widetilde{\sigma}(l(v_1+\ker\sigma)+h(v_2+\ker\sigma))=\widetilde{\sigma}(lv_1+hv_2+\ker\sigma)=\sigma(lv_1+hv_2)=l\sigma(v_1)+h\sigma(v_2)=l\widetilde{\sigma}(v_1+\ker\sigma)+h\widetilde{\sigma}(v_2+\ker\sigma)$，这表明$\widetilde{\sigma}$是$V/(\ker\sigma)$到$W$上的线性映射.
                \item 对于$v_1+\ker\sigma,v_2+\ker\sigma$，若有$\sigma(v_1)=\sigma(v_2)$，则$v_1-v_2\in \ker\sigma$，从而$v_1+\ker\sigma=v_2+\ker\sigma$.因此$\widetilde{\sigma}$为单射.
                \item 对于$\forall u\in \operatorname{im}\widetilde{\sigma}$，有$u \in \operatorname{im}\sigma$，故$\operatorname{im} \widetilde{\sigma} \subseteq \operatorname{im}\sigma$.

                对于$\forall v=\operatorname{im}(w) \in\operatorname{im}\sigma$，$v=\sigma(w)=\widetilde{\sigma}(w+\ker\sigma)\in\operatorname{im}\widetilde{\sigma}$，从而$\operatorname{im}\sigma\subseteq \operatorname{im} \widetilde{\sigma}$.综上有$\operatorname{im} \widetilde{\sigma}=\operatorname{im}\sigma$.
                \item 由$(2)(3)$可知$\widetilde{\sigma}$实际上给出了一个从$V/(\ker \sigma)$到$\operatorname{im}\sigma$的一一映射，从而$V/(\ker\sigma)\cong \operatorname{im}\sigma$.
            \end{enumerate}
\end{solve}

\begin{exercise}\label{prob:L14-B-2}
设$T\in \mathcal{L}(V)$，证明：$T$是数乘变换（即$T=cI_V$，其中$c$为非零常数，$I_V$为$V$上的恒等映射）的充要条件是$V$的每一个一维子空间都是$T$的不变子空间.
\end{exercise}
\noindent 来源：\href{https://github.com/yhwu-is/Linear-Algebra-Left-Undone/blob/new/%E8%AE%B2%E4%B9%89/%E4%B8%93%E9%A2%98/14%20%E7%9B%B8%E4%BC%BC%E6%A0%87%E5%87%86%E5%BD%A2%EF%BC%9A%E5%8A%A8%E6%9C%BA%E4%B8%8E%E5%9F%BA%E7%A1%80.tex}{L14-B-2}。
\begin{solve}
必要性是显然的.下证充分性.
            设$v_1,v_2,\ldots,v_n$为$V$的一组基，则对于每个$1\leqslant i \leqslant n$，存在$c_i\in \mathbf{F}$，使得$T(v_i)=c_iv_i$.对于$\forall i\neq j$，$T(v_i+v_j)=c_iv_i+c_jv_j\in \operatorname{span}(v_i+v_j)$，这表明$c_i=c_j$.从而所有的$c_i$均相等，$T$为数乘变换.
\end{solve}

\begin{exercise}\label{prob:L14-B-3}
设$\sigma\in \mathcal{L}(V)$，证明：
        \begin{enumerate}
            \item $\sigma/(\operatorname{im} \sigma)=0$；

            \item $\sigma/(\ker \sigma)$是单射$\iff \ker \sigma\cap\operatorname{im} \sigma=\{0\}$.
        \end{enumerate}
\end{exercise}
\noindent 来源：\href{https://github.com/yhwu-is/Linear-Algebra-Left-Undone/blob/new/%E8%AE%B2%E4%B9%89/%E4%B8%93%E9%A2%98/14%20%E7%9B%B8%E4%BC%BC%E6%A0%87%E5%87%86%E5%BD%A2%EF%BC%9A%E5%8A%A8%E6%9C%BA%E4%B8%8E%E5%9F%BA%E7%A1%80.tex}{L14-B-3}。
\begin{solve}
\begin{enumerate}
                \item 对于$\forall v+\operatorname{im}\sigma$，$(\sigma/(\operatorname{im}\sigma))(v+\operatorname{im}\sigma)=\sigma(v)+\operatorname{im}\sigma=0+\operatorname{im}\sigma$，因此$\sigma/(\operatorname{im}\sigma)=0$.
                \item $\sigma/(\ker\sigma)$为单射$\iff \ker(\sigma/(\ker\sigma))=0+\ker\sigma\iff \ker\sigma\cap\operatorname{im}\sigma=\{0\}$.
            \end{enumerate}
\end{solve}

\begin{exercise}\label{prob:L14-B-4}
设$V$是有限维的，$T\in \mathcal{L}(V)$且$U$在$T$下不变. 证明：$T/U$的每个特征值均为$T$的特征值.
\end{exercise}
\noindent 来源：\href{https://github.com/yhwu-is/Linear-Algebra-Left-Undone/blob/new/%E8%AE%B2%E4%B9%89/%E4%B8%93%E9%A2%98/14%20%E7%9B%B8%E4%BC%BC%E6%A0%87%E5%87%86%E5%BD%A2%EF%BC%9A%E5%8A%A8%E6%9C%BA%E4%B8%8E%E5%9F%BA%E7%A1%80.tex}{L14-B-4}。
\begin{solve}
对于$T/U$的特征值$\lambda$，设$(T/U)(v+U)=\lambda(v+U)=\lambda v+U$.则$(T/U)(v+U)=Tv+U=\lambda v+U$.
\end{solve}

\begin{exercise}\label{prob:L14-B-5}
设$V$为$n$维复向量空间，$T\in \mathcal{L}(V)$，$T$在$V$的一组基$e_1,e_2,\ldots,e_n$下的矩阵为对角矩阵$\operatorname{diag}\{d_1,\ldots,d_n\}$，且$d_i\neq d_j\enspace(i\neq j)$.
        \begin{enumerate}
            \item 求$T$的所有一维不变子空间；

            \item 求$T$的所有不变子空间与个数.
        \end{enumerate}
\end{exercise}
\noindent 来源：\href{https://github.com/yhwu-is/Linear-Algebra-Left-Undone/blob/new/%E8%AE%B2%E4%B9%89/%E4%B8%93%E9%A2%98/14%20%E7%9B%B8%E4%BC%BC%E6%A0%87%E5%87%86%E5%BD%A2%EF%BC%9A%E5%8A%A8%E6%9C%BA%E4%B8%8E%E5%9F%BA%E7%A1%80.tex}{L14-B-5}。
\begin{solve}
\begin{enumerate}
                \item $T$的所有一维不变子空间为$\operatorname{span}(e_i),\enspace i=1,2,\ldots,n$.
                \item $T$的所有不变子空间为所有$\operatorname{span}(e_{i_1},e_{i_2},\ldots,e_{i_k})$，其中$(i_1,i_2,\ldots,i_k)\subset\{1,2,\ldots,n\}$.

                $T$的不变子空间个数为$2^n$.
            \end{enumerate}
\end{solve}

\begin{exercise}\label{prob:L14-B-6}
给定$\mathbf{R}$上的2维线性空间$V$上的算子$T$，其在一组基$\alpha_1,\alpha_2$下的矩阵为
        \[\begin{pmatrix}
                0 & 1 \\ 1-a & 0
            \end{pmatrix}.\]
        求$T$的所有不变子空间.
\end{exercise}
\noindent 来源：\href{https://github.com/yhwu-is/Linear-Algebra-Left-Undone/blob/new/%E8%AE%B2%E4%B9%89/%E4%B8%93%E9%A2%98/14%20%E7%9B%B8%E4%BC%BC%E6%A0%87%E5%87%86%E5%BD%A2%EF%BC%9A%E5%8A%A8%E6%9C%BA%E4%B8%8E%E5%9F%BA%E7%A1%80.tex}{L14-B-6}。
\begin{solve}
令$|\lambda I-T|=0$，若$a<1$，则$\lambda_1=\sqrt{1-a},\lambda_2=-\sqrt{1-a}$，相应的特征向量为$(\alpha_1+\sqrt{1-a}\alpha_2),(\alpha_1-\sqrt{1-a}\alpha_2)$.从而$T$的所有不变子空间为$\{0\},\operatorname{span}(\alpha_1+\sqrt{1-a}\alpha_2),\operatorname{span}(\alpha_1-\sqrt{1-a}\alpha_2),V$.

            若$a=1$，则$\lambda=0$，相应的特征向量为$\alpha_1$，从而$T$的所有不变子空间为$\{0\},\operatorname{span}(\alpha_1),V$.

            若$a>1$，则$T$在$\mathbf{R}$上无特征值，$T$的不变子空间只有$\{0\}$和$V$.
\end{solve}

\begin{exercise}\label{prob:L14-B-7}
设$A,B$都是$n$阶矩阵，且$r(A)+r(B)<n$，证明：$A,B$有相同的特征值和特征向量.
\end{exercise}
\noindent 来源：\href{https://github.com/yhwu-is/Linear-Algebra-Left-Undone/blob/new/%E8%AE%B2%E4%B9%89/%E4%B8%93%E9%A2%98/14%20%E7%9B%B8%E4%BC%BC%E6%A0%87%E5%87%86%E5%BD%A2%EF%BC%9A%E5%8A%A8%E6%9C%BA%E4%B8%8E%E5%9F%BA%E7%A1%80.tex}{L14-B-7}。
\begin{solve}
由于$r(A)+r(B)<n$，故$A$和$B$都有特征值$0$.

            由于$\dim(N(A)\cap N(B))=\dim(N(A))+\dim(N(B))-\dim(N(A)+N(B))\geqslant\dim(N(A))+\dim(N(B))-n=n-r(A)+n-r(B)-n=n-(r(A)+r(B))\ge 0$，这表明$N(A)\cap N(B)\neq \{0\}$，从而存在$A$和$B$有相同的特征向量（对应于特征值$0$）.
\end{solve}

\begin{exercise}\label{prob:L14-B-8}
设$\lambda_1,\lambda_2,\ldots,\lambda_n$是矩阵$A=(a_{ij})_{n\times n}$的$n$个特征值，证明：$\lambda_1^2,\lambda_2^2,\ldots,\lambda_n^2$是$A^2$的$n$个特征值，且$\displaystyle\sum_{i=1}^{n}\lambda_i^2=\sum_{j=1}^{n}\displaystyle\sum_{k=1}^{n}a_{jk}a_{kj}$.
\end{exercise}
\noindent 来源：\href{https://github.com/yhwu-is/Linear-Algebra-Left-Undone/blob/new/%E8%AE%B2%E4%B9%89/%E4%B8%93%E9%A2%98/14%20%E7%9B%B8%E4%BC%BC%E6%A0%87%E5%87%86%E5%BD%A2%EF%BC%9A%E5%8A%A8%E6%9C%BA%E4%B8%8E%E5%9F%BA%E7%A1%80.tex}{L14-B-8}。
\begin{solve}
对于$\lambda_i,i=1,2,\ldots,n$，存在$u_i\in V,u_i\neq 0$，使得$Au_i=\lambda_iu_i$，从而$A^2u_i=A(\lambda_iu_i)=\lambda_iA(u_i)=\lambda_i^2u_i$，这表明$\lambda_i^2$为$A^2$的特征值，从而$\lambda_1^2,\lambda_2^2,\ldots,\lambda_n^2$为$A^2$的$n$个特征值.进而知$\displaystyle \sum_{i=1}^{n}\lambda_i^2=tr(A^2)=\sum_{j=1}^{n}\displaystyle\sum_{k=1}^{n}a_{jk}a_{kj}$.
\end{solve}

\begin{exercise}\label{prob:L14-B-9}
设$A$为$n$阶矩阵，$X_1,X_2,X_3$为$n$元列向量，且$AX_1=kX_1\enspace(X_1\neq 0),AX_2=lX_1+kX_2,AX_3=lX_2+kX_3\enspace(l\neq 0)$. 证明：$X_1,X_2,X_3$线性无关.
\end{exercise}
\noindent 来源：\href{https://github.com/yhwu-is/Linear-Algebra-Left-Undone/blob/new/%E8%AE%B2%E4%B9%89/%E4%B8%93%E9%A2%98/14%20%E7%9B%B8%E4%BC%BC%E6%A0%87%E5%87%86%E5%BD%A2%EF%BC%9A%E5%8A%A8%E6%9C%BA%E4%B8%8E%E5%9F%BA%E7%A1%80.tex}{L14-B-9}。
\begin{solve}
设$a_1X_1+a_2X_2+a_3X_3=0$.在等式两边同时左乘$(A-kE)$可知：$a_2lX_1+a_3lX_2=0$.再左乘$(A-kE)$知$a_3l^2=0$.结合$l\neq 0$，有$a_3=0$.回代知$a_1=a_2=0$，从而$X_1,X_2,X_3$线性无关.
\end{solve}

\begin{exercise}\label{prob:L14-C-1}
设$A,B\in \mathbf{M}_n(\mathbf{C})$，$B$的特征多项式$f(\lambda)=|\lambda E-B|$. 证明：$f(A)$可逆的充要条件是$B$的任一特征值都不是$A$的特征值.
\end{exercise}
\noindent 来源：\href{https://github.com/yhwu-is/Linear-Algebra-Left-Undone/blob/new/%E8%AE%B2%E4%B9%89/%E4%B8%93%E9%A2%98/14%20%E7%9B%B8%E4%BC%BC%E6%A0%87%E5%87%86%E5%BD%A2%EF%BC%9A%E5%8A%A8%E6%9C%BA%E4%B8%8E%E5%9F%BA%E7%A1%80.tex}{L14-C-1}。
\begin{solve}
$f(\lambda)=|\lambda E-B|=(\lambda-\lambda_1)^{\dim E(\lambda_1,B)}\cdots(\lambda-\lambda_m)^{\dim E(\lambda_m,B)}$，故$f(A)$可逆$\iff (A-\lambda_1E)^{\dim E(\lambda_1,B)}\cdots(A-\lambda_mE)^{\dim E(\lambda_m,B)}$可逆$\iff \mid(A-\lambda_1E)^{\dim E(\lambda_1,B)}\cdots(A-\lambda_mE)^{\dim E(\lambda_m,B)}\mid \neq 0 \iff \mid(A-\lambda_1E)\mid^{\dim E(\lambda_1,B)}\cdots\mid(A-\lambda_mE)\mid^{\dim E(\lambda_m,B)}\neq 0 \iff \forall i=1,2,\ldots,m,|A-\lambda_mE|\neq 0 \iff B$的任一特征值都不是$A$的特征值.
\end{solve}

\begin{exercise}\label{prob:L14-C-2}
证明：若$AB=BA$，则$A$和$B$至少有一个共同的特征向量.
\end{exercise}
\noindent 来源：\href{https://github.com/yhwu-is/Linear-Algebra-Left-Undone/blob/new/%E8%AE%B2%E4%B9%89/%E4%B8%93%E9%A2%98/14%20%E7%9B%B8%E4%BC%BC%E6%A0%87%E5%87%86%E5%BD%A2%EF%BC%9A%E5%8A%A8%E6%9C%BA%E4%B8%8E%E5%9F%BA%E7%A1%80.tex}{L14-C-2}。
\begin{solve}
命题等价于$\sigma,\tau \in \mathcal{L}(V),\sigma\tau=\tau\sigma$，则$\sigma$和$\tau$至少有一个公共的特征向量.

            考虑$\sigma$的不变子空间$E(\lambda,\sigma)$，对于$v\in E(\lambda,\sigma)$，有$\sigma v=\lambda v \implies \sigma(\tau v)=\tau\sigma v=\lambda (\tau v)$，这表明$E(\lambda,\sigma)$是$\tau$的不变子空间.考虑限制映射$\tau|_E(\lambda,\sigma)$，它一定有特征值和相应的特征向量$u\in E(\lambda,\sigma),u \neq 0$，从而$u$为$\sigma$和$\tau$的公共特征向量.
\end{solve}

\begin{exercise}\label{prob:L14-例-4}
设$A=\begin{pmatrix}
            1 & -1 & 0 \\ 2 & 0 & 1 \\ 1 & a & 0
        \end{pmatrix}$，且存在非零向量$\alpha$使得$A\alpha=2\alpha$，求$a$.
\end{exercise}
\noindent 来源：\href{https://github.com/yhwu-is/Linear-Algebra-Left-Undone/blob/new/%E8%AE%B2%E4%B9%89/%E4%B8%93%E9%A2%98/14%20%E7%9B%B8%E4%BC%BC%E6%A0%87%E5%87%86%E5%BD%A2%EF%BC%9A%E5%8A%A8%E6%9C%BA%E4%B8%8E%E5%9F%BA%E7%A1%80.tex}{L14-例-4}。
\begin{solve}
由题意知2是矩阵$A$的特征值，因此我们有
    \[|2E-A|=\begin{vmatrix}
            1 & 1 & 0 \\ -2 & 2 & -1 \\ -1 & -a & 2
        \end{vmatrix}=9-a=0,\]
    因此$a=9$.
\end{solve}

\begin{exercise}\label{prob:L14-例-6}
对下列矩阵$A$的特征值，能做出怎样的断言？
    \begin{enumerate}
        \item $A$可逆/$A$不可逆/$E+A$可逆/$4E+A$不可逆；

        \item $|E - A^2| = 0$；

        \item $A^2 = E$（对合）/$A^2 = A$（幂等）/$A^k = 0$（幂零）；

        \item $A=\lambda_0E+B$（$\lambda_0$为常数，且已知$B$的$n$个特征值为$\lambda_1,\lambda_2,\ldots,\lambda_n$）；

        \item $A$为对角块矩阵，即$A=\operatorname{diag}(A_1,A_2,\ldots,A_m)$.
    \end{enumerate}
\end{exercise}
\noindent 来源：\href{https://github.com/yhwu-is/Linear-Algebra-Left-Undone/blob/new/%E8%AE%B2%E4%B9%89/%E4%B8%93%E9%A2%98/14%20%E7%9B%B8%E4%BC%BC%E6%A0%87%E5%87%86%E5%BD%A2%EF%BC%9A%E5%8A%A8%E6%9C%BA%E4%B8%8E%E5%9F%BA%E7%A1%80.tex}{L14-例-6}。
\begin{solve}
\begin{enumerate}
        \item $A$可逆时有$|A|=\lambda_1\cdots\lambda_n\neq 0$，因此$A$的特征值都不为0. 同理，$A$不可逆表明存在特征值等于0，$E+A$可逆表明$-1$不是$A$的特征值，$4E+A$不可逆表明$-4$是$A$的特征值.

        \item $|E - A^2| = |E - A||E + A| = 0$，因此 $A$ 的特征值可能为 $1$ 或 $-1$ 或二者兼有.

        \item 首先考虑对合矩阵，接下来的同理可以得到类似结论. 由于$A^2=E$，设$AX=\lambda X$，则$A^2X=\lambda^2X=X$，因此$\lambda^2=1$，即$\lambda=\pm 1$，因此$1$或$-1$是$A$的特征值.

              但这里需要强调的是，不同于前两问，前两问中我们都是说某些值是$A$的特征值，但无法保证$A$的特征值只能是某些值，但在本题这样给出矩阵方程的情况下，我们可以得到$A$的特征值只能是$\pm 1$，没有其他值. 我们用反证法，假设存在$\lambda_0\neq\pm 1$是$A$的特征值，即$AX=\lambda_0X$，则$A^2X=\lambda_0^2X\neq X$（因为$X$不是零向量），导出矛盾. 当然有同学可能会思考，$A$的特征值一定兼有$\pm 1$吗，事实上并非如此，例如$E$满足$E^2=E$，但其特征值只有1，$-E$满足$(-E)^2=E$，但其特征值只有$-1$. 并且利用下一讲对角化的结论可以知道（我们放在下一讲习题中供读者练习），满足$A^2=E$且特征值只有1的矩阵只能是$E$，特征值只有$-1$的矩阵只能是$-E$.

              注：本题解决过程中告诉我们一个解题技巧，如果看到$A$的多项式$f(A)=O$这种形式的表达式，事实上$A$的特征值只能是$f(\lambda)=0$的根，如上题中$f(A)=A^2-E$，则$f(\lambda)=\lambda^2-1$，因此$A$的特征值只能是$\pm 1$. 同理，幂等矩阵的特征值只能是0和1，幂零矩阵的特征值只能是0（这是一个重要的幂零矩阵等价条件，未来我们会再次遇到）. 这个技巧的证明放在接下来的原讲义“多项式零点为特征值”处 中.

        \item 设$BX=\lambda_iX_i(X_i\neq 0,\enspace i=1,\ldots,n)$，则$AX_i=\lambda_0X_i+BX_i=\lambda_0X_i+\lambda_iX_i=(\lambda_0+\lambda_i)X_i$，因此$\lambda_0+\lambda_i\enspace(i=1,\ldots,n)$都是$A$的特征值.

        \item \begin{align*}
                  |\lambda E-A| & =\begin{vmatrix}
                                       \lambda E_1-A_1 & 0               & \cdots & 0               \\
                                       0               & \lambda E_2-A_2 & \cdots & 0               \\
                                       \vdots          & \vdots          & \ddots & \vdots          \\
                                       0               & 0               & \cdots & \lambda E_m-A_m
                                   \end{vmatrix}
                                & =\prod_{i=1}^{m}|\lambda E_i-A_i|=0
              \end{align*}
              因此，$A_i\enspace(i=1,\ldots,m)$的特征值都是$A$的特征值.
    \end{enumerate}
\end{solve}

\begin{exercise}\label{prob:L14-例-8}
回答以下问题：
    \begin{enumerate}
        \item 设$\alpha=(1,0,-1)^\mathrm{T}$，且$A=\alpha\alpha^\mathrm{T}$，求$|6E-A^n|$；

        \item 设$A$为三阶矩阵，其特征值为$1,-2,-1$，求$|A|$，$\operatorname{adj}(A)+3E$的特征值，$(A^{-1})^2+2E$的特征值以及$|A^2-A+E|$；

        \item 设$A$为三阶矩阵，$A^2-A-2E=O$，$|A|=2$，求$|\operatorname{adj}(A)+3E|$；

        \item 设$A$为三阶矩阵，其特征值为$-1,-1,5$，求$A_{11}+A_{22}+A_{33}$；
    \end{enumerate}
\end{exercise}
\noindent 来源：\href{https://github.com/yhwu-is/Linear-Algebra-Left-Undone/blob/new/%E8%AE%B2%E4%B9%89/%E4%B8%93%E9%A2%98/14%20%E7%9B%B8%E4%BC%BC%E6%A0%87%E5%87%86%E5%BD%A2%EF%BC%9A%E5%8A%A8%E6%9C%BA%E4%B8%8E%E5%9F%BA%E7%A1%80.tex}{L14-例-8}。
\begin{solve}
\begin{enumerate}
        \item 事实上$A=\alpha\alpha^\mathrm{T}=\begin{pmatrix}
                      1 & 0 & -1 \\ 0 & 0 & 0 \\ -1 & 0 & 1
                  \end{pmatrix}$，由$|\lambda E-A|=0$解得$A$的特征值为$\lambda_1=\lambda_2=0,\lambda_3=2$，而根据$A^n$的特征值性质和原讲义“特征值的性质”处 可知，$6E-A^n$的特征值为$6-\lambda_1^n,6-\lambda_2^n,6-\lambda_3^n$，即$6,6,6-2^n$，因此$|6E-A^n|=6^2(6-2^n)=36(6-2^n)$.

        \item 由于$A$的特征值为$1,-2,-1$，因此$|A|=1\times(-2)\times(-1)=2$，而$\operatorname{adj}(A)$的特征值为$|A|\lambda^{-1}$，因此$\operatorname{adj}(A)$的特征值为$2,-1,-2$，故$\operatorname{adj}(A)+3E$的特征值为$\operatorname{adj}(A)$的特征值加3（根据原讲义“特征值的性质”处），即为$5,2,1$，又根据$A^{-1}$和$A^2$特征值的性质可知，$(A^{-1})^2+2E$的特征值为$1^2+2,(-1/2)^2+2,(-1)^2+2$，即为$3,9/4,3$，而$A^2-A+E$的特征值根据$f(\sigma)$特征值性质的讨论可知为$1^2-1+1,(-2)^2-(-2)+1,(-1)^2-(-1)+1$，即为$1,7,3$，因此$|A^2-A+E|=1\times 7\times 3=21$.

        \item 根据原讲义“多项式零点为特征值”处，$A$的特征值只能出自$-1$和$2$. 又$|A|=2$，且$A$为3阶矩阵，因此$A$的3个特征值必为$-1,-1,2$.

              又$\operatorname{adj}(A)$的特征值为$|A|\lambda^{-1}$，因此$\operatorname{adj}(A)$的特征值为$1,-2,-2$，又根据原讲义“特征值的性质”处 的结论，$\operatorname{adj}(A)+3E$的特征值为$\operatorname{adj}(A)$的特征值加3，即$\lambda_1=\lambda_2=1,\lambda_3=4$，故$|\operatorname{adj}(A)+3E|=\lambda_1\lambda_2\lambda_3=4$.

        \item 由题意知$|A|=5$，故$\operatorname{adj}(A)$的特征值为$|A|\lambda^{-1}$即为$\mu_1=\mu_2=-5,\mu_3=1$，而$A_{11}+A_{22}+A_{33}$就是$\operatorname{adj}(A)$的迹（即矩阵对角线元素之和），因此$A_{11}+A_{22}+A_{33}=\mu_1+\mu_2+\mu_3=-9$.
    \end{enumerate}
\end{solve}

\begin{exercise}\label{prob:L14-例-13}
设$\sigma\in\mathcal{L}(\mathbf{F}^2)$定义为$\sigma(w,z)=(-z,w)$.
    \begin{enumerate}
        \item 当$\mathbf{F}=\mathbf{R}$时，求$\sigma$的特征值和特征向量；

        \item 当$\mathbf{F}=\mathbf{C}$时，求$\sigma$的特征值和特征向量.
    \end{enumerate}
\end{exercise}
\noindent 来源：\href{https://github.com/yhwu-is/Linear-Algebra-Left-Undone/blob/new/%E8%AE%B2%E4%B9%89/%E4%B8%93%E9%A2%98/14%20%E7%9B%B8%E4%BC%BC%E6%A0%87%E5%87%86%E5%BD%A2%EF%BC%9A%E5%8A%A8%E6%9C%BA%E4%B8%8E%E5%9F%BA%E7%A1%80.tex}{L14-例-13}。
\begin{solve}
我们首先写出$\sigma$在任意一组基下的矩阵表示，为了方便，我们选取标准基$e_1=(1,0),e_2=(0,1)$，则矩阵表示为
    \[ A=\begin{pmatrix}
            0 & -1 \\ 1 & 0
        \end{pmatrix}. \]
    则其特征多项式$f(\lambda)=|\lambda E-A|=\lambda^2+1$，因此
    \begin{enumerate}
        \item 在实数域上无特征值和特征向量；

        \item 复数域上特征值为$\pm\mathrm{i}$，其中$\mathrm{i}$对应的特征向量为$(w,-w\mathrm{i})$，$-\mathrm{i}$对应的特征向量为$(w,w\mathrm{i})$，其中$w\in\mathbf{R}$且$w\neq 0$.
    \end{enumerate}
\end{solve}

\section{不变子空间与对角化}\label{sec:topic15}
本节收录 46 道题，其中 6 道为编者补充计算题。可对角化要求有特征向量基，并依赖所选数域；讨论多个矩阵时要区分分别可对角化与可同时对角化。
\begin{exercise}\label{prob:L15-A-1}
请举例：存在两个矩阵相抵但不相似.
\end{exercise}
\noindent 来源：\href{https://github.com/yhwu-is/Linear-Algebra-Left-Undone/blob/new/%E8%AE%B2%E4%B9%89/%E4%B8%93%E9%A2%98/15%20%E7%9B%B8%E4%BC%BC%E6%A0%87%E5%87%86%E5%BD%A2%EF%BC%9A%E5%A4%8D%E6%95%B0%E5%9F%9F%E4%B8%8A%E7%9A%84%E5%B0%9D%E8%AF%95%E4%B8%8E%E7%90%86%E8%AE%BA.tex}{L15-A-1}。
\begin{solve}
非常简单，可以举$\begin{pmatrix}
                1 & 0 \\ 0 & 1
            \end{pmatrix}$和$\begin{pmatrix}
                1 & 0 \\ 0 & 2
            \end{pmatrix}$的例子，也可以举一个可对角化另一个不可对角化的例子，因为上述情况矩阵对应的相似标准形不一样，故不相似.
\end{solve}

\begin{exercise}\label{prob:L15-A-2}
求矩阵
        \[A = \begin{pmatrix}
                0  & -1 & 1 \\
                -1 & 0  & 1 \\
                1  & 1  & 0
            \end{pmatrix}\]
        的所有特征值，对应的特征子空间，以及与 $A$ 相似的一个对角矩阵.
\end{exercise}
\noindent 来源：\href{https://github.com/yhwu-is/Linear-Algebra-Left-Undone/blob/new/%E8%AE%B2%E4%B9%89/%E4%B8%93%E9%A2%98/15%20%E7%9B%B8%E4%BC%BC%E6%A0%87%E5%87%86%E5%BD%A2%EF%BC%9A%E5%A4%8D%E6%95%B0%E5%9F%9F%E4%B8%8A%E7%9A%84%E5%B0%9D%E8%AF%95%E4%B8%8E%E7%90%86%E8%AE%BA.tex}{L15-A-2}。
\begin{solve}
解$|\lambda E-A|=0$可得特征值为$1,1,-2$. 然后解$(E-A)X=0$得$X=t_1(-1,1,0)^\mathrm{T}+t_2(1,0,1)^\mathrm{T}$；解$(-2E-A)X=0$，得$X=t(-1,-1,1)^\mathrm{T}$. 可见特征值1对应的特征子空间为$\operatorname{span}((-1,1,0)^\mathrm{T},(1,0,1)^\mathrm{T})$，特征值$-2$对应的特征子空间为 $\operatorname{span}((-1,-1,1)^\mathrm{T})$. 可知与$A$相似的对角矩阵为 $\operatorname{diag}(1,1,-2)$.
\end{solve}

\begin{exercise}\label{prob:L15-A-3}
设$A=\begin{pmatrix}
                a & b \\ c & d
            \end{pmatrix}$为二阶实矩阵.
        \begin{enumerate}
            \item 若$|A|<0$，问：$A$与对角矩阵是否相似；

            \item 若$ad-bc=1$，$|a+d|>2$，问：$A$是否可对角化.
        \end{enumerate}
\end{exercise}
\noindent 来源：\href{https://github.com/yhwu-is/Linear-Algebra-Left-Undone/blob/new/%E8%AE%B2%E4%B9%89/%E4%B8%93%E9%A2%98/15%20%E7%9B%B8%E4%BC%BC%E6%A0%87%E5%87%86%E5%BD%A2%EF%BC%9A%E5%A4%8D%E6%95%B0%E5%9F%9F%E4%B8%8A%E7%9A%84%E5%B0%9D%E8%AF%95%E4%B8%8E%E7%90%86%E8%AE%BA.tex}{L15-A-3}。
\begin{solve}
\begin{enumerate}
                \item $f(\lambda)=|\lambda E-A|=\begin{vmatrix}
                              \lambda-a & -b \\ -c & \lambda-d
                          \end{vmatrix}=\lambda^2-(a+d)\lambda+(ad-bc)=\lambda^2-(a+d)\lambda+|A|$. 由于$|A|<0$，因此判别式$\Delta=(a+d)^2-4|A|>0$，因此二阶矩阵$A$有两个互异特征值，故可对角化，因此与对角矩阵相似；

                \item 同 (1)，判别式$\Delta=(a+d)^2-4|A|>0$，因此二阶矩阵$A$有两个不同的特征值，故可对角化.
            \end{enumerate}
\end{solve}

\begin{exercise}\label{prob:L15-A-4}
设$A=\begin{pmatrix}
                1 & 1 & a \\
                1 & a & 1 \\
                a & 1 & 1
            \end{pmatrix}$，$\beta=\begin{pmatrix}
                1 \\ 1 \\ -2
            \end{pmatrix}$，方程组$AX=\beta$有解但不唯一.
        \begin{enumerate}
            \item 求$a$的值；

            \item 求可逆矩阵$P$使得$P^{-1}AP$为对角矩阵.
        \end{enumerate}
\end{exercise}
\noindent 来源：\href{https://github.com/yhwu-is/Linear-Algebra-Left-Undone/blob/new/%E8%AE%B2%E4%B9%89/%E4%B8%93%E9%A2%98/15%20%E7%9B%B8%E4%BC%BC%E6%A0%87%E5%87%86%E5%BD%A2%EF%BC%9A%E5%A4%8D%E6%95%B0%E5%9F%9F%E4%B8%8A%E7%9A%84%E5%B0%9D%E8%AF%95%E4%B8%8E%E7%90%86%E8%AE%BA.tex}{L15-A-4}。
\begin{solve}
\begin{enumerate}
              \item 由Cramer法则，$A$为方阵且$AX=\beta$有解但不唯一，即$|A|\neq 0$，解得$a=-2$或$a=1$，分别代入可知$a=1$时增广矩阵的秩为2，而$r(A)=1$，故无解；$a=-2$时增广矩阵的秩为2，而$r(A)=2$，故有解，故$a=-2$.

              \item 容易求得特征值为$0,-3,3$，求特征向量可知$P=\begin{pmatrix}
                            1 & -1 & 1 \\ 1 & 0 & -2 \\ 1 & 1 & 1
                        \end{pmatrix}$.
          \end{enumerate}
\end{solve}

\begin{exercise}\label{prob:L15-A-5}
设$A$为三阶矩阵，$\alpha_1,\alpha_2,\alpha_3$线性无关，且$A\alpha_1=\alpha_1,A\alpha_2=\alpha_1+\alpha_2-2\alpha_3,A\alpha_3=\alpha_1-2\alpha_2+\alpha_3$，求$A$的特征值.
\end{exercise}
\noindent 来源：\href{https://github.com/yhwu-is/Linear-Algebra-Left-Undone/blob/new/%E8%AE%B2%E4%B9%89/%E4%B8%93%E9%A2%98/15%20%E7%9B%B8%E4%BC%BC%E6%A0%87%E5%87%86%E5%BD%A2%EF%BC%9A%E5%A4%8D%E6%95%B0%E5%9F%9F%E4%B8%8A%E7%9A%84%E5%B0%9D%E8%AF%95%E4%B8%8E%E7%90%86%E8%AE%BA.tex}{L15-A-5}。
\begin{solve}
由分块矩阵乘法，
\[A(\alpha_1,\alpha_2,\alpha_3)=(A\alpha_1,A\alpha_2,A\alpha_3)=(\alpha_1,\alpha_1+\alpha_2-2\alpha_3,\alpha_1-2\alpha_2+\alpha_3)=(\alpha_1,\alpha_2,\alpha_3)\begin{pmatrix}
                1 & 1 & 1 \\ 0 & 1 & -2 \\ 0 & -2 & 1
            \end{pmatrix}\]
. 令$P=(\alpha_1,\alpha_2,\alpha_3)$，则$AP=P\begin{pmatrix}
                1 & 1 & 1 \\ 0 & 1 & -2 \\ 0 & -2 & 1
            \end{pmatrix}=PB$，即$P^{-1}AP=B$，故$A$与$B$相似，故两矩阵的特征值相同，简单求解$B$的特征值即得$A$的特征值为$1,-1,3$.
\end{solve}

\begin{exercise}\label{prob:L15-A-6}
设三阶实对称矩阵$A$的各行元素之和为3，向量$\alpha_1=(-1,2,-1)^\mathrm{T}$，$\alpha_2=(0,-1,1)^\mathrm{T}$是方程组$AX=0$的两个解，求矩阵$P$使得$P^{-1}AP$为对角矩阵.
\end{exercise}
\noindent 来源：\href{https://github.com/yhwu-is/Linear-Algebra-Left-Undone/blob/new/%E8%AE%B2%E4%B9%89/%E4%B8%93%E9%A2%98/15%20%E7%9B%B8%E4%BC%BC%E6%A0%87%E5%87%86%E5%BD%A2%EF%BC%9A%E5%A4%8D%E6%95%B0%E5%9F%9F%E4%B8%8A%E7%9A%84%E5%B0%9D%E8%AF%95%E4%B8%8E%E7%90%86%E8%AE%BA.tex}{L15-A-6}。
\begin{solve}
每行元素之和为3，则我们知道$\alpha=(1,1,1)^\mathrm{T}$是3对应的特征向量（具体理由只需要验证$A\alpha=3\alpha$即可，这是很常用的性质），而$AX=0$的解是对应特征值0的特征向量，故由题意可知取$P=\begin{pmatrix}
                1 & -1 & 0 \\ 1 & 2 & -1 \\ 1 & -1 & 1
            \end{pmatrix}$，则$P^{-1}AP=\begin{pmatrix}
                3 & 0 & 0 \\ 0 & 0 & 0 \\ 0 & 0 & 0
            \end{pmatrix}$.
\end{solve}

\begin{exercise}\label{prob:L15-B-1}
设$a\neq b$，且$(aE-A)(bE-A)=O$. 证明：$A$可对角化（特例：对合矩阵）；
\end{exercise}
\noindent 来源：\href{https://github.com/yhwu-is/Linear-Algebra-Left-Undone/blob/new/%E8%AE%B2%E4%B9%89/%E4%B8%93%E9%A2%98/15%20%E7%9B%B8%E4%BC%BC%E6%A0%87%E5%87%86%E5%BD%A2%EF%BC%9A%E5%A4%8D%E6%95%B0%E5%9F%9F%E4%B8%8A%E7%9A%84%E5%B0%9D%E8%AF%95%E4%B8%8E%E7%90%86%E8%AE%BA.tex}{L15-B-1}。
\begin{solve}
回忆例题中$A^2=2A$，即$A(A-2E)=O$的证明，我们知道$A$的特征值只能为$a$或$b$，然后我们可以证明$r(aE-A)+r(bE-A)\leqslant n$和$r(aE-A)+r(bE-A)\geqslant n$，因此$r(aE-A)+r(bE-A)=n$，即$A$可对角化（由于和例题完全类似，这里不再展开具体做法）.
\end{solve}

\begin{exercise}\label{prob:L15-B-2}
证明：满足$A^2=E$且特征值只有1的矩阵只能是$E$，特征值只有$-1$的矩阵只能是$-E$.
\end{exercise}
\noindent 来源：\href{https://github.com/yhwu-is/Linear-Algebra-Left-Undone/blob/new/%E8%AE%B2%E4%B9%89/%E4%B8%93%E9%A2%98/15%20%E7%9B%B8%E4%BC%BC%E6%A0%87%E5%87%86%E5%BD%A2%EF%BC%9A%E5%A4%8D%E6%95%B0%E5%9F%9F%E4%B8%8A%E7%9A%84%E5%B0%9D%E8%AF%95%E4%B8%8E%E7%90%86%E8%AE%BA.tex}{L15-B-2}。
\begin{solve}
若$A^2=E$，则$A$可对角化.设$A=P^{-1}BP$，$B$为对角矩阵，若$A$的特征值只有$1$，则$B=E$，从而$A=P^{-1}EP=E$.若$A$的特征值只有$-1$，则$B=-E$，从而$A=-E$.
\end{solve}

\begin{exercise}\label{prob:L15-B-3}
设$A$为三阶矩阵，$A^2=A$且$r(A)=r$，求$|A-2E|$.
\end{exercise}
\noindent 来源：\href{https://github.com/yhwu-is/Linear-Algebra-Left-Undone/blob/new/%E8%AE%B2%E4%B9%89/%E4%B8%93%E9%A2%98/15%20%E7%9B%B8%E4%BC%BC%E6%A0%87%E5%87%86%E5%BD%A2%EF%BC%9A%E5%A4%8D%E6%95%B0%E5%9F%9F%E4%B8%8A%E7%9A%84%E5%B0%9D%E8%AF%95%E4%B8%8E%E7%90%86%E8%AE%BA.tex}{L15-B-3}。
\begin{solve}
根据例题中关于$A^2=A$的讨论可知，满足$A^2=A$的矩阵$A$可对角化（设过渡矩阵为$P$），特征值0和1的重数分别为$n-r$和$r$，因此$|A-2E|=|P^{-1}||A-2E||P|=|\operatorname{diag}(-1,\ldots,-1,-2,\ldots,-2)|=(-1)^r(-2)^{n-r}$.
\end{solve}

\begin{exercise}\label{prob:L15-B-4}
设$T\in \mathcal{L}(\mathbf{C}^3)$使得6和7是$T$的特征值，且$T$不可对角化. 证明：存在$(x,y,z)\in\mathbf{C}^3$使得$T(x,y,z)=(17+8x,\sqrt{5}+8y,2\pi+8z)$.
\end{exercise}
\noindent 来源：\href{https://github.com/yhwu-is/Linear-Algebra-Left-Undone/blob/new/%E8%AE%B2%E4%B9%89/%E4%B8%93%E9%A2%98/15%20%E7%9B%B8%E4%BC%BC%E6%A0%87%E5%87%86%E5%BD%A2%EF%BC%9A%E5%A4%8D%E6%95%B0%E5%9F%9F%E4%B8%8A%E7%9A%84%E5%B0%9D%E8%AF%95%E4%B8%8E%E7%90%86%E8%AE%BA.tex}{L15-B-4}。
\begin{solve}
三阶矩阵如果有三个特征值则一定可对角化，故由题意$T$只能有6和7两个特征值. 故$\lambda\neq 6,7$时$T-\lambda I$都是可逆的（回顾正文中的定理），故$T-8I$可逆（故是满射），因此一定存在$(x,y,z)\in\mathbf{C}^3$使得$(T-8I)(x,y,z)=(17,\sqrt{5},2\pi)$，因此存在$(x,y,z)\in\mathbf{C}^3$使得$T(x,y,z)=(17+8x,\sqrt{5}+8y,2\pi+8z)$.
\end{solve}

\begin{exercise}\label{prob:L15-B-5}
证明：两个可对角化的同阶矩阵特征值相同（包括重数）等价于它们相似. 对于不可对角化的矩阵来说，这一结论还成立吗？
\end{exercise}
\noindent 来源：\href{https://github.com/yhwu-is/Linear-Algebra-Left-Undone/blob/new/%E8%AE%B2%E4%B9%89/%E4%B8%93%E9%A2%98/15%20%E7%9B%B8%E4%BC%BC%E6%A0%87%E5%87%86%E5%BD%A2%EF%BC%9A%E5%A4%8D%E6%95%B0%E5%9F%9F%E4%B8%8A%E7%9A%84%E5%B0%9D%E8%AF%95%E4%B8%8E%E7%90%86%E8%AE%BA.tex}{L15-B-5}。
\begin{solve}
可对角化矩阵特征值（包括重数）相等则相似标准形（对角矩阵）相等，故一定相似. 不可对角化的举例说明不成立：$\begin{pmatrix}
                0 & 1 & 0 \\ 0 & 0 & 1 \\ 0 & 0 & 0
            \end{pmatrix}$和$\begin{pmatrix}
                0 & 1 & 0 \\ 0 & 0 & 0 \\ 0 & 0 & 0
            \end{pmatrix}$的特征值全为0（重数均为3），但它们对应不同的若当标准形（之后章节会介绍），因此不相似.
\end{solve}

\begin{exercise}\label{prob:L15-B-6}
设$A=\begin{pmatrix}
                2 & 2 & 0 \\ 8 & 2 & a \\ 0 & 0 & 6
            \end{pmatrix}$相似于对角矩阵，求常数$a$，并求可逆矩阵$P$使得$P^{-1}AP$为对角矩阵.
\end{exercise}
\noindent 来源：\href{https://github.com/yhwu-is/Linear-Algebra-Left-Undone/blob/new/%E8%AE%B2%E4%B9%89/%E4%B8%93%E9%A2%98/15%20%E7%9B%B8%E4%BC%BC%E6%A0%87%E5%87%86%E5%BD%A2%EF%BC%9A%E5%A4%8D%E6%95%B0%E5%9F%9F%E4%B8%8A%E7%9A%84%E5%B0%9D%E8%AF%95%E4%B8%8E%E7%90%86%E8%AE%BA.tex}{L15-B-6}。
\begin{solve}
直接计算$|\lambda E-A|=0$，得到$A$的特征值为$-2,6,6$，因为$A$可对角化，则特征值6对应的特征子空间为2维，即$(6E-A)X=0$的解空间为2维，根据齐次线性方程组一般理论，$r(6E-A)=3-2=1$，即$r\begin{pmatrix}
                4 & -2 & 0 \\ -8 & 4 & -a \\ 0 & 0 & 0
            \end{pmatrix}=1$，显然$a=0$. 接下来求解-2和6对应的特征向量即可，得到可行的解$P=\begin{pmatrix}
                -1 & 1 & 0 \\ 2 & 2 & 0 \\ 0 & 0 & 1
            \end{pmatrix}$.
\end{solve}

\begin{exercise}\label{prob:L15-B-7}
设$A=(a_{ij})_{n\times n}$是上三角矩阵.
        \begin{enumerate}
            \item 求$A$的全部特征值；

            \item 若$A$主对角元互不相等，证明：$A$与对角阵相似；

            \item 若$n$个主对角元相等且$A$不为对角矩阵，证明：$A$不与对角阵相似.
        \end{enumerate}
\end{exercise}
\noindent 来源：\href{https://github.com/yhwu-is/Linear-Algebra-Left-Undone/blob/new/%E8%AE%B2%E4%B9%89/%E4%B8%93%E9%A2%98/15%20%E7%9B%B8%E4%BC%BC%E6%A0%87%E5%87%86%E5%BD%A2%EF%BC%9A%E5%A4%8D%E6%95%B0%E5%9F%9F%E4%B8%8A%E7%9A%84%E5%B0%9D%E8%AF%95%E4%B8%8E%E7%90%86%E8%AE%BA.tex}{L15-B-7}。
\begin{solve}
\begin{enumerate}
                \item 上三角矩阵特征值就是所有主对角元$a_{ii},\enspace i=1,2,\ldots,n$；

                \item 由（1）知主对角元互不相等表示$A$有$n$个互不相等的特征值，故可对角化；

                \item 由（1）知此时$A$有一个$n$重特征值，记为$c$，则所有特征向量都在$(cE-A)X=0$的解空间中，而$r(cE-A)=1$，故解空间维数为$n-1<n$，故$A$不可对角化.
            \end{enumerate}
\end{solve}

\begin{exercise}\label{prob:L15-B-8}
已知$\mathbf{R}^3$的一个线性变换
        \[\sigma(x_1,x_2,x_3)=(2x_1-2x_2,-2x_1+x_2-2x_3,-2x_2).\]
        \begin{enumerate}
            \item 求$\sigma$关于自然基$\{e_1,e_2,e_3\}$所对应的矩阵$A$；

            \item 求$\sigma$关于基$\{(1,1,1),(0,1,1),(0,0,1)\}$所对应的矩阵$B$；

            \item 求矩阵$C_1$，使$C_1^{-1}BC_1=A$.
        \end{enumerate}
\end{exercise}
\noindent 来源：\href{https://github.com/yhwu-is/Linear-Algebra-Left-Undone/blob/new/%E8%AE%B2%E4%B9%89/%E4%B8%93%E9%A2%98/15%20%E7%9B%B8%E4%BC%BC%E6%A0%87%E5%87%86%E5%BD%A2%EF%BC%9A%E5%A4%8D%E6%95%B0%E5%9F%9F%E4%B8%8A%E7%9A%84%E5%B0%9D%E8%AF%95%E4%B8%8E%E7%90%86%E8%AE%BA.tex}{L15-B-8}。
\begin{solve}
\begin{enumerate}
                \item $A=\begin{pmatrix}
                              2 & -2 & 0 \\ -2 & 1 & -2 \\ 0 & -2 & 0
                          \end{pmatrix}$；

                \item $B=\begin{pmatrix}
                              0 & -2 & 0 \\ -3 & 1 & -2 \\ 1 & -1 & 2
                          \end{pmatrix}$；

                \item 根据讲义本讲开头的定理（线性变换在不同基下的表示），$C_1$实际上就是基$\alpha_1,\alpha_2,\alpha_3$变为自然基$e_1,e_2,e_3$的过渡矩阵，简单计算可得$\begin{pmatrix}
                              1 & 0 & 0 \\ -1 & 1 & 0 \\ 0 & -1 & 1
                          \end{pmatrix}$.
            \end{enumerate}
\end{solve}

\begin{exercise}\label{prob:L15-B-9}
设$A=\begin{pmatrix}
                0 & 0 & 1 \\ 0 & 0 & 0 \\ 1 & 0 & 0
            \end{pmatrix},B=\begin{pmatrix}
                1 & 0 & 0 \\ 0 & 1 & 2 \\ 0 & -1 & -2
            \end{pmatrix}$，证明：$A\sim B$，并求可逆矩阵$P$使得$P^{-1}AP=B$.
\end{exercise}
\noindent 来源：\href{https://github.com/yhwu-is/Linear-Algebra-Left-Undone/blob/new/%E8%AE%B2%E4%B9%89/%E4%B8%93%E9%A2%98/15%20%E7%9B%B8%E4%BC%BC%E6%A0%87%E5%87%86%E5%BD%A2%EF%BC%9A%E5%A4%8D%E6%95%B0%E5%9F%9F%E4%B8%8A%E7%9A%84%E5%B0%9D%E8%AF%95%E4%B8%8E%E7%90%86%E8%AE%BA.tex}{L15-B-9}。
\begin{solve}
不难解得$A$和$B$的特征值均为$-1,0,1$，因此它们都与$\begin{pmatrix}
                -1 & 0 & 0 \\ 0 & 0 & 0 \\ 0 & 0 & 1
            \end{pmatrix}$相似，即$A$与$B$相似. 我们也可以解得$P_1=\begin{pmatrix}
                -1 & 0 & 1 \\ 0 & 1 & 0 \\ 1 & 0 & 1
            \end{pmatrix}$和$P_2=\begin{pmatrix}
                0 & 0 & 1 \\ -1 & -2 & 0 \\ 1 & 1 & 0
            \end{pmatrix}$使得$P_1^{-1}AP_1=P_2^{-1}BP_2=\begin{pmatrix}
                -1 & 0 & 0 \\ 0 & 0 & 0 \\ 0 & 0 & 1
            \end{pmatrix}$，故$(P_1P_2^{-1})^{-1}A(P_1P_2^{-1})=B$，即题目要求的$P=P_1P_2^{-1}=\begin{pmatrix}
                1 & -1 & -2 \\ 0 & -1 & -1 \\ 1 & 1 & 2
            \end{pmatrix}$.
\end{solve}

\begin{exercise}\label{prob:L15-B-10}
已知$A=\begin{pmatrix}
                2 & 0 & 0 \\ 0 & 0 & 1 \\ 0 & 1 & x
            \end{pmatrix}$与$B=\begin{pmatrix}
                2 & 0 & 0 \\ 0 & y & 0 \\ 0 & 0 & -1
            \end{pmatrix}$相似.
        \begin{enumerate}
            \item 求$x$和$y$；

            \item 求一个可逆矩阵$P$，使$P^{-1}AP$为对角矩阵.
        \end{enumerate}
\end{exercise}
\noindent 来源：\href{https://github.com/yhwu-is/Linear-Algebra-Left-Undone/blob/new/%E8%AE%B2%E4%B9%89/%E4%B8%93%E9%A2%98/15%20%E7%9B%B8%E4%BC%BC%E6%A0%87%E5%87%86%E5%BD%A2%EF%BC%9A%E5%A4%8D%E6%95%B0%E5%9F%9F%E4%B8%8A%E7%9A%84%E5%B0%9D%E8%AF%95%E4%B8%8E%E7%90%86%E8%AE%BA.tex}{L15-B-10}。
\begin{solve}
\begin{enumerate}
                \item $A$与$B$相似则特征值相等，我们知道矩阵的对角线元素之和等于特征值之和，行列式等于特征值之积，则它们也一定相等. 根据这一原理非常容易解得$x=0,y=1$.

                \item 事实上$B$就是对角矩阵，实际上这里就是求过渡矩阵$P$使$A$对角化，具体步骤不再赘述，得到$P=\begin{pmatrix}
                              1 & 0 & 0 \\ 0 & 1 & 1 \\ 0 & 1 & -1
                          \end{pmatrix}$.
            \end{enumerate}
\end{solve}

\begin{exercise}\label{prob:L15-B-11}
设$A=\begin{pmatrix}
                1 & 2 & 0 & 0  & 0 \\ 4 & 3 & 0 & 0 & 0 \\ 0 & 0 & 1 & -3 & 3 \\ 0 & 0 & 3 & -5 & 3 \\
                0 & 0 & 6 & -6 & 4
            \end{pmatrix}$，求$A$的特征值. 若$A$可对角化，求可逆矩阵$P$，使$P^{-1}AP$为对角矩阵.
\end{exercise}
\noindent 来源：\href{https://github.com/yhwu-is/Linear-Algebra-Left-Undone/blob/new/%E8%AE%B2%E4%B9%89/%E4%B8%93%E9%A2%98/15%20%E7%9B%B8%E4%BC%BC%E6%A0%87%E5%87%86%E5%BD%A2%EF%BC%9A%E5%A4%8D%E6%95%B0%E5%9F%9F%E4%B8%8A%E7%9A%84%E5%B0%9D%E8%AF%95%E4%B8%8E%E7%90%86%E8%AE%BA.tex}{L15-B-11}。
\begin{solve}
令$A_1=\begin{pmatrix}
                1 & 2 \\ 4 & 3
            \end{pmatrix}$，$A_2=\begin{pmatrix}
                1 & -3 & 3 \\ 3 & -5 & 3 \\ 6 & -6 & 4
            \end{pmatrix}$，则$A=\begin{pmatrix}
                A_1 & O \\ O & A_2
            \end{pmatrix}$. 首先得到$A_1$和$A_2$的特征值分别为$-1,5$和$-2,-2,4$. 分别求过渡矩阵$P_1$和$P_2$使得$A_1$，$A_2$对角化，容易解得$P_1=\begin{pmatrix}
                -1 & 1 \\ 1 & 2
            \end{pmatrix}$，$P_2=\begin{pmatrix}
                1 & 1 & 1 \\ 1 & 0 & 1 \\ 0 & -1 & 2
            \end{pmatrix}$. 令$P=\begin{pmatrix}
                P_1 & O \\ O & P_2
            \end{pmatrix}$，则
        \[P^{-1}AP=\begin{pmatrix}
                P_1^{-1}A_1P_1 & O \\ O & P_2^{-1}A_2P_2
            \end{pmatrix}=\begin{pmatrix}
                -1 & 0 & 0 & 0 & 0 \\ 0 & 5 & 0 & 0 & 0 \\ 0 & 0 & -2 & 0 & 0 \\ 0 & 0 & 0 & -2 & 0 \\ 0 & 0 & 0 & 0 & 4
            \end{pmatrix},\]
        即$P^{-1}AP=\operatorname{diag}(-1,5,-2,-2,4)$，故$A$与$\operatorname{diag}(-1,5,-2,-2,4)$相似. 总之，$A$的特征值为$-1,5,-2,-2,4$，过渡矩阵
\[P=\begin{pmatrix}
                P_1 & O \\ O & P_2
            \end{pmatrix}=\begin{pmatrix}
                -1 & 1 & 0 & 0 & 0 \\ 1 & 2 & 0 & 0 & 0 \\ 0 & 0 & 1 & 1 & 1 \\ 0 & 0 & 1 & 0 & 1 \\ 0 & 0 & 0 & -1 & 2
            \end{pmatrix}\]
.

        （注：通过本题也可以看出分块对角矩阵的特征值是两个对角块矩阵的并集，过渡矩阵就是两个过渡矩阵按同样分块方式排列得到的矩阵.）
\end{solve}

\begin{exercise}\label{prob:L15-B-12}
设三阶矩阵$A$的特征值为$\lambda_1=1,\lambda_2=2,\lambda_3=3$，它们对应的特征向量为$\xi_1=(1,1,1)^\mathrm{T}, \xi_2=(1,2,4)^\mathrm{T},\xi_3=(1,3,9)^\mathrm{T}$，又$\beta=(1,1,3)^\mathrm{T}$，计算$A^n\beta$.
\end{exercise}
\noindent 来源：\href{https://github.com/yhwu-is/Linear-Algebra-Left-Undone/blob/new/%E8%AE%B2%E4%B9%89/%E4%B8%93%E9%A2%98/15%20%E7%9B%B8%E4%BC%BC%E6%A0%87%E5%87%86%E5%BD%A2%EF%BC%9A%E5%A4%8D%E6%95%B0%E5%9F%9F%E4%B8%8A%E7%9A%84%E5%B0%9D%E8%AF%95%E4%B8%8E%E7%90%86%E8%AE%BA.tex}{L15-B-12}。
\begin{solve}
回顾对角化过程可知，令$P=(\xi_1,\xi_2,\xi_3)=\begin{pmatrix}
                1 & 1 & 1 \\ 1 & 2 & 3 \\ 1 & 4 & 9
            \end{pmatrix}$，则$P^{-1}AP=\begin{pmatrix}
                1 & 0 & 0 \\ 0 & 2 & 0 \\ 0 & 0 & 3
            \end{pmatrix}$，即$P^{-1}A^nP=\begin{pmatrix}
                1 & 0 & 0 \\ 0 & 2^n & 0 \\ 0 & 0 & 3^n
            \end{pmatrix}$，故$A^n=P\begin{pmatrix}
                1 & 0 & 0 \\ 0 & 2^n & 0 \\ 0 & 0 & 3^n
            \end{pmatrix}P^{-1}$，最后可以计算得到$A^n\beta=\begin{pmatrix}
                2-2^{n+1}+3^n \\ 2-2^{n+2}+3^{n+1} \\ 2-2^{n+3}+3^{n+2}
            \end{pmatrix}$.
\end{solve}

\begin{exercise}\label{prob:L15-B-13}
设$A=\begin{pmatrix}
                3 & 4 & 0 & 0 \\ 4 & -3 & 0 & 0 \\ 0 & 0 & 2 & 4 \\ 0 & 0 & 0 & 2
            \end{pmatrix}$，求$A^n(n\in\mathbf{N}_+)$.
\end{exercise}
\noindent 来源：\href{https://github.com/yhwu-is/Linear-Algebra-Left-Undone/blob/new/%E8%AE%B2%E4%B9%89/%E4%B8%93%E9%A2%98/15%20%E7%9B%B8%E4%BC%BC%E6%A0%87%E5%87%86%E5%BD%A2%EF%BC%9A%E5%A4%8D%E6%95%B0%E5%9F%9F%E4%B8%8A%E7%9A%84%E5%B0%9D%E8%AF%95%E4%B8%8E%E7%90%86%E8%AE%BA.tex}{L15-B-13}。
\begin{solve}
设$A=\begin{pmatrix}
                A_1 & O \\ O & A_2
            \end{pmatrix}$，其中$A_1=\begin{pmatrix}
                3 & 4 \\ 4 & -3
            \end{pmatrix}$，$A_2=\begin{pmatrix}
                2 & 4 \\ 0 & 2
            \end{pmatrix}$. 回顾矩阵运算进阶中介绍的求幂方法，我们发现$A_2=2E+B$，其中$B=\begin{pmatrix}
                0 & 4 \\ 0 & 0
            \end{pmatrix}$，且$B^2=O$，因此$A_2^n=(2E+B)^n=2^nE+n2^{n-1}B=\begin{pmatrix}
                2^n & n2^{n+1} \\ 0 & 2^n
            \end{pmatrix}$. 再看$A_1$，我们用对角化方法求幂，容易求得$P=\begin{pmatrix}
                2 & -1 \\ 1 & 2
            \end{pmatrix}$使得$P^{-1}A_1P=\begin{pmatrix}
                5 & 0 \\ 0 & -5
            \end{pmatrix}$，故
\[A_1^n=P\begin{pmatrix}
                5^n & 0 \\ 0 & (-5)^n
            \end{pmatrix}P^{-1}=5^{n-1}\begin{pmatrix}
                4+(-1)^n & 2+2(-1)^{n+1} \\ 2+2(-1)^{n+1} & 1+4(-1)^n
            \end{pmatrix}\]
，最后结合上面的讨论得到$A^n=\begin{pmatrix}
                A_1^n & O \\ O & A_2^n
            \end{pmatrix}$.
\end{solve}

\begin{exercise}\label{prob:L15-B-14}
已知三阶矩阵$A$和三元列向量$X$，使得向量组$X,AX,A^2X$线性无关，且满足
        \[A^3X=3AX-2A^2X.\]
        \begin{enumerate}
            \item 记$P=(X,AX,A^2X)$，求三阶矩阵$B$使得$A=PBP^{-1}$；

            \item 计算行列式$|A+E|$.
        \end{enumerate}
\end{exercise}
\noindent 来源：\href{https://github.com/yhwu-is/Linear-Algebra-Left-Undone/blob/new/%E8%AE%B2%E4%B9%89/%E4%B8%93%E9%A2%98/15%20%E7%9B%B8%E4%BC%BC%E6%A0%87%E5%87%86%E5%BD%A2%EF%BC%9A%E5%A4%8D%E6%95%B0%E5%9F%9F%E4%B8%8A%E7%9A%84%E5%B0%9D%E8%AF%95%E4%B8%8E%E7%90%86%E8%AE%BA.tex}{L15-B-14}。
\begin{solve}
\begin{enumerate}
              \item 由题意，$PB=AP$，利用分块矩阵乘法可知$PB=AP=(AX,A^2X,A^3X)=(AX,A^2X,3AX-2A^2X)$（注意$(AX,A^2X,A^3X)B=(AXB,A^2XB,A^3XB)$是错误的，因为这是$1\times 3$分块和$1\times 1$分块相乘，显然不能这么乘）.

                    则$PB=(AX,A^2X,3AX-2A^2X)=(X,AX,A^X)\begin{pmatrix}
                            0 & 0 & 0 \\ 1 & 0 & 3 \\ 0 & 1 & -2
                        \end{pmatrix}$. 等号两边左乘$P^{-1}$有$B=\begin{pmatrix}
                            0 & 0 & 0 \\ 1 & 0 & 3 \\ 0 & 1 & -2
                        \end{pmatrix}$.

              \item 由于$A$与$B$相似，故有相同的特征值，易求得$B$的特征值为$0,-3,1$，则$A$的特征值也为$0,-3,1$，则$A+E$的特征值根据正文例题可知为$1,-2,2$，则$A+E$的行列式为其特征值之积，即为$-4$.
          \end{enumerate}
\end{solve}

\begin{exercise}\label{prob:L15-C-2}
设$B=\alpha\alpha^\mathrm{T}$，其中$\alpha=(a_1,\ldots,a_n)^\mathrm{T}\neq 0\enspace(a_i\in\mathbf{R},\enspace i=1,2,\ldots,n)$.
        \begin{enumerate}
            \item 证明：$B^k=tB$，其中$k$为正整数，$t$为常数，并求$t$；

            \item 求可逆阵$P$使得$P^{-1}BP$为对角矩阵，并写出该对角矩阵.
        \end{enumerate}
\end{exercise}
\noindent 来源：\href{https://github.com/yhwu-is/Linear-Algebra-Left-Undone/blob/new/%E8%AE%B2%E4%B9%89/%E4%B8%93%E9%A2%98/15%20%E7%9B%B8%E4%BC%BC%E6%A0%87%E5%87%86%E5%BD%A2%EF%BC%9A%E5%A4%8D%E6%95%B0%E5%9F%9F%E4%B8%8A%E7%9A%84%E5%B0%9D%E8%AF%95%E4%B8%8E%E7%90%86%E8%AE%BA.tex}{L15-C-2}。
\begin{solve}
\begin{enumerate}
                \item $B^2=\alpha(\alpha^\mathrm{T}\alpha)\alpha^\mathrm{T}=m\alpha\alpha^\mathrm{T}=mB$，其中$m=\alpha^\mathrm{T}\alpha=a_1^2+\cdot+a_n^2$，且易知这就是$\operatorname{tr}(B)$. 用归纳法可知$B^k=m^{k-1}B$，令$t=m^{k-1}$则得证，其中$t=(\operatorname{tr}(B))^{k-1}$.

                \item 这是经典的秩1矩阵，由于$B$的秩为1，则$BX=0$的解空间维数为$n-1$，因此0是$B$的$n-1$重特征值. 而特征值之和等于矩阵的迹. 则剩下一个一重特征值为$\operatorname{tr}(B)$（由题意迹不为0）.

                      然后求特征向量，这里需要一些观察. 特征值0对应的特征向量即为$BX=0$的解，设解为$(x_1,\ldots,x_n)$，不失一般性地，设$a_1\neq 0$，由于$B$的秩为1，因此将$BX=0$展开为线性方程组后，每一行（除去全0的）代表的方程是完全一致的（因为它们成比例），因此我们只需考虑第一行方程
                      \[a_1^2x_1+a_1a_2x_2+\cdots+a_1a_nx_n=0,\]
                      解得$n-1$个线性无关特征向量为
                      \[X_1=(-a_2,a_1,0,\ldots,0),X_2=(-a_3,0,a_1,0,\ldots,0),\ldots,X_{n-1}=(-a_n,0,\ldots,0,a_1).\]
                      对于特征值$\operatorname{tr}(B)=\alpha^\mathrm{T}\alpha$，我们可以设特征向量为$X$，则$(\operatorname{tr}(B)E-B)X=(\alpha^\mathrm{T}\alpha E-\alpha\alpha^\mathrm{T})X=0$，即$\alpha^\mathrm{T}X=\alpha\alpha^\mathrm{T}X$，我们可以观察发现$\alpha$就是一个解，则$X_n=\alpha$是其特征向量. 则
\[P=\begin{pmatrix}
                              -a_2 & -a_3 & \cdots & -a_n & a_1 \\ a_1 & 0 & \cdots & 0 & a_2 \\ 0 & a_1 & \cdots & 0 & a_3 \\ \vdots & \vdots & \ddots & \vdots & \vdots \\ 0 & 0 & \cdots & a_1 & a_n
                          \end{pmatrix}\]
，且对角矩阵为$\operatorname{diag}(0,\ldots,0,\alpha^\mathrm{T}\alpha)$.
            \end{enumerate}
\end{solve}

\begin{exercise}\label{prob:L15-C-3}
（秩为1的矩阵）设$n$阶矩阵$A$的元素均为1.
        \begin{enumerate}
            \item 求$A$的特征值，并求矩阵$P$使得$P^{-1}AP$为对角矩阵；

            \item 若$f(x)$是$x$的$m$次多项式，且常数项为0，证明：存在$k\in\mathbf{R}$使得$f(A)=kA$，并求出$k$；

            \item 设$B$是$n$阶实对称矩阵，每行元素之和都为$b$，若$b$是$f(\lambda)=|\lambda E-B|$的单根，求$B$属于$b$的特征向量；当$f(\lambda)=(\lambda-b)g(\lambda)$时（其中$f(B)=0$），证明：$g(B)=kA$，其中$k$为常数，$A$为元素全部为1的$n$阶矩阵.
        \end{enumerate}
\end{exercise}
\noindent 来源：\href{https://github.com/yhwu-is/Linear-Algebra-Left-Undone/blob/new/%E8%AE%B2%E4%B9%89/%E4%B8%93%E9%A2%98/15%20%E7%9B%B8%E4%BC%BC%E6%A0%87%E5%87%86%E5%BD%A2%EF%BC%9A%E5%A4%8D%E6%95%B0%E5%9F%9F%E4%B8%8A%E7%9A%84%E5%B0%9D%E8%AF%95%E4%B8%8E%E7%90%86%E8%AE%BA.tex}{L15-C-3}。
\begin{solve}
\begin{enumerate}
                \item 由题意知$A$的秩为1，因此特征值0对应的特征子空间为$AX=0$的解空间，维数为$n-1$，故特征值0的重数至少为$n-1$. 剩下可能还有一个特征值，我们知道特征值之和等于对角线元素之和，因此最后一个特征值等于$n$，$n$对应的特征子空间维数至少为1，结合0对应特征子空间维数为$n-1$可知$n$对应的特征子空间维数就是1，二者可以张成整个$\mathbf{R}^n$，因此可以对角化.

                      下面求过渡矩阵，$n$对应的特征向量只需解方程$(nE-A)X=0$，即得$X_1=(1,1,\ldots,1)^\mathrm{T}$. 0对应的$n-1$个线性无关特征向量就是$AX=0$的基础解系，即$X_2=(-1,1,0,\ldots,0)^\mathrm{T},X_3=(-1,0,1,0,\ldots,0)^\mathrm{T},\ldots,X_n=(-1,0,\ldots,0,1)^\mathrm{T}$，令
\[P=(X_1,X_2,\ldots,X_n)=\begin{pmatrix}
                              1 & -1 & -1 & \cdots & -1 \\ 1 & 1 & 0 & \cdots & 0 \\ 1 & 0 & 1 & \cdots & 0 \\ \vdots & \vdots & \vdots & \ddots & \vdots \\ 1 & 0 & 0 & \cdots & 1
                          \end{pmatrix}\]
，则有$P^{-1}AP=\Lambda=\operatorname{diag}(n,0,\ldots,0)$.

                \item 由上一小问可知$A=P\Lambda P^{-1}$，则$A^k=P\Lambda^kP^{-1}$，故
                      \[A^k=P\Lambda^kP^{-1}=P\operatorname{diag}(n^k,0,\ldots,0)P^{-1},\]
                      故
                      \begin{align*}
                          f(A) & =a_mA^m+a_{m-1}A^{m-1}+\cdots+a_1A                                                                                                   \\
                               & =P(a_m\Lambda^m+a_{m-1}\Lambda^{m-1}+\cdots+a_1\Lambda)P^{-1}                                                               \\
                               & =P\begin{pmatrix}
                                       a_mn^m+a_{m-1}n^{m-1}+\cdots+a_1n & 0 & \cdots & 0 \\ 0 & 0 & \cdots & 0 \\ \vdots & \vdots & \ddots & \vdots \\ 0 & 0 & \cdots & 0
                                   \end{pmatrix}P^{-1} \\
                               & =\begin{pmatrix}
                                      f(n) & 0 & \cdots & 0 \\ 0 & 0 & \cdots & 0 \\ \vdots & \vdots & \ddots & \vdots \\ 0 & 0 & \cdots & 0
                                  \end{pmatrix}=\dfrac{f(n)}{n}P\begin{pmatrix}
                                                                    n & 0 & \cdots & 0 \\ 0 & 0 & \cdots & 0 \\ \vdots & \vdots & \ddots & \vdots \\ 0 & 0 & \cdots & 0
                                                                \end{pmatrix}P^{-1}    \\
                               & =\dfrac{f(n)}{n}P\Lambda P^{-1}=kA.
                      \end{align*}
                      其中$k=\dfrac{f(n)}{n}$. 故$f(A)=kA$.

                \item 显然$X=(1,1,x\cdots,1)^\mathrm{T}$是$B$关于特征值$b$的特征向量，由于$b$是特征多项式的单根，则特征子空间维数也为1，故特征向量就是$kX$（其中$k$为非零常数）.

                      又$f(B)=(B-bE)g(B)=O$，则$Bg(B)=bg(B)$. 令$g(B)$的列向量组为$\alpha_1,\ldots,\alpha_n$，则上式可以写为$B(\alpha_1,\ldots,\alpha_n)=(b\alpha_1,\ldots,b\alpha_n)$，即$B\alpha_i=b\alpha_i$，即$\alpha_i$是$B$关于特征值$b$的特征向量，故$\alpha_i=k_iX=(k_i,k_i,\ldots,k_i)^{\mathrm{T}}$，即
\[g(B)=\begin{pmatrix}
                              k_1 & k_2 & \cdots & k_n \\ k_1 & k_2 & \cdots & k_n \\ \vdots & \vdots & \ddots & \vdots \\ k_1 & k_2 & \cdots & k_n
                          \end{pmatrix}\]
. 由于$B$是实对称矩阵，故$g(B)$也是实对称矩阵（很容易验证实对称矩阵经过幂次、加法运算后仍是实对称矩阵），因此$g(B)=(g(B))^\mathrm{T}$，故$k_1=k_2=\cdots=k_n=k$，即$g(B)=kA$.
            \end{enumerate}
\end{solve}

\begin{exercise}\label{prob:L15-C-4}
设$A,B$为$n$阶矩阵，且$A+B=AB$，求证：
        \begin{enumerate}
            \item $A,B$的特征向量是公共的；

            \item $A$相似于对角矩阵当且仅当$B$相似于对角矩阵；

            \item $r(A)=r(B)$.
        \end{enumerate}
\end{exercise}
\noindent 来源：\href{https://github.com/yhwu-is/Linear-Algebra-Left-Undone/blob/new/%E8%AE%B2%E4%B9%89/%E4%B8%93%E9%A2%98/15%20%E7%9B%B8%E4%BC%BC%E6%A0%87%E5%87%86%E5%BD%A2%EF%BC%9A%E5%A4%8D%E6%95%B0%E5%9F%9F%E4%B8%8A%E7%9A%84%E5%B0%9D%E8%AF%95%E4%B8%8E%E7%90%86%E8%AE%BA.tex}{L15-C-4}。
\begin{solve}
\begin{enumerate}
                \item 设$\alpha$是$B$的特征向量，对应的特征值为$\lambda$，即$B\alpha=\lambda\alpha$，则$A\alpha+B\alpha-AB\alpha=0$，即$A\alpha+\lambda\alpha-A\lambda\alpha=0$，即$(1-\lambda)A\alpha=-\lambda\alpha$. 若$\lambda=1$，则$\alpha=0$，与它是特征向量矛盾，故$\lambda\neq 1$，从而$A\alpha=\dfrac{-\lambda}{1-\lambda}\alpha$，这就说明了$B$的特征向量就是$A$的特征向量.

                      由$A+B=AB$可知，$(A-E)(B-E)=E$，故$A-E$可逆，且$(A-E)^{-1}=B-E$，从而$E=(B-E)(A-E)=BA-A-B+E$，故$BA=A+B$，同上理可得$A$的特征向量也是$B$的，得证.

                \item \begin{enumerate}
                          \item 必要性：若$A$可对角化，则$\mathbf{R}^n$有一组由$A$的特征向量张成的基，由（1）知$B$与$A$特征向量一致，故$\mathbf{R}^n$也有一组由$B$的特征向量张成的基，即$B$可对角化.

                          \item 充分性与必要性类似.
                      \end{enumerate}

                \item 由$A+B=AB$可知$A=(A-E)B$且$B=A(B-E)$，由此可得$r(A)\leqslant r(B)\leqslant r(A)$，故$r(A)=r(B)$.
            \end{enumerate}
\end{solve}

\begin{exercise}\label{prob:L15-C-5}
设$A,B\in \mathbf{M}_n(\mathbf{R})$，证明$A$与$B$在$\mathbf{R}$上相似当且仅当在$\mathbf{C}$上相似.

        {\kaishu 注：实际上相似这一性质与数域无关，本题是这一结论的特例.}
\end{exercise}
\noindent 来源：\href{https://github.com/yhwu-is/Linear-Algebra-Left-Undone/blob/new/%E8%AE%B2%E4%B9%89/%E4%B8%93%E9%A2%98/15%20%E7%9B%B8%E4%BC%BC%E6%A0%87%E5%87%86%E5%BD%A2%EF%BC%9A%E5%A4%8D%E6%95%B0%E5%9F%9F%E4%B8%8A%E7%9A%84%E5%B0%9D%E8%AF%95%E4%B8%8E%E7%90%86%E8%AE%BA.tex}{L15-C-5}。
\begin{solve}
\begin{enumerate}
                \item 必要性：$A$和$B$在实数域上相似，则存在可逆实矩阵$C$使得$B=C^{-1}AC$，将$A,B,C$视为复数域上的矩阵，则可知$A$和$B$在复数域上也相似.

                \item 充分性：若$A$和$B$在复数域上相似，则存在可逆复矩阵$P=P_1+\mathrm{i} P_2$使得$B=P^{-1}AP$，其中$P_1,P_2$是实矩阵，即$AP=PB$，从而$A(P_1+\mathrm{i} P_2)=(P_1+\mathrm{i} P_2)B$，比较实部虚部可知$AP_1=P_1B$且$AP_2=P_2B$. 构造实系数多项式$f(t)=|P_1+tP_2|$，显然$f(\mathrm{i})=|P|\neq 0$，则$f(t)$是非零多项式，所以$f(t)$的根只有有限多个，从而存在$t_0\in\mathbf{R}$使得$f(t_0)\neq 0$，则$|P_1+t_0P_2|\neq 0$，从而$P_1+t_0P_2$可逆，从而$A(P_1+t_0P_2)=AP_1+t_0AP_2=P_1B+t_0P_2B=(P_1+t_0P_2)B$，即$(P_1+t_0P_2)^{-1}A(P_1+t_0P_2)=B$，即$A$和$B$在实数域上相似.
            \end{enumerate}
\end{solve}

\begin{exercise}\label{prob:L15-C-6}
设 $A,B\in\mathbf F^{n\times n}$。
\begin{enumerate}
\item 若 $A$ 在 $\mathbf F$ 中有 $n$ 个不同特征值，证明：$AB=BA$ 当且仅当 $A$ 的每个特征向量也是 $B$ 的特征向量。
\item 若 $A,B$ 都在 $\mathbf F$ 上可对角化且 $AB=BA$，证明它们可以用同一个过渡矩阵对角化。
\end{enumerate}
\end{exercise}
\noindent 来源：\href{https://github.com/yhwu-is/Linear-Algebra-Left-Undone/blob/new/%E8%AE%B2%E4%B9%89/%E4%B8%93%E9%A2%98/15%20%E7%9B%B8%E4%BC%BC%E6%A0%87%E5%87%86%E5%BD%A2%EF%BC%9A%E5%A4%8D%E6%95%B0%E5%9F%9F%E4%B8%8A%E7%9A%84%E5%B0%9D%E8%AF%95%E4%B8%8E%E7%90%86%E8%AE%BA.tex}{L15-C-6}。
\begin{solve}
\begin{enumerate}
\item 若 $AB=BA$，对 $Ax=\lambda x$ 有 $A(Bx)=\lambda Bx$，故 $B$ 保持 $A$ 的每个特征子空间。由于各特征子空间为一维，$Bx=\mu x$，所以每个非零特征向量 $x$ 也是 $B$ 的特征向量，其中允许 $\mu=0$。反过来，取 $A$ 的特征向量基，它也由 $B$ 的特征向量组成，两矩阵在此基下都是对角矩阵，因而交换。
\item 记 $A$ 的互异特征值为 $\lambda_1,\ldots,\lambda_s$，则 $\mathbf F^n=\bigoplus_i E_{\lambda_i}(A)$，每个直和分量均为 $B$ 的不变子空间。因 $B$ 可对角化，其最小多项式分裂为互异一次因子的乘积；每个限制 $B|_{E_{\lambda_i}(A)}$ 的最小多项式整除它，故该限制也可对角化。在各分量内选择 $B$ 的特征向量基，再将这些基合并，就得到 $A,B$ 的共同特征向量基。
\end{enumerate}
\end{solve}

\begin{exercise}\label{prob:L15-C-7}
设 $A,B\in\mathbf F^{n\times n}$，$A$ 在 $\mathbf F$ 中有 $n$ 个不同的特征值，且 $AB=BA$。证明：存在唯一的次数不超过 $n-1$ 的多项式 $f\in\mathbf F[x]$，使 $B=f(A)$。
\end{exercise}
\noindent 来源：\href{https://github.com/yhwu-is/Linear-Algebra-Left-Undone/blob/new/%E8%AE%B2%E4%B9%89/%E4%B8%93%E9%A2%98/15%20%E7%9B%B8%E4%BC%BC%E6%A0%87%E5%87%86%E5%BD%A2%EF%BC%9A%E5%A4%8D%E6%95%B0%E5%9F%9F%E4%B8%8A%E7%9A%84%E5%B0%9D%E8%AF%95%E4%B8%8E%E7%90%86%E8%AE%BA.tex}{L15-C-7}。
\begin{solve}
由习题~\ref{prob:L15-C-6}，存在可逆矩阵 $P$ 使 $P^{-1}AP=\operatorname{diag}(\lambda_1,\ldots,\lambda_n)$、$P^{-1}BP=\operatorname{diag}(\mu_1,\ldots,\mu_n)$。取 Lagrange 插值多项式
\[f(x)=\sum_{i=1}^n\mu_i\prod_{j\ne i}\frac{x-\lambda_j}{\lambda_i-\lambda_j},\]
则次数不超过 $n-1$，且 $f(\lambda_i)=\mu_i$，故 $P^{-1}f(A)P=\operatorname{diag}(\mu_1,\ldots,\mu_n)$，即 $f(A)=B$。若还有 $g$ 满足同样的次数界和矩阵等式，则 $f-g$ 在 $n$ 个不同的 $\lambda_i$ 上为零，而次数小于 $n$，所以 $f=g$。
\end{solve}

\begin{exercise}\label{prob:L15-C-8}
设$T\in \mathcal{L}(V)$，$\lambda\in\mathbf{F}$. 证明：对$V$的每个使得$T$有上三角矩阵的基，$\lambda$出现在$T$的矩阵的对角线上的次数等于$\lambda$作为$T$的特征值的重数.
\end{exercise}
\noindent 来源：\href{https://github.com/yhwu-is/Linear-Algebra-Left-Undone/blob/new/%E8%AE%B2%E4%B9%89/%E4%B8%93%E9%A2%98/15%20%E7%9B%B8%E4%BC%BC%E6%A0%87%E5%87%86%E5%BD%A2%EF%BC%9A%E5%A4%8D%E6%95%B0%E5%9F%9F%E4%B8%8A%E7%9A%84%E5%B0%9D%E8%AF%95%E4%B8%8E%E7%90%86%E8%AE%BA.tex}{L15-C-8}。
\begin{solve}
设 $T$ 在题设基下的上三角矩阵为 $A=(a_{ij})$。用独立变量 $t$ 表示特征多项式，则
\[\chi_T(t)=\det(tI-A)=\prod_{j=1}^n(t-a_{jj}).\]
因此固定的 $\lambda$ 在对角线上出现几次，因子 $t-\lambda$ 就在该乘积中出现几次，这恰是 $\lambda$ 的代数重数。这里不能把代数重数换成特征子空间的维数（几何重数）；例如 $\begin{pmatrix}\lambda&1\\0&\lambda\end{pmatrix}$ 的代数重数为 $2$，几何重数只有 $1$。
\end{solve}

\begin{exercise}\label{prob:L15-C-9}
证明：$A$为幂零矩阵$\iff \forall k \in \mathbf{N}_+,\enspace\operatorname{tr}(A^k)=0$.
\end{exercise}
\noindent 来源：\href{https://github.com/yhwu-is/Linear-Algebra-Left-Undone/blob/new/%E8%AE%B2%E4%B9%89/%E4%B8%93%E9%A2%98/15%20%E7%9B%B8%E4%BC%BC%E6%A0%87%E5%87%86%E5%BD%A2%EF%BC%9A%E5%A4%8D%E6%95%B0%E5%9F%9F%E4%B8%8A%E7%9A%84%E5%B0%9D%E8%AF%95%E4%B8%8E%E7%90%86%E8%AE%BA.tex}{L15-C-9}。
\begin{solve}
在复数域上考察 $A$ 的特征值，不改变矩阵幂或迹。若 $A$ 幂零，则其所有特征值为零，因而每个 $A^k$ 的迹为零。
反过来，设所有正整数 $k$ 都满足 $\operatorname{tr}(A^k)=0$。若存在非零特征值，设互异非零特征值为 $\lambda_1,\ldots,\lambda_r$，代数重数为正整数 $m_1,\ldots,m_r$。由复矩阵可上三角化，有
\[\sum_{j=1}^r m_j\lambda_j^k=0,\qquad k=1,\ldots,r.\]
这一关于 $m_j$ 的线性方程组的系数行列式为
\[\left(\prod_{j=1}^r\lambda_j\right)\prod_{1\le i<j\le r}(\lambda_j-\lambda_i)\ne0,\]
故所有 $m_j=0$，矛盾。因此全部特征值为零，特征多项式为 $t^n$，由 Cayley--Hamilton 定理得 $A^n=O$，即 $A$ 幂零。
\end{solve}

\begin{exercise}\label{prob:L15-例-1}
证明$r$阶上三角矩阵$(r>1)$
    \[J_0=\begin{pmatrix}
            \lambda_0 & 1         &        &           \\
                      & \lambda_0 & \ddots &           \\
                      &           & \ddots & 1         \\
                      &           &        & \lambda_0
        \end{pmatrix}\]
    不与对角阵相似.
\end{exercise}
\noindent 来源：\href{https://github.com/yhwu-is/Linear-Algebra-Left-Undone/blob/new/%E8%AE%B2%E4%B9%89/%E4%B8%93%E9%A2%98/15%20%E7%9B%B8%E4%BC%BC%E6%A0%87%E5%87%86%E5%BD%A2%EF%BC%9A%E5%A4%8D%E6%95%B0%E5%9F%9F%E4%B8%8A%E7%9A%84%E5%B0%9D%E8%AF%95%E4%B8%8E%E7%90%86%E8%AE%BA.tex}{L15-例-1}。
\begin{solve}
首先求出特征多项式为 $f(\lambda) = |\lambda E - J_0| = (\lambda - \lambda_0)^r$，因此 $J_0$ 只有一个特征值 $\lambda_0$，且代数重数为 $r$.

    接下来求几何重数，即 $J_0X = \lambda_0X$ 的解空间维数，即 $(\lambda_0 E - J_0)X = 0$的解空间维数，事实上由于 $r(\lambda_0 E - J_0) = r - 1$，因此解空间维数为 $r - (r - 1) = 1$，即几何重数为 $1 < r$，因此不可对角化.
\end{solve}

\begin{exercise}\label{prob:L15-例-2}
求矩阵
    \[A=\begin{pmatrix}
            0  & -1 & 1 \\
            -1 & 0  & 1 \\
            1  & 1  & 0
        \end{pmatrix}\]
    的所有特征值，对应的特征子空间，以及与 $A$ 相似的一个对角矩阵.
\end{exercise}
\noindent 来源：\href{https://github.com/yhwu-is/Linear-Algebra-Left-Undone/blob/new/%E8%AE%B2%E4%B9%89/%E4%B8%93%E9%A2%98/15%20%E7%9B%B8%E4%BC%BC%E6%A0%87%E5%87%86%E5%BD%A2%EF%BC%9A%E5%A4%8D%E6%95%B0%E5%9F%9F%E4%B8%8A%E7%9A%84%E5%B0%9D%E8%AF%95%E4%B8%8E%E7%90%86%E8%AE%BA.tex}{L15-例-2}。
\begin{solve}
对于求解矩阵的对角化问题，首先求出其特征多项式（具体步骤不展开，实际上就是三阶行列式的计算，可以使用按行（列）展开、公式法或者初等变换化为三角矩阵等方法）$f(\lambda)=|\lambda E-A|=(\lambda-1)^2(\lambda+2)$，令$f(\lambda)=0$，解得特征值为 $\lambda_1=\lambda_2=1,\lambda_3=-2$.

    接下来求解特征向量和特征子空间，即求解$(E-A)x=0$和解$(-2E-A)x=0$，得到特征值1对应的特征子空间为$\operatorname{span}((-1,1,0)^{\mathrm{T}},(1,0,1)^{\mathrm{T}})$，特征值-2对应的特征子空间为$\operatorname{span}((-1,-1,1)^{\mathrm{T}})$.

    与$A$相似的对角矩阵实际上就是特征值排列在对角线上的结果，即 $\operatorname{diag}(1,1,-2)$.
\end{solve}

\begin{exercise}\label{prob:L15-例-4}
已知$A=\begin{pmatrix}
            0 & \dfrac{1}{2}  & \dfrac{1}{2} \\[2ex]
            1 & -\dfrac{1}{2} & \dfrac{1}{2} \\[2ex]
            1 & -\dfrac{1}{2} & \dfrac{1}{2}
        \end{pmatrix}$，求$A^n$.
\end{exercise}
\noindent 来源：\href{https://github.com/yhwu-is/Linear-Algebra-Left-Undone/blob/new/%E8%AE%B2%E4%B9%89/%E4%B8%93%E9%A2%98/15%20%E7%9B%B8%E4%BC%BC%E6%A0%87%E5%87%86%E5%BD%A2%EF%BC%9A%E5%A4%8D%E6%95%B0%E5%9F%9F%E4%B8%8A%E7%9A%84%E5%B0%9D%E8%AF%95%E4%B8%8E%E7%90%86%E8%AE%BA.tex}{L15-例-4}。
\begin{solve}
首先求出$A$的特征多项式$f(\lambda)=|\lambda E-A|=\lambda(\lambda-1)(\lambda+1)$，令$f(\lambda)=0$，解得特征值为 $\lambda_1=0,\lambda_2=1,\lambda_3=-1$.

    接下来求解特征向量和特征子空间，实际上就是求解$(0E-A)x=0,(-E-A)x=0,(E-A)x=0$，得到特征向量为
    \[\eta_1=\begin{pmatrix}
            1 \\ 1 \\ -1
        \end{pmatrix},\enspace \eta_2=\begin{pmatrix}
            1 \\ 1 \\ 1
        \end{pmatrix},\enspace \eta_3=\begin{pmatrix}
            1 \\ -1 \\ -1
        \end{pmatrix}.\]
    所以记$P=(\eta_1,\eta_2,\eta_3)$，则$A=P\operatorname{diag}(0,1,-1)P^{-1}$，因此
    \[A^n=P\operatorname{diag}(0^n,1^n,(-1)^n)P^{-1},\]
    进一步计算得到
    \[A^n=\frac{1}{2}\begin{pmatrix}
            1+(-1)^n     & (-1)^{n+1} & 1 \\
            1+(-1)^{n+1} & (-1)^n     & 1 \\
            1+(-1)^{n+1} & (-1)^n     & 1
        \end{pmatrix}.\]
\end{solve}

\begin{exercise}\label{prob:L15-例-5}
线性变换 $T : \mathbf{R}^3 \to \mathbf{R}^3$ 的定义是：
    \[T(x_1,x_2,x_3)=(4x_1+x_3,2x_1+3x_2+2x_3,x_1+4x_3)\]
    \begin{enumerate}
        \item 求出$T$的特征多项式及特征值；

        \item 判断$T$是否可对角化，并给出理由.
    \end{enumerate}
\end{exercise}
\noindent 来源：\href{https://github.com/yhwu-is/Linear-Algebra-Left-Undone/blob/new/%E8%AE%B2%E4%B9%89/%E4%B8%93%E9%A2%98/15%20%E7%9B%B8%E4%BC%BC%E6%A0%87%E5%87%86%E5%BD%A2%EF%BC%9A%E5%A4%8D%E6%95%B0%E5%9F%9F%E4%B8%8A%E7%9A%84%E5%B0%9D%E8%AF%95%E4%B8%8E%E7%90%86%E8%AE%BA.tex}{L15-例-5}。
\begin{solve}
\begin{enumerate}
        \item 我们知道，线性变换与其在任意一组基下的矩阵有相同的特征值和特征多项式，因此我们可以任意选取一组基（为方便我们取$\mathbf{R}^3$的自然基，然后求出$T$在这组基下的矩阵，即
              \[A=\begin{pmatrix}
                      4 & 0 & 1 \\
                      2 & 3 & 2 \\
                      1 & 0 & 4
                  \end{pmatrix}\]
              故$T$的特征多项式为$f(\lambda)=|A-\lambda E|=\begin{vmatrix}
                      4-\lambda & 0         & 1         \\
                      2         & 3-\lambda & 2         \\
                      1         & 0         & 4-\lambda
                  \end{vmatrix}=-(\lambda-3)^2(\lambda-5)$，$T$的特征值为$f(\lambda)=0$的解，即为$\lambda_1=\lambda_2=3,\lambda_3=5$.

        \item 对于$\lambda=3$，我们求解$(T-3I)\alpha$得到特征向量（其中令$\alpha=(x_1,x_2,x_3)$），解得$\alpha_1=(1,0,-1)^\mathrm{T},\alpha_2=(0,1,0)^\mathrm{T}$，因此$T$在$\lambda=3$时的特征子空间为$V_3=\operatorname{span}(\alpha_1,\alpha_2)$，维数为2，因此$\lambda=3$时代数重数与几何重数相等. 而$\lambda=5$是一重特征值，因此也必有代数重数与几何重数相等（因为一定有特征向量，那么特征子空间维数至少为1，再结合代数重数大于等于几何重数可知，特征子空间维数恰为1）. 因此$T$可对角化.
    \end{enumerate}
\end{solve}

\begin{exercise}\label{prob:L15-例-6}
设$\alpha$和$\beta$是$\mathbf{R}^n\enspace (n>1)$中两个列向量，$A=\alpha\beta^\mathrm{T}\neq O$.
    \begin{enumerate}
        \item 求$A$的特征值；

        \item 证明：$\alpha^\mathrm{T}\beta=0\iff A$不可对角化.
    \end{enumerate}
\end{exercise}
\noindent 来源：\href{https://github.com/yhwu-is/Linear-Algebra-Left-Undone/blob/new/%E8%AE%B2%E4%B9%89/%E4%B8%93%E9%A2%98/15%20%E7%9B%B8%E4%BC%BC%E6%A0%87%E5%87%86%E5%BD%A2%EF%BC%9A%E5%A4%8D%E6%95%B0%E5%9F%9F%E4%B8%8A%E7%9A%84%E5%B0%9D%E8%AF%95%E4%B8%8E%E7%90%86%E8%AE%BA.tex}{L15-例-6}。
\begin{solve}
\begin{enumerate}
        \item 我们知道，$r(A)\leqslant\min{\{r(\alpha),r(\beta)\}}=1$，并且$A\neq O$因此$r(A)>0$，故$A$的秩为1. 而$n>1$，因此$A$一定不可逆，故0一定是$A$的特征值，且对应的特征子空间维数为$AX=0$的解空间维数，即为$n-1$.

              由此我们知道$A$最多有两个特征值，因为0的代数重数（即作为$n$次特征多项式的零点次数）必然大于等于其几何重数$n-1$，当其代数重数为$n-1$时可能还有一个一重特征值. 我们利用特征值之和等于$A$的迹来找出可能的第二个特征值. 我们设$\alpha=(a_1,a_2,\ldots,a_n)^\mathrm{T},\beta=(b_1,b_2,\ldots,b_n)^\mathrm{T}$，则
\[A=\alpha\beta^\mathrm{T}=\begin{pmatrix}
                      a_1b_1 & a_1b_2 & \cdots & a_1b_n \\
                      a_2b_1 & a_2b_2 & \cdots & a_2b_n \\
                      \vdots & \vdots & \ddots & \vdots \\
                      a_nb_1 & a_nb_2 & \cdots & a_nb_n
                  \end{pmatrix}\]
，因此$A$的迹为$\sum\limits_{i=1}^na_ib_i=\alpha^\mathrm{T}\beta=\sum\limits_{i=1}^n\lambda_i$，其中$\lambda_i$为$A$的特征值. 若$\alpha^\mathrm{T}\beta\neq 0$，则$\lambda_i=0,\enspace i=1,\ldots,n-1$，$\lambda_n=\alpha^\mathrm{T}\beta$. 若$\alpha^\mathrm{T}\beta=0$，则$A$的所有特征值均为0.

        \item 由上一小问可知，若$\alpha^\mathrm{T}\beta=0$即$A$的全部特征值为0，因此只有一个$n-1$维的特征子空间，故特征子空间直和不等于$V$，故不可对角化.

              反之，若$A$不可对角化，我们用反证法. 假设$\alpha^\mathrm{T}\beta\neq 0$，则$A$有两个特征值，一个为0，一个为$\alpha^\mathrm{T}\beta$，因此$A$有两个特征子空间，一个是0对应的$n-1$维特征子空间，一个是$\alpha^\mathrm{T}\beta$对应的一维特征子空间，因此$V$可分解为两个特征子空间的直和，与$A$不可对角化矛盾，因此$\alpha^\mathrm{T}\beta=0$.
    \end{enumerate}
\end{solve}

\begin{exercise}\label{prob:L15-例-8}
设三阶矩阵$A$的特征值为$\lambda_1=-2,\lambda_2=1,\lambda_3=2$，对应的特征向量分别为$\alpha_1=(1,1,0)^\mathrm{T},\alpha_2=(1,0,1)^\mathrm{T},\alpha_3=(1,1,1)^\mathrm{T}$，求矩阵$A$.
\end{exercise}
\noindent 来源：\href{https://github.com/yhwu-is/Linear-Algebra-Left-Undone/blob/new/%E8%AE%B2%E4%B9%89/%E4%B8%93%E9%A2%98/15%20%E7%9B%B8%E4%BC%BC%E6%A0%87%E5%87%86%E5%BD%A2%EF%BC%9A%E5%A4%8D%E6%95%B0%E5%9F%9F%E4%B8%8A%E7%9A%84%E5%B0%9D%E8%AF%95%E4%B8%8E%E7%90%86%E8%AE%BA.tex}{L15-例-8}。
\begin{solve}
根据特征值分解可知，$A=P\Lambda P^{-1}$，其中$P=(\alpha_1,\alpha_2,\alpha_3),\Lambda=\begin{pmatrix}
            -2 & 0 & 0 \\
            0  & 1 & 0 \\
            0  & 0 & 2
        \end{pmatrix}$，因此
    \[A=P\Lambda P^{-1}=\begin{pmatrix}
            1 & 1 & 1 \\
            1 & 0 & 1 \\
            0 & 1 & 1
        \end{pmatrix}\begin{pmatrix}
            -2 & 0 & 0 \\
            0  & 1 & 0 \\
            0  & 0 & 2
        \end{pmatrix}\begin{pmatrix}
            1  & 0  & -1 \\
            1  & -1 & 0  \\
            -1 & 1  & 1
        \end{pmatrix}=\begin{pmatrix}
            -3 & 1 & 4 \\
            -4 & 2 & 4 \\
            -1 & 1 & 2
        \end{pmatrix}.\]
\end{solve}

\begin{exercise}\label{prob:L15-例-10}
设$\sigma\in \mathcal{L}(\mathbf{C}^3)$定义为
    \[\sigma(z_1,z_2,z_3)=(6z_1+3z_2+4z_3,6z_2+2z_3,7z_3),\]求一组基使其有分块对角矩阵并写出对应的分块对角矩阵.
\end{exercise}
\noindent 来源：\href{https://github.com/yhwu-is/Linear-Algebra-Left-Undone/blob/new/%E8%AE%B2%E4%B9%89/%E4%B8%93%E9%A2%98/15%20%E7%9B%B8%E4%BC%BC%E6%A0%87%E5%87%86%E5%BD%A2%EF%BC%9A%E5%A4%8D%E6%95%B0%E5%9F%9F%E4%B8%8A%E7%9A%84%E5%B0%9D%E8%AF%95%E4%B8%8E%E7%90%86%E8%AE%BA.tex}{L15-例-10}。
\begin{solve}
不难写出 $\sigma$ 在$\mathcal{L}(\mathbf{C}^3)$的自然基下的矩阵是
    \[ \begin{pmatrix}
        6 & 3 & 4 \\
        0 & 6 & 2 \\
        0 & 0 & 7
    \end{pmatrix}, \]
    则 $\sigma$的特征值为 $\lambda_1 = 6,\lambda_2 = 7$. 解方程可以得到两个广义特征子空间 \[G(\lambda_1,\sigma) = \operatorname{span}\left((1,0,0),\left(0,\dfrac{1}{3},0\right)\right)\]，$G(\lambda_2,\sigma)=\operatorname{span}((10,2,1))$.

    取$\alpha_1=(1,0,0),\alpha_2=\left(0,\dfrac{1}{3},0\right),\alpha_3=(10,2,1)$，则$\sigma$在基$\{\alpha_1,\alpha_2,\alpha_3\}$下有分块对角矩阵
    \[ \begin{pmatrix}
        6 & 1 & 0 \\
        0 & 6 & 0 \\
        0 & 0 & 7
    \end{pmatrix}. \]
\end{solve}

\begin{exercise}\label{prob:L15-例-11}
设$\sigma\in \mathcal{L}(\mathbf{C}^4)$定义为
    \[\sigma(z_1,z_2,z_3,z_4)=(-z_1+z_2,-z_1+z_2+z_3,z_1-z_2+2z_3+z_4,-z_1+z_2-z_3+3z_4),\]
    求$\mathbf{C}^4$的一组基，使得$\sigma$在这组基下的矩阵为上三角矩阵.
\end{exercise}
\noindent 来源：\href{https://github.com/yhwu-is/Linear-Algebra-Left-Undone/blob/new/%E8%AE%B2%E4%B9%89/%E4%B8%93%E9%A2%98/15%20%E7%9B%B8%E4%BC%BC%E6%A0%87%E5%87%86%E5%BD%A2%EF%BC%9A%E5%A4%8D%E6%95%B0%E5%9F%9F%E4%B8%8A%E7%9A%84%E5%B0%9D%E8%AF%95%E4%B8%8E%E7%90%86%E8%AE%BA.tex}{L15-例-11}。
\begin{solve}
记$\mathbf{C}^4$的自然基为$e_1,e_2,e_3,e_4$，则$\sigma$在自然基下的矩阵为
    \[
        A=\begin{pmatrix}
            -1 & 1  & 0  & 0 \\
            -1 & 1  & 1  & 0 \\
            1  & -1 & 2  & 1 \\
            -1 & 1  & -1 & 3
        \end{pmatrix}.
    \]
    下面按照原讲义“上三角矩阵存在”处 证明的法二逐步构造相似变换矩阵. 注意到 $A(1,1,0,0)^\mathrm{T} = 0$，因此 $0$ 是 $A$ 的一个特征值，对应的特征向量可取 $\alpha_1=(1,1,0,0)^\mathrm{T}$. 令 $P_1 = (\alpha_1,e_2,e_3,e_4)$，则
    \[ P_1^{-1}AP_1 = \begin{pmatrix}
            0 & 1  & 0  & 0 \\
            0 & 0  & 1  & 0 \\
            0 & -1 & 2  & 1 \\
            0 & 1  & -1 & 3
        \end{pmatrix}. \]
    其右下角三阶矩阵记为 $B_1$. 由于$B_1(1,1,0)^\mathrm{T} = (1,1,0)^\mathrm{T}$，故 $1$ 是 $B_1$ 的一个特征值. 取 $Q_2 = \begin{pmatrix}
            1 & 0 & 0 \\
            1 & 1 & 0 \\
            0 & 0 & 1
        \end{pmatrix}$，则
    \[
        Q_2^{-1}B_1Q_2
        =
        \begin{pmatrix}
            1 & 1  & 0 \\
            0 & 1  & 1 \\
            0 & -1 & 3
        \end{pmatrix}.
    \]
    再看 $Q_2^{-1}B_1Q_2$ 的右下角二阶矩阵，记为 $B_2$，有 $B_2(1,1)^\mathrm{T} = 2(1,1)^\mathrm{T}$，故 $2$ 是 $B_2$ 的一个特征值. 取 $Q_3= \begin{pmatrix}
            1 & 0 \\
            1 & 1
        \end{pmatrix}$，则
    \[
        Q_3^{-1}B_2Q_3
        =
        \begin{pmatrix}
            2 & 1 \\
            0 & 2
        \end{pmatrix}.
    \]
    至此，可以逐步构造出相似变换矩阵. 首先得到
    \[
        P_2 = Q_2
        \begin{pmatrix}
            1 & O \\ O & Q_3
        \end{pmatrix}
        =
        \begin{pmatrix}
            1 & 0 & 0 \\
            1 & 1 & 0 \\
            0 & 1 & 1
        \end{pmatrix}.
    \]
    则 $P_2$ 使得 $P_2^{-1}B_1P_2$ 为上三角矩阵. 进而得到
    \[
        P = P_1 \begin{pmatrix}
            1 & O \\ O & P_2
        \end{pmatrix} = \begin{pmatrix}
            1 & 0 & 0 & 0 \\
            1 & 1 & 0 & 0 \\
            0 & 1 & 1 & 0 \\
            0 & 0 & 1 & 1
        \end{pmatrix}. \]
    也就是说，取 $\alpha_1=(1,1,0,0)^\mathrm{T}, \alpha_2=(0,1,1,0)^\mathrm{T}, \alpha_3=(0,0,1,1)^\mathrm{T}, \alpha_4=(0,0,0,1)^\mathrm{T}$，则$\sigma$在基$\{\alpha_1,\alpha_2,\alpha_3,\alpha_4\}$下的矩阵为上三角矩阵
    \[
        P^{-1}AP = \begin{pmatrix}
            0 & 1 & 0 & 0 \\
            0 & 1 & 1 & 0 \\
            0 & 0 & 2 & 1 \\
            0 & 0 & 0 & 2
        \end{pmatrix}. \]
\end{solve}

\begin{exercise}\label{prob:L15-例-15}
定义$N\in \mathcal{L}(\mathbf{F}^5)$为
    \[N(x_1,x_2,x_3,x_4,x_5)=(2x_2,3x_3,-x_4,4x_5,0)\]
    求$(I+N)$的一个平方根.
\end{exercise}
\noindent 来源：\href{https://github.com/yhwu-is/Linear-Algebra-Left-Undone/blob/new/%E8%AE%B2%E4%B9%89/%E4%B8%93%E9%A2%98/15%20%E7%9B%B8%E4%BC%BC%E6%A0%87%E5%87%86%E5%BD%A2%EF%BC%9A%E5%A4%8D%E6%95%B0%E5%9F%9F%E4%B8%8A%E7%9A%84%E5%B0%9D%E8%AF%95%E4%B8%8E%E7%90%86%E8%AE%BA.tex}{L15-例-15}。
\begin{solve}
由于$N^4=0$，我们设$\sqrt{I+N}=I+a_1N+a_2N^2+a_3N^3$，则
    \begin{align*}
        I+N&=(I+a_1N+a_2N^2+a_3N^3)^2 \\
        &=I+2a_1N+(a_1^2+2a_2)N^2+(2a_3+2a_1a_2)N^3,
    \end{align*}
    可解得$a_1=\dfrac{1}{2},a_2=-\dfrac{1}{8},a_3=\dfrac{1}{16}$. 代入即可算出$(I+N)$的一个平方根.
\end{solve}

\begin{exercise}\label{prob:L15-例-18}
找出有限维实向量空间的一个线性变换$\sigma$，使得0是$\sigma$仅有的特征值但$\sigma$不是幂零线性变换.
\end{exercise}
\noindent 来源：\href{https://github.com/yhwu-is/Linear-Algebra-Left-Undone/blob/new/%E8%AE%B2%E4%B9%89/%E4%B8%93%E9%A2%98/15%20%E7%9B%B8%E4%BC%BC%E6%A0%87%E5%87%86%E5%BD%A2%EF%BC%9A%E5%A4%8D%E6%95%B0%E5%9F%9F%E4%B8%8A%E7%9A%84%E5%B0%9D%E8%AF%95%E4%B8%8E%E7%90%86%E8%AE%BA.tex}{L15-例-18}。
\begin{solve}
我们定义$\sigma\in\mathcal{L}({\mathbf{R}^3})$，$\sigma$在自然基下的矩阵为
            \[\begin{pmatrix}
                0 & 1 & 0 \\
                -1 & 0 & 0 \\
                0 & 0 & 0
            \end{pmatrix},\]
            则$0$是$\sigma$仅有的特征值，但$\sigma^3=-\sigma\neq 0$，$\sigma$不为幂零线性变换.
\end{solve}

\begin{exercise}\label{prob:L15-例-19}
找出一个$\sigma\in \mathcal{L}(\mathbf{R}^2)$使得$\sigma^4=-I$.
\end{exercise}
\noindent 来源：\href{https://github.com/yhwu-is/Linear-Algebra-Left-Undone/blob/new/%E8%AE%B2%E4%B9%89/%E4%B8%93%E9%A2%98/15%20%E7%9B%B8%E4%BC%BC%E6%A0%87%E5%87%86%E5%BD%A2%EF%BC%9A%E5%A4%8D%E6%95%B0%E5%9F%9F%E4%B8%8A%E7%9A%84%E5%B0%9D%E8%AF%95%E4%B8%8E%E7%90%86%E8%AE%BA.tex}{L15-例-19}。
\begin{solve}
从旋转的角度考虑可以得到
            \[
            \sigma(x,y)=(x,y)\begin{pmatrix}
                \dfrac{\sqrt{2}}{2} & -\dfrac{\sqrt{2}}{2} \\[1ex]
                \dfrac{\sqrt{2}}{2} & \dfrac{\sqrt{2}}{2}
            \end{pmatrix}
            \]
            满足要求.
\end{solve}

\begin{exercise}\label{prob:L15-例-20}
证明或给出反例：$V$上的幂零线性变换的集合是$\mathcal{L}(V)$的子空间.
\end{exercise}
\noindent 来源：\href{https://github.com/yhwu-is/Linear-Algebra-Left-Undone/blob/new/%E8%AE%B2%E4%B9%89/%E4%B8%93%E9%A2%98/15%20%E7%9B%B8%E4%BC%BC%E6%A0%87%E5%87%86%E5%BD%A2%EF%BC%9A%E5%A4%8D%E6%95%B0%E5%9F%9F%E4%B8%8A%E7%9A%84%E5%B0%9D%E8%AF%95%E4%B8%8E%E7%90%86%E8%AE%BA.tex}{L15-例-20}。
\begin{solve}
反例：$V=\mathbf{R}^2,\sigma,\tau$在自然基下的矩阵分别为
            \[ \begin{pmatrix}
                0 & 1 \\
                0 & 0
            \end{pmatrix},
            \begin{pmatrix}
                0 & 0 \\
                1 & 0
            \end{pmatrix}\]
            则$\sigma+\tau$在自然基下的矩阵为
            \[ \begin{pmatrix}
                0 & 1 \\
                1 & 0
            \end{pmatrix} \]
            表明$\sigma+\tau$不是幂零线性变换. 从而命题不成立.
\end{solve}

\begin{exercise}\label{prob:S-EIG-1}
设 $A=\begin{pmatrix}-1&-6&0\\0&1&0\\-4&-12&3\end{pmatrix}$。求特征值与一组特征向量基，写出对角化分解，并计算 $A^4$ 与 $\det A$。
\end{exercise}
\noindent 编者补充：S-EIG-1；相关主题见原讲义对应讲，题目数据与解答由本选编另行构造。
\begin{solve}
计算特征多项式，或逐一核对下面的列向量，可得三个不同特征值 $-1,1,3$。对应特征向量按列排列为
\[P=\begin{pmatrix}1&-3&0\\0&1&0\\1&0&1\end{pmatrix},\qquad D=\operatorname{diag}(-1,1,3).\]
直接计算 $AP=PD$ 且 $\det P=1$，从而 $A=PDP^{-1}$。因此
\[A^4=PD^4P^{-1}=\begin{pmatrix}1&0&0\\0&1&0\\-80&-240&81\end{pmatrix},\qquad \det A=-3.\]
不同特征值对应的特征向量线性无关，所以这三个向量构成一组基。该方法把高次矩阵乘法转化为标量幂；一般地 $A^m=P\operatorname{diag}((-1)^m,1,3^m)P^{-1}$，对所有非负整数 $m$ 成立。
\end{solve}

\begin{exercise}\label{prob:S-EIG-2}
设 $A=\begin{pmatrix}-2&-3&-6\\0&1&6\\0&0&4\end{pmatrix}$。求特征值与一组特征向量基，写出对角化分解，并计算 $A^4$ 与 $\det A$。
\end{exercise}
\noindent 编者补充：S-EIG-2；相关主题见原讲义对应讲，题目数据与解答由本选编另行构造。
\begin{solve}
计算特征多项式，或逐一核对下面的列向量，可得三个不同特征值 $-2,1,4$。对应特征向量按列排列为
\[P=\begin{pmatrix}1&-1&-2\\0&1&2\\0&0&1\end{pmatrix},\qquad D=\operatorname{diag}(-2,1,4).\]
直接计算 $AP=PD$ 且 $\det P=1$，从而 $A=PDP^{-1}$。因此
\[A^4=PD^4P^{-1}=\begin{pmatrix}16&15&-510\\0&1&510\\0&0&256\end{pmatrix},\qquad \det A=-8.\]
不同特征值对应的特征向量线性无关，所以这三个向量构成一组基。该方法把高次矩阵乘法转化为标量幂；一般地 $A^m=P\operatorname{diag}((-2)^m,1,4^m)P^{-1}$，对所有非负整数 $m$ 成立。
\end{solve}

\begin{exercise}\label{prob:S-EIG-3}
设 $A=\begin{pmatrix}5&-8&-16\\4&-7&-8\\0&0&5\end{pmatrix}$。求特征值与一组特征向量基，写出对角化分解，并计算 $A^4$ 与 $\det A$。
\end{exercise}
\noindent 编者补充：S-EIG-3；相关主题见原讲义对应讲，题目数据与解答由本选编另行构造。
\begin{solve}
计算特征多项式，或逐一核对下面的列向量，可得三个不同特征值 $-3,1,5$。对应特征向量按列排列为
\[P=\begin{pmatrix}1&-2&-4\\1&-1&-2\\0&0&1\end{pmatrix},\qquad D=\operatorname{diag}(-3,1,5).\]
直接计算 $AP=PD$ 且 $\det P=1$，从而 $A=PDP^{-1}$。因此
\[A^4=PD^4P^{-1}=\begin{pmatrix}-79&160&-2496\\-80&161&-1248\\0&0&625\end{pmatrix},\qquad \det A=-15.\]
不同特征值对应的特征向量线性无关，所以这三个向量构成一组基。该方法把高次矩阵乘法转化为标量幂；一般地 $A^m=P\operatorname{diag}((-3)^m,1,5^m)P^{-1}$，对所有非负整数 $m$ 成立。
\end{solve}

\begin{exercise}\label{prob:S-EIG-4}
设 $A=\begin{pmatrix}-4&0&-20\\0&1&-10\\0&0&6\end{pmatrix}$。求特征值与一组特征向量基，写出对角化分解，并计算 $A^4$ 与 $\det A$。
\end{exercise}
\noindent 编者补充：S-EIG-4；相关主题见原讲义对应讲，题目数据与解答由本选编另行构造。
\begin{solve}
计算特征多项式，或逐一核对下面的列向量，可得三个不同特征值 $-4,1,6$。对应特征向量按列排列为
\[P=\begin{pmatrix}1&0&-2\\0&1&-2\\0&0&1\end{pmatrix},\qquad D=\operatorname{diag}(-4,1,6).\]
直接计算 $AP=PD$ 且 $\det P=1$，从而 $A=PDP^{-1}$。因此
\[A^4=PD^4P^{-1}=\begin{pmatrix}256&0&-2080\\0&1&-2590\\0&0&1296\end{pmatrix},\qquad \det A=-24.\]
不同特征值对应的特征向量线性无关，所以这三个向量构成一组基。该方法把高次矩阵乘法转化为标量幂；一般地 $A^m=P\operatorname{diag}((-4)^m,1,6^m)P^{-1}$，对所有非负整数 $m$ 成立。
\end{solve}

\begin{exercise}\label{prob:S-EIG-5}
设 $A=\begin{pmatrix}-5&-48&36\\0&1&0\\0&-12&7\end{pmatrix}$。求特征值与一组特征向量基，写出对角化分解，并计算 $A^4$ 与 $\det A$。
\end{exercise}
\noindent 编者补充：S-EIG-5；相关主题见原讲义对应讲，题目数据与解答由本选编另行构造。
\begin{solve}
计算特征多项式，或逐一核对下面的列向量，可得三个不同特征值 $-5,1,7$。对应特征向量按列排列为
\[P=\begin{pmatrix}1&4&3\\0&1&0\\0&2&1\end{pmatrix},\qquad D=\operatorname{diag}(-5,1,7).\]
直接计算 $AP=PD$ 且 $\det P=1$，从而 $A=PDP^{-1}$。因此
\[A^4=PD^4P^{-1}=\begin{pmatrix}625&-13152&5328\\0&1&0\\0&-4800&2401\end{pmatrix},\qquad \det A=-35.\]
不同特征值对应的特征向量线性无关，所以这三个向量构成一组基。该方法把高次矩阵乘法转化为标量幂；一般地 $A^m=P\operatorname{diag}((-5)^m,1,7^m)P^{-1}$，对所有非负整数 $m$ 成立。
\end{solve}

\begin{exercise}\label{prob:S-EIG-6}
设 $A=\begin{pmatrix}-6&0&0\\0&1&0\\14&7&8\end{pmatrix}$。求特征值与一组特征向量基，写出对角化分解，并计算 $A^4$ 与 $\det A$。
\end{exercise}
\noindent 编者补充：S-EIG-6；相关主题见原讲义对应讲，题目数据与解答由本选编另行构造。
\begin{solve}
计算特征多项式，或逐一核对下面的列向量，可得三个不同特征值 $-6,1,8$。对应特征向量按列排列为
\[P=\begin{pmatrix}1&0&0\\0&1&0\\-1&-1&1\end{pmatrix},\qquad D=\operatorname{diag}(-6,1,8).\]
直接计算 $AP=PD$ 且 $\det P=1$，从而 $A=PDP^{-1}$。因此
\[A^4=PD^4P^{-1}=\begin{pmatrix}1296&0&0\\0&1&0\\2800&4095&4096\end{pmatrix},\qquad \det A=-48.\]
不同特征值对应的特征向量线性无关，所以这三个向量构成一组基。该方法把高次矩阵乘法转化为标量幂；一般地 $A^m=P\operatorname{diag}((-6)^m,1,8^m)P^{-1}$，对所有非负整数 $m$ 成立。
\end{solve}

\section{若当标准形}\label{sec:topic16}
本节收录 14 道题，其中 6 道为编者补充计算题。由特征值、核维数的增长或矩阵幂的秩确定若当块，再通过若当链核对过渡矩阵；这里默认特征多项式在所用域上分裂。
\begin{exercise}\label{prob:L16-例-3}
设$V$是由两个变量$x,y$生成的次数不大于2的多项式函数（例如$xy^2$不是，因为次数等于1+2=3，$x^2+y$是）关于一般的多项式加法和数乘构成的线性空间，于是我们可以取一组基为$\{1,x,y,x^2,y^2,xy\}$，定义$\sigma\in \mathcal{L}(V)$为
    \[\sigma(f(x,y))=\dfrac{\partial}{\partial x}f(x,y).\]
    求$V$的一组基使得$\sigma$在这组基下的矩阵是若当标准形.
\end{exercise}
\noindent 来源：\href{https://github.com/yhwu-is/Linear-Algebra-Left-Undone/blob/new/%E8%AE%B2%E4%B9%89/%E4%B8%93%E9%A2%98/16%20%E8%8B%A5%E5%BD%93%E6%A0%87%E5%87%86%E5%BD%A2.tex}{L16-例-3}。
\begin{solve}
由于对多项式求导会降低次数，因此$\sigma$显然是幂零的（也可以直接写出$\sigma$在题目中给出的这组基的矩阵表示，不难得知其为对角线全为$0$的上三角矩阵），因此唯一特征值为$0$. 为了表示简便，这里用 $G_j$ 代替之前的记号 $G_j(\lambda,\sigma)=\ker (\sigma-\lambda I)^j$.

    不难得到：
    \begin{itemize}
        \item $G_1 = \ker \sigma = \operatorname{span}\{1, y, y^2\}$，$\dim G_1 = 3$；
        \item $G_2 = \ker \sigma^2 = \operatorname{span}\{1, y, y^2, x, xy\}$，$\dim G_2 = 5$；
        \item $G_3 = \ker \sigma^3 = \operatorname{span}\{1, y, y^2, x, xy, x^2\} = V$，$\dim G_3 = 6$.
    \end{itemize}

    从而根据原讲义“Young图行长度”处，Young图的行长度分别为 $r_1 = 3, r_2 = 2, r_3 = 1$，我们按照刚刚的讨论填充Young图.

    \begin{enumerate}
        \item 倒数第一层中，我们需要找到一个向量 $v_1 \in G_3 \backslash G_2$，因为这里$G_1, G_2, G_3$的基满足了包含关系，所以作差变得非常简单：直接挑出 $v_1=x^2$ 即可.
        \item 倒数第二层中，我们需要找到一个向量 $v_2 \in G_2 \backslash G_1$，且 $v_2$ 与该层已有向量 $\sigma(v_1)=2x$ 线性无关. 注意到 $G_2 = \operatorname{span}\{G_1, x, xy\}$，在 $G_2$ 中，$\{1, y, y^2, 2x, xy\}$ 线性无关且张成 $G_2$，故可取 $v_2 = xy$.
        \item 第一层中，我们需要找到一个向量 $v_3 \in G_1$，且 $v_3$ 与该层已有向量 $\sigma^2(v_1)=2, \sigma(v_2)=y$ 线性无关. 注意到 $G_1 = \operatorname{span}\{1, y, y^2\}$. 显然 $\{2, y, y^2\}$ 是 $G_1$ 的一组基，故可取 $v_3 = y^2$.
    \end{enumerate}

    从而我们得到了若当基 $\{2, 2x, x^2, y, xy, y^2\}$.
\end{solve}

\begin{exercise}\label{prob:L16-例-4}
设$V$为$n$维复向量空间，$\sigma\in \mathcal{L}(V)$，$\sigma$在基$v_1,\ldots,v_n$下的矩阵是一个若当块
    \[\begin{pmatrix}
            \lambda & 1       & 0      & \cdots & 0       \\
            0       & \lambda & 1      & \cdots & 0       \\
            \vdots  & \vdots  & \vdots & \ddots & \vdots  \\
            0       & 0       & 0      & \cdots & \lambda
        \end{pmatrix},\]
    证明：
    \begin{enumerate}
        \item $V$中包含$v_n$的不变子空间只有$V$自身；

        \item $V$中任意非零不变子空间都包含$v_1$；

        \item $V$不能分解为两个非平凡的不变子空间的直和；

        \item $V$中有且仅有$n+1$个不变子空间，它们分别是
              \[\{0\},\operatorname{span}(v_1),\operatorname{span}(v_1,v_2),\ldots,\operatorname{span}(v_1,\ldots,v_{n-1},v_n)\]
    \end{enumerate}
\end{exercise}
\noindent 来源：\href{https://github.com/yhwu-is/Linear-Algebra-Left-Undone/blob/new/%E8%AE%B2%E4%B9%89/%E4%B8%93%E9%A2%98/16%20%E8%8B%A5%E5%BD%93%E6%A0%87%E5%87%86%E5%BD%A2.tex}{L16-例-4}。
\begin{solve}
\begin{enumerate}
        \item 设$W$是包含$v_n$的不变子空间，根据线性映射矩阵表示的定义，$\sigma(v_n)=\lambda v_n+v_{n-1}\in W$，故$v_{n-1}\in W$，同理可得$v_{n-2},\ldots,v_1\in W$，即$W=V$.

        \item 设$U$为任一非零不变子空间，那么$U$中存在一个非零向量$v$，则$v$一定可以表达为
              \[v=k_1v_1+\cdots+k_sv_s(s\leqslant n),\]
              其中$k_s\neq 0$. 由$\sigma(v)\in U$以及线性映射矩阵表示有
              \begin{align*}
                  \sigma(v) & = \lambda k_1v_1+k_2(\lambda v_2+v_1)+\cdots+k_s(\lambda v_s+v_{s-1}) \\
                            & = \lambda v+k_2v_1+\cdots+k_sv_{s-1}\in U,
              \end{align*}
              因此$k_2v_1+\cdots+k_sv_{s-1}\in U$，同理进一步作用$\sigma$有$k_3v_1+\cdots+k_sv_{s-2}\in U$，以此类推有$k_sv_1\in U$. 又$k_s\neq 0$，故$v_1\in U$，命题得证.

        \item 利用上一点显然.

        \item 显然$\{0\},\operatorname{span}(v_1),\operatorname{span}(v_1,v_2),\ldots,\operatorname{span}(v_1,\ldots,v_{n-1},v_n)$都是不变子空间，下面证明不变子空间一定在它们当中.

              设$U$为任一不变子空间，若$U=\{0\}$，则也符合结论. 我们考虑非零的情况，设$\beta_1,\ldots,\beta_s$是$U$的一组基，并且假设
              \[\beta_i=a_{i1}v_1+a_{i2}v_2+\cdots+a_{it}v_t,\enspace i=1,2,\ldots,s\]
              其中满足至少有一个$i$使得$a_{it}\neq 0$（把所有基下表示末尾的0去除即可）. 对于满足$a_{it}\neq 0$的$i$，由于$U$是不变子空间，故$\sigma(\beta_i)\in U$，与第二点证明完全相同，我们得到
              \begin{align*}
                  a_{i2}v_1+a_{i3}v_2+\cdots+a_{it}v_{t-1}\in U, \\
                  a_{i3}v_1+a_{i4}v_2+\cdots+a_{it}v_{t-2}\in U, \\
                  \cdots                                         \\
                  a_{i,t-1}v_1+a_{it}v_2\in U,                   \\
                  a_{it}v_1\in U.
              \end{align*}
              由于$a_{it}\neq 0$，故由最后的式子知道$v_1\in U$. 代入倒数第二个式子知道$v_2\in U$，以此类推，最终得到$\operatorname{span}(v_1,\ldots,v_t)\subset U$，而$U\subset\operatorname{span}(v_1,\ldots,v_t)$显然（因为每个$\beta_i$都在其中），故命题得证.
    \end{enumerate}
\end{solve}

\begin{exercise}\label{prob:L16-例-5}
设$V$是复数域上的$n$维线性空间，$\sigma\in\mathcal{L}(V)$，若$\sigma$有$n$个不同的特征值$\lambda_1,\ldots,\lambda_n$，求$\sigma$的不变子空间的个数.
\end{exercise}
\noindent 来源：\href{https://github.com/yhwu-is/Linear-Algebra-Left-Undone/blob/new/%E8%AE%B2%E4%B9%89/%E4%B8%93%E9%A2%98/16%20%E8%8B%A5%E5%BD%93%E6%A0%87%E5%87%86%E5%BD%A2.tex}{L16-例-5}。
\begin{solve}
显然$\sigma$可对角化，每个对角元素就是一阶若当块，因此我们需要考虑组合这些若当块对应的基构成新的不变子空间. 由可对角化我们有
    \[V=V_{\lambda_1}\oplus\cdots\oplus V_{\lambda_n},\]
    结合$\dim V=n$可知每个特征子空间都是1维的. 根据我们组合的思想，我们知道对于任意的$1\leqslant j_1<j_2<\cdots<j_r\leqslant n$，$r=1,2,\ldots,n$，$V_{\lambda_{j_1}}\oplus\cdots\oplus V_{\lambda_{j_r}}$都是不变子空间. 这样的不变子空间加上零空间的个数总和为$\mathrm{C}_n^0+\mathrm{C}_n^1+\cdots+\mathrm{C}_n^n=2^n$.

    下面我们需要说明这些不变子空间是全部的不变子空间. 设$U$为任一不变子空间，若为零空间则符合要求，若不是零空间，设$\dim U=m>0$，$U$的一组基为$\beta_1,\ldots,\beta_m$，将其扩充为$V$的一组基$\beta_1,\ldots,\beta_m,\beta_{m+1},\ldots,\beta_n$，则在这组基下$\sigma$的矩阵为
    \[B=\begin{pmatrix}
            A & C \\ O & D
        \end{pmatrix},\]
    于是$\sigma\vert_U$的矩阵为$A$，对应特征多项式$|\lambda E-A|$，而$\sigma$的特征多项式$|\lambda E-B|=|\lambda E-A||\lambda E-D|$. 由于$\sigma$有$n$个不同的特征值，因此$\sigma\vert_U$必定有$m$个不同的特征值$\lambda_{l_1},\ldots,\lambda_{l_m}$，且$1\leqslant l_1<\cdots<l_m\leqslant n$. 而对每个特征值$\lambda_{l_i}$都有特征向量$\alpha_i\in U$，满足$\sigma\vert_U(\alpha_i)=\lambda_{l_i}\alpha_i$，因此$\sigma(\alpha_i)=\lambda_{l_i}\alpha_i$，即$\alpha_i\in V_{\lambda_{l_i}}$，故$V_{\lambda_{l_1}}\oplus\cdots\oplus V_{\lambda_{l_m}}\subset U$. 再结合
    \[\dim(V_{\lambda_{l_1}}\oplus\cdots\oplus V_{\lambda_{l_m}})=m=\dim U,\]
    故$U=V_{\lambda_{l_1}}\oplus\cdots\oplus V_{\lambda_{l_m}}$，故所有不变子空间都是之前讨论的$2^n$种之一，命题得证.
\end{solve}

\begin{exercise}\label{prob:L16-例-6}
设$A=\begin{pmatrix}
            2 & 6 & -15 \\ 1 & 1 & -5 \\ 1 & 2 & -6
        \end{pmatrix}$，求$A^{m}(m\geqslant 1)$.
\end{exercise}
\noindent 来源：\href{https://github.com/yhwu-is/Linear-Algebra-Left-Undone/blob/new/%E8%AE%B2%E4%B9%89/%E4%B8%93%E9%A2%98/16%20%E8%8B%A5%E5%BD%93%E6%A0%87%E5%87%86%E5%BD%A2.tex}{L16-例-6}。
\begin{solve}
首先求 $A$ 的若当标准形，$|\lambda E - A|=(\lambda + 1)^3$，$-1$为唯一特征值，几何重数为$3$，而$A+E=\begin{pmatrix}
            3 & 6 & -15 \\ 1 & 2 & -5 \\ 1 & 2 & -5
        \end{pmatrix}$为秩为$1$的矩阵，因此由原讲义“矩阵若当标准形唯一”处，若当块有$3-1=2$个，只可能是$J_2(-1)$和$J_1(-1)$，因此$A=PJP^{-1}$，其中$J=\begin{pmatrix}
            -1 & 1 & 0 \\ 0 & -1 & 0 \\ 0 & 0 & -1
        \end{pmatrix}$，用原讲义“求若当标准形过渡矩阵”处 中的方法，可得$P=\begin{pmatrix}
            3 & 1 & 2 \\ 1 & 0 & -1 \\ 1 & 0 & 0
        \end{pmatrix}$，$P^{-1}=\begin{pmatrix}
            0 & 0 & 1 \\ 1 & 2 & -5 \\ 0 & -1 & 1
        \end{pmatrix}$，故$A^m=PJ^mP^{-1}$，其中
        \[
            J^m=\begin{pmatrix}
                (-1)^m E + m(-1)^{m-1}J_2(0) & O \\ O & (-1)^m
            \end{pmatrix}=(-1)^m\begin{pmatrix}
                1 & -m & 0 \\ 0 & 1 & 0 \\ 0 & 0 & 1
            \end{pmatrix},
        \]
        整理可得 $A^m=(-1)^m\begin{pmatrix}
            1-3m & -6m & 15m \\ -m & 1-2m & 5m \\ -m & -2m & 1+5m
        \end{pmatrix}$.
\end{solve}

\begin{exercise}\label{prob:L16-例-7}
设$J=J_n(a)$是特征值为$a$的$n$阶若当块，求$J^m$的若当标准形，其中$m$为非零整数.当 $a=0$ 时，仅考虑正整数 $m$，因为负整数次幂无定义。
\end{exercise}
\noindent 来源：\href{https://github.com/yhwu-is/Linear-Algebra-Left-Undone/blob/new/%E8%AE%B2%E4%B9%89/%E4%B8%93%E9%A2%98/16%20%E8%8B%A5%E5%BD%93%E6%A0%87%E5%87%86%E5%BD%A2.tex}{L16-例-7}。
\begin{solve}
我们需要分别考虑$a\neq 0$和$a=0$的情况：
    \begin{enumerate}
        \item 当$a\neq 0$时，首先显然$J^m$的特征值都为$a^m$，我们先处理$m\geqslant 1$的情况，我们有
              \[J^m=(aE+J(0))^m=a^mE+\mathrm{C}_m^1a^{m-1}J_n(0)+\cdots+\mathrm{C}_m^mJ_n(0)^m.\]
              故$r(J^m-a^mE_n)=r(\mathrm{C}_m^1a^{m-1}J_n(0)+\cdots+\mathrm{C}_m^mJ_n(0)^m)=n-1$（很容易通过$J_n(0)$的幂次的特点写出这一矩阵，然后通过行列式秩判断），因此根据原讲义“若当标准形唯一”处 可知$J^m$的若当块个数为$n-(n-1)=1$，因此$J^m$的若当标准形就是$J_n(a^m)$.

              然后我们考虑$m=-1$的情况，显然$J^{-1}$的特征值为$a^{-1}$，这里我们借用原讲义“分式求逆1”处 的分式思想，得到
              \[J^{-1}=\dfrac{1}{a}E_n-\dfrac{1}{a^2}J_n(0)+\cdots+(-1)^{n-1}\dfrac{1}{a^n}J_n(0)^{n-1}.\]
              故$r(J^{-1}-\dfrac{1}{a}E_n)=r(-\dfrac{1}{a^2}J_n(0)+\cdots+(-1)^{n-1}\dfrac{1}{a^n}J_n(0)^{n-1})=n-1$，因此$J^{-1}$的若当块个数为$n-(n-1)=1$，因此$J^{-1}$的若当标准形就是$J_n(a^{-1})$.

              最后处理$m\leqslant -2$的情况，注意到$J^m=(J^{-1})^{-m}$，故根据前面$m\geqslant 1$的讨论，$J^m$的若当标准形就是$J_n((a^{-1})^{-m})=J_n(a^m)$.

              综上所述，无论$m$取什么整数值，$J^m$的若当标准形就是$J_n(a^m)$.

        \item 当$a=0$时，$J$不可逆，因此只考虑$m>0$的情形. 我们有$m\geqslant n$时$J^m=O$，此时若当标准形也就是零矩阵. 我们考虑$m<n$的情况，我们需要使用原讲义“若当标准形唯一”处 的结论来辅助我们求解$J^m$的若当标准形，因此我们需要研究$(J^m)^k$的秩. 我们做带余除法，$n=mq+r(0\leqslant r<m)$，则根据$J_n(0)$的幂次的特点，我们有
              \[r((J^m)^k)=n-mk,\enspace 0\leqslant k\leqslant q,\enspace r((J^m)^k)=0,\enspace k\geqslant q+1.\]
              因此
              \begin{enumerate}
                  \item $1\leqslant k<q$时，$J_k(0)$的个数为$r((J^m)^{k-1})+r((J^m)^{k+1})-2r((J^m)^k)=(n-m(k-1))+(n-m(k+1))-2(n-mk)=0$；
                  \item $k=q$时，$J_k(0)$的个数为$r((J^m)^{q-1})+r((J^m)^{q+1})-2r((J^m)^q)=(n-m(q-1))+0-2(n-mq)=m-r$；
                  \item $k=q+1$时，$J_k(0)$的个数为$r((J^m)^q)+r((J^m)^{q+2})-2r((J^m)^{q+1})=(n-mq)+0-0=r$.
                  \item $k>q+1$时，$J_k(0)$的个数为0.
              \end{enumerate}
              综上所述，$J^m$的若当标准形是$\operatorname{diag}(J_q(0),\ldots,J_q(0),J_{q+1}(0),\ldots,J_{q+1}(0))$，其中$J_q(0)$的个数为$m-r$，$J_{q+1}(0)$的个数为$r$.
    \end{enumerate}
\end{solve}

\begin{exercise}\label{prob:L16-例-9}
利用若当标准形求解初值问题
    \[
        X'(t)=\begin{pmatrix}
            2 & 3 & 2 \\ 1 & 8 & 2 \\ -2 & -14 & -3
        \end{pmatrix}X(t),\qquad
        X(0)=\begin{pmatrix}
            -1 \\ 0 \\ 0
        \end{pmatrix}.
    \]
\end{exercise}
\noindent 来源：\href{https://github.com/yhwu-is/Linear-Algebra-Left-Undone/blob/new/%E8%AE%B2%E4%B9%89/%E4%B8%93%E9%A2%98/16%20%E8%8B%A5%E5%BD%93%E6%A0%87%E5%87%86%E5%BD%A2.tex}{L16-例-9}。
\begin{solve}
记
    \[
        A=\begin{pmatrix}
            2 & 3 & 2 \\ 1 & 8 & 2 \\ -2 & -14 & -3
        \end{pmatrix}.
    \]
    由原讲义“求若当标准形过渡矩阵”处 可知，$A$的若当标准形为
    \[
        J=\begin{pmatrix}
            3 & 1 & 0 \\ 0 & 3 & 0 \\ 0 & 0 & 1
        \end{pmatrix},
        \qquad
        P=\begin{pmatrix}
            1 & -1 & 2 \\ -1 & 0 & 0 \\ 2 & 0 & -1
        \end{pmatrix},
        \qquad
        P^{-1}AP=J.
    \]

    令$X(t)=PY(t)$，则
    \[
        PY'(t)=X'(t)=AX(t)=APY(t)=PJY(t).
    \]
    由于$P$可逆，故$Y'(t)=JY(t)$. 又注意到初值向量$X(0)=(-1,0,0)^\mathrm{T}$恰为$P$的第二列，因此
    \[
        Y(0)=P^{-1}X(0)=\begin{pmatrix}
            0 \\ 1 \\ 0
        \end{pmatrix}.
    \]

    现在只需求解若当标准形对应的方程组. 由于
    \[
        J=\operatorname{diag}(J_2(3),1),
    \]
    因而
    \[
        \mathrm{e}^{Jt}=\begin{pmatrix}
            \mathrm{e}^{3t} & t\mathrm{e}^{3t} & 0 \\
            0 & \mathrm{e}^{3t} & 0 \\
            0 & 0 & \mathrm{e}^{t}
        \end{pmatrix}.
    \]
    故
    \[
        Y(t)=\mathrm{e}^{Jt}Y(0)=\begin{pmatrix}
            t\mathrm{e}^{3t} \\ \mathrm{e}^{3t} \\ 0
        \end{pmatrix}.
    \]

    最后代回$X(t)=PY(t)$，得到
    \[
        X(t)=\begin{pmatrix}
            1 & -1 & 2 \\ -1 & 0 & 0 \\ 2 & 0 & -1
        \end{pmatrix}\begin{pmatrix}
            t\mathrm{e}^{3t} \\ \mathrm{e}^{3t} \\ 0
        \end{pmatrix}
        =\mathrm{e}^{3t}\begin{pmatrix}
            t-1 \\ -t \\ 2t
        \end{pmatrix}.
    \]
    这就是所求的解.
\end{solve}

\begin{exercise}\label{prob:L16-例-10}
求解线性递推方程$a_n = 5a_{n - 1} - 8a_{n - 2} + 4a_{n - 3}$，其中$a_0 = 2, a_1 = 4, a_2 = 9$.
\end{exercise}
\noindent 来源：\href{https://github.com/yhwu-is/Linear-Algebra-Left-Undone/blob/new/%E8%AE%B2%E4%B9%89/%E4%B8%93%E9%A2%98/16%20%E8%8B%A5%E5%BD%93%E6%A0%87%E5%87%86%E5%BD%A2.tex}{L16-例-10}。
\begin{solve}
设这一递推方程的特征多项式为
    \[
        p(\lambda)=\lambda^3-5\lambda^2+8\lambda-4=(\lambda-1)(\lambda-2)^2.
    \]
    因此特征值为$1$与$2$，其中$2$是二重根. 根据前面对重根情形的讨论，数列$\{a_n\}$的通项应具有形式
    \[
        a_n=c_1\cdot 1^n+c_2\cdot 2^n+c_3\binom{n}{1}2^{n-1}
        =c_1+c_2\cdot 2^n+c_3n2^{n-1},
    \]
    其中$c_1,c_2,c_3$为待定常数.

    将初值代入上式，得到方程组
    \[
        \begin{cases}
            c_1+c_2=2, \\
            c_1+2c_2+c_3=4, \\
            c_1+4c_2+4c_3=9.
        \end{cases}
    \]
    解得
    \[
        c_1=1,\qquad c_2=1,\qquad c_3=1.
    \]
    从而
    \[
        a_n=1+2^n+n2^{n-1}.
    \]
    这就是所求的通项公式.
\end{solve}

\begin{exercise}\label{prob:L16-例-12}
求矩阵
    \[
        A=\begin{pmatrix}
            0 & 0 & 4 \\
            1 & 0 & -6 \\
            0 & 1 & 4
        \end{pmatrix}
    \]
    在实数域上的若当标准形.
\end{exercise}
\noindent 来源：\href{https://github.com/yhwu-is/Linear-Algebra-Left-Undone/blob/new/%E8%AE%B2%E4%B9%89/%E4%B8%93%E9%A2%98/16%20%E8%8B%A5%E5%BD%93%E6%A0%87%E5%87%86%E5%BD%A2.tex}{L16-例-12}。
\begin{solve}
按照前面的总结，我们先将$A$视为复矩阵来求其复数域上的若当标准形. 计算特征多项式：
    \[
        |\lambda E-A|
        =\begin{vmatrix}
            \lambda & 0 & -4 \\
            -1 & \lambda & 6 \\
            0 & -1 & \lambda-4
        \end{vmatrix}
        =\lambda(\lambda^2-4\lambda+6)-4
        =\lambda^3-4\lambda^2+6\lambda-4.
    \]
    不难得到 $A$ 的特征值为 $2,\qquad 1+\mathrm{i},\qquad 1-\mathrm{i}$. 这三个特征值两两不同，所以$A$在复数域上的若当标准形就是对角矩阵 $J_{\mathbf{C}}=\operatorname{diag}(2,1+\mathrm{i},1-\mathrm{i})$. 现在回到实数域上. 其中实特征值$2$对应一个通常的一阶若当块$J_1(2)=(2)$；一对共轭复特征值$1\pm\mathrm{i}$对应一个二维实块 $R(1,1)=\begin{pmatrix}
            1 & 1 \\
            -1 & 1
        \end{pmatrix}$. 因而$A$在实数域上的若当标准形可取为
    \[
        J_{\mathbf{R}}=\operatorname{diag}(2,R(1,1))
        =\begin{pmatrix}
            2 & 0 & 0 \\
            0 & 1 & 1 \\
            0 & -1 & 1
        \end{pmatrix}.
    \]
\end{solve}

\begin{exercise}\label{prob:S-JORDAN-1}
设 $A=\begin{pmatrix}5&10&-3\\-1&-1&1\\0&1&2\end{pmatrix}$。求若当标准形和最小多项式，计算 $\dim\ker(A-2I)^r$（$r=1,2,3$），并给出 $A^m$ 的通式。
\end{exercise}
\noindent 编者补充：S-JORDAN-1；相关主题见原讲义对应讲，题目数据与解答由本选编另行构造。
\begin{solve}
取 $P=\begin{pmatrix}1&-3&0\\0&1&0\\1&0&1\end{pmatrix}$，逐列相乘可验证
\[AP=PJ,\qquad J=\begin{pmatrix}2&1&0\\0&2&1\\0&0&2\end{pmatrix}.\]
所以若当标准形是一个三阶块 $J_3(2)$。记 $N=A-2I$，有 $N^3=O$、$N^2\ne O$，最小多项式为 $(x-2)^3$。在若当基下，$N$ 依次把第三个基向量送到第二个、第二个送到第一个、第一个送到零，故三个核维数依次为 $1,2,3$。因为 $N$ 与 $I$ 可交换，二项式公式给出
\[A^m=2^mI+m2^{m-1}N+\binom m22^{m-2}N^2\quad(m\ge2).\]
$m=0,1$ 时分别为 $I,A$。最大若当块为三阶，说明该矩阵不能对角化。
\end{solve}

\begin{exercise}\label{prob:S-JORDAN-2}
设 $A=\begin{pmatrix}3&1&-3\\0&3&1\\0&0&3\end{pmatrix}$。求若当标准形和最小多项式，计算 $\dim\ker(A-3I)^r$（$r=1,2,3$），并给出 $A^m$ 的通式。
\end{exercise}
\noindent 编者补充：S-JORDAN-2；相关主题见原讲义对应讲，题目数据与解答由本选编另行构造。
\begin{solve}
取 $P=\begin{pmatrix}1&-1&-2\\0&1&2\\0&0&1\end{pmatrix}$，逐列相乘可验证
\[AP=PJ,\qquad J=\begin{pmatrix}3&1&0\\0&3&1\\0&0&3\end{pmatrix}.\]
所以若当标准形是一个三阶块 $J_3(3)$。记 $N=A-3I$，有 $N^3=O$、$N^2\ne O$，最小多项式为 $(x-3)^3$。在若当基下，$N$ 依次把第三个基向量送到第二个、第二个送到第一个、第一个送到零，故三个核维数依次为 $1,2,3$。因为 $N$ 与 $I$ 可交换，二项式公式给出
\[A^m=3^mI+m3^{m-1}N+\binom m23^{m-2}N^2\quad(m\ge2).\]
$m=0,1$ 时分别为 $I,A$。最大若当块为三阶，说明该矩阵不能对角化。
\end{solve}

\begin{exercise}\label{prob:S-JORDAN-3}
设 $A=\begin{pmatrix}3&1&-4\\-1&5&-3\\0&0&4\end{pmatrix}$。求若当标准形和最小多项式，计算 $\dim\ker(A-4I)^r$（$r=1,2,3$），并给出 $A^m$ 的通式。
\end{exercise}
\noindent 编者补充：S-JORDAN-3；相关主题见原讲义对应讲，题目数据与解答由本选编另行构造。
\begin{solve}
取 $P=\begin{pmatrix}1&-2&-4\\1&-1&-2\\0&0&1\end{pmatrix}$，逐列相乘可验证
\[AP=PJ,\qquad J=\begin{pmatrix}4&1&0\\0&4&1\\0&0&4\end{pmatrix}.\]
所以若当标准形是一个三阶块 $J_3(4)$。记 $N=A-4I$，有 $N^3=O$、$N^2\ne O$，最小多项式为 $(x-4)^3$。在若当基下，$N$ 依次把第三个基向量送到第二个、第二个送到第一个、第一个送到零，故三个核维数依次为 $1,2,3$。因为 $N$ 与 $I$ 可交换，二项式公式给出
\[A^m=4^mI+m4^{m-1}N+\binom m24^{m-2}N^2\quad(m\ge2).\]
$m=0,1$ 时分别为 $I,A$。最大若当块为三阶，说明该矩阵不能对角化。
\end{solve}

\begin{exercise}\label{prob:S-JORDAN-4}
设 $A=\begin{pmatrix}5&1&2\\0&5&1\\0&0&5\end{pmatrix}$。求若当标准形和最小多项式，计算 $\dim\ker(A-5I)^r$（$r=1,2,3$），并给出 $A^m$ 的通式。
\end{exercise}
\noindent 编者补充：S-JORDAN-4；相关主题见原讲义对应讲，题目数据与解答由本选编另行构造。
\begin{solve}
取 $P=\begin{pmatrix}1&0&-2\\0&1&-2\\0&0&1\end{pmatrix}$，逐列相乘可验证
\[AP=PJ,\qquad J=\begin{pmatrix}5&1&0\\0&5&1\\0&0&5\end{pmatrix}.\]
所以若当标准形是一个三阶块 $J_3(5)$。记 $N=A-5I$，有 $N^3=O$、$N^2\ne O$，最小多项式为 $(x-5)^3$。在若当基下，$N$ 依次把第三个基向量送到第二个、第二个送到第一个、第一个送到零，故三个核维数依次为 $1,2,3$。因为 $N$ 与 $I$ 可交换，二项式公式给出
\[A^m=5^mI+m5^{m-1}N+\binom m25^{m-2}N^2\quad(m\ge2).\]
$m=0,1$ 时分别为 $I,A$。最大若当块为三阶，说明该矩阵不能对角化。
\end{solve}

\begin{exercise}\label{prob:S-JORDAN-5}
设 $A=\begin{pmatrix}6&-7&4\\0&4&1\\0&-4&8\end{pmatrix}$。求若当标准形和最小多项式，计算 $\dim\ker(A-6I)^r$（$r=1,2,3$），并给出 $A^m$ 的通式。
\end{exercise}
\noindent 编者补充：S-JORDAN-5；相关主题见原讲义对应讲，题目数据与解答由本选编另行构造。
\begin{solve}
取 $P=\begin{pmatrix}1&4&3\\0&1&0\\0&2&1\end{pmatrix}$，逐列相乘可验证
\[AP=PJ,\qquad J=\begin{pmatrix}6&1&0\\0&6&1\\0&0&6\end{pmatrix}.\]
所以若当标准形是一个三阶块 $J_3(6)$。记 $N=A-6I$，有 $N^3=O$、$N^2\ne O$，最小多项式为 $(x-6)^3$。在若当基下，$N$ 依次把第三个基向量送到第二个、第二个送到第一个、第一个送到零，故三个核维数依次为 $1,2,3$。因为 $N$ 与 $I$ 可交换，二项式公式给出
\[A^m=6^mI+m6^{m-1}N+\binom m26^{m-2}N^2\quad(m\ge2).\]
$m=0,1$ 时分别为 $I,A$。最大若当块为三阶，说明该矩阵不能对角化。
\end{solve}

\begin{exercise}\label{prob:S-JORDAN-6}
设 $A=\begin{pmatrix}7&1&0\\1&8&1\\-1&-2&6\end{pmatrix}$。求若当标准形和最小多项式，计算 $\dim\ker(A-7I)^r$（$r=1,2,3$），并给出 $A^m$ 的通式。
\end{exercise}
\noindent 编者补充：S-JORDAN-6；相关主题见原讲义对应讲，题目数据与解答由本选编另行构造。
\begin{solve}
取 $P=\begin{pmatrix}1&0&0\\0&1&0\\-1&-1&1\end{pmatrix}$，逐列相乘可验证
\[AP=PJ,\qquad J=\begin{pmatrix}7&1&0\\0&7&1\\0&0&7\end{pmatrix}.\]
所以若当标准形是一个三阶块 $J_3(7)$。记 $N=A-7I$，有 $N^3=O$、$N^2\ne O$，最小多项式为 $(x-7)^3$。在若当基下，$N$ 依次把第三个基向量送到第二个、第二个送到第一个、第一个送到零，故三个核维数依次为 $1,2,3$。因为 $N$ 与 $I$ 可交换，二项式公式给出
\[A^m=7^mI+m7^{m-1}N+\binom m27^{m-2}N^2\quad(m\ge2).\]
$m=0,1$ 时分别为 $I,A$。最大若当块为三阶，说明该矩阵不能对角化。
\end{solve}

\section{最小多项式}\label{sec:topic17}
本节收录 2 道题。最小多项式是使矩阵或算子为零的首一多项式中次数最低者。它整除特征多项式，根的指数由最大若当块决定。
\begin{exercise}\label{prob:L17-例-3}
求矩阵
    \[
        A=\begin{pmatrix}
                2 & 3 & 2 \\
                1 & 8 & 2 \\
                -2 & -14 & -3
            \end{pmatrix}
    \]
    的初等因子组.
\end{exercise}
\noindent 来源：\href{https://github.com/yhwu-is/Linear-Algebra-Left-Undone/blob/new/%E8%AE%B2%E4%B9%89/%E4%B8%93%E9%A2%98/17%20%E5%A4%9A%E9%A1%B9%E5%BC%8F%E7%9A%84%E8%BF%9B%E4%B8%80%E6%AD%A5%E8%AE%A8%E8%AE%BA.tex}{L17-例-3}。
\begin{solve}
这一矩阵正是原讲义“求若当标准形过渡矩阵”处中讨论过的矩阵. 我们已经求得它的若当标准形为
    \[
        J=\begin{pmatrix}
                3 & 1 & 0 \\
                0 & 3 & 0 \\
                0 & 0 & 1
            \end{pmatrix},
    \]
    即它有一个关于特征值$3$的二阶若当块$J_2(3)$，以及一个关于特征值$1$的一阶若当块$J_1(1)$. 根据前面的讨论可知，每一个若当块$J_r(\lambda_i)$都对应一个初等因子$(\lambda-\lambda_i)^r$，因此矩阵$A$的初等因子组为
    \[
        (\lambda-3)^2,\ \lambda-1.
    \]
\end{solve}

\begin{exercise}\label{prob:L17-例-4}
求矩阵 $A = \begin{pmatrix}
            0 & 0 & 0 \\ 1 & 0 & 2 \\ 2 & 1 & -1
        \end{pmatrix}$ 的极小多项式.
\end{exercise}
\noindent 来源：\href{https://github.com/yhwu-is/Linear-Algebra-Left-Undone/blob/new/%E8%AE%B2%E4%B9%89/%E4%B8%93%E9%A2%98/17%20%E5%A4%9A%E9%A1%B9%E5%BC%8F%E7%9A%84%E8%BF%9B%E4%B8%80%E6%AD%A5%E8%AE%A8%E8%AE%BA.tex}{L17-例-4}。
\begin{solve}
对于矩阵$A$，我们依次考察方程组
    \[
        a_0I+A=O,\qquad a_0I+a_1A+A^2=O,\qquad a_0I+a_1A+a_2A^2+A^3=O.
    \]
    先计算
    \[
        A^2=\begin{pmatrix}
                0 & 0  & 0  \\
                4 & 2  & -2 \\
                -1 & -1 & 3
            \end{pmatrix},
        \qquad
        A^3=\begin{pmatrix}
                0 & 0  & 0  \\
                -2 & -2 & 6  \\
                5 & 3  & -5
            \end{pmatrix}.
    \]
    读者不难验证，前两个方程组均无解. 根据三个方程组可以解得 $a_0 = 0, a_1 = -2, a_2 = 1$. 于是 $A^3 + A^2 - 2A=O$. 因为次数小于 $3$ 的首一多项式都不能化零 $A$，所以 $A$ 的极小多项式为 $\lambda^3+\lambda^2-2\lambda$.
\end{solve}

\section{内积空间与正交化}\label{sec:topic19}
本节收录 18 道题，其中 12 道为编者补充计算题。先确认内积约定，再进行正交投影和施密特正交化；复数情形必须保留共轭，归一化前检查向量非零。
\begin{exercise}\label{prob:L19-例-1}
验证下面的例子确实满足内积的定义：
    \begin{enumerate}
        \item $\mathbf{F}^{n}$ 上的欧几里得内积定义为：
              \[\left\langle (w_1, \ldots, w_n), (z_1, \ldots, z_n)\right\rangle = w_1\overline{z_1} + \cdots + w_n\overline{z_n} = \vec{w}\overline{\vec{z}}^{\mathrm{T}}.\]
              对应范数：
              \[\left\lVert (z_1, \ldots, z_n) \right\rVert  = \sqrt{\lvert z_1 \rvert^2 + \cdots + \lvert z_n \rvert^2}.\]

        \item \label{item:23:内积和范数:2}
              定义在 $ \left[-1, 1\right] $ 上的连续实值函数构成的向量空间可定义内积如下：
              \[\left\langle f, g\right\rangle = \int_{-1}^1f(x)g(x)\,\mathrm{d}x.\]
              对应范数：
              \[\left\lVert f \right\rVert = \sqrt{\int_{-1}^1(f(x))^2\,\mathrm{d}x}.\]
    \end{enumerate}
\end{exercise}
\noindent 来源：\href{https://github.com/yhwu-is/Linear-Algebra-Left-Undone/blob/new/%E8%AE%B2%E4%B9%89/%E4%B8%93%E9%A2%98/19%20%E5%86%85%E7%A7%AF%E7%A9%BA%E9%97%B4.tex}{L19-例-1}。
\begin{solve}
\begin{enumerate}
        \item 对于$u=(w_1,\ldots,w_n),v=(z_1,\ldots,z_n),w=(t_1,\ldots,t_n)\in\mathbf{F}^n$，$\lambda\in\mathbf{F}$，有
            \begin{enumerate}
                \item 正定性：\[\langle v,v\rangle = \sum_{i=1}^n z_i\overline{z_i} = \sum_{i=1}^n \lvert z_i\rvert^2\geqslant 0,\]
                且
                \[\langle v,v\rangle = 0 \iff z_i=0\enspace(i=1,2,\ldots,n) \iff v=\vec{0}.\]
                故满足正定性.
                \item 第一个位置的加性：\[\langle u+v,w\rangle = \sum_{i=1}^n(w_i+z_i)\overline{t_i}=\sum_{i=1}^n w_i\overline{t_i}+\sum_{i=1}^n z_i\overline{t_i}=\langle u,w\rangle+\langle v,w\rangle.\]
                故满足第一个位置的加性；
                \item 第一个位置的齐性：\[\langle \lambda u,v\rangle = \sum_{i=1}^n(\lambda w_i\overline{z_i})=\lambda\sum_{i=1}^nw_i\overline{z_i}=\lambda\langle u,v\rangle.\]
                故满足第一个位置的齐性；
                \item 共轭对称性：\[\langle u,v\rangle = \sum_{i=1}^nw_i\overline{z_i}=\overline{\sum_{i=1}^nz_i\overline{w_i}}=\overline{\langle v,u\rangle}.\]
                故满足共轭对称性.
            \end{enumerate}
            从而确实满足内积定义.
        \item 对于 $f, g, h \in C[-1,1]$，$\lambda \in \mathbf{R}$，有
            \begin{enumerate}
                \item 正定性：\[\langle f,f\rangle = \int_{-1}^1(f(x))^2\,\mathrm{d}x\geqslant 0,\]
                且
                \[\langle f,f\rangle = 0 \iff f(x)=0\enspace(\forall x\in[-1,1]) \iff f=\vec{0}.\]
                故满足正定性.
                \item 第一个位置的加性：
                \begin{align*}
                    \langle f+g,h\rangle &= \int_{-1}^1(f(x)+g(x))h(x)\,\mathrm{d}x=\int_{-1}^1f(x)h(x)\,\mathrm{d}x+\int_{-1}^1g(x)h(x)\,\mathrm{d}x \\
                    &=\langle f,h\rangle+\langle g,h\rangle.
                \end{align*}
                故满足第一个位置的加性；
                \item 第一个位置的齐性：\[\langle \lambda f,g\rangle = \int_{-1}^1(\lambda f(x))g(x)\,\mathrm{d}x=\lambda\int_{-1}^1f(x)g(x)\,\mathrm{d}x=\lambda\langle f,g\rangle.\]
                故满足第一个位置的齐性；
                \item 共轭对称性：\[\langle f,g\rangle = \int_{-1}^1f(x)g(x)\,\mathrm{d}x=\int_{-1}^1g(x)f(x)\,\mathrm{d}x=\langle g,f\rangle.\]
                故满足共轭对称性（注意到是实值函数，因此共轭对称性退化为对称性）.
            \end{enumerate}
            从而确实满足内积定义.
    \end{enumerate}
\end{solve}

\begin{exercise}\label{prob:L19-例-2}
证明下列式子成立：
    \begin{enumerate}
        \item $\mathbf{F} = \mathbf{R}$ 时：
              \begin{gather}
                  \label{eq:23:内积和范数的性质1}
                  \langle u, v\rangle = \frac{1}{4}\left( \lVert u + v \rVert^2 - \lVert u - v \rVert^2\right) \\
                  \label{eq:23:内积和范数的性质2}
                  \langle Tu, v\rangle + \langle Tv, u\rangle = \frac{1}{2}\left(\langle T(u + v), u + v\rangle - \langle T(u - v), u - v\rangle\right)
              \end{gather}

        \item $\mathbf{F} = \mathbf{C}$ 时：
              \begin{gather}
                  \label{eq:23:内积和范数的性质3}
                  \langle u, v\rangle = \frac{1}{4}\left((\lVert u + v \rVert^2 - \lVert u - v \rVert^2) + \mathrm{i}(\lVert u + \mathrm{i} v \rVert^2 - \lVert u - \mathrm{i} v \rVert^2)\right) \\
                  \label{eq:23:内积和范数的性质4}
                  \begin{aligned}
                      \langle Tu, v\rangle ={} & \frac{1}{4}  ((\langle T(u + v), u + v\rangle - \langle T(u - v), u - v\rangle)     \\
                                               & + \mathrm{i}(\langle T(u + \mathrm{i} v), u + \mathrm{i} v\rangle + \langle T(u - \mathrm{i} v), u - \mathrm{i} v\rangle)).
                  \end{aligned}
              \end{gather}
    \end{enumerate}
\end{exercise}
\noindent 来源：\href{https://github.com/yhwu-is/Linear-Algebra-Left-Undone/blob/new/%E8%AE%B2%E4%B9%89/%E4%B8%93%E9%A2%98/19%20%E5%86%85%E7%A7%AF%E7%A9%BA%E9%97%B4.tex}{L19-例-2}。
\begin{solve}
\begin{enumerate}
        \item 利用原讲义“勾股定理展开”处 将右侧展开得证（由于是实内积空间，因此 $\operatorname{Re}(\langle u, v \rangle) = \langle u, v \rangle$）.

        \item 我们展开等式右侧的两项，利用线性算子 $T$ 的线性性：
        \begin{align*}
            & \quad \langle T(u + v), u + v\rangle - \langle T(u - v), u - v\rangle \\
            & = \left(\langle Tu, u \rangle + \langle Tu, v \rangle + \langle Tv, u \rangle + \langle Tv, v \rangle\right) - \left(\langle Tu, u \rangle - \langle Tu, v \rangle - \langle Tv, u \rangle + \langle Tv, v \rangle\right) \\
            & = 2 (\langle Tu, v \rangle + \langle Tv, u \rangle),
        \end{align*}
        两边同除以 $2$ 即得.

        \item 利用原讲义“勾股定理展开”处 将右侧展开：
        \begin{gather*}
            \lVert u + v \rVert^2 - \lVert u - v \rVert^2 = 4 \operatorname{Re}(\langle u, v \rangle), \\
            \lVert u + \mathrm{i} v \rVert^2 - \lVert u - \mathrm{i} v \rVert^2 = 4 \operatorname{Re}(\langle u, \mathrm{i} v \rangle) = 4 \operatorname{Re}(-\mathrm{i} \langle u, v \rangle) = 4 \operatorname{Im}(\langle u, v \rangle).
        \end{gather*}
        故等号右侧为 $ \operatorname{Re}(\langle u, v \rangle) + \mathrm{i} \operatorname{Im}(\langle u, v \rangle) = \langle u, v \rangle$. 得证.

        \item
        \begin{gather*}
            \langle T(u + v), u + v\rangle - \langle T(u - v), u - v\rangle = 2\langle Tu, v \rangle + 2\langle Tv, u \rangle,\\
            \langle T(u + \mathrm{i} v), u + \mathrm{i} v\rangle - \langle T(u - \mathrm{i} v), u - \mathrm{i} v\rangle = 2\langle Tu, \mathrm{i} v \rangle + 2\langle \mathrm{i} Tv, u \rangle = 2\mathrm{i} (\langle Tv, u \rangle - \langle Tu, v \rangle),
        \end{gather*}
        代入等式右边可得
        \begin{align*}
            &\quad \frac{1}{4}\left(\left(2\langle Tu, v \rangle + 2\langle Tv, u \rangle\right) + \mathrm{i} \left(2\mathrm{i} (\langle Tv, u \rangle - \langle Tu, v \rangle)\right)\right) \\
            & = \frac{1}{4}\left(2\langle Tu, v \rangle + 2\langle Tv, u \rangle - 2\langle Tv, u \rangle + 2\langle Tu, v \rangle\right) \\
            & = \langle Tu, v \rangle.
        \end{align*}
        得证.
    \end{enumerate}
\end{solve}

\begin{exercise}\label{prob:L19-例-3}
设 $e_1, \ldots , e_m$ 是 $ V $ 的标准正交组. 设 $ v \in V $. 证明
    \[ \lVert v \rVert^2 = \lvert \langle v, e_1\rangle \rvert^2 + \cdots + \lvert \langle v, e_m\rangle \rvert^2 \]
    当且仅当 $ v \in \operatorname{span}(e_1, \ldots , e_m)$.
\end{exercise}
\noindent 来源：\href{https://github.com/yhwu-is/Linear-Algebra-Left-Undone/blob/new/%E8%AE%B2%E4%B9%89/%E4%B8%93%E9%A2%98/19%20%E5%86%85%E7%A7%AF%E7%A9%BA%E9%97%B4.tex}{L19-例-3}。
\begin{solve}
记 $W = \operatorname{span}(e_1, \ldots , e_m)$.
    \begin{enumerate}
        \item 必要性：构造向量 $u = \displaystyle\sum_{i=1}^{m} \langle v, e_i \rangle e_i$. 显然 $u \in W$.

        令 $w = v - u$，注意到：
        \[
            \langle w, e_j \rangle = \left\langle v - \sum_{i=1}^{m} \langle v, e_i \rangle e_i, e_j \right\rangle = \langle v, e_j \rangle - \sum_{i=1}^{m} \langle v, e_i \rangle \delta_{ij} = \langle v, e_j \rangle - \langle v, e_j \rangle = 0.
        \]
        因此 $w$ 与所有 $e_j$ 正交，由于 $u$ 为 $e_1, \ldots , e_m$ 的线性组合，故 $w$ 与 $u$ 也正交.

        计算 $\lVert w \rVert^2$:
        \begin{align*}
            \|w\|^2 &= \langle w, v-u \rangle = \langle w, v \rangle - \langle w, u \rangle \\
            &= \langle v-u, v \rangle = \langle v, v \rangle - \langle u, v \rangle \\
            &= \|v\|^2 - \sum_{i=1}^{m} \langle v, e_i \rangle \langle e_i, v \rangle \\
            &= \|v\|^2 - \sum_{i=1}^{m} \langle v, e_i \rangle \overline{\langle v, e_i \rangle} \\
            &= \|v\|^2 - \sum_{i=1}^{m} |\langle v, e_i \rangle|^2.
        \end{align*}
        根据假设 $\lVert v \rVert^2 = \displaystyle\sum_{i=1}^{m} \lvert \langle v, e_i \rangle \rvert^2$，故 $\lVert w \rVert^2 = 0$.
        由范数正定性 $w = \vec{0}$，即 $v = u$.

        由于 $u \in W$，所以 $v \in \operatorname{span}(e_1, \ldots , e_m)$. 必要性得证.

        \item 充分性：对子空间 $W$ 使用原讲义“Fourier展开”处 即可. 充分性得证.
    \end{enumerate}
\end{solve}

\begin{exercise}\label{prob:L19-例-4}
求 $\mathbf{R}[x]_2$ 的一组标准正交基，内积定义为 $\langle p, q \rangle = \displaystyle\int_{-1}^1 p(x)q(x)\,\mathrm{d}x$.
\end{exercise}
\noindent 来源：\href{https://github.com/yhwu-is/Linear-Algebra-Left-Undone/blob/new/%E8%AE%B2%E4%B9%89/%E4%B8%93%E9%A2%98/19%20%E5%86%85%E7%A7%AF%E7%A9%BA%E9%97%B4.tex}{L19-例-4}。
\begin{solve}
取 $\mathbf{R}[x]_3$ 的一组基 $\{v_1, v_2, v_3\} = \{1, x, x^2\}$.

    我们使用 Gram-Schmidt 过程.

    首先计算 $e_1$:
    \[ \lVert v_1 \rVert^2 = \langle 1, 1 \rangle = \int_{-1}^1 1 \,\mathrm{d}x = 2, e_1 = \frac{v_1}{\lVert v_1 \rVert} = \frac{1}{\sqrt{2}}. \]

    正交化：
    \[ u_2 = v_2 - \langle v_2, e_1 \rangle e_1, \]
    计算投影系数：
    \[ \langle v_2, e_1 \rangle = \left\langle x, \frac{1}{\sqrt{2}} \right\rangle = \frac{1}{\sqrt{2}} \int_{-1}^1 x \,\mathrm{d}x = 0, \]
    因此 $u_2 = x - 0 \cdot e_1 = x$.

    单位化：
    \[ \lVert u_2 \rVert^2 = \langle x, x \rangle = \int_{-1}^1 x^2 \,\mathrm{d}x = \left.\left(\frac{1}{3}x^3\right)\right\vert_{-1}^1= \frac{2}{3}, e_2 = \frac{u_2}{\lVert u_2 \rVert} = \frac{x}{\sqrt{2/3}} = \sqrt{\frac{3}{2}} x. \]

    正交化：
    \[ u_3 = v_3 - \langle v_3, e_1 \rangle e_1 - \langle v_3, e_2 \rangle e_2, \]
    计算投影系数：
    \begin{gather*}
        \langle v_3, e_1 \rangle = \left\langle x^2, \frac{1}{\sqrt{2}} \right\rangle = \frac{1}{\sqrt{2}} \int_{-1}^1 x^2 \,\mathrm{d}x = \frac{\sqrt{2}}{3}, \\
        \langle v_3, e_2 \rangle = \left\langle x^2, \sqrt{\frac{3}{2}} x \right\rangle = \sqrt{\frac{3}{2}} \int_{-1}^1 x^3 \,\mathrm{d}x = 0,
    \end{gather*}
    计算 $u_3$:
    \[ u_3 = x^2 - \left(\frac{\sqrt{2}}{3}\right) e_1 - 0 = x^2 - \frac{\sqrt{2}}{3} \cdot \frac{1}{\sqrt{2}} = x^2 - \frac{1}{3}. \]

    单位化：
    \[ \lVert u_3 \rVert^2 = \left\langle x^2 - \frac{1}{3}, x^2 - \frac{1}{3} \right\rangle = \int_{-1}^1 \left(x^4 - \frac{2}{3}x^2 + \frac{1}{9}\right) \,\mathrm{d}x = \frac{8}{45}, \]
    \[ e_3 = \frac{u_3}{\lVert u_3 \rVert} = \frac{x^2 - 1/3}{\sqrt{8/45}} = \sqrt{\frac{45}{8}} \left(x^2 - \frac{1}{3}\right) = \frac{\sqrt{10}}{4} (3x^2 - 1). \]

    综上，$\mathbf{R}[x]_3$ 的一组标准正交基是 $\left\{e_1, e_2, e_3\right\}$，其中
    \[ e_1(x) = \frac{1}{\sqrt{2}}, \quad e_2(x) = \sqrt{\frac{3}{2}} x, \quad e_3(x) = \frac{\sqrt{10}}{4} (3x^2 - 1). \]
\end{solve}

\begin{exercise}\label{prob:L19-例-5}
求矩阵
    \[
        A = \begin{pmatrix}
            1 & 1 & 0 \\
            1 & 0 & 1 \\
            0 & 1 & 1
        \end{pmatrix}
    \]
    的 QR 分解.
\end{exercise}
\noindent 来源：\href{https://github.com/yhwu-is/Linear-Algebra-Left-Undone/blob/new/%E8%AE%B2%E4%B9%89/%E4%B8%93%E9%A2%98/19%20%E5%86%85%E7%A7%AF%E7%A9%BA%E9%97%B4.tex}{L19-例-5}。
\begin{solve}
记 $A$ 的列向量为
    \[
        \alpha_1 = \begin{pmatrix}
            1 \\
            1 \\
            0
        \end{pmatrix}, \enspace
        \alpha_2 = \begin{pmatrix}
            1 \\
            0 \\
            1
        \end{pmatrix}, \enspace
        \alpha_3 = \begin{pmatrix}
            0 \\
            1 \\
            1
        \end{pmatrix}.
    \]
    对 $\alpha_1, \alpha_2, \alpha_3$ 施行 Gram-Schmidt 过程.

    先有
    \[
        u_1 = \alpha_1, \enspace q_1 = \frac{\alpha_1}{\lVert \alpha_1 \rVert}
        = \frac{1}{\sqrt{2}}
        \begin{pmatrix}
            1 \\
            1 \\
            0
        \end{pmatrix}.
    \]

    再看第二个向量，
    \[
        \langle \alpha_2, q_1 \rangle = \frac{1}{\sqrt{2}}, \enspace
        u_2 = \alpha_2 - \langle \alpha_2, q_1 \rangle q_1
        =
        \begin{pmatrix}
            \frac{1}{2} \\
            -\frac{1}{2} \\
            1
        \end{pmatrix},
    \]
    从而
    \[
        q_2 = \frac{u_2}{\lVert u_2 \rVert}
        = \frac{1}{\sqrt{6}}
        \begin{pmatrix}
            1 \\
            -1 \\
            2
        \end{pmatrix}.
    \]

    对第三个向量，
    \[
        \langle \alpha_3, q_1 \rangle = \frac{1}{\sqrt{2}}, \enspace
        \langle \alpha_3, q_2 \rangle = \frac{1}{\sqrt{6}},
    \]
    因而
    \[
        u_3 = \alpha_3 - \langle \alpha_3, q_1 \rangle q_1 - \langle \alpha_3, q_2 \rangle q_2
        = \frac{2}{3}
        \begin{pmatrix}
            -1 \\
            1 \\
            1
        \end{pmatrix},
    \]
    从而
    \[
        q_3 = \frac{u_3}{\lVert u_3 \rVert}
        = \frac{1}{\sqrt{3}}
        \begin{pmatrix}
            -1 \\
            1 \\
            1
        \end{pmatrix}.
    \]

    因而
    \[
        Q = \begin{pmatrix}
            \dfrac{1}{\sqrt{2}} & \dfrac{1}{\sqrt{6}}  & -\dfrac{1}{\sqrt{3}} \\
            \dfrac{1}{\sqrt{2}} & -\dfrac{1}{\sqrt{6}} & \dfrac{1}{\sqrt{3}}  \\
            0                   & \dfrac{2}{\sqrt{6}}  & \dfrac{1}{\sqrt{3}}
        \end{pmatrix}, \enspace
        R = \begin{pmatrix}
            \sqrt{2} & \dfrac{1}{\sqrt{2}} & \dfrac{1}{\sqrt{2}} \\
            0        & \dfrac{\sqrt{6}}{2} & \dfrac{1}{\sqrt{6}} \\
            0        & 0                   & \dfrac{2}{\sqrt{3}}
        \end{pmatrix}.
    \]
    这里
    \[
        r_{11} = \lVert \alpha_1 \rVert = \sqrt{2}, \enspace
        r_{22} = \lVert u_2 \rVert = \frac{\sqrt{6}}{2}, \enspace
        r_{33} = \lVert u_3 \rVert = \frac{2}{\sqrt{3}}.
    \]
    直接验证可得 $A = QR$，故这就是所求的 QR 分解.
\end{solve}

\begin{exercise}\label{prob:L19-例-7}
定义在 $ V = \mathbf{R}^3 $ 上的运算
    \[ \langle \vec{x}, \vec{y} \rangle_V = x_1 y_1 + x_2 y_2 + (x_2 + x_3)(y_2 + y_3) \]
    其中 $ \vec{x} = (x_1, x_2, x_3),\enspace \vec{y} = (y_1, y_2, y_3) $.
    \begin{enumerate}
        \item 验证 $ \langle \cdot, \cdot \rangle_V $ 是 $ \mathbf{R}^3 $ 上的一个内积；

        \item 求 $ \mathbf{R}^3 $ 在 $ \langle \cdot, \cdot \rangle_V $ 下的一组标准正交基；

        \item 求 $ \vec{\beta} \in V $ 使得 $ \forall \vec{x} \in V,\enspace x_1 + 2x_2 = \langle \vec{x}, \vec{\beta} \rangle_V $.
    \end{enumerate}
\end{exercise}
\noindent 来源：\href{https://github.com/yhwu-is/Linear-Algebra-Left-Undone/blob/new/%E8%AE%B2%E4%B9%89/%E4%B8%93%E9%A2%98/19%20%E5%86%85%E7%A7%AF%E7%A9%BA%E9%97%B4.tex}{L19-例-7}。
\begin{solve}
\begin{enumerate}
        \item 设 $\vec{x}, \vec{y}, \vec{z} \in \mathbf{R}^3$，$\lambda \in \mathbf{R}$. 先验证线性性与对称性：
        \begin{align*}
            \langle \vec{x} + \vec{y}, \vec{z} \rangle_V
             & = (x_1 + y_1)z_1 + (x_2 + y_2)z_2 + (x_2 + y_2 + x_3 + y_3)(z_2 + z_3) \\
             & = \langle \vec{x}, \vec{z} \rangle_V + \langle \vec{y}, \vec{z} \rangle_V, \\
            \langle \lambda \vec{x}, \vec{y} \rangle_V
             & = \lambda x_1 y_1 + \lambda x_2 y_2 + (\lambda x_2 + \lambda x_3)(y_2 + y_3)
            = \lambda \langle \vec{x}, \vec{y} \rangle_V.
        \end{align*}
        交换 $\vec{x}, \vec{y}$ 后定义式不变，因此 $\langle \vec{x}, \vec{y} \rangle_V = \langle \vec{y}, \vec{x} \rangle_V$.

        再看正定性：$\langle \vec{x}, \vec{x} \rangle_V = x_1^2 + x_2^2 + (x_2 + x_3)^2 \geqslant 0$. 若 $\langle \vec{x}, \vec{x} \rangle_V = 0$，则 $x_1^2 = 0$，$x_2^2 = 0$，$(x_2 + x_3)^2 = 0$，从而 $x_1 = 0$，$x_2 = 0$，$x_3 = 0$. 故 $\vec{x} = \vec{0}$.

        综上，$\langle \cdot, \cdot \rangle_V$ 是 $\mathbf{R}^3$ 上的一个内积.

        \item 取 $\mathbf{R}^3$ 的标准基 $v_1 = (1, 0, 0)^{\mathrm{T}}, \enspace v_2 = (0, 1, 0)^{\mathrm{T}}, \enspace v_3 = (0, 0, 1)^{\mathrm{T}}$.
        对其施行 Gram-Schmidt 过程（具体过程不展开），得到在内积 $\langle \cdot, \cdot \rangle_V$ 下的一组标准正交基可以取为
        \[
            \left(1, 0, 0\right)^{\mathrm{T}}, \enspace \left(0, \frac{1}{\sqrt{2}}, 0\right)^{\mathrm{T}}, \enspace \left(0, -\frac{1}{\sqrt{2}}, \sqrt{2}\right)^{\mathrm{T}}.
        \]

        \item 设 $\varphi \colon V \to \mathbf{R}$ 定义为
        \[
            \varphi(\vec{x}) = x_1 + 2x_2.
        \]
        由原讲义“Riesz”处，存在唯一的 $\vec{\beta} \in V$ 使得对任意 $\vec{x} \in V$ 都有
        \[
            \varphi(\vec{x}) = \langle \vec{x}, \vec{\beta} \rangle_V.
        \]
        利用上一步得到的标准正交基 $e_1, e_2, e_3$，可得
        \[
            \vec{\beta} = \varphi(e_1)e_1 + \varphi(e_2)e_2 + \varphi(e_3)e_3 = (1, 2, -2)^{\mathrm{T}}.
        \]
    \end{enumerate}
\end{solve}

\begin{exercise}\label{prob:S-QR-1}
对矩阵 $A=\begin{pmatrix}1&0&0\\0&1&0\\0&-4&1\end{pmatrix}$ 的列依次进行施密特正交化，给出正交基、标准正交基与 $A=QR$ 分解，并说明为何每一步均不会得到零向量。
\end{exercise}
\noindent 编者补充：S-QR-1；相关主题见原讲义对应讲，题目数据与解答由本选编另行构造。
\begin{solve}
设列为 $a_1,a_2,a_3$。按
\[u_j=a_j-\sum_{i<j}\frac{u_i^{\mathrm T}a_j}{u_i^{\mathrm T}u_i}u_i\]
依次计算，得到 $U=(u_1,u_2,u_3)=\begin{pmatrix}1&0&0\\0&1&\frac{4}{17}\\0&-4&\frac{1}{17}\end{pmatrix}$，且
\[U^{\mathrm T}U=\operatorname{diag}(1,17,\frac{1}{17}).\]
所以 $Q=\left(u_1/\sqrt{1},u_2/\sqrt{17},u_3/\sqrt{\frac{1}{17}}\right)$ 是标准正交矩阵，而
\[R=Q^{\mathrm T}A=\begin{pmatrix}\frac{1}{\sqrt{1}}&0&0\\0&\frac{17}{\sqrt{17}}&\frac{-4}{\sqrt{17}}\\0&0&\frac{\frac{1}{17}}{\sqrt{\frac{1}{17}}}\end{pmatrix}.\]
由构造可得 $A=QR$。这里投影系数的分母是未归一化向量的长度平方，不能省略。因为 $\det A=1$，原列向量线性无关；若某个 $u_j=0$，则 $a_j$ 属于前面的列的张成空间，矛盾。
\end{solve}

\begin{exercise}\label{prob:S-LS-1}
对数据点 $(-1,1),(0,2),(2,-1)$，用直线 $y=c_0+c_1t$ 作最小二乘拟合。写出设计矩阵、正规方程、拟合值与残差，并计算最小残差平方和。
\end{exercise}
\noindent 编者补充：S-LS-1；相关主题见原讲义对应讲，题目数据与解答由本选编另行构造。
\begin{solve}
设计矩阵与观测向量为 $A=\begin{pmatrix}1&-1\\1&0\\1&2\end{pmatrix}$、$b=\begin{pmatrix}1\\2\\-1\end{pmatrix}$。目标为 $\min_c\|Ac-b\|_2^2$，正规方程为
\[A^{\mathrm T}Ac=A^{\mathrm T}b,\qquad \begin{pmatrix}3&1\\1&5\end{pmatrix}c=\begin{pmatrix}2\\-3\end{pmatrix}.\]
因为两列线性无关，$A^{\mathrm T}A$ 正定且可逆，得到唯一解 $c=\begin{pmatrix}\frac{13}{14}\\-\frac{11}{14}\end{pmatrix}$。进一步，
\[\widehat b=Ac=\begin{pmatrix}\frac{12}{7}\\\frac{13}{14}\\-\frac{9}{14}\end{pmatrix},\qquad e=b-\widehat b=\begin{pmatrix}-\frac{5}{7}\\\frac{15}{14}\\-\frac{5}{14}\end{pmatrix},\qquad \|e\|_2^2=\frac{25}{14}.\]
直接检验 $A^{\mathrm T}e=0$，说明残差正交于列空间。对任意 $d$，有 $\|b-A(c+d)\|^2=\|e\|^2+\|Ad\|^2$，从而上述驻点确为全局最小点，并由满列秩得到唯一性。
\end{solve}

\begin{exercise}\label{prob:S-QR-2}
对矩阵 $A=\begin{pmatrix}1&1&-1\\0&1&0\\-2&-2&3\end{pmatrix}$ 的列依次进行施密特正交化，给出正交基、标准正交基与 $A=QR$ 分解，并说明为何每一步均不会得到零向量。
\end{exercise}
\noindent 编者补充：S-QR-2；相关主题见原讲义对应讲，题目数据与解答由本选编另行构造。
\begin{solve}
设列为 $a_1,a_2,a_3$。按
\[u_j=a_j-\sum_{i<j}\frac{u_i^{\mathrm T}a_j}{u_i^{\mathrm T}u_i}u_i\]
依次计算，得到 $U=(u_1,u_2,u_3)=\begin{pmatrix}1&0&\frac{2}{5}\\0&1&0\\-2&0&\frac{1}{5}\end{pmatrix}$，且
\[U^{\mathrm T}U=\operatorname{diag}(5,1,\frac{1}{5}).\]
所以 $Q=\left(u_1/\sqrt{5},u_2/\sqrt{1},u_3/\sqrt{\frac{1}{5}}\right)$ 是标准正交矩阵，而
\[R=Q^{\mathrm T}A=\begin{pmatrix}\frac{5}{\sqrt{5}}&\frac{5}{\sqrt{5}}&\frac{-7}{\sqrt{5}}\\0&\frac{1}{\sqrt{1}}&0\\0&0&\frac{\frac{1}{5}}{\sqrt{\frac{1}{5}}}\end{pmatrix}.\]
由构造可得 $A=QR$。这里投影系数的分母是未归一化向量的长度平方，不能省略。因为 $\det A=1$，原列向量线性无关；若某个 $u_j=0$，则 $a_j$ 属于前面的列的张成空间，矛盾。
\end{solve}

\begin{exercise}\label{prob:S-LS-2}
对数据点 $(-1,2),(0,2),(3,-2)$，用直线 $y=c_0+c_1t$ 作最小二乘拟合。写出设计矩阵、正规方程、拟合值与残差，并计算最小残差平方和。
\end{exercise}
\noindent 编者补充：S-LS-2；相关主题见原讲义对应讲，题目数据与解答由本选编另行构造。
\begin{solve}
设计矩阵与观测向量为 $A=\begin{pmatrix}1&-1\\1&0\\1&3\end{pmatrix}$、$b=\begin{pmatrix}2\\2\\-2\end{pmatrix}$。目标为 $\min_c\|Ac-b\|_2^2$，正规方程为
\[A^{\mathrm T}Ac=A^{\mathrm T}b,\qquad \begin{pmatrix}3&2\\2&10\end{pmatrix}c=\begin{pmatrix}2\\-8\end{pmatrix}.\]
因为两列线性无关，$A^{\mathrm T}A$ 正定且可逆，得到唯一解 $c=\begin{pmatrix}\frac{18}{13}\\-\frac{14}{13}\end{pmatrix}$。进一步，
\[\widehat b=Ac=\begin{pmatrix}\frac{32}{13}\\\frac{18}{13}\\-\frac{24}{13}\end{pmatrix},\qquad e=b-\widehat b=\begin{pmatrix}-\frac{6}{13}\\\frac{8}{13}\\-\frac{2}{13}\end{pmatrix},\qquad \|e\|_2^2=\frac{8}{13}.\]
直接检验 $A^{\mathrm T}e=0$，说明残差正交于列空间。对任意 $d$，有 $\|b-A(c+d)\|^2=\|e\|^2+\|Ad\|^2$，从而上述驻点确为全局最小点，并由满列秩得到唯一性。
\end{solve}

\begin{exercise}\label{prob:S-QR-3}
对矩阵 $A=\begin{pmatrix}1&0&0\\-4&1&0\\-4&2&1\end{pmatrix}$ 的列依次进行施密特正交化，给出正交基、标准正交基与 $A=QR$ 分解，并说明为何每一步均不会得到零向量。
\end{exercise}
\noindent 编者补充：S-QR-3；相关主题见原讲义对应讲，题目数据与解答由本选编另行构造。
\begin{solve}
设列为 $a_1,a_2,a_3$。按
\[u_j=a_j-\sum_{i<j}\frac{u_i^{\mathrm T}a_j}{u_i^{\mathrm T}u_i}u_i\]
依次计算，得到 $U=(u_1,u_2,u_3)=\begin{pmatrix}1&\frac{4}{11}&-\frac{4}{21}\\-4&-\frac{5}{11}&-\frac{2}{21}\\-4&\frac{6}{11}&\frac{1}{21}\end{pmatrix}$，且
\[U^{\mathrm T}U=\operatorname{diag}(33,\frac{7}{11},\frac{1}{21}).\]
所以 $Q=\left(u_1/\sqrt{33},u_2/\sqrt{\frac{7}{11}},u_3/\sqrt{\frac{1}{21}}\right)$ 是标准正交矩阵，而
\[R=Q^{\mathrm T}A=\begin{pmatrix}\frac{33}{\sqrt{33}}&\frac{-12}{\sqrt{33}}&\frac{-4}{\sqrt{33}}\\0&\frac{\frac{7}{11}}{\sqrt{\frac{7}{11}}}&\frac{\frac{6}{11}}{\sqrt{\frac{7}{11}}}\\0&0&\frac{\frac{1}{21}}{\sqrt{\frac{1}{21}}}\end{pmatrix}.\]
由构造可得 $A=QR$。这里投影系数的分母是未归一化向量的长度平方，不能省略。因为 $\det A=1$，原列向量线性无关；若某个 $u_j=0$，则 $a_j$ 属于前面的列的张成空间，矛盾。
\end{solve}

\begin{exercise}\label{prob:S-LS-3}
对数据点 $(-1,3),(0,2),(4,-3)$，用直线 $y=c_0+c_1t$ 作最小二乘拟合。写出设计矩阵、正规方程、拟合值与残差，并计算最小残差平方和。
\end{exercise}
\noindent 编者补充：S-LS-3；相关主题见原讲义对应讲，题目数据与解答由本选编另行构造。
\begin{solve}
设计矩阵与观测向量为 $A=\begin{pmatrix}1&-1\\1&0\\1&4\end{pmatrix}$、$b=\begin{pmatrix}3\\2\\-3\end{pmatrix}$。目标为 $\min_c\|Ac-b\|_2^2$，正规方程为
\[A^{\mathrm T}Ac=A^{\mathrm T}b,\qquad \begin{pmatrix}3&3\\3&17\end{pmatrix}c=\begin{pmatrix}2\\-15\end{pmatrix}.\]
因为两列线性无关，$A^{\mathrm T}A$ 正定且可逆，得到唯一解 $c=\begin{pmatrix}\frac{79}{42}\\-\frac{17}{14}\end{pmatrix}$。进一步，
\[\widehat b=Ac=\begin{pmatrix}\frac{65}{21}\\\frac{79}{42}\\-\frac{125}{42}\end{pmatrix},\qquad e=b-\widehat b=\begin{pmatrix}-\frac{2}{21}\\\frac{5}{42}\\-\frac{1}{42}\end{pmatrix},\qquad \|e\|_2^2=\frac{1}{42}.\]
直接检验 $A^{\mathrm T}e=0$，说明残差正交于列空间。对任意 $d$，有 $\|b-A(c+d)\|^2=\|e\|^2+\|Ad\|^2$，从而上述驻点确为全局最小点，并由满列秩得到唯一性。
\end{solve}

\begin{exercise}\label{prob:S-QR-4}
对矩阵 $A=\begin{pmatrix}1&0&-3\\-1&1&0\\0&0&1\end{pmatrix}$ 的列依次进行施密特正交化，给出正交基、标准正交基与 $A=QR$ 分解，并说明为何每一步均不会得到零向量。
\end{exercise}
\noindent 编者补充：S-QR-4；相关主题见原讲义对应讲，题目数据与解答由本选编另行构造。
\begin{solve}
设列为 $a_1,a_2,a_3$。按
\[u_j=a_j-\sum_{i<j}\frac{u_i^{\mathrm T}a_j}{u_i^{\mathrm T}u_i}u_i\]
依次计算，得到 $U=(u_1,u_2,u_3)=\begin{pmatrix}1&\frac{1}{2}&0\\-1&\frac{1}{2}&0\\0&0&1\end{pmatrix}$，且
\[U^{\mathrm T}U=\operatorname{diag}(2,\frac{1}{2},1).\]
所以 $Q=\left(u_1/\sqrt{2},u_2/\sqrt{\frac{1}{2}},u_3/\sqrt{1}\right)$ 是标准正交矩阵，而
\[R=Q^{\mathrm T}A=\begin{pmatrix}\frac{2}{\sqrt{2}}&\frac{-1}{\sqrt{2}}&\frac{-3}{\sqrt{2}}\\0&\frac{\frac{1}{2}}{\sqrt{\frac{1}{2}}}&\frac{-\frac{3}{2}}{\sqrt{\frac{1}{2}}}\\0&0&\frac{1}{\sqrt{1}}\end{pmatrix}.\]
由构造可得 $A=QR$。这里投影系数的分母是未归一化向量的长度平方，不能省略。因为 $\det A=1$，原列向量线性无关；若某个 $u_j=0$，则 $a_j$ 属于前面的列的张成空间，矛盾。
\end{solve}

\begin{exercise}\label{prob:S-LS-4}
对数据点 $(-1,4),(0,2),(5,-4)$，用直线 $y=c_0+c_1t$ 作最小二乘拟合。写出设计矩阵、正规方程、拟合值与残差，并计算最小残差平方和。
\end{exercise}
\noindent 编者补充：S-LS-4；相关主题见原讲义对应讲，题目数据与解答由本选编另行构造。
\begin{solve}
设计矩阵与观测向量为 $A=\begin{pmatrix}1&-1\\1&0\\1&5\end{pmatrix}$、$b=\begin{pmatrix}4\\2\\-4\end{pmatrix}$。目标为 $\min_c\|Ac-b\|_2^2$，正规方程为
\[A^{\mathrm T}Ac=A^{\mathrm T}b,\qquad \begin{pmatrix}3&4\\4&26\end{pmatrix}c=\begin{pmatrix}2\\-24\end{pmatrix}.\]
因为两列线性无关，$A^{\mathrm T}A$ 正定且可逆，得到唯一解 $c=\begin{pmatrix}\frac{74}{31}\\-\frac{40}{31}\end{pmatrix}$。进一步，
\[\widehat b=Ac=\begin{pmatrix}\frac{114}{31}\\\frac{74}{31}\\-\frac{126}{31}\end{pmatrix},\qquad e=b-\widehat b=\begin{pmatrix}\frac{10}{31}\\-\frac{12}{31}\\\frac{2}{31}\end{pmatrix},\qquad \|e\|_2^2=\frac{8}{31}.\]
直接检验 $A^{\mathrm T}e=0$，说明残差正交于列空间。对任意 $d$，有 $\|b-A(c+d)\|^2=\|e\|^2+\|Ad\|^2$，从而上述驻点确为全局最小点，并由满列秩得到唯一性。
\end{solve}

\begin{exercise}\label{prob:S-QR-5}
对矩阵 $A=\begin{pmatrix}1&-2&0\\-2&5&2\\-1&2&1\end{pmatrix}$ 的列依次进行施密特正交化，给出正交基、标准正交基与 $A=QR$ 分解，并说明为何每一步均不会得到零向量。
\end{exercise}
\noindent 编者补充：S-QR-5；相关主题见原讲义对应讲，题目数据与解答由本选编另行构造。
\begin{solve}
设列为 $a_1,a_2,a_3$。按
\[u_j=a_j-\sum_{i<j}\frac{u_i^{\mathrm T}a_j}{u_i^{\mathrm T}u_i}u_i\]
依次计算，得到 $U=(u_1,u_2,u_3)=\begin{pmatrix}1&\frac{1}{3}&\frac{1}{2}\\-2&\frac{1}{3}&0\\-1&-\frac{1}{3}&\frac{1}{2}\end{pmatrix}$，且
\[U^{\mathrm T}U=\operatorname{diag}(6,\frac{1}{3},\frac{1}{2}).\]
所以 $Q=\left(u_1/\sqrt{6},u_2/\sqrt{\frac{1}{3}},u_3/\sqrt{\frac{1}{2}}\right)$ 是标准正交矩阵，而
\[R=Q^{\mathrm T}A=\begin{pmatrix}\frac{6}{\sqrt{6}}&\frac{-14}{\sqrt{6}}&\frac{-5}{\sqrt{6}}\\0&\frac{\frac{1}{3}}{\sqrt{\frac{1}{3}}}&\frac{\frac{1}{3}}{\sqrt{\frac{1}{3}}}\\0&0&\frac{\frac{1}{2}}{\sqrt{\frac{1}{2}}}\end{pmatrix}.\]
由构造可得 $A=QR$。这里投影系数的分母是未归一化向量的长度平方，不能省略。因为 $\det A=1$，原列向量线性无关；若某个 $u_j=0$，则 $a_j$ 属于前面的列的张成空间，矛盾。
\end{solve}

\begin{exercise}\label{prob:S-LS-5}
对数据点 $(-1,5),(0,2),(6,-5)$，用直线 $y=c_0+c_1t$ 作最小二乘拟合。写出设计矩阵、正规方程、拟合值与残差，并计算最小残差平方和。
\end{exercise}
\noindent 编者补充：S-LS-5；相关主题见原讲义对应讲，题目数据与解答由本选编另行构造。
\begin{solve}
设计矩阵与观测向量为 $A=\begin{pmatrix}1&-1\\1&0\\1&6\end{pmatrix}$、$b=\begin{pmatrix}5\\2\\-5\end{pmatrix}$。目标为 $\min_c\|Ac-b\|_2^2$，正规方程为
\[A^{\mathrm T}Ac=A^{\mathrm T}b,\qquad \begin{pmatrix}3&5\\5&37\end{pmatrix}c=\begin{pmatrix}2\\-35\end{pmatrix}.\]
因为两列线性无关，$A^{\mathrm T}A$ 正定且可逆，得到唯一解 $c=\begin{pmatrix}\frac{249}{86}\\-\frac{115}{86}\end{pmatrix}$。进一步，
\[\widehat b=Ac=\begin{pmatrix}\frac{182}{43}\\\frac{249}{86}\\-\frac{441}{86}\end{pmatrix},\qquad e=b-\widehat b=\begin{pmatrix}\frac{33}{43}\\-\frac{77}{86}\\\frac{11}{86}\end{pmatrix},\qquad \|e\|_2^2=\frac{121}{86}.\]
直接检验 $A^{\mathrm T}e=0$，说明残差正交于列空间。对任意 $d$，有 $\|b-A(c+d)\|^2=\|e\|^2+\|Ad\|^2$，从而上述驻点确为全局最小点，并由满列秩得到唯一性。
\end{solve}

\begin{exercise}\label{prob:S-QR-6}
对矩阵 $A=\begin{pmatrix}2&-3&0\\-1&2&0\\0&0&1\end{pmatrix}$ 的列依次进行施密特正交化，给出正交基、标准正交基与 $A=QR$ 分解，并说明为何每一步均不会得到零向量。
\end{exercise}
\noindent 编者补充：S-QR-6；相关主题见原讲义对应讲，题目数据与解答由本选编另行构造。
\begin{solve}
设列为 $a_1,a_2,a_3$。按
\[u_j=a_j-\sum_{i<j}\frac{u_i^{\mathrm T}a_j}{u_i^{\mathrm T}u_i}u_i\]
依次计算，得到 $U=(u_1,u_2,u_3)=\begin{pmatrix}2&\frac{1}{5}&0\\-1&\frac{2}{5}&0\\0&0&1\end{pmatrix}$，且
\[U^{\mathrm T}U=\operatorname{diag}(5,\frac{1}{5},1).\]
所以 $Q=\left(u_1/\sqrt{5},u_2/\sqrt{\frac{1}{5}},u_3/\sqrt{1}\right)$ 是标准正交矩阵，而
\[R=Q^{\mathrm T}A=\begin{pmatrix}\frac{5}{\sqrt{5}}&\frac{-8}{\sqrt{5}}&0\\0&\frac{\frac{1}{5}}{\sqrt{\frac{1}{5}}}&0\\0&0&\frac{1}{\sqrt{1}}\end{pmatrix}.\]
由构造可得 $A=QR$。这里投影系数的分母是未归一化向量的长度平方，不能省略。因为 $\det A=1$，原列向量线性无关；若某个 $u_j=0$，则 $a_j$ 属于前面的列的张成空间，矛盾。
\end{solve}

\begin{exercise}\label{prob:S-LS-6}
对数据点 $(-1,6),(0,2),(7,-6)$，用直线 $y=c_0+c_1t$ 作最小二乘拟合。写出设计矩阵、正规方程、拟合值与残差，并计算最小残差平方和。
\end{exercise}
\noindent 编者补充：S-LS-6；相关主题见原讲义对应讲，题目数据与解答由本选编另行构造。
\begin{solve}
设计矩阵与观测向量为 $A=\begin{pmatrix}1&-1\\1&0\\1&7\end{pmatrix}$、$b=\begin{pmatrix}6\\2\\-6\end{pmatrix}$。目标为 $\min_c\|Ac-b\|_2^2$，正规方程为
\[A^{\mathrm T}Ac=A^{\mathrm T}b,\qquad \begin{pmatrix}3&6\\6&50\end{pmatrix}c=\begin{pmatrix}2\\-48\end{pmatrix}.\]
因为两列线性无关，$A^{\mathrm T}A$ 正定且可逆，得到唯一解 $c=\begin{pmatrix}\frac{194}{57}\\-\frac{26}{19}\end{pmatrix}$。进一步，
\[\widehat b=Ac=\begin{pmatrix}\frac{272}{57}\\\frac{194}{57}\\-\frac{352}{57}\end{pmatrix},\qquad e=b-\widehat b=\begin{pmatrix}\frac{70}{57}\\-\frac{80}{57}\\\frac{10}{57}\end{pmatrix},\qquad \|e\|_2^2=\frac{200}{57}.\]
直接检验 $A^{\mathrm T}e=0$，说明残差正交于列空间。对任意 $d$，有 $\|b-A(c+d)\|^2=\|e\|^2+\|Ad\|^2$，从而上述驻点确为全局最小点，并由满列秩得到唯一性。
\end{solve}

\section{内积空间上的算子}\label{sec:topic20}
本节收录 13 道题，其中 12 道为编者补充计算题。区分伴随、自伴、正交或酉以及正规算子；谱分解和奇异值分解使用的假设与基的选择各不相同。
\begin{exercise}\label{prob:L20-例-8}
求矩阵
    \[
        A = \begin{pmatrix}
            1 & 1  \\
            1 & 1  \\
            1 & -1
        \end{pmatrix}
    \]
    的奇异值分解.
\end{exercise}
\noindent 来源：\href{https://github.com/yhwu-is/Linear-Algebra-Left-Undone/blob/new/%E8%AE%B2%E4%B9%89/%E4%B8%93%E9%A2%98/20%20%E5%86%85%E7%A7%AF%E7%A9%BA%E9%97%B4%E4%B8%8A%E7%9A%84%E7%AE%97%E5%AD%90.tex}{L20-例-8}。
\begin{solve}
我们先总结一套求矩阵奇异值分解的通用步骤，然后在本题中逐步执行.

    \begin{enumerate}
        \item 先计算 $A^{\mathrm{T}}A$（或复数域中的 $A^*A$），求出其全部特征值与一组标准正交特征向量.

        \item 将这些特征值开平方，得到奇异值，并按从大到小排序.

        \item 与各个正奇异值 $s_i$ 对应的单位特征向量 $v_i$ 就是右奇异向量，再由公式
              \[
                  u_i = \frac{1}{s_i}Av_i
              \]
              求出左奇异向量.

        \item 若 $U$ 或 $V$ 的列向量还不够构成一组标准正交基，就将其补齐.

        \item 以左奇异向量为列组成 $U$，以右奇异向量为列组成 $V$，再把奇异值按顺序填入 $\Sigma$ 的主对角线，即得
              \[
                  A = U\Sigma V^{\mathrm{T}}.
              \]
    \end{enumerate}

    下面按此步骤计算.

    \begin{enumerate}
        \item 先求 $A^{\mathrm{T}}A$ 的特征值与特征向量，不难得到其特征值为 $4$ 和 $2$，对应的特征向量分别为 $(1, 1)^{\mathrm{T}}$ 和 $(1, -1)^{\mathrm{T}}$. 归一化后得到一组标准正交特征向量
              \[
                  v_1 = \frac{1}{\sqrt{2}}\begin{pmatrix}
                      1 \\
                      1
                  \end{pmatrix},\qquad
                  v_2 = \frac{1}{\sqrt{2}}\begin{pmatrix}
                      1  \\
                      -1
                  \end{pmatrix}.
              \]

        \item 因此 $A$ 的奇异值为 $s_1 = \sqrt{4} = 2, s_2 = \sqrt{2}$.

        \item 由 $u_i = \dfrac{1}{s_i}Av_i$ 求出左奇异向量为
              \[ u_1 = \frac{1}{2}\begin{pmatrix}
                      1 \\
                      1 \\
                      0
                  \end{pmatrix}, \qquad u_2 = \frac{1}{\sqrt{2}}\begin{pmatrix}
                      0 \\
                      0 \\
                      1
                  \end{pmatrix}. \]

        \item 由于 $A$ 是 $3 \times 2$ 矩阵，$U$ 还需要第三个列向量. 我们只需补一个与 $u_1, u_2$ 都正交的单位向量即可，例如取
              \[
                  u_3 = \frac{1}{\sqrt{2}}
                  \begin{pmatrix}
                      1  \\
                      -1 \\
                      0
                  \end{pmatrix}.
              \]
              于是
              \[
                  U =
                  \begin{pmatrix}
                      \frac{1}{\sqrt{2}} & 0 & \frac{1}{\sqrt{2}}  \\
                      \frac{1}{\sqrt{2}} & 0 & -\frac{1}{\sqrt{2}} \\
                      0                  & 1 & 0
                  \end{pmatrix},\qquad
                  V =
                  \frac{1}{\sqrt{2}}
                  \begin{pmatrix}
                      1 & 1  \\
                      1 & -1
                  \end{pmatrix}.
              \]

        \item 最后构造
              \[
                  \Sigma =
                  \begin{pmatrix}
                      2 & 0        \\
                      0 & \sqrt{2} \\
                      0 & 0
                  \end{pmatrix}.
              \]
              因而得到 $A = U\Sigma V^{\mathrm{T}}$，至此，奇异值分解求得.
    \end{enumerate}
\end{solve}

\begin{exercise}\label{prob:S-SPECTRAL-1}
设实对称矩阵 $A=\begin{pmatrix}\frac{23}{25}&-\frac{36}{25}&0\\-\frac{36}{25}&\frac{2}{25}&0\\0&0&4\end{pmatrix}$。求正交对角化，判断其正定性，求二次型 $x^{\mathrm T}Ax$ 在 $x^{\mathrm T}x=1$ 下的最大值、最小值及取到这些值的向量。
\end{exercise}
\noindent 编者补充：S-SPECTRAL-1；相关主题见原讲义对应讲，题目数据与解答由本选编另行构造。
\begin{solve}
取正交矩阵 $Q=\begin{pmatrix}\frac{3}{5}&-\frac{4}{5}&0\\\frac{4}{5}&\frac{3}{5}&0\\0&0&1\end{pmatrix}$，相乘可验证
\[Q^{\mathrm T}AQ=\operatorname{diag}(-1,2,4).\]
由于既有正特征值又有负特征值，矩阵不定，正、负惯性指数分别为 $2,1$。令 $y=Q^{\mathrm T}x$，正交变换保持单位球面，因此
\[x^{\mathrm T}Ax=(-1)y_1^2+2y_2^2+4y_3^2,\qquad y_1^2+y_2^2+y_3^2=1.\]
这是三个特征值的凸组合，最小值为 $-1$，最大值为 $4$。分别在 $x=\pm Qe_1$ 与 $x=\pm Qe_3$ 处取到；因为这些特征值互异，等号条件没有额外方向。
\end{solve}

\begin{exercise}\label{prob:S-SVD-1}
对 $A=\begin{pmatrix}\frac{36}{25}&\frac{48}{25}\\\frac{48}{25}&\frac{64}{25}\\-\frac{8}{5}&\frac{6}{5}\end{pmatrix}$ 求一个紧致奇异值分解，给出秩、两个奇异值、谱范数及 Moore--Penrose 广义逆，并验证 $A^+A=I_2$。
\end{exercise}
\noindent 编者补充：S-SVD-1；相关主题见原讲义对应讲，题目数据与解答由本选编另行构造。
\begin{solve}
取
\[U=\begin{pmatrix}\frac{3}{5}&0\\\frac{4}{5}&0\\0&1\end{pmatrix},\quad\Sigma=\operatorname{diag}(4,2),\quad V=\begin{pmatrix}\frac{3}{5}&-\frac{4}{5}\\\frac{4}{5}&\frac{3}{5}\end{pmatrix}.\]
直接核对 $U^{\mathrm T}U=I_2$、$V^{\mathrm T}V=I_2$ 和 $A=U\Sigma V^{\mathrm T}$，所以这就是紧致奇异值分解。两个奇异值均为正，故秩为 $2$；谱范数为较大的奇异值 $4$。广义逆为
\[A^+=V\Sigma^{-1}U^{\mathrm T}=\begin{pmatrix}\frac{9}{100}&\frac{3}{25}&-\frac{2}{5}\\\frac{3}{25}&\frac{4}{25}&\frac{3}{10}\end{pmatrix}.\]
于是 $A^+A=V\Sigma^{-1}\Sigma V^{\mathrm T}=I_2$，而 $AA^+=UU^{\mathrm T}$ 是到 $A$ 列空间的正交投影。后者通常不等于 $I_3$，这与矩阵只有两列相符。
\end{solve}

\begin{exercise}\label{prob:S-SPECTRAL-2}
设实对称矩阵 $A=\begin{pmatrix}\frac{6}{5}&-\frac{12}{5}&0\\-\frac{12}{5}&-\frac{1}{5}&0\\0&0&5\end{pmatrix}$。求正交对角化，判断其正定性，求二次型 $x^{\mathrm T}Ax$ 在 $x^{\mathrm T}x=1$ 下的最大值、最小值及取到这些值的向量。
\end{exercise}
\noindent 编者补充：S-SPECTRAL-2；相关主题见原讲义对应讲，题目数据与解答由本选编另行构造。
\begin{solve}
取正交矩阵 $Q=\begin{pmatrix}\frac{3}{5}&-\frac{4}{5}&0\\\frac{4}{5}&\frac{3}{5}&0\\0&0&1\end{pmatrix}$，相乘可验证
\[Q^{\mathrm T}AQ=\operatorname{diag}(-2,3,5).\]
由于既有正特征值又有负特征值，矩阵不定，正、负惯性指数分别为 $2,1$。令 $y=Q^{\mathrm T}x$，正交变换保持单位球面，因此
\[x^{\mathrm T}Ax=(-2)y_1^2+3y_2^2+5y_3^2,\qquad y_1^2+y_2^2+y_3^2=1.\]
这是三个特征值的凸组合，最小值为 $-2$，最大值为 $5$。分别在 $x=\pm Qe_1$ 与 $x=\pm Qe_3$ 处取到；因为这些特征值互异，等号条件没有额外方向。
\end{solve}

\begin{exercise}\label{prob:S-SVD-2}
对 $A=\begin{pmatrix}\frac{9}{5}&\frac{12}{5}\\\frac{12}{5}&\frac{16}{5}\\-\frac{12}{5}&\frac{9}{5}\end{pmatrix}$ 求一个紧致奇异值分解，给出秩、两个奇异值、谱范数及 Moore--Penrose 广义逆，并验证 $A^+A=I_2$。
\end{exercise}
\noindent 编者补充：S-SVD-2；相关主题见原讲义对应讲，题目数据与解答由本选编另行构造。
\begin{solve}
取
\[U=\begin{pmatrix}\frac{3}{5}&0\\\frac{4}{5}&0\\0&1\end{pmatrix},\quad\Sigma=\operatorname{diag}(5,3),\quad V=\begin{pmatrix}\frac{3}{5}&-\frac{4}{5}\\\frac{4}{5}&\frac{3}{5}\end{pmatrix}.\]
直接核对 $U^{\mathrm T}U=I_2$、$V^{\mathrm T}V=I_2$ 和 $A=U\Sigma V^{\mathrm T}$，所以这就是紧致奇异值分解。两个奇异值均为正，故秩为 $2$；谱范数为较大的奇异值 $5$。广义逆为
\[A^+=V\Sigma^{-1}U^{\mathrm T}=\begin{pmatrix}\frac{9}{125}&\frac{12}{125}&-\frac{4}{15}\\\frac{12}{125}&\frac{16}{125}&\frac{1}{5}\end{pmatrix}.\]
于是 $A^+A=V\Sigma^{-1}\Sigma V^{\mathrm T}=I_2$，而 $AA^+=UU^{\mathrm T}$ 是到 $A$ 列空间的正交投影。后者通常不等于 $I_3$，这与矩阵只有两列相符。
\end{solve}

\begin{exercise}\label{prob:S-SPECTRAL-3}
设实对称矩阵 $A=\begin{pmatrix}\frac{37}{25}&-\frac{84}{25}&0\\-\frac{84}{25}&-\frac{12}{25}&0\\0&0&6\end{pmatrix}$。求正交对角化，判断其正定性，求二次型 $x^{\mathrm T}Ax$ 在 $x^{\mathrm T}x=1$ 下的最大值、最小值及取到这些值的向量。
\end{exercise}
\noindent 编者补充：S-SPECTRAL-3；相关主题见原讲义对应讲，题目数据与解答由本选编另行构造。
\begin{solve}
取正交矩阵 $Q=\begin{pmatrix}\frac{3}{5}&-\frac{4}{5}&0\\\frac{4}{5}&\frac{3}{5}&0\\0&0&1\end{pmatrix}$，相乘可验证
\[Q^{\mathrm T}AQ=\operatorname{diag}(-3,4,6).\]
由于既有正特征值又有负特征值，矩阵不定，正、负惯性指数分别为 $2,1$。令 $y=Q^{\mathrm T}x$，正交变换保持单位球面，因此
\[x^{\mathrm T}Ax=(-3)y_1^2+4y_2^2+6y_3^2,\qquad y_1^2+y_2^2+y_3^2=1.\]
这是三个特征值的凸组合，最小值为 $-3$，最大值为 $6$。分别在 $x=\pm Qe_1$ 与 $x=\pm Qe_3$ 处取到；因为这些特征值互异，等号条件没有额外方向。
\end{solve}

\begin{exercise}\label{prob:S-SVD-3}
对 $A=\begin{pmatrix}\frac{54}{25}&\frac{72}{25}\\\frac{72}{25}&\frac{96}{25}\\-\frac{16}{5}&\frac{12}{5}\end{pmatrix}$ 求一个紧致奇异值分解，给出秩、两个奇异值、谱范数及 Moore--Penrose 广义逆，并验证 $A^+A=I_2$。
\end{exercise}
\noindent 编者补充：S-SVD-3；相关主题见原讲义对应讲，题目数据与解答由本选编另行构造。
\begin{solve}
取
\[U=\begin{pmatrix}\frac{3}{5}&0\\\frac{4}{5}&0\\0&1\end{pmatrix},\quad\Sigma=\operatorname{diag}(6,4),\quad V=\begin{pmatrix}\frac{3}{5}&-\frac{4}{5}\\\frac{4}{5}&\frac{3}{5}\end{pmatrix}.\]
直接核对 $U^{\mathrm T}U=I_2$、$V^{\mathrm T}V=I_2$ 和 $A=U\Sigma V^{\mathrm T}$，所以这就是紧致奇异值分解。两个奇异值均为正，故秩为 $2$；谱范数为较大的奇异值 $6$。广义逆为
\[A^+=V\Sigma^{-1}U^{\mathrm T}=\begin{pmatrix}\frac{3}{50}&\frac{2}{25}&-\frac{1}{5}\\\frac{2}{25}&\frac{8}{75}&\frac{3}{20}\end{pmatrix}.\]
于是 $A^+A=V\Sigma^{-1}\Sigma V^{\mathrm T}=I_2$，而 $AA^+=UU^{\mathrm T}$ 是到 $A$ 列空间的正交投影。后者通常不等于 $I_3$，这与矩阵只有两列相符。
\end{solve}

\begin{exercise}\label{prob:S-SPECTRAL-4}
设实对称矩阵 $A=\begin{pmatrix}\frac{44}{25}&-\frac{108}{25}&0\\-\frac{108}{25}&-\frac{19}{25}&0\\0&0&7\end{pmatrix}$。求正交对角化，判断其正定性，求二次型 $x^{\mathrm T}Ax$ 在 $x^{\mathrm T}x=1$ 下的最大值、最小值及取到这些值的向量。
\end{exercise}
\noindent 编者补充：S-SPECTRAL-4；相关主题见原讲义对应讲，题目数据与解答由本选编另行构造。
\begin{solve}
取正交矩阵 $Q=\begin{pmatrix}\frac{3}{5}&-\frac{4}{5}&0\\\frac{4}{5}&\frac{3}{5}&0\\0&0&1\end{pmatrix}$，相乘可验证
\[Q^{\mathrm T}AQ=\operatorname{diag}(-4,5,7).\]
由于既有正特征值又有负特征值，矩阵不定，正、负惯性指数分别为 $2,1$。令 $y=Q^{\mathrm T}x$，正交变换保持单位球面，因此
\[x^{\mathrm T}Ax=(-4)y_1^2+5y_2^2+7y_3^2,\qquad y_1^2+y_2^2+y_3^2=1.\]
这是三个特征值的凸组合，最小值为 $-4$，最大值为 $7$。分别在 $x=\pm Qe_1$ 与 $x=\pm Qe_3$ 处取到；因为这些特征值互异，等号条件没有额外方向。
\end{solve}

\begin{exercise}\label{prob:S-SVD-4}
对 $A=\begin{pmatrix}\frac{63}{25}&\frac{84}{25}\\\frac{84}{25}&\frac{112}{25}\\-4&3\end{pmatrix}$ 求一个紧致奇异值分解，给出秩、两个奇异值、谱范数及 Moore--Penrose 广义逆，并验证 $A^+A=I_2$。
\end{exercise}
\noindent 编者补充：S-SVD-4；相关主题见原讲义对应讲，题目数据与解答由本选编另行构造。
\begin{solve}
取
\[U=\begin{pmatrix}\frac{3}{5}&0\\\frac{4}{5}&0\\0&1\end{pmatrix},\quad\Sigma=\operatorname{diag}(7,5),\quad V=\begin{pmatrix}\frac{3}{5}&-\frac{4}{5}\\\frac{4}{5}&\frac{3}{5}\end{pmatrix}.\]
直接核对 $U^{\mathrm T}U=I_2$、$V^{\mathrm T}V=I_2$ 和 $A=U\Sigma V^{\mathrm T}$，所以这就是紧致奇异值分解。两个奇异值均为正，故秩为 $2$；谱范数为较大的奇异值 $7$。广义逆为
\[A^+=V\Sigma^{-1}U^{\mathrm T}=\begin{pmatrix}\frac{9}{175}&\frac{12}{175}&-\frac{4}{25}\\\frac{12}{175}&\frac{16}{175}&\frac{3}{25}\end{pmatrix}.\]
于是 $A^+A=V\Sigma^{-1}\Sigma V^{\mathrm T}=I_2$，而 $AA^+=UU^{\mathrm T}$ 是到 $A$ 列空间的正交投影。后者通常不等于 $I_3$，这与矩阵只有两列相符。
\end{solve}

\begin{exercise}\label{prob:S-SPECTRAL-5}
设实对称矩阵 $A=\begin{pmatrix}\frac{51}{25}&-\frac{132}{25}&0\\-\frac{132}{25}&-\frac{26}{25}&0\\0&0&8\end{pmatrix}$。求正交对角化，判断其正定性，求二次型 $x^{\mathrm T}Ax$ 在 $x^{\mathrm T}x=1$ 下的最大值、最小值及取到这些值的向量。
\end{exercise}
\noindent 编者补充：S-SPECTRAL-5；相关主题见原讲义对应讲，题目数据与解答由本选编另行构造。
\begin{solve}
取正交矩阵 $Q=\begin{pmatrix}\frac{3}{5}&-\frac{4}{5}&0\\\frac{4}{5}&\frac{3}{5}&0\\0&0&1\end{pmatrix}$，相乘可验证
\[Q^{\mathrm T}AQ=\operatorname{diag}(-5,6,8).\]
由于既有正特征值又有负特征值，矩阵不定，正、负惯性指数分别为 $2,1$。令 $y=Q^{\mathrm T}x$，正交变换保持单位球面，因此
\[x^{\mathrm T}Ax=(-5)y_1^2+6y_2^2+8y_3^2,\qquad y_1^2+y_2^2+y_3^2=1.\]
这是三个特征值的凸组合，最小值为 $-5$，最大值为 $8$。分别在 $x=\pm Qe_1$ 与 $x=\pm Qe_3$ 处取到；因为这些特征值互异，等号条件没有额外方向。
\end{solve}

\begin{exercise}\label{prob:S-SVD-5}
对 $A=\begin{pmatrix}\frac{72}{25}&\frac{96}{25}\\\frac{96}{25}&\frac{128}{25}\\-\frac{24}{5}&\frac{18}{5}\end{pmatrix}$ 求一个紧致奇异值分解，给出秩、两个奇异值、谱范数及 Moore--Penrose 广义逆，并验证 $A^+A=I_2$。
\end{exercise}
\noindent 编者补充：S-SVD-5；相关主题见原讲义对应讲，题目数据与解答由本选编另行构造。
\begin{solve}
取
\[U=\begin{pmatrix}\frac{3}{5}&0\\\frac{4}{5}&0\\0&1\end{pmatrix},\quad\Sigma=\operatorname{diag}(8,6),\quad V=\begin{pmatrix}\frac{3}{5}&-\frac{4}{5}\\\frac{4}{5}&\frac{3}{5}\end{pmatrix}.\]
直接核对 $U^{\mathrm T}U=I_2$、$V^{\mathrm T}V=I_2$ 和 $A=U\Sigma V^{\mathrm T}$，所以这就是紧致奇异值分解。两个奇异值均为正，故秩为 $2$；谱范数为较大的奇异值 $8$。广义逆为
\[A^+=V\Sigma^{-1}U^{\mathrm T}=\begin{pmatrix}\frac{9}{200}&\frac{3}{50}&-\frac{2}{15}\\\frac{3}{50}&\frac{2}{25}&\frac{1}{10}\end{pmatrix}.\]
于是 $A^+A=V\Sigma^{-1}\Sigma V^{\mathrm T}=I_2$，而 $AA^+=UU^{\mathrm T}$ 是到 $A$ 列空间的正交投影。后者通常不等于 $I_3$，这与矩阵只有两列相符。
\end{solve}

\begin{exercise}\label{prob:S-SPECTRAL-6}
设实对称矩阵 $A=\begin{pmatrix}\frac{58}{25}&-\frac{156}{25}&0\\-\frac{156}{25}&-\frac{33}{25}&0\\0&0&9\end{pmatrix}$。求正交对角化，判断其正定性，求二次型 $x^{\mathrm T}Ax$ 在 $x^{\mathrm T}x=1$ 下的最大值、最小值及取到这些值的向量。
\end{exercise}
\noindent 编者补充：S-SPECTRAL-6；相关主题见原讲义对应讲，题目数据与解答由本选编另行构造。
\begin{solve}
取正交矩阵 $Q=\begin{pmatrix}\frac{3}{5}&-\frac{4}{5}&0\\\frac{4}{5}&\frac{3}{5}&0\\0&0&1\end{pmatrix}$，相乘可验证
\[Q^{\mathrm T}AQ=\operatorname{diag}(-6,7,9).\]
由于既有正特征值又有负特征值，矩阵不定，正、负惯性指数分别为 $2,1$。令 $y=Q^{\mathrm T}x$，正交变换保持单位球面，因此
\[x^{\mathrm T}Ax=(-6)y_1^2+7y_2^2+9y_3^2,\qquad y_1^2+y_2^2+y_3^2=1.\]
这是三个特征值的凸组合，最小值为 $-6$，最大值为 $9$。分别在 $x=\pm Qe_1$ 与 $x=\pm Qe_3$ 处取到；因为这些特征值互异，等号条件没有额外方向。
\end{solve}

\begin{exercise}\label{prob:S-SVD-6}
对 $A=\begin{pmatrix}\frac{81}{25}&\frac{108}{25}\\\frac{108}{25}&\frac{144}{25}\\-\frac{28}{5}&\frac{21}{5}\end{pmatrix}$ 求一个紧致奇异值分解，给出秩、两个奇异值、谱范数及 Moore--Penrose 广义逆，并验证 $A^+A=I_2$。
\end{exercise}
\noindent 编者补充：S-SVD-6；相关主题见原讲义对应讲，题目数据与解答由本选编另行构造。
\begin{solve}
取
\[U=\begin{pmatrix}\frac{3}{5}&0\\\frac{4}{5}&0\\0&1\end{pmatrix},\quad\Sigma=\operatorname{diag}(9,7),\quad V=\begin{pmatrix}\frac{3}{5}&-\frac{4}{5}\\\frac{4}{5}&\frac{3}{5}\end{pmatrix}.\]
直接核对 $U^{\mathrm T}U=I_2$、$V^{\mathrm T}V=I_2$ 和 $A=U\Sigma V^{\mathrm T}$，所以这就是紧致奇异值分解。两个奇异值均为正，故秩为 $2$；谱范数为较大的奇异值 $9$。广义逆为
\[A^+=V\Sigma^{-1}U^{\mathrm T}=\begin{pmatrix}\frac{1}{25}&\frac{4}{75}&-\frac{4}{35}\\\frac{4}{75}&\frac{16}{225}&\frac{3}{35}\end{pmatrix}.\]
于是 $A^+A=V\Sigma^{-1}\Sigma V^{\mathrm T}=I_2$，而 $AA^+=UU^{\mathrm T}$ 是到 $A$ 列空间的正交投影。后者通常不等于 $I_3$，这与矩阵只有两列相符。
\end{solve}

\section{线性代数与解析几何}\label{sec:topic21}
本节收录 6 道题。用线性方程和参数表示描述几何对象，区分方向空间与位置向量；距离、垂直和夹角的讨论依赖给定内积。
\begin{exercise}\label{prob:L21-例-1}
若已知一个平面上有三点
    \[
        A(1, 2, 0), \quad B(0, 1, -1), \quad C(1, 1, 1),
    \]
    求该平面的一般方程.
\end{exercise}
\noindent 来源：\href{https://github.com/yhwu-is/Linear-Algebra-Left-Undone/blob/new/%E8%AE%B2%E4%B9%89/%E4%B8%93%E9%A2%98/21%20%E7%BA%BF%E6%80%A7%E4%BB%A3%E6%95%B0%E4%B8%8E%E5%87%A0%E4%BD%95.tex}{L21-例-1}。
\begin{solve}
取平面内两个方向向量
    \[
        \overrightarrow{AB} = (-1, -1, -1), \qquad \overrightarrow{AC} = (0, -1, 1).
    \]
    因此平面的一个法向量为
    \[
        \vec{n} = \overrightarrow{AB}\times\overrightarrow{AC} = (-2, 1, 1).
    \]
    因此平面的点法式可写成
    \[
        -2(x - 1) + (y - 2) + z = 0,
    \]
    化简得 $-2x + y + z = 0$.
\end{solve}

\begin{exercise}\label{prob:L21-例-2}
将直线
    \[
        \begin{cases}
            x + y - z = 2, \\
            2x - y + z = 1
        \end{cases}
    \]
    化为参数方程.
\end{exercise}
\noindent 来源：\href{https://github.com/yhwu-is/Linear-Algebra-Left-Undone/blob/new/%E8%AE%B2%E4%B9%89/%E4%B8%93%E9%A2%98/21%20%E7%BA%BF%E6%80%A7%E4%BB%A3%E6%95%B0%E4%B8%8E%E5%87%A0%E4%BD%95.tex}{L21-例-2}。
\begin{solve}
两平面的法向量分别为
    \[
        \vec{n}_1 = (1, 1, -1), \qquad \vec{n}_2 = (2, -1, 1).
    \]
    因此交线的一个方向向量为
    \[
        \vec{l} = \vec{n}_1 \times \vec{n}_2 = (0, -3, -3),
    \]
    可化简为 $(0, 1, 1)$. 取 $z = 0$，解方程组
    \[
        \begin{cases}
            x + y = 2, \\
            2x - y = 1
        \end{cases}
    \]
    得到直线上一点 $(1, 1, 0)$. 故所求参数方程为
    \[
        \begin{cases}
            x = 1, \\
            y = 1 + t, \\
            z = t.
        \end{cases}
    \]
\end{solve}

\begin{exercise}\label{prob:L21-例-3}
已知直线
    \[
        L_1 =
        \begin{cases}
            x + y + z - 1 = 0, \\
            x - 2y + 2 = 0
        \end{cases},
        \qquad
        L_2 =
        \begin{cases}
            x = 2t, \\
            y = t + a, \\
            z = bt + 1
        \end{cases},
    \]
    试确定 $a, b$ 的值，使得 $L_1, L_2$ 分别是：
    \begin{enumerate}
        \item 平行直线；
        \item 异面直线.
    \end{enumerate}
\end{exercise}
\noindent 来源：\href{https://github.com/yhwu-is/Linear-Algebra-Left-Undone/blob/new/%E8%AE%B2%E4%B9%89/%E4%B8%93%E9%A2%98/21%20%E7%BA%BF%E6%80%A7%E4%BB%A3%E6%95%B0%E4%B8%8E%E5%87%A0%E4%BD%95.tex}{L21-例-3}。
\begin{solve}
先把 $L_1$ 化为参数方程. 只需要将 $L_1$ 方程视为线性方程组来求解即可得到：
    \[
        L_1:\ (x, y, z) = (-2, 0, 3) + s (2, 1, -3).
    \]
    因而 $L_1$ 的方向向量可取
    \[
        \vec{u} = (2, 1, -3),
    \]
    而 $L_2$ 的方向向量为
    \[
        \vec{v} = (2, 1, b).
    \]

    \begin{enumerate}
        \item 若 $L_1, L_2$ 平行，则 $\vec{u}, \vec{v}$ 平行. 由于前两坐标已经相同，必有 $b = -3$. 此时两直线不可能重合，因为把 $L_2$ 上点 $(0, a, 1)$ 代入 $L_1$ 的两个平面方程分别得到
        \[
            a = 0, \qquad 2 - 2a = 0,
        \]
        二者矛盾. 故当且仅当 $b = -3$ 时，$L_1, L_2$ 为平行直线，而 $a$ 任意.

        \item 若 $L_1, L_2$ 异面，则必须不平行且不共面. 取
        \[
            \vec{r}_1 = (-2, 0, 3), \qquad \vec{r}_2 = (0, a, 1),
        \]
        则
        \[
            \vec{r}_2 - \vec{r}_1 = (2, a, -2).
        \]
        计算混合积
        \[
            [\vec{r}_2 - \vec{r}_1, \vec{u}, \vec{v}]
            =
            \begin{vmatrix}
                2 & a & -2 \\
                2 & 1 & -3 \\
                2 & 1 & b
            \end{vmatrix}
            = 2(b + 3)(1 - a).
        \]
        因而异面的充分必要条件是 $2(b + 3)(1 - a) \neq 0$，即 $a \neq 1, \qquad b \neq -3$.
    \end{enumerate}
\end{solve}

\begin{exercise}\label{prob:L21-例-4}
求球面
    \[
        x^2 + y^2 + z^2 - 2x + 4y - 2z - 20 = 0
    \]
    在点 $P(3, -1, 2)$ 处的切平面.
\end{exercise}
\noindent 来源：\href{https://github.com/yhwu-is/Linear-Algebra-Left-Undone/blob/new/%E8%AE%B2%E4%B9%89/%E4%B8%93%E9%A2%98/21%20%E7%BA%BF%E6%80%A7%E4%BB%A3%E6%95%B0%E4%B8%8E%E5%87%A0%E4%BD%95.tex}{L21-例-4}。
\begin{solve}
将球面配方得
    \[
        (x - 1)^2 + (y + 2)^2 + (z - 1)^2 = 26,
    \]
    故球心为 $C(1, -2, 1)$. 点 $P$ 处的法向量可取
    \[
        \overrightarrow{CP} = (2, 1, 1).
    \]
    因而切平面方程为
    \[
        2(x - 3) + (y + 1) + (z - 2) = 0,
    \]
    即
    \[
        2x + y + z - 7 = 0.
    \]
\end{solve}

\begin{exercise}\label{prob:L21-例-5}
设 $F:\mathbb{R}^2 \to \mathbb{R}^2$ 由
    \[
        F(x, y) = (x + 2y + 1, y - 1)
    \]
    给出. 说明它的几何意义.
\end{exercise}
\noindent 来源：\href{https://github.com/yhwu-is/Linear-Algebra-Left-Undone/blob/new/%E8%AE%B2%E4%B9%89/%E4%B8%93%E9%A2%98/21%20%E7%BA%BF%E6%80%A7%E4%BB%A3%E6%95%B0%E4%B8%8E%E5%87%A0%E4%BD%95.tex}{L21-例-5}。
\begin{solve}
该映射可分解为
    \[
        (x, y) \mapsto (x + 2y, y)
    \]
    与
    \[
        (x, y) \mapsto (x + 1, y - 1)
    \]
    的复合. 前者是沿 $x$ 方向的错切，后者是平移. 因而 $F$ 是一个典型的仿射变换.
\end{solve}

\begin{exercise}\label{prob:L21-例-6}
将二次曲面
    \[
        x^2 + y^2 + 2 x y - 4 z + 2 = 0
    \]
    化为标准型，并判定其类型.
\end{exercise}
\noindent 来源：\href{https://github.com/yhwu-is/Linear-Algebra-Left-Undone/blob/new/%E8%AE%B2%E4%B9%89/%E4%B8%93%E9%A2%98/21%20%E7%BA%BF%E6%80%A7%E4%BB%A3%E6%95%B0%E4%B8%8E%E5%87%A0%E4%BD%95.tex}{L21-例-6}。
\begin{solve}
二次项对应矩阵为
    \[
        A =
        \begin{pmatrix}
            1 & 1 & 0 \\
            1 & 1 & 0 \\
            0 & 0 & 0
        \end{pmatrix}.
    \]
    它的二次型可直接写成
    \[
        x^2 + y^2 + 2 x y = (x + y)^2.
    \]
    取正交变换
    \[
        u = \frac{x + y}{\sqrt{2}}, \qquad
        v = \frac{x - y}{\sqrt{2}}, \qquad
        w = z.
    \]
    则原方程化为
    \[
        2 u^2 - 4 w + 2 = 0,
    \]
    令 $w' = w - \frac{1}{2}$，则 $S$ 的方程化为
    \[
    u^2 = 2 w',
    \]
    这就是抛物柱面型的标准方程.
\end{solve}

\section{二次型}\label{sec:topic22}
本节收录 22 道题，其中 12 道为编者补充计算题。二次型的变量替换对应合同变换；正交对角化还可给出特征值。判断正定性、惯性指数或约束极值时应分别使用相应条件。
\begin{exercise}\label{prob:L22-例-1}
回答以下问题：
    \begin{enumerate}
        \item 已知$A$是一个$n$阶矩阵，则$A$为反对称矩阵的充要条件是对任意$n$元列向量$X$都有$X^\mathrm{T}AX=0$；

        \item 若二次型$f(x_1,x_2,\ldots,x_n)=X^\mathrm{T}AX$对任意的$n$元列向量$X$都有$f(x_1,x_2,\ldots,x_n)=0$，证明：$A=O$；

        \item 设二次型$f(x_1,x_2,\ldots,x_n)=X^\mathrm{T}AX,\enspace g(x_1,x_2,\ldots,x_n)=X^\mathrm{T}BX$.\\
              证明：若$f(x_1,x_2,\ldots,x_n)=g(x_1,x_2,\ldots,x_n)$，则$A=B$.
    \end{enumerate}
\end{exercise}
\noindent 来源：\href{https://github.com/yhwu-is/Linear-Algebra-Left-Undone/blob/new/%E8%AE%B2%E4%B9%89/%E4%B8%93%E9%A2%98/22%20%E4%BA%8C%E6%AC%A1%E5%9E%8B.tex}{L22-例-1}。
\begin{solve}
\begin{enumerate}
        \item 先证充分性. 若 $A^\mathrm{T}=-A$，则对任意列向量 $X$ 都有
              \[
                  (X^\mathrm{T}AX)^\mathrm{T}=X^\mathrm{T}A^\mathrm{T}X=-X^\mathrm{T}AX.
              \]
              由于 $X^\mathrm{T}AX$ 是一个 $1$ 阶矩阵，也就是一个数，它等于自身的转置，于是
              \[
                  X^\mathrm{T}AX=-X^\mathrm{T}AX,
              \]
              从而 $X^\mathrm{T}AX=0$.

              再证必要性. 若对任意 $X$ 都有 $X^\mathrm{T}AX=0$，取 $X=e_i$ 得 $a_{ii}=0$. 再取 $X=e_i+e_j\enspace(i\neq j)$，则
              \[
                  0=(e_i+e_j)^\mathrm{T}A(e_i+e_j)=a_{ii}+a_{ij}+a_{ji}+a_{jj}=a_{ij}+a_{ji}.
              \]
              故 $a_{ji}=-a_{ij}$ 对任意 $i,j$ 都成立，即 $A^\mathrm{T}=-A$.

        \item 因为 $f(X)=X^\mathrm{T}AX$ 是二次型，所以这里的 $A$ 是对称矩阵. 已知对任意 $X$ 都有 $X^\mathrm{T}AX=0$，取 $X=e_i$，得到 $a_{ii}=0$. 再取 $X=e_i+e_j\enspace(i\neq j)$，则
              \[
                  0=(e_i+e_j)^\mathrm{T}A(e_i+e_j)=a_{ii}+2a_{ij}+a_{jj}=2a_{ij}.
              \]
              故 $a_{ij}=0$. 因而 $A=O$.

        \item 由 $f(X)=g(X)$ 可知对任意 $X$ 都有
              \[
                  X^\mathrm{T}(A-B)X=0.
              \]
              又 $A,B$ 都是二次型的相伴矩阵，所以它们都是对称矩阵，从而 $A-B$ 也是对称矩阵. 由第 2 问立刻得到 $A-B=O$，即 $A=B$.
    \end{enumerate}
\end{solve}

\begin{exercise}\label{prob:L22-例-2}
已知二次型
    \[f(X)=(x_1,x_2,x_3,x_4)\begin{pmatrix}
            1 & 2 & 3 & -4 \\ 3 & 2 & 1 & 4 \\ -4 & 3 & -7 & 2 \\ 0 & -6 & 8 & 4
        \end{pmatrix}\begin{pmatrix}
            x_1 \\ x_2 \\ x_3 \\ x_4
        \end{pmatrix}\]
    写出二次型$f(X)$的相伴矩阵.
\end{exercise}
\noindent 来源：\href{https://github.com/yhwu-is/Linear-Algebra-Left-Undone/blob/new/%E8%AE%B2%E4%B9%89/%E4%B8%93%E9%A2%98/22%20%E4%BA%8C%E6%AC%A1%E5%9E%8B.tex}{L22-例-2}。
\begin{solve}
设题中给出的矩阵为
    \[
        M=\begin{pmatrix}
            1 & 2 & 3 & -4 \\
            3 & 2 & 1 & 4 \\
            -4 & 3 & -7 & 2 \\
            0 & -6 & 8 & 4
        \end{pmatrix}.
    \]
    因为
    \[
        X^\mathrm{T}MX=X^\mathrm{T}\left(\frac{M+M^\mathrm{T}}{2}\right)X+X^\mathrm{T}\left(\frac{M-M^\mathrm{T}}{2}\right)X,
    \]
    而 $ \dfrac{M-M^\mathrm{T}}{2} $ 是反对称矩阵，故对任意 $X$ 都有
    \[
        X^\mathrm{T}\left(\frac{M-M^\mathrm{T}}{2}\right)X=0.
    \]
    因此二次型 $f(X)$ 的相伴矩阵就是 $M$ 的对称部分
    \[
        A=\frac{M+M^\mathrm{T}}{2}
        =\begin{pmatrix}
            1 & \dfrac{5}{2} & -\dfrac{1}{2} & -2 \\
            \dfrac{5}{2} & 2 & 2 & -1 \\
            -\dfrac{1}{2} & 2 & -7 & 5 \\
            -2 & -1 & 5 & 4
        \end{pmatrix}.
    \]
\end{solve}

\begin{exercise}\label{prob:L22-例-3}
利用初等变换法将二次型 $f(x_1,x_2,x_3)=2x_1x_2-2x_1x_3+2x_2x_3$ 化为标准形，并且求出坐标变换矩阵$C$.
\end{exercise}
\noindent 来源：\href{https://github.com/yhwu-is/Linear-Algebra-Left-Undone/blob/new/%E8%AE%B2%E4%B9%89/%E4%B8%93%E9%A2%98/22%20%E4%BA%8C%E6%AC%A1%E5%9E%8B.tex}{L22-例-3}。
\begin{solve}
类似于利用初等变换求矩阵的逆，我们将相伴矩阵和单位矩阵拼接为 $\begin{pmatrix}
        A \\ E
    \end{pmatrix}$，对 $A$ 进行初等行列变换，使其变为对角矩阵，同时对 $E$ 进行对应的列变换，最终得到的 $\begin{pmatrix}
        \Lambda \\ C
    \end{pmatrix}$ 就是我们要求的结果.

    \begin{align*}
        \left(\begin{array}{c}
            A \\ \hdashline E
        \end{array}\right) ={} & \left(\begin{array}{ccc}
            0 & 1 & -1 \\
            1 & 0 & 1  \\
            -1 & 1 & 0 \\
            \hdashline
            1 & 0 & 0 \\
            0 & 1 & 0 \\
            0 & 0 & 1
        \end{array}\right)
        \xrightarrow[r_1+r_2]{c_1+c_2} \left(\begin{array}{ccc}
            2 & 1 & 0 \\
            1 & 0 & 1  \\
            0 & 1 & 0 \\
            \hdashline
            1 & 0 & 0 \\
            1 & 1 & 0 \\
            0 & 0 & 1
        \end{array}\right) & \\
        \xrightarrow[r_2-\dfrac{1}{2}r_1]{c_2-\dfrac{1}{2}c_1} & \left(\begin{array}{ccc}
            2 & 0 & 0 \\
            0 & -\dfrac{1}{2} & 1  \\
            0 & 1 & 0 \\
            \hdashline
            1 & -\dfrac{1}{2} & 0 \\
            1 & \dfrac{1}{2} & 0 \\
            0 & 0 & 1
        \end{array}\right)
        \xrightarrow[r_3+2r_2]{c_3+2c_2} \left(\begin{array}{ccc}
            2 & 0 & 0 \\
            0 & -\dfrac{1}{2} & 0  \\
            0 & 0 & 2 \\
            \hdashline
            1 & -\dfrac{1}{2} & -1 \\
            1 & \dfrac{1}{2} & 1 \\
            0 & 0 & 1
        \end{array}\right) & \\
    \end{align*}
\end{solve}

\begin{exercise}\label{prob:L22-例-4}
用配方法把三元二次型
    \[f(x_1,x_2,x_3)=2x_1^2+3x_2^2+x_3^2+4x_1x_2-4x_1x_3-8x_2x_3\]
    化为标准形，并求所用的坐标变换$X=CY$即变换矩阵$C$.
\end{exercise}
\noindent 来源：\href{https://github.com/yhwu-is/Linear-Algebra-Left-Undone/blob/new/%E8%AE%B2%E4%B9%89/%E4%B8%93%E9%A2%98/22%20%E4%BA%8C%E6%AC%A1%E5%9E%8B.tex}{L22-例-4}。
\begin{solve}
直接配方可得
    \begin{align*}
        f(x_1,x_2,x_3)
         & = 2x_1^2+4x_1x_2-4x_1x_3+3x_2^2-8x_2x_3+x_3^2 \\
         & = 2(x_1+x_2-x_3)^2+x_2^2-4x_2x_3-3x_3^2         \\
         & = 2(x_1+x_2-x_3)^2+(x_2-2x_3)^2-5x_3^2.
    \end{align*}
    令
    \[
        \begin{cases}
            y_1 = x_1+x_2-x_3, \\
            y_2 = x_2-2x_3,    \\
            y_3 = x_3,
        \end{cases}
    \]
    则
    \[
        f(y_1,y_2,y_3)=2y_1^2+y_2^2-5y_3^2.
    \]
    由上式反解得
    \[
        \begin{cases}
            x_3 = y_3,             \\
            x_2 = y_2+2y_3,        \\
            x_1 = y_1-y_2-y_3,
        \end{cases}
    \]
    因而所求坐标变换 $X=CY$ 为
    \[
        \begin{pmatrix}
            x_1 \\
            x_2 \\
            x_3
        \end{pmatrix}
        =
        \begin{pmatrix}
            1 & -1 & -1 \\
            0 & 1 & 2   \\
            0 & 0 & 1
        \end{pmatrix}
        \begin{pmatrix}
            y_1 \\
            y_2 \\
            y_3
        \end{pmatrix}.
    \]
\end{solve}

\begin{exercise}\label{prob:L22-例-5}
用配方法把二次型
    \[f(x_1,x_2,x_3)=2x_1x_2-2x_1x_3+2x_2x_3\]
    化为标准形，并求所用的坐标变换$X=CY$即变换矩阵$C$.
\end{exercise}
\noindent 来源：\href{https://github.com/yhwu-is/Linear-Algebra-Left-Undone/blob/new/%E8%AE%B2%E4%B9%89/%E4%B8%93%E9%A2%98/22%20%E4%BA%8C%E6%AC%A1%E5%9E%8B.tex}{L22-例-5}。
\begin{solve}
对于 $f(x_1,x_2,x_3)=2x_1x_2-2x_1x_3+2x_2x_3$，原式中没有平方项，因此不能直接像通常那样配方. 这时先作一个可逆线性变换
    \[
        \begin{cases}
            y_1 = x_1+x_2, \\
            y_2 = x_1-x_2, \\
            y_3 = x_3.
        \end{cases}
    \]
    由此反解得
    \[
        \begin{cases}
            x_1 = \dfrac{y_1+y_2}{2}, \\
            x_2 = \dfrac{y_1-y_2}{2}, \\
            x_3 = y_3.
        \end{cases}
    \]
    代入原二次型，有
    \[ f(y_1, y_2, y_3) = \frac{1}{2}y_1^2-\frac{1}{2}y_2^2-2y_2y_3. \]
    经过这一步，二次型中已经出现了平方项，于是便可以继续配方：
    \[ f(y_1, y_2, y_3) = \frac{1}{2}y_1^2-\frac{1}{2}(y_2+2y_3)^2+2y_3^2. \]
    再令
    \[
        z_1=y_1,\qquad z_2=y_2+2y_3,\qquad z_3=y_3,
    \]
    便得到标准形
    \[
        f(z_1,z_2,z_3)=\frac{1}{2}z_1^2-\frac{1}{2}z_2^2+2z_3^2.
    \]
    如果写成从新坐标到旧坐标的形式，则由
    \[
        y_1=z_1,\qquad y_2=z_2-2z_3,\qquad y_3=z_3
    \]
    以及
    \[
        x_1=\frac{y_1+y_2}{2},\qquad x_2=\frac{y_1-y_2}{2},\qquad x_3=y_3
    \]
    可得
    \[
        \begin{cases}
            x_1 = \dfrac{1}{2}z_1+\dfrac{1}{2}z_2-z_3, \\
            x_2 = \dfrac{1}{2}z_1-\dfrac{1}{2}z_2+z_3, \\
            x_3 = z_3.
        \end{cases}
    \]
    因此坐标变换矩阵可取为
    \[
        X=
        \begin{pmatrix}
            \dfrac{1}{2} & \dfrac{1}{2} & -1 \\
            \dfrac{1}{2} & -\dfrac{1}{2} & 1 \\
            0 & 0 & 1
        \end{pmatrix}
        Z.
    \]
    这说明：当原二次型没有平方项、无法直接配方时，可以先通过一个简单的可逆线性变换人为造出平方项，再继续做配方法.
\end{solve}

\begin{exercise}\label{prob:L22-例-6}
用正交变换法将二次型 $f(x_1,x_2) = 2x_1^2 + 5x_2^2 + 4x_1x_2$ 化为标准形，并给出变换矩阵 $P$.
\end{exercise}
\noindent 来源：\href{https://github.com/yhwu-is/Linear-Algebra-Left-Undone/blob/new/%E8%AE%B2%E4%B9%89/%E4%B8%93%E9%A2%98/22%20%E4%BA%8C%E6%AC%A1%E5%9E%8B.tex}{L22-例-6}。
\begin{solve}
二次型的相伴矩阵为
    \[
        A=\begin{pmatrix}
            2 & 2 \\
            2 & 5
        \end{pmatrix}.
    \]
    其特征多项式为
    \[
        |\lambda E-A|=\begin{vmatrix}
            \lambda-2 & -2        \\
            -2        & \lambda-5
        \end{vmatrix}
        =\lambda^2-7\lambda+6=(\lambda-1)(\lambda-6),
    \]
    故特征值为 $1,6$.

    对于特征值 $1$，解方程 $(A-E)X=0$ 得特征向量可取为 $(-2,1)^\mathrm{T}$；对于特征值 $6$，解方程 $(A-6E)X=0$ 得特征向量可取为 $(1,2)^\mathrm{T}$. 将它们单位化，得到两个单位特征向量
    \[
        u_1=\frac{1}{\sqrt{5}}\begin{pmatrix}
            -2 \\
            1
        \end{pmatrix},
        \qquad
        u_2=\frac{1}{\sqrt{5}}\begin{pmatrix}
            1 \\
            2
        \end{pmatrix}.
    \]
    令
    \[
        P=\begin{pmatrix}
            -\dfrac{2}{\sqrt{5}} & \dfrac{1}{\sqrt{5}} \\
            \dfrac{1}{\sqrt{5}}  & \dfrac{2}{\sqrt{5}}
        \end{pmatrix},
    \]
    则 $P$ 为正交矩阵，且
    \[
        P^\mathrm{T}AP=\operatorname{diag}(1,6).
    \]
    因此在正交变换 $X=PY$ 下，
    \[
        f(X)=Y^\mathrm{T}(P^\mathrm{T}AP)Y=y_1^2+6y_2^2.
    \]
    这就是所求标准形.
\end{solve}

\begin{exercise}\label{prob:L22-例-7}
确定二次型$f(x_1,x_2,\ldots,x_{10})=x_1x_2+x_3x_4+x_5x_6+x_7x_8+x_9x_{10}$的秩以及正、负惯性指数.
\end{exercise}
\noindent 来源：\href{https://github.com/yhwu-is/Linear-Algebra-Left-Undone/blob/new/%E8%AE%B2%E4%B9%89/%E4%B8%93%E9%A2%98/22%20%E4%BA%8C%E6%AC%A1%E5%9E%8B.tex}{L22-例-7}。
\begin{solve}
对每一对变量分别配方：
    \[
        x_{2k-1}x_{2k}=\frac{1}{4}\bigl((x_{2k-1}+x_{2k})^2-(x_{2k-1}-x_{2k})^2\bigr), \quad k=1,2,\ldots,5.
    \]
    因而
    \[
        f=\frac{1}{4}\sum_{k=1}^5\bigl((x_{2k-1}+x_{2k})^2-(x_{2k-1}-x_{2k})^2\bigr).
    \]
    令
    \[
        y_{2k-1}=x_{2k-1}+x_{2k}, \qquad y_{2k}=x_{2k-1}-x_{2k}, \qquad k=1,2,\ldots,5,
    \]
    则这是一个可逆线性变换，并且
    \[
        f=\frac{1}{4}(y_1^2-y_2^2+y_3^2-y_4^2+y_5^2-y_6^2+y_7^2-y_8^2+y_9^2-y_{10}^2).
    \]
    因而该二次型的相合规范形中有 $5$ 个正项和 $5$ 个负项，没有零项，所以其秩为 $10$，正惯性指数为 $5$，负惯性指数为 $5$.
\end{solve}

\begin{exercise}\label{prob:L22-例-8}
求$t$的取值范围，使得下面的二次型是正定的：
    \[f(x_1,x_2,x_3,x_4) = x_1^2+4x_2^2+4x_3^2+3x_4^2+2tx_1x_2-2x_1x_3+4x_1x_2.\]
\end{exercise}
\noindent 来源：\href{https://github.com/yhwu-is/Linear-Algebra-Left-Undone/blob/new/%E8%AE%B2%E4%B9%89/%E4%B8%93%E9%A2%98/22%20%E4%BA%8C%E6%AC%A1%E5%9E%8B.tex}{L22-例-8}。
\begin{solve}
该二次型的相伴矩阵为
    \[
        A=\begin{pmatrix}
            1 & t+2 & -1 & 0 \\
            t+2 & 4 & 0 & 0 \\
            -1 & 0 & 4 & 0 \\
            0 & 0 & 0 & 3
        \end{pmatrix}.
    \]
    由 Sylvester 判别法，只需计算顺序主子式：
    \[
        D_1=1>0,
    \]
    \[
        D_2=\begin{vmatrix}
            1 & t+2 \\
            t+2 & 4
        \end{vmatrix}
        =4-(t+2)^2,
    \]
    \[
        D_3 = \begin{vmatrix}
                 1 & t+2 & -1 \\
                 t+2 & 4 & 0 \\
                 -1 & 0 & 4
             \end{vmatrix} = 4\bigl(3-(t+2)^2\bigr),
    \]
    \[
        D_4=3D_3=12\bigl(3-(t+2)^2\bigr).
    \]
    因而正定的充要条件是 $D_2>0,\qquad D_3>0,\qquad D_4>0$. 其中 $D_3>0$ 已经推出 $D_2>0$ 和 $D_4>0$，故最终只需 $3-(t+2)^2>0$. 即 $-2-\sqrt{3}<t<-2+\sqrt{3}$.
\end{solve}

\begin{exercise}\label{prob:L22-例-9}
证明：每个秩为$r$的$n\enspace(r<n)$阶实对称矩阵均可表示为$n-r$个秩为$n-1$的实对称矩阵的乘积.
\end{exercise}
\noindent 来源：\href{https://github.com/yhwu-is/Linear-Algebra-Left-Undone/blob/new/%E8%AE%B2%E4%B9%89/%E4%B8%93%E9%A2%98/22%20%E4%BA%8C%E6%AC%A1%E5%9E%8B.tex}{L22-例-9}。
\begin{solve}
设 $A$ 是秩为 $r$ 的 $n$ 阶实对称矩阵，且 $r<n$. 由原讲义“实谱定理”处，存在正交矩阵 $Q$ 使得 $A=Q^\mathrm{T}\Lambda Q$，其中 $\Lambda=\operatorname{diag}(\lambda_1,\ldots,\lambda_r,0,\ldots,0)$ 且 $\lambda_1,\ldots,\lambda_r\neq 0$.

    下面将对角矩阵 $\Lambda$ 分解成 $n-r$ 个秩为 $n-1$ 的对角矩阵的乘积. 定义
    \[
        D_1=\operatorname{diag}(\lambda_1,\ldots,\lambda_r,0,1,\ldots,1),
    \]
    对于 $s=2,\ldots,n-r$，定义
    \[
        D_s=\operatorname{diag}(\underbrace{1,\ldots,1}_{r+s-1},0,1,\ldots,1).
    \]
    每个 $D_s$ 都是实对称矩阵，且恰有一个对角元为 $0$，因此 $r(D_s)=n-1$ 并且显然有 $\Lambda=D_1D_2\cdots D_{n-r}$. 令 $A_s=Q^\mathrm{T}D_sQ$（$s=1,2,\ldots,n-r$），则每个 $A_s$ 都是实对称矩阵，并且由于正交相似保持秩不变，有 $r(A_s)=n-1$. 同时
    \[
        A_1A_2\cdots A_{n-r}
        =Q^\mathrm{T}D_1Q\cdots Q^\mathrm{T}D_{n-r}Q
        =Q^\mathrm{T}(D_1D_2\cdots D_{n-r})Q
        =Q^\mathrm{T}\Lambda Q
        =A.
    \]
    因而 $A$ 可表示为 $n-r$ 个秩为 $n-1$ 的实对称矩阵的乘积.
\end{solve}

\begin{exercise}\label{prob:L22-例-10}
求二次型
    \[
        f(x_1,x_2)=2x_1^2+5x_2^2+4x_1x_2
    \]
    在约束 $x_1^2+x_2^2=1$ 下的最大值与最小值.
\end{exercise}
\noindent 来源：\href{https://github.com/yhwu-is/Linear-Algebra-Left-Undone/blob/new/%E8%AE%B2%E4%B9%89/%E4%B8%93%E9%A2%98/22%20%E4%BA%8C%E6%AC%A1%E5%9E%8B.tex}{L22-例-10}。
\begin{solve}
该二次型的相伴矩阵为
    \[
        A=\begin{pmatrix}
            2 & 2 \\
            2 & 5
        \end{pmatrix}.
    \]
    在原讲义“正交变换法例1”处已经算出它的特征值为 $1,6$. 由于约束条件 $x_1^2+x_2^2=1$ 等价于 $\|X\|=1$，故 $f(X)=X^\mathrm{T}AX=R_A(X)$. 由原讲义“Rayleigh-Ritz”处定理立刻得到
    \[
        \min_{\|X\|=1}f(X)=1,\qquad \max_{\|X\|=1}f(X)=6.
    \]
\end{solve}

\begin{exercise}\label{prob:S-QUAD-1}
设二次型 $f(x)=x^{\mathrm T}Ax$，$A=\begin{pmatrix}2&2&2\\2&0&4\\2&4&0\end{pmatrix}$。用可逆线性变量替换化为标准形，写明 $x=Py$ 的矩阵 $P$，求秩、正负惯性指数与零指数。
\end{exercise}
\noindent 编者补充：S-QUAD-1；相关主题见原讲义对应讲，题目数据与解答由本选编另行构造。
\begin{solve}
观察配方可取 $y=Cx$，其中
\[C=\begin{pmatrix}1&1&1\\0&1&-1\\0&0&1\end{pmatrix},\qquad P=C^{-1}=\begin{pmatrix}1&-1&-2\\0&1&1\\0&0&1\end{pmatrix}.\]
直接核对 $A=C^{\mathrm T}DC$，其中 $D=\operatorname{diag}(2,-2,0)$。因此 $x=Py$ 时，
\[P^{\mathrm T}AP=D,\qquad f(x)=2y_1^2-2y_2^2+0y_3^2.\]
正、负、零惯性指数依次为 $1,1,1$，秩为 $2$。因为 $\det C=1$，变换可逆，Sylvester 惯性定理保证这些指数与所选坐标无关。这里使用的是合同变换 $P^{\mathrm T}AP$；不能把它直接替换成相似变换 $P^{-1}AP$。
\end{solve}

\begin{exercise}\label{prob:S-CHOLESKY-1}
设 $A=\begin{pmatrix}1&1&1\\1&5&-1\\1&-1&11\end{pmatrix}$。用顺序主子式判定正定性，求对角元为正的 Cholesky 分解，并利用两次三角方程求解 $Ax=\begin{pmatrix}1\\7\\-11\end{pmatrix}$。
\end{exercise}
\noindent 编者补充：S-CHOLESKY-1；相关主题见原讲义对应讲，题目数据与解答由本选编另行构造。
\begin{solve}
三个顺序主子式依次为 $1,4,36$，均为正，故 $A$ 正定。按 Cholesky 递推逐列计算得到
\[A=LL^{\mathrm T},\qquad L=\begin{pmatrix}1&0&0\\1&2&0\\1&-1&3\end{pmatrix}.\]
先求解 $Ly=b$，再求解 $L^{\mathrm T}x=y$。这两步分别是前代和回代，得到
\[y=\begin{pmatrix}1\\3\\-3\end{pmatrix},\qquad x=\begin{pmatrix}1\\1\\-1\end{pmatrix}.\]
因为 $L$ 的三个对角元均为正且非零，两次三角方程都有唯一解；将所得 $x$ 代回原式可核验。另由 $z^{\mathrm T}Az=\|L^{\mathrm T}z\|^2$，对任意非零 $z$ 此值严格为正，这从分解角度再次证明了正定性。
\end{solve}

\begin{exercise}\label{prob:S-QUAD-2}
设二次型 $f(x)=x^{\mathrm T}Ax$，$A=\begin{pmatrix}3&6&3\\6&10&8\\3&8&4\end{pmatrix}$。用可逆线性变量替换化为标准形，写明 $x=Py$ 的矩阵 $P$，求秩、正负惯性指数与零指数。
\end{exercise}
\noindent 编者补充：S-QUAD-2；相关主题见原讲义对应讲，题目数据与解答由本选编另行构造。
\begin{solve}
观察配方可取 $y=Cx$，其中
\[C=\begin{pmatrix}1&2&1\\0&1&-1\\0&0&1\end{pmatrix},\qquad P=C^{-1}=\begin{pmatrix}1&-2&-3\\0&1&1\\0&0&1\end{pmatrix}.\]
直接核对 $A=C^{\mathrm T}DC$，其中 $D=\operatorname{diag}(3,-2,3)$。因此 $x=Py$ 时，
\[P^{\mathrm T}AP=D,\qquad f(x)=3y_1^2-2y_2^2+3y_3^2.\]
正、负、零惯性指数依次为 $2,1,0$，秩为 $3$。因为 $\det C=1$，变换可逆，Sylvester 惯性定理保证这些指数与所选坐标无关。这里使用的是合同变换 $P^{\mathrm T}AP$；不能把它直接替换成相似变换 $P^{-1}AP$。
\end{solve}

\begin{exercise}\label{prob:S-CHOLESKY-2}
设 $A=\begin{pmatrix}1&2&1\\2&8&0\\1&0&11\end{pmatrix}$。用顺序主子式判定正定性，求对角元为正的 Cholesky 分解，并利用两次三角方程求解 $Ax=\begin{pmatrix}4\\18\\-10\end{pmatrix}$。
\end{exercise}
\noindent 编者补充：S-CHOLESKY-2；相关主题见原讲义对应讲，题目数据与解答由本选编另行构造。
\begin{solve}
三个顺序主子式依次为 $1,4,36$，均为正，故 $A$ 正定。按 Cholesky 递推逐列计算得到
\[A=LL^{\mathrm T},\qquad L=\begin{pmatrix}1&0&0\\2&2&0\\1&-1&3\end{pmatrix}.\]
先求解 $Ly=b$，再求解 $L^{\mathrm T}x=y$。这两步分别是前代和回代，得到
\[y=\begin{pmatrix}4\\5\\-3\end{pmatrix},\qquad x=\begin{pmatrix}1\\2\\-1\end{pmatrix}.\]
因为 $L$ 的三个对角元均为正且非零，两次三角方程都有唯一解；将所得 $x$ 代回原式可核验。另由 $z^{\mathrm T}Az=\|L^{\mathrm T}z\|^2$，对任意非零 $z$ 此值严格为正，这从分解角度再次证明了正定性。
\end{solve}

\begin{exercise}\label{prob:S-QUAD-3}
设二次型 $f(x)=x^{\mathrm T}Ax$，$A=\begin{pmatrix}4&12&4\\12&34&14\\4&14&2\end{pmatrix}$。用可逆线性变量替换化为标准形，写明 $x=Py$ 的矩阵 $P$，求秩、正负惯性指数与零指数。
\end{exercise}
\noindent 编者补充：S-QUAD-3；相关主题见原讲义对应讲，题目数据与解答由本选编另行构造。
\begin{solve}
观察配方可取 $y=Cx$，其中
\[C=\begin{pmatrix}1&3&1\\0&1&-1\\0&0&1\end{pmatrix},\qquad P=C^{-1}=\begin{pmatrix}1&-3&-4\\0&1&1\\0&0&1\end{pmatrix}.\]
直接核对 $A=C^{\mathrm T}DC$，其中 $D=\operatorname{diag}(4,-2,0)$。因此 $x=Py$ 时，
\[P^{\mathrm T}AP=D,\qquad f(x)=4y_1^2-2y_2^2+0y_3^2.\]
正、负、零惯性指数依次为 $1,1,1$，秩为 $2$。因为 $\det C=1$，变换可逆，Sylvester 惯性定理保证这些指数与所选坐标无关。这里使用的是合同变换 $P^{\mathrm T}AP$；不能把它直接替换成相似变换 $P^{-1}AP$。
\end{solve}

\begin{exercise}\label{prob:S-CHOLESKY-3}
设 $A=\begin{pmatrix}1&3&1\\3&13&1\\1&1&11\end{pmatrix}$。用顺序主子式判定正定性，求对角元为正的 Cholesky 分解，并利用两次三角方程求解 $Ax=\begin{pmatrix}9\\41\\-7\end{pmatrix}$。
\end{exercise}
\noindent 编者补充：S-CHOLESKY-3；相关主题见原讲义对应讲，题目数据与解答由本选编另行构造。
\begin{solve}
三个顺序主子式依次为 $1,4,36$，均为正，故 $A$ 正定。按 Cholesky 递推逐列计算得到
\[A=LL^{\mathrm T},\qquad L=\begin{pmatrix}1&0&0\\3&2&0\\1&-1&3\end{pmatrix}.\]
先求解 $Ly=b$，再求解 $L^{\mathrm T}x=y$。这两步分别是前代和回代，得到
\[y=\begin{pmatrix}9\\7\\-3\end{pmatrix},\qquad x=\begin{pmatrix}1\\3\\-1\end{pmatrix}.\]
因为 $L$ 的三个对角元均为正且非零，两次三角方程都有唯一解；将所得 $x$ 代回原式可核验。另由 $z^{\mathrm T}Az=\|L^{\mathrm T}z\|^2$，对任意非零 $z$ 此值严格为正，这从分解角度再次证明了正定性。
\end{solve}

\begin{exercise}\label{prob:S-QUAD-4}
设二次型 $f(x)=x^{\mathrm T}Ax$，$A=\begin{pmatrix}5&20&5\\20&78&22\\5&22&6\end{pmatrix}$。用可逆线性变量替换化为标准形，写明 $x=Py$ 的矩阵 $P$，求秩、正负惯性指数与零指数。
\end{exercise}
\noindent 编者补充：S-QUAD-4；相关主题见原讲义对应讲，题目数据与解答由本选编另行构造。
\begin{solve}
观察配方可取 $y=Cx$，其中
\[C=\begin{pmatrix}1&4&1\\0&1&-1\\0&0&1\end{pmatrix},\qquad P=C^{-1}=\begin{pmatrix}1&-4&-5\\0&1&1\\0&0&1\end{pmatrix}.\]
直接核对 $A=C^{\mathrm T}DC$，其中 $D=\operatorname{diag}(5,-2,3)$。因此 $x=Py$ 时，
\[P^{\mathrm T}AP=D,\qquad f(x)=5y_1^2-2y_2^2+3y_3^2.\]
正、负、零惯性指数依次为 $2,1,0$，秩为 $3$。因为 $\det C=1$，变换可逆，Sylvester 惯性定理保证这些指数与所选坐标无关。这里使用的是合同变换 $P^{\mathrm T}AP$；不能把它直接替换成相似变换 $P^{-1}AP$。
\end{solve}

\begin{exercise}\label{prob:S-CHOLESKY-4}
设 $A=\begin{pmatrix}1&4&1\\4&20&2\\1&2&11\end{pmatrix}$。用顺序主子式判定正定性，求对角元为正的 Cholesky 分解，并利用两次三角方程求解 $Ax=\begin{pmatrix}16\\82\\-2\end{pmatrix}$。
\end{exercise}
\noindent 编者补充：S-CHOLESKY-4；相关主题见原讲义对应讲，题目数据与解答由本选编另行构造。
\begin{solve}
三个顺序主子式依次为 $1,4,36$，均为正，故 $A$ 正定。按 Cholesky 递推逐列计算得到
\[A=LL^{\mathrm T},\qquad L=\begin{pmatrix}1&0&0\\4&2&0\\1&-1&3\end{pmatrix}.\]
先求解 $Ly=b$，再求解 $L^{\mathrm T}x=y$。这两步分别是前代和回代，得到
\[y=\begin{pmatrix}16\\9\\-3\end{pmatrix},\qquad x=\begin{pmatrix}1\\4\\-1\end{pmatrix}.\]
因为 $L$ 的三个对角元均为正且非零，两次三角方程都有唯一解；将所得 $x$ 代回原式可核验。另由 $z^{\mathrm T}Az=\|L^{\mathrm T}z\|^2$，对任意非零 $z$ 此值严格为正，这从分解角度再次证明了正定性。
\end{solve}

\begin{exercise}\label{prob:S-QUAD-5}
设二次型 $f(x)=x^{\mathrm T}Ax$，$A=\begin{pmatrix}6&30&6\\30&148&32\\6&32&4\end{pmatrix}$。用可逆线性变量替换化为标准形，写明 $x=Py$ 的矩阵 $P$，求秩、正负惯性指数与零指数。
\end{exercise}
\noindent 编者补充：S-QUAD-5；相关主题见原讲义对应讲，题目数据与解答由本选编另行构造。
\begin{solve}
观察配方可取 $y=Cx$，其中
\[C=\begin{pmatrix}1&5&1\\0&1&-1\\0&0&1\end{pmatrix},\qquad P=C^{-1}=\begin{pmatrix}1&-5&-6\\0&1&1\\0&0&1\end{pmatrix}.\]
直接核对 $A=C^{\mathrm T}DC$，其中 $D=\operatorname{diag}(6,-2,0)$。因此 $x=Py$ 时，
\[P^{\mathrm T}AP=D,\qquad f(x)=6y_1^2-2y_2^2+0y_3^2.\]
正、负、零惯性指数依次为 $1,1,1$，秩为 $2$。因为 $\det C=1$，变换可逆，Sylvester 惯性定理保证这些指数与所选坐标无关。这里使用的是合同变换 $P^{\mathrm T}AP$；不能把它直接替换成相似变换 $P^{-1}AP$。
\end{solve}

\begin{exercise}\label{prob:S-CHOLESKY-5}
设 $A=\begin{pmatrix}1&5&1\\5&29&3\\1&3&11\end{pmatrix}$。用顺序主子式判定正定性，求对角元为正的 Cholesky 分解，并利用两次三角方程求解 $Ax=\begin{pmatrix}25\\147\\5\end{pmatrix}$。
\end{exercise}
\noindent 编者补充：S-CHOLESKY-5；相关主题见原讲义对应讲，题目数据与解答由本选编另行构造。
\begin{solve}
三个顺序主子式依次为 $1,4,36$，均为正，故 $A$ 正定。按 Cholesky 递推逐列计算得到
\[A=LL^{\mathrm T},\qquad L=\begin{pmatrix}1&0&0\\5&2&0\\1&-1&3\end{pmatrix}.\]
先求解 $Ly=b$，再求解 $L^{\mathrm T}x=y$。这两步分别是前代和回代，得到
\[y=\begin{pmatrix}25\\11\\-3\end{pmatrix},\qquad x=\begin{pmatrix}1\\5\\-1\end{pmatrix}.\]
因为 $L$ 的三个对角元均为正且非零，两次三角方程都有唯一解；将所得 $x$ 代回原式可核验。另由 $z^{\mathrm T}Az=\|L^{\mathrm T}z\|^2$，对任意非零 $z$ 此值严格为正，这从分解角度再次证明了正定性。
\end{solve}

\begin{exercise}\label{prob:S-QUAD-6}
设二次型 $f(x)=x^{\mathrm T}Ax$，$A=\begin{pmatrix}7&42&7\\42&250&44\\7&44&8\end{pmatrix}$。用可逆线性变量替换化为标准形，写明 $x=Py$ 的矩阵 $P$，求秩、正负惯性指数与零指数。
\end{exercise}
\noindent 编者补充：S-QUAD-6；相关主题见原讲义对应讲，题目数据与解答由本选编另行构造。
\begin{solve}
观察配方可取 $y=Cx$，其中
\[C=\begin{pmatrix}1&6&1\\0&1&-1\\0&0&1\end{pmatrix},\qquad P=C^{-1}=\begin{pmatrix}1&-6&-7\\0&1&1\\0&0&1\end{pmatrix}.\]
直接核对 $A=C^{\mathrm T}DC$，其中 $D=\operatorname{diag}(7,-2,3)$。因此 $x=Py$ 时，
\[P^{\mathrm T}AP=D,\qquad f(x)=7y_1^2-2y_2^2+3y_3^2.\]
正、负、零惯性指数依次为 $2,1,0$，秩为 $3$。因为 $\det C=1$，变换可逆，Sylvester 惯性定理保证这些指数与所选坐标无关。这里使用的是合同变换 $P^{\mathrm T}AP$；不能把它直接替换成相似变换 $P^{-1}AP$。
\end{solve}

\begin{exercise}\label{prob:S-CHOLESKY-6}
设 $A=\begin{pmatrix}1&6&1\\6&40&4\\1&4&11\end{pmatrix}$。用顺序主子式判定正定性，求对角元为正的 Cholesky 分解，并利用两次三角方程求解 $Ax=\begin{pmatrix}36\\242\\14\end{pmatrix}$。
\end{exercise}
\noindent 编者补充：S-CHOLESKY-6；相关主题见原讲义对应讲，题目数据与解答由本选编另行构造。
\begin{solve}
三个顺序主子式依次为 $1,4,36$，均为正，故 $A$ 正定。按 Cholesky 递推逐列计算得到
\[A=LL^{\mathrm T},\qquad L=\begin{pmatrix}1&0&0\\6&2&0\\1&-1&3\end{pmatrix}.\]
先求解 $Ly=b$，再求解 $L^{\mathrm T}x=y$。这两步分别是前代和回代，得到
\[y=\begin{pmatrix}36\\13\\-3\end{pmatrix},\qquad x=\begin{pmatrix}1\\6\\-1\end{pmatrix}.\]
因为 $L$ 的三个对角元均为正且非零，两次三角方程都有唯一解；将所得 $x$ 代回原式可核验。另由 $z^{\mathrm T}Az=\|L^{\mathrm T}z\|^2$，对任意非零 $z$ 此值严格为正，这从分解角度再次证明了正定性。
\end{solve}

\section{线性代数与微积分}\label{sec:topic24}
本节收录 1 道题。将局部方程求解与 Jacobi 矩阵的可逆性联系起来，明确函数的定义域、陪域、光滑性及局部唯一性的范围。
\begin{exercise}\label{prob:L24-例-2}
设 $U\subseteq\mathbf R^2$ 为开集，$(x_0,y_0)\in U$，$f\colon U\to\mathbf R$ 为 $C^k$ 函数，$k\ge1$，且 $f_y(x_0,y_0)\ne0$。证明方程 $f(x,y)=f(x_0,y_0)$ 在该点附近可唯一地写成 $y=\psi(x)$，并求 $\psi'(x)$。
\end{exercise}
\noindent 来源：\href{https://github.com/yhwu-is/Linear-Algebra-Left-Undone/blob/new/%E8%AE%B2%E4%B9%89/%E4%B8%93%E9%A2%98/24%20%E7%BA%BF%E6%80%A7%E4%BB%A3%E6%95%B0%E4%B8%8E%E5%BE%AE%E7%A7%AF%E5%88%86.tex}{L24-例-2}。
\begin{solve}
令 $c=f(x_0,y_0)$，构造 $F\colon U\to\mathbf R^2$，$F(x,y)=(x,f(x,y))$。其 Jacobi 矩阵为
\[DF(x_0,y_0)=\begin{pmatrix}1&0\\f_x(x_0,y_0)&f_y(x_0,y_0)\end{pmatrix},\]
行列式为 $f_y(x_0,y_0)\ne0$。由反函数定理，$F$ 在该点的某邻域上有 $C^k$ 逆映射。对足够接近 $x_0$ 的 $x$，写
\[F^{-1}(x,c)=(x,\psi(x));\]
第一分量确为 $x$，是因为 $F$ 保持第一坐标。于是 $\psi(x_0)=y_0$，且 $f(x,\psi(x))=c$。缩小为适当的乘积邻域后，逆映射的唯一性保证其中的全部解恰为这条图像。
由于 $f_y$ 连续，可进一步缩小邻域使其不为零。对恒等式求导得到
\[f_x(x,\psi(x))+f_y(x,\psi(x))\psi'(x)=0,\qquad
\psi'(x)=-\frac{f_x(x,\psi(x))}{f_y(x,\psi(x))}.\]
此处 $f$ 必须为实值函数；若把陪域误写成 $\mathbf R^2$，上述二维 Jacobi 矩阵及隐函数结论均不再适用。
\end{solve}

\clearpage
\appendix
\section{概念与证明训练}\label{app:proofs}
本附编保留方而静题解中原书 1A 至 3C 的 154 道题，按九个小节组织，题后给出可编辑的参考解答。P 标记保留原书小节与题号，内部互引改用本书编号。它与正文共享数学环境；所用多项式下标、复化记号与正文的区别已在编选说明中说明。

题目及解答由原 Typst 源码转写，原作者署名及许可见编选说明；转写中补充了必要假设，并修正若干索引、符号和证明缺口。原书数字编号仅用于标明引用出处，不代表本书对应编号。
\subsection{实数向量与复数向量}\label{app:section-1A}
\begin{exercise}\label{prob:P-1A-1}
证明：$\alpha + \beta = \beta + \alpha$ 对所有 $\alpha , \beta \in \mathbf C$ 成立。
\end{exercise}
\noindent 来源：\href{https://github.com/szdytom/LADRSolutions/blob/main/sections/1A.typ}{P-1A-1}（原书 1A 第 1 题；方而静题解）。
\begin{solve}
根据定义，令 $\alpha = a + b \mathrm i$，$\beta = c + d \mathrm i$（其中 $a , b , c , d \in \mathbf R$），则根据实数的加法交换律，有

\[\begin{aligned}\alpha + \beta & = (a + c) + (b + d) \mathrm i \\ & = (c + a) + (d + b) \mathrm i \\ & = \beta + \alpha\end{aligned}\]
\end{solve}

\begin{exercise}\label{prob:P-1A-2}
证明：$(\alpha + \beta) + \lambda = \alpha + (\beta + \lambda)$ 对所有 $\alpha , \beta , \lambda \in \mathbf C$ 成立。
\end{exercise}
\noindent 来源：\href{https://github.com/szdytom/LADRSolutions/blob/main/sections/1A.typ}{P-1A-2}（原书 1A 第 2 题；方而静题解）。
\begin{solve}
根据定义，令 $\alpha = a + b \mathrm i$，$\beta = c + d \mathrm i$，$\lambda = e + f \mathrm i$（其中 $a , b , c , d , e , f \in \mathbf R$），则有

\[\begin{aligned}(\alpha + \beta) + \lambda & = ((a + c) + (b + d) \mathrm i) + (e + f \mathrm i) \\ & = (a + c + e) + (b + d + f) \mathrm i \\ & = \alpha + (\beta + \lambda)\end{aligned}\]
\end{solve}

\begin{exercise}\label{prob:P-1A-3}
证明：$(\alpha \beta) \lambda = \alpha (\beta \lambda)$ 对所有 $\alpha , \beta , \lambda \in \mathbf C$ 成立。
\end{exercise}
\noindent 来源：\href{https://github.com/szdytom/LADRSolutions/blob/main/sections/1A.typ}{P-1A-3}（原书 1A 第 3 题；方而静题解）。
\begin{solve}
根据定义，令 $\alpha = a + b \mathrm i$，$\beta = c + d \mathrm i$，$\lambda = e + f \mathrm i$（其中 $a , b , c , d , e , f \in \mathbf R$），则有

\[\begin{aligned}(\alpha \beta) \lambda & = ((a + b \mathrm i) (c + d \mathrm i)) (e + f \mathrm i) \\ & = (a c - b d + (a d + b c) \mathrm i) (e + f \mathrm i) \\ & = (a c e - b d e - (a d + b c) f) + ((a d + b c) e + (a c - b d) f) \mathrm i \\ & = \alpha (\beta \lambda)\end{aligned}\]
\end{solve}

\begin{exercise}\label{prob:P-1A-4}
证明：$\lambda (\alpha + \beta) = \lambda \alpha + \lambda \beta$ 对所有 $\lambda , \alpha , \beta \in \mathbf C$ 成立。
\end{exercise}
\noindent 来源：\href{https://github.com/szdytom/LADRSolutions/blob/main/sections/1A.typ}{P-1A-4}（原书 1A 第 4 题；方而静题解）。
\begin{solve}
根据定义，令 $\alpha = a + b \mathrm i$，$\beta = c + d \mathrm i$，$\lambda = e + f \mathrm i$（其中 $a , b , c , d , e , f \in \mathbf R$），则有

\[\begin{aligned}\lambda (\alpha + \beta) & = (e + f \mathrm i) ((a + c) + (b + d) \mathrm i) \\ & = (e (a + c) - f (b + d)) + (f (a + c) + e (b + d)) \mathrm i \\ & = (e a - f b + e c - f d) + (f a + e b + f c + e d) \mathrm i \\ & = \lambda \alpha + \lambda \beta\end{aligned}\]
\end{solve}

\begin{exercise}\label{prob:P-1A-5}
证明：对于任意 $\alpha \in \mathbf C$，都存在唯一的 $\beta \in \mathbf C$ 使得 $\alpha + \beta = 0$。
\end{exercise}
\noindent 来源：\href{https://github.com/szdytom/LADRSolutions/blob/main/sections/1A.typ}{P-1A-5}（原书 1A 第 5 题；方而静题解）。
\begin{solve}
根据定义，令 $\alpha = a + b \mathrm i$（其中 $a , b \in \mathbf R$），则取 $\beta = (- a) + (- b) \mathrm i$，则有

\[\begin{aligned}\alpha + \beta & = (a + b \mathrm i) + ((- a) + (- b) \mathrm i) \\ & = (a - a) + (b - b) \mathrm i \\ & = 0\end{aligned}\]

因此，这样的 $\beta$ 存在。为了说明其唯一性，假设存在另一个 $\beta'$，也满足 $\alpha + \beta' = 0$，则有

\[\beta = \beta + 0 = \beta + (\alpha + \beta') = (\beta + \alpha) + \beta' = 0 + \beta' = \beta'\]
\end{solve}

\begin{exercise}\label{prob:P-1A-6}
证明：对于任意 $\alpha \in \mathbf C$（$\alpha \ne 0$），都存在唯一的 $\beta \in \mathbf C$ 使得 $\alpha \beta = 1$。
\end{exercise}
\noindent 来源：\href{https://github.com/szdytom/LADRSolutions/blob/main/sections/1A.typ}{P-1A-6}（原书 1A 第 6 题；方而静题解）。
\begin{solve}
根据定义，令 $\alpha = a + b \mathrm i$（其中 $a , b \in \mathbf R$），则取 $\beta = (\frac{a}{a^{2} + b^{2}}) - (\frac{b}{a^{2} + b^{2}}) \mathrm i$，则有

\[\begin{aligned}\alpha \beta & = (a + b \mathrm i) (\frac{a}{a^{2} + b^{2}} - \frac{b}{a^{2} + b^{2}} \mathrm i) \\ & = \frac{a^{2} + b^{2}}{a^{2} + b^{2}} \\ & = 1\end{aligned}\]

因此，这样的 $\beta$ 存在。为了说明其唯一性，我们假设存在另一个 $\beta'$，也满足 $\alpha \beta' = 1$，则有

\[\beta = \beta \cdot 1 = \beta (\alpha \beta') = (\beta \alpha) \beta' = 1 \cdot \beta' = \beta'\]
\end{solve}

\begin{exercise}\label{prob:P-1A-7}
证明：
\[\frac{- 1 + \sqrt{3} \mathrm i}{2}\]
是 $1$ 的立方根（意即它的立方等于 $1$）。
\end{exercise}
\noindent 来源：\href{https://github.com/szdytom/LADRSolutions/blob/main/sections/1A.typ}{P-1A-7}（原书 1A 第 7 题；方而静题解）。
\begin{solve}
\[\begin{aligned}(\frac{- 1 + \sqrt{3} \mathrm i}{2})^{3} & = \frac{(- 1 + \sqrt{3} \mathrm i)^{3}}{8} \\ & = \frac{1}{8} ((- 1)^{3} + 3 (- 1)^{2} (\sqrt{3} \mathrm i) + 3 (- 1) (\sqrt{3} \mathrm i)^{2} + (\sqrt{3} \mathrm i)^{3}) \\ & = \frac{1}{8} (- 1 + 3 \sqrt{3} \mathrm i + 9 - 3 \sqrt{3} \mathrm i) \\ & = 1\end{aligned}\]
\end{solve}

\begin{exercise}\label{prob:P-1A-8}
求 $\mathrm i$ 的两个相异的平方根。
\end{exercise}
\noindent 来源：\href{https://github.com/szdytom/LADRSolutions/blob/main/sections/1A.typ}{P-1A-8}（原书 1A 第 8 题；方而静题解）。
\begin{solve}
令 $x = a + b \mathrm i$（其中 $a , b \in \mathbf R$），满足 $x^{2} = \mathrm i$，即 $a^{2} - b^{2} + 2 a b \mathrm i = 0 + 1 \mathrm i$。因此，我们有两个方程

\[\begin{cases}a^{2} - b^{2} = 0\\2 a b = 1\end{cases}\]

解这两个方程组，我们得到 $a = \frac{\sqrt{2}}{2}$ 和 $b = \frac{\sqrt{2}}{2}$ 或者 $a = - \frac{\sqrt{2}}{2}$ 和 $b = - \frac{\sqrt{2}}{2}$。因此，$\mathrm i$ 的两个相异的平方根是：

\[x_{1} = \frac{\sqrt{2}}{2} + \frac{\sqrt{2}}{2} \mathrm i \text{，} \quad x_{2} = - \frac{\sqrt{2}}{2} - \frac{\sqrt{2}}{2} \mathrm i\]
\end{solve}

\begin{exercise}\label{prob:P-1A-9}
求 $x \in \mathbf R^{4}$，使得
\[(4 , - 3 , 1 , 7) + 2 x = (5 , 9 , - 6 , 8)\]
\end{exercise}
\noindent 来源：\href{https://github.com/szdytom/LADRSolutions/blob/main/sections/1A.typ}{P-1A-9}（原书 1A 第 9 题；方而静题解）。
\begin{solve}
根据定义，令 $x = (x_{1} , x_{2} , x_{3} , x_{4})$，则

\[(4 , - 3 , 1 , 7) + 2 x = (4 + 2 x_{1} , - 3 + 2 x_{2} , 1 + 2 x_{3} , 7 + 2 x_{4})\]

令其等于 $(5 , 9 , - 6 , 8)$，则我们有四个方程：

\[\begin{cases}4 + 2 x_{1} = 5\\- 3 + 2 x_{2} = 9\\1 + 2 x_{3} = - 6\\7 + 2 x_{4} = 8\end{cases}\]

解这四个方程，我们得到 $x_{1} = \frac{1}{2} , x_{2} = 6 , x_{3} = - \frac{7}{2} , x_{4} = \frac{1}{2}$。因此，

\[x = (\frac{1}{2} , 6 , - \frac{7}{2} , \frac{1}{2})\]
\end{solve}

\begin{exercise}\label{prob:P-1A-10}
解释为什么不存在 $\lambda \in \mathbf C$，使得
\[\lambda (2 - 3 \mathrm i , 5 + 4 \mathrm i , - 6 + 7 \mathrm i) = (12 - 5 \mathrm i , 7 + 22 \mathrm i , - 32 - 9 \mathrm i)\]
\end{exercise}
\noindent 来源：\href{https://github.com/szdytom/LADRSolutions/blob/main/sections/1A.typ}{P-1A-10}（原书 1A 第 10 题；方而静题解）。
\begin{solve}
注意到，

\[\frac{12 - 5 \mathrm i}{2 - 3 \mathrm i} = 3 + 2 \mathrm i\]

而

\[\frac{- 32 - 9 \mathrm i}{- 6 + 7 \mathrm i} \ne 3 + 2 \mathrm i\]

因此，这样的 $\lambda$ 不存在。
\end{solve}

\begin{exercise}\label{prob:P-1A-11}\label{app:E-ffn-add-assoc}
证明：$(x + y) + z = x + (y + z)$ 对所有 $x , y , z \in \mathbf F^{n}$ 成立。

\begin{remark}
沿用原书记号1.6与记号1.10，即 $\mathbf F$ 表示 $\mathbf R$ 或 $\mathbf C$，$n$ 表示某一固定的正整数。下文不再赘述。
\end{remark}
\end{exercise}
\noindent 来源：\href{https://github.com/szdytom/LADRSolutions/blob/main/sections/1A.typ}{P-1A-11}（原书 1A 第 11 题；方而静题解）。
\begin{solve}
根据定义，令 $x = (x_{1} , \ldots , x_{n})$，$y = (y_{1} , \ldots , y_{n})$，$z = (z_{1} , \ldots , z_{n})$，则有

\[\begin{aligned}(x + y) + z & = ((x_{1} + y_{1} , \ldots , x_{n} + y_{n}) + (z_{1} , \ldots , z_{n})) \\ & = (x_{1} + y_{1} + z_{1} , \ldots , x_{n} + y_{n} + z_{n}) \\ & = ((x_{1} , \ldots , x_{n}) + (y_{1} + z_{1} , \ldots , y_{n} + z_{n})) \\ & = x + (y + z)\end{aligned}\]
\end{solve}

\begin{exercise}\label{prob:P-1A-12}\label{app:E-ffn-mul-assoc}
证明：$(a b) x = a (b x)$ 对所有 $x \in \mathbf F^{n}$ 和 $a , b \in \mathbf F$ 成立。
\end{exercise}
\noindent 来源：\href{https://github.com/szdytom/LADRSolutions/blob/main/sections/1A.typ}{P-1A-12}（原书 1A 第 12 题；方而静题解）。
\begin{solve}
设 $x=(x_1,\ldots,x_n)$。逐坐标使用标量乘法的结合律可得
\[(ab)x=((ab)x_1,\ldots,(ab)x_n)=(a(bx_1),\ldots,a(bx_n))=a(bx).\]
\end{solve}

\begin{exercise}\label{prob:P-1A-13}\label{app:E-ffn-mul-scalar-id}
证明：$1 x = x$ 对所有 $x \in \mathbf F^{n}$ 成立。
\end{exercise}
\noindent 来源：\href{https://github.com/szdytom/LADRSolutions/blob/main/sections/1A.typ}{P-1A-13}（原书 1A 第 13 题；方而静题解）。
\begin{solve}
根据定义，令 $x = (x_{1} , \ldots , x_{n})$，则有

\[\begin{aligned}1 x & = 1 (x_{1} , \ldots , x_{n}) \\ & = (1 \cdot x_{1} , \ldots , 1 \cdot x_{n}) \\ & = (x_{1} , \ldots , x_{n}) \\ & = x\end{aligned}\]
\end{solve}

\begin{exercise}\label{prob:P-1A-14}\label{app:E-ffn-distributivity-vector-add}
证明：$\lambda (x + y) = \lambda x + \lambda y$ 对所有 $\lambda \in \mathbf F$ 和 $x , y \in \mathbf F^{n}$ 成立。
\end{exercise}
\noindent 来源：\href{https://github.com/szdytom/LADRSolutions/blob/main/sections/1A.typ}{P-1A-14}（原书 1A 第 14 题；方而静题解）。
\begin{solve}
根据定义，令 $x = (x_{1} , \ldots , x_{n})$，$y = (y_{1} , \ldots , y_{n})$，则有

\[\begin{aligned}\lambda (x + y) & = \lambda ((x_{1} + y_{1} , \ldots , x_{n} + y_{n})) \\ & = (\lambda (x_{1} + y_{1}) , \ldots , \lambda (x_{n} + y_{n})) \\ & = (\lambda x_{1} + \lambda y_{1} , \ldots , \lambda x_{n} + \lambda y_{n}) \\ & = \lambda x + \lambda y\end{aligned}\]
\end{solve}

\begin{exercise}\label{prob:P-1A-15}\label{app:E-ffn-distributivity-scalar-add}
证明：$(a + b) x = a x + b x$ 对所有 $a , b \in \mathbf F$ 和 $x \in \mathbf F^{n}$ 成立。
\end{exercise}
\noindent 来源：\href{https://github.com/szdytom/LADRSolutions/blob/main/sections/1A.typ}{P-1A-15}（原书 1A 第 15 题；方而静题解）。
\begin{solve}
根据定义，令 $x = (x_{1} , \ldots , x_{n})$，则有

\[\begin{aligned}(a + b) x & = (a + b) (x_{1} , \ldots , x_{n}) \\ & = (a x_{1} + b x_{1} , \ldots , a x_{n} + b x_{n}) \\ & = (a x_{1} , \ldots , a x_{n}) + (b x_{1} , \ldots , b x_{n}) \\ & = a x + b x\end{aligned}\]
\end{solve}

\subsection{向量空间的定义}\label{app:section-1B}
\begin{exercise}\label{prob:P-1B-1}
证明：$- (- v) = v$ 对任意 $v \in V$ 都成立。

\begin{remark}
沿用原书记号1.29，即 $V$ 表示 $\mathbf F$ 上的向量空间。下文不再赘述。
\end{remark}
\end{exercise}
\noindent 来源：\href{https://github.com/szdytom/LADRSolutions/blob/main/sections/1B.typ}{P-1B-1}（原书 1B 第 1 题；方而静题解）。
\begin{solve}
根据定义 $v + (- v) = 0$，由可交换性，得 $(- v) + v = 0$，即 $v$ 是 $- v$ 的加法逆元。由加法逆元的唯一性（原书1.27），得 $- (- v) = v$。
\end{solve}

\begin{exercise}\label{prob:P-1B-2}\label{app:E-vector-trivial-annihilation}
设 $a \in \mathbf F$，$v \in V$ 且 $a v = 0$，证明：$a = 0$ 或 $v = 0$。
\end{exercise}
\noindent 来源：\href{https://github.com/szdytom/LADRSolutions/blob/main/sections/1B.typ}{P-1B-2}（原书 1B 第 2 题；方而静题解）。
\begin{solve}
我们使用反证法，假设 $a \ne 0$ 且 $v \ne 0$，同时 $a v = 0$。由于 $a \ne 0$ 故存在 $a^{- 1} \in \mathbf F$（$a^{- 1} \ne 0$），使得 $a^{- 1} a = 1$。因此有

\[1 v = a^{- 1} a v = a^{- 1} (a v) = a^{- 1} 0 = 0\]

这与 $v \ne 0$ 矛盾，假设不成立。因此，$a = 0$ 或 $v = 0$。
\end{solve}

\begin{exercise}\label{prob:P-1B-3}
设 $v,w\in V$，解释为什么存在唯一的 $x\in V$ 使得 $v+3x=w$。
\end{exercise}
\noindent 来源：\href{https://github.com/szdytom/LADRSolutions/blob/main/sections/1B.typ}{P-1B-3}（原书 1B 第 3 题；方而静题解）。
\begin{solve}
取 $x = \frac{1}{3} (w - v)$，则有 $v + 3 x = v + 3 (\frac{1}{3} (w - v)) = v + (w - v) = w$。由于向量空间的加法和数乘满足封闭性，因此这样的 $x$ 存在于 $V$ 中。

为了说明其唯一性，假设存在另一个 $x' \in V$ 使得 $v + 3 x' = w$，则有

\[v + 3 x = w = v + 3 x'\]

在等式两边同时加上 $- v$，化简得 $3 x = 3 x'$，等式两边同时乘以 $\frac{1}{3}$，得 $x = x'$。因此，这样的 $x$ 是唯一的。
\end{solve}

\begin{exercise}\label{prob:P-1B-4}
空集不是向量空间。对于在向量空间的定义（原书1.20）中列出的要求，空集仅不满足其中的一条。是哪一条？
\end{exercise}
\noindent 来源：\href{https://github.com/szdytom/LADRSolutions/blob/main/sections/1B.typ}{P-1B-4}（原书 1B 第 4 题；方而静题解）。
\begin{solve}
空集中不存在加法单位元。
\end{solve}

\begin{exercise}\label{prob:P-1B-5}
证明：在向量空间的定义（原书1.20）中，加法逆元条件可以替换为这一条件——

$0 v = 0$ 对所有 $v \in V$ 成立。

这里，左侧的 $0$ 是数 $0$，而右侧的 $0$ 是 $V$ 中的加法单位元。

\begin{remark}
在定义中“条件可以替换”，指原来的条件替换成新条件后，满足定义的对象还是原来的那些。
\end{remark}
\end{exercise}
\noindent 来源：\href{https://github.com/szdytom/LADRSolutions/blob/main/sections/1B.typ}{P-1B-5}（原书 1B 第 5 题；方而静题解）。
\begin{solve}
采用原有定义时，新条件成立的证明由原书1.30给出。我们现在采用替换后的新定义，并以此证明加法逆元条件，即“对于任意 $v \in V$，都存在 $w \in V$ 使得 $v + w = 0$”。

更具体地，我们说明 $v + (- 1) v = 0$。这是由于

\[v + (- 1) v = 1 v + (- 1) v = (1 + (- 1)) v = 0 v = 0\]

所以对于任意 $v \in V$，它的加法逆元都存在，即 $(- 1) v$。故两个条件可以相互替换。
\end{solve}

\begin{exercise}\label{prob:P-1B-6}
令 $\infty$ 和 $- \infty$ 是不在 $\mathbf R$ 中的不同对象。以最符合直觉的方式定义 $\mathbf R \cup \{\infty , - \infty\}$ 上的加法和标量乘法。具体而言，两个实数的和和积照常定义，而对于 $t \in \mathbf R$，我们定义

\[t \infty = \begin{cases}- \infty \quad & t < 0\\0 & t = 0\\\infty & t > 0\end{cases} \quad t (- \infty) = \begin{cases}\infty \quad & t < 0\\0 & t = 0\\- \infty & t > 0\end{cases}\]

以及

\[\begin{aligned}t + \infty & = \infty + t = \infty + \infty = \infty \\ t + (- \infty) & = (- \infty) + t = (- \infty) + (- \infty) = - \infty \\ \infty + (- \infty) & = (- \infty) + \infty = 0\end{aligned}\]

具有这样的加法和标量乘法的 $\mathbf R \cup \{\infty , - \infty\}$ 是 $\mathbf R$ 上的向量空间吗？解释一下。
\end{exercise}
\noindent 来源：\href{https://github.com/szdytom/LADRSolutions/blob/main/sections/1B.typ}{P-1B-6}（原书 1B 第 6 题；方而静题解）。
\begin{solve}
不是。任取 $t \in \mathbf R$（$t \ne 0$），注意到，

\[\begin{aligned}(t + \infty) + (- \infty) & = \infty + (- \infty) = 0 \\ t + (\infty + (- \infty)) & = t + 0 = t\end{aligned}\]

这违背了加法的可结合性要求。因此，$\mathbf R \cup \{\infty , - \infty\}$ 不是 $\mathbf R$ 上的向量空间。
\end{solve}

\begin{exercise}\label{prob:P-1B-7}
设 $S$ 是非空集合，令 $V^{S}$ 表示所有从 $S$ 到 $V$ 的函数构成的集合。请在 $V^{S}$ 定义一种自然的加法和标量乘法，并证明：具有这些定义的 $V^{S}$ 是向量空间。
\end{exercise}
\noindent 来源：\href{https://github.com/szdytom/LADRSolutions/blob/main/sections/1B.typ}{P-1B-7}（原书 1B 第 7 题；方而静题解）。
\begin{solve}
我们按如下方式定义 $V^{S}$ 上的加法和标量乘法：

\begin{itemize}
\item 对于 $f , g \in V^{S}$，和 $f + g \in V^{S}$ 是由下式定义的函数：对于任意 $x \in S$，
\[(f + g) (x) = f (x) + g (x)\]
\item 对于 $\lambda \in \mathbf F$ 和 $f \in V^{S}$，乘积 $\lambda f \in V^{S}$ 是由下式定义的函数：对于任意 $x \in S$，
\[(\lambda f) (x) = \lambda f (x)\]

我们现在证明 $V^{S}$ 是 $\mathbf F$ 上的向量空间。具体而言，我们逐条验证向量空间的定义（原书1.20）中的要求：

\end{itemize}
\begin{description}
\item[可交换性] 对于任意 $f , g \in V^{S}$，都有 $f + g = g + f$。
证明：设 $x \in S$，有
\[(f + g) (x) = f (x) + g (x) = g (x) + f (x) = (g + f) (x)\]
因此 $f + g = g + f$。

\item[可结合性] 对于任意 $f , g , h \in V^{S}$ 以及 $a , b \in \mathbf F$，都有 $(f + g) + h = f + (g + h)$ 且 $(a b) f = a (b f)$。
证明：设 $x \in S$，则对于加法的结合性，有
\[\begin{aligned}((f + g) + h) (x) & = (f + g) (x) + h (x) \\ & = f (x) + g (x) + h (x) \\ & = f (x) + (g + h) (x) \\ & = (f + (g + h)) (x)\end{aligned}\]
因此 $(f + g) + h = f + (g + h)$。对于标量乘法的结合性，有
\[\begin{aligned}((a b) f) (x) & = (a b) (f (x)) \\ & = a (b f (x)) \\ & = a ((b f) (x)) = (a (b f)) (x)\end{aligned}\]
因此 $(a b) f = a (b f)$。

\item[加法单位元] 存在 $0 \in V^{S}$ 使得对于任意 $f \in V^{S}$，都有 $f + 0 = f$。
证明：取 $0 : x \mapsto 0$ 为 $V^{S}$ 中的加法单位元。设 $x \in S$，对于任意 $f \in V^{S}$，都有
\[(f + 0) (x) = f (x) + 0 = f (x) = (f) (x)\]
因此 $f + 0 = f$。

\item[加法逆元] 对于任意 $f \in V^{S}$，存在 $g \in V^{S}$ 使得 $f + g = 0$。
证明：取 $g : x \mapsto - f (x)$ 为 $f$ 的加法逆元。设 $x \in S$，都有
\[(f + g) (x) = f (x) + g (x) = f (x) - f (x) = 0\]
因此 $f + g = 0$。

\item[乘法单位元] 对于任意 $f \in V^{S}$，都有 $1 f = f$。
证明：设 $x \in S$，都有
\[(1 f) (x) = 1 f (x) = f (x)\]
因此 $1 f = f$。

\item[分配性质] 对于任意 $f , g \in V^{S}$ 以及所有 $a , b \in \mathbf F$，都有 $a (f + g) = a f + a g$ 且 $(a + b) f = a f + b f$。
证明：设 $x \in S$，则对于第一个分配性质，有
\[\begin{aligned}(a (f + g)) (x) & = a ((f + g) (x)) = a (f (x) + g (x)) \\ & = a f (x) + a g (x) = (a f) (x) + (a g) (x) \\ & = (a f + a g) (x)\end{aligned}\]
因此 $a (f + g) = a f + a g$。对于第二个分配性质，有
\[\begin{aligned}((a + b) f) (x) & = (a + b) f (x) \\ & = a f (x) + b f (x) \\ & = (a f) (x) + (b f) (x) \\ & = (a f + b f) (x)\end{aligned}\]
因此 $(a + b) f = a f + b f$。

综上所述，$V^{S}$ 满足向量空间的所有要求，因此 $V^{S}$ 是 $\mathbf F$ 上的向量空间。
\end{description}
\end{solve}

\begin{exercise}\label{prob:P-1B-8}\label{app:E-vector-dspace-complexification}
设 $V$ 是实向量空间。

\begin{itemize}
\item $V$ 的\emph{复化（complexification）}记为 $V_{\mathbf C}$，等于 $V \times V$。$V_{\mathbf C}$ 中的所有元素为有序对 $(u , v)$，其中 $u , v \in V$，不过我们将其记作 $u + \mathrm i v$。

\item $V_{\mathbf C}$ 上的加法定义为
\[(u_{1} + \mathrm i v_{1}) + (u_{2} + \mathrm i v_{2}) = (u_{1} + u_{2}) + \mathrm i (v_{1} + v_{2})\]
对所有 $u_{1} , v_{1} , u_{2} , v_{2} \in V$ 都成立。

\item $V_{\mathbf C}$ 上的标量乘法定义为
\[(a + b \mathrm i) (u + \mathrm i v) = (a u - b v) + \mathrm i (a v + b u)\]
对所有 $a , b \in \mathbf R$ 和所有 $u , v \in V$ 都成立。

\end{itemize}
证明：具有如上加法和标量乘法定义的 $V_{\mathbf C}$ 是向量空间。

\begin{remark}
将 $u \in V$ 等同于 $u + \mathrm i 0$，从而将 $V$ 视为 $V_{\mathbf C}$ 的一个子集。这样一来，由 $V$ 构造 $V_{\mathbf C}$ 就可以视作由 $\mathbf R^{n}$ 构造 $\mathbf C^{n}$ 的推广。
\end{remark}
\end{exercise}
\noindent 来源：\href{https://github.com/szdytom/LADRSolutions/blob/main/sections/1B.typ}{P-1B-8}（原书 1B 第 8 题；方而静题解）。
\begin{solve}
把 $u+\mathrm iv$ 看作有序对 $(u,v)$，则加法就是直积空间 $V\times V$ 的逐分量加法，因此满足交换律、结合律，零元为 $(0,0)$，逆元为 $(-u,-v)$。对 $\lambda=a+b\mathrm i\in\mathbf C$，记
\[M(\lambda)=\begin{pmatrix}a&-b\\b&a\end{pmatrix}.\]
题设数乘正是用实系数矩阵 $M(\lambda)$ 作用于 $(u,v)$。若 $\mu=c+d\mathrm i$，则直接相乘得
\[M(\lambda)M(\mu)=\begin{pmatrix}ac-bd&-(ad+bc)\\ad+bc&ac-bd\end{pmatrix}=M(\lambda\mu).\]
此外，$M(\lambda+\mu)=M(\lambda)+M(\mu)$ 且 $M(1)=I_2$。这些恒等式分别给出数乘结合律、对标量加法的分配律和单位元性质；矩阵作用对 $(u,v)$ 的加法保持线性，给出另一条分配律。所有运算均封闭，故 $V_{\mathbf C}$ 是复向量空间。
\end{solve}

\subsection{子空间}\label{app:section-1C}
\begin{exercise}\label{prob:P-1C-1}
对于 $\mathbf F^{3}$ 的下列各个子集，判断其是否是 $\mathbf F^{3}$ 的子空间：

\begin{enumerate}[label=(\alph*)]
\item $\{(x_{1} , x_{2} , x_{3}) \in \mathbf F^{3} : x_{1} + 2 x_{2} + 3 x_{3} = 0\}$
\item $\{(x_{1} , x_{2} , x_{3}) \in \mathbf F^{3} : x_{1} + 2 x_{2} + 3 x_{3} = 4\}$
\item $\{(x_{1} , x_{2} , x_{3}) \in \mathbf F^{3} : x_{1} x_{2} x_{3} = 0\}$
\item $\{(x_{1} , x_{2} , x_{3}) \in \mathbf F^{3} : x_{1} = 5 x_{3}\}$
\end{enumerate}
\end{exercise}
\noindent 来源：\href{https://github.com/szdytom/LADRSolutions/blob/main/sections/1C.typ}{P-1C-1}（原书 1C 第 1 题；方而静题解）。
\begin{solve}
方便起见，将这四个集合分别记作 $S_{1} , S_{2} , S_{3} , S_{4}$。为了验证它们是否是 $\mathbf F^{3}$ 的子空间，我们需要验证它们是否满足子空间的条件（原书1.34）。

对于 $S_{1}$，我们验证以下三个条件：

\begin{description}
\item[加法单位元] $0 \in S_{1}$。
证明：注意到 $0 + 2 \cdot 0 + 3 \cdot 0 = 0$，故 $0 \in S_{1}$。

\item[加法封闭性] $u , w \in S_{1}$ 意味着 $u + w \in S_{1}$。
证明：设 $u = (u_{1} , u_{2} , u_{3})$，$w = (w_{1} , w_{2} , w_{3})$，则 $u + w = (u_{1} + w_{1} , u_{2} + w_{2} , u_{3} + w_{3})$。由于 $u_{1} + 2 u_{2} + 3 u_{3} = 0$ 且 $w_{1} + 2 w_{2} + 3 w_{3} = 0$，因此
\[\begin{aligned}(u + w)_{1} + 2 (u + w)_{2} + 3 (u + w)_{3} & = (u_{1} + w_{1}) + 2 (u_{2} + w_{2}) + 3 (u_{3} + w_{3}) \\ & = (u_{1} + 2 u_{2} + 3 u_{3}) + (w_{1} + 2 w_{2} + 3 w_{3}) \\ & = 0 + 0 = 0\end{aligned}\]
故 $u + w \in S_{1}$。

\item[数乘封闭性] $a \in \mathbf F$ 且 $u \in S_{1}$ 意味着 $a u \in S_{1}$。
证明：设 $u = (u_{1} , u_{2} , u_{3})$，则 $a u = (a u_{1} , a u_{2} , a u_{3})$。由于 $u_{1} + 2 u_{2} + 3 u_{3} = 0$，因此
\[\begin{aligned}(a u)_{1} + 2 (a u)_{2} + 3 (a u)_{3} & = a u_{1} + 2 a u_{2} + 3 a u_{3} \\ & = a (u_{1} + 2 u_{2} + 3 u_{3}) \\ & = a \cdot 0 = 0\end{aligned}\]
故 $a u \in S_{1}$。

综上所述，$S_{1}$ 是 $\mathbf F^{3}$ 的子空间。

对于 $S_{2}$，注意到 $0 + 2 \cdot 0 + 3 \cdot 0 = 0 \ne 4$，故 $0 \notin S_{2}$。这违反了“加法单位元”的要求，说明 $S_{2}$ 不是 $\mathbf F^{3}$ 的子空间。

对于 $S_{3}$，取 $u = (1 , 1 , 0) \in S_{3}$，$v = (0 , 0 , 1) \in S_{3}$，注意到 $u + v = (1 , 1 , 1) \notin S_{3}$，这违反了“加法封闭性”的要求，说明 $S_{3}$ 不是 $\mathbf F^{3}$ 的子空间。

对于 $S_{4}$，我们验证以下三个条件：

\item[加法单位元] $0 \in S_{4}$。
证明：注意到 $0 = 5 \cdot 0$，故 $0 \in S_{4}$。

\item[加法封闭性] $u , w \in S_{4}$ 意味着 $u + w \in S_{4}$。
证明：设 $u = (u_{1} , u_{2} , u_{3})$，$w = (w_{1} , w_{2} , w_{3})$，则 $u + w = (u_{1} + w_{1} , u_{2} + w_{2} , u_{3} + w_{3})$。由于 $u_{1} = 5 u_{3}$ 且 $w_{1} = 5 w_{3}$，因此
\[u_{1} + w_{1} = 5 u_{3} + 5 w_{3} = 5 (u_{3} + w_{3})\]
故 $u + w \in S_{4}$。

\item[数乘封闭性] $a \in \mathbf F$ 且 $u \in S_{4}$ 意味着 $a u \in S_{4}$。
证明：设 $u = (u_{1} , u_{2} , u_{3})$，则 $a u = (a u_{1} , a u_{2} , a u_{3})$。由于 $u_{1} = 5 u_{3}$，因此
\[a u_{1} = a (5 u_{3}) = 5 (a u_{3})\]
故 $a u \in S_{4}$。

综上所述，$S_{4}$ 是 $\mathbf F^{3}$ 的子空间。
\end{description}
\end{solve}

\begin{exercise}\label{prob:P-1C-2}
验证下面这些有关子空间的结论。

\begin{enumerate}[label=(\alph*)]
\item 如果 $b \in \mathbf F$，那么当且仅当 $b = 0$ 时，
\[\{(x_{1} , x_{2} , x_{3} , x_{4}) \in \mathbf F^{4} : x_{3} = 5 x_{4} + b\}\]
是 $\mathbf F^{4}$ 的子空间；
\item 定义在区间 $[0 , 1]$ 上的全体连续实值函数构成的集合是 $\mathbf R^{[0 , 1]}$ 的子空间；
\item 定义在 $\mathbf R$ 上的全体可微实值函数构成的集合是 $\mathbf R^{\mathbf R}$ 的子空间；
\item 当且仅当 $b = 0$ 时，定义在区间 $(0 , 3)$ 上且满足 $f' (2) = b$ 的全体可微实值函数 $f$ 构成的集合是 $\mathbf R^{(0 , 3)}$ 的子空间；
\item 极限为 $0$ 的所有复数序列所构成的集合是 $\mathbf C^{\infty}$ 的子空间。

\end{enumerate}
\begin{remark}
本题原文为“验证例1.35中关于子空间的所有结论”。出于完整性考虑，这里将原书例1.35的所有结论摘录在上面。
\end{remark}
\end{exercise}
\noindent 来源：\href{https://github.com/szdytom/LADRSolutions/blob/main/sections/1C.typ}{P-1C-2}（原书 1C 第 2 题；方而静题解）。
\begin{solve}
使用子空间判别法：非空子集是子空间，当且仅当对任意两个元素及任意两个标量的线性组合封闭。
\begin{enumerate}[label=(\alph*)]
\item 零向量满足 $x_3=5x_4+b$ 当且仅当 $b=0$；此时条件是齐次线性方程，解集对线性组合封闭。
\item 零函数连续，连续函数的线性组合仍连续，故为子空间。
\item 零函数可微，可微函数的线性组合仍可微，故为子空间。
\item 零函数满足 $f'(2)=b$ 当且仅当 $b=0$；此时由 $(af+cg)'(2)=af'(2)+cg'(2)$ 可知解集对线性组合封闭。
\item 零序列极限为零；若 $z_n\to0$、$w_n\to0$，则对任意 $a,c\in\mathbf C$，有 $az_n+cw_n\to0$，故为子空间。
\end{enumerate}
\end{solve}

\begin{exercise}\label{prob:P-1C-3}
证明：在区间 $(- 4 , 4)$ 上的满足 $f' (- 1) = 3 f (2)$ 的可微实值函数 $f$ 构成的集合是 $\mathbf R^{(- 4 , 4)}$ 的子空间。
\end{exercise}
\noindent 来源：\href{https://github.com/szdytom/LADRSolutions/blob/main/sections/1C.typ}{P-1C-3}（原书 1C 第 3 题；方而静题解）。
\begin{solve}
记 $S$ 为题目所说的子集，即 $S = \{f \in \mathbf R^{(- 4 , 4)} : f' (- 1) = 3 f (2)\}$。我们逐条验证其满足子空间的条件（原书1.34）：

\begin{description}
\item[加法单位元] $0 \in S$。
证明：注意到 $0' (- 1) = 0 = 3 \cdot 0 = 3 \cdot 0 (2)$，故 $0 \in S$。

\item[加法封闭性] $u , w \in S$ 意味着 $u + w \in S$。
证明：设 $u , w \in S$，则 $u' (- 1) = 3 u (2)$ 且 $w' (- 1) = 3 w (2)$。因此
\[(u + w)' (- 1) = u' (- 1) + w' (- 1) = 3 u (2) + 3 w (2) = 3 (u (2) + w (2))\]
故 $u + w \in S$。

\item[数乘封闭性] $a \in \mathbf R$ 且 $u \in S$ 意味着 $a u \in S$。
证明：设 $u \in S$，则 $u' (- 1) = 3 u (2)$。因此
\[(a u)' (- 1) = a u' (- 1) = a \cdot 3 u (2) = 3 (a u (2))\]
故 $a u \in S$。

综上所述，满足条件的集合是 $\mathbf R^{(- 4 , 4)}$ 的子空间。
\end{description}
\end{solve}

\begin{exercise}\label{prob:P-1C-4}
设 $b \in \mathbf R$，证明：在区间 $[0 , 1]$ 上满足 $\int_{0}^{1} f = b$ 的连续实值函数 $f$ 构成的集合是 $\mathbf R^{[0 , 1]}$ 的子空间，当且仅当 $b = 0$。
\end{exercise}
\noindent 来源：\href{https://github.com/szdytom/LADRSolutions/blob/main/sections/1C.typ}{P-1C-4}（原书 1C 第 4 题；方而静题解）。
\begin{solve}
我们首先说明其充分性。假设 $b = 0$，此时 $S = \{f \in \mathbf R^{[0 , 1]} : \int_{0}^{1} f = 0\}$，我们逐条验证其满足子空间的条件（原书1.34）：

\begin{description}
\item[加法单位元] $0 \in S$。
证明：注意到 $\int_{0}^{1} 0 = 0$，故 $0 \in S$。

\item[加法封闭性] $u , w \in S$ 意味着 $u + w \in S$。
证明：设 $u , w \in S$，则
\[\int_{0}^{1} (u + w) = \int_{0}^{1} u + \int_{0}^{1} w = 0 + 0 = 0\]
故 $u + w \in S$。

\item[数乘封闭性] $a \in \mathbf R$ 且 $u \in S$ 意味着 $a u \in S$。
证明：设 $u \in S$，则
\[\int_{0}^{1} (a u) = a \int_{0}^{1} u = a \cdot 0 = 0\]
故 $a u \in S$。

综上所述，当 $b = 0$ 时，$S$ 是 $\mathbf R^{[0 , 1]}$ 的子空间。
\end{description}
\end{solve}

\begin{exercise}\label{prob:P-1C-5}
$\mathbf R^{2}$ 是不是复向量空间 $\mathbf C^{2}$ 的子空间？
\end{exercise}
\noindent 来源：\href{https://github.com/szdytom/LADRSolutions/blob/main/sections/1C.typ}{P-1C-5}（原书 1C 第 5 题；方而静题解）。
\begin{solve}
不是。注意到，取 $v = (1 , 0) \in \mathbf R^{2}$，则

\[\mathrm i v = (\mathrm i , 0) \notin \mathbf R^{2}\]

这违反了子空间的条件（原书1.34）中对“数乘封闭性”的要求。由此，$\mathbf R^{2}$ 不是 $\mathbf C^{2}$ 的子空间。
\end{solve}

\begin{exercise}\label{prob:P-1C-6}
\begin{enumerate}[label=(\alph*)]
\item $\{(a , b , c) \in \mathbf R^{3} : a^{3} = b^{3}\}$ 是不是 $\mathbf R^{3}$ 的子空间？
\item $\{(a , b , c) \in \mathbf C^{3} : a^{3} = b^{3}\}$ 是不是 $\mathbf C^{3}$ 的子空间？
\end{enumerate}
\end{exercise}
\noindent 来源：\href{https://github.com/szdytom/LADRSolutions/blob/main/sections/1C.typ}{P-1C-6}（原书 1C 第 6 题；方而静题解）。
\begin{solve}
我们首先来看第一个集合 $S_{\mathbf R} = \{(a , b , c) \in \mathbf R^{3} : a^{3} = b^{3}\}$。当 $a^{3} = b^{3}$ 时，$a^{3} - b^{3} = 0$，即

\[(a - b) (a^{2} + a b + b^{2}) = 0\]

当 $a b > 0$ 时，

\[a^{2} + a b + b^{2} > a^{2} - 2 a b + b^{2} = (a - b)^{2} \ge 0\]

当 $a b < 0$ 时，

\[a^{2} + a b + b^{2} > a^{2} + 2 a b + b^{2} = (a + b)^{2} \ge 0\]

即当 $a \ne 0$ 且 $b \ne 0$ 时，$a^{2} + a b + b^{2} > 0$。此时，必然有 $a - b = 0$，即 $a = b$。而 $a = b = 0$ 时，自然也有 $a = b$。综上所述，$a^{3} = b^{3}$ 意味着 $a = b$。

由此，我们可以将 $S_{\mathbf R}$ 重写为 $S_{\mathbf R} = \{(a , b , c) \in \mathbf R^{3} : a = b\}$。我们逐条验证其满足子空间的条件（原书1.34）：

\begin{description}
\item[加法单位元] $0 \in S_{\mathbf R}$。
证明：注意到 $0 = 0$，故 $0 \in S_{\mathbf R}$。

\item[加法封闭性] $u , w \in S_{\mathbf R}$ 意味着
$u + w \in S_{\mathbf R}$。
证明：设 $u = (u_{1} , u_{2} , u_{3})$，$w = (w_{1} , w_{2} , w_{3})$，
则 $u + w = (u_{1} + w_{1} , u_{2} + w_{2} , u_{3} + w_{3})$。由于 $u_{1} = u_{2}$ 且 $w_{1} = w_{2}$，因此
\[u_{1} + w_{1} = u_{2} + w_{2}\]
故 $u + w \in S_{\mathbf R}$。

\item[数乘封闭性] $a \in \mathbf R$ 且 $u \in S_{\mathbf R}$ 意味着
$a u \in S_{\mathbf R}$。
证明：设 $u = (u_{1} , u_{2} , u_{3})$，则 $a u = (a u_{1} , a u_{2} , a u_{3})$。由于 $u_{1} = u_{2}$，因此
\[a u_{1} = a u_{2}\]
故 $a u \in S_{\mathbf R}$。

综上所述，$S_{\mathbf R}$ 是 $\mathbf R^{3}$ 的子空间。

对于第二个集合 $S_{\mathbf C} = \{(a , b , c) \in \mathbf C^{3} : a^{3} = b^{3}\}$，注意到，取

\end{description}
\[u = (\frac{- 1 + \sqrt{3} \mathrm i}{2} , 1 , 0) \text{， } v = (\frac{- 1 - \sqrt{3} \mathrm i}{2} , 1 , 0)\]

容易验证 $u , v \in S_{\mathbf C}$，而

\[u + v = (- 1 , 2 , 0) \notin S_{\mathbf C}\]

这违反了子空间的条件（原书1.34）中对“加法封闭性”的要求。由此，$S_{\mathbf C}$ 不是 $\mathbf C^{3}$ 的子空间。
\end{solve}

\begin{exercise}\label{prob:P-1C-7}
证明或证伪：如果 $U$ 是 $\mathbf R^{2}$ 的非空子集，满足对加法封闭和对“取加法逆元”封闭（即 $u \in U$ 意味着 $- u \in U$），那么 $U$ 是 $\mathbf R^{2}$ 的子空间。
\end{exercise}
\noindent 来源：\href{https://github.com/szdytom/LADRSolutions/blob/main/sections/1C.typ}{P-1C-7}（原书 1C 第 7 题；方而静题解）。
\begin{solve}
取 $U = \{(1 , 0) , (0 , 0) , (- 1 , 0)\}$，容易验证 $U$ 满足对加法封闭和对“取加法逆元”封闭。但是，取 $u = (1 , 0) \in U$, $2 u = (2 , 0) \notin U$，这违反了子空间的条件（原书1.34）中对“数乘封闭性”的要求。由此，$U$ 不是 $\mathbf R^{2}$ 的子空间。

我们找到了一个反例，这说明题目中的命题不成立。
\end{solve}

\begin{exercise}\label{prob:P-1C-8}
给出一例：$\mathbf R^{2}$ 的非空子集 $U$，满足对标量数乘封闭，但不是 $\mathbf R^{2}$ 的子空间。
\end{exercise}
\noindent 来源：\href{https://github.com/szdytom/LADRSolutions/blob/main/sections/1C.typ}{P-1C-8}（原书 1C 第 8 题；方而静题解）。
\begin{solve}
取

\[U = \{(a , 0) : a \in \mathbf R\} \cup \{(0 , a) : a \in \mathbf R\} \subseteq \mathbf R^{2}\]

这个集合满足对标量数乘封闭，但不满足对加法封闭。比如，取 $u = (1 , 0) \in U$，$v = (0 , 1) \in U$，则 $u + v = (1 , 1) \notin U$。因此 $U$ 不是 $\mathbf R^{2}$ 的子空间。
\end{solve}

\begin{exercise}\label{prob:P-1C-9}
函数 $f : \mathbf R \to \mathbf R$ 被称为\emph{周期的（periodic）}，是指存在一正数 $p$，使得 $f (x) = f (x + p)$ 对所有 $x \in \mathbf R$ 成立。$\mathbf R \to \mathbf R$ 上的周期函数构成的集合是不是 $\mathbf R^{\mathbf R}$ 的子空间？请作解释。
\end{exercise}
\noindent 来源：\href{https://github.com/szdytom/LADRSolutions/blob/main/sections/1C.typ}{P-1C-9}（原书 1C 第 9 题；方而静题解）。
\begin{solve}
不是。取 $f (x) = \sin(x)$，$g (x) = \sin(\sqrt{2} x)$。容易验证，对于任意 $x \in \mathbf R$，

\[\begin{aligned}f (x + 2 \pi) & = f (x) \\ g (x + \sqrt{2} \pi) & = g (x)\end{aligned}\]

因此 $f$ 和 $g$ 都是 $\mathbf R^{\mathbf R}$ 中的周期函数。现在令 $h = f + g$。下面说明 $h$ 不是周期函数。

使用反证法，假设存在实数 $p > 0$，满足 $h (x) = h (x + p)$ 对所有 $x \in \mathbf R$ 成立，即

\begin{equation}\label{app:eq:1C-h-periodic-assume-eq}\sin(x) + \sin(\sqrt{2} x) = \sin(x + p) + \sin(\sqrt{2} x + \sqrt{2} p)\end{equation}

对式~\eqref{app:eq:1C-h-periodic-assume-eq} 两边同时求导两次，得到

\begin{equation}\label{app:eq:1C-h-periodic-assume-eq-dd}- \sin(x) - 2 \sin(\sqrt{2} x) = - \sin(x + p) - 2 \sin(\sqrt{2} x + \sqrt{2} p)\end{equation}

将式~\eqref{app:eq:1C-h-periodic-assume-eq} 与式~\eqref{app:eq:1C-h-periodic-assume-eq-dd} 相加并化简，得到

\begin{equation}\label{app:eq:1C-h-periodic-assume-eq-res-sqrt2}\sin(\sqrt{2} x) = \sin(\sqrt{2} x + \sqrt{2} p)\end{equation}

进一步将式~\eqref{app:eq:1C-h-periodic-assume-eq} 减去式~\eqref{app:eq:1C-h-periodic-assume-eq-res-sqrt2}，得到

\begin{equation}\label{app:eq:1C-h-periodic-assume-eq-res-1}\sin(x) = \sin(x + p)\end{equation}

向式~\eqref{app:eq:1C-h-periodic-assume-eq-res-sqrt2} 与式~\eqref{app:eq:1C-h-periodic-assume-eq-res-1} 中代入 $x = 0$，得到

\[\sin(p) = \sin(\sqrt{2} p) = 0\]

这意味着存在 $k_{1} , k_{2} \in \mathbf Z$，使得 $\sqrt{2} p = 2 k_{1} \pi$ 且 $p = 2 k_{2} \pi$。联立消去 $p$，得到 $\sqrt{2} = k_{1} / k_{2}$，这与我们熟知的 $\sqrt{2} \notin \mathbf Q$ 矛盾，故假设不成立。

综上所述，$h$ 不是 $\mathbf R^{\mathbf R}$ 上的周期函数。这表明 $\mathbf R \to \mathbf R$ 上的周期函数构成的集合并不符合子空间的条件（原书1.34）中对“加法封闭性”的要求，因此其不是 $\mathbf R^{\mathbf R}$ 的子空间。
\end{solve}

\begin{exercise}\label{prob:P-1C-10}\label{app:E-inter-of-subspace-is-subspace}
设 $V_{1}$ 和 $V_{2}$ 都是 $V$ 的子空间，证明：交集 $V_{1} \cap V_{2}$ 是 $V$ 的子空间。
\end{exercise}
\noindent 来源：\href{https://github.com/szdytom/LADRSolutions/blob/main/sections/1C.typ}{P-1C-10}（原书 1C 第 10 题；方而静题解）。
\begin{solve}
记 $S = V_{1} \cap V_{2}$，我们逐条验证其满足子空间的条件（原书1.34）：

\begin{description}
\item[加法单位元] $0 \in S$。
证明：由于 $V_{1}$ 和 $V_{2}$ 都是 $V$ 的子空间，因此 $0 \in V_{1}$ 且 $0 \in V_{2}$，从而 $0 \in S$。

\item[加法封闭性] $u , w \in S$ 意味着 $u + w \in S$。
证明：设 $u , w \in S$，则
\[u \in V_{1} , w \in V_{1} , u \in V_{2} , w \in V_{2}\]
因此 $u + w \in V_{1}$ 且 $u + w \in V_{2}$，从而 $u + w \in S$。

\item[数乘封闭性] $a \in V$ 且 $u \in S$ 意味着 $a u \in S$。
证明：设 $u \in S$，则
\[u \in V_{1} , u \in V_{2}\]
因此 $a u \in V_{1}$ 且 $a u \in V_{2}$，从而 $a u \in S$。

综上所述，$V_{1} \cap V_{2}$ 是 $V$ 的子空间。
\end{description}
\end{solve}

\begin{exercise}\label{prob:P-1C-11}
证明：$V$ 的任意一族子空间的交集是 $V$ 的子空间。
\end{exercise}
\noindent 来源：\href{https://github.com/szdytom/LADRSolutions/blob/main/sections/1C.typ}{P-1C-11}（原书 1C 第 11 题；方而静题解）。
\begin{solve}
设 $V_{1} , \ldots , V_{n}$ 都是 $V$ 的子空间，记 $S = V_{1} \cap \cdots \cap V_{n}$。我们关于 $n$ 使用数学归纳法。

\begin{description}
\item[第 $1$ 步]
当 $n = 1$ 时，$S = V_{1}$，显然是 $V$ 的子空间。

\item[第 $k$ 步]
假设当 $n = k - 1$ 时，结论成立，即 $V_{1} \cap \cdots \cap V_{k - 1}$ 是 $V$ 的子空间。又因为 $V_{k}$ 是 $V$ 的子空间，由习题~\ref{app:E-inter-of-subspace-is-subspace} 可知，$(V_{1} \cap \cdots \cap V_{k - 1}) \cap V_{k}$ 也是 $V$ 的子空间。由此，我们证明了当 $n = k$ 时，结论也成立。

综上所述，$V$ 的任意一族子空间的交集是 $V$ 的子空间。
\end{description}
\end{solve}

\begin{exercise}\label{prob:P-1C-12}\label{app:E-when-union-of-two-subspaces-is-subspace}
证明：$V$ 的两个子空间的并集是 $V$ 的子空间，当且仅当其中一个子空间是另一个的子集。
\end{exercise}
\noindent 来源：\href{https://github.com/szdytom/LADRSolutions/blob/main/sections/1C.typ}{P-1C-12}（原书 1C 第 12 题；方而静题解）。
\begin{solve}
设 $V_{1}$ 和 $V_{2}$ 都是 $V$ 的子空间。我们首先说明充分性。不妨设 $V_{1} \subseteq V_{2}$，则 $V_{1} \cup V_{2} = V_{2}$ 是 $V$ 的子空间。

下面说明必要性。使用反证法，假设 $V_{1} \cup V_{2}$ 是 $V$ 的子空间，以及 $V_{1} \nsubseteq V_{2}$ 且 $V_{2} \nsubseteq V_{1}$。则可以找到 $v_{1} , v_{2} \in V$，使得 $v_{1} \in V_{1}$ 且 $v_{1} \notin V_{2}$，以及 $v_{2} \in V_{2}$ 且 $v_{2} \notin V_{1}$。

令 $u = v_{1} + v_{2}$。由于 $v_{1} , v_{2} \in V_{1} \cup V_{2}$，因此 $u \in V_{1} \cup V_{2}$。不妨设 $u \in V_{1}$，则 $v_{2} = u - v_{1} \in V_{1}$，矛盾，故假设不成立。

综上所述，$V$ 的两个子空间的并集是 $V$ 的子空间，当且仅当其中一个子空间是另一个的子集。
\end{solve}

\begin{exercise}\label{prob:P-1C-13}
证明：$V$ 的三个子空间的并集是 $V$ 的子空间，当且仅当其中一个包含另外两个。

\begin{remark}
令人惊讶的是，这道习题比习题~\ref{app:E-when-union-of-two-subspaces-is-subspace} 难不少，也许是因为如果我们把 $\mathbf F$ 换成只包含两个元素的域，这道习题的结论就不成立了。
\end{remark}
\end{exercise}
\noindent 来源：\href{https://github.com/szdytom/LADRSolutions/blob/main/sections/1C.typ}{P-1C-13}（原书 1C 第 13 题；方而静题解）。
\begin{solve}
设 $V_{1} , V_{2} , V_{3}$ 都是 $V$ 的子空间。我们首先说明充分性。不妨设 $V_{1} \subseteq V_{3}$ 且 $V_{2} \subseteq V_{3}$，则 $V_{1} \cup V_{2} \cup V_{3} = V_{3}$ 是 $V$ 的子空间。

下面说明必要性。使用反证法，假设 $V_{1} \cup V_{2} \cup V_{3}$ 是 $V$ 的子空间，以及任意一个 $V_{j}$ 都不包含另外两个。

我们首先说明，任意一个 $V_{j}$ 都不是另外两个的并集的子集。否则，不妨设 $V_{1} \subseteq V_{2} \cup V_{3}$，则 $V_{1} \cup V_{2} \cup V_{3} = V_{2} \cup V_{3}$ 是 $V$ 的子空间。应用习题~\ref{app:E-when-union-of-two-subspaces-is-subspace}，可以推出 $V_{2} \subseteq V_{3}$ 或 $V_{3} \subseteq V_{2}$，这说明 $V_{2}$ 或 $V_{3}$ 包含另外两个，矛盾，故假设不成立。因此，

\[V_{1} \nsubseteq V_{2} \cup V_{3} \quad \text{且} \quad V_{2} \cup V_{3} \nsubseteq V_{1}\]

所以可以找到 $u , v \in V$ 使得 $u \in V_{1}$ 且 $u \notin V_{2} \cup V_{3}$，以及 $v \in V_{2} \cup V_{3}$ 且 $v \notin V_{1}$。由于 $V_{1}$，$V_{2}$ 和 $V_{3}$ 都包含 $0$，因此 $u \ne 0$ 且 $v \ne 0$。我们取集合 $v + \operatorname{span} (u)$\footnote{记号 $\operatorname{span}$ 在后续的2A节中由原书2.19定义，而记号 $v + V$ 表示平移，由后续3E节原书3.97定义。但是这里无需明白这些定义，将其当作一个集合的名字即可。}，

\[v + \operatorname{span} (u) = \{v + \lambda u : \lambda \in \mathbf F\}\]

下面我们说明 $V_{1} \cap (v + \operatorname{span} (u)) = $。使用反证法，假设存在 $w = v + k_{1} u \in V_{1}$，则 $v = w - k_{1} u$。又因为 $w \in V_{1}$ 且 $u \in V_{1}$，因此 $v \in V_{1}$，矛盾，故假设不成立，即 $V_{1} \cap (v + \operatorname{span} (u)) = $。

根据反证假设 $V_{1} \cup V_{2} \cup V_{3}$ 是向量空间，因此 $v + \operatorname{span} (u) \subseteq V_{1} \cup V_{2} \cup V_{3}$。又因为 $V_{1} \cap (v + \operatorname{span} (u)) = $，所以

\[v + \operatorname{span} (u) \subseteq V_{2} \cup V_{3}\]

现在，我们说明 $v + \operatorname{span} (u)$ 中包含至少 $3$ 个元素。我们取函数

\[\begin{aligned}f : & \mathbf F \to v + \operatorname{span} (u) \\ & x \mapsto v + x u\end{aligned}\]

对于任意的 $x_{1} , x_{2} \in \mathbf F$，我们有

\[f (x_{1}) - f (x_{2}) = (v + x_{1} u) - (v + x_{2} u) = (x_{1} - x_{2}) u\]

由于 $u \ne 0$，因此 $f (x_{1}) = f (x_{2})$ 当且仅当 $x_{1} = x_{2}$。这说明 $f$ 是单射，即 $v + \operatorname{span} (u)$ 至少和 $\mathbf F$ 一样大，因此 $v + \operatorname{span} (u)$ 至少包含 $3$ 个元素。

根据抽屉原理\footnote{抽屉原理的一种通俗的说法是：若将 $n$ 个物品放在 $r$ 个盒子里，$r < n$，那么至少有一个盒子包含多于一个物品。在这里，相当于是将 $v + \operatorname{span} (u)$ 中的大于等于 $3$ 个物品放入 $V_{2}$ 和 $V_{3}$ 两个盒子中。}，在 $V_{2}$ 与 $V_{3}$ 中至少有一个包含 $v + \operatorname{span} (u)$ 中的两个元素。不妨设 $V_{2}$ 包含 $v + \operatorname{span} (u)$ 中的两个元素 $w_{1} = v + \mu_{1} u$ 和 $w_{2} = v + \mu_{2} u$，其中 $\mu_{1} , \mu_{2} \in \mathbf F$ 且 $\mu_{1} \ne \mu_{2}$。由于 $V_{2}$ 是向量空间，故

\[w_{1} - w_{2} = (v + \mu_{1} u) - (v + \mu_{2} u) = (\mu_{1} - \mu_{2}) u \in V_{2}\]

由于 $\mu_{1} \ne \mu_{2}$，我们立即得到 $u \in V_{2}$，而这与 $u \notin V_{2} \cup V_{3}$ 矛盾，故假设不成立。

综上所述，$V$ 的三个子空间的并集是 $V$ 的子空间，当且仅当其中一个包含另外两个。
\end{solve}

\begin{exercise}\label{prob:P-1C-14}
令

\[U = \{(x , - x , 2 x) \in \mathbf F^{3} : x \in \mathbf F\} \quad \text{与} \quad W = \{(x , x , 2 x) \in \mathbf F^{3} : x \in \mathbf F\}\]

用符号描述 $U + W$，并给出不使用符号的描述。
\end{exercise}
\noindent 来源：\href{https://github.com/szdytom/LADRSolutions/blob/main/sections/1C.typ}{P-1C-14}（原书 1C 第 14 题；方而静题解）。
\begin{solve}
我们声称

\[S = \{(x , y , 2 x) \in \mathbf F^{3} : x , y \in \mathbf F^{3}\} = U + W\]

为了证明这一点，我们将论证 $S \subseteq U + W$ 以及 $U + W \subseteq S$。现在设 $u = (a , b , 2 a)$，其中 $a , b \in \mathbf F$，即 $u \in S$。注意到 $u = v_{1} + v_{2}$，其中

\[\begin{aligned}v_{1} & = (\frac{a - b}{2} , - \frac{a - b}{2} , a - b) \in U \\ v_{2} & = (\frac{a + b}{2} , \frac{a + b}{2} , a + b) \in W\end{aligned}\]

故 $u \in U + W$，这表明 $S \subseteq U + W$。另一方面，设 $u = (a + b , - a + b , 2 a + 2 b)$，其中 $a , b \in \mathbf F$，即 $u \in U + W$。注意到 $u = (x , y , 2 x)$，其中

\[\begin{cases}x = a + b\\y = - a + b\end{cases}\]

这表明 $u \in S$，即 $U + W \subseteq S$。综上所述，$S = U + W$。

不使用符号的描述：$U + W$ 是 $\mathbf F^{3}$ 中第三个分量是第一个分量的两倍的所有元素构成的集合。
\end{solve}

\begin{exercise}\label{prob:P-1C-15}
设 $U$ 是 $V$ 的子空间，那么 $U + U$ 是什么？
\end{exercise}
\noindent 来源：\href{https://github.com/szdytom/LADRSolutions/blob/main/sections/1C.typ}{P-1C-15}（原书 1C 第 15 题；方而静题解）。
\begin{solve}
由子空间加法封闭性可知，$U + U \subseteq U$，另一方面，对于 $u \in U$，我们有 $u + 0 = u$，因此 $U \subseteq U + U$。综上所述，$U + U = U$。
\end{solve}

\begin{exercise}\label{prob:P-1C-16}
$V$ 的“子空间求和”运算可交换吗？换句话说，设 $U$ 和 $W$ 都是 $V$ 的子空间，那么 $U + W = W + U$ 是否成立？
\end{exercise}
\noindent 来源：\href{https://github.com/szdytom/LADRSolutions/blob/main/sections/1C.typ}{P-1C-16}（原书 1C 第 16 题；方而静题解）。
\begin{solve}
是的。我们有

\[U + W = \{u + w : u \in U , w \in W\} = \{w + u : u \in U , w \in W\} = W + U\]

这说明 $U + W = W + U$ 成立。
\end{solve}

\begin{exercise}\label{prob:P-1C-17}
$V$ 的“子空间求和”运算可结合吗？换句话说，设 $V_{1} , V_{2} , V_{3}$ 都是 $V$ 的子空间，那么 $(V_{1} + V_{2}) + V_{3} = V_{1} + (V_{2} + V_{3})$ 是否成立？
\end{exercise}
\noindent 来源：\href{https://github.com/szdytom/LADRSolutions/blob/main/sections/1C.typ}{P-1C-17}（原书 1C 第 17 题；方而静题解）。
\begin{solve}
是的。我们有

\[\begin{aligned}(V_{1} + V_{2}) + V_{3} & = \{u + w : u \in V_{1} + V_{2} , w \in V_{3}\} \\ & = \{u + w : u \in \{v_{1} + v_{2} : v_{1} \in V_{1} , v_{2} \in V_{2}\} , w \in V_{3}\} \\ & = \{v_{1} + v_{2} + w : v_{1} \in V_{1} , v_{2} \in V_{2} , w \in V_{3}\} \\ & = V_{1} + (V_{2} + V_{3})\end{aligned}\]

这说明 $(V_{1} + V_{2}) + V_{3} = V_{1} + (V_{2} + V_{3})$ 成立。
\end{solve}

\begin{exercise}\label{prob:P-1C-18}
$V$ 的“子空间求和”运算有没有加法单位元？哪些子空间有加法逆元？
\end{exercise}
\noindent 来源：\href{https://github.com/szdytom/LADRSolutions/blob/main/sections/1C.typ}{P-1C-18}（原书 1C 第 18 题；方而静题解）。
\begin{solve}
“子空间求和”运算的加法单位元是 $\{0\}$。设 $U$ 是 $V$ 的子空间，对于任意 $u \in U$,都有 $u + 0 = u$，故 $U + \{0\} = U$。

由于 $\{0\} + \{0\} = \{0\}$，故 $\{0\}$ 有加法逆元 $\{0\}$。我们下面说明，除此之外，没有任何子空间有加法逆元。使用反证法，假设 $U$ 是 $V$ 的子空间，满足 $U \ne \{0\}$ 且 $W$ 是 $U$ 的加法逆元，注意到 $U \subseteq U + W = \{0\}$，矛盾，故假设不成立。

综上所述，“子空间求和”运算的加法单位元是 $\{0\}$，而只有 $\{0\}$ 有加法逆元。
\end{solve}

\begin{exercise}\label{prob:P-1C-19}
证明或证伪：如果 $V_{1}$，$V_{2}$，$U$ 都是 $V$ 的子空间，且

\[V_{1} + U = V_{2} + U\]

则有 $V_{1} = V_{2}$。
\end{exercise}
\noindent 来源：\href{https://github.com/szdytom/LADRSolutions/blob/main/sections/1C.typ}{P-1C-19}（原书 1C 第 19 题；方而静题解）。
\begin{solve}
令

\[\begin{aligned}V & = \mathbf R^{3} \\ U & = \{(0 , x , y) \in \mathbf R^{3} : x , y \in \mathbf R\} \\ V_{1} & = \{0\} \\ V_{2} & = \{(0 , x , 0) \in \mathbf R^{3} : x \in \mathbf R\}\end{aligned}\]

注意到 $V_{1} + U = U = V_{2} + U$，但 $V_{1} \ne V_{2}$。

该反例说明，题目中的命题不成立。
\end{solve}

\begin{exercise}\label{prob:P-1C-20}
令

\[U = \{(x , x , y , y) \in \mathbf F^{4} : x , y \in \mathbf F\}\]

求 $\mathbf F^{4}$ 的一个子空间 $W$，使得 $\mathbf F^{4} = U \oplus W$。
\end{exercise}
\noindent 来源：\href{https://github.com/szdytom/LADRSolutions/blob/main/sections/1C.typ}{P-1C-20}（原书 1C 第 20 题；方而静题解）。
\begin{solve}
取

\[W = \{(x , 0 , y , 0) \in \mathbf F^{4} : x , y \in \mathbf F\}\]

我们首先证明 $\mathbf F^{4} = U + W$。任取 $u = (a , b , c , d) \in \mathbf F^{4}$，注意到 $u = v_{1} + v_{2}$，其中，

\[\begin{aligned}v_{1} & = (b , b , d , d) \in U \\ v_{2} & = (a - b , 0 , c - d , 0) \in W\end{aligned}\]

进一步地，我们说明这个和是直和。根据两个子空间的直和的条件（原书1.46），我们只需说明 $U \cap W = \{0\}$。设 $v \in U \cap W$，那么存在 $a , b , c , d \in \mathbf F$，使得

\[(a , a , b , b) = v = (c , 0 , d , 0)\]

解得 $a = b = c = d = 0$，故 $U \cap W = \{0\}$。

综上所述，$\mathbf F^{4} = U \oplus W$。
\end{solve}

\begin{exercise}\label{prob:P-1C-21}
令

\[U = \{(x , y , x + y , x - y , 2 x) \in \mathbf F^{5} : x , y \in \mathbf F\}\]

求 $\mathbf F^{5}$ 的一个子空间 $W$，使得 $\mathbf F^{5} = U \oplus W$。
\end{exercise}
\noindent 来源：\href{https://github.com/szdytom/LADRSolutions/blob/main/sections/1C.typ}{P-1C-21}（原书 1C 第 21 题；方而静题解）。
\begin{solve}
可取 $W=\{(0,0,a,b,c):a,b,c\in\mathbf F\}$。任意 $(x_1,x_2,x_3,x_4,x_5)$ 唯一地分解为
\[(x_1,x_2,x_1+x_2,x_1-x_2,2x_1)+(0,0,x_3-x_1-x_2,x_4-x_1+x_2,x_5-2x_1).\]
第一项属于 $U$，第二项属于 $W$，故两空间的和为 $\mathbf F^5$。若一个向量同时属于两者，则前两个坐标为零，$U$ 的定义随即使其他坐标也为零。因此交为 $\{0\}$，和是直和。
\end{solve}

\begin{exercise}\label{prob:P-1C-22}
令

\[U = \{(x , y , x + y , x - y , 2 x) \in \mathbf F^{5} : x , y \in \mathbf F\}\]

求 $\mathbf F^{5}$ 的三个都不为 $\{0\}$ 的子空间 $W_{1}$，$W_{2}$ 和 $W_{3}$，使得 $\mathbf F^{5} = U \oplus W_{1} \oplus W_{2} \oplus W_{3}$。
\end{exercise}
\noindent 来源：\href{https://github.com/szdytom/LADRSolutions/blob/main/sections/1C.typ}{P-1C-22}（原书 1C 第 22 题；方而静题解）。
\begin{solve}
对于 $i \in \{1 , 2 , 3\}$，取

\[\begin{aligned}W_{1} & = \{(x , 0 , 0 , 0 , 0) \in \mathbf F^{5} : x \in \mathbf F\} \\ W_{2} & = \{(0 , x , 0 , 0 , 0) \in \mathbf F^{5} : x \in \mathbf F\} \\ W_{3} & = \{(0 , 0 , x , 0 , 0) \in \mathbf F^{5} : x \in \mathbf F\}\end{aligned}\]

我们首先说明，$\mathbf F^{5} = U + W_{1} + W_{2} + W_{3}$。任取 $u = (a , b , c , d , e , f) \in \mathbf F^{5}$，注意到 $u = w + v_{1} + v_{2} + v_{3}$，其中

\[\begin{aligned}w & = (\frac{e}{2} , - d + \frac{e}{2} , - d + e , d , e) \\ v_{1} & = (a - \frac{e}{2} , 0 , 0 , 0 , 0) \\ v_{2} & = (0 , b + d - \frac{e}{2} , 0 , 0 , 0) \\ v_{3} & = (0 , 0 , c + d - e , 0 , 0)\end{aligned}\]

进一步地，我们说明这个和是直和。我们将 $0$ 分解为四个向量的和，使得每个向量都来自于一个子空间。具体地，我们设

\[0 = (a , b , a + b , a - b , 2 a) + (c , 0 , 0 , 0 , 0) + (0 , d , 0 , 0 , 0) + (0 , 0 , e , 0 , 0)\]

其中 $a , b , c , d , e \in \mathbf F$。这给出一个五元一次方程组

\[\begin{cases}a + c = 0\\b + d = 0\\a + b + e = 0\\a - b = 0\\2 a = 0\end{cases}\]

这个方程组的唯一解是 $a = b = c = d = e = 0$，这说明 $0$ 只有唯一的表示方式。根据直和的条件（原书1.45），我们确认 $\mathbf F^{5} = U \oplus W_{1} \oplus W_{2} \oplus W_{3}$。
\end{solve}

\begin{exercise}\label{prob:P-1C-23}
证明或证伪：如果 $V_{1}$，$V_{2}$，$U$ 都是 $V$ 的子空间，且

\[V = V_{1} \oplus U \quad \text{且} \quad V = V_{2} \oplus U\]

则有 $V_{1} = V_{2}$。

\begin{remark}
在尝试确认线性代数中的一个命题是否成立时，先在 $\mathbf F^{2}$ 中试试，通常时很有帮助的。
\end{remark}
\end{exercise}
\noindent 来源：\href{https://github.com/szdytom/LADRSolutions/blob/main/sections/1C.typ}{P-1C-23}（原书 1C 第 23 题；方而静题解）。
\begin{solve}
令

\[\begin{aligned}V & = \mathbf F^{2} \\ U & = \{(0 , x) \in \mathbf F^{2} : x \in \mathbf F\} \\ V_{1} & = \{(x , 0) \in \mathbf F^{2} : x \in \mathbf F\} \\ V_{2} & = \{(x , x) \in \mathbf F^{2} : x \in \mathbf F\}\end{aligned}\]

容易验证，$V = V_{1} \oplus U$ 且 $V = V_{2} \oplus U$，但 $V_{1} \ne V_{2}$。

该反例说明，题目中的命题不成立。
\end{solve}

\begin{exercise}\label{prob:P-1C-24}
函数 $f : \mathbf R \to \mathbf R$ 被称为\emph{偶的（even）}，是指

\[f (- x) = f (x)\]

对所有 $x \in \mathbf R$ 成立。函数 $f : \mathbf R \to \mathbf R$ 被称为\emph{奇的（odd）}，是指

\[f (- x) = - f (x)\]

对所有 $x \in \mathbf R$ 成立。令 $V_{\text{e}}$ 代表 $\mathbf R$ 上的实值偶函数构成的集合，$V_{\text{o}}$ 代表 $\mathbf R$ 上的实值奇函数构成的集合。证明：$\mathbf R^{\mathbf R} = V_{\text{e}} + V_{\text{o}}$。
\end{exercise}
\noindent 来源：\href{https://github.com/szdytom/LADRSolutions/blob/main/sections/1C.typ}{P-1C-24}（原书 1C 第 24 题；方而静题解）。
\begin{solve}
对任意 $f\colon\mathbf R\to\mathbf R$，令
\[f_{\mathrm e}(x)=\frac{f(x)+f(-x)}2,\qquad f_{\mathrm o}(x)=\frac{f(x)-f(-x)}2.\]
则 $f=f_{\mathrm e}+f_{\mathrm o}$，且 $f_{\mathrm e}(-x)=f_{\mathrm e}(x)$、$f_{\mathrm o}(-x)=-f_{\mathrm o}(x)$，故 $\mathbf R^{\mathbf R}=V_{\mathrm e}+V_{\mathrm o}$。若 $g$ 同时是奇函数和偶函数，则 $g(x)=g(-x)=-g(x)$，于是 $g=0$。因此还得到更强的结论 $\mathbf R^{\mathbf R}=V_{\mathrm e}\oplus V_{\mathrm o}$，上述分解唯一。
\end{solve}

\subsection{张成空间与线性无关性}\label{app:section-2A}
\begin{exercise}\label{prob:P-2A-1}
求 $\mathbf F^{3}$ 中的四个不同的向量，其张成空间为

\[\{(x , y , z) \in \mathbf F^{3} : x + y + z = 0\}\]
\end{exercise}
\noindent 来源：\href{https://github.com/szdytom/LADRSolutions/blob/main/sections/2A.typ}{P-2A-1}（原书 2A 第 1 题；方而静题解）。
\begin{solve}
取 $u = (1 , 0 , - 1)$，$v = (0 , 1 , - 1)$。根据张成空间的定义，

\[\operatorname{span} (u , v) = \{(x , y , - x - y) \in \mathbf F^{3} : x , y \in \mathbf F\}\]

这已经是题目所要求的张成空间了。为了补足四个不同的向量，我们可以取 $w_{1} = 2 u$，$w_{2} = 2 v$。

综上所述，题目要求的四个不同的向量可以是

\[(1 , 0 , - 1) , (0 , 1 , - 1) , (2 , 0 , - 2) , (0 , 2 , - 2)\]
\end{solve}

\begin{exercise}\label{prob:P-2A-2}
证明或证伪：如果向量组 $v_{1} , v_{2} , v_{3} , v_{4}$ 张成 $V$，那么向量组

\[v_{1} - v_{2} , v_{2} - v_{3} , v_{3} - v_{4} , v_{4}\]

也张成 $V$。
\end{exercise}
\noindent 来源：\href{https://github.com/szdytom/LADRSolutions/blob/main/sections/2A.typ}{P-2A-2}（原书 2A 第 2 题；方而静题解）。
\begin{solve}
对于任意 $v \in V$，可以将其表示为

\[v = a_{1} v_{1} + a_{2} v_{2} + a_{3} v_{3} + a_{4} v_{4}\]

其中 $a_{i} \in \mathbf F$。我们可以将其改写为

\[v = a_{1} (v_{1} - v_{2}) + (a_{1} + a_{2}) (v_{2} - v_{3}) + (a_{2} + a_{3}) (v_{3} - v_{4}) + (a_{3} + a_{4}) v_{4}\]

这说明 $v$ 可以用 $v_{1} - v_{2}$，$v_{2} - v_{3}$，$v_{3} - v_{4}$ 和 $v_{4}$ 线性表示，这表明 $V \subseteq \operatorname{span} (v_{1} - v_{2} , v_{2} - v_{3} , v_{3} - v_{4} , v_{4})$。

另一方面，设 $v \in \operatorname{span} (v_{1} - v_{2} , v_{2} - v_{3} , v_{3} - v_{4} , v_{4})$，则 $v$ 可以表示为

\[v = b_{1} (v_{1} - v_{2}) + b_{2} (v_{2} - v_{3}) + b_{3} (v_{3} - v_{4}) + b_{4} v_{4}\]

其中 $b_{i} \in \mathbf F$。我们可以将其改写为

\[v = (b_{1} + b_{2}) v_{1} + (b_{2} + b_{3}) v_{2} + (b_{3} + b_{4}) v_{3} + b_{4} v_{4}\]

这说明 $v$ 可以用 $v_{1}$，$v_{2}$，$v_{3}$ 和 $v_{4}$ 线性表示，这表明 $\operatorname{span} (v_{1} - v_{2} , v_{2} - v_{3} , v_{3} - v_{4} , v_{4}) \subseteq V$。

综上所述，$V = \operatorname{span} (v_{1} - v_{2} , v_{2} - v_{3} , v_{3} - v_{4} , v_{4})$，即向量组 $v_{1} - v_{2} , v_{2} - v_{3} , v_{3} - v_{4} , v_{4}$ 也张成 $V$。
\end{solve}

\begin{exercise}\label{prob:P-2A-3}\label{app:E-cumulative-sum-span}
设 $v_{1} , \ldots , v_{m}$ 是 $V$ 中的一组向量。对于 $k \in \{1 , \ldots , m\}$，令

\[w_{k} = v_{1} + \cdots + v_{k}\]

证明：$\operatorname{span} (v_{1} , \ldots , v_{m}) = \operatorname{span} (w_{1} , \ldots , w_{m})$。
\end{exercise}
\noindent 来源：\href{https://github.com/szdytom/LADRSolutions/blob/main/sections/2A.typ}{P-2A-3}（原书 2A 第 3 题；方而静题解）。
\begin{solve}
设 $u \in \operatorname{span} (v_{1} , \ldots , v_{m})$，则 $u$ 可以表示为

\[u = a_{1} v_{1} + \cdots + a_{m} v_{m}\]

其中 $a_{i} \in \mathbf F$。我们可以将其改写为

\[u = b_{1} w_{1} + \cdots + b_{m} w_{m}\]

其中对于 $k \in \{1 , \ldots , m\}$，$b_{k} = a_{k} - a_{k + 1}$（方便起见，我们设 $a_{m + 1} = 0$。为了验证这一点，我们带入 $b_{i}$ 和 $w_{i}$ 的定义，得到

\[\begin{aligned}b_{1} w_{1} + \cdots + b_{m} w_{m} & = \sum_{i = 1}^{m} (a_{i} - a_{i + 1}) \sum_{j = 1}^{i} v_{j} \\ & = \sum_{i = 1}^{m} \sum_{j = 1}^{i} (a_{i} - a_{i + 1}) v_{j} \\ & = \sum_{j = 1}^{m} \sum_{i = j}^{m} (a_{i} - a_{i + 1}) v_{j} \\ & = \sum_{j = 1}^{m} (\sum_{i = j}^{m} a_{i} - \sum_{i = j}^{m} a_{i + 1}) v_{j} \\ & = \sum_{j = 1}^{m} a_{j} v_{j} = u\end{aligned}\]

这说明 $u$ 可以用 $w_{1} , \ldots , w_{m}$ 线性表示，因此 $\operatorname{span} (v_{1} , \ldots , v_{m}) \subseteq \operatorname{span} (w_{1} , \ldots , w_{m})$。另一方面，设 $u \in \operatorname{span} (w_{1} , \ldots , w_{m})$，则 $u$ 可以表示为

\[u = a_{1} w_{1} + \cdots + a_{m} w_{m}\]

其中 $a_{i} \in \mathbf F$。我们可以将其改写为

\[u = b_{1} v_{1} + \cdots + b_{m} v_{m}\]

其中对于 $k \in \{1 , \ldots , m\}$，$b_{k} = a_{k} + \cdots + a_{m}$。为了验证这一点，我们带入 $b_{i}$ 和 $w_{i}$ 的定义，得到

\[\begin{aligned}b_{1} v_{1} + \cdots + b_{m} v_{m} & = \sum_{i = 1}^{m} (\sum_{j = i}^{m} a_{j}) v_{i} \\ & = \sum_{i = 1}^{m} \sum_{j = i}^{m} a_{j} v_{i} \\ & = \sum_{j = 1}^{m} \sum_{i = 1}^{j} a_{j} v_{i} \\ & = \sum_{j = 1}^{m} a_{j} \sum_{i = 1}^{j} v_{i} \\ & = \sum_{j = 1}^{m} a_{j} w_{j} = u\end{aligned}\]

这说明 $u$ 可以用 $v_{1} , \ldots , v_{m}$ 线性表示，因此 $\operatorname{span} (w_{1} , \ldots , w_{m}) \subseteq \operatorname{span} (v_{1} , \ldots , v_{m})$。

综上所述，$\operatorname{span} (v_{1} , \ldots , v_{m}) = \operatorname{span} (w_{1} , \ldots , w_{m})$。
\end{solve}

\begin{exercise}\label{prob:P-2A-4}\label{app:E-when-1-or-2-vectors-indep}
\begin{enumerate}[label=(\alph*)]
\item 证明：向量空间中长度为 $1$ 的组线性无关，当且仅当组中的该向量不是 $0$；
\item 证明：向量空间中长度为 $2$ 的组线性无关，当且仅当组中两个向量的任意一个不是另一个的标量倍。
\end{enumerate}
\end{exercise}
\noindent 来源：\href{https://github.com/szdytom/LADRSolutions/blob/main/sections/2A.typ}{P-2A-4}（原书 2A 第 4 题；方而静题解）。
\begin{solve}
设 $V$ 是向量空间。对于第一个命题，设 $v \in V$。首先说明充分性。假设 $v \ne 0$，根据习题~\ref{app:E-vector-trivial-annihilation}，使得 $a v = 0$ 成立的 $a \in \mathbf F$ 的唯一选取方式是 $a = 0$，根据线性无关的定义（原书2.15），这表明向量组 $v$ 是线性无关的。

然后说明必要性，如果 $v = 0$，则我们有

\[0 v = 1 v = 0\]

根据线性无关的定义（原书2.15），向量组 $0$ 不是线性无关的。

对于第二个命题，设 $v_{1} , v_{2} \in V$。首先说明充分性：使用反证法，假设 $v_{1}$ 和 $v_{2}$ 不是线性无关的，即存在 $a_{1} , a_{2} \in \mathbf F$，使得

\[a_{1} v_{1} + a_{2} v_{2} = 0\]

其中 $a_{1}$ 和 $a_{2}$ 中至少有一个向量不为 $0$。不妨设 $a_{1} \ne 0$，那么，可以整理得

\[v_{1} = - \frac{a_{2}}{a_{1}} v_{2}\]

这表明 $v_{1}$ 是 $v_{2}$ 的标量倍，与题目条件矛盾。这说明，$v_{1}$ 和 $v_{2}$ 线性无关。

然后说明必要性：假设 $v_{1}$ 和 $v_{2}$ 线性无关，使用反证法，假设 $v_{1}$ 和 $v_{2}$ 中有一个向量是另一个向量的标量倍，不妨设 $v_{1} = k v_{2}$，其中 $k \in \mathbf F$。注意到，

\[v_{1} + (- k) v_{2} = 0\]

这与线性无关的定义矛盾。因此，$v_{1}$ 和 $v_{2}$ 线性无关当且仅当组中两个向量的任意一个不是另一个的标量倍。
\end{solve}

\begin{exercise}\label{prob:P-2A-5}
求一数 $t \in \mathbf R$，使得向量组

\[(3 , 1 , 4) , (2 , - 3 , 5) , (5 , 9 , t)\]

在 $\mathbf R^{3}$ 中不是线性无关的。
\end{exercise}
\noindent 来源：\href{https://github.com/szdytom/LADRSolutions/blob/main/sections/2A.typ}{P-2A-5}（原书 2A 第 5 题；方而静题解）。
\begin{solve}
$t = 2$。为了说明这一点，注意到，

\[(- 3) (3 , 1 , 4) + 2 (2 , - 3 , 5) + 1 (5 , 9 , 2) = 0\]

根据线性无关的定义（原书2.15），这表明向量组 $(3 , 1 , 4) , (2 , - 3 , 5) , (5 , 9 , 2)$ 在 $\mathbf R^{3}$ 中不是线性无关的。
\end{solve}

\begin{exercise}\label{prob:P-2A-6}
证明：向量组

\[(2 , 3 , 1) , (1 , - 1 , 2) , (7 , 3 , c)\]

在 $\mathbf F^{3}$ 中线性相关，当且仅当 $c = 8$。
\end{exercise}
\noindent 来源：\href{https://github.com/szdytom/LADRSolutions/blob/main/sections/2A.typ}{P-2A-6}（原书 2A 第 6 题；方而静题解）。
\begin{solve}
首先说明充分性：当 $c = 8$ 时，注意到

\[2 (2 , 3 , 1) + 3 (1 , - 1 , 2) - (7 , 3 , 8) = 0\]

根据线性相关的定义（原书2.17），这表明向量组 $(2 , 3 , 1)$，$(1 , - 1 , 2)$，$(7 , 3 , 8)$ 在 $\mathbf F^{3}$ 中线性相关。

然后说明必要性：使用反证法，假设 $c \ne 8$ 且向量组 $(2 , 3 , 1) , (1 , - 1 , 2) , (7 , 3 , c)$ 线性相关。根据定义，存在 $a_{1} , a_{2} , a_{3} \in \mathbf F$，使得

\[a_{1} (2 , 3 , 1) + a_{2} (1 , - 1 , 2) + a_{3} (7 , 3 , c) = 0\]

其中 $a_{1} , a_{2} , a_{3} \in \mathbf F$ 中至少有一个不为 $0$。将其展开，得到下面方程组

\[
\begin{cases}
2a_1+a_2+7a_3=0,\\
3a_1-a_2+3a_3=0,\\
a_1+2a_2+ca_3=0.
\end{cases}
\]

由前两个方程，我们可以得到 $a_{2} = \frac{3}{2} a_{1}$ 且 $a_{3} = - \frac{1}{2} a_{1}$，代入第三个方程中，化简得

\[(c - 8) a_{1} = 0\]

由于 $c \ne 8$，只能有 $a_{1} = 0$，而这将给出 $a_{1} = a_{2} = a_{3} = 0$，与反证假设矛盾，故假设不成立。

综上所述，向量组 $(2 , 3 , 1) , (1 , - 1 , 2) , (7 , 3 , c)$ 在 $\mathbf F^{3}$ 中线性相关当且仅当 $c = 8$。
\end{solve}

\begin{exercise}\label{prob:P-2A-7}
\begin{enumerate}[label=(\alph*)]
\item 证明：如果我们将 $\mathbf C$ 视为 $\mathbf R$ 上的向量空间，那么向量组 $1 + \mathrm i , 1 - \mathrm i$ 是线性无关的；
\item 证明：如果我们将 $\mathbf C$ 视为 $\mathbf C$ 上的向量空间，那么向量组 $1 + \mathrm i , 1 - \mathrm i$ 是线性相关的。
\end{enumerate}
\end{exercise}
\noindent 来源：\href{https://github.com/szdytom/LADRSolutions/blob/main/sections/2A.typ}{P-2A-7}（原书 2A 第 7 题；方而静题解）。
\begin{solve}
利用习题~\ref{app:E-when-1-or-2-vectors-indep} 中的结论，我们只需注意到，

\[\frac{1 + i}{1 - i} = \mathrm i\]

考虑到 $\mathrm i \in \mathbf C$ 但是 $\mathrm i \notin \mathbf R$，因此 $1 + i$ 只在 $\mathbf C$ 上是 $1 - \mathrm i$ 的标量倍。这表明，如果我们将 $\mathbf C$ 视为 $\mathbf R$ 上的向量空间，那么向量组 $1 + \mathrm i , 1 - \mathrm i$ 是线性无关的，而如果我们将 $\mathbf C$ 视为 $\mathbf C$ 上的向量空间，那么向量组 $1 + \mathrm i , 1 - \mathrm i$ 是线性相关的。
\end{solve}

\begin{exercise}\label{prob:P-2A-8}
设 $v_{1} , v_{2} , v_{3} , v_{4}$ 是 $V$ 中的线性无关向量组。证明：向量组

\[v_{1} - v_{2} , v_{2} - v_{3} , v_{3} - v_{4} , v_{4}\]

也线性无关。
\end{exercise}
\noindent 来源：\href{https://github.com/szdytom/LADRSolutions/blob/main/sections/2A.typ}{P-2A-8}（原书 2A 第 8 题；方而静题解）。
\begin{solve}
设 $a_{1} , a_{2} , a_{3} , a_{4} \in \mathbf F$，使得

\[a_{1} (v_{1} - v_{2}) + a_{2} (v_{2} - v_{3}) + a_{3} (v_{3} - v_{4}) + a_{4} v_{4} = 0\]

整理得到

\[a_{1} v_{1} + (a_{2} - a_{1}) v_{2} + (a_{3} - a_{2}) v_{3} + (a_{4} - a_{3}) v_{4} = 0\]

由于 $v_{1} , v_{2} , v_{3} , v_{4}$ 线性无关，根据线性无关的定义（原书2.15），只能有

\[\begin{cases}a_{1} = 0\\a_{2} - a_{1} = 0\\a_{3} - a_{2} = 0\\a_{4} - a_{3} = 0\end{cases}\]

这说明 $a_{1} = a_{2} = a_{3} = a_{4} = 0$。因此，根据线性无关的定义（原书2.15），向量组 $v_{1} - v_{2} , v_{2} - v_{3} , v_{3} - v_{4} , v_{4}$ 线性无关。
\end{solve}

\begin{exercise}\label{prob:P-2A-9}
设 $m\ge2$，$v_1,\ldots,v_m$ 是 $V$ 中的线性无关向量组。证明或证伪：$5v_1-4v_2,v_2,v_3,\ldots,v_m$ 也线性无关。
\end{exercise}
\noindent 来源：\href{https://github.com/szdytom/LADRSolutions/blob/main/sections/2A.typ}{P-2A-9}（原书 2A 第 9 题；方而静题解）。
\begin{solve}
设 $a_{1} , \ldots , a_{m} \in \mathbf F$，使得

\[a_{1} (5 v_{1} - 4 v_{2}) + a_{2} v_{2} + \cdots + a_{m} v_{m} = 0\]

整理得到

\[5 a_{1} v_{1} + (a_{2} - 4 a_{1}) + a_{3} v_{3} + a_{4} v_{4} + \cdots + a_{m} v_{m} = 0\]

由于 $v_{1} , \ldots , v_{m}$ 线性无关，根据线性无关的定义（原书2.15），只能有

\[\begin{cases}5 a_{1} = 0\\a_{2} - 4 a_{1} = 0\\a_{3} = 0\\a_{4} = 0\\\cdots\\a_{m} = 0\end{cases}\]

这说明 $a_{1} = \cdots = a_{m} = 0$。因此，根据线性无关的定义（原书2.15），向量组 $5 v_{1} - 4 v_{2} , v_{2} , v_{3} , \ldots , v_{m}$ 线性无关。
\end{solve}

\begin{exercise}\label{prob:P-2A-10}
证明或证伪：设 $v_{1} , \ldots , v_{m}$ 是 $V$ 中的线性无关向量组，$\lambda \in \mathbf F$（$\lambda \ne 0$）。则向量组

\[\lambda v_{1} , \ldots , \lambda v_{m}\]

也线性无关。
\end{exercise}
\noindent 来源：\href{https://github.com/szdytom/LADRSolutions/blob/main/sections/2A.typ}{P-2A-10}（原书 2A 第 10 题；方而静题解）。
\begin{solve}
设 $a_{1} , \ldots , a_{m} \in \mathbf F$，使得

\[a_{1} (\lambda v_{1}) + \cdots + a_{m} (\lambda v_{m}) = 0\]

整理得到

\[\lambda (a_{1} v_{1} + \cdots + a_{m} v_{m}) = 0\]

由于 $\lambda \ne 0$，因此 $a_{1} v_{1} + \cdots + a_{m} v_{m} = 0$。由于 $v_{1} , \ldots , v_{m}$ 线性无关，根据线性无关的定义（原书2.15），只能有 $a_{1} = \cdots = a_{m} = 0$。因此，根据线性无关的定义（原书2.15），向量组 $\lambda v_{1} , \ldots , \lambda v_{m}$ 线性无关。
\end{solve}

\begin{exercise}\label{prob:P-2A-11}
证明或证伪：设 $v_{1} , \ldots v_{m}$ 和 $w_{1} , \ldots , w_{m}$ 都是 $V$ 中的线性无关向量组。则向量组

\[v_{1} + w_{1} , \ldots , v_{m} + w_{m}\]

也线性无关。
\end{exercise}
\noindent 来源：\href{https://github.com/szdytom/LADRSolutions/blob/main/sections/2A.typ}{P-2A-11}（原书 2A 第 11 题；方而静题解）。
\begin{solve}
取 $V = \mathbf R^{2}$，并令

\[\begin{aligned}v_{1} & = (1 , 0) , \quad & v_{2} = (0 , 1) \\ w_{1} & = (0 , 1) , & w_{2} = (1 , 0)\end{aligned}\]

容易验证这两个向量组都是 $\mathbf R^{2}$ 中的线性无关向量组。然而，注意到

\[1 (v_{1} + w_{1}) + (- 1) (v_{2} + w_{2}) = 1 (1 , 1) + (- 1) (1 , 1) = 0\]

根据线性无关的定义（原书2.15），这表明向量组 $v_{1} + w_{1} , v_{2} + w_{2}$ 不是线性无关的，因此原命题不成立。
\end{solve}

\begin{exercise}\label{prob:P-2A-12}\label{app:E-when-shared-vec-add-become-dep}
设 $v_{1} , \ldots , v_{m}$ 是 $V$ 中的线性无关向量组，且 $w \in V$。证明：若向量组

\[v_{1} + w , \ldots , v_{m} + w\]

线性相关，则 $w \in \operatorname{span} (v_{1} , \ldots , v_{m})$。
\end{exercise}
\noindent 来源：\href{https://github.com/szdytom/LADRSolutions/blob/main/sections/2A.typ}{P-2A-12}（原书 2A 第 12 题；方而静题解）。
\begin{solve}
由于向量组 $v_{1} + w , \ldots , v_{m} + w$ 线性相关，根据线性相关的定义（原书2.17），存在 $a_{1} , \ldots , a_{m} \in \mathbf F$，使得

\begin{equation}\label{app:eq:2A-vi-plus-w-is-dependent-def}a_{1} (v_{1} + w) + \cdots + a_{m} (v_{m} + w) = 0\end{equation}

其中 $a_{1} , \ldots , a_{m}$ 中至少有一个不为 $0$。

下面我们说明 $a_{1} + \cdots + a_{m} \ne 0$。整理式~\eqref{app:eq:2A-vi-plus-w-is-dependent-def} 可得

\[a_{1} v_{1} + \cdots + a_{m} v_{m} + (a_{1} + \cdots + a_{m}) w = 0\]

反证假设 $a_{1} + \cdots + a_{m} = 0$，则 $a_{1} v_{1} + \cdots + a_{m} v_{m} = 0$，而这与线性无关的定义（原书2.15）矛盾。因此，我们只能有 $a_{1} + \cdots + a_{m} \ne 0$。

所以，我们可以将式~\eqref{app:eq:2A-vi-plus-w-is-dependent-def} 改写为

\[w = - \frac{a_{1} v_{1} + \cdots + a_{m} v_{m}}{a_{1} + \cdots + a_{m}}\]

更进一步地，对于 $k \in \{1 , \ldots , m\}$，令

\[b_{k} = - \frac{a_{k}}{a_{1} + \cdots + a_{m}}\]

则有

\[w = b_{1} v_{1} + \cdots + b_{m} v_{m}\]

根据张成空间的定义（原书2.4），这表明 $w \in \operatorname{span} (v_{1} , \ldots , v_{m})$。
\end{solve}

\begin{exercise}\label{prob:P-2A-13}\label{app:E-when-vector-list-append-remains-indep}
设 $v_{1} , \ldots , v_{m}$ 是 $V$ 中的线性无关向量组，且 $w \in V$。证明：

\[v_{1} , \ldots , v_{m} , w \text{线性无关} \quad \iff \quad w \notin \operatorname{span} (v_{1} , \ldots , v_{m})\]
\end{exercise}
\noindent 来源：\href{https://github.com/szdytom/LADRSolutions/blob/main/sections/2A.typ}{P-2A-13}（原书 2A 第 13 题；方而静题解）。
\begin{solve}
首先说明充分性：现在 $v_{1} , \ldots , v_{m} , w$ 线性无关。反证假设 $w \in \operatorname{span} (v_{1} , \ldots , v_{m})$，则根据张成空间的定义（原书2.4），存在 $a_{1} , \ldots , a_{m} \in \mathbf F$，使得

\[w = a_{1} v_{1} + \cdots + a_{m} v_{m}\]

整理得到

\[a_{1} v_{1} + \cdots + a_{m} v_{m} + (- 1) w = 0\]

根据线性无关的定义（原书2.15），这与 $v_{1} , \ldots , v_{m} , w$ 线性无关矛盾，因此，$w \notin \operatorname{span} (v_{1} , \ldots , v_{m})$。

然后说明必要性：现在 $w \notin \operatorname{span} (v_{1} , \ldots , v_{m})$。反证假设 $v_{1} , \ldots , v_{m} , w$ 线性相关。根据线性相关的定义（原书2.17），存在 $a_{1} , \ldots , a_{m + 1} \in \mathbf F$，使得

\begin{equation}\label{app:eq:2A-v-union-w-is-dependent-def}a_{1} v_{1} + \cdots + a_{m} v_{m} + a_{m + 1} w = 0\end{equation}

其中 $a_{1} , \ldots , a_{m + 1}$ 中至少有一个不为 $0$。我们有 $a_{m + 1} \ne 0$。这是因为，如果 $a_{m + 1} = 0$，则我们可以将式~\eqref{app:eq:2A-v-union-w-is-dependent-def} 改写为

\[a_{1} v_{1} + \cdots + a_{m} v_{m} = 0\]

这与题目条件中 $v_{1} , \ldots , v_{m}$ 线性无关矛盾。因此，$a_{m + 1} \ne 0$。

所以，我们可以将式~\eqref{app:eq:2A-v-union-w-is-dependent-def} 改写为

\[w = - \frac{a_{1} v_{1} + \cdots + a_{m} v_{m}}{a_{m + 1}}\]

更进一步地，对于 $k \in \{1 , \ldots , m\}$，令

\[b_{k} = - \frac{a_{k}}{a_{m + 1}}\]

则有

\[w = b_{1} v_{1} + \cdots + b_{m} v_{m}\]

这表明 $w \in \operatorname{span} (v_{1} , \ldots , v_{m})$，与反证假设 $w \notin \operatorname{span} (v_{1} , \ldots , v_{m})$ 矛盾。因此，$v_{1} , \ldots , v_{m} , w$ 线性无关。

综上所述，$v_{1} , \ldots , v_{m} , w$ 线性无关当且仅当 $w \notin \operatorname{span} (v_{1} , \ldots , v_{m})$。
\end{solve}

\begin{exercise}\label{prob:P-2A-14}\label{app:E-cumulative-sum-independence}
设 $v_{1} , \ldots , v_{m}$ 是 $V$ 中的向量组。对于 $k \in \{1 , \ldots , m\}$，令

\[w_{k} = v_{1} + \cdots + v_{k}\]

证明：如果向量组 $v_{1} , \ldots , v_{m}$ 线性无关，当且仅当向量组 $w_{1} , \ldots , w_{m}$ 线性无关。
\end{exercise}
\noindent 来源：\href{https://github.com/szdytom/LADRSolutions/blob/main/sections/2A.typ}{P-2A-14}（原书 2A 第 14 题；方而静题解）。
\begin{solve}
首先说明充分性：现在 $w_{1} , \ldots , w_{m}$ 线性无关。设 $a_{1} , \ldots , a_{m} \in \mathbf F$ 使得

\[a_{1} v_{1} + \cdots + a_{m} v_{m} = 0\]

我们可以将其改写为

\[b_{1} w_{1} + \cdots + b_{m} w_{m} = 0\]

其中对于 $k \in \{1 , \ldots , m\}$，$b_{k} = a_{k} - a_{k + 1}$（方便起见，我们令 $a_{m + 1} = 0$）。为了验证这一点，我们带入 $b_{i}$ 和 $w_{i}$ 的定义，得到

\[\begin{aligned}b_{1} w_{1} + \cdots + b_{m} w_{m} & = \sum_{i = 1}^{m} (a_{i} - a_{i + 1}) \sum_{j = 1}^{i} v_{j} \\ & = \sum_{i = 1}^{m} \sum_{j = 1}^{i} (a_{i} - a_{i + 1}) v_{j} \\ & = \sum_{j = 1}^{m} \sum_{i = j}^{m} (a_{i} - a_{i + 1}) v_{j} \\ & = \sum_{j = 1}^{m} (\sum_{i = j}^{m} a_{i} - \sum_{i = j}^{m} a_{i + 1}) v_{j} \\ & = \sum_{j = 1}^{m} a_{j} v_{j} = 0\end{aligned}\]

这说明 $b_{1} = \cdots = b_{m} = 0$，也即

\[\begin{cases}a_{1} - a_{2} = 0\\\cdots\\a_{m - 1} - a_{m} = 0\\a_{m} = 0\end{cases}\]

这解得 $a_{1} = \cdots = a_{m} = 0$，于是根据线性无关的定义（原书2.15），向量组 $v_{1} , \ldots , v_{m}$ 线性无关。

然后说明必要性：现在 $v_{1} , \ldots , v_{m}$ 线性无关。设 $a_{1} , \ldots , a_{m} \in \mathbf F$ 使得

\[a_{1} w_{1} + \cdots + a_{m} w_{m} = 0\]

我们可以将其改写为

\[b_{1} v_{1} + \cdots + b_{m} v_{m} = 0\]

其中对于 $k \in \{1 , \ldots , m\}$，$b_{k} = a_{k} + \cdots + a_{m}$。为了验证这一点，我们带入 $b_{i}$ 和 $w_{i}$ 的定义，得到

\[\begin{aligned}b_{1} v_{1} + \cdots + b_{m} v_{m} & = \sum_{i = 1}^{m} (\sum_{j = i}^{m} a_{j}) v_{i} \\ & = \sum_{i = 1}^{m} \sum_{j = i}^{m} a_{j} v_{i} \\ & = \sum_{j = 1}^{m} \sum_{i = 1}^{j} a_{j} v_{i} \\ & = \sum_{j = 1}^{m} a_{j} \sum_{i = 1}^{j} v_{i} \\ & = \sum_{j = 1}^{m} a_{j} w_{j} = 0\end{aligned}\]

这说明 $b_{1} = \cdots = b_{m} = 0$，也即

\[\begin{cases}a_{1} + \cdots + a_{m} = 0\\a_{2} + \cdots + a_{m} = 0\\\cdots\\a_{m - 1} + a_{m} = 0\\a_{m} = 0\end{cases}\]

这解得 $a_{1} = \cdots = a_{m} = 0$，于是根据线性无关的定义（原书2.15），向量组 $w_{1} , \ldots , w_{m}$ 线性无关。

综上所述，向量组 $v_{1} , \ldots , v_{m}$ 线性无关当且仅当向量组 $w_{1} , \ldots , w_{m}$ 线性无关。
\end{solve}

\begin{exercise}\label{prob:P-2A-15}
解释为什么在 $\mathcal P_{4} (\mathbf F)$ 上不存在由六个多项式组成的线性无关组。
\end{exercise}
\noindent 来源：\href{https://github.com/szdytom/LADRSolutions/blob/main/sections/2A.typ}{P-2A-15}（原书 2A 第 15 题；方而静题解）。
\begin{solve}
对于 $k \in \{0 , \ldots , 4\}$，令

\[\begin{aligned}p_{k} : & \mathbf F \to \mathbf F \\ & z \mapsto z^{k}\end{aligned}\]

我们现在论证向量组 $p_{0} , \ldots , p_{4}$ 张成 $\mathcal P_{4} (\mathbf F)$：对于任意 $p \in \mathcal P_{4} (\mathbf F)$，设对于任意 $z \in \mathbf F$，$p (z) = a_{0} + a_{1} z + a_{2} z^{2} + a_{3} z^{3} + a_{4} z^{4}$，则可以将其表示为

\[p = a_{0} p_{0} + a_{1} p_{1} + a_{2} p_{2} + a_{3} p_{3} + a_{4} p_{4}\]

这说明 $\mathcal P_{4} (\mathbf F) = \operatorname{span} (p_{0} , p_{1} , p_{2} , p_{3} , p_{4})$。因此，根据“线性无关组的长度 $\le$ 张成组的长度”（原书2.22），我们可以得出结论，$\mathcal P_{4} (\mathbf F)$ 上的线性无关组的长度不能超过 $5$。

所以，在 $\mathcal P_{4} (\mathbf F)$ 上不存在由六个多项式组成的线性无关组。
\end{solve}

\begin{exercise}\label{prob:P-2A-16}
解释为什么由四个多项式构成的向量组不可能张成 $\mathcal P_{4} (\mathbf F)$。
\end{exercise}
\noindent 来源：\href{https://github.com/szdytom/LADRSolutions/blob/main/sections/2A.typ}{P-2A-16}（原书 2A 第 16 题；方而静题解）。
\begin{solve}
对于 $k \in \{0 , \ldots , 4\}$，令

\[\begin{aligned}p_{k} : & \mathbf F \to \mathbf F \\ & z \mapsto z^{k}\end{aligned}\]

我们现在论证向量组 $p_{0} , \ldots , p_{4}$ 是线性无关的：设 $a_{0} , \ldots , a_{4} \in \mathbf F$，满足

\[a_{0} p_{0} + a_{1} p_{1} + a_{2} p_{2} + a_{3} p_{3} + a_{4} p_{4} = 0\]

即对于任意 $z \in \mathbf F$，有

\[a_{0} + a_{1} z + a_{2} z^{2} + a_{3} z^{3} + a_{4} z^{4} = 0\]

分别取 $z \in \{0 , \ldots , 4\}$，得到方程组

\[\begin{cases}a_{0} = 0\\a_{0} + a_{1} + a_{2} + a_{3} + a_{4} = 0\\a_{0} + 2 a_{1} + 4 a_{2} + 8 a_{3} + 16 a_{4} = 0\\a_{0} + 3 a_{1} + 9 a_{2} + 27 a_{3} + 81 a_{4} = 0\\a_{0} + 4 a_{1} + 16 a_{2} + 64 a_{3} + 256 a_{4} = 0\end{cases}\]

解得 $a_{0} = a_{1} = a_{2} = a_{3} = a_{4} = 0$，因此向量组 $p_{0} , \ldots , p_{4}$ 是线性无关的。

因此，根据“线性无关组的长度 $\le$ 张成组的长度”（原书2.22），$\mathcal P_{4} (\mathbf F)$ 上的张成组的长度不少于 $5$。因此，由四个多项式构成的向量组不可能张成 $\mathcal P_{4} (\mathbf F)$。
\end{solve}

\begin{exercise}\label{prob:P-2A-17}\label{app:E-inf-dim-space-seq-characterization}
证明：$V$ 是无限维的，当且仅当 $V$ 中存在一个序列 $v_{1} , v_{2} , \ldots$ 使得对于任意正整数 $m$，均有向量组 $v_{1} , \ldots , v_{m}$ 线性无关。
\end{exercise}
\noindent 来源：\href{https://github.com/szdytom/LADRSolutions/blob/main/sections/2A.typ}{P-2A-17}（原书 2A 第 17 题；方而静题解）。
\begin{solve}
首先说明充分性：现在假设 $V$ 中存在一个序列 $v_{1} , v_{2} , \ldots$ 使得对于任意正整数 $m$，均有向量组 $v_{1} , \ldots , v_{m}$ 线性无关。反证假设 $V$ 是有限维的，即存在一个向量组 $u_{1} , \ldots , u_{\ell}$ 张成 $V$。根据“线性无关组的长度 $\le$ 张成组的长度”（原书2.22），必然有向量组 $v_{1} , \ldots , v_{\ell + 1}$ 线性相关，这与条件矛盾。因此，$V$ 是无限维的。

然后说明必要性：现在假设 $V$ 是无限维的。我们现在构造题目所要求的序列 $v_{1} , v_{2} , \ldots$ 如下

\begin{description}
\item[第 $1$ 步]
任取 $v_{1} \in V$，使得 $v_{1} \ne 0$。根据习题~\ref{app:E-when-1-or-2-vectors-indep} 中的结论，向量组 $v_{1}$ 是线性无关的。

\item[第 $k$ 步]
由于 $V$ 是无限维的，存在一个向量 $v_{k} \in V$，使得 $v_{k} \notin \operatorname{span} (v_{1} , \ldots , v_{k - 1})$。根据习题~\ref{app:E-when-vector-list-append-remains-indep} 中的结论，向量组 $v_{1} , \ldots , v_{k}$ 线性无关。

所以，$V$ 中存在一个序列 $v_{1} , v_{2} , \ldots$ 使得对于任意正整数 $m$，均有向量组 $v_{1} , \ldots , v_{m}$ 线性无关。必要性得证。

综上所述，$V$ 是无限维的，当且仅当 $V$ 中存在一个序列 $v_{1} , v_{2} , \ldots$ 使得对于任意正整数 $m$，均有向量组 $v_{1} , \ldots , v_{m}$ 线性无关。
\end{description}
\end{solve}

\begin{exercise}\label{prob:P-2A-18}
证明：$\mathbf F^{\infty}$ 是无限维的。
\end{exercise}
\noindent 来源：\href{https://github.com/szdytom/LADRSolutions/blob/main/sections/2A.typ}{P-2A-18}（原书 2A 第 18 题；方而静题解）。
\begin{solve}
对于 $k \in \mathbf N^{+}$，我们令

\[v_{k} = ( \underbrace{0 , \ldots , 0}_{k - 1 \text{个} 0} , 1 , 0 , \ldots )\]

其中 $v_{k} \in \mathbf F^{\infty}$。我们现在论证，对于任意正整数 $m$，向量组 $v_{1} , \ldots , v_{m}$ 是线性无关的。设 $a_{1} , \ldots a_{m} \in \mathbf F$，使得

\[a_{1} v_{1} + \cdots + a_{m} v_{m} = 0\]

这立即给出 $a_{1} = \cdots = a_{m} = 0$，于是根据线性无关的定义（原书2.15），向量组 $v_{1} , \ldots , v_{m}$ 是线性无关的。

所以，根据习题~\ref{app:E-inf-dim-space-seq-characterization}，$\mathbf F^{\infty}$ 是无限维的。
\end{solve}

\begin{exercise}\label{prob:P-2A-19}
证明：由区间 $[0 , 1]$ 上的所有连续实值函数构成的 $\mathbf R$ 上的向量空间是无限维的。
\end{exercise}
\noindent 来源：\href{https://github.com/szdytom/LADRSolutions/blob/main/sections/2A.typ}{P-2A-19}（原书 2A 第 19 题；方而静题解）。
\begin{solve}
对每个整数 $k\ge0$，取 $f_k(x)=x^k$，这些函数都属于 $C([0,1],\mathbf R)$。若有限线性组合 $\sum_{k=0}^m a_kf_k=0$，则多项式 $\sum_{k=0}^m a_kx^k$ 在区间 $[0,1]$ 的每一点为零。非零的 $m$ 次多项式至多有 $m$ 个不同的根，因此该多项式必须为零，多项式系数的唯一性给出 $a_0=\cdots=a_m=0$。于是对任意 $m$ 都有 $m+1$ 个线性无关的连续函数，空间不可能是有限维的。
\end{solve}

\begin{exercise}\label{prob:P-2A-20}
设 $p_{0} , \ldots , p_{m}$ 是 $\mathcal P_{m} (\mathbf F)$ 中的多项式，其满足对任意 $k \in \{0 , \ldots , m\}$ 都有 $p_{k} (2) = 0$。证明：$p_{0} , \ldots , p_{m}$ 在 $\mathcal P_{m} (\mathbf F)$ 中不是线性无关的。
\end{exercise}
\noindent 来源：\href{https://github.com/szdytom/LADRSolutions/blob/main/sections/2A.typ}{P-2A-20}（原书 2A 第 20 题；方而静题解）。
\begin{solve}
考虑求值映射 $E_2\colon\mathcal P_m(\mathbf F)\to\mathbf F$，$E_2(p)=p(2)$。它是线性满射，因为任意常数都可取到。由 $\dim\mathcal P_m(\mathbf F)=m+1$ 及秩与零度定理，$\dim\ker E_2=m$。题给 $m+1$ 个多项式都属于这个 $m$ 维核空间，所以线性相关。
\end{solve}

\subsection{基}\label{app:section-2B}
\begin{exercise}\label{prob:P-2B-1}
求出所有恰好有一个基的向量空间。
\end{exercise}
\noindent 来源：\href{https://github.com/szdytom/LADRSolutions/blob/main/sections/2B.typ}{P-2B-1}（原书 2B 第 1 题；方而静题解）。
\begin{solve}
$\{0\}$ 是唯一满足要求的向量空间，它的基是空集。对于任何其他向量空间 $V$，不妨设其的一个基为 $v_{1} , \ldots , v_{m}$，则由基的判定准则（原书2.28）可知，$V$ 中的任意向量都 $v$ 都可以唯一地表示为

\[v = a_{1} v_{1} + \cdots + a_{m} v_{m}\]

其中 $a_{1} , \ldots , a_{m} \in \mathbf F$。现在取向量组 $2 v_{1} , \ldots , 2 v_{m}$，则 $v$，可以被表示为

\[v = b_{1} (2 v_{1}) + \cdots + b_{m} (2 v_{m})\]

则对于 $k \in \{1 , \ldots , m\}$，$a_{k} = 2 b_{k}$。这只有唯一的解，即 $b_{k} = a_{k} / 2$，因此 $v$ 可以唯一地被向量组 $2 v_{1} , \ldots , 2 v_{m}$ 的线性组合表示，这表明向量组 $2 v_{1} , \ldots , 2 v_{m}$ 也是 $V$ 的一个基。由此可知，$V$ 中的任意向量都可以被表示为两个不同的基的线性组合，因此 $V$ 不可能只有一个基。

综上所述，只有 $\{0\}$ 满足题目要求。
\end{solve}

\begin{exercise}\label{prob:P-2B-2}
验证下面这些的结论。

\begin{enumerate}[label=(\alph*)]
\item 向量组 $(1 , 0 , \ldots , 0) , (0 , 1 , 0 , \ldots , 0) , \ldots , (0 , \ldots , 0 , 1)$ 是 $\mathbf F^{n}$ 的基；
\item 向量组 $(1 , 2) , (3 , 5)$ 是 $\mathbf F^{2}$ 的基；
\item 向量组 $(1 , 2 , - 4) , (7 , - 5 , 6)$ 在 $\mathbf F^{3}$ 中是线性无关的，但不是 $\mathbf F^{3}$ 的基，因为它不张成 $\mathbf F^{3}$；
\item 向量组 $(1 , 2) , (3 , 5) , (4 , 13)$ 张成 $\mathbf F^{2}$，但不是 $\mathbf F^{2}$ 的基，因为它们是线性相关的；
\item 向量组 $(1 , 1 , 0) , (0 , 0 , 1)$ 是 $\{(x , x , y) \in \mathbf F^{3} : x , y \in \mathbf F\}$ 的基；
\item 向量组 $(1 , - 1 , 0) , (1 , 0 , - 1)$ 是 $\{(x , y , z) \in \mathbf F^{3} : x + y + z = 0\}$ 的基；
\item 向量组 $1 , z , \ldots , z^{m}$ 是 $\mathcal P_{m} (\mathbf F)$ 的基。

\end{enumerate}
\begin{remark}
本题原文为“验证例2.27中的所有结论”。出于完整性考虑，这里将原书例2.27的所有结论摘录在上面。
\end{remark}
\end{exercise}
\noindent 来源：\href{https://github.com/szdytom/LADRSolutions/blob/main/sections/2B.typ}{P-2B-2}（原书 2B 第 2 题；方而静题解）。
\begin{solve}
对于 (a)，记这些向量为 $v_{1} , \ldots , v_{n}$，设 $a_{1} , \ldots , a_{n} \in \mathbf F$，使得

\[a_{1} v_{1} + \cdots + a_{n} v_{n} = 0\]

这立即给出 $a_{1} = \cdots = a_{n} = 0$，根据线性无关的定义（原书2.15），可知向量组 $v_{1} , \ldots , v_{n}$ 是线性无关的。设 $v = (x_{1} , \ldots , x_{n}) \in \mathbf F^{n}$，则

\[v = x_{1} v_{1} + \cdots + x_{n} v_{n}\]

因此，$v_{1} , \ldots , v_{n}$ 张成 $\mathbf F^{n}$。根据基的定义（原书2.26），可知向量组 $v_{1} , \ldots , v_{n}$ 是 $\mathbf F^{n}$ 的基。

对于 (b)，设 $a_{1} , a_{2} \in \mathbf F$，$v = (x_{1} , x_{2}) \in \mathbf F^{2}$，满足

\[v = a_{1} (1 , 2) + a_{2} (3 , 5)\]

求解 $a_{1} , a_{2}$，得到唯一的一组解是

\[\begin{cases}a_{1} = - 5 x_{1} + 3 x_{2}\\a_{2} = 2 x_{1} - x_{2}\end{cases}\]

这表明 $\mathbf F^{2}$ 中的每个向量都可以唯一地被表示为向量组 $(1 , 2) , (3 , 5)$ 的线性组合。所以，根据基的判定准则（原书2.28），向量组 $(1 , 2) , (3 , 5)$ 是 $\mathbf F^{2}$ 的基。

对于 (c)，设 $a_{1} , a_{2} \in \mathbf F$，满足

\[a_{1} (1 , 2 , - 4) + a_{2} (7 , - 5 , 6) = 0\]

这给出 $a_{1} = a_{2} = 0$，因此向量组 $(1 , 2 , - 4) , (7 , - 5 , 6)$ 是线性无关的。反证假设向量组 $(1 , 2 , - 4) , (7 , - 5 , 6)$ 张成 $\mathbf F^{3}$，则存在 $a_{1} , a_{2} \in \mathbf F$，使得

\[(1 , 2 , 0) = a_{1} (1 , 2 , - 4) + a_{2} (7 , - 5 , 6)\]

然而，该方程组无解。因此，向量组 $(1 , 2 , - 4) , (7 , - 5 , 6)$ 不张成 $\mathbf F^{3}$，所以它们不是 $\mathbf F^{3}$ 的基。

对于 (d)，由 (b) 可知，向量组 $(1 , 2) , (3 , 5)$ 张成 $\mathbf F^{2}$，所以向量组 $(1 , 2) , (3 , 5) , (4 , 13)$ 也张成 $\mathbf F^{2}$。然而，注意到

\[(- 19) (1 , 2) + 5 (3 , 5) + 1 (4 , 13) = 0\]

所以，向量组 $(1 , 2) , (3 , 5) , (4 , 13)$ 不是线性无关的，因此它们不是 $\mathbf F^{2}$ 的基。

对于 (e)，设 $a_{1} , a_{2} \in \mathbf F$，$v = (x , x , y) \in \mathbf F^{3}$，满足

\[v = a_{1} (1 , 1 , 0) + a_{2} (0 , 0 , 1)\]

求解 $a_{1} , a_{2}$，得到唯一的一组解是

\[\begin{cases}a_{1} = x\\a_{2} = y\end{cases}\]

这表明 $\{(x , x , y) \in \mathbf F^{3} : x , y \in \mathbf F\}$ 中的每个向量都可以唯一地被表示为向量组 $(1 , 1 , 0) , (0 , 0 , 1)$ 的线性组合。所以，根据基的判定准则，向量组 $(1 , 1 , 0) , (0 , 0 , 1)$ 是 $\{(x , x , y) \in \mathbf F^{3} : x , y \in \mathbf F\}$ 的基。

对于 (f)，设 $a_{1} , a_{2} \in \mathbf F$，$v = (x , y , - x - y) \in \mathbf F^{3}$，满足

\[v = a_{1} (1 , - 1 , 0) + a_{2} (1 , 0 , - 1)\]

求解 $a_{1} , a_{2}$，得到唯一的一组解是

\[\begin{cases}a_{1} = - y\\a_{2} = x + y\end{cases}\]

这表明 $\{(x , y , z) \in \mathbf F^{3} : x + y + z = 0\}$ 中的每个向量都可以唯一地被表示为向量组 $(1 , - 1 , 0) , (1 , 0 , - 1)$ 的线性组合。所以，根据基的判定准则，向量组 $(1 , - 1 , 0) , (1 , 0 , - 1)$ 是 $\{(x , y , z) \in \mathbf F^{3} : x + y + z = 0\}$ 的基。

对于 (g)，根据多项式的次数的定义（原书2.11），立即可得 $1 , z , \ldots , z^{m}$ 张成 $\mathcal P_{m} (\mathbf F)$。现在反证假设 $1 , z , \ldots , z^{m}$ 不是线性无关的。即存在 $a_{0} , \ldots , a_{m} \in \mathbf F$，其中至少有一个不为 $0$，使得对于任意 $z \in \mathbf F$，有

\[a_{0} + a_{1} z + \cdots + a_{m} z^{m} = 0\]

现在找到编号最大的不为 $0$ 的系数 $\ell$，即 $a_{\ell} \ne 0$，且 $a_{k} = 0$ 对于 $\ell < k \le m$ 成立。取

\[z = \frac{\lvert a_{0}\rvert + \cdots + \lvert a_{\ell - 1}\rvert}{\lvert a_{\ell}\rvert} + 1\]

注意到 $z \ge 1$，于是对 $j \in \{0 , \ldots , \ell - 1\}$，有 $z^{j} \le z^{\ell - 1}$。使用三角不等式\footnote{见原书4.4。一般而言，我们不应该引用后面的定理，因为这将带来循环论证的风险。但是复数的性质这个定理完全独立，因此从逻辑上说，这里引用原书4.4是没有问题的。}，我们有

\[\lvert a_{0} + a_{1} z + \cdots + a_{\ell - 1} z^{\ell - 1}\rvert \le (\lvert a_{0}\rvert + \cdots + \lvert a_{\ell - 1}\rvert) z^{\ell - 1} < \lvert a_{\ell} z^{\ell}\rvert\]

于是 $a_{0} + a_{1} z + \cdots + a_{\ell - 1} z^{\ell - 1} \ne - a_{\ell} z^{\ell}$，这表明

\[a_{0} + a_{1} z + \cdots + a_{m} z^{m} \ne 0\]

矛盾，因此 $1 , z , \ldots , z^{m}$ 是线性无关的。根据基的定义，$1 , z , \ldots , z^{m}$ 是 $\mathcal P_{m} (\mathbf F)$ 的基。
\end{solve}

\begin{exercise}\label{prob:P-2B-3}
\begin{enumerate}[label=(\alph*)]
\item 设 $U$ 为 $\mathbf R^{5}$ 的子空间，定义为
\[U = \{(x_{1} , x_{2} , x_{3} , x_{4} , x_{5}) \in \mathbf R^{5} : x_{1} = 3 x_{2} \text{且} x_{3} = 7 x_{4}\}\]
求 $U$ 的一个基；

\item 将 (a) 中的基扩充为 $\mathbf R^{5}$ 的一个基；

\item 求 $\mathbf R^{5}$ 的一个子空间 $W$，使得 $\mathbf R^{5} = U \oplus W$。
\end{enumerate}
\end{exercise}
\noindent 来源：\href{https://github.com/szdytom/LADRSolutions/blob/main/sections/2B.typ}{P-2B-3}（原书 2B 第 3 题；方而静题解）。
\begin{solve}
对于 (a)，令

\[u_{1} = (3 , 1 , 0 , 0 , 0) , \quad u_{2} = (0 , 0 , 7 , 1 , 0) , \quad u_{3} = (0 , 0 , 0 , 0 , 1)\]

下面说明 $u_{1} , u_{2} , u_{3}$ 是 $U$ 的一个基。设 $a_{1} , a_{2} , a_{3} \in \mathbf R$，$v = (3 x , x , 7 y , y , z) \in U$，满足

\[v = a_{1} u_{1} + a_{2} u_{2} + a_{3} u_{3}\]

求解 $a_{1} , a_{2} , a_{3}$，得到唯一的一组解是

\[\begin{cases}a_{1} = x\\a_{2} = y\\a_{3} = z\end{cases}\]

这表明 $U$ 中的每个向量都可以唯一地被表示为向量组 $u_{1} , u_{2} , u_{3}$ 的线性组合。所以，根据基的判定准则（原书2.28），向量组 $u_{1} , u_{2} , u_{3}$ 是 $U$ 的基。

对于 (b)，令

\[u_{4} = (1 , 0 , 0 , 0 , 0) , \quad u_{5} = (0 , 0 , 1 , 0 , 0)\]

下面说明 $u_{1} , \ldots , u_{5}$ 是 $\mathbf R^{5}$ 的一个基。设 $a_{1} , \ldots , a_{5} \in \mathbf R$，$v = (x_{1} , \ldots , x_{5}) \in \mathbf R^{5}$，满足

\[v = a_{1} u_{1} + \cdots + a_{5} u_{5}\]

求解 $a_{1} , \ldots , a_{5}$，得到唯一的一组解是

\[\begin{cases}a_{1} = x_{2}\\a_{2} = x_{4}\\a_{3} = x_{5}\\a_{4} = x_{1} - 3 x_{2}\\a_{5} = x_{3} - 7 x_{4}\end{cases}\]

这表明 $\mathbf R^{5}$ 中的每个向量都可以唯一地被表示为向量组 $u_{1} , \ldots , u_{5}$ 的线性组合。所以，根据基的判定准则，向量组 $u_{1} , \ldots , u_{5}$ 是 $\mathbf R^{5}$ 的基。

对于 (c)，令

\[w = \operatorname{span} (u_{4} , u_{5})\]

我们首先说明，$\mathbf R^{5} = U + W$。由于向量组 $u_{1} , \ldots , u_{5}$ 张成 $\mathbf R^{5}$，因此任意向量 $v \in \mathbf R^{5}$ 都可以被表示为

\[v = (a_{1} u_{1} + a_{2} u_{2} + a_{3} u_{3}) + (a_{4} u_{4} + a_{5} u_{5})\]

注意到 $a_{1} u_{1} + a_{2} u_{2} + a_{3} u_{3} \in U$，且 $a_{4} u_{4} + a_{5} u_{5} \in W$，故 $\mathbf R^{5} = U + W$。

设 $v \in U \cap W$。则存在标量 $a_{1} , \ldots , a_{5}$，满足

\[a_{1} u_{1} + a_{2} u_{2} + a_{3} u_{3} = v = a_{4} u_{4} + a_{5} u_{5}\]

于是

\[a_{1} u_{1} + a_{2} u_{2} + a_{3} u_{3} - a_{4} u_{4} - a_{5} u_{5} = 0\]

由于 $u_{1} , \ldots , u_{5}$ 是线性无关的（见上面 (b) 的证明），因此 $a_{1} = \cdots = a_{5} = 0$。这表明 $v = 0$，因此 $U \cap W = \{0\}$。根据“两个子空间的直和”（原书1.46），我们得到 $\mathbf R^{5} = U \oplus W$。
\end{solve}

\begin{exercise}\label{prob:P-2B-4}
\begin{enumerate}[label=(\alph*)]
\item 设 $U$ 为 $\mathbf C^{5}$ 的子空间，定义为
\[U = \{(z_{1} , z_{2} , z_{3} , z_{4} , z_{5}) \in \mathbf C^{5} : 6 z_{1} = z_{2} \text{且} z_{3} + 2 z_{4} + 3 z_{5} = 0\}\]
求 $U$ 的一个基；

\item 将 (a) 中的基扩充为 $\mathbf C^{5}$ 的一个基；

\item 求 $\mathbf C^{5}$ 的一个子空间 $W$，使得 $\mathbf C^{5} = U \oplus W$。
\end{enumerate}
\end{exercise}
\noindent 来源：\href{https://github.com/szdytom/LADRSolutions/blob/main/sections/2B.typ}{P-2B-4}（原书 2B 第 4 题；方而静题解）。
\begin{solve}
对于 (a)，令

\[u_{1} = (1 , 6 , 0 , 0 , 0) , \quad u_{2} = (0 , 0 , - 2 , 1 , 0) , \quad u_{3} = (0 , 0 , - 3 , 0 , 1)\]

下面说明 $u_{1} , u_{2} , u_{3}$ 是 $U$ 的一个基。设 $a_{1} , a_{2} , a_{3} \in \mathbf C$，$v = (x , 6 x , y , z , - \frac{1}{3} (y + 2 z)) \in U$，满足

\[v = a_{1} u_{1} + a_{2} u_{2} + a_{3} u_{3}\]

求解 $a_{1} , a_{2} , a_{3}$，得到唯一的一组解是

\[\begin{cases}a_{1} = x\\a_{2} = z\\a_{3} = - \frac{1}{3} (y + 2 z)\end{cases}\]

这表明 $U$ 中的每个向量都可以唯一地被表示为向量组 $u_{1} , u_{2} , u_{3}$ 的线性组合。所以，根据基的判定准则（原书2.28），向量组 $u_{1} , u_{2} , u_{3}$ 是 $U$ 的基。

对于 (b)，令

\[u_{4} = (1 , 0 , 0 , 0 , 0) , \quad u_{5} = (0 , 0 , 0 , 1 , 0)\]

下面说明 $u_{1} , \ldots , u_{5}$ 是 $\mathbf C^{5}$ 的一个基。设 $a_{1} , \ldots , a_{5} \in \mathbf C$，$v = (z_{1} , \ldots , z_{5}) \in \mathbf C^{5}$，满足

\[v = a_{1} u_{1} + \cdots + a_{5} u_{5}\]

求解 $a_{1} , \ldots , a_{5}$，得到唯一的一组解是

\[\begin{cases}a_{1} = \frac{1}{6} z_{2}\\a_{2} = z_{4}\\a_{3} = z_{5}\\a_{4} = z_{1} - \frac{1}{6} z_{2}\\a_{5} = z_{3} + 2 z_{4} + 3 z_{5}\end{cases}\]

这表明 $\mathbf C^{5}$ 中的每个向量都可以唯一地被表示为向量组 $u_{1} , \ldots , u_{6}$ 的线性组合。所以，根据基的判定准则，向量组 $u_{1} , \ldots , u_{6}$ 是 $\mathbf C^{5}$ 的基。

对于 (c)，令

\[W = \operatorname{span} (u_{4} , u_{5})\]

我们首先说明，$\mathbf C^{5} = U + W$。由于向量组 $u_{1} , \ldots , u_{5}$ 张成 $\mathbf C^{5}$，因此任意向量 $v \in \mathbf C^{5}$ 都可以被表示为

\[v = (a_{1} u_{1} + a_{2} u_{2} + a_{3} u_{3}) + (a_{4} u_{4} + a_{5} u_{5})\]

注意到 $a_{1} u_{1} + a_{2} u_{2} + a_{3} u_{3} \in U$，且 $a_{4} u_{4} + a_{5} u_{5} \in W$，故 $\mathbf C^{5} = U + W$。

设 $v \in U \cap W$。则存在标量 $a_{1} , \ldots , a_{5}$，满足

\[a_{1} u_{1} + a_{2} u_{2} + a_{3} u_{3} = v = a_{4} u_{4} + a_{5} u_{5}\]

于是

\[a_{1} u_{1} + a_{2} u_{2} + a_{3} u_{3} - a_{4} u_{4} - a_{5} u_{5} = 0\]

由于 $u_{1} , \ldots , u_{5}$ 是线性无关的（见上面 (b) 的证明），因此 $a_{1} = \cdots = a_{5} = 0$。这表明 $v = 0$，因此 $U \cap W = \{0\}$。根据“两个子空间的直和”（原书1.46），我们得到 $\mathbf C^{5} = U \oplus W$。
\end{solve}

\begin{exercise}\label{prob:P-2B-5}
设 $V$ 是有限维向量空间，$U , W$ 是 $V$ 的子空间，且 $V = U + W$。证明：$V$ 有一个由 $U \cup W$ 中的向量组成的基。
\end{exercise}
\noindent 来源：\href{https://github.com/szdytom/LADRSolutions/blob/main/sections/2B.typ}{P-2B-5}（原书 2B 第 5 题；方而静题解）。
\begin{solve}
设 $u_{1} , \ldots , u_{m} \in U$ 是 $U$ 的一组基，$w_{1} , \ldots , w_{\ell} \in W$ 是 $W$ 的一组基。由于 $V = U + W$，因此任意向量 $v \in V$ 都可以被表示为

\[v = (a_{1} u_{1} + \cdots + a_{m} u_{m}) + (b_{1} w_{1} + \cdots + b_{n} w_{\ell})\]

其中 $a_{1} , \ldots , a_{m} , b_{1} , \ldots , b_{\ell} \in \mathbf F$。这表明

\[V = \operatorname{span} (u_{1} , \ldots , u_{m} , w_{1} , \ldots w_{\ell})\]

由于每个张成组都包含基（原书2.30），因此 $V$ 有一个由 $U \cup W$ 中的向量组成的基。
\end{solve}

\begin{exercise}\label{prob:P-2B-6}
证明或证伪：如果 $p_{0} , p_{1} , p_{2} , p_{3}$ 是 $\mathcal P_{3} (\mathbf F)$ 中的向量组，该组中的每个多项式次数都不为 $2$，那么 $p_{0} , p_{1} , p_{2} , p_{3}$ 不是 $\mathcal P_{3} (\mathbf F)$ 的基。
\end{exercise}
\noindent 来源：\href{https://github.com/szdytom/LADRSolutions/blob/main/sections/2B.typ}{P-2B-6}（原书 2B 第 6 题；方而静题解）。
\begin{solve}
取 $\mathbf F \to \mathbf F$ 上的函数 $p_{0} , \ldots , p_{3}$

\[\begin{aligned}p_{0} : & z \to 1 \\ p_{1} : & z \to z \\ p_{2} : & z \to z^{2} + z^{3} \\ p_{3} : & z \to z^{3}\end{aligned}\]

设 $p \in \mathcal P_{3} (\mathbf F)$，则根据多项式的次数的定义（原书2.11），存在 $a_{0} , \ldots , a_{3} \in \mathbf F$，使得对于任意 $z \in \mathbf F$，有

\[p (z) = a_{0} + a_{1} z + a_{2} z^{2} + a_{3} z^{3}\]

于是，

\[p = a_{0} p_{0} + a_{1} p_{1} + a_{2} p_{2} + (a_{3} - a_{2}) p_{3}\]

这表明 $p$ 可以被表示为向量组 $p_{0} , p_{1} , p_{2} , p_{3}$ 的线性组合，即 $p_{0} , p_{1} , p_{2} , p_{3}$ 张成 $\mathcal P_{3} (\mathbf F)$。

设 $a_{0} , a_{1} , a_{2} , a_{3} \in \mathbf F$，满足

\[a_{0} p_{0} + a_{1} p_{1} + a_{2} p_{2} + a_{3} p_{3} = 0\]

即对于任意 $z \in \mathbf F$，有

\[a_{0} + a_{1} z + a_{2} z^{2} + (a_{2} + a_{3}) z^{3} = 0\]

根据多项式系数的唯一性，我们有 $a_{0} = a_{1} = a_{2} = a_{3} = 0$。这表明向量组 $p_{0} , p_{1} , p_{2} , p_{3}$ 是线性无关的。

综上所述，$p_{0} , p_{1} , p_{2} , p_{3}$ 是 $\mathcal P_{3} (\mathbf F)$ 的基，故原命题不成立。
\end{solve}

\begin{exercise}\label{prob:P-2B-7}
设 $v_{1} , v_{2} , v_{3} , v_{4}$ 是 $V$ 的基，证明：向量组

\[v_{1} + v_{2} , v_{2} + v_{3} , v_{3} + v_{4} , v_{4}\]

也是 $V$ 的基。
\end{exercise}
\noindent 来源：\href{https://github.com/szdytom/LADRSolutions/blob/main/sections/2B.typ}{P-2B-7}（原书 2B 第 7 题；方而静题解）。
\begin{solve}
设 $v \in V$，注意到

\[v = a_{1} v_{1} + a_{2} v_{2} + a_{3} v_{3} + a_{4} v_{4}\]

其中 $a_{1} , a_{2} , a_{3} , a_{4} \in \mathbf F$。另一方面，设 $b_{1} , b_{2} , b_{3} , b_{4} \in \mathbf F$，满足

\[v = b_{1} (v_{1} + v_{2}) + b_{2} (v_{2} + v_{3}) + b_{3} (v_{3} + v_{4}) + b_{4} v_{4}\]

则

\[v = (b_{1} + b_{2}) v_{1} + (b_{2} + b_{3}) v_{2} + (b_{3} + b_{4}) v_{3} + b_{4} v_{4}\]

由于 $v_{1} , v_{2} , v_{3} , v_{4}$ 是基，根据基的判定准则（原书2.28），$v_{1} , v_{2} , v_{3} , v_{4}$ 的系数只能对应相等，即

\[\begin{cases}a_{1} = b_{1} + b_{2}\\a_{2} = b_{2} + b_{3}\\a_{3} = b_{3} + b_{4}\\a_{4} = b_{4}\end{cases}\]

求解 $b_{1} , b_{2} , b_{3} , b_{4}$，得到唯一的一组解是

\[\begin{cases}b_{1} = a_{1} - a_{2}\\b_{2} = a_{2} - a_{3}\\b_{3} = a_{3} - a_{4}\\b_{4} = a_{4}\end{cases}\]

这表明 $v$ 可以唯一地被表示为向量组 $v_{1} + v_{2} , v_{2} + v_{3} , v_{3} + v_{4} , v_{4}$ 的线性组合。因此，向量组 $v_{1} + v_{2} , v_{2} + v_{3} , v_{3} + v_{4} , v_{4}$ 是 $V$ 的基。
\end{solve}

\begin{exercise}\label{prob:P-2B-8}
证明或证伪：设向量组 $v_{1} , v_{2} , v_{3} , v_{4}$ 是 $V$ 的基，且子空间 $U$ 满足 $v_{1} , v_{2} \in U$，而 $v_{3} \notin U$ 和 $v_{4} \notin U$，则 $v_{1} , v_{2}$ 是 $U$ 的基。
\end{exercise}
\noindent 来源：\href{https://github.com/szdytom/LADRSolutions/blob/main/sections/2B.typ}{P-2B-8}（原书 2B 第 8 题；方而静题解）。
\begin{solve}
令 $V = \mathbf R^{4}$，且

\[v_{1} = (1 , 0 , 0 , 0) , \quad v_{2} = (0 , 1 , 0 , 0) , \quad v_{3} = (0 , 0 , 1 , 0) , \quad v_{4} = (0 , 0 , 0 , 1)\]

注意到，集合

\[U = \{(x , y , z , z) \in \mathbf R^{4} : x , y , z \in \mathbf R\}\]

是 $\mathbf R^{4}$ 的子空间,且满足 $v_{1} , v_{2} \in U$，而 $v_{3} \notin U$ 和 $v_{4} \notin U$。然而，$v_{1} , v_{2}$ 不是 $U$ 的基，因为向量 $(0 , 0 , 1 , 1) \in U$ 不是 $\operatorname{span} (u_{1} , u_{2})$ 中的元素。由此可知，原命题不成立。
\end{solve}

\begin{exercise}\label{prob:P-2B-9}
设 $v_{1} , \ldots , v_{m}$ 是 $V$ 上的向量组，对于 $k \in \{1 , \ldots , m\}$，定义

\[w_{k} = v_{1} + \cdots + v_{k}\]

证明：向量组 $v_{1} , \ldots , v_{m}$ 是 $V$ 的基，当且仅当向量组 $w_{1} , \ldots , w_{m}$ 是 $V$ 的基。
\end{exercise}
\noindent 来源：\href{https://github.com/szdytom/LADRSolutions/blob/main/sections/2B.typ}{P-2B-9}（原书 2B 第 9 题；方而静题解）。
\begin{solve}
由习题~\ref{app:E-cumulative-sum-span} 和习题~\ref{app:E-cumulative-sum-independence} 得证。
\end{solve}

\begin{exercise}\label{prob:P-2B-10}
设 $U$ 和 $W$ 是 $V$ 的子空间，且 $V = U \oplus W$。又设 $u_{1} , \ldots , u_{m}$ 是 $U$ 的基，$w_{1} , \ldots , w_{n}$ 是 $W$ 的基。证明：向量组

\[u_{1} , \ldots , u_{m} , w_{1} , \ldots , w_{n}\]

是 $V$ 的基。
\end{exercise}
\noindent 来源：\href{https://github.com/szdytom/LADRSolutions/blob/main/sections/2B.typ}{P-2B-10}（原书 2B 第 10 题；方而静题解）。
\begin{solve}
设 $v \in V$，由于 $V = U \oplus W$，因此存在唯一的 $u_{1} , \ldots , u_{m} \in U$ 和 $w_{1} , \ldots , w_{n} \in W$，使得

\[v = (a_{1} u_{1} + \cdots + a_{m} u_{m}) + (b_{1} w_{1} + \cdots + b_{n} w_{n})\]

其中 $a_{1} , \ldots , a_{m} , b_{1} , \ldots , b_{n} \in \mathbf F$，这表明 $u_{1} , \ldots , u_{m} , w_{1} , \ldots w_{n}$ 张成 $V$。

另一方面，设 $a_{1} , \ldots , a_{m} , b_{1} , \ldots , b_{n} \in \mathbf F$，满足

\[(a_{1} u_{1} + \cdots + a_{m} u_{m}) + (b_{1} w_{1} + \cdots + b_{n} w_{n}) = 0\]

根据直和的条件（原书1.45），必须有

\[\begin{aligned}a_{1} u_{1} & + \cdots + a_{m} u_{m} & = 0 \\ b_{1} w_{1} & + \cdots + b_{n} w_{n} & = 0\end{aligned}\]

由于 $u_{1} , \ldots , u_{m}$ 是 $U$ 的基，$w_{1} , \ldots , w_{n}$ 是 $W$ 的基，根据基的定义（原书2.26），我们有 $a_{1} = \cdots = a_{m} = b_{1} = \cdots = b_{n} = 0$。这表明向量组 $u_{1} , \ldots , u_{m} , w_{1} , \ldots , w_{n}$ 是线性无关的。

综上所述，向量组 $u_{1} , \ldots , u_{m} , w_{1} , \ldots , w_{n}$ 是 $V$ 的基。
\end{solve}

\begin{exercise}\label{prob:P-2B-11}
设 $V$ 是实向量空间，证明：若 $v_{1} , \ldots , v_{n}$ 是 $V$（视为实向量空间）的基，则 $v_{1} , \ldots , v_{n}$ 也是其复化 $V_{\mathbf C}$（视为复向量空间）的基。

\begin{remark}
复化 $V_{\mathbf C}$ 的定义见习题~\ref{app:E-vector-dspace-complexification}。
\end{remark}
\end{exercise}
\noindent 来源：\href{https://github.com/szdytom/LADRSolutions/blob/main/sections/2B.typ}{P-2B-11}（原书 2B 第 11 题；方而静题解）。
\begin{solve}
对于 $u + \mathrm i v \in V_{\mathbf C}$，由于 $v_{1} , \ldots , v_{n}$ 是 $V$ 的基，可以找到 $a_{1} , \ldots , a_{n} \in \mathbf R$ 和 $b_{1} , \ldots , b_{n} \in \mathbf R$，使得

\[\begin{aligned}u & = a_{1} v_{1} & + & \cdots + a_{n} v_{n} \\ v & = b_{1} v_{1} & + & \cdots + b_{n} v_{n}\end{aligned}\]

于是

\[u + \mathrm i v = (a_{1} + \mathrm i b_{1}) v_{1} + \cdots + (a_{n} + \mathrm i b_{n}) v_{n}\]

这表明 $v_{1} , \ldots , v_{n}$ 张成 $V_{\mathbf C}$。

另一方面，设 $a_{1} + \mathrm i b_{1} , \ldots , a_{n} + \mathrm i b_{n} \in \mathbf C$，满足

\[(a_{1} + \mathrm i b_{1}) v_{1} + \cdots + (a_{n} + \mathrm i b_{n}) v_{n} = 0 + \mathrm i 0\]

根据 $V_{\mathbf C}$ 上标量乘法的定义，这相当于

\[(a_{1} v_{1} + \cdots + a_{n} v_{n}) + \mathrm i (b_{1} v_{1} + \cdots + b_{n} v_{n}) = 0 + \mathrm i 0\]

更进一步，必须有

\[\begin{aligned}a_{1} v_{1} & + \cdots + a_{n} v_{n} & = 0 \\ b_{1} v_{1} & + \cdots + b_{n} v_{n} & = 0\end{aligned}\]

由于 $v_{1} , \ldots , v_{n}$ 是 $V$ 的基，根据基的定义（原书2.26），我们有 $a_{1} = \cdots = a_{n} = b_{1} = \cdots = b_{n} = 0$。这表明向量组 $v_{1} , \ldots , v_{n}$ 是线性无关的。

综上所述，向量组 $v_{1} , \ldots , v_{n}$ 是 $V_{\mathbf C}$ 的基。
\end{solve}

\subsection{维数}\label{app:section-2C}
\begin{exercise}\label{prob:P-2C-1}
证明：$\mathbf R^{2}$ 的子空间恰有 $\{0\}$，$\mathbf R^{2}$ 中所有过原点的直线，以及 $\mathbf R^{2}$ 本身。
\end{exercise}
\noindent 来源：\href{https://github.com/szdytom/LADRSolutions/blob/main/sections/2C.typ}{P-2C-1}（原书 2C 第 1 题；方而静题解）。
\begin{solve}
设 $U$ 是 $\mathbf R^{2}$ 的子空间，根据子空间的维数性质（原书2.37），有 $\dim U \in \{0 , 1 , 2\}$。

\begin{itemize}
\item 如果 $\dim U = 0$，则 $U = \operatorname{span} () = \{0\}$；

\item 如果 $\dim U = 1$，则存在一个非零向量 $v$，使得 $U = \operatorname{span} (v)$，即 $U$ 是过原点的直线；

\item 如果 $\dim U = 2$，则根据满维数的子空间等于整个空间（原书2.39），$U = \mathbf R^{2}$。

综上所述，$\mathbf R^{2}$ 的子空间恰有 $\{0\}$，$\mathbf R^{2}$ 中所有过原点的直线，以及 $\mathbf R^{2}$ 本身。
\end{itemize}
\end{solve}

\begin{exercise}\label{prob:P-2C-2}
证明：$\mathbf R^{3}$ 的子空间恰有 $\{0\}$，$\mathbf R^{3}$ 中所有过原点的直线，$\mathbf R^{3}$ 中所有过原点的平面，以及 $\mathbf R^{3}$ 本身。
\end{exercise}
\noindent 来源：\href{https://github.com/szdytom/LADRSolutions/blob/main/sections/2C.typ}{P-2C-2}（原书 2C 第 2 题；方而静题解）。
\begin{solve}
设 $U$ 是 $\mathbf R^{3}$ 的子空间，根据子空间的维数性质（原书2.37），有 $\dim U \in \{0 , 1 , 2 , 3\}$。

\begin{itemize}
\item 如果 $\dim U = 0$，则 $U = \operatorname{span} () = \{0\}$；

\item 如果 $\dim U = 1$，则存在一个非零向量 $v$，使得 $U = \operatorname{span} (v)$，即 $U$ 是过原点的直线；

\item 如果 $\dim U = 2$，则存在两个线性无关的向量 $v_{1} , v_{2}$，使得 $U = \operatorname{span} (v_{1} , v_{2})$，即 $U$ 是过原点的平面；

\item 如果 $\dim U = 3$，则根据满维数的子空间等于整个空间（原书2.39），$U = \mathbf R^{3}$。

综上所述，$\mathbf R^{3}$ 的子空间恰有 $\{0\}$，$\mathbf R^{3}$ 中所有过原点的直线，$\mathbf R^{3}$ 中所有过原点的平面，以及 $\mathbf R^{3}$ 本身。
\end{itemize}
\end{solve}

\begin{exercise}\label{prob:P-2C-3}
\begin{enumerate}[label=(\alph*)]
\item 令 $U = \{p \in \mathcal P_{4} (\mathbf F) : p (6) = 0\}$，求 $U$ 的一组基；
\item 将 (a) 中的基扩充为 $\mathcal P_{4} (\mathbf F)$ 的基；
\item 求 $\mathcal P_{4} (\mathbf F)$ 的一个子空间 $W$，使得 $\mathcal P_{4} (\mathbf F) = U \oplus W$。
\end{enumerate}
\end{exercise}
\noindent 来源：\href{https://github.com/szdytom/LADRSolutions/blob/main/sections/2C.typ}{P-2C-3}（原书 2C 第 3 题；方而静题解）。
\begin{solve}
对于 (a)，我们的思路是，考虑到 $p (6) = 0$ 意味着任意 $p \in U$ 都可以表示为 $p (z) = (z - 6) q (z)$，其中 $q \in \mathcal P_{3} (\mathbf F)$。因此，我们很容易猜测 $z - 6 , z^{2} - 6 z , z^{3} - 6 z^{2} , z^{4} - 6 z^{3}$\footnote{将这些表达式看作关于 $z$ 的函数，即这里表示的是 $z \mapsto \cdots$，下同。}是 $U$ 的一组基。

为了说明这一点，我们首先证明 $z - 6 , z^{2} - 6 z , z^{3} - 6 z^{2} , z^{4} - 6 z^{3}$ 是线性无关的。设 $a , b , c , d \in \mathbf F$，满足对于任意 $z \in \mathbf F$，

\[a (z - 6) + b (z^{2} - 6 z) + c (z^{3} - 6 z^{2}) + d (z^{4} - 6 z^{3}) = 0\]

整理得到

\[- 6 a + (a - 6 b) z + (b - 6 c) z^{2} + (c - 6 d) z^{3} + d z^{4} = 0\]

根据多项式系数的唯一性，我们有

\[\begin{cases}- 6 a = 0\\a - 6 b = 0\\b - 6 c = 0\\c - 6 d = 0\\d = 0\end{cases}\]

由此可得 $a = b = c = d = 0$，因此 $z - 6 , z^{2} - 6 z , z^{3} - 6 z^{2} , z^{4} - 6 z^{3}$ 是线性无关的。

根据子空间的维数的性质，我们知道 $\dim U \le \dim \mathcal P_{4} (\mathbf F) = 5$。注意到 $z \mapsto z \notin U$，故 $U \ne \mathcal P_{4} (\mathbf F)$，于是 $\dim U \le 4$。另一方面，将 $U$ 的一组基看作张成组，根据“线性无关组的长度 $\le$ 张成组的长度”（原书2.22），我们得到 $\dim U \ge 4$。于是只能有 $\dim U = 4$。

代入 $z = 6$，立即可以验证 $z - 6 , z^{2} - 6 z , z^{3} - 6 z^{2} , z^{4} - 6 z^{3} \in U$，根据长度恰当的线性无关组是基（原书2.38），我们得出 $z - 6 , z^{2} - 6 z , z^{3} - 6 z^{2} , z^{4} - 6 z^{3}$ 是 $U$ 的一组基。

对于 (b)，注意到 $1 \notin U$，根据习题~\ref{app:E-when-vector-list-append-remains-indep}，可得向量组 $z - 6 , z^{2} - 6 z , z^{3} - 6 z^{2} , z^{4} - 6 z^{3} , 1$ 是线性无关的，进一步地，向量组 $z - 6 , z^{2} - 6 z , z^{3} - 6 z^{2} , z^{4} - 6 z^{3} , 1$ 是长度恰当（$\dim \mathcal P_{4} (\mathbf F) = 5$）的线性无关组，因此它是 $\mathcal P_{4} (\mathbf F)$ 的一组基。

对于 (c)，我们可以取 $W = \operatorname{span} (1)$。由 (b) 可知 $\mathcal P_{4} (\mathbf F) = U + W$，根据子空间之和的维数（原书2.43），我们有

\[\dim \mathcal P_{4} (\mathbf F) = \dim U + \dim W - \dim(U \cap W)\]

解得 $\dim(U \cap W) = 0$，即 $U \cap W = \{0\}$。根据“两个子空间的直和”（原书1.46），得到 $\mathcal P_{4} (\mathbf F) = U \oplus W$。
\end{solve}

\begin{exercise}\label{prob:P-2C-4}
\begin{enumerate}[label=(\alph*)]
\item 令 $U = \{p \in \mathcal P_{4} (\mathbf F) : p'' (6) = 0\}$，求 $U$ 的一组基；
\item 将 (a) 中的基扩充为 $\mathcal P_{4} (\mathbf F)$ 的基；
\item 求 $\mathcal P_{4} (\mathbf F)$ 的一个子空间 $W$，使得 $\mathcal P_{4} (\mathbf F) = U \oplus W$。
\end{enumerate}
\end{exercise}
\noindent 来源：\href{https://github.com/szdytom/LADRSolutions/blob/main/sections/2C.typ}{P-2C-4}（原书 2C 第 4 题；方而静题解）。
\begin{solve}
对于 (a)，向量组 $1 , z , z^{3} - 18 z^{2} , z^{4} - 216 z^{2}$ 是 $U$ 的一组基。为了说明这一点，我们首先证明 $1 , z , z^{3} - 18 z^{2} , z^{4} - 216 z^{2}$ 是线性无关的。设 $a , b , c , d \in \mathbf F$，满足对于任意 $z \in \mathbf F$，

\[a + b z + c (z^{3} - 18 z^{2}) + d (z^{4} - 216 z^{2}) = 0\]

整理得到

\[a + b z + (- 18 c - 216 d) z^{2} + c z^{3} + d z^{4} = 0\]

根据多项式系数的唯一性，我们有

\[\begin{cases}a = 0\\b = 0\\- 18 c - 216 d = 0\\c = 0\\d = 0\end{cases}\]

由此可得 $a = b = c = d = 0$，因此 $1 , z , z^{3} - 18 z^{2} , z^{4} - 216 z^{2}$ 是线性无关的。

根据子空间的维数的性质，我们知道 $\dim U \le \dim \mathcal P_{4} (\mathbf F) = 5$。注意到 $z \mapsto z^{2} \notin U$，故 $U \ne \mathcal P_{4} (\mathbf F)$，于是 $\dim U \le 4$。另一方面，将 $U$ 的一组基看作张成组，根据“线性无关组的长度 $\le$ 张成组的长度”（原书2.22），我们得到 $\dim U \ge 4$。于是只能有 $\dim U = 4$。

分别求二阶导，再代入 $z = 6$，可以验证 $1 , z , z^{3} - 18 z^{2} , z^{4} - 216 z^{2} \in U$，根据长度恰当的线性无关组是基（原书2.38），我们得出 $1 , z , z^{3} - 18 z^{2} , z^{4} - 216 z^{2}$ 是 $U$ 的一组基。

对于 (b)，注意到 $z^{2} \notin U$，根据习题~\ref{app:E-when-vector-list-append-remains-indep}，可得向量组 $1 , z , z^{3} - 18 z^{2} , z^{4} - 216 z^{2} , z^{2}$ 是线性无关的，进一步地，向量组 $1 , z , z^{3} - 18 z^{2} , z^{4} - 216 z^{2} , z^{2}$ 是长度恰当（$\dim \mathcal P_{4} (\mathbf F) = 5$）的线性无关组，因此它是 $\mathcal P_{4} (\mathbf F)$ 的一组基。

对于 (c)，我们可以取 $W = \operatorname{span} (z^{2})$。由 (b) 可知 $\mathcal P_{4} (\mathbf F) = U + W$，
根据子空间之和的维数（原书2.43），我们有

\[\dim \mathcal P_{4} (\mathbf F) = \dim U + \dim W - \dim(U \cap W)\]

解得 $\dim(U \cap W) = 0$，即 $U \cap W = \{0\}$。根据“两个子空间的直和”（原书1.46），得到 $\mathcal P_{4} (\mathbf F) = U \oplus W$。
\end{solve}

\begin{exercise}\label{prob:P-2C-5}\label{app:E-p2-eq-p5-poly4}
\begin{enumerate}[label=(\alph*)]
\item 令 $U = \{p \in \mathcal P_{4} (\mathbf F) : p (2) = p (5)\}$，求 $U$ 的一组基；
\item 将 (a) 中的基扩充为 $\mathcal P_{4} (\mathbf F)$ 的基；
\item 求 $\mathcal P_{4} (\mathbf F)$ 的一个子空间 $W$，使得 $\mathcal P_{4} (\mathbf F) = U \oplus W$。
\end{enumerate}
\end{exercise}
\noindent 来源：\href{https://github.com/szdytom/LADRSolutions/blob/main/sections/2C.typ}{P-2C-5}（原书 2C 第 5 题；方而静题解）。
\begin{solve}
对于 (a)，向量组 $1 , z^{2} - 7 z , z^{3} - 39 z , z^{4} - 203 z$ 是 $U$ 的一组基。为了说明这一点，我们首先证明 $1 , z^{2} - 7 z , z^{3} - 39 z , z^{4} - 203 z$ 是线性无关的。设 $a , b , c , d \in \mathbf F$，满足对于任意 $z \in \mathbf F$，

\[a + b (z^{2} - 7 z) + c (z^{3} - 39 z) + d (z^{4} - 203 z) = 0\]

整理得到

\[a + (- 7 b - 39 c - 203 d) z + b z^{2} + c z^{3} + d z^{4} = 0\]

根据多项式系数的唯一性，我们有

\[\begin{cases}a = 0\\- 7 b - 39 c - 203 d = 0\\b = 0\\c = 0\\d = 0\end{cases}\]

由此可得 $a = b = c = d = 0$，因此 $1 , z^{2} - 7 z , z^{3} - 39 z , z^{4} - 203 z$ 是线性无关的。

根据子空间的维数的性质，我们知道 $\dim U \le \dim \mathcal P_{4} (\mathbf F) = 5$。注意到 $z \mapsto z \notin U$，故 $U \ne \mathcal P_{4} (\mathbf F)$，于是 $\dim U \le 4$。另一方面，将 $U$ 的一组基看作张成组，根据“线性无关组的长度 $\le$ 张成组的长度”（原书2.22），我们得到 $\dim U \ge 4$。于是只能有 $\dim U = 4$。

分别代入 $z = 2$ 和 $z = 5$，可以验证 $1 , z^{2} - 7 z , z^{3} - 39 z , z^{4} - 203 z \in U$，根据长度恰当的线性无关组是基（原书2.38），我们得出 $1 , z^{2} - 7 z , z^{3} - 39 z , z^{4} - 203 z$ 是 $U$ 的一组基。

对于 (b)，注意到 $z \notin U$，根据习题~\ref{app:E-when-vector-list-append-remains-indep}，可得向量组 $1 , z^{2} - 7 z , z^{3} - 39 z , z^{4} - 203 z , z$ 是线性无关的，进一步地，向量组 $1 , z^{2} - 7 z , z^{3} - 39 z , z^{4} - 203 z , z$ 是长度恰当（$\dim \mathcal P_{4} (\mathbf F) = 5$）的线性无关组，因此它是 $\mathcal P_{4} (\mathbf F)$ 的一组基。

对于 (c)，我们可以取 $W = \operatorname{span} (z)$。由 (b) 可知 $\mathcal P_{4} (\mathbf F) = U + W$，根据子空间之和的维数（原书2.43），我们有

\[\dim \mathcal P_{4} (\mathbf F) = \dim U + \dim W - \dim(U \cap W)\]

解得 $\dim(U \cap W) = 0$，即 $U \cap W = \{0\}$。根据“两个子空间的直和”（原书1.46），得到 $\mathcal P_{4} (\mathbf F) = U \oplus W$。
\end{solve}

\begin{exercise}\label{prob:P-2C-6}
\begin{enumerate}[label=(\alph*)]
\item 令 $U = \{p \in \mathcal P_{4} (\mathbf F) : p (2) = p (5) = p (6)\}$，求 $U$ 的一组基；
\item 将 (a) 中的基扩充为 $\mathcal P_{4} (\mathbf F)$ 的基；
\item 求 $\mathcal P_{4} (\mathbf F)$ 的一个子空间 $W$，使得 $\mathcal P_{4} (\mathbf F) = U \oplus W$。
\end{enumerate}
\end{exercise}
\noindent 来源：\href{https://github.com/szdytom/LADRSolutions/blob/main/sections/2C.typ}{P-2C-6}（原书 2C 第 6 题；方而静题解）。
\begin{solve}
令 $h(z)=(z-2)(z-5)(z-6)$。条件 $p(2)=p(5)=p(6)$ 等价于存在 $c,a,b\in\mathbf F$ 使
\[p(z)=c+h(z)(a+bz).\]
事实上，减去公共函数值 $c$ 后的多项式有三个不同的根，故被 $h$ 整除；又因次数不超过 $4$，商的次数不超过 $1$。
\begin{enumerate}[label=(\alph*)]
\item $1,h,zh$ 是 $U$ 的一组基：它们张成 $U$，且次数分别为 $0,3,4$，比较最高次项可证线性无关。
\item 加入 $z,z^2$ 得到 $\mathcal P_4(\mathbf F)$ 的一组基，因为五个多项式的次数恰为 $0,1,2,3,4$。
\item 可取 $W=\operatorname{span}(z,z^2)$。由上述基的分组，$\mathcal P_4(\mathbf F)=U\oplus W$。
\end{enumerate}
\end{solve}

\begin{exercise}\label{prob:P-2C-7}
\begin{enumerate}[label=(\alph*)]
\item 令 $U = \{p \in \mathcal P_{4} (\mathbf F) : \int_{- 1}^{1} p = 0\}$，求 $U$ 的一组基；
\item 将 (a) 中的基扩充为 $\mathcal P_{4} (\mathbf F)$ 的基；
\item 求 $\mathcal P_{4} (\mathbf F)$ 的一个子空间 $W$，使得 $\mathcal P_{4} (\mathbf F) = U \oplus W$。
\end{enumerate}
\end{exercise}
\noindent 来源：\href{https://github.com/szdytom/LADRSolutions/blob/main/sections/2C.typ}{P-2C-7}（原书 2C 第 7 题；方而静题解）。
\begin{solve}
我们的思路是，考虑到

\begin{equation}\label{app:eq:2C-F-intt-poly4--1-to-1}\int_{- 1}^{1} a_{0} + a_{1} z + a_{2} z^{2} + a_{3} z^{3} + a_{4} z^{4} \,\mathrm d z = 2 a_{0} + \frac{2}{3} a_{2} + \frac{2}{5} a_{4}\end{equation}

我们很容易猜测 $z , 3 z^{2} - 1 , z^{3} , 5 z^{4} - 1$ 是 $U$ 的一组基。为了严格说明这一点，我们首先证明 $z , 3 z^{2} - 1 , z^{3} , 5 z^{4} - 1$ 是线性无关的。设 $a , b , c , d \in \mathbf F$，满足对于任意 $z \in \mathbf F$，

\[a z + b (3 z^{2} - 1) + c z^{3} + d (5 z^{4} - 1) = 0\]

整理得到

\[(- b - d) + a z + (3 b) z^{2} + c z^{3} + (5 d) z^{4} = 0\]

根据多项式系数的唯一性，我们有

\[\begin{cases}- b - d = 0\\a = 0\\3 b = 0\\c = 0\\5 d = 0\end{cases}\]

由此可得 $a = b = c = d = 0$，因此 $z , 3 z^{2} - 1 , z^{3} , 5 z^{4} - 1$ 是线性无关的。

根据子空间的维数的性质，我们知道 $\dim U \le \dim \mathcal P_{4} (\mathbf F) = 5$。注意到 $z \mapsto 1 \notin U$，故 $U \ne \mathcal P_{4} (\mathbf F)$，于是 $\dim U \le 4$。另一方面，将 $U$ 的一组基看作张成组，根据“线性无关组的长度 $\le$ 张成组的长度”（原书2.22），我们得到 $\dim U \ge 4$。于是只能有 $\dim U = 4$。

将系数代入式~\eqref{app:eq:2C-F-intt-poly4--1-to-1} 的积分结果，发现 $z , 3 z^{2} - 1 , z^{3} , 5 z^{4} - 1 \in U$，根据长度恰当的线性无关组是基（原书2.38），我们得出 $z , 3 z^{2} - 1 , z^{3} , 5 z^{4} - 1$ 是 $U$ 的一组基。

对于 (b)，注意到 $1 \notin U$，根据习题~\ref{app:E-when-vector-list-append-remains-indep}，可得向量组 $z , 3 z^{2} - 1 , z^{3} , 5 z^{4} - 1 , 1$ 是线性无关的，进一步地，向量组 $z , 3 z^{2} - 1 , z^{3} , 5 z^{4} - 1 , 1$ 是长度恰当（$\dim \mathcal P_{4} (\mathbf F) = 5$）的线性无关组，因此它是 $\mathcal P_{4} (\mathbf F)$ 的一组基。

对于 (c)，我们可以取 $W = \operatorname{span} (1)$。由 (b) 可知 $\mathcal P_{4} (\mathbf F) = U + W$，根据子空间之和的维数（原书2.43），我们有

\[\dim \mathcal P_{4} (\mathbf F) = \dim U + \dim W - \dim(U \cap W)\]

解得 $\dim(U \cap W) = 0$，即 $U \cap W = \{0\}$。根据“两个子空间的直和”（原书1.46），得到 $\mathcal P_{4} (\mathbf F) = U \oplus W$。
\end{solve}

\begin{exercise}\label{prob:P-2C-8}
设向量组 $v_{1} , \ldots , v_{m}$ 在 $V$ 中线性无关，$w \in V$，证明

\[\dim \operatorname{span} (v_{1} + w , \ldots , v_{m} + w) \ge m - 1\]
\end{exercise}
\noindent 来源：\href{https://github.com/szdytom/LADRSolutions/blob/main/sections/2C.typ}{P-2C-8}（原书 2C 第 8 题；方而静题解）。
\begin{solve}
当 $m=1$ 时结论是维数非负。设 $m\ge2$。向量 $v_2-v_1,\ldots,v_m-v_1$ 线性无关：若 $\sum_{j=2}^m a_j(v_j-v_1)=0$，则由 $v_1,\ldots,v_m$ 线性无关，所有 $a_j=0$。又
\[v_j-v_1=(v_j+w)-(v_1+w),\]
所以这 $m-1$ 个向量都在题给张成空间中，维数至少为 $m-1$。
进一步，维数恰为 $m-1$ 当且仅当 $w=\sum_{j=1}^m c_jv_j$ 且 $\sum_jc_j=-1$。若存在非平凡关系 $\sum_j a_j(v_j+w)=0$，则 $s=\sum_j a_j$ 不能为零，否则原向量组线性相关；于是 $w=-\sum_j(a_j/s)v_j$。反过来，所给条件使 $\sum_j c_j(v_j+w)=0$ 成为非平凡关系，结合下界即可取等。
\end{solve}

\begin{exercise}\label{prob:P-2C-9}\label{app:E-Poly-space-basis}
设 $m$ 是正整数，$p_{0} , \ldots , p_{m} \in \mathcal P (\mathbf F)$，其中 $p_{k}$ 的次数为 $k$，证明： $p_{0} , \ldots , p_{m}$ 是 $\mathcal P_{m} (\mathbf F)$ 的基。
\end{exercise}
\noindent 来源：\href{https://github.com/szdytom/LADRSolutions/blob/main/sections/2C.typ}{P-2C-9}（原书 2C 第 9 题；方而静题解）。
\begin{solve}
设 $\sum_{k=0}^m a_kp_k=0$。若系数不全为零，令 $j$ 是 $a_j\ne0$ 的最大下标；由于 $\deg p_j=j$ 而其他出现的多项式次数均小于 $j$，最高次项不可能抵消，矛盾。因此 $p_0,\ldots,p_m$ 线性无关。它们共有 $m+1=\dim\mathcal P_m(\mathbf F)$ 个，所以构成一组基。
\end{solve}

\begin{exercise}\label{prob:P-2C-10}
设 $m$ 是正整数，对于 $k \in \{0 , \ldots , m\}$，令

\[\begin{aligned}p_{k} : & \mathbf F \to \mathbf F \\ & z \mapsto z^{k} (1 - z)^{m - k}\end{aligned}\]

证明：$p_{0} , \ldots , p_{m}$ 是 $\mathcal P_{m} (\mathbf F)$ 的基。

\begin{remark}
这道习题中的基能够引出所谓\emph{伯恩斯坦多项式（Bernstein polynomials）}。你可以上网搜索，了解伯恩斯坦多项式如何用于逼近 $[0 , 1]$ 上的连续函数。
\end{remark}
\end{exercise}
\noindent 来源：\href{https://github.com/szdytom/LADRSolutions/blob/main/sections/2C.typ}{P-2C-10}（原书 2C 第 10 题；方而静题解）。
\begin{solve}
对每个 $j=0,\ldots,m$，由二项式定理有
\[z^j=z^j\bigl(z+(1-z)\bigr)^{m-j}=\sum_{k=j}^m\binom{m-j}{k-j}z^k(1-z)^{m-k}.\]
所以自然基 $1,z,\ldots,z^m$ 的每个元素都在 $p_0,\ldots,p_m$ 的张成空间内；这些多项式自身又属于 $\mathcal P_m(\mathbf F)$，故它们张成该空间。共有 $m+1=\dim\mathcal P_m(\mathbf F)$ 个生成向量，因此构成一组基。
\end{solve}

\begin{exercise}\label{prob:P-2C-11}
设 $U$ 和 $W$ 都是 $\mathbf C^{6}$ 的四维子空间，证明：在 $U \cap W$ 中存在两个向量，其中任意一个都不是另一个的标量倍。
\end{exercise}
\noindent 来源：\href{https://github.com/szdytom/LADRSolutions/blob/main/sections/2C.typ}{P-2C-11}（原书 2C 第 11 题；方而静题解）。
\begin{solve}
根据子空间之和的维数（原书2.43），我们有

\[\dim(U + W) = \dim U + \dim W - \dim(U \cap W)\]

另一方面，根据子空间的维数性质（原书2.37），$\dim(U + W) \le \dim \mathbf C^{6} = 6$。因此，

\[\dim(U \cap W) \ge 2\]

设 $u_{1} , \ldots , u_{m}$ 是 $U \cap W$ 的一组基，其中 $m \ge 2$。则 $u_{1} , u_{2}$ 是线性无关的。根据习题~\ref{app:E-when-1-or-2-vectors-indep}, $u_{1} , u_{2}$ 中任意一个都不是另一个的标量倍。由此得证。
\end{solve}

\begin{exercise}\label{prob:P-2C-12}
设 $U$ 和 $W$ 都是 $\mathbf R^{8}$ 的子空间，$\dim U = 3$，$\dim W = 5$，且 $U + W = \mathbf R^{8}$，证明：$\mathbf R^{8} = U \oplus W$。
\end{exercise}
\noindent 来源：\href{https://github.com/szdytom/LADRSolutions/blob/main/sections/2C.typ}{P-2C-12}（原书 2C 第 12 题；方而静题解）。
\begin{solve}
根据子空间之和的维数（原书2.43），我们有

\[\dim(U + W) = \dim U + \dim W - \dim(U \cap W)\]

解得 $\dim(U \cap W) = 0$，即 $U \cap W = \{0\}$。根据“两个子空间的直和”（原书1.46），得到 $\mathbf R^{8} = U \oplus W$。
\end{solve}

\begin{exercise}\label{prob:P-2C-13}
设 $U$ 和 $W$ 都是 $\mathbf R^{9}$ 的五维子空间，证明：$U \cap W \ne \{0\}$。
\end{exercise}
\noindent 来源：\href{https://github.com/szdytom/LADRSolutions/blob/main/sections/2C.typ}{P-2C-13}（原书 2C 第 13 题；方而静题解）。
\begin{solve}
根据子空间之和的维数（原书2.43），我们有

\[\dim(U + W) = \dim U + \dim W - \dim(U \cap W)\]

另一方面，根据子空间的维数性质（原书2.37），$\dim(U + W) \le \dim \mathbf R^{9} = 9$。因此，

\[\dim(U \cap W) \ge 1\]

由此得证 $U \cap W \ne \{0\}$。
\end{solve}

\begin{exercise}\label{prob:P-2C-14}
设 $V$ 是 $10$ 维向量空间，$V_{1} , V_{2} , V_{3}$ 都是 $V$ 的子空间，$\dim V_{1} = \dim V_{2} = \dim V_{3} = 7$，证明：$V_{1} \cap V_{2} \cap V_{3} \ne \{0\}$。
\end{exercise}
\noindent 来源：\href{https://github.com/szdytom/LADRSolutions/blob/main/sections/2C.typ}{P-2C-14}（原书 2C 第 14 题；方而静题解）。
\begin{solve}
根据子空间之和的维数（原书2.43），我们有

\[\begin{aligned}\dim(V_{1} + V_{2}) & = \dim V_{1} + \dim V_{2} - \dim(V_{1} \cap V_{2}) \\ \dim((V_{1} + V_{2}) + V_{3}) & = \dim(V_{1} + V_{2}) + \dim V_{3} - \dim((V_{1} + V_{2}) \cap V_{3})\end{aligned}\]

上下两式相加，化简得

\[\begin{aligned}\dim(V_{1} + V_{2} + V_{3}) = & \dim V_{1} + \dim V_{2} + \dim V_{3} \\ & - \dim(V_{1} \cap V_{2}) - \dim((V_{1} + V_{2}) \cap V_{3})\end{aligned}\]

另一方面，根据子空间的维数性质（原书2.37），$\dim(V_{1} + V_{2} + V_{3}) \le \dim V = 10$。同时，考虑到 $(V_{1} + V_{2}) \cap V_{3} \subseteq V_{3}$，我们有 $\dim((V_{1} + V_{2}) \cap V_{3}) \le \dim V_{3} = 7$，因此

\[\dim(V_{1} \cap V_{2}) \ge 4\]

同时

\[\dim((V_{1} \cap V_{2}) + V_{3}) = \dim(V_{1} \cap V_{2}) + \dim V_{3} - \dim(V_{1} \cap V_{2} \cap V_{3})\]

结合 $\dim((V_{1} \cap V_{2}) + V_{3}) \le \dim V = 10$，于是

\[4 \le \dim(V_{1} \cap V_{2}) \le 3 + \dim(V_{1} \cap V_{2} \cap V_{3})\]

所以 $\dim(V_{1} \cap V_{2} \cap V_{3}) \ge 1$，由此得证 $V_{1} \cap V_{2} \cap V_{3} \ne \{0\}$。
\end{solve}

\begin{exercise}\label{prob:P-2C-15}
设 $V$ 是有限维向量空间，$V_{1} , V_{2} , V_{3}$ 都是 $V$ 的子空间，$\dim V_{1} + \dim V_{2} + \dim V_{3} > 2 \dim V$，证明：$V_{1} \cap V_{2} \cap V_{3} \ne \{0\}$。
\end{exercise}
\noindent 来源：\href{https://github.com/szdytom/LADRSolutions/blob/main/sections/2C.typ}{P-2C-15}（原书 2C 第 15 题；方而静题解）。
\begin{solve}
根据子空间之和的维数（原书2.43），我们有

\[\begin{aligned}\dim((V_{1} \cap V_{2}) + V_{3}) & = \dim(V_{1} \cap V_{2}) + \dim V_{3} - \dim(V_{1} \cap V_{2} \cap V_{3}) \\ \dim(V_{1} + V_{2}) & = \dim V_{1} + \dim V_{2} - \dim(V_{1} \cap V_{2})\end{aligned}\]

两式向加，移项得

\[\begin{aligned}\dim((V_{1} \cap V_{2}) + V_{3}) & = & & \dim V_{1} + \dim V_{2} + \dim V_{3} \\ & & & - \dim(V_{1} + V_{2}) - \dim(V_{1} \cap V_{2} \cap V_{3}) \\ & > & & 2 \dim V - \dim(V_{1} + V_{2}) - \dim(V_{1} \cap V_{2} \cap V_{3})\end{aligned}\]

另一方面，根据子空间的维数性质（原书2.37），$\dim((V_{1} \cap V_{2}) + V_{3}) \le \dim V$，以及 $\dim(V_{1} + V_{2}) \le \dim V$，因此

\[\dim(V_{1} \cap V_{2} \cap V_{3}) > 0\]

由此得证 $V_{1} \cap V_{2} \cap V_{3} \ne \{0\}$。
\end{solve}

\begin{exercise}\label{prob:P-2C-16}
设 $V$ 是有限维向量空间，$U$ 是 $V$ 的子空间（$U \ne V$）。令 $n = \dim V$，$m = \dim U$，证明：$V$ 中存在这样 $n - m$ 个子空间，其中每个子空间的维数都是 $n - 1$，且它们的交集是 $U$。
\end{exercise}
\noindent 来源：\href{https://github.com/szdytom/LADRSolutions/blob/main/sections/2C.typ}{P-2C-16}（原书 2C 第 16 题；方而静题解）。
\begin{solve}
设 $u_{1} , \ldots , u_{m}$ 是 $U$ 的一组基，将其看作 $V$ 中的线性无关组，根据每个线性无关组都可被扩充为基（原书2.32），我们可以将 $u_{1} , \ldots , u_{m}$ 扩充为 $V$ 的一组基 $u_{1} , \ldots , u_{n}$。现在，对于 $k \in \{1 , \ldots , n - m\}$ 和 $j \in \{1 , \ldots , n - 1\}$，令

\[W_{k} = \operatorname{span} (u_{1} , \ldots , u_{m + k - 1} , u_{m + k + 1} , \ldots , u_{n})\]

即 $W_{k}$ 是向量组 $u_{1} , \ldots , u_{m}$ 中去掉第 $m + k$ 个向量后得到的向量组所张成的子空间。注意到 $\dim W_{k} = n - 1$，因此 $W_{1} , \ldots , W_{n - m}$ 是 $V$ 的 $n - m$ 个维数为 $n - 1$ 的子空间。同时，注意到 $U = \operatorname{span} (u_{1} , \ldots , u_{m}) \subseteq W_{k}$，因此

\[U \subseteq W_{1} \cap \cdots \cap W_{n - m}\]

另一方面，我们说明 $W_{1} \cap \cdots \cap W_{n - m} \subseteq U$。当 $n - m = 1$ 时，情况是平凡的。现在设 $v \in W_{1} \cap \cdots \cap W_{n - m}$，则对于 $k \in \{1 , \ldots , n - m\}$ 和 $i \in \{1 , \ldots , n\}$，存在 $a_{k , i}$，使得对于任意 $k \in \{1 , \ldots , n - m\}$，有

\[v = a_{k , 1} u_{1} + \cdots + a_{k , n} u_{n}\]

且 $a_{k , m + k} = 0$。现在，对于任意 $j \in \{2 , \ldots , n - m\}$，我们将 $v$ 在 $W_{1}$ 和 $W_{j}$ 中线性组合的表达式相减，得到

\[0 = (a_{1 , 1} - a_{j , 1}) u_{1} + \cdots + (a_{1 , n} - a_{j , n}) u_{n}\]

由于 $u_{1} , \ldots , u_{n}$ 是线性无关的，因此，对于任意 $i \in \{1 , \ldots , n\}$，$a_{1 , i} = a_{j , i}$。特别注意到，$a_{1 , m + j} = a_{j , m + j} = 0$，考虑到 $a_{1 , m + 1} = 0$，以及 $j \in \{2 , \ldots , n - m\}$ 选取的任意性，我们得到，对于任意 $j \in \{1 , \ldots , n - m\}$，$a_{1 , m + j} = 0$，即

\[v = a_{1 , 1} u_{1} + \cdots + a_{1 , m} u_{m} \in U\]

这表明 $W_{1} \cap \cdots \cap W_{n - m} \subseteq U$。

综上所述 $U = W_{1} \cap \cdots \cap W_{n - m}$。也就是说，$W_{1} , \ldots , W_{n - m}$ 就是所求的 $n - m$ 个子空间。
\end{solve}

\begin{exercise}\label{prob:P-2C-17}
设 $V_{1} , \ldots , V_{m}$ 都是 $V$ 的有限维子空间，证明：$V_{1} + \cdots + V_{m}$ 是有限维的，且

\[\dim(V_{1} + \cdots + V_{m}) \le \dim V_{1} + \cdots + \dim V_{m}\]

\begin{remark}
以上不等式取等，当且仅当 $V_{1} + \cdots + V_{m}$ 是直和，这将在原书3.94中得到证明。
\end{remark}
\end{exercise}
\noindent 来源：\href{https://github.com/szdytom/LADRSolutions/blob/main/sections/2C.typ}{P-2C-17}（原书 2C 第 17 题；方而静题解）。
\begin{solve}
我们关于 $m$ 使用数学归纳法。

\begin{description}
\item[第 $1$ 步]
当 $m = 1$ 时，$V_{1}$ 根据题目条件自然是有限维的，且 $\dim V_{1} \le \dim V_{1}$。

\item[第 $k$ 步]
假设 $V_{1} + \cdots + V_{k - 1}$ 是有限维的，且 $\dim(V_{1} + \cdots + V_{k - 1}) \le \dim V_{1} + \cdots + \dim V_{k - 1}$。设 $u_{1} , \ldots , u_{n}$ 是其的一组基。同时，$V_{k}$ 是有限维的，设 $v_{1} , \ldots , v_{\ell}$ 是其的一组基。则
\[V_{1} + \cdots + V_{k} = \operatorname{span} (u_{1} , \ldots , u_{n} , v_{1} , \ldots , v_{\ell})\]
于是，$V_{1} + \cdots + V_{k}$ 是有限维的。根据子空间之和的维数（原书2.43），我们有
\[\begin{aligned}\dim(V_{1} + \cdots + V_{k}) = & \dim(V_{1} + \cdots + V_{k - 1}) + \dim V_{k} \\ & - \dim((V_{1} + \cdots + V_{k - 1}) \cap V_{k})\end{aligned}\]
根据归纳假设，
\[\begin{aligned}\dim(V_{1} + \cdots + V_{k}) & \le \dim(V_{1} + \cdots + V_{k - 1}) + \dim V_{k} \\ & \le \dim V_{1} + \cdots + \dim V_{k}\end{aligned}\]

综上所述，$V_{1} + \cdots + V_{m}$ 是有限维的，且 $\dim(V_{1} + \cdots + V_{m}) \le \dim V_{1} + \cdots + \dim V_{m}$。
\end{description}
\end{solve}

\begin{exercise}\label{prob:P-2C-18}
设 $V$ 是有限维向量空间，$\dim V = n \ge 1$。证明：存在 $V$ 的一维子空间 $V_{1} , \ldots , V_{n}$，使得

\[V = V_{1} \oplus \cdots \oplus V_{n}\]
\end{exercise}
\noindent 来源：\href{https://github.com/szdytom/LADRSolutions/blob/main/sections/2C.typ}{P-2C-18}（原书 2C 第 18 题；方而静题解）。
\begin{solve}
设 $v_{1} , \ldots , v_{n}$ 是 $V$ 的一组基。我们可以取 $V_{k} = \operatorname{span} (v_{k})$，其中 $k \in \{1 , \ldots , n\}$。注意到 $\dim V_{k} = 1$，因此 $V_{1} , \ldots , V_{n}$ 都是 $V$ 的一维子空间。根据基的判定准则（原书2.28）和直和的定义（原书1.41），我们立即得到 $V = V_{1} \oplus \cdots \oplus V_{n}$。
\end{solve}

\begin{exercise}\label{prob:P-2C-19}
受三个有限集的并集的元素数量公式的启发，你可能会这样猜测：若 $V_{1} , V_{2} , V_{3}$ 是一个有限维向量空间的子空间，那么有

\[\begin{aligned}\dim(V_{1} + V_{2} + V_{3}) = & \dim V_{1} + \dim V_{2} + \dim V_{3} \\ & - \dim(V_{1} \cap V_{2}) - \dim(V_{1} \cap V_{3}) - \dim(V_{2} \cap V_{3}) \\ & + \dim(V_{1} \cap V_{2} \cap V_{3})\end{aligned}\]

解释一下为什么这样猜测，然后证明以上公式，或给出反例。
\end{exercise}
\noindent 来源：\href{https://github.com/szdytom/LADRSolutions/blob/main/sections/2C.typ}{P-2C-19}（原书 2C 第 19 题；方而静题解）。
\begin{solve}
集合的容斥公式通过减去重复计数的交集、再补回三重交集，计算并集的元素数，这正是题中猜想的来源。但子空间的和包含跨子空间相加得到的向量，维数不能直接照搬集合的元素数。
在 $\mathbf R^2$ 中取
\[V_1=\operatorname{span}((1,0)),\quad V_2=\operatorname{span}((0,1)),\quad V_3=\operatorname{span}((1,1)).\]
三者维数都为 $1$，任意两个的交及三者的交都为 $\{0\}$，而 $V_1+V_2+V_3=\mathbf R^2$。因此猜想左端为 $2$、右端为 $3$，不成立。问题在于 $(V_1+V_2)\cap V_3=V_3$，但 $(V_1\cap V_3)+(V_2\cap V_3)=\{0\}$，交与和一般不满足这种分配律。
\end{solve}

\begin{exercise}\label{prob:P-2C-20}
已知 $V_{1}$，$V_{2}$ 和 $V_{3}$ 都是一个有限维向量空间的子空间，证明：

\[\begin{aligned}\dim(V_{1} + V_{2} + V_{3}) = & \dim V_{1} + \dim V_{2} + \dim V_{3} \\ & - \frac{\dim(V_{1} \cap V_{2})}{3} - \frac{\dim(V_{1} \cap V_{3})}{3} - \frac{\dim(V_{2} \cap V_{3})}{3} \\ & - \frac{\dim((V_{1} + V_{2}) \cap V_{3})}{3} \\ & - \frac{\dim((V_{1} + V_{3}) \cap V_{2})}{3} \\ & - \frac{\dim((V_{2} + V_{3}) \cap V_{1})}{3}\end{aligned}\]

\begin{remark}
上面公式初看可能有些奇怪，因为等式右边是否是整数令人怀疑。
\end{remark}
\end{exercise}
\noindent 来源：\href{https://github.com/szdytom/LADRSolutions/blob/main/sections/2C.typ}{P-2C-20}（原书 2C 第 20 题；方而静题解）。
\begin{solve}
等式两边同时乘以 $3$，并利用子空间之和的维数（原书2.43）将形如 $\dim(A \cap B)$ 的项代换，我们得到\footnote{$\mathrm{R.H.S.}$ 是 Right Hand Side 的缩写，表示等式右边的表达式。对称地，$\mathrm{L.H.S.}$ 是 Left Hand Side 的缩写，表示等式左边的表达式。}

\[\begin{aligned}3 \mathrm{R.H.S.} = & 3 \dim V_{1} + 3 \dim V_{2} + 3 \dim V_{3} \\ & - (\dim V_{1} + \dim V_{2} - \dim(V_{1} + V_{2})) \\ & - (\dim V_{1} + \dim V_{3} - \dim(V_{1} + V_{3})) \\ & - (\dim V_{2} + \dim V_{3} - \dim(V_{2} + V_{3})) \\ & - (\dim(V_{1} + V_{2}) + \dim V_{3} - \dim(V_{1} + V_{2} + V_{3})) \\ & - (\dim(V_{1} + V_{3}) + \dim V_{2} - \dim(V_{1} + V_{3} + V_{2})) \\ & - (\dim(V_{2} + V_{3}) + \dim V_{1} - \dim(V_{2} + V_{3} + V_{1}))\end{aligned}\]

整理后，立即给出了等式两边相等的结果。
\end{solve}

\subsection{线性映射构成的空间}\label{app:section-3A}
\begin{exercise}\label{prob:P-3A-1}
设 $b , c \in \mathbf R$。定义

\[\begin{aligned}T : & \mathbf R^{3} \to \mathbf R^{2} \\ & (x , y , z) \mapsto (2 x - 4 y + 3 z + b , 6 x + c x y z)\end{aligned}\]

证明：$T$ 是线性映射，当且仅当，$b = c = 0$。
\end{exercise}
\noindent 来源：\href{https://github.com/szdytom/LADRSolutions/blob/main/sections/3A.typ}{P-3A-1}（原书 3A 第 1 题；方而静题解）。
\begin{solve}
首先，假设 $b = c = 0$。则 $T (x , y , z) = (2 x - 4 y + 3 z , 6 x)$，我们逐条验证线性映射的定义（原书3.1）中的要求：

\begin{description}
\item[可加性] 对任意 $u , v \in \mathbf R^{3}$，$T (u + v) = T u + T v$。
证明：设 $u = (u_{1} , u_{2} , u_{3})$，$v = (v_{1} , v_{2} , v_{3})$，则
\[\begin{aligned}T (u + v) & = T (u_{1} + v_{1} , u_{2} + v_{2} , u_{3} + v_{3}) \\ & = (2 (u_{1} + v_{1}) - 4 (u_{2} + v_{2}) + 3 (u_{3} + v_{3}) , 6 (u_{1} + v_{1})) \\ & = (2 u_{1} - 4 u_{2} + 3 u_{3} , 6 u_{1}) + (2 v_{1} - 4 v_{2} + 3 v_{3} , 6 v_{1}) \\ & = T u + T v\end{aligned}\]

\item[齐次性] 对任意 $u \in \mathbf R^{3}$ 和任意 $\lambda \in \mathbf R$，$T (\lambda u) = \lambda T u$。
证明：设 $u = (u_{1} , u_{2} , u_{3})$，则
\[\begin{aligned}T (\lambda u) & = T (\lambda u_{1} , \lambda u_{2} , \lambda u_{3}) \\ & = (2 (\lambda u_{1}) - 4 (\lambda u_{2}) + 3 (\lambda u_{3}) , 6 (\lambda u_{1})) \\ & = \lambda (2 u_{1} - 4 u_{2} + 3 u_{3} , 6 u_{1}) \\ & = \lambda T u\end{aligned}\]

综上，$T$ 满足线性映射的定义。

另一方面，假设 $T$ 是线性映射。则根据线性映射将 $0$ 映射到 $0$（原书3.10），有

\end{description}
\[\begin{aligned}T (0 , 0 , 0) & = (2 \cdot 0 - 4 \cdot 0 + 3 \cdot 0 + b , 6 \cdot 0 + c \cdot 0) \\ & = (b , 0) = (0 , 0)\end{aligned}\]

因此 $b = 0$。另一方面，设 $u = (1 , 1 , 1)$，则根据齐次性的要求，有

\[T (2 u) = (2 , 12 + 8 c) = (2 , 12 + 2 c) = 2 T (u)\]

于是，$c = 0$。

综上所述，$T$ 是线性映射，当且仅当 $b = c = 0$。
\end{solve}

\begin{exercise}\label{prob:P-3A-2}
设 $b , c \in \mathbf R$。定义

\[\begin{aligned}T : & \mathcal P (\mathbf R) \to \mathbf R^{2} \\ & p \mapsto \begin{pmatrix}3 p (4) + 5 p' (6) + b p (1) p (2)\\\int_{- 1}^{2} x^{3} p (x) \,\mathrm d x + c \sin p (0)\end{pmatrix}\end{aligned}\]

证明：$T$ 是线性映射，当且仅当，$b = c = 0$。
\end{exercise}
\noindent 来源：\href{https://github.com/szdytom/LADRSolutions/blob/main/sections/3A.typ}{P-3A-2}（原书 3A 第 2 题；方而静题解）。
\begin{solve}
首先，假设 $b = c = 0$。则 $T (p) = (3 p (4) + 5 p' (6) , \int_{- 1}^{2} x^{3} p (x) \,\mathrm d x)$，我们逐条验证线性映射的定义（原书3.1）中的要求：

\begin{description}
\item[可加性] 对任意 $p , q \in \mathcal P (\mathbf R)$，$T (p + q) = T p + T q$。
证明：设 $p , q \in \mathcal P (\mathbf R)$，则
\[\begin{aligned}T (p + q) & = (3 (p + q) (4) + 5 (p + q)' (6) , \int_{- 1}^{2} x^{3} (p + q) (x) \,\mathrm d x) \\ & = (3 p (4) + 3 q (4) + 5 p' (6) + 5 q' (6) , \int_{- 1}^{2} x^{3} p (x) \,\mathrm d x + \int_{- 1}^{2} x^{3} q (x) \,\mathrm d x) \\ & = (3 p (4) + 5 p' (6) , \int_{- 1}^{2} x^{3} p (x) \,\mathrm d x) + (3 q (4) + 5 q' (6) , \int_{- 1}^{2} x^{3} q (x) \,\mathrm d x) \\ & = T p + T q\end{aligned}\]

\item[齐次性] 对任意 $p \in \mathcal P (\mathbf R)$ 和任意 $\lambda \in \mathbf R$，$T (\lambda p) = \lambda T p$。
证明：设 $p \in \mathcal P (\mathbf R)$，则
\[\begin{aligned}T (\lambda p) & = (3 (\lambda p) (4) + 5 (\lambda p)' (6) , \int_{- 1}^{2} x^{3} (\lambda p) (x) \,\mathrm d x) \\ & = (\lambda (3 p (4)) + \lambda (5 p' (6)) , \lambda \int_{- 1}^{2} x^{3} p (x) \,\mathrm d x) \\ & = \lambda (3 p (4) , \int_{- 1}^{2} x^{3} p (x) \,\mathrm d x) \\ & = \lambda T p\end{aligned}\]

综上，$T$ 满足线性映射的定义。

另一方面，假设 $T$ 是线性映射。设 $p : x \mapsto x + 1$。则根据齐次性的要求，有

\end{description}
\[T (2 p) = (40 + 24 b , \frac{207}{10} + c \sin 2) = (40 + 12 b , \frac{207}{10} + 2 c \sin 1) = 2 T (p)\]

解得 $b = c = 0$。这就是说 $T$ 是线性映射，当且仅当 $b = c = 0$。
\end{solve}

\begin{exercise}\label{prob:P-3A-3}
设 $T \in \mathcal L (\mathbf F^{n} , \mathbf F^{m})$。证明：对于 $j \in \{1 , \ldots , m\}$ 和 $k \in \{1 , \ldots , n\}$，存在标量 $A_{j , k} \in \mathbf F$，使得对于任意 $(x_{1} , \ldots , x_{n}) \in \mathbf F^{n}$，

\[T (x_{1} , \ldots , x_{n}) = (A_{1 , 1} x_{1} + \cdots + A_{1 , n} x_{n} , \ldots , A_{m , 1} x_{1} + \cdots + A_{m , n} x_{n})\]

\begin{remark}
此习题表明，线性映射 $T$ 具有在原书例3.3的倒数第二个例子中预示的形式。
\end{remark}
\end{exercise}
\noindent 来源：\href{https://github.com/szdytom/LADRSolutions/blob/main/sections/3A.typ}{P-3A-3}（原书 3A 第 3 题；方而静题解）。
\begin{solve}
令 $e_1,\ldots,e_n$、$f_1,\ldots,f_m$ 分别为定义域与陪域的标准基。对每个 $k$，唯一地写出
\[Te_k=\sum_{j=1}^m A_{jk}f_j.\]
于是由线性性，
\[T\Bigl(\sum_{k=1}^n x_ke_k\Bigr)=\sum_{j=1}^m\Bigl(\sum_{k=1}^n A_{jk}x_k\Bigr)f_j,\]
逐坐标即得题中公式；这些系数也由各个 $Te_k$ 唯一确定。
\end{solve}

\begin{exercise}\label{prob:P-3A-4}
设 $T \in \mathcal L (V , W)$ 且 $v_{1} , \ldots , v_{m}$ 是 $V$ 中的一组向量，使得向量组 $T v_{1} , \ldots , T v_{m}$ 在 $W$ 中是线性无关的。证明：向量组 $v_{1} , \ldots , v_{m}$ 是线性无关的。
\end{exercise}
\noindent 来源：\href{https://github.com/szdytom/LADRSolutions/blob/main/sections/3A.typ}{P-3A-4}（原书 3A 第 4 题；方而静题解）。
\begin{solve}
我们证明其逆否命题，即若 $v_{1} , \ldots , v_{m}$ 是线性相关的，则 $T v_{1} , \ldots , T v_{m}$ 也是线性相关的。现在假设 $v_{1} , \ldots , v_{m}$ 是线性相关的。

根据线性相关的定义（原书2.17），存在 $a_{1} , \ldots , a_{m} \in \mathbf F$，使得

\[a_{1} v_{1} + \cdots + a_{m} v_{m} = 0\]

其中 $a_{1} , \ldots , a_{m}$ 不全为 $0$。由于 $T$ 是线性映射，根据线性映射将 $0$ 映射到 $0$（原书3.10），我们有

\[\begin{aligned}0 = T 0 & = T (a_{1} v_{1} + \cdots + a_{m} v_{m}) \\ & = a_{1} T v_{1} + \cdots + a_{m} T v_{m} \\\end{aligned}\]

这立即说明 $T v_{1} , \ldots , T v_{m}$ 是线性相关的。

一个命题成立，当且仅当其逆否命题成立。因此，原命题得证。
\end{solve}

\begin{exercise}\label{prob:P-3A-5}
证明：$\mathcal L (V , W)$ 是向量空间，即原书 3.6 的结论。
\end{exercise}
\noindent 来源：\href{https://github.com/szdytom/LADRSolutions/blob/main/sections/3A.typ}{P-3A-5}（原书 3A 第 5 题；方而静题解）。
\begin{solve}
我们逐条验证 $\mathcal L (V , W)$ 满足向量空间的定义（原书1.20）中的要求。

\begin{description}
\item[可交换性] 对任意 $T , S \in \mathcal L (V , W)$，$T + S = S + T$。
证明：设 $v \in V$，则
\[\begin{aligned}(T + S) v & = T v + S v \\ & = S v + T v \\ & = (S + T) v\end{aligned}\]
因此 $T + S = S + T$。

\item[可结合性] 对任意 $T , S , R \in \mathcal L (V , W)$ 以及 $a , b \in \mathbf F$，都有 $(T + S) + R = T + (S + R)$，以及 $(a b) T = a (b T)$。
证明：设 $v \in V$，则对于加法的结合性，有
\[\begin{aligned}((T + S) + R) v & = (T + S) v + R v \\ & = T v + S v + R v \\ & = T v + (S + R) v \\ & = T v + S v + R v \\ & = (T + (S + R)) v\end{aligned}\]
因此 $(T + S) + R = T + (S + R)$。另一方面，对于标量乘法的结合性，有
\[\begin{aligned}((a b) T) v & = (a b) (T v) \\ & = a (b (T v)) \\ & = a ((b T) v) = (a (b T)) v\end{aligned}\]
因此 $(a b) T = a (b T)$。

\item[加法单位元] 存在 $0 \in \mathcal L (V , W)$，使得对任意 $T \in \mathcal L (V , W)$，$T + 0 = T$。
证明：取 $0 : v \mapsto 0$，设 $v \in V$，则
\[\begin{aligned}(T + 0) v & = T v + 0 v \\ & = T v + 0 \\ & = T v\end{aligned}\]
因此 $T + 0 = T$。

\item[加法逆元] 对任意 $T \in \mathcal L (V , W)$，存在 $- T \in \mathcal L (V , W)$，使得 $T + (- T) = 0$。
证明：取 $- T : v \mapsto - T v$，设 $v \in V$，则
\[\begin{aligned}(T + (- T)) v & = T v + (- T) v = T v - T v = 0 = 0 v\end{aligned}\]
因此 $T + (- T) = 0$。

\item[乘法单位元] 对于任意 $T \in \mathcal L (V , W)$，$1 T = T$。
证明：设 $v \in V$，则
\[(1 T) v = 1 (T v) = T v\]
因此 $1 T = T$。

\item[分配性质] 对于任意 $T , S \in \mathcal L (V , W)$ 和 $a , b \in \mathbf F$，都有 $a (T + S) = a T + a S$，以及 $(a + b) T = a T + b T$。
证明：设 $v \in V$，则对于第一个分配性质，有
\[\begin{aligned}(a (T + S)) v & = a ((T + S) v) \\ & = a (T v + S v) \\ & = a T v + a S v \\ & = (a T + a S) v\end{aligned}\]
因此 $a (T + S) = a T + a S$。另一方面，对于第二个分配性质，有
\[\begin{aligned}((a + b) T) v & = (a + b) (T v) \\ & = a (T v) + b (T v) \\ & = (a T) v + (b T) v \\ & = (a T + b T) v\end{aligned}\]
因此 $(a + b) T = a T + b T$。
\end{description}
\end{solve}

\begin{exercise}\label{prob:P-3A-6}
证明：线性映射的乘法具有可结合性、单位元和分配性质，即原书 3.8 的结论。
\end{exercise}
\noindent 来源：\href{https://github.com/szdytom/LADRSolutions/blob/main/sections/3A.typ}{P-3A-6}（原书 3A 第 6 题；方而静题解）。
\begin{solve}
我们逐条验证线性映射满足这些性质。

\begin{description}
\item[可结合性] 对于任意使乘积有意义的线性映射 $T_{1}$，$T_{2}$ 和 $T_{3}$（即 $T_{3}$ 映射到 $T_{2}$ 的定义空间中，$T_{2}$ 映射到 $T_{1}$ 的定义空间中），都有 $(T_{1} T_{2}) T_{3} = T_{1} (T_{2} T_{3})$。
证明：设 $v$ 是 $T_{3}$ 的定义空间中的任意向量，则
\[((T_{1} T_{2}) T_{3}) v = T_{1} (T_{2} (T_{3} v)) = (T_{1} (T_{2} T_{3})) v\]
因此 $(T_{1} T_{2}) T_{3} = T_{1} (T_{2} T_{3})$。

\item[单位元] 对于任意 $T \in \mathcal L (V , W)$，都有 $T I = I T = T$。这里的第一个 $I$ 是 $V$ 上的恒等变换，而第二个 $I$ 是 $W$ 上的恒等变换。
证明：设 $v \in V$，对于 $I T = T$，有
\[(I T) v = I (T v) = T v\]
因此 $I T = T$。另一方面，对于 $T I = T$，
\[(T I) v = T (I v) = T v\]
因此 $T I = T$。

\item[分配性质] 对于任意 $T , T_{1} , T_{2} \in \mathcal L (U , V)$ 和 $S , S_{1} , S_{2} \in \mathcal L (V , W)$，有 $(S_{1} + S_{2}) T = S_{1} T + S_{2} T$ 且 $S (T_{1} + T_{2}) = S T_{1} + S T_{2}$。
证明：设 $v \in U$，则对于第一个分配性质，有
\[\begin{aligned}((S_{1} + S_{2}) T) v & = (S_{1} + S_{2}) (T v) \\ & = S_{1} (T v) + S_{2} (T v) \\ & = (S_{1} T + S_{2} T) v\end{aligned}\]
因此 $(S_{1} + S_{2}) T = S_{1} T + S_{2} T$。另一方面，对于第二个分配性质，有
\[\begin{aligned}(S (T_{1} + T_{2})) v & = S ((T_{1} + T_{2}) v) \\ & = S (T_{1} v + T_{2} v) \\ & = S (T_{1} v) + S (T_{2} v) \\ & = (S T_{1} + S T_{2}) v\end{aligned}\]
因此 $S (T_{1} + T_{2}) = S T_{1} + S T_{2}$。

综上所述，线性映射的乘法具有可结合性、单位元和分配性质。
\end{description}
\end{solve}

\begin{exercise}\label{prob:P-3A-7}\label{app:E-linear-map-space-of-dim-1-is-dim-1}
证明：任何从一维向量空间到其自身的线性映射，就是标量乘法。形式化地说，即若 $\dim V = 1$ 且 $T \in \mathcal L (V)$，则存在 $\lambda \in \mathbf F$，使得 $T v = \lambda v$ 对任意 $v \in V$ 成立。
\end{exercise}
\noindent 来源：\href{https://github.com/szdytom/LADRSolutions/blob/main/sections/3A.typ}{P-3A-7}（原书 3A 第 7 题；方而静题解）。
\begin{solve}
设 $w$ 是 $V$ 的一组基。由于 $T w \in V$，根据基的性质，存在 $\lambda \in \mathbf F$，使得 $T w = \lambda w$。现在考虑任意 $v \in V$。根据基的性质，存在唯一的 $a \in \mathbf F$，使得 $v = a w$。因此

\[T v = T (a w) = a T w = a (\lambda w) = \lambda (a w) = \lambda v\]

综上所述，$T v = \lambda v$ 对任意 $v \in V$ 成立。
\end{solve}

\begin{exercise}\label{prob:P-3A-8}
给出一例：函数 $\phi : \mathbf R^{2} \to \mathbf R$，使得对于任意 $a \in \mathbf R$ 和 $v \in \mathbf R^{2}$，有

\[\phi (a v) = a \phi (v)\]

但是 $\phi$ 不是线性映射。

\begin{remark}
本题和下一题表明，仅有齐次性或仅有可加性，都不足以推导出一个函数是线性映射。
\end{remark}
\end{exercise}
\noindent 来源：\href{https://github.com/szdytom/LADRSolutions/blob/main/sections/3A.typ}{P-3A-8}（原书 3A 第 8 题；方而静题解）。
\begin{solve}
定义
\[\phi(x,y)=\begin{cases}x,&|x|\ge|y|,\\y,&|x|<|y|.\end{cases}\]
当 $a=0$ 时，$\phi(ax,ay)=0=a\phi(x,y)$。当 $a\ne0$ 时，乘以 $|a|$ 保持 $|x|$ 与 $|y|$ 的大小关系，故两边取同一分支，仍有 $\phi(ax,ay)=a\phi(x,y)$。但
\[\phi(1,0)+\phi(0,-1)=0,\qquad \phi(1,-1)=1,\]
所以 $\phi$ 不可加，因而不是线性映射。
\end{solve}

\begin{exercise}\label{prob:P-3A-9}
给出一例：函数 $\phi : \mathbf C \to \mathbf C$，使得对于任意 $w , z \in \mathbf C$，有

\[\phi (w + z) = \phi (w) + \phi (z)\]

但是 $\phi$ 不是线性映射。（此处将 $\mathbf C$ 视为复向量空间。）

\begin{remark}
满足上述可加性条件，而不是线性映射的函数 $\phi : \mathbf R \to \mathbf R$ 也是存在的。然而，证明这样的一个函数存在，要用到高得多的数学工具。\footnote{这有很强的分析背景，感兴趣的读者可以搜索“柯西方程的不连续解”“哈默尔基“等关键词。明确构造出这样的函数一般而言需要承认选择公理，构造的基本思路是将 $\mathbf R$ 视为 $\mathbf Q$ 上的一个无限维向量空间。}
\end{remark}
\end{exercise}
\noindent 来源：\href{https://github.com/szdytom/LADRSolutions/blob/main/sections/3A.typ}{P-3A-9}（原书 3A 第 9 题；方而静题解）。
\begin{solve}
设 $z \in \mathbf C$，令\footnote{这里的 $\operatorname{Re} z$ 表示 $z$ 的实部，类似地，记号 $\operatorname{Im} z$ 表示 $z$ 的虚部。}

\[\phi (z) = \operatorname{Re} z\]

设 $w = (a + b \mathrm i) , z = (c + d \mathrm i) \in \mathbf C$，其中 $a , b , c , d \in \mathbf R$。则

\[\phi (w + z) = \operatorname{Re} ((a + c) + (b + d) \mathrm i) = a + c = \operatorname{Re} w + \operatorname{Re} z = \phi (w) + \phi (z)\]

因此，对于任意 $w , z \in \mathbf C$，都有 $\phi (w + z) = \phi (w) + \phi (z)$。

另一方面，注意到

\[\phi (\mathrm i 2) = 0 \ne 2 \mathrm i = \mathrm i \phi (2)\]

这说明 $\phi$ 不满足线性映射的定义（原书3.1）中对齐次性的要求。因此 $\phi$ 不是线性映射。
\end{solve}

\begin{exercise}\label{prob:P-3A-10}
证明或证伪：如果 $q \in \mathcal P (\mathbf R)$，$T : \mathcal P (\mathbf R) \to \mathcal P (\mathbf R)$ 定义为 $T p = q \circ p$，那么 $T$ 是线性映射。

\begin{remark}
这里定义的函数 $T$，不同于原书3.3中最后一个例子定义的函数 $T$，区别在于复合的次序。
\end{remark}
\end{exercise}
\noindent 来源：\href{https://github.com/szdytom/LADRSolutions/blob/main/sections/3A.typ}{P-3A-10}（原书 3A 第 10 题；方而静题解）。
\begin{solve}
设 $x \in \mathbf R$。令 $q \in \mathcal P (\mathbf R)$ 为 $x \mapsto x^{2}$。注意到，取 $p \in \mathcal P (\mathbf R)$ 为 $x \mapsto x$，则

\[(T 2 p) (x) = q (2 p (x)) = 4 x^{2} \ne 2 x^{2} = 2 q (p (x)) = 2 T p\]

这说明 $T$ 不满足线性映射的定义（原书3.1）中对齐次性的要求。因此 $T$ 不是线性映射。
\end{solve}

\begin{exercise}\label{prob:P-3A-11}
设 $V$ 是有限维向量空间，$T \in \mathcal L (V)$。证明：$T$ 是恒等算子的标量倍，当且仅当，对于任意 $S \in \mathcal L (V)$，$S T = T S$ 成立。
\end{exercise}
\noindent 来源：\href{https://github.com/szdytom/LADRSolutions/blob/main/sections/3A.typ}{P-3A-11}（原书 3A 第 11 题；方而静题解）。
\begin{solve}
若 $T=\lambda I$，则对任意 $S$ 都有 $ST=TS$。反过来，设 $T$ 与所有线性变换交换。若 $V=\{0\}$，结论成立；以下设基为 $e_1,\ldots,e_n$。令 $P_j$ 是保留第 $j$ 个坐标、其余坐标置零的投影。由 $P_jT=TP_j$ 得 $Te_j=P_jTe_j\in\operatorname{span}(e_j)$，故 $Te_j=\lambda_je_j$。
对任意 $i,j$，令 $E_{ij}e_j=e_i$、$E_{ij}e_k=0$（$k\ne j$）。将 $TE_{ij}=E_{ij}T$ 作用于 $e_j$，得到 $\lambda_ie_i=\lambda_je_i$，所以所有 $\lambda_j$ 相等。记公共值为 $\lambda$，则 $T=\lambda I$。
\end{solve}

\begin{exercise}\label{prob:P-3A-12}
设 $U$ 是 $V$ 的子空间，$U \ne V$。设 $S \in \mathcal L (U , W)$，且 $S \ne 0$（即存在 $u \in U$，使得 $S u \ne 0$）。定义 $T : V \to W$，使得

\[T v = \begin{cases}S v \quad & v \in U\\0 \quad & v \in V \quad \text{且} \quad v \notin U\end{cases}\]

证明：$T$ 不是 $V$ 上的线性映射。
\end{exercise}
\noindent 来源：\href{https://github.com/szdytom/LADRSolutions/blob/main/sections/3A.typ}{P-3A-12}（原书 3A 第 12 题；方而静题解）。
\begin{solve}
设 $u \in U$，使得 $S u \ne 0$，$v \in V$ 且 $v \notin V$。由于 $u + (v - u) = v \notin U$，根据子空间的条件（原书1.34）中对加法封闭性的要求，只能有 $v - u \notin U$。注意到

\[T u = S u \ne 0 = 0 + 0 = T u + T (v - u)\]

这说明 $T$ 不满足线性映射的定义（原书3.1）中对可加性的要求。因此 $T$ 不是 $V$ 上的线性映射。
\end{solve}

\begin{exercise}\label{prob:P-3A-13}\label{app:E-extend-linear-map}
设 $V$ 是有限维的向量空间。证明：$V$ 的子空间上的任意一个线性映射都可以扩充为 $V$ 上的线性映射。形式化地说，设 $U$ 是 $V$ 的子空间，$S \in \mathcal L (U , W)$。则存在 $T \in \mathcal L (V , W)$，使得 $T u = S u$ 对任意 $u \in U$ 成立。

\begin{remark}
原书 3.125 的证明将会用到本题的结果。
\end{remark}
\end{exercise}
\noindent 来源：\href{https://github.com/szdytom/LADRSolutions/blob/main/sections/3A.typ}{P-3A-13}（原书 3A 第 13 题；方而静题解）。
\begin{solve}
设 $v_{1} , \ldots , v_{m}$ 是 $U$ 的一组基，根据每一个线性无关组都可被扩充为基（原书2.32），我们可以找到 $v_{m + 1} , \ldots , v_{n} \in V$，使得 $v_{1} , \ldots , v_{n}$ 是 $V$ 的一组基。

现在，根据线性映射引理（原书3.4），我们可以找到线性映射 $T : V \to W$，使得对于任意 $i \in \{1 , \ldots , m\}$，都有 $T v_{i} = S v_{i}$，且对于 $j \in \{m + 1 , \ldots , n\}$，$T v_{j} = 0$。对于任意 $u \in U$，由于 $u$ 可以唯一地表示为

\[u = a_{1} v_{1} + \cdots + a_{m} v_{m}\]

因此

\[T u = a_{1} T v_{1} + \cdots + a_{m} T v_{m} = a_{1} S v_{1} + \cdots + a_{m} S v_{m} = S u\]

综上所述，$T$ 满足 $T u = S u$ 对任意 $u \in U$ 成立，即 $T$ 是 $S$ 的扩充。
\end{solve}

\begin{exercise}\label{prob:P-3A-14}
设 $V$ 是有限维向量空间，$\dim V > 0$，$W$ 是无限维向量空间。证明：$\mathcal L (V , W)$ 是无限维的。
\end{exercise}
\noindent 来源：\href{https://github.com/szdytom/LADRSolutions/blob/main/sections/3A.typ}{P-3A-14}（原书 3A 第 14 题；方而静题解）。
\begin{solve}
设 $v \in V$，$v \ne 0$。根据习题~\ref{app:E-inf-dim-space-seq-characterization}，$W$ 中存在一个序列 $w_{1} , w_{2} , \ldots$ 使得对于任意 $m \in \mathbf N^{+}$，向量组 $w_{1} , \ldots , w_{m}$ 线性无关。

根据线性映射引理（原书3.4），我们可以找到线性映射 $S_{k} : \operatorname{span} (v) \to W$，使得 $S_{k} v = w_{k}$。然后根据习题~\ref{app:E-extend-linear-map}，存在 $T_{k} \in \mathcal L (V , W)$，使得对于任意 $u \in \operatorname{span} (v)$，都有 $T_{k} u = S_{k} u$。

对于任意 $m \in \mathbf N^{+}$，设 $a_{1} , \ldots a_{m} \in \mathbf F$，使得

\[a_{1} T_{1} + \cdots + a_{m} T_{m} = 0\]

等式两边同时代入 $v$，我们得到

\[a_{1} w_{1} + \cdots + a_{m} w_{m} = 0\]

由于 $w_{1} , \ldots , w_{m}$ 是线性无关的，因此 $a_{1} = \cdots = a_{m} = 0$。这说明 $T_{1} , \ldots , T_{m}$ 是线性无关的。

所以，根据习题~\ref{app:E-inf-dim-space-seq-characterization}，$\mathcal L (V , W)$ 是无限维的。
\end{solve}

\begin{exercise}\label{prob:P-3A-15}
设 $v_{1} , \ldots , v_{m}$ 是 $V$ 中的线性相关向量组，$\dim W > 0$。证明：存在 $w_{1} , \ldots , w_{m} \in W$，使得不存在 $T \in \mathcal L (V , W)$ 对于任意 $k \in \{1 , \ldots , m\}$，都有 $T v_{k} = w_{k}$。
\end{exercise}
\noindent 来源：\href{https://github.com/szdytom/LADRSolutions/blob/main/sections/3A.typ}{P-3A-15}（原书 3A 第 15 题；方而静题解）。
\begin{solve}
由于向量组 $v_{1} , \ldots , v_{m}$ 是线性相关的，根据线性相关性引理（原书2.19），存在 $k \in \{1 , \ldots , m\}$，使得 $v_{k} \in \operatorname{span} (v_{1} , \ldots , v_{k - 1})$。设

\[v_{k} = a_{1} v_{1} + \cdots + a_{k - 1} v_{k - 1}\]

其中 $a_{1} , \ldots , a_{k - 1} \in \mathbf F$。任取 $w_{k} \ne 0$，并令 $w_{1} = \cdots = w_{k - 1} = 0$。于是

\[T v_{k} = a_{1} T v_{1} + \cdots + a_{k - 1} T v_{k - 1} = 0 \ne w_{k}\]

这说明不存在 $T \in \mathcal L (V , W)$，使得对于任意 $k \in \{1 , \ldots , m\}$，都有 $T v_{k} = w_{k}$。
\end{solve}

\begin{exercise}\label{prob:P-3A-16}
设 $V$ 是有限维向量空间，$\dim V > 1$。证明：存在 $S , T \in \mathcal L (V)$ 使得 $S T \ne T S$。
\end{exercise}
\noindent 来源：\href{https://github.com/szdytom/LADRSolutions/blob/main/sections/3A.typ}{P-3A-16}（原书 3A 第 16 题；方而静题解）。
\begin{solve}
设 $v_{1} , v_{2} \in V$ 是线性无关的向量组，令 $U = \operatorname{span} (v_{1} , v_{2})$。设 $u \in U$，则 $u$ 可以唯一地表示为

\[u = a_{1} v_{1} + a_{2} v_{2}\]

其中 $a_{1} , a_{2} \in \mathbf F$。根据线性映射引理（原书3.4），我们可以找到线性映射 $S|_U : U \to W$ 和 $T|_U : U \to W$\footnote{这里使用的是限制算子的记号，表示定义域限制在 $U$ 上，原书第三版的5.14定义了此记号。但是这里无需明白这些定义，将其分别当作两个函数的名字即可。}，使得 $S|_U (v_{1}) = v_{2}$ 且 $S|_U (v_{2}) = v_{1}$，以及 $T|_U (v_{1}) = T|_U (v_{2}) = v_{1}$。

现在，根据习题~\ref{app:E-extend-linear-map}，存在 $S , T \in \mathcal L (V)$，使得对于任意 $u \in U$，都有 $S u = S|_U (u)$ 且 $T u = T|_U (u)$。注意到

\[S T v_{2} = S v_{1} = v_{2} \ne v_{1} = T v_{1} = T S v_{2}\]

于是 $S T \ne T S$，这立即完成了证明。
\end{solve}

\begin{exercise}\label{prob:P-3A-17}\label{app:E-two-sided-ideal}
设 $V$ 是有限维向量空间。

\begin{itemize}
\item $\mathcal L (V)$ 的子空间 $\mathcal E$ 被称为\emph{双边理想（two-sided ideal）}，是指 $T E \in \mathcal E$ 和 $E T \in \mathcal E$，对于任意 $E \in \mathcal E$ 和 $T \in \mathcal L (V)$ 都成立。

\end{itemize}
证明：$\mathcal L (V)$ 仅有的双边理想是 $\{0\}$ 和 $\mathcal L (V)$。
\end{exercise}
\noindent 来源：\href{https://github.com/szdytom/LADRSolutions/blob/main/sections/3A.typ}{P-3A-17}（原书 3A 第 17 题；方而静题解）。
\begin{solve}
零子空间与整个 $\mathcal L(V)$ 满足双边理想的定义。若 $V=\{0\}$，它们相同，结论成立。设 $\mathcal E$ 是非零双边理想，取 $A\in\mathcal E$、$u\in V$ 使 $Au=w\ne0$。把 $w$ 扩充为基并定义线性泛函 $f$，使 $f(w)=1$。
固定 $V$ 的基 $e_1,\ldots,e_n$ 及其对偶基 $e^1,\ldots,e^n$。对任意 $i,j$，定义
\[S_i(x)=f(x)e_i,\qquad R_j(x)=e^j(x)u.\]
则 $S_iAR_j(x)=e^j(x)e_i$，这是矩阵单位对应的线性变换 $E_{ij}$。由于 $\mathcal E$ 是双边理想，$E_{ij}=S_iAR_j\in\mathcal E$。任意线性变换都能写成 $\sum_{i,j}a_{ij}E_{ij}$，由子空间对线性组合封闭可知 $\mathcal E=\mathcal L(V)$。
\end{solve}

\subsection{核与像}\label{app:section-3B}
\begin{exercise}\label{prob:P-3B-1}
给出一例：满足 $\dim \ker T = 3$ 且 $\dim \operatorname{im} T = 2$ 的线性映射 $T$。
\end{exercise}
\noindent 来源：\href{https://github.com/szdytom/LADRSolutions/blob/main/sections/3B.typ}{P-3B-1}（原书 3B 第 1 题；方而静题解）。
\begin{solve}
令

\[\begin{aligned}T : & \mathbf R^{5} \to \mathbf R^{2} \\ & (x_{1} , x_{2} , x_{3} , x_{4} , x_{5}) \mapsto (x_{1} , x_{2})\end{aligned}\]

根据定义

\[\begin{aligned}\operatorname{im} T & = \mathbf R^{2} \\ \ker T & = \{(0 , 0 , x , y , z) \in \mathbf R^{5} : x , y , z \in \mathbf R\}\end{aligned}\]

于是 $\dim \ker T = 3$ 且 $\dim \operatorname{im} T = 2$。
\end{solve}

\begin{exercise}\label{prob:P-3B-2}
设 $S , T \in \mathcal L (V)$ 使得 $\operatorname{im} S \subseteq \ker T$，证明：$(S T)^{2} = 0$。
\end{exercise}
\noindent 来源：\href{https://github.com/szdytom/LADRSolutions/blob/main/sections/3B.typ}{P-3B-2}（原书 3B 第 2 题；方而静题解）。
\begin{solve}
设 $v \in V$。考虑到 $S (T v) \in \operatorname{im} S \subseteq \ker T$，根据定义，$T S T v = 0$。根据线性映射将 $0$ 映射到 $0$（原书3.10），

\[(S T)^{2} v = S (T S T) v = S 0 = 0\]

因此 $(S T)^{2} = 0$。
\end{solve}

\begin{exercise}\label{prob:P-3B-3}
设向量组 $v_{1} , \ldots , v_{m} \in V$，定义 $T \in \mathcal L (\mathbf F^{m} , V)$ 为

\[(z_{1} , \ldots , z_{m}) \mapsto z_{1} v_{1} + \cdots + z_{m} v_{m}\]

\begin{enumerate}[label=(\alph*)]
\item $T$ 的什么性质对应于向量组 $v_{1} , \ldots , v_{m}$ 张成 $V$？
\item $T$ 的什么性质对应于向量组 $v_{1} , \ldots , v_{m}$ 线性无关？
\end{enumerate}
\end{exercise}
\noindent 来源：\href{https://github.com/szdytom/LADRSolutions/blob/main/sections/3B.typ}{P-3B-3}（原书 3B 第 3 题；方而静题解）。
\begin{solve}
由定义，$\operatorname{im} T=\operatorname{span}(v_1,\ldots,v_m)$，所以 $T$ 满射当且仅当该向量组张成 $V$。另一方面，
\[\ker T=\{(z_1,\ldots,z_m):z_1v_1+\cdots+z_mv_m=0\}.\]
因此 $\ker T=\{0\}$ 当且仅当向量组线性无关；线性映射的单射性又等价于核为零，得到第二个结论。整个论证无需预先假定向量组线性无关，也无需假定 $V$ 有限维。
\end{solve}

\begin{exercise}\label{prob:P-3B-4}
证明：$\{T \in \mathcal L (\mathbf R^{5} , \mathbf R^{4}) : \dim \ker T > 2\}$ 不是 $\mathcal L (\mathbf R^{5} , \mathbf R^{4})$ 的子空间。
\end{exercise}
\noindent 来源：\href{https://github.com/szdytom/LADRSolutions/blob/main/sections/3B.typ}{P-3B-4}（原书 3B 第 4 题；方而静题解）。
\begin{solve}
记 $S = \{T \in \mathcal L (\mathbf R^{5} , \mathbf R^{4}) : \dim \ker T > 2\}$。取 $T_{1} , T_{2} \in \mathcal L (\mathbf R^{5} , \mathbf R^{4})$，使得对于任意 $x_{1} , x_{2} , x_{3} , x_{4} , x_{5} \in \mathbf R$，有

\[\begin{aligned}T_{1} (x_{1} , x_{2} , x_{3} , x_{4} , x_{5}) & = (x_{1} , x_{2} , 0 , 0) \\ T_{2} (x_{1} , x_{2} , x_{3} , x_{4} , x_{5}) & = (0 , 0 , x_{3} , x_{4})\end{aligned}\]

容易验证 $\dim \ker T_{1} = \dim \ker T_{2} = 3 > 2$，即 $T_{1}$ 和 $T_{2}$ 都是 $S$ 中的元素。然而，注意到 $\dim \operatorname{im} (T_{1} + T_{2}) = 4$，即根据线性映射基本定理（原书3.21），

\[\dim \ker (T_{1} + T_{2}) = \dim \mathbf R^{5} - \dim \operatorname{im} (T_{1} + T_{2}) = 5 - 4 = 1\]

因此 $T_{1} + T_{2} \notin S$。这说明 $S$ 违反了子空间的条件（原书1.34）中对加法封闭性的要求，故 $S$ 不是 $\mathcal L (\mathbf R^{5} , \mathbf R^{4})$ 的子空间。
\end{solve}

\begin{exercise}\label{prob:P-3B-5}
给出一例：线性映射 $T \in \mathcal L (\mathbf R^{4})$，满足 $\operatorname{im} T = \ker T$。
\end{exercise}
\noindent 来源：\href{https://github.com/szdytom/LADRSolutions/blob/main/sections/3B.typ}{P-3B-5}（原书 3B 第 5 题；方而静题解）。
\begin{solve}
设 $v_{1} , \ldots , v_{4}$ 是 $\mathbf R^{4}$ 的一组基。根据线性映射引理（原书3.4），存在 $T \in \mathcal L (\mathbf R^{4})$，使得 $T v_{1} = T v_{2} = 0$，$T v_{3} = v_{1}$，$T v_{4} = v_{2}$。因此

\[\operatorname{im} T = \ker T = \operatorname{span} (v_{1} , v_{2})\]

因此 $T$ 满足题目要求。
\end{solve}

\begin{exercise}\label{prob:P-3B-6}
证明：不存在 $T \in \mathcal L (\mathbf R^{5})$，使得 $\operatorname{im} T = \ker T$。
\end{exercise}
\noindent 来源：\href{https://github.com/szdytom/LADRSolutions/blob/main/sections/3B.typ}{P-3B-6}（原书 3B 第 6 题；方而静题解）。
\begin{solve}
假设存在 $T \in \mathcal L (\mathbf R^{5})$，使得 $\operatorname{im} T = \ker T$。根据线性映射基本定理（原书3.21），有

\[\dim \mathbf R^{5} = \dim \ker T + \dim \operatorname{im} T\]

由于 $\operatorname{im} T = \ker T$，因此 $\dim \operatorname{im} T = \dim \ker T$。设 $\dim \operatorname{im} T = n$，则

\[\dim \mathbf R^{5} = n + n = 2 n\]

由于 $\dim \mathbf R^{5} = 5$，因此 $2 n = 5$，这与 $n$ 是整数矛盾。因此不存在这样的 $T$。
\end{solve}

\begin{exercise}\label{prob:P-3B-7}
设 $V$ 和 $W$ 是有限维向量空间，$2 \le \dim V \le \dim W$。证明：$\{T \in \mathcal L (V , W) : T \text{不是单射}\}$ 不是 $\mathcal L (V , W)$ 的子空间。
\end{exercise}
\noindent 来源：\href{https://github.com/szdytom/LADRSolutions/blob/main/sections/3B.typ}{P-3B-7}（原书 3B 第 7 题；方而静题解）。
\begin{solve}
记 $S = \{T \in \mathcal L (V , W) : T \text{不是单射}\}$。设 $v_{1} , \ldots , v_{m}$ 是 $V$ 的一组基，$w_{1} , \ldots , w_{n}$ 是 $W$ 的一组基，其中 $n \ge m = \dim V \ge 2$。

根据线性映射引理（原书3.4），存在 $T \in \mathcal L (V , W)$，使得对于任意 $i \in \{1 , \ldots , m\}$，有 $T v_{i} = w_{i}$。设 $v \in V$。假设 $T v = 0$，将 $v$ 表示为 $v = z_{1} v_{1} + \cdots + z_{m} v_{m}$，其中 $z_{1} , \ldots , z_{m} \in \mathbf F$。则

\[T v = z_{1} T v_{1} + \cdots + z_{m} T v_{m} = z_{1} w_{1} + \cdots + z_{m} w_{m} = 0\]

由于 $w_{1} , \ldots , w_{m}$ 是线性无关的，因此 $z_{1} = \cdots = z_{m} = 0$，即 $v = 0$。因此 $\ker T = \{0\}$，根据“单射性 $\iff$ 零空间为 $\{0\}$”（原书3.15），$T$ 是单射，即 $T \notin S$。

再次利用线性映射引理（原书3.4），存在 $R \in \mathcal L (V , W)$，使得 $R v_{1} = w_{1}$，且对于任意 $i \in \{2 , \ldots , m\}$，$R v_{i} = 0$。由于 $m \ge 2$，我们至少有 $R v_{2} = R 0 = 0$，故 $R$ 不是单射。

注意到

\[(T - R) v_{1} = T v_{1} - R v_{1} = w_{1} - w_{1} = 0 = (T - R) 0\]

故 $T - R$ 不是单射。即 $R , T - R \in S$，注意到

\[R + (T - R) = T \notin S\]

这说明 $S$ 违反了子空间的条件（原书1.34）中对加法封闭性的要求，故 $S$ 不是 $\mathcal L (V , W)$ 的子空间。
\end{solve}

\begin{exercise}\label{prob:P-3B-8}
设 $V$ 和 $W$ 是有限维向量空间，$2 \le \dim W \le \dim V$。证明：$\{T \in \mathcal L (V , W) : T \text{不是满射}\}$ 不是 $\mathcal L (V , W)$ 的子空间。
\end{exercise}
\noindent 来源：\href{https://github.com/szdytom/LADRSolutions/blob/main/sections/3B.typ}{P-3B-8}（原书 3B 第 8 题；方而静题解）。
\begin{solve}
记 $S = \{T \in \mathcal L (V , W) : T \text{不是满射}\}$。设 $v_{1} , \ldots , v_{m}$ 是 $V$ 的一组基，$w_{1} , \ldots , w_{n}$ 是 $W$ 的一组基，其中 $m \ge n = \dim W \ge 2$。

根据线性映射引理（原书3.4），存在 $T \in \mathcal L (V , W)$，使得对于任意 $i \in \{1 , \ldots , n\}$，有 $T v_{i} = w_{i}$，且对于任意 $i \in \{n + 1 , \ldots , m\}$，$T v_{i} = 0$。设 $w \in W$。将 $w$ 表示为 $w = z_{1} w_{1} + \cdots + z_{n} w_{n}$，其中 $z_{1} , \ldots , z_{m} \in \mathbf F$。则

\[w = z_{1} T v_{1} + \cdots + z_{n} T v_{n} = T (z_{1} v_{1} + \cdots + z_{n} v_{n})\]

故 $\operatorname{im} T = W$，即 $T$ 是满射，因此 $T \notin S$。

再次利用线性映射引理（原书3.4），存在 $R \in \mathcal L (V , W)$，使得 $R v_{1} = w_{1}$，且对于任意 $i \in \{2 , \ldots , n\}$，$R v_{i} = 0$。由于 $n \ge 2$，$w_{1} , w_{2}$ 线性无关，根据习题~\ref{app:E-when-1-or-2-vectors-indep}，不存在 $\lambda \in \mathbf F$，使得 $w_{2} = \lambda w_{1}$，即 $w \notin \operatorname{im} R$，故 $R$ 不是满射。

反证假设 $T - R$ 是满射，则存在 $v \in V$ 使得 $(T - R) v = w_{1}$。将 $v$ 表示为 $v = a_{1} v_{1} + \cdots + a_{m} v_{m}$，其中 $a_{1} , \ldots , a_{m} \in \mathbf F$。则

\[\begin{aligned}w_{1} = (T - R) v & = T v - R v \\ & = a_{1} T v_{1} + \cdots + a_{m} T v_{m} + a_{1} R v_{1} + \cdots + a_{m} R v_{m} \\ & = a_{1} w_{1} + \cdots + a_{n} w_{n} - a_{1} w_{1} \\ & = a_{2} w_{2} + \cdots + a_{n} w_{n} \in \operatorname{span} (w_{2} , \ldots , w_{m})\end{aligned}\]

根据习题~\ref{app:E-when-vector-list-append-remains-indep}，$w_{2} , \ldots , w_{n} , w_{1}$ 不是线性无关的，矛盾。故 $T - R$ 不是满射。现在 $R , T - R \in S$，注意到

\[R + (T - R) = T \notin S\]

这说明 $S$ 违反了子空间的条件（原书1.34）中对加法封闭性的要求，故 $S$ 不是 $\mathcal L (V , W)$ 的子空间。
\end{solve}

\begin{exercise}\label{prob:P-3B-9}\label{app:E-indep-preservance-under-inj}
设 $T \in \mathcal L (V , W)$ 是单射，向量组 $v_{1} , \ldots , v_{n}$ 在 $V$ 中线性无关。证明：向量组 $T v_{1} , \ldots , T v_{n}$ 在 $W$ 中线性无关。
\end{exercise}
\noindent 来源：\href{https://github.com/szdytom/LADRSolutions/blob/main/sections/3B.typ}{P-3B-9}（原书 3B 第 9 题；方而静题解）。
\begin{solve}
设 $a_{1} , \ldots , a_{n} \in \mathbf F$ 使得

\[a_{1} T v_{1} + \cdots + a_{n} T v_{n} = 0\]

根据线性映射的定义，有

\[T (a_{1} v_{1} + \cdots + a_{n} v_{n}) = 0\]

由于 $T$ 是单射，根据“单射性 $\iff$ 零空间为 $\{0\}$”（原书3.15），$\ker T = \{0\}$，因此 $a_{1} v_{1} + \cdots + a_{n} v_{n} = 0$。由于向量组 $v_{1} , \ldots , v_{n}$ 在 $V$ 中线性无关，只能有 $a_{1} = \cdots = a_{n} = 0$。

这说明向量组 $T v_{1} , \ldots , T v_{n}$ 在 $W$ 中线性无关。
\end{solve}

\begin{exercise}\label{prob:P-3B-10}\label{app:E-domain-span-to-range-span}
设向量组 $v_{1} , \ldots , v_{n}$ 张成 $V$，$T \in \mathcal L (V , W)$，证明：$T v_{1} , \ldots , T v_{n}$ 张成 $\operatorname{im} T$。
\end{exercise}
\noindent 来源：\href{https://github.com/szdytom/LADRSolutions/blob/main/sections/3B.typ}{P-3B-10}（原书 3B 第 10 题；方而静题解）。
\begin{solve}
设 $w \in \operatorname{im} T$，则可设 $v \in V$，使得 $T v = w$。将 $v$ 表示为 $v = a_{1} v_{1} + \cdots + a_{n} v_{n}$，其中 $a_{1} , \ldots , a_{n} \in \mathbf F$。则

\[w = T v = T (a_{1} v_{1} + \cdots + a_{n} v_{n}) = a_{1} T a_{1} + \cdots + a_{n} T a_{n}\]

这说明 $w$ 可以表示为向量组 $T v_{1} , \ldots , T v_{n}$ 的线性组合，根据张成的定义（原书2.7），可得 $T v_{1} , \ldots , T v_{n}$ 张成 $\operatorname{im} T$。
\end{solve}

\begin{exercise}\label{prob:P-3B-11}
设 $V$ 是有限维向量空间，$T \in \mathcal L (V , W)$。证明：存在 $V$ 的子空间 $U$，使得

\[U \cap \ker T = \{0\} \quad \text{且} \quad \operatorname{im} T = \{T u : u \in U\}\]
\end{exercise}
\noindent 来源：\href{https://github.com/szdytom/LADRSolutions/blob/main/sections/3B.typ}{P-3B-11}（原书 3B 第 11 题；方而静题解）。
\begin{solve}
设 $v_{1} , \ldots , v_{n}$ 是 $V$ 的一组基，则根据习题~\ref{app:E-domain-span-to-range-span}，$T v_{1} , \ldots , T v_{n}$ 张成 $\operatorname{im} T$。在根据“每个张成组都包含基”（原书2.30），不妨设 $T v_{1} , \ldots , T v_{m}$ 是 $\operatorname{im} T$ 的一组基。

令 $U = \operatorname{span} (v_{1} , \ldots , v_{m})$。将 $T$ 视作 $U \to W$ 的线性映射，则根据习题~\ref{app:E-domain-span-to-range-span}，$\operatorname{im} T = \{T u : u \in U\}$。现在证明 $U \cap \ker T = \{0\}$。设 $u \in U \cap \ker T$，则存在 $a_{1} , \ldots , a_{m} \in \mathbf F$，使得

\[u = a_{1} v_{1} + \cdots + a_{m} v_{m}\]

由于 $u \in \ker T$，根据线性映射的定义，有

\[0 = T u = a_{1} T v_{1} + \cdots + a_{m} T v_{m}\]

由于 $T v_{1} , \ldots , T v_{m}$ 是线性无关的，故 $a_{1} = \cdots = a_{m} = 0$，即 $u = 0$。因此 $U \cap \ker T = \{0\}$。
\end{solve}

\begin{exercise}\label{prob:P-3B-12}
设 $T$ 是 $\mathbf F^{4} \to \mathbf F^{2}$ 的线性映射，且

\[\ker T = \{(x_{1} , x_{2} , x_{3} , x_{4}) \in \mathbf F^{4} : x_{1} = 5 x_{2} \text{且} x_{3} = 7 x_{4}\}\]

证明：$T$ 是满射。
\end{exercise}
\noindent 来源：\href{https://github.com/szdytom/LADRSolutions/blob/main/sections/3B.typ}{P-3B-12}（原书 3B 第 12 题；方而静题解）。
\begin{solve}
注意到，取

\[\begin{aligned}v_{1} & = (5 , 1 , 0 , 0) \\ v_{2} & = (0 , 0 , 7 , 1)\end{aligned}\]

则 $\ker T = \operatorname{span} (v_{1} , v_{2})$，因此 $\dim \ker T = 2$。根据线性映射基本定理（原书3.21），

\[\dim \mathbf F^{4} = \dim \ker T + \dim \operatorname{im} T\]

解得 $\dim \operatorname{im} T = 2$，即 $\dim \operatorname{im} T = \dim \mathbf F^{2}$，根据“某空间中与之维数相同的子空间即为该空间本身”（原书2.39），$\operatorname{im} T = \mathbf F^{2}$，即 $T$ 是满射。
\end{solve}

\begin{exercise}\label{prob:P-3B-13}
设 $U$ 是 $\mathbf R^{8}$ 的 $3$ 维子空间，$T$ 是 $\mathbf R^{8} \to \mathbf R^{5}$ 的线性映射，使得 $\ker T = U$。证明：$T$ 是满射。
\end{exercise}
\noindent 来源：\href{https://github.com/szdytom/LADRSolutions/blob/main/sections/3B.typ}{P-3B-13}（原书 3B 第 13 题；方而静题解）。
\begin{solve}
根据线性映射基本定理（原书3.21），有

\[\dim \mathbf R^{8} = \dim U + \dim \operatorname{im} T\]

解得 $\dim \operatorname{im} T = 5$。根据“某空间中与之维数相同的子空间即为该空间本身”（原书2.39），$\operatorname{im} T = \mathbf R^{5}$，即 $T$ 是满射。
\end{solve}

\begin{exercise}\label{prob:P-3B-14}
证明：不存在 $\mathbf F^{5} \to \mathbf F^{2}$ 的线性映射，使得其零空间等于

\[\{(x_{1} , x_{2} , x_{3} , x_{4} , x_{5}) \in \mathbf F^{5} : x_{1} = 3 x_{2} \text{且} x_{3} = x_{4} = x_{5}\}\]
\end{exercise}
\noindent 来源：\href{https://github.com/szdytom/LADRSolutions/blob/main/sections/3B.typ}{P-3B-14}（原书 3B 第 14 题；方而静题解）。
\begin{solve}
假设这样的线性映射 $T$ 存在。注意到，取

\[\begin{aligned}v_{1} & = (3 , 1 , 0 , 0 , 0) \\ v_{2} & = (0 , 0 , 1 , 1 , 1)\end{aligned}\]

则 $\ker T = \operatorname{span} (v_{1} , v_{2})$，因此 $\dim \ker T = 2$。根据线性映射基本定理（原书3.21），

\[\dim \mathbf F^{5} = \dim \ker T + \dim \operatorname{im} T\]

解得 $\dim \operatorname{im} T = 3$，然而 $\dim \mathbf F^{2} = 2$，根据“子空间的维数不超过该空间的维数”（原书2.37），$\operatorname{im} T \subseteq \mathbf F^{2}$，因此 $\dim \operatorname{im} T \le \dim \mathbf F^{2} = 2$，矛盾。因此不存在这样的线性映射。
\end{solve}

\begin{exercise}\label{prob:P-3B-15}
设 $V$ 上存在一个线性映射，使得其零空间和值域都是有限维的，证明：$V$ 是有限维的。
\end{exercise}
\noindent 来源：\href{https://github.com/szdytom/LADRSolutions/blob/main/sections/3B.typ}{P-3B-15}（原书 3B 第 15 题；方而静题解）。
\begin{solve}
设 $T \in \mathcal L (V)$，使得 $\dim \ker T = m$ 且 $\dim \operatorname{im} T = n$。进一步可设 $u_{1} , \ldots , u_{m}$ 是 $\ker T$ 的一组基，$T v_{1} , \ldots , T v_{n}$ 是 $\operatorname{im} T$ 的一组基，其中 $v_{1} , \ldots , v_{n} \in V$。

设 $w \in V$，则 $T w \in \operatorname{im} T$，因此存在 $a_{1} , \ldots , a_{n} \in \mathbf F$，使得

\[T w = a_{1} T v_{1} + \cdots + a_{n} T v_{n} = T (a_{1} v_{1} + \cdots + a_{n} v_{n})\]

记 $v = a_{1} v_{1} + \cdots + a_{n} v_{n}$，则 $T (w - v) = 0$，故 $w - v \in \ker T$，即存在 $b_{1} , \ldots , b_{m} \in \mathbf F$，使得

\[w - v = b_{1} u_{1} + \cdots + b_{m} u_{m}\]

这说明 $w$ 可以表示为

\[w = (a_{1} v_{1} + \cdots + a_{n} v_{n}) + (b_{1} u_{1} + \cdots + b_{m} u_{m})\]

即 $w \in \operatorname{span} (u_{1} , \ldots , u_{m} , v_{1} , \ldots , v_{n})$。$V$ 可以被有限个向量张成，因此 $V$ 是有限维的。
\end{solve}

\begin{exercise}\label{prob:P-3B-16}
设 $V$ 和 $W$ 都是有限维向量空间，证明：存在 $V \to W$ 的单的线性映射，当且仅当 $\dim V \le \dim W$。
\end{exercise}
\noindent 来源：\href{https://github.com/szdytom/LADRSolutions/blob/main/sections/3B.typ}{P-3B-16}（原书 3B 第 16 题；方而静题解）。
\begin{solve}
首先假设存在 $T \in \mathcal L (V , W)$ 是单射。根据“单射性 $\iff$ 零空间为 $\{0\}$”（原书3.15），$\ker T = \{0\}$，因此 $\dim \ker T = 0$。根据线性映射基本定理（原书3.21），有

\[\dim V = \dim \ker T + \dim \operatorname{im} T = \dim \operatorname{im} T\]

由于 $\operatorname{im} T \subseteq W$，根据子空间的维数（原书2.37），$\dim \operatorname{im} T \le \dim W$，因此 $\dim V \le \dim W$。

现在假设 $\dim V \le \dim W$。设 $v_{1} , \ldots , v_{n}$ 是 $V$ 的一组基，$w_{1} , \ldots , w_{n}$ 是 $W$ 上的一个线性无关组。根据线性映射引理（原书3.4），存在 $T \in \mathcal L (V , W)$，使得对于任意 $i \in \{1 , \ldots , n\}$，有 $T v_{i} = w_{i}$。

设 $v \in V$，使得 $T v = 0$。将 $v$ 表示为 $v = a_{1} v_{1} + \cdots + a_{n} v_{n}$，其中 $a_{1} , \ldots , a_{n} \in \mathbf F$。则

\[0 = T v = a_{1} T v_{1} + \cdots + a_{n} T v_{n} = a_{1} w_{1} + \cdots + a_{n} w_{n}\]

由于 $w_{1} , \ldots , w_{n}$ 是线性无关的，故 $a_{1} = \cdots = a_{n} = 0$，即 $v = 0$。因此 $\ker T = \{0\}$，根据“单射性 $\iff$ 零空间为 $\{0\}$”（原书3.15），$T$ 是单射。
\end{solve}

\begin{exercise}\label{prob:P-3B-17}
设 $V$ 和 $W$ 都是有限维向量空间，证明：存在 $V \to W$ 的满的线性映射，当且仅当 $\dim V \ge \dim W$。
\end{exercise}
\noindent 来源：\href{https://github.com/szdytom/LADRSolutions/blob/main/sections/3B.typ}{P-3B-17}（原书 3B 第 17 题；方而静题解）。
\begin{solve}
首先假设存在 $T \in \mathcal L (V , W)$ 是满射。根据线性映射基本定理（原书3.21），有

\[\dim V = \dim \ker T + \dim \operatorname{im} T\]

由于 $\operatorname{im} T = W$，因此 $\dim \operatorname{im} T = \dim W$，解得 $\dim V = \dim \ker T + \dim W \ge \dim W$。

现在假设 $\dim V \ge \dim W$。设 $w_{1} , \ldots , w_{n}$ 是 $W$ 的一组基，$v_{1} , \ldots , v_{n}$ 是 $V$ 上的一个线性无关向量组。根据线性映射引理（原书3.4），存在 $S \in \mathcal L (\operatorname{span} (v_{1} , \ldots , v_{n}) , W)$，使得对于任意 $i \in \{1 , \ldots , n\}$，有 $S v_{i} = w_{i}$。进一步，根据习题~\ref{app:E-extend-linear-map}，存在 $T \in \mathcal L (V , W)$，使得对于任意 $v \in \operatorname{span} (v_{1} , \ldots , v_{n})$，有 $T v = S v$。

设 $w \in W$，则可以将 $w$ 表示为 $w = a_{1} w_{1} + \cdots + a_{n} w_{n}$，其中 $a_{1} , \ldots , a_{n} \in \mathbf F$。则

\[w = a_{1} T v_{1} + \cdots + a_{n} T v_{n} = T (a_{1} v_{1} + \cdots + a_{n} v_{n})\]

这说明对于任意 $w \in W$，都存在一个向量在 $V$ 中，使得其通过 $T$ 映射到该向量。因此 $T$ 是满射。
\end{solve}

\begin{exercise}\label{prob:P-3B-18}
设 $V$ 和 $W$ 都是有限维向量空间，$U$ 是 $V$ 的子空间。证明：存在 $T \in \mathcal L (V , W)$，使得 $\ker T = U$ 当且仅当 $\dim U \ge \dim V - \dim W$。
\end{exercise}
\noindent 来源：\href{https://github.com/szdytom/LADRSolutions/blob/main/sections/3B.typ}{P-3B-18}（原书 3B 第 18 题；方而静题解）。
\begin{solve}
首先假设存在 $T \in \mathcal L (V , W)$，使得 $\ker T = U$。根据线性映射基本定理（原书3.21），有

\[\dim V = \dim \ker T + \dim \operatorname{im} T\]

由于 $\ker T = U$，因此 $\dim \ker T = \dim U$。由于 $\operatorname{im} T \subseteq W$，根据子空间的维数（原书2.37），$\dim \operatorname{im} T \le \dim W$，因此

\[\dim V \le \dim U + \dim W\]

解得 $\dim U \ge \dim V - \dim W$。

现在假设 $\dim U \ge \dim V - \dim W$。设 $u_{1} , \ldots , u_{m}$ 是 $U$ 的一组基。根据“每个线性无关组都可以扩展为基”（原书2.31），存在 $v_{1} , \ldots , v_{n} \in V$，使得向量组 $u_{1} , \ldots , u_{m} , v_{1} , \ldots , v_{n}$ 是 $V$ 的一组基。其中 $m = \dim U$，$n = \dim V - \dim U$。由于 $\dim U \ge \dim V - \dim W$，解得 $n \le \dim W$。

设 $w_{1} , \ldots , w_{n} \in W$ 是线性无关组，根据线性映射引理（原书3.4），存在 $T \in \mathcal L (V , W)$，使得对于任意 $i \in \{1 , \ldots , m\}$，有 $T u_{i} = 0$，且对于任意 $i \in \{1 , \ldots , n\}$，有 $T v_{i} = w_{i}$。

设 $v \in \ker T$，则存在 $a_{1} , \ldots , a_{m} , b_{1} , \ldots , b_{n} \in \mathbf F$，使得

\[v = a_{1} u_{1} + \cdots + a_{m} u_{m} + b_{1} v_{1} + \cdots + b_{n} v_{n}\]

其中

\[0 = T v = b_{1} T v_{1} + \cdots + b_{n} T v_{n} = b_{1} w_{1} + \cdots + b_{n} w_{n}\]

由于 $w_{1} , \ldots , w_{n}$ 是线性无关的，故 $b_{1} = \cdots = b_{n} = 0$，即 $v \in U$。因此 $\ker T \subseteq U$。另一方面，显然 $U \subseteq \ker T$，因此 $\ker T = U$。
\end{solve}

\begin{exercise}\label{prob:P-3B-19}
设 $W$ 是有限维向量空间，$T \in \mathcal L (V , W)$。证明：$T$ 是单射，当且仅当，存在 $S \in \mathcal L (W , V)$，使得 $S T$ 是 $V$ 上的恒等映射。
\end{exercise}
\noindent 来源：\href{https://github.com/szdytom/LADRSolutions/blob/main/sections/3B.typ}{P-3B-19}（原书 3B 第 19 题；方而静题解）。
\begin{solve}
首先假设 $T$ 是单射。根据值域是子空间（原书3.18），$\operatorname{im} T$ 是 $W$ 的子空间，进一步根据子空间的维数（原书2.37），可得 $\operatorname{im} T$ 是有限维的。设 $T u_{1} , \ldots , T u_{n}$ 是 $\operatorname{im} T$ 的一组基，其中 $u_{1} , \ldots , u_{n} \in V$。

设 $u \in V$，则可以将 $T u$ 表示为

\[T u = a_{1} T u_{1} + \cdots + a_{n} T u_{n} = T (a_{1} u_{1} + \cdots + a_{n} u_{n})\]

其中 $a_{1} , \ldots , a_{n} \in \mathbf F$。根据单射的定义（原书3.14），$T$ 是单射，因此 $u \in \operatorname{span} (u_{1} , \ldots , u_{n})$。因此向量组 $u_{1} , \ldots , u_{n}$ 张成 $V$。故 $V$ 是有限维向量空间。

故可设 $v_{1} , \ldots , v_{n}$ 是 $V$ 的一组基。对于 $i \in \{1 , \ldots , n\}$，令 $w_{i} = T v_{i}$。根据习题~\ref{app:E-indep-preservance-under-inj}，$w_{1} , \ldots , w_{n}$ 是线性无关的。

根据线性映射引理（原书3.4），存在 $R \in \mathcal L (\operatorname{span} (w_{1} , \ldots , w_{n}) , V)$，使得对于任意 $i \in \{1 , \ldots , n\}$，有 $R w_{i} = v_{i}$。进一步，根据习题~\ref{app:E-extend-linear-map}，存在 $S \in \mathcal L (W , V)$，使得对于任意 $w \in \operatorname{span} (w_{1} , \ldots , w_{n})$，有 $S w = R w$。

设 $v \in V$，则可以将 $v$ 表示为 $v = a_{1} v_{1} + \cdots + a_{n} v_{n}$，其中 $a_{1} , \ldots , a_{n} \in \mathbf F$。则

\[\begin{aligned}S T v & = S (a_{1} T v_{1} + \cdots + a_{n} T v_{n}) \\ & = S (a_{1} w_{1} + \cdots + a_{n} w_{n}) \\ & = a_{1} S w_{1} + \cdots + a_{n} S w_{n} \\ & = a_{1} v_{1} + \cdots + a_{n} v_{n} = v\end{aligned}\]

这说明 $S T$ 确实是 $V$ 上的恒等映射。

现在假设存在 $S \in \mathcal L (W , V)$，使得 $S T$ 是 $V$ 上的恒等映射。设 $v \in \ker T$，根据“线性映射将 $0$ 映射到 $0$”（原书3.10），有

\[v = S T v = S 0 = 0\]

这说明 $\ker T = \{0\}$，根据“单射性 $\iff$ 零空间为 $\{0\}$”（原书3.15），$T$ 是单射。
\end{solve}

\begin{exercise}\label{prob:P-3B-20}
设 $W$ 是有限维的向量空间，$T \in \mathcal L (V , W)$。证明：$T$ 是满射，当且仅当，存在 $S \in \mathcal L (W , V)$，使得 $T S$ 是 $W$ 上的恒等映射。
\end{exercise}
\noindent 来源：\href{https://github.com/szdytom/LADRSolutions/blob/main/sections/3B.typ}{P-3B-20}（原书 3B 第 20 题；方而静题解）。
\begin{solve}
首先假设 $T$ 是满射。令 $n = \dim W$。如果 $V$ 是有限维的，则根据“映射到更高维空间上的线性映射不是满射”（原书3.24）的逆否命题，$\dim V \ge \dim W$。因此可设 $v_{1} , \ldots , v_{n}$ 是 $V$ 上的一个线性无关向量组。如果 $V$ 是无限维的，自然也可以取 $v_{1} , \ldots , v_{n}$ 为 $V$ 上的线性无关向量组。

现在对于 $i \in \{1 , \ldots , n\}$，令 $w_{i} = T v_{i}$。根据习题~\ref{app:E-indep-preservance-under-inj}，$w_{1} , \ldots , w_{n}$ 是线性无关的。进一步，根据“长度恰当的线性无关组是基”（原书2.38），$w_{1} , \ldots , w_{n}$ 是 $W$ 的一组基。于是，根据线性映射引理（原书3.4），存在 $S \in \mathcal L (W , V)$，使得对于任意 $i \in \{1 , \ldots , n\}$，有 $S w_{i} = v_{i}$。

设 $w \in W$，则可以将 $w$ 表示为 $w = a_{1} w_{1} + \cdots + a_{n} w_{n}$，其中 $a_{1} , \ldots , a_{n} \in \mathbf F$。则

\[\begin{aligned}T S w & = T (a_{1} S w_{1} + \cdots + a_{n} S w_{n}) \\ & = T (a_{1} v_{1} + \cdots + a_{n} v_{n}) \\ & = a_{1} T v_{1} + \cdots + a_{n} T v_{n} \\ & = a_{1} w_{1} + \cdots + a_{n} w_{n} = w\end{aligned}\]

这说明 $T S$ 确实是 $W$ 上的恒等映射。

现在假设存在 $S \in \mathcal L (W , V)$，使得 $T S$ 是 $W$ 上的恒等映射。设 $w \in W$，则可取 $v = S w \in V$，使得 $T v = T S w = w$。这说明对于任意 $w \in W$，都存在一个向量在 $V$ 中，使得其通过 $T$ 映射到该向量。因此 $T$ 是满射。
\end{solve}

\begin{exercise}\label{prob:P-3B-21}
设 $V$ 是有限维向量空间，$T \in \mathcal L (V , W)$，$U$ 是 $W$ 的子空间。证明：$\{v \in V : T v \in U\}$ 是 $V$ 的子空间，且

\[\dim \{v \in V : T v \in U\} = \dim \ker T + \dim(U \cap \operatorname{im} T)\]
\end{exercise}
\noindent 来源：\href{https://github.com/szdytom/LADRSolutions/blob/main/sections/3B.typ}{P-3B-21}（原书 3B 第 21 题；方而静题解）。
\begin{solve}
记 $V_{U} = \{v \in V : T v \in U\}$。为了说明 $V_{U}$ 是 $V$ 的子空间，我们逐条验证“子空间的条件”（原书1.34）中给出的要求：

\begin{description}
\item[加法单位元] $0 \in V_{U}$。
证明：根据“线性映射将 $0$ 映射到 $0$”（原书3.10），有 $T 0 = 0 \in U$，因此 $0 \in V_{U}$。

\item[加法封闭性] $u , w \in V_{U}$ 意味着 $u + w \in V_{U}$。
证明：设 $u , w \in V_{U}$，则 $T u \in U$ 且 $T w \in U$。因此 $T (u + w) = T u + T w \in U$，即 $u + w \in V_{U}$。

\item[乘法封闭性] $\lambda \in \mathbf F$ 且 $v \in V_{U}$ 意味着 $\lambda v \in V_{U}$。
证明：设 $v \in V_{U}$，则 $T v \in U$。因此 $T (\lambda v) = \lambda T v \in U$，即 $\lambda v \in V_{U}$。

这说明 $V_{U}$ 是 $V$ 的子空间。

现在证明 $\dim V_{U} = \dim \ker T + \dim(U \cap \operatorname{im} T)$。设 $v_{1} , \ldots , v_{m}$ 是 $V_{U}$ 的一组基，根据“每个线性无关组都可以扩展为基”（原书2.31），存在 $v_{m + 1} , \ldots , v_{n} \in V$，使得向量组 $v_{1} , \ldots , v_{n}$ 是 $V$ 的一组基。

根据线性映射引理（原书3.4），存在 $T|_{V_{U}} \in \mathcal L (V_{U} , U)$，使得对于任意 $i \in \{1 , \ldots , m\}$，有 $T|_{V_{U}} v_{i} = T v_{i}$。

首先，我们证明 $\dim \ker T|_{V_{U}} = \dim \ker T$。设 $v \in \ker T|_{V_{U}}$，则 $T|_{V_{U}} v = 0$，因此 $T v = 0$，即 $v \in \ker T$。这说明 $\ker T|_{V_{U}} \subseteq \ker T$。另一方面，设 $v \in \ker T$，则 $T v = 0 \in U$，因此 $v \in V_{U}$，故 $T|_{V_{U}} v = 0$，即 $v \in \ker T|_{V_{U}}$。这说明 $\ker T \subseteq \ker T|_{V_{U}}$。因此 $\dim \ker T|_{V_{U}} = \dim \ker T$。

其次，我们证明 $\dim \operatorname{im} T|_{V_{U}} = \dim(U \cap \operatorname{im} T)$。设 $w \in \operatorname{im} T|_{V_{U}}$，则存在 $v \in V_{U}$，使得 $T|_{V_{U}} v = w$。因此 $T v = w$，且 $T v \in U$，即 $w \in U \cap \operatorname{im} T$。这说明 $\operatorname{im} T|_{V_{U}} \subseteq U \cap \operatorname{im} T$。

另一方面，设 $w \in U \cap \operatorname{im} T$，则存在 $v \in V$，使得 $T v = w$ 且 $w \in U$。由于 $v \in V_{U}$，因此 $T|_{V_{U}} v = w$。这说明 $U \cap \operatorname{im} T \subseteq \operatorname{im} T|_{V_{U}}$。因此 $\dim \operatorname{im} T|_{V_{U}} = \dim(U \cap \operatorname{im} T)$。

现在根据线性映射基本定理（原书3.21），有

\end{description}
\[\dim V_{U} = \dim \ker T|_{V_{U}} + \dim \operatorname{im} T|_{V_{U}}\]

代入上面的结果，得到 $\dim V_{U} = \dim \ker T + \dim(U \cap \operatorname{im} T)$。
\end{solve}

\begin{exercise}\label{prob:P-3B-22}
设 $U$ 和 $V$ 都是有限维向量空间，$S \in \mathcal L (V , W)$，$T \in \mathcal L (U , V)$. 证明：

\[\dim \ker S T \le \dim \ker S + \dim \ker T\]
\end{exercise}
\noindent 来源：\href{https://github.com/szdytom/LADRSolutions/blob/main/sections/3B.typ}{P-3B-22}（原书 3B 第 22 题；方而静题解）。
\begin{solve}
令 $N = \ker S T$。由于 $S T \in \mathcal L (U , W)$，故 $N$ 是 $U$ 的子空间。设 $u_{1} , \ldots , u_{m}$ 是 $N$ 的一组基。根据线性映射引理（原书3.4），存在 $T|_N \in \mathcal L (N , V)$，使得 $T|_N u_{i} = T u_{i}$。设 $u \in N$，根据零空间的定义（原书3.11），有 $S T u = 0$，故 $\operatorname{im} T|_N \subseteq \ker S$，即

\begin{equation}\label{app:eq:3B-c-range-TN-leq-null-S}\dim \operatorname{im} T|_N \le \dim \ker S\end{equation}

根据线性映射基本定理（原书3.21），有

\begin{equation}\label{app:eq:3B-c-dim-N-eq-null-TN-plus-range-TN}\dim N = \dim \ker T|_N + \dim \operatorname{im} T|_N\end{equation}

注意到

\begin{equation}\label{app:eq:3B-c-null-TN-eq-N-inter-null-T}\ker T|_N = \{u \in N : T|_N u = 0\} = \{u \in N : T u = 0\} = N \cap \ker T\end{equation}

将式~\eqref{app:eq:3B-c-range-TN-leq-null-S} 和式~\eqref{app:eq:3B-c-null-TN-eq-N-inter-null-T} 代入式~\eqref{app:eq:3B-c-dim-N-eq-null-TN-plus-range-TN}，得到

\[\dim N \le \dim(N \cap \ker T) + \dim \ker S\]

由于 $N \cap \ker T \subseteq \ker T$，根据子空间的维数（原书2.37），$\dim(N \cap \ker T) \le \dim \ker T$，因此 $\dim \ker S T \le \dim \ker T + \dim \ker S$。
\end{solve}

\begin{exercise}\label{prob:P-3B-23}
设 $U$ 和 $V$ 都是有限维向量空间，$S \in \mathcal L (V , W)$，$T \in \mathcal L (U , V)$. 证明：

\[\dim \operatorname{im} S T \le \min \{\dim \operatorname{im} S , \dim \operatorname{im} T\}\]
\end{exercise}
\noindent 来源：\href{https://github.com/szdytom/LADRSolutions/blob/main/sections/3B.typ}{P-3B-23}（原书 3B 第 23 题；方而静题解）。
\begin{solve}
首先证明 $\dim \operatorname{im} S T \le \dim \operatorname{im} S$。设 $u \in U$，则 $S T u = S (T u) \in \operatorname{im} S$，故 $\operatorname{im} S T \subseteq \operatorname{im} S$，即 $\dim \operatorname{im} S T \le \dim \operatorname{im} S$。

现在证明 $\dim \operatorname{im} S T \le \dim \operatorname{im} T$。令 $I = \operatorname{im} T$，则 $I$ 是 $V$ 的子空间。根据线性映射引理（原书3.4），存在 $S|_I \in \mathcal L (I , W)$，使得对于任意 $v \in I$，有 $S|_I v = S v$。设 $u \in U$，则 $T u \in I$，因此 $S T u = S|_I (T u)$。故 $\operatorname{im} S T = \operatorname{im} S|_I T \subseteq \operatorname{im} S|_I$，即 $\dim \operatorname{im} S T \le \dim \operatorname{im} S|_I$。

由于 $S|_I$ 是从 $I$ 到 $W$ 的线性映射，根据线性映射基本定理（原书3.21），有

\[\dim I = \dim \ker S|_I + \dim \operatorname{im} S|_I\]

即 $\dim I \ge \dim \operatorname{im} S|_I$，又因为 $\dim \operatorname{im} S T \le \dim \operatorname{im} S|_I$，故 $\dim \operatorname{im} S T \le \dim \operatorname{im} T$。

综上所述，$\dim \operatorname{im} S T \le \min \{\dim \operatorname{im} S , \dim \operatorname{im} T\}$。
\end{solve}

\begin{exercise}\label{prob:P-3B-24}
\begin{enumerate}[label=(\alph*)]
\item 设 $\dim V = 5$，且 $S , T \in \mathcal L (V)$，使得 $S T = 0$。证明 $\dim \operatorname{im} T S \le 2$；
\item 给出一例：$S , T \in \mathcal L (\mathbf F^{5})$，使得 $S T = 0$ 且 $\dim \operatorname{im} T S = 2$。
\end{enumerate}
\end{exercise}
\noindent 来源：\href{https://github.com/szdytom/LADRSolutions/blob/main/sections/3B.typ}{P-3B-24}（原书 3B 第 24 题；方而静题解）。
\begin{solve}
对于（a），根据线性映射基本定理（原书3.21），有

\[\begin{aligned}\dim V & = \dim \ker S + \dim \operatorname{im} S \\ \dim V & = \dim \ker T + \dim \operatorname{im} T\end{aligned}\]

设 $v \in V$，则 $S (T v) = S T v = 0$，故 $\operatorname{im} T \subseteq \ker S$，即 $\dim \operatorname{im} T \le \dim \ker S$。再代入 $\dim V = 5$，整理得

\[\dim \operatorname{im} T + \dim \operatorname{im} S \le \dim \ker S + \dim \operatorname{im} S = 5\]

注意到 $3 + 3 = 6 > 5$，故 $\dim \operatorname{im} T < 3$ 或 $\dim \operatorname{im} S < 3$。分类讨论：当 $\dim \operatorname{im} T < 3$ 时，即 $\dim \operatorname{im} T \le 2$。考虑到

\[\operatorname{im} T S = \{T v : v \in \operatorname{im} S\} \subseteq \operatorname{im} T\]

这说明 $\dim \operatorname{im} T S \le \dim \operatorname{im} T \le 2$。

另一方面，当 $\dim \operatorname{im} S < 3$ 时，即 $\dim \operatorname{im} S \le 2$，故 $\dim \ker S \ge 3$。根据“线性映射将 $0$ 映射到 $0$”（原书3.10），$0 \in \ker T$，故

\[\ker S = \{v \in V : S v = 0\} \subseteq \{v \in V : S v \in \ker T\} = \ker T S\]

这说明 $3 \le \dim \ker S \le \dim \ker T S$，再根据线性映射基本定理（原书3.21），可得 $\dim \operatorname{im} T S \le 2$。

对于（b），设 $v_{1} , \ldots , v_{5}$ 是 $\mathbf F^{5}$ 的一组基。根据线性映射引理（原书3.4），存在 $S , T \in \mathcal L (V)$，使得 $S v_{1} = S v_{2} = S v_{3} = 0$，$S v_{4} = v_{4}$，$S v_{5} = v_{5}$，以及 $T v_{1} = T v_{2} = T v_{3} = 0$，$T v_{4} = v_{1}$，$T v_{5} = v_{2}$。

设 $v \in \mathbf F^{5}$，将 $v$ 表示为 $v = a_{1} v_{1} + \cdots + a_{5} v_{5}$，其中 $a_{1} , \ldots , a_{5} \in \mathbf F$。则

\[\begin{aligned}S T v & = S (a_{1} T v_{1} + \cdots + a_{5} T v_{5}) \\ & = S (a_{4} v_{1} + a_{5} v_{2}) \\ & = a_{4} S v_{1} + a_{5} S v_{2} \\ & = 0\end{aligned}\]

这说明 $S T = 0$。而

\[\begin{aligned}T S v & = T (a_{1} S v_{1} + \cdots + a_{5} S v_{5}) \\ & = T (a_{4} v_{4} + a_{5} v_{5}) \\ & = a_{4} T v_{4} + a_{5} T v_{5} \\ & = a_{4} v_{1} + a_{5} v_{2}\end{aligned}\]

这说明 $\operatorname{im} T S = \operatorname{span} (v_{1} , v_{2})$，因此 $\dim \operatorname{im} T S = 2$。综上所述，$S , T \in \mathcal L (\mathbf F^{5})$，使得 $S T = 0$ 且 $\dim \operatorname{im} T S = 2$。
\end{solve}

\begin{exercise}\label{prob:P-3B-25}
设 $W$ 是有限维向量空间，$S , T \in \mathcal L (V , W)$。证明：$\ker S \subseteq \ker T$，当且仅当，存在 $E \in \mathcal L (W)$，使得 $T = E S$。
\end{exercise}
\noindent 来源：\href{https://github.com/szdytom/LADRSolutions/blob/main/sections/3B.typ}{P-3B-25}（原书 3B 第 25 题；方而静题解）。
\begin{solve}
若 $T=ES$，则 $Sv=0$ 蕴含 $Tv=E0=0$，即 $\ker S\subseteq\ker T$。
反过来，设该包含关系成立，在 $\operatorname{im} S$ 上定义 $E_0(Sv)=Tv$。此定义与原像选择无关：若 $Sv=Sv'$，则 $v-v'\in\ker S\subseteq\ker T$，故 $Tv=Tv'$。由 $S,T$ 的线性性，$E_0$ 线性。将 $\operatorname{im} S$ 的一组基扩充为有限维空间 $W$ 的基，在补入的基向量上令像为零，即得到 $E\in\mathcal L(W)$。此时对所有 $v\in V$ 都有 $ESv=E_0Sv=Tv$，故 $T=ES$。
\end{solve}

\begin{exercise}\label{prob:P-3B-26}
设 $V$ 是有限维向量空间，$S , T \in \mathcal L (V)$。证明：$\operatorname{im} S \subseteq \operatorname{im} T$，当且仅当，存在 $E \in \mathcal L (V)$，使得 $S = T E$。
\end{exercise}
\noindent 来源：\href{https://github.com/szdytom/LADRSolutions/blob/main/sections/3B.typ}{P-3B-26}（原书 3B 第 26 题；方而静题解）。
\begin{solve}
首先假设 $\operatorname{im} S \subseteq \operatorname{im} T$。设 $v_{1} , \ldots , v_{m}$ 是 $V$ 的一组基。对于每个 $i \in \{1 , \ldots , m\}$，由于 $S v_{i} \in \operatorname{im} S \subseteq \operatorname{im} T$，因此存在 $u_{i} \in V$，使得 $T u_{i} = S v_{i}$。

根据线性映射引理（原书3.4），存在 $E \in \mathcal L (V , V)$，使得对于任意 $i \in \{1 , \ldots , m\}$，有 $E v_{i} = u_{i}$。设 $v \in V$，将 $v$ 表示为 $v = a_{1} v_{1} + \cdots + a_{m} v_{m}$，其中 $a_{1} , \ldots , a_{m} \in \mathbf F$。则

\[S v = S (\sum_{k = 1}^{m} a_{k} v_{k}) = T (\sum_{k = 1}^{m} a_{k} u_{k}) = T (\sum_{k = 1}^{m} a_{k} E v_{k}) = T E v\]

这说明 $S = T E$。

另一方面，现在假设存在 $E \in \mathcal L (V)$，使得 $S = T E$。设 $w \in \operatorname{im} S$，则存在 $v \in V$，使得 $S v = w$。因此

\[w = S v = T E v = T (E v) \in \operatorname{im} T\]

这说明 $\operatorname{im} S \subseteq \operatorname{im} T$。

综上所述，$\operatorname{im} S \subseteq \operatorname{im} T$，当且仅当，存在 $E \in \mathcal L (V)$，使得 $S = T E$。
\end{solve}

\begin{exercise}\label{prob:P-3B-27}
设 $P \in \mathcal L (V)$，且 $P^{2} = P$。证明：$V = \ker P \oplus \operatorname{im} P$。
\end{exercise}
\noindent 来源：\href{https://github.com/szdytom/LADRSolutions/blob/main/sections/3B.typ}{P-3B-27}（原书 3B 第 27 题；方而静题解）。
\begin{solve}
设 $v \in V$，则 $P v = P (P v)$，故 $P (v - P v) = 0$，即 $v - P v \in \ker P$，另一方面，$P v \in \operatorname{im} P$，即

\[v = (v - P v) + P v\]

其中 $v - P v \in \ker N$ 且 $P v \in \operatorname{im} P$，故 $V = \ker P + \operatorname{im} P$。

下面说明这个和是直和。设 $v \in \ker P \cap \operatorname{im} P$，则 $P v = 0$ 且存在 $w \in V$，使得 $v = P w$。故

\[0 = P v = P^{2} w = P w = v\]

即 $\ker P \cap \operatorname{im} P = \{0\}$。因此，根据“两个子空间的直和”（原书1.46），得 $V = \ker P \oplus \operatorname{im} P$。
\end{solve}

\begin{exercise}\label{prob:P-3B-28}\label{app:E-Poly-lower-const-degree-surj}
设 $D \in \mathcal L (\mathcal P (\mathbf R))$，满足对于任意非常数多项式 $p \in \mathcal P (\mathbf R)$，都有 $\deg D p = (\deg p) - 1$。证明：$D$ 是满射。

\begin{remark}
上面的记号 $D$ 是用来让你想起微分映射\footnote{注意，微分映射是满足题设条件的映射，但并非满足题设条件的映射就一定是微分映射。}，它将多项式 $p$ 映射到其导数 $p'$。
\end{remark}
\end{exercise}
\noindent 来源：\href{https://github.com/szdytom/LADRSolutions/blob/main/sections/3B.typ}{P-3B-28}（原书 3B 第 28 题；方而静题解）。
\begin{solve}
对 $k\ge0$ 定义 $q_k=D(x^{k+1})$。题设直接给出 $\deg q_k=k$。若 $p\ne0$ 且 $\deg p=m$，由习题~\ref{app:E-Poly-space-basis}，$q_0,\ldots,q_m$ 构成 $\mathcal P_m(\mathbf R)$ 的一组基。因此可写成 $p=\sum_{k=0}^m a_kq_k$，从而
\[p=D\left(\sum_{k=0}^m a_kx^{k+1}\right).\]
零多项式也有原像 $0$，所以 $D$ 满射。
\end{solve}

\begin{exercise}\label{prob:P-3B-29}
设 $p \in \mathcal P (\mathbf R)$。证明：存在多项式 $q \in \mathcal P (\mathbf R)$，使得 $5 q'' + 3 q' = p$。

\begin{remark}
这道题不一定要用线性代数，但是线性代数的解答更有趣。
\end{remark}
\end{exercise}
\noindent 来源：\href{https://github.com/szdytom/LADRSolutions/blob/main/sections/3B.typ}{P-3B-29}（原书 3B 第 29 题；方而静题解）。
\begin{solve}
令 $\mathcal P (\mathbf R)$ 上的映射 $T$ 为 $p \mapsto 5 p'' + 3 p'$，容易验证 $T$ 为线性映射，且对于任意 $p \in \mathcal P (\mathbf R)$，$\deg T p = \deg p - 1$。于是，根据习题~\ref{app:E-Poly-lower-const-degree-surj}，$T$ 是满射。

这说明，对于任意 $p \in \mathcal P (\mathbf R)$，都存在 $q \in \mathcal P (\mathbf R)$，使得 $T q = p$，即 $5 q'' + 3 q' = p$。
\end{solve}

\begin{exercise}\label{prob:P-3B-30}
设 $\phi \in \mathcal L (V , \mathbf F)$，$u \in V$，满足 $\phi \ne 0$ 且 $u \notin \ker \phi$。证明：

\[V = \ker \phi \oplus \{a u : a \in \mathbf F\}\]
\end{exercise}
\noindent 来源：\href{https://github.com/szdytom/LADRSolutions/blob/main/sections/3B.typ}{P-3B-30}（原书 3B 第 30 题；方而静题解）。
\begin{solve}
设 $v \in V$，满足 $\phi v \ne 0$。令 $k = \frac{\phi v}{\phi u} \in \mathbf F$，则 $\phi (k u) = \phi v$，故 $\phi (v - k u) = 0$，即 $v - k u \in \ker \phi$。因此

\[v = (v - k u) + k u\]

其中 $v - k u \in \ker \phi$ 且 $k u \in \{a u : a \in \mathbf F\}$，故 $V = \ker \phi + \{a u : a \in \mathbf F\}$。现在说明这个和是直和：将 $0$ 表示为 $0 = v + w$，其中 $v \in \ker \phi$ 且 $w \in \{a u : a \in \mathbf F\}$，于是

\[\phi 0 = \phi (v + w) = \phi w\]

设 $w = a u$，则 $\phi w = a \phi u$，由于 $\phi u \ne 0$，故 $a = 0$，即 $w = 0$。进一步，由于 $0 = v + w$，故 $v = 0$。根据直和的条件（原书1.45），得 $V = \ker \phi \oplus \{a u : a \in \mathbf F\}$。
\end{solve}

\begin{exercise}\label{prob:P-3B-31}
设 $V$ 是有限维向量空间，$X$ 是 $V$ 的子空间，$Y$ 是 $W$ 的有限维子空间。证明：存在 $T \in \mathcal L (V , W)$ 使得 $\ker T = X$ 且 $\operatorname{im} T = Y$，当且仅当，$\dim X + \dim Y = \dim V$。
\end{exercise}
\noindent 来源：\href{https://github.com/szdytom/LADRSolutions/blob/main/sections/3B.typ}{P-3B-31}（原书 3B 第 31 题；方而静题解）。
\begin{solve}
首先，假设存在 $T \in \mathcal L (V , W)$ 使得 $\ker T = X$ 且 $\operatorname{im} T = Y$。根据线性映射基本定理（原书3.21），立即可得 $\dim X + \dim Y = \dim V$。

现在假设 $\dim X + \dim Y = \dim V$。令 $n = \dim V$。设 $w_{1} , \ldots , w_{m}$ 是 $Y$ 的一组基，$v_{m + 1} , \ldots , v_{n}$ 是 $X$ 的一组基，根据“每个线性无关组都可被扩充成基”（原书2.32），可以找到 $v_{1} , \ldots , v_{m} \in V$，使得 $v_{1} , \ldots , v_{n}$ 是 $V$ 的一组基。

根据线性映射引理（原书3.4），存在 $T \in \mathcal L (V , W)$，使得对于任意 $i \in \{1 , \ldots , m\}$，有 $T v_{i} = w_{i}$，且对于任意 $i \in \{m + 1 , \ldots , n\}$，有 $T v_{i} = 0$。

设 $v \in X$，则存在 $a_{m + 1} , \ldots , a_{n} \in \mathbf F$，使得 $v = a_{m + 1} v_{m + 1} + \cdots + a_{n} v_{n}$。因此

\[T v = T (\sum_{k = m + 1}^{n} a_{k} v_{k}) = \sum_{k = m + 1}^{n} a_{k} T v_{k} = 0\]

这说明 $X \subseteq \ker T$。另一方面，设 $v \in \ker T$，将 $v$ 表示为 $v = a_{1} v_{1} + \cdots + a_{n} v_{n}$，其中 $a_{1} , \ldots , a_{n} \in \mathbf F$。由于 $T v = 0$，因此

\[0 = T v = T (\sum_{k = 1}^{n} a_{k} v_{k}) = \sum_{k = 1}^{n} a_{k} T v_{k} = \sum_{k = 1}^{m} a_{k} w_{k}\]

由于 $w_{1} , \ldots , w_{m}$ 是线性无关的，故 $a_{1} = \cdots = a_{m} = 0$，即 $v \in X$。因此 $\ker T \subseteq X$，故 $\ker T = X$。

设 $w \in Y$，则存在 $b_{1} , \ldots , b_{m} \in \mathbf F$，使得 $w = b_{1} w_{1} + \cdots + b_{m} w_{m}$。注意到

\[T (\sum_{k = 1}^{m} b_{k} v_{k}) = \sum_{k = 1}^{m} b_{k} T v_{k} = \sum_{k = 1}^{m} b_{k} w_{k} = w\]

这说明 $Y \subseteq \operatorname{im} T$。另一方面，设 $w \in \operatorname{im} T$，则存在 $v \in V$，使得 $T v = w$。将 $v$ 表示为 $v = a_{1} v_{1} + \cdots + a_{n} v_{n}$，其中 $a_{1} , \ldots , a_{n} \in \mathbf F$。则

\[w = T v = T (\sum_{k = 1}^{n} a_{k} v_{k}) = \sum_{k = 1}^{m} a_{k} T v_{k} + \sum_{k = m + 1}^{n} a_{k} T v_{k} = \sum_{k = 1}^{m} a_{k} w_{k}\]

这说明 $w \in \operatorname{span} (w_{1} , \ldots , w_{m}) = Y$。因此 $\operatorname{im} T \subseteq Y$，故 $\operatorname{im} T = Y$。

综上所述，存在 $T \in \mathcal L (V , W)$ 使得 $\ker T = X$ 且 $\operatorname{im} T = Y$，当且仅当，$\dim X + \dim Y = \dim V$。
\end{solve}

\begin{exercise}\label{prob:P-3B-32}
设 $V$ 是有限维向量空间，$\dim V > 1$。证明：若 $\phi : \mathcal L (V) \to \mathbf F$ 是线性映射，使得对于任意 $S , T \in \mathcal L (V)$，$\phi (S T) = \phi (S) \phi (T)$，则 $\phi = 0$。

\begin{remark}
习题~\ref{app:E-two-sided-ideal} 中给出了关于 $\mathcal L (V)$ 的双边理想的描述，或许有用。
\end{remark}
\end{exercise}
\noindent 来源：\href{https://github.com/szdytom/LADRSolutions/blob/main/sections/3B.typ}{P-3B-32}（原书 3B 第 32 题；方而静题解）。
\begin{solve}
核 $\ker\phi$ 是 $\mathcal L(V)$ 的线性子空间。若 $A\in\ker\phi$、$B\in\mathcal L(V)$，则
\[\phi(AB)=\phi(A)\phi(B)=0,\qquad\phi(BA)=\phi(B)\phi(A)=0,\]
所以 $\ker\phi$ 是双边理想。设 $n=\dim V>1$，则 $\dim\mathcal L(V)=n^2>1=\dim\mathbf F$，线性映射 $\phi$ 不可能单射，故 $\ker\phi\ne\{0\}$。由习题~\ref{app:E-two-sided-ideal}，只能有 $\ker\phi=\mathcal L(V)$，即 $\phi=0$。
\end{solve}

\begin{exercise}\label{prob:P-3B-33}
设 $V$ 和 $W$ 都是实向量空间，$T \in \mathcal L (V , W)$。定义 $T_{\mathbf C} : V_{\mathbf C} \to W_{\mathbf C}$ 为对于任意 $u , v \in V$，

\[T_{\mathbf C} (u + \mathrm i v) = T u + \mathrm i T v\]

\begin{enumerate}[label=(\alph*)]
\item 证明：$T_{\mathbf C}$ 是 $V_{\mathbf C} \to W_{\mathbf C}$ 的（复）线性映射；
\item 证明：$T_{\mathbf C}$ 是单射，当且仅当 $T$ 是单射；
\item 证明：$\operatorname{im} T_{\mathbf C} = W_{\mathbf C}$，当且仅当 $\operatorname{im} T = W$。

\end{enumerate}
\begin{remark}
复化 $V_{\mathbf C}$ 定义于习题~\ref{app:E-vector-dspace-complexification}，线性映射 $T_{\mathbf C}$ 被称为线性映射 $T$ 的\emph{复化（complexification）}。
\end{remark}
\end{exercise}
\noindent 来源：\href{https://github.com/szdytom/LADRSolutions/blob/main/sections/3B.typ}{P-3B-33}（原书 3B 第 33 题；方而静题解）。
\begin{solve}
对于（a），我们逐条验证线性映射的定义（原书3.1）中给出的要求：

\begin{description}
\item[可加性] 对于任意 $u , v \in V_{\mathbf C}$，均有 $T_{\mathbf C} (u + v) = T_{\mathbf C} u + T_{\mathbf C} v$。
证明：设 $u = u_{1} + \mathrm i u_{2}$，$v = v_{1} + \mathrm i v_{2}$，其中 $u_{1} , u_{2} , v_{1} , v_{2} \in V$。则
\[\begin{aligned}T_{\mathbf C} (u + v) & = T_{\mathbf C} ((u_{1} + v_{1}) + \mathrm i (u_{2} + v_{2})) \\ & = T (u_{1} + v_{1}) + \mathrm i T (u_{2} + v_{2}) \\ & = T u_{1} + \mathrm i T u_{2} + T v_{1} + \mathrm i T v_{2} \\ & = T_{\mathbf C} u + T_{\mathbf C} v\end{aligned}\]

\item[齐次性] 对于任意 $\lambda \in \mathbf C$，$u \in V_{\mathbf C}$，均有 $T_{\mathbf C} (\lambda u) = \lambda T_{\mathbf C} u$。
证明：设 $u = u_{1} + \mathrm i u_{2}$，其中 $u_{1} , u_{2} \in V$。则
\[\begin{aligned}T_{\mathbf C} (\lambda u) & = T_{\mathbf C} (\lambda (u_{1} + \mathrm i u_{2})) \\ & = T (\lambda u_{1}) + \mathrm i T (\lambda u_{2}) \\ & = \lambda T u_{1} + \mathrm i \lambda T u_{2} \\ & = \lambda (T u_{1} + \mathrm i T u_{2}) \\ & = \lambda T_{\mathbf C} u\end{aligned}\]

这说明 $T_{\mathbf C}$ 确实是 $V_{\mathbf C} \to W_{\mathbf C}$ 的线性映射。

对于（b），首先假设 $T_{\mathbf C}$ 是单射。根据“单射性 $\iff$ 零空间为 $\{0\}$”（原书3.15），可得 $\ker T_{\mathbf C} = \{0\}$。设 $v \in \ker T$，则根据“线性映射将 $0$ 映射为 $0$”（原书3.10），可得 $0 = T v = T v + \mathrm i T 0 = T_{\mathbf C} (v + \mathrm i 0)$，因此 $v = 0$，即 $\ker T = \{0\}$，故 $T$ 是单射。

另一方面，假设 $T$ 是单射。设 $u + \mathrm i v \in \ker T_{\mathbf C}$，则 $0 = T_{\mathbf C} (u + \mathrm i v) = T u + \mathrm i T v$，故 $T u = T v = 0$。又因为 $\ker T = \{0\}$，只能有 $u = v = 0$，即 $\ker T_{\mathbf C} = \{0\}$，因此 $T_{\mathbf C}$ 是单射。

对于（c），首先假设 $\operatorname{im} T_{\mathbf C} = W_{\mathbf C}$。设 $w \in W$，则存在 $u + \mathrm i v \in V_{\mathbf C}$，使得 $T (u + \mathrm i v) = w + \mathrm i 0$，即 $T u + \mathrm i T v = w + \mathrm i 0$，故 $T u = w$。于是 $W \subseteq \operatorname{im} T$，即 $\operatorname{im} T = W$。

另一方面，假设 $\operatorname{im} T = W$。设 $w + \mathrm i v \in W_{\mathbf C}$，则存在 $u \in V$，使得 $T u = w$。因此 $T_{\mathbf C} (u + \mathrm i v) = T u + \mathrm i T v = w + \mathrm i T v$，即 $w + \mathrm i v \in \operatorname{im} T_{\mathbf C}$。这说明 $W_{\mathbf C} \subseteq \operatorname{im} T_{\mathbf C}$，即 $\operatorname{im} T_{\mathbf C} = W_{\mathbf C}$。
\end{description}
\end{solve}

\subsection{矩阵表示}\label{app:section-3C}
\begin{exercise}\label{prob:P-3C-1}
设 $V,W$ 是有限维向量空间，$T\in\mathcal L(V,W)$。证明：在任意一对基下，$T$ 的矩阵至少有 $\dim\operatorname{im} T$ 个非零元素。
\end{exercise}
\noindent 来源：\href{https://github.com/szdytom/LADRSolutions/blob/main/sections/3C.typ}{P-3C-1}（原书 3C 第 1 题；方而静题解）。
\begin{solve}
设 $v_{1} , \ldots , v_{n}$ 是 $V$ 的一组基，$w_{1} , \ldots , w_{m}$ 是 $W$ 的一组基，$A$ 为 $T$ 在这两组基下的矩阵表示。则对于 $i \in \{1 , \ldots , m\}$ 和 $j \in \{1 , \ldots , n\}$，

\[T v_{j} = \sum_{k = 1}^{m} A_{k , j} w_{k}\]

设 $p = \dim \ker T$，且 $q = \dim \operatorname{im} T$，则根据线性映射基本定理（原书3.21），$p + q = n$。则在 $v_{1} , \ldots , v_{n}$ 中至多有 $p$ 个向量属于 $\ker T$，也即至少有 $q$ 个向量不属于 $\ker T$。不妨设 $v_{1} , \ldots , v_{q}$ 不属于 $\ker T$，则对于 $j \in \{1 , \ldots , q\}$，我们有

\[A_{1 , j} w_{1} + \cdots + A_{m , j} w_{m} = T v_{j} \ne 0\]

考虑到 $w_{1} , \ldots , w_{m}$ 线性无关，只能有至少一个 $k \in \{1 , \ldots , m\}$ 使得 $A_{k , j} \ne 0$。因此，矩阵 $A$ 至少有 $q$ 个非零元素。
\end{solve}

\begin{exercise}\label{prob:P-3C-2}
设 $T \in \mathcal L (V , W)$，其中 $V$ 和 $W$ 都是有限维的非零向量空间。证明：$\dim \operatorname{im} T = 1$，当且仅当，存在 $V$ 的一组基和 $W$ 的一组基，使得 $\mathcal M (T)$ 的所有元素都是 $1$。
\end{exercise}
\noindent 来源：\href{https://github.com/szdytom/LADRSolutions/blob/main/sections/3C.typ}{P-3C-2}（原书 3C 第 2 题；方而静题解）。
\begin{solve}
若某对基下的矩阵所有元素均为 $1$，由于 $V,W$ 非零，这个矩阵的列均相同且非零，故秩为 $1$。
反过来，设 $\dim\operatorname{im} T=1$。取 $u\in V$ 使 $Tu=w\ne0$，再取 $\ker T$ 的基 $k_2,\ldots,k_n$。由秩与零度定理，$u,k_2,\ldots,k_n$ 是 $V$ 的一组基；令 $v_1=u$，$v_j=u+k_j$（$j\ge2$）。比较各个 $k_j$ 的系数可知 $v_1,\ldots,v_n$ 仍是一组基，而且 $Tv_j=w$ 对所有 $j$ 成立。
将 $w$ 扩充为 $W$ 的基 $w,z_2,\ldots,z_m$。令 $w_1=w-\sum_{j=2}^m z_j$、$w_j=z_j$（$j\ge2$），则它们也是一组基，且 $\sum_jw_j=w$。因此每个 $Tv_j$ 在这组基下的坐标都是 $(1,\ldots,1)^{\mathrm T}$，所求矩阵的全部元素均为 $1$。维数为 $1$ 时，上述空和按零处理。
\end{solve}

\begin{exercise}\label{prob:P-3C-3}
设 $v_{1} , \ldots , v_{n}$ 是 $V$ 的一组基，$w_{1} , \ldots , w_{m}$ 是 $W$ 的一组基。证明：

\begin{enumerate}[label=(\alph*)]
\item 若 $S , T \in \mathcal L (V , W)$，则 $M (S + T) = M (S) + M (T)$；
\item 若 $\lambda \in \mathbf F$，$T \in \mathcal L (V , W)$，则 $M (\lambda T) = \lambda M (T)$。

\end{enumerate}
\begin{remark}
本题是在让你验证原书3.35和3.38。
\end{remark}
\end{exercise}
\noindent 来源：\href{https://github.com/szdytom/LADRSolutions/blob/main/sections/3C.typ}{P-3C-3}（原书 3C 第 3 题；方而静题解）。
\begin{solve}
对于（a），记 $A = M (S)$，$B = M (T)$，$C = M (S + T)$。则对于 $j \in \{1 , \ldots , n\}$，有

\[\sum_{k = 1}^{m} C_{k , j} w_{k} = (S + T) v_{j} = S v_{j} + T v_{j} = \sum_{k = 1}^{m} (A_{k , j} + B_{k , j}) w_{k}\]

故 $C_{k , j} = A_{k , j} + B_{k , j}$，即 $C = A + B$。

对于（b），记 $A = M (T)$，$B = M (\lambda T)$。则对于 $j \in \{1 , \ldots , n\}$，有

\[\sum_{k = 1}^{m} B_{k , j} w_{k} = (\lambda T) v_{j} = \lambda (T v_{j}) = \sum_{k = 1}^{m} (\lambda A_{k , j}) w_{k}\]

故 $B_{k , j} = \lambda A_{k , j}$，即 $B = \lambda A$。
\end{solve}

\begin{exercise}\label{prob:P-3C-4}\label{app:E-mat-of-derivation-of-poly3}
设 $D \in \mathcal L (\mathcal P_{3} (\mathbf R) , \mathcal P_{2} (\mathbf R))$ 是微分映射，定义为 $p \mapsto p'$。求 $\mathcal P_{3} (\mathbf R)$ 的一个基和 $\mathcal P_{2} (\mathbf R)$ 的一个基，使得 $\mathcal M (D)$ 为

\[\begin{pmatrix}1&0&0&0\\0&1&0&0\\0&0&1&0\end{pmatrix}\]

\begin{remark}
和原书3.33比较一下。下一题拓展了本题。
\end{remark}
\end{exercise}
\noindent 来源：\href{https://github.com/szdytom/LADRSolutions/blob/main/sections/3C.typ}{P-3C-4}（原书 3C 第 4 题；方而静题解）。
\begin{solve}
定义域取有序基 $(x^3,x^2,x,1)$，陪域取有序基 $(3x^2,2x,1)$。两组多项式的次数互异且数目等于各自空间维数，故确为基。由
\[D(x^3)=3x^2,\quad D(x^2)=2x,\quad D(x)=1,\quad D(1)=0,\]
四个像的坐标依次为前三个标准基向量及零向量，所得矩阵正是题给矩阵。
\end{solve}

\begin{exercise}\label{prob:P-3C-5}
设 $V$ 和 $W$ 都是有限维向量空间，$T \in \mathcal L (V , W)$。证明：存在 $V$ 的一组基和 $W$ 的一组基，使得关于这些基的 $\mathcal M (T)$ 的第 $k$ 行第 $k$ 列元素为 $1$，其余元素为 $0$，其中 $k \in \{1 , \ldots , \dim \operatorname{im} T\}$。
\end{exercise}
\noindent 来源：\href{https://github.com/szdytom/LADRSolutions/blob/main/sections/3C.typ}{P-3C-5}（原书 3C 第 5 题；方而静题解）。
\begin{solve}
记 $r=\dim\operatorname{im} T$。先取 $\ker T$ 的基 $k_1,\ldots,k_s$，扩充为 $V$ 的基 $k_1,\ldots,k_s,u_1,\ldots,u_r$。向量 $Tu_1,\ldots,Tu_r$ 张成 $\operatorname{im} T$；若 $\sum_ja_jTu_j=0$，则 $\sum_ja_ju_j\in\ker T$，由扩充后基的线性无关性知所有 $a_j=0$，所以它们是像空间的一组基。
将其扩充为 $W$ 的基 $(Tu_1,\ldots,Tu_r,w_{r+1},\ldots,w_m)$，并将定义域基按 $(u_1,\ldots,u_r,k_1,\ldots,k_s)$ 排列。此时矩阵为
\[\begin{pmatrix}I_r&O\\O&O\end{pmatrix},\]
即仅前 $r$ 个对角元为 $1$。当 $r=0$ 时为零矩阵，零维空间的基按空组理解。
\end{solve}

\begin{exercise}\label{prob:P-3C-6}
设 $v_1,\ldots,v_m$ 是 $V$ 的一组基，$m\ge1$，$W$ 是有限维非零向量空间，$T\in\mathcal L(V,W)$。证明：存在 $W$ 的一组基，使得 $T$ 在给定定义域基及该陪域基下的矩阵第一列至多有一个非零元素，且若存在，则它是第一行的 $1$。
\end{exercise}
\noindent 来源：\href{https://github.com/szdytom/LADRSolutions/blob/main/sections/3C.typ}{P-3C-6}（原书 3C 第 6 题；方而静题解）。
\begin{solve}
分类讨论，当 $T v_{1} = 0$ 时，任取 $W$ 的一组基为 $w_{1} , \ldots , w_{n}$，则

\[T v_{1} = 0 = 0 \cdot w_{1} + \cdots + 0 \cdot w_{n}\]

于是 $\mathcal M (T)$ 的第一列的元素均为 $0$。

当 $T v_{1} \ne 0$ 时，取 $w_{1} = T v_{1}$。根据“每个线性无关组都可被扩充成基”（原书2.32），可以将 $w_{1}$ 扩充为 $W$ 的一组基 $w_{1} , \ldots , w_{n}$。则

\[T v_{1} = w_{1} = 1 \cdot w_{1} + 0 \cdot w_{2} + \cdots + 0 \cdot w_{n}\]

于是 $\mathcal M (T)$ 的第一行第一列的元素为 $1$，且第一列的其余元素均为 $0$。

综上所述，存在 $W$ 的一组基 $w_{1} , \ldots , w_{n}$，使得关于基 $v_{1} , \ldots , v_{m}$ 和 $w_{1} , \ldots , w_{n}$ 的 $\mathcal M (T)$ 的第一行第一列的元素为 $1$ 或 $0$，且第一列的其余元素均为 $0$。
\end{solve}


\section*{本次整理与修订说明}
\addcontentsline{toc}{section}{本次整理与修订说明}
2026 年 10 月 10 日的整理将正文 493 道题及附编 154 道题合入单一源文件，保留题目来源标记、原作者署名与许可，取消外部 PDF 插页和资源目录依赖。原目录完整备份留存于项目临时目录；备份不参与本文件的使用。

主要数学修订包括：区分伴随矩阵、对偶空间与伴随算子的记号；规范等价类和商集的表示；补足空间公理判别、极大线性无关组、有限子空间覆盖、仿射子集交集等论证；补全基础解系扩充与幂等映射分解；为正交矩阵结论补入实数域和阶数条件；修复多项式整除、插值次数界、代数重数及幂零矩阵的迹判据；修正隐函数题的陪域。附编进一步修复复化空间的数乘证明、奇偶分解、秩一映射的基构造及多处索引和零维边界问题。具体修订题目标记记于源码注释，既有修订仍保留在正文中。

本次交付只涉及源文件，未编译或更新 PDF。源码检查及针对性数学核验不能替代对全部题解的独立审定。
\end{document}
